% Options for packages loaded elsewhere
\PassOptionsToPackage{unicode}{hyperref}
\PassOptionsToPackage{hyphens}{url}
\PassOptionsToPackage{dvipsnames,svgnames,x11names}{xcolor}
\documentclass[
  spanish,
  10pt,
]{report}
\usepackage{xcolor}
\usepackage[a4paper,margin=20mm,headsep=7mm]{geometry}
\usepackage{amsmath,amssymb}
\setcounter{secnumdepth}{5}
\usepackage{iftex}
\ifPDFTeX
  \usepackage[T1]{fontenc}
  \usepackage[utf8]{inputenc}
  \usepackage{textcomp} % provide euro and other symbols
\else % if luatex or xetex
  \usepackage{unicode-math} % this also loads fontspec
  \defaultfontfeatures{Scale=MatchLowercase}
  \defaultfontfeatures[\rmfamily]{Ligatures=TeX,Scale=1}
\fi
\usepackage{lmodern}
\ifPDFTeX\else
  % xetex/luatex font selection
  \setmainfont[]{DejaVu Serif}
  \setsansfont[]{DejaVu Sans}
  \setmonofont[]{DejaVu Sans Mono}
\fi
% Use upquote if available, for straight quotes in verbatim environments
\IfFileExists{upquote.sty}{\usepackage{upquote}}{}
\IfFileExists{microtype.sty}{% use microtype if available
  \usepackage[]{microtype}
  \UseMicrotypeSet[protrusion]{basicmath} % disable protrusion for tt fonts
}{}
\usepackage{setspace}
\makeatletter
\@ifundefined{KOMAClassName}{% if non-KOMA class
  \IfFileExists{parskip.sty}{%
    \usepackage{parskip}
  }{% else
    \setlength{\parindent}{0pt}
    \setlength{\parskip}{6pt plus 2pt minus 1pt}}
}{% if KOMA class
  \KOMAoptions{parskip=half}}
\makeatother
\usepackage{color}
\usepackage{fancyvrb}
\newcommand{\VerbBar}{|}
\newcommand{\VERB}{\Verb[commandchars=\\\{\}]}
\DefineVerbatimEnvironment{Highlighting}{Verbatim}{commandchars=\\\{\}}
% Add ',fontsize=\small' for more characters per line
\newenvironment{Shaded}{}{}
\newcommand{\AlertTok}[1]{\textcolor[rgb]{1.00,0.00,0.00}{\textbf{#1}}}
\newcommand{\AnnotationTok}[1]{\textcolor[rgb]{0.38,0.63,0.69}{\textbf{\textit{#1}}}}
\newcommand{\AttributeTok}[1]{\textcolor[rgb]{0.49,0.56,0.16}{#1}}
\newcommand{\BaseNTok}[1]{\textcolor[rgb]{0.25,0.63,0.44}{#1}}
\newcommand{\BuiltInTok}[1]{\textcolor[rgb]{0.00,0.50,0.00}{#1}}
\newcommand{\CharTok}[1]{\textcolor[rgb]{0.25,0.44,0.63}{#1}}
\newcommand{\CommentTok}[1]{\textcolor[rgb]{0.38,0.63,0.69}{\textit{#1}}}
\newcommand{\CommentVarTok}[1]{\textcolor[rgb]{0.38,0.63,0.69}{\textbf{\textit{#1}}}}
\newcommand{\ConstantTok}[1]{\textcolor[rgb]{0.53,0.00,0.00}{#1}}
\newcommand{\ControlFlowTok}[1]{\textcolor[rgb]{0.00,0.44,0.13}{\textbf{#1}}}
\newcommand{\DataTypeTok}[1]{\textcolor[rgb]{0.56,0.13,0.00}{#1}}
\newcommand{\DecValTok}[1]{\textcolor[rgb]{0.25,0.63,0.44}{#1}}
\newcommand{\DocumentationTok}[1]{\textcolor[rgb]{0.73,0.13,0.13}{\textit{#1}}}
\newcommand{\ErrorTok}[1]{\textcolor[rgb]{1.00,0.00,0.00}{\textbf{#1}}}
\newcommand{\ExtensionTok}[1]{#1}
\newcommand{\FloatTok}[1]{\textcolor[rgb]{0.25,0.63,0.44}{#1}}
\newcommand{\FunctionTok}[1]{\textcolor[rgb]{0.02,0.16,0.49}{#1}}
\newcommand{\ImportTok}[1]{\textcolor[rgb]{0.00,0.50,0.00}{\textbf{#1}}}
\newcommand{\InformationTok}[1]{\textcolor[rgb]{0.38,0.63,0.69}{\textbf{\textit{#1}}}}
\newcommand{\KeywordTok}[1]{\textcolor[rgb]{0.00,0.44,0.13}{\textbf{#1}}}
\newcommand{\NormalTok}[1]{#1}
\newcommand{\OperatorTok}[1]{\textcolor[rgb]{0.40,0.40,0.40}{#1}}
\newcommand{\OtherTok}[1]{\textcolor[rgb]{0.00,0.44,0.13}{#1}}
\newcommand{\PreprocessorTok}[1]{\textcolor[rgb]{0.74,0.48,0.00}{#1}}
\newcommand{\RegionMarkerTok}[1]{#1}
\newcommand{\SpecialCharTok}[1]{\textcolor[rgb]{0.25,0.44,0.63}{#1}}
\newcommand{\SpecialStringTok}[1]{\textcolor[rgb]{0.73,0.40,0.53}{#1}}
\newcommand{\StringTok}[1]{\textcolor[rgb]{0.25,0.44,0.63}{#1}}
\newcommand{\VariableTok}[1]{\textcolor[rgb]{0.10,0.09,0.49}{#1}}
\newcommand{\VerbatimStringTok}[1]{\textcolor[rgb]{0.25,0.44,0.63}{#1}}
\newcommand{\WarningTok}[1]{\textcolor[rgb]{0.38,0.63,0.69}{\textbf{\textit{#1}}}}
\usepackage{longtable,booktabs,array}
\newcounter{none} % for unnumbered tables
\usepackage{calc} % for calculating minipage widths
% Correct order of tables after \paragraph or \subparagraph
\usepackage{etoolbox}
\makeatletter
\patchcmd\longtable{\par}{\if@noskipsec\mbox{}\fi\par}{}{}
\makeatother
% Allow footnotes in longtable head/foot
\IfFileExists{footnotehyper.sty}{\usepackage{footnotehyper}}{\usepackage{footnote}}
\makesavenoteenv{longtable}
\ifLuaTeX
\usepackage[bidi=basic,shorthands=off]{babel}
\else
\usepackage[bidi=default,shorthands=off]{babel}
\fi
\ifPDFTeX
\else
\babelfont{rm}[]{DejaVu Serif}
\fi
\ifLuaTeX
  \usepackage{selnolig} % disable illegal ligatures
\fi
\setlength{\emergencystretch}{3em} % prevent overfull lines
\providecommand{\tightlist}{%
  \setlength{\itemsep}{0pt}\setlength{\parskip}{0pt}}
\usepackage{amsmath,amssymb}
\usepackage{tikz}
\usetikzlibrary{positioning,arrows.meta}
\usepackage{xcolor}
\definecolor{brand}{HTML}{1F5A43}
\definecolor{surface}{HTML}{F6F7F4}
\usepackage{fancyhdr}
\usepackage{fvextra}
\usepackage{seqsplit}
\usepackage{pdfpages}
\usepackage{xurl}
\usepackage{titlesec}
\titleformat{\chapter}[display]{\sffamily\bfseries\color{brand}\Huge}{\chaptertitlename\ \thechapter}{12pt}{\huge}
\titleformat{\section}{\sffamily\bfseries\color{brand}\Large}{\thesection}{.7em}{}
\titleformat{\subsection}{\sffamily\bfseries\large}{\thesubsection}{.7em}{}
\pagestyle{fancy}
\fancyhf{}
\fancyhead[L]{\sffamily\small BoviCierre}
\fancyhead[R]{\sffamily\small Manual técnico integral}
\fancyfoot[L]{\sffamily\scriptsize Revisión 92d3d41 · 12 octubre 2026}
\fancyfoot[R]{\sffamily\small\thepage}
\setlength{\headheight}{15pt}
\setlength{\emergencystretch}{3em}
\setlength{\parskip}{6pt}
\setlength{\parindent}{0pt}
\setcounter{tocdepth}{1}
\setcounter{secnumdepth}{2}
\AtBeginDocument{\hypersetup{colorlinks=true,linkcolor=brand,urlcolor=brand}}
\AtBeginDocument{\RecustomVerbatimEnvironment{Highlighting}{Verbatim}{breaklines,breakanywhere,fontsize=\scriptsize,commandchars=\\\{\}}}
\widowpenalty=10000
\clubpenalty=10000
\fancyhead[L]{\color{brand}\sffamily\small BoviCierre}
\fancyhead[R]{\color{black!70}\sffamily\small Manual técnico integral}
\fancyfoot[L]{\color{black!70}\sffamily\footnotesize Revisión 92d3d41 · 12 octubre 2026}
\fancyfoot[R]{\color{brand}\sffamily\small\thepage}
\renewcommand{\maketitle}{%
\hypersetup{pageanchor=false}
\begin{titlepage}
\centering
\vspace*{28mm}
{\sffamily\bfseries\color{brand}\fontsize{38}{45}\selectfont BoviCierre\par}
\vspace{12mm}
{\sffamily\fontsize{21}{28}\selectfont Manual integral del software\par}
\vspace{9mm}
{\large Arquitectura, código, datos y operación local\par}
\vspace{20mm}
{\normalsize Una guía para comprender el negocio, seguir las peticiones\par
 y mantener el sistema desde sus fuentes reales\par}
\vfill
{\sffamily\normalsize Preparado para Aldo y Jared\par}
\vspace{5mm}
{\normalsize Revisión del repositorio 92d3d41\par
Actualización de campaña UI/UX y QA (V14--V18)\par
12 de octubre de 2026\par}
\vspace{13mm}
{\small PostgreSQL · Java · React · Docker Compose\par}
\end{titlepage}
\pagenumbering{roman}
\hypersetup{pageanchor=true}
}

\AtBeginDocument{\special{pdf:minorversion 7}}
\DeclareRobustCommand{\softcode}[1]{\texttt{\seqsplit{#1}}}
\DeclareRobustCommand{\softtext}[1]{\seqsplit{#1}}

\sloppy
\makeatletter
\renewcommand*\l@section{\@dottedtocline{1}{1.5em}{3.5em}}
\makeatother
\makeatletter
\renewcommand{\@pnumwidth}{2.5em}
\renewcommand{\@tocrmarg}{3.5em}
\makeatother
\makeatletter
\patchcmd{\l@chapter}{\setlength\@tempdima{1.5em}}{\setlength\@tempdima{2.5em}}{}{}
\makeatother
\usepackage{bookmark}
\IfFileExists{xurl.sty}{\usepackage{xurl}}{} % add URL line breaks if available
\urlstyle{same}
\hypersetup{
  pdftitle={BoviCierre},
  pdfauthor={Documentación técnica preparada para Aldo y Jared},
  pdflang={es},
  colorlinks=true,
  linkcolor={Maroon},
  filecolor={Maroon},
  citecolor={Blue},
  urlcolor={Blue},
  pdfcreator={LaTeX via pandoc}}

\title{BoviCierre}
\usepackage{etoolbox}
\makeatletter
\providecommand{\subtitle}[1]{% add subtitle to \maketitle
  \apptocmd{\@title}{\par {\large #1 \par}}{}{}
}
\makeatother
\subtitle{Manual integral del software y su funcionamiento}
\author{Documentación técnica preparada para Aldo y Jared}
\date{12 de octubre de 2026}

\begin{document}
\maketitle

\renewcommand*\contentsname{Índice de contenido}
{
\setcounter{tocdepth}{1}
\tableofcontents
}
\setstretch{1.08}
\clearpage
\pagenumbering{arabic}

\chapter{Cómo leer y entender
BoviCierre}\label{cuxf3mo-leer-y-entender-bovicierre}

\section{Propósito y revisión de este
manual}\label{propuxf3sito-y-revisiuxf3n-de-este-manual}

Este manual explica BoviCierre desde el problema del negocio hasta sus
archivos \texorpdfstring{\softtext{ejecutables.}}{ejecutables.} Está
dirigido a Aldo, Jared y a una persona que necesite aprender a usar,
mantener y explicar el sistema en una exposición
\texorpdfstring{\softtext{académica.}}{académica.} No exige conocer
Spring, React o PostgreSQL de antemano: introduce sus responsabilidades
antes de entrar a las firmas de métodos y estructuras de datos. Una
lectura completa permite seguir una petición desde el clic hasta el
registro que queda en
\texorpdfstring{\softtext{PostgreSQL,}}{PostgreSQL,} identificar quién
puede modificarlo y entender por qué una cifra cambia o permanece
\texorpdfstring{\softtext{histórica.}}{histórica.}

La revisión documentada es
\texorpdfstring{\softcode{92d3d4103a174eb11f4ccf21eff2deed8cc77682}}{92d3d4103a174eb11f4ccf21eff2deed8cc77682},
publicada en la rama main de
\texorpdfstring{\softcode{https://github.com/ALDOOGOAT/bovicierre}}{https://github.com/ALDOOGOAT/bovicierre}.
La fecha de elaboración es el 12 de octubre de 2026. Los resultados de
implementación y pruebas reproducidos corresponden a sus ejecuciones
registradas del 11 de octubre; elaborar este PDF no significa haber
repetido todas las suites. Las fuentes principales son el código
versionado, las migraciones, OpenAPI y la evidencia guardada. Los
archivos del futuro piloto y los datos comerciales reales no existen en
esta revisión.

El manual distingue tres planos. El \textbf{negocio} describe qué
significado tiene comprar, entregar o pagar. La \textbf{implementación}
identifica archivos y métodos que realizan esas acciones. La
\textbf{evidencia} indica qué escenarios se ejecutaron, con qué
resultados y qué límites conservan. Tener una documentación o un
endpoint descrito no equivale a probar todas sus
\texorpdfstring{\softtext{combinaciones.}}{combinaciones.} Un escenario
DEMO es un dato sintético persistido en el backend; no es información de
un productor real.

\section{Ruta recomendada de lectura}\label{ruta-recomendada-de-lectura}

Si primero necesitas arrancarlo, lee la guía de ejecución local y luego
el recorrido de compra y venta. Si deseas entender el código, estudia la
visión de arquitectura, la comunicación HTTP, el backend y el frontend
en ese orden. Si debes explicar la asignatura de bases de datos,
comienza con identidad y ciclo, continúa con modelos, funciones SQL,
roles, transacciones, replicación y
\texorpdfstring{\softtext{restauración.}}{restauración.} Para revisar
una fórmula, consulta su explicación y su caso GOLD: el caso contiene
esperados fijos que sirven como oráculo.

Los anexos son material de consulta, no una segunda
\texorpdfstring{\softtext{aplicación.}}{aplicación.} El catálogo de API
permite ubicar una ruta concreta; el diccionario reúne campos y
relaciones; las reglas numeradas conservan los identificadores del
\texorpdfstring{\softtext{repositorio.}}{repositorio.} Las referencias
de la forma \texorpdfstring{\softcode{ruta:línea}}{ruta:línea} apuntan
al corte \texorpdfstring{\softtext{versionado.}}{versionado.} Una línea
puede cambiar en commits posteriores, por lo que debe contrastarse
siempre con el nombre del método y el archivo. Los nombres técnicos
aparecen en tipografía monoespaciada, por ejemplo
\texorpdfstring{\softcode{posicion\_inventario}}{posicion\_inventario}.

\section{Qué problema resuelve}\label{quuxe9-problema-resuelve}

Una operación ganadera tiene varias dimensiones que rara vez ocurren a
la vez. Se puede pactar una compra hoy, recibir los animales mañana,
pagar parte después y recibir una factura de flete la semana siguiente.
También puede venderse una parte del lote antes de conocer todos sus
costos. Si se utiliza una única etiqueta como ``completado'', se pierde
la posibilidad de responder cuánto dinero falta, quién tiene físicamente
los animales y cuánto dejó cada ciclo.

BoviCierre conserva esas dimensiones separadas. La compra o venta tiene
condiciones y una versión. La posición representa la adquisición
concreta de un animal. La recepción y entrega son movimientos físicos.
La liquidación representa una
\texorpdfstring{\softtext{obligación.}}{obligación.} El pago representa
dinero validado; su aplicación reduce una
\texorpdfstring{\softtext{obligación.}}{obligación.} El gasto puede
aumentar el costo de inventario, el costo de venta o la pérdida, según
el estado del ciclo. La evidencia conserva documentos, pesajes y
revisiones que explican el resultado.

El cliente objetivo es un comercializador o acopiador que compra y vende
bovinos por cuenta propia, con una o dos sedes. El dueño o administrador
compra el software; operaciones, caja y supervisión lo utilizan. El
producto no es un marketplace, no custodia dinero y no emite una
autorización sanitaria. Registrar un identificador o revisar un archivo
administrativamente no certifica su autenticidad oficial.

\section{Conceptos que no debes
confundir}\label{conceptos-que-no-debes-confundir}

{\def\LTcaptype{none} % do not increment counter
\begin{longtable}[]{@{}
  >{\raggedright\arraybackslash}p{(\linewidth - 4\tabcolsep) * \real{0.3333}}
  >{\raggedright\arraybackslash}p{(\linewidth - 4\tabcolsep) * \real{0.3333}}
  >{\raggedright\arraybackslash}p{(\linewidth - 4\tabcolsep) * \real{0.3333}}@{}}
\toprule\noalign{}
\begin{minipage}[b]{\linewidth}\raggedright
Concepto
\end{minipage} & \begin{minipage}[b]{\linewidth}\raggedright
Significado
\end{minipage} & \begin{minipage}[b]{\linewidth}\raggedright
Ejemplo
\end{minipage} \\
\midrule\noalign{}
\endhead
\bottomrule\noalign{}
\endlastfoot
Animal & Identidad \texorpdfstring{\softtext{permanente.}}{permanente.}
& El mismo animal puede aparecer en una recompra. \\
Posición o ciclo & Una adquisición con costo y cierre propios. & Primera
compra de 1100 y segunda compra de 1300. \\
Operación & Compra o venta con selección y
\texorpdfstring{\softtext{condiciones.}}{condiciones.} & Una venta de
cuatro \texorpdfstring{\softtext{posiciones.}}{posiciones.} \\
Liquidación & Obligación y conceptos que explican su importe. & El
comprador debe 6400 antes de
\texorpdfstring{\softtext{aplicaciones.}}{aplicaciones.} \\
Pago o cobro & Movimiento monetario registrado y validado. & Cobro de
1200 que puede tener 200
\texorpdfstring{\softtext{disponibles.}}{disponibles.} \\
Aplicación & Vínculo de dinero validado a una
\texorpdfstring{\softtext{obligación.}}{obligación.} & Se aplican 1000 a
una deuda de 1000. \\
Entrega & Hecho físico distinto del reconocimiento
\texorpdfstring{\softtext{económico.}}{económico.} & La venta ya se
confirmó pero aún no salió el animal. \\
Margen & Resultado gerencial de venta, costo y gastos directos. & Venta
6400 menos costo 4400 menos gasto 200. \\
Periodo & Intervalo mensual de
\texorpdfstring{\softtext{reconocimiento.}}{reconocimiento.} & Un costo
viejo se reconoce en el mes abierto actual. \\
\end{longtable}
}

\section{Una historia completa antes de ver el
código}\label{una-historia-completa-antes-de-ver-el-cuxf3digo}

Imagina una compra DEMO de diez animales a 1000 cada uno. Confirmar la
compra reconoce 10000 y crea diez
\texorpdfstring{\softtext{posiciones.}}{posiciones.} Un gasto de
adquisición de 1000 distribuido por cabeza agrega 100 a cada posición:
el costo unitario pasa a 1100. La deuda del ganado y el gasto tienen su
propia procedencia; el gasto no cambia silenciosamente el precio pactado
con el vendedor.

Después se reciben los animales. Se seleccionan cuatro posiciones
disponibles para vender a 1600 cada una y se reservan para impedir que
otra venta las comprometa
\texorpdfstring{\softtext{simultáneamente.}}{simultáneamente.} Confirmar
la venta reconoce 6400, consume reservas, transfiere 4400 del inventario
al costo de venta y deja los animales vendidos pendientes de entrega. Un
gasto directo de venta de 200 produce margen 1800. Quedan seis
posiciones con inventario de 6600.

Un cobro posterior reduce cartera mediante aplicaciones, pero no crea
otra venta. Una entrega posterior cierra físicamente los ciclos
vendidos, pero no vuelve a pasar el costo a
\texorpdfstring{\softtext{resultados.}}{resultados.} Un flete tardío de
500, distribuido entre los diez animales originales, agrega 50 al costo
de cada ciclo: 200 corresponden a los cuatro vendidos y 300 a los seis
todavía en \texorpdfstring{\softtext{inventario.}}{inventario.} El
margen actual baja a 1600 y el inventario aumenta a 6900. Si ya existía
un cierre histórico, ese cierre no se reescribe para esconder el cambio.

La utilidad del sistema aparece al abrir el expediente: cada importe
remite al acuerdo, a sus animales, a un gasto, una aplicación o un
ajuste. El operador no tiene que deducir retrospectivamente qué ocurrió
a partir de un total mutable.

\chapter{Mapa general de arquitectura y
comunicación}\label{mapa-general-de-arquitectura-y-comunicaciuxf3n}

\section{Las capas y sus
responsabilidades}\label{las-capas-y-sus-responsabilidades}

El navegador ejecuta React. React presenta formularios, conserva el
estado de interacción y envía HTTP. El navegador no tiene la contraseña
de PostgreSQL ni puede ejecutar SQL; tampoco es la autoridad del cálculo
monetario. Caddy recibe la conexión del usuario en
\texorpdfstring{\softtext{localhost:15180.}}{localhost:15180.} Para
\texorpdfstring{\softcode{/api/*}}{/api/*} utiliza el backend Java; para
las páginas y archivos estáticos utiliza el servidor del frontend.

Java recibe las solicitudes, valida su forma, comprueba sesión y
autorización, elige el dueño de la operación y ejecuta el caso de uso.
Los pools JDBC proporcionan conexiones a PostgreSQL A, B y
\texorpdfstring{\softtext{autenticación.}}{autenticación.} Una
transacción económica utiliza solamente el PostgreSQL
\texorpdfstring{\softtext{propietario.}}{propietario.} PostgreSQL
conserva las restricciones e implementa los comandos canónicos con SQL y
PL/pgSQL. El resultado comprometido vuelve como JSON a Java y luego a
React.

Dos bases de negocio son clústeres distintos con volúmenes distintos. A
administra maestros y sus hechos locales; B administra los hechos de su
sede. Las publicaciones y suscripciones distribuyen filas entre A y B.
La replicación es asíncrona: una confirmación local puede responder
antes de que la copia remota la reciba. El tercer PostgreSQL es un
proveedor simulado para demostrar FDW; no es un organismo oficial ni una
integración comercial.

\begin{figure}[htbp]
\centering
\begin{tikzpicture}[node distance=8mm and 13mm,every node/.style={font=\small},box/.style={draw=brand,rounded corners,align=center,minimum width=30mm,minimum height=10mm,fill=surface},>=stealth]
\node[box] (u) {Usuario\\Navegador React};
\node[box,below=of u] (c) {Caddy\\localhost:15180};
\node[box,below left=of c] (f) {Frontend estático\\React y CSS};
\node[box,below right=of c] (j) {Java API\\sesión y comandos};
\node[box,below left=of j] (pa) {Pgpool A\\raw};
\node[box,below right=of j] (pb) {Pgpool B\\raw};
\node[box,below=of pa] (a) {PostgreSQL A\\maestros y sede A};
\node[box,below=of pb] (b) {PostgreSQL B\\sede B};
\node[box,below=of a] (p) {PostgreSQL proveedor\\FDW sólo lectura};
\draw[->] (u)--(c);\draw[->] (c)--(f);\draw[->] (c)--node[right]{/api}(j);
\draw[->] (j)--(pa);\draw[->] (j)--(pb);\draw[->] (pa)--(a);\draw[->] (pb)--(b);
\draw[<->,dashed] (a)--node[above,font=\scriptsize]{réplica lógica}(b);\draw[->,dashed] (a)--(p);
\end{tikzpicture}
\caption{Comunicación del perfil académico. Autenticación accede a A; archivos privados se almacenan fuera del frontend.}
\end{figure}

\section{Diseño inicial e implementación
real}\label{diseuxf1o-inicial-e-implementaciuxf3n-real}

El proyecto contiene una arquitectura objetivo con diez módulos y
nombres de interfaces como
\texorpdfstring{\softcode{OwnerDatabaseRouter}}{OwnerDatabaseRouter} o
\texorpdfstring{\softcode{BranchTransactionRunner}}{BranchTransactionRunner}.
Esos nombres describen una separación planeada, no un permiso para
afirmar que existen clases con ese nombre. En la revisión documentada,
las responsabilidades ejecutables se concentran en los paquetes
\texorpdfstring{\softcode{acceso}}{acceso},
\texorpdfstring{\softcode{operaciones}}{operaciones},
\texorpdfstring{\softcode{evidencia}}{evidencia},
\texorpdfstring{\softcode{reportes}}{reportes},
\texorpdfstring{\softcode{integraciones}}{integraciones} y
\texorpdfstring{\softcode{shared}}{shared}. El capítulo Java identifica
la clase concreta que realiza cada paso. El controlador central y el
servicio de comandos comparten validación y ejecución; PostgreSQL reúne
reglas persistentes en funciones
\texorpdfstring{\softtext{canónicas.}}{canónicas.}

Esta concentración reduce código duplicado, pero hace importante
comprender el despacho de comandos y consultas. Añadir una función al
SQL sin autorizarla en el backend no crea automáticamente una acción
pública. Añadir una pantalla tampoco concede un permiso. La arquitectura
se comprende por responsabilidades y fronteras, no contando carpetas con
nombres ideales.

\section{Quién escribe y quién sólo
consulta}\label{quiuxe9n-escribe-y-quiuxe9n-suxf3lo-consulta}

La regla ``un escritor por agregado'' significa que no se actualiza la
misma fila comercial indistintamente en A y B. El propietario está
fijado por el tipo de dato y la sede. A conserva los maestros centrales;
cada sede produce sus hechos locales en su dueño. Pgpool no toma esta
decisión y no promociona B cuando A falla. Cada Pgpool raw tiene un
único backend fijo.

Una consulta puede leer una copia replicada, pero esa copia tiene una
fecha y una \texorpdfstring{\softtext{procedencia.}}{procedencia.} Un
consolidado A/B no es una transacción global atómica. Por eso el reporte
debe indicar fuente, corte y antigüedad conocida o
\texorpdfstring{\softtext{desconocida.}}{desconocida.} Un
\texorpdfstring{\softcode{null}}{null} en antigüedad no debe convertirse
en cero sólo para aparentar
\texorpdfstring{\softtext{sincronización.}}{sincronización.}

La sesión y la autorización vigentes son centrales en A. Si cae A, las
nuevas solicitudes autenticadas a ambas sedes reciben un error
\texorpdfstring{\softtext{controlado.}}{controlado.} Una transacción B
que ya había sido autorizada puede terminar. Esto es una dependencia
decidida del MVP y debe explicarse en la exposición: dos escritores
académicos sobre conjuntos disjuntos no equivalen a disponibilidad
continua de toda la \texorpdfstring{\softtext{aplicación.}}{aplicación.}

\section{Qué pasa al pulsar
confirmar}\label{quuxe9-pasa-al-pulsar-confirmar}

Primero React conserva una intención con su clave idempotente y la
versión revisada. Envía el body tipado,
\texorpdfstring{\softcode{Idempotency-Key}}{Idempotency-Key},
\texorpdfstring{\softcode{If-Match}}{If-Match} cuando corresponde y el
token CSRF. Caddy conserva el ETag de la API: la compresión gzip se
aplica al frontend estático, no a la respuesta API que participa en la
\texorpdfstring{\softtext{precondición.}}{precondición.}

Java autentica, valida la sede y los permisos, reconoce el comando
admitido y determina el nodo. La transacción del dueño reclama la clave,
resuelve un replay terminado antes de exigir una versión nueva, bloquea
periodo y recursos en orden y valida las precondiciones
\texorpdfstring{\softtext{persistentes.}}{persistentes.} Si todo
procede, registra liquidación, inventario o dinero, auditoría, outbox y
resultado \texorpdfstring{\softtext{idempotente.}}{idempotente.} Si algo
falla antes del commit, se revierte todo ese conjunto. La respuesta
aparece sólo después del resultado de la
\texorpdfstring{\softtext{transacción.}}{transacción.}

Si la respuesta se pierde después del commit, el navegador no sabe
todavía si se registró la
\texorpdfstring{\softtext{operación.}}{operación.} Mostrar ``no se
guardó'' sería \texorpdfstring{\softtext{incorrecto.}}{incorrecto.} La
UI conserva la intención, consulta el resultado o reintenta exactamente
la misma clave y contenido. Una nueva clave podría crear un segundo
hecho; por eso se distingue un error de validación de un resultado
pendiente de comprobar.

\chapter{Guía completa para ejecutar en
local}\label{guuxeda-completa-para-ejecutar-en-local}

\section{Requisitos y versiones}\label{requisitos-y-versiones}

La aplicación usa Java 21, Maven Wrapper 3.9.12, Spring Boot 4.1.1,
React 19.2.4, TypeScript 5.9.3, Vite 8.0.3, PostgreSQL servidor 18.6 y
Pgpool 4.7.3. El cliente \texorpdfstring{\softcode{psql}}{psql} del host
puede tener versión menor distinta; en la evidencia original era 18.3.
El frontend exige Node compatible con
\texorpdfstring{\softcode{\textasciicircum{}20.19.0 || >=22.12.0}}{\textasciicircum{}20.19.0 || >=22.12.0},
según su \texorpdfstring{\softtext{manifiesto.}}{manifiesto.} Docker y
Compose deben funcionar para el usuario que ejecuta los comandos.

El wrapper evita depender de la versión global de Maven. Para las
herramientas Python se utilizan entornos privados y los requisitos
fijados; no se recomienda instalar librerías de pruebas dentro del
runtime de la \texorpdfstring{\softtext{aplicación.}}{aplicación.} Las
imágenes de infraestructura están fijadas por etiqueta y digest donde el
Compose lo \texorpdfstring{\softtext{especifica.}}{especifica.} Las
imágenes \texorpdfstring{\softtext{Java/frontend}}{Java/frontend} se
construyen \texorpdfstring{\softtext{localmente.}}{localmente.}

\section{Comprobar el checkout y el
entorno}\label{comprobar-el-checkout-y-el-entorno}

Ejecuta desde la raíz del
\texorpdfstring{\softtext{repositorio.}}{repositorio.} No copies rutas
privadas de otra persona. Si estás en la computadora original, puedes
entrar así:

\begin{Shaded}
\begin{Highlighting}[]
\BuiltInTok{cd}\NormalTok{ /home/aldo/Descargas/PROYECTO\_ZACA\_TOPICOS\_AVANZADOS}
\FunctionTok{git}\NormalTok{ status }\AttributeTok{{-}{-}short} \AttributeTok{{-}{-}branch}
\FunctionTok{make}\NormalTok{ doctor}
\end{Highlighting}
\end{Shaded}

\texorpdfstring{\softcode{make doctor}}{make doctor} consulta versiones,
CPU, memoria y puertos. Un puerto ocupado no demuestra un error: puede
pertenecer a BoviCierre ya iniciado. Identifica primero el servicio. El
script no mata procesos, no borra volúmenes y no muestra archivos de
secretos. Si Java o Docker faltan, resuelve ese requisito antes de
interpretar los errores siguientes como fallos del dominio.

No existe un instalador universal de
\texorpdfstring{\softtext{Docker/Java}}{Docker/Java} para todos los
sistemas en este proyecto. Las instrucciones que siguen corresponden al
entorno Linux probado. En otra máquina deben verificarse permisos de
Docker, disponibilidad de imágenes y puertos. El motor PostgreSQL vive
en contenedores: instalar \texorpdfstring{\softcode{psql}}{psql} sirve
como cliente, no sustituye el servidor de Compose.

\section{Arranque académico
integrado}\label{arranque-acaduxe9mico-integrado}

\begin{Shaded}
\begin{Highlighting}[]
\FunctionTok{make}\NormalTok{ academic}
\FunctionTok{make}\NormalTok{ demo}
\end{Highlighting}
\end{Shaded}

\texorpdfstring{\softcode{make academic}}{make academic} ejecuta
\texorpdfstring{\softcode{scripts/academic.py}}{scripts/academic.py}.
Verifica que no exista una API manual concurrente en18080, prepara el
directorio de documentos y comprueba
\texorpdfstring{\softtext{propietario/permisos.}}{propietario/permisos.}
Después llama al bootstrap de BD, aplica migraciones, configura las
relaciones académicas y
\texorpdfstring{\softtext{construye/levanta}}{construye/levanta}
backend, frontend y Caddy. Los ocho servicios pertenecen al proyecto
\texorpdfstring{\softcode{bovicierre-academic}}{bovicierre-academic}. El
arranque normal conserva datos y no contiene un
\texorpdfstring{\softcode{down -v}}{down -v}.

\texorpdfstring{\softcode{make demo}}{make demo} instala las
dependencias Python de BD en el entorno privado y ejecuta
\texorpdfstring{\softcode{db/seed/demo.py}}{db/seed/demo.py}. La semilla
utiliza comandos de negocio y claves reproducibles; no inserta totales
financieros a mano ni borra una BD ajena. Se comprueban el proyecto,
base y empresa DEMO. Ejecutar de nuevo la semilla debe recuperar los
resultados anteriores sin duplicar compras o pagos. El modo de
producción rechaza bootstrap y seed
\texorpdfstring{\softtext{académicos.}}{académicos.}

Abre
\texorpdfstring{\softcode{http://localhost:15180}}{http://localhost:15180}.
El acceso comienza obteniendo un CSRF anónimo, realizando login y
seleccionando una sede permitida. Los usuarios demo son
\texorpdfstring{\softcode{demo.admin}}{demo.admin},
\texorpdfstring{\softcode{demo.operador}}{demo.operador},
\texorpdfstring{\softcode{demo.caja}}{demo.caja} y
\texorpdfstring{\softcode{demo.auditor}}{demo.auditor}. Las contraseñas
se generan en archivos locales privados; este manual no contiene sus
valores.

Para consultar tu propia contraseña de ADMIN, usa únicamente la terminal
local y no pegues el resultado en capturas, chats o Git:

\begin{Shaded}
\begin{Highlighting}[]
\FunctionTok{cat} \StringTok{"}\VariableTok{$HOME}\StringTok{/.local/share/bovicierre{-}academic/secrets/demo\_password"}
\end{Highlighting}
\end{Shaded}

Los otros archivos son
\texorpdfstring{\softcode{demo\_operador\_password}}{demo\_operador\_password},
\texorpdfstring{\softcode{demo\_caja\_password}}{demo\_caja\_password} y
\texorpdfstring{\softcode{demo\_auditor\_password}}{demo\_auditor\_password}
dentro del mismo \texorpdfstring{\softtext{directorio.}}{directorio.} El
directorio debe tener permisos0700 y los archivos
\texorpdfstring{\softtext{sensibles0600.}}{sensibles0600.} En producción
futura no se reutilizan estas credenciales ni datos DEMO.

\section{Puertos y direcciones}\label{puertos-y-direcciones}

{\def\LTcaptype{none} % do not increment counter
\begin{longtable}[]{@{}
  >{\raggedright\arraybackslash}p{(\linewidth - 4\tabcolsep) * \real{0.3333}}
  >{\raggedright\arraybackslash}p{(\linewidth - 4\tabcolsep) * \real{0.3333}}
  >{\raggedright\arraybackslash}p{(\linewidth - 4\tabcolsep) * \real{0.3333}}@{}}
\toprule\noalign{}
\begin{minipage}[b]{\linewidth}\raggedright
Servicio
\end{minipage} & \begin{minipage}[b]{\linewidth}\raggedright
Dirección en el host
\end{minipage} & \begin{minipage}[b]{\linewidth}\raggedright
Función
\end{minipage} \\
\midrule\noalign{}
\endhead
\bottomrule\noalign{}
\endlastfoot
Interfaz integrada &
\texorpdfstring{\softtext{http://localhost:15180}}{http://localhost:15180}
& Entrada recomendada del usuario. \\
PostgreSQL A &
\texorpdfstring{\softtext{127.0.0.1:15432}}{127.0.0.1:15432} & Maestros,
acceso y hechos sede A. \\
PostgreSQL B &
\texorpdfstring{\softtext{127.0.0.1:15433}}{127.0.0.1:15433} & Hechos
sede B. \\
Proveedor simulado &
\texorpdfstring{\softtext{127.0.0.1:15434}}{127.0.0.1:15434} & Datos FDW
\texorpdfstring{\softtext{sintéticos.}}{sintéticos.} \\
Pgpool A & \texorpdfstring{\softtext{127.0.0.1:19999}}{127.0.0.1:19999}
& Canal JDBC propietario A. \\
Pgpool B & \texorpdfstring{\softtext{127.0.0.1:19998}}{127.0.0.1:19998}
& Canal JDBC propietario B. \\
API manual &
\texorpdfstring{\softtext{http://127.0.0.1:18080}}{http://127.0.0.1:18080}
& Alternativa de \texorpdfstring{\softtext{desarrollo.}}{desarrollo.} \\
Vite manual &
\texorpdfstring{\softtext{http://127.0.0.1:15173}}{http://127.0.0.1:15173}
& Alternativa de \texorpdfstring{\softtext{desarrollo.}}{desarrollo.} \\
\end{longtable}
}

Los puertos publicados son loopback. Dentro de la red de Compose se
emplean nombres como \texorpdfstring{\softcode{pg-a}}{pg-a},
\texorpdfstring{\softcode{pgpool-a}}{pgpool-a} y
\texorpdfstring{\softcode{backend}}{backend}; el puerto interno puede
ser distinto del publicado. \texorpdfstring{\softcode{0.0.0.0}}{0.0.0.0}
dentro del contenedor permite recibir al proxy, pero no implica que se
haya expuesto PostgreSQL a Internet.

\section{Comprobar servicios sin modificar
datos}\label{comprobar-servicios-sin-modificar-datos}

\begin{Shaded}
\begin{Highlighting}[]
\ExtensionTok{docker}\NormalTok{ compose }\AttributeTok{{-}p}\NormalTok{ bovicierre{-}academic }\DataTypeTok{\textbackslash{}}
  \AttributeTok{{-}f}\NormalTok{ infra/compose/compose.academic.yml }\DataTypeTok{\textbackslash{}}
  \AttributeTok{{-}f}\NormalTok{ infra/compose/application.yml ps}
\ExtensionTok{curl} \AttributeTok{{-}{-}fail}\NormalTok{ http://localhost:15180/api/v1/salud}
\end{Highlighting}
\end{Shaded}

El endpoint de salud diferencia el proceso y las
\texorpdfstring{\softtext{dependencias.}}{dependencias.} La información
de replicación y mensajes debe consultarse con el alcance autorizado; no
se deduce de que el HTML de login aparezca. Una caída de B no debe
confundirse con una caída de A: las consecuencias sobre comandos,
sesiones y lecturas son distintas.

\section{Detener y volver a iniciar}\label{detener-y-volver-a-iniciar}

\begin{Shaded}
\begin{Highlighting}[]
\FunctionTok{make}\NormalTok{ stop}
\FunctionTok{make}\NormalTok{ academic}
\end{Highlighting}
\end{Shaded}

\texorpdfstring{\softcode{make stop}}{make stop} detiene únicamente los
servicios del proyecto declarado, sin eliminar
\texorpdfstring{\softtext{volúmenes.}}{volúmenes.} Si elegiste la
alternativa de desarrollo manual, sus procesos se detienen desde la
terminal que ejecuta \texorpdfstring{\softcode{make dev}}{make dev},
mediante Ctrl+C. No utilices
\texorpdfstring{\softcode{killall java}}{killall java},
\texorpdfstring{\softcode{docker system prune}}{docker system prune} o
borrados de volúmenes para resolver una ocupación de puerto: la
computadora puede contener otros proyectos.

\section{Desarrollo con Java y Vite en primer
plano}\label{desarrollo-con-java-y-vite-en-primer-plano}

Primero detén el perfil integrado para evitar API simultáneas y el
puerto15180 ocupado. Mantén o vuelve a iniciar la infraestructura de BD,
instala el frontend y ejecuta el ciclo manual:

\begin{Shaded}
\begin{Highlighting}[]
\FunctionTok{make}\NormalTok{ stop}
\FunctionTok{sh}\NormalTok{ db/bootstrap/start.sh}
\ExtensionTok{npm} \AttributeTok{{-}{-}prefix}\NormalTok{ frontend ci}
\FunctionTok{make}\NormalTok{ dev}
\end{Highlighting}
\end{Shaded}

\texorpdfstring{\softcode{make dev}}{make dev} ejecuta
\texorpdfstring{\softcode{scripts/dev.py}}{scripts/dev.py}, que crea dos
procesos hijos:
\texorpdfstring{\softcode{backend/run-local.sh}}{backend/run-local.sh} y
\texorpdfstring{\softcode{npm run dev}}{npm run dev} en frontend.
Rechaza
\texorpdfstring{\softtext{puertos18080/15173/15180}}{puertos18080/15173/15180}
ocupados y conserva la relación con sus hijos. Ctrl+C los termina
gradualmente; una espera vencida sólo afecta a esos hijos
\texorpdfstring{\softtext{identificados.}}{identificados.} Vite
sirve15173 y envía las peticiones API al backend local conforme a su
\texorpdfstring{\softtext{configuración.}}{configuración.}

El script
\texorpdfstring{\softcode{backend/run-local.sh}}{backend/run-local.sh}
lee los secretos externos, exporta las variables que requiere Spring y
escribe \texorpdfstring{\softtext{PID/registro}}{PID/registro} en un
directorio privado. No se utiliza una contraseña fija en
\texorpdfstring{\softcode{application.yml}}{application.yml}. No debes
iniciar otra instancia Java encima del backend integrado: ambos podrían
despachar outbox y competir por el presupuesto de
\texorpdfstring{\softtext{conexiones.}}{conexiones.}

\section{Variables de configuración que debes
conocer}\label{variables-de-configuraciuxf3n-que-debes-conocer}

Una variable de entorno es un dato que el proceso recibe al iniciarse.
No es una columna de BD ni una constante
\texorpdfstring{\softtext{TypeScript.}}{TypeScript.} Las variables
\texorpdfstring{\softcode{BOVI\_DB\_*\_PASSWORD}}{BOVI\_DB\_*\_PASSWORD}
son secretos; los nombres y funciones pueden documentarse, sus valores
no. El Compose utiliza archivos montados en
\texorpdfstring{\softcode{/run/secrets}}{/run/secrets} y el entrypoint
los transforma para el proceso Java.

{\def\LTcaptype{none} % do not increment counter
\begin{longtable}[]{@{}
  >{\raggedright\arraybackslash}p{(\linewidth - 4\tabcolsep) * \real{0.3333}}
  >{\raggedright\arraybackslash}p{(\linewidth - 4\tabcolsep) * \real{0.3333}}
  >{\raggedright\arraybackslash}p{(\linewidth - 4\tabcolsep) * \real{0.3333}}@{}}
\toprule\noalign{}
\begin{minipage}[b]{\linewidth}\raggedright
Variable
\end{minipage} & \begin{minipage}[b]{\linewidth}\raggedright
Propósito
\end{minipage} & \begin{minipage}[b]{\linewidth}\raggedright
Valor o tratamiento
\end{minipage} \\
\midrule\noalign{}
\endhead
\bottomrule\noalign{}
\endlastfoot
\texorpdfstring{\softcode{BOVI\_SECRET\_DIR}}{BOVI\_SECRET\_DIR} &
Directorio de secretos de
\texorpdfstring{\softtext{scripts/Compose.}}{scripts/Compose.} & Por
defecto fuera del checkout. \\
\texorpdfstring{\softcode{BOVI\_PRIVATE\_DIR}}{BOVI\_PRIVATE\_DIR} &
Raíz privada de documentos y respaldo. &
\texorpdfstring{\softcode{\textasciitilde{}/.local/share/bovicierre-academic}}{\textasciitilde{}/.local/share/bovicierre-academic}. \\
\texorpdfstring{\softcode{BOVI\_DOCUMENTS\_DIR}}{BOVI\_DOCUMENTS\_DIR} &
Almacenamiento del adaptador
\texorpdfstring{\softtext{documental.}}{documental.} & Volumen privado,
no webroot. \\
\texorpdfstring{\softcode{BOVI\_UID/BOVI\_GID}}{BOVI\_UID/BOVI\_GID} &
Identidad de ejecución de
\texorpdfstring{\softtext{contenedores.}}{contenedores.} & Usuario
local; no root por comodidad. \\
\texorpdfstring{\softcode{BOVI\_API\_PORT}}{BOVI\_API\_PORT} & Puerto
API manual o \texorpdfstring{\softtext{restaurada.}}{restaurada.} &
18080 normal;18180 ensayo HTTP aislado. \\
\texorpdfstring{\softcode{BOVI\_BIND\_ADDRESS}}{BOVI\_BIND\_ADDRESS} &
Interfaz de escucha de Java. & Loopback manual; interfaz interna en
Docker. \\
\texorpdfstring{\softcode{BOVI\_DB\_A\_URL/BOVI\_DB\_B\_URL}}{BOVI\_DB\_A\_URL/BOVI\_DB\_B\_URL}
& JDBC de los dueños. & Pgpool A/B; nunca intercambiar por fallo. \\
\texorpdfstring{\softcode{BOVI\_AUTH\_URL}}{BOVI\_AUTH\_URL} & JDBC
dedicado de
\texorpdfstring{\softtext{acceso/sesiones.}}{acceso/sesiones.} &
PostgreSQL A central. \\
\texorpdfstring{\softcode{BOVI\_OUTBOX\_ENABLED}}{BOVI\_OUTBOX\_ENABLED}
& Activación del worker. & Configurable; no altera dueño ni
\texorpdfstring{\softtext{deduplicación.}}{deduplicación.} \\
\texorpdfstring{\softcode{BOVI\_RUNTIME\_DIR}}{BOVI\_RUNTIME\_DIR} &
\texorpdfstring{\softtext{PID/registro}}{PID/registro} de API manual. &
Directorio separado por
\texorpdfstring{\softtext{origen/restauración.}}{origen/restauración.} \\
\texorpdfstring{\softcode{BOVI\_TEST\_VENV}}{BOVI\_TEST\_VENV} & Entorno
Python auxiliar. & Caché privada de pruebas SQL. \\
\texorpdfstring{\softcode{BOVI\_API\_IT}}{BOVI\_API\_IT} & Activación de
pruebas HTTP reales. & \texorpdfstring{\softcode{true}}{true} sólo con
API preparada. \\
\texorpdfstring{\softcode{BOVI\_API\_TEST\_URL}}{BOVI\_API\_TEST\_URL} &
Base API de pruebas HTTP. &
\texorpdfstring{\softtext{localhost15180/api/v1}}{localhost15180/api/v1}
en evidencia final. \\
\texorpdfstring{\softcode{BOVI\_TEST\_PASSWORD\_FILE}}{BOVI\_TEST\_PASSWORD\_FILE}
& Archivo de contraseña de fixture. & Se lee localmente, no se
versiona. \\
\texorpdfstring{\softcode{SPRING\_PROFILES\_ACTIVE}}{SPRING\_PROFILES\_ACTIVE}
& Perfil Spring. & Producción exige Secure; local permite HTTP. \\
\end{longtable}
}

\texorpdfstring{\softcode{.env.example}}{.env.example} es una
orientación, no una lectura automática por todos los programas. El
propio comentario pide exportar las variables para make o usar
\texorpdfstring{\softcode{--env-file}}{--env-file} con Compose. No
supongas que copiar un archivo \texorpdfstring{\softcode{.env}}{.env}
cambia un proceso Java ya arrancado. Algunas variables son de scripts y
otras de Spring; su consumidor exacto aparece en el atlas de
\texorpdfstring{\softtext{configuración.}}{configuración.}

\section{Verificación y su alcance}\label{verificaciuxf3n-y-su-alcance}

\begin{Shaded}
\begin{Highlighting}[]
\FunctionTok{make}\NormalTok{ verify{-}docs}
\FunctionTok{make}\NormalTok{ verify}
\FunctionTok{make}\NormalTok{ verify{-}api}
\FunctionTok{make}\NormalTok{ verify{-}sql}
\ExtensionTok{npm} \AttributeTok{{-}{-}prefix}\NormalTok{ frontend run typecheck}
\ExtensionTok{npm} \AttributeTok{{-}{-}prefix}\NormalTok{ frontend run build}
\ExtensionTok{npm} \AttributeTok{{-}{-}prefix}\NormalTok{ frontend test}
\ExtensionTok{npm} \AttributeTok{{-}{-}prefix}\NormalTok{ frontend run generate:api}
\end{Highlighting}
\end{Shaded}

\texorpdfstring{\softcode{verify-docs}}{verify-docs} comprueba IDs
documentales, referencias del esquema lógico, cinco tests de
herramientas de diagramas y OpenAPI. No confirma compras ni carreras.
\texorpdfstring{\softcode{verify}}{verify} agrega Maven, frontend y
comparación de catálogo A/B.
\texorpdfstring{\softcode{verify-api}}{verify-api} activa los tests HTTP
contra una API viva, utilizando la contraseña externa.
\texorpdfstring{\softcode{verify-sql}}{verify-sql} prepara namespaces QA
y ejecuta el runner SQL: incluye casos de negocio, guardas y
\texorpdfstring{\softtext{pesajes/reportes.}}{pesajes/reportes.} Debes
distinguir sus bases de prueba de la base de negocio principal.

El generador frontend produce tipos desde OpenAPI. Después de modificar
el contrato, genera de nuevo y comprueba
\texorpdfstring{\softtext{tipos/build;}}{tipos/build;} el navegador
conserva su adaptador HTTP propio. Tener un tipo generado no realiza una
llamada ni prueba que Java devuelva ese formato.

Los E2E requieren el navegador y los fixtures configurados por
\texorpdfstring{\softcode{frontend/playwright.config.ts}}{frontend/playwright.config.ts}
y
\texorpdfstring{\softcode{tests\_e2e/runtime.ts}}{tests\_e2e/runtime.ts}.
El script es
\texorpdfstring{\softcode{npm --prefix frontend run test:e2e}}{npm --prefix frontend run test:e2e};
revisa primero sus requisitos de API, contraseña y estado del periodo.
Algunos casos provocan incidencias controladas; no ejecutes
indiscriminadamente pruebas de caída mientras alguien usa la demo. La
evidencia final separó15 casos integrados de una prueba física adicional
del proveedor, no una única suite invisible de16 casos
\texorpdfstring{\softtext{concurrentes.}}{concurrentes.}

\section{Respaldo y recuperación
aislada}\label{respaldo-y-recuperaciuxf3n-aislada}

\begin{Shaded}
\begin{Highlighting}[]
\FunctionTok{make}\NormalTok{ backup BACKUP\_DIR=}\StringTok{"}\VariableTok{$HOME}\StringTok{/.local/share/bovicierre{-}academic/backups/ensayo{-}nuevo"}
\FunctionTok{make}\NormalTok{ restore{-}check BACKUP\_DIR=}\StringTok{"}\VariableTok{$HOME}\StringTok{/.local/share/bovicierre{-}academic/backups/ensayo{-}nuevo"}
\end{Highlighting}
\end{Shaded}

El directorio de respaldo debe ser nuevo y externo al checkout. El
procedimiento pausa la API y sus workers, espera que no queden
\texorpdfstring{\softtext{transacciones/mensajes}}{transacciones/mensajes}
en vuelo y comprueba convergencia A/B. Copia archivos y secretos con
hashes, crea dumps de
\texorpdfstring{\softtext{A/B/proveedor}}{A/B/proveedor} y excluye
\texorpdfstring{\softtext{sesiones/publicaciones/suscripciones}}{sesiones/publicaciones/suscripciones}
que podrían apuntar al origen. Finalmente reanuda la fuente; la prueba
de restauración no la deja pausada durante todo el ensayo.

La restauración crea un proyecto distinto
\texorpdfstring{\softcode{bovicierre-restore-<UTC>}}{bovicierre-restore-<UTC>},
con otros puertos y \texorpdfstring{\softtext{volúmenes.}}{volúmenes.}
Verifica integridad antes de configurar suscripciones sin copia inicial,
reconstruye FDW hacia su propio proveedor y conserva
\texorpdfstring{\softtext{inbox/idempotencia.}}{inbox/idempotencia.} Se
invalidan sesiones. El ensayo HTTP requiere una API adicional
únicamente18180, con
\texorpdfstring{\softtext{runtime/documentos/secrets}}{runtime/documentos/secrets}
propios y JDBC del destino;
\texorpdfstring{\softcode{verify-http}}{verify-http} no arranca esa API
por sí solo.

\begin{Shaded}
\begin{Highlighting}[]
\ExtensionTok{python3}\NormalTok{ scripts/recovery.py verify{-}http }\DataTypeTok{\textbackslash{}}
  \AttributeTok{{-}{-}directory} \StringTok{"}\VariableTok{$HOME}\StringTok{/.local/share/bovicierre{-}academic/backups/ensayo{-}nuevo"}
\end{Highlighting}
\end{Shaded}

La restauración V13 registrada
\texorpdfstring{\softtext{pausó8.038}}{pausó8.038} segundos el origen y
\texorpdfstring{\softtext{tardó28.616}}{tardó28.616} segundos en el
ensayo de \texorpdfstring{\softtext{restauración.}}{restauración.} Se
comprobaron archivos, datos e intención económica nueva por HTTP. Son
medidas de una recuperación local en ese equipo; no representan un SLA,
un RTO contractual ni resistencia a perder la laptop.

\chapter{Infraestructura explicada desde sus
archivos}\label{infraestructura-explicada-desde-sus-archivos}

\section{Docker Compose y separación de
servicios}\label{docker-compose-y-separaciuxf3n-de-servicios}

La topología se describe en dos archivos:
\texorpdfstring{\softcode{infra/compose/compose.academic.yml}}{infra/compose/compose.academic.yml}
contiene PostgreSQL y Pgpool;
\texorpdfstring{\softcode{infra/compose/application.yml}}{infra/compose/application.yml}
agrega Java, frontend y Caddy. El nombre de proyecto identifica redes,
contenedores y volúmenes propios. No basta con que un contenedor se
llame \texorpdfstring{\softtext{PostgreSQL:}}{PostgreSQL:} los scripts
comprueban etiquetas de proyecto para evitar administrar un servicio
ajeno.

El backend se construye desde
\texorpdfstring{\softcode{backend/Dockerfile}}{backend/Dockerfile} y el
frontend desde
\texorpdfstring{\softcode{frontend/Dockerfile}}{frontend/Dockerfile}.
Los secretos se montan como archivos, no se incorporan en las capas de
una imagen. Caddy utiliza configuración local de sólo lectura y
directorios temporales acotados. El montaje de documentos permite que
Java escriba en un directorio privado del usuario; el frontend no sirve
ese \texorpdfstring{\softtext{directorio.}}{directorio.}

En el perfil académico inicial, A/B tienen1GiB de límite cada uno, el
\texorpdfstring{\softtext{proveedor512MiB}}{proveedor512MiB} y Java1GiB.
Frontend y Caddy
\texorpdfstring{\softtext{tienen128MiB.}}{tienen128MiB.} Esas cifras son
límites de contenedor, no consumo medido de toda la
\texorpdfstring{\softtext{aplicación.}}{aplicación.} El host también
ejecuta otras \texorpdfstring{\softtext{herramientas.}}{herramientas.}
La asignación debe contrastarse con muestras de memoria, conexiones y
carga antes de prometer capacidad.

\section{Pools y presupuesto de
conexiones}\label{pools-y-presupuesto-de-conexiones}

\texorpdfstring{\softcode{DatabaseConfiguration.pool}}{DatabaseConfiguration.pool}
crea un \texorpdfstring{\softtext{HikariDataSource}}{HikariDataSource} a
partir del prefijo
\texorpdfstring{\softcode{bovi.database.<nombre>}}{bovi.database.<nombre>}.
Los pools de dueños y autenticación tienen máximo4 conexiones; cada
worker tiene2. \texorpdfstring{\softcode{minimumIdle=0}}{minimumIdle=0}
permite no reservar conexiones ociosas
\texorpdfstring{\softtext{obligatorias.}}{obligatorias.}
\texorpdfstring{\softcode{connectionTimeout=3000}}{connectionTimeout=3000}
y
\texorpdfstring{\softcode{validationTimeout=1000}}{validationTimeout=1000}
acotan esperas iniciales;
\texorpdfstring{\softcode{initializationFailTimeout=-1}}{initializationFailTimeout=-1}
evita exigir toda la infraestructura en la construcción del bean. Esa
última opción no hace que una consulta a una BD caída tenga éxito.

Cada Pgpool tiene16 procesos y
\texorpdfstring{\softcode{max\_pool=1}}{max\_pool=1}, apuntando a un
solo dueño. PostgreSQL tiene
\texorpdfstring{\softcode{max\_connections=50}}{max\_connections=50}. Al
presupuestar deben incluirse los clientes directos de auth, workers de
replicación, administración, migración y
\texorpdfstring{\softtext{observación.}}{observación.} No se cuentan dos
veces una misma conexión lógica, pero tampoco se omiten las conexiones
físicas que Pgpool mantiene abiertas. El informe
\texorpdfstring{\softcode{connection-budget.json}}{connection-budget.json}
guarda la observación concreta del entorno probado.

\section{Proxy, ETag y origen común}\label{proxy-etag-y-origen-comuxfan}

Caddy maneja \texorpdfstring{\softcode{/api/*}}{/api/*} con
\texorpdfstring{\softcode{reverse\_proxy backend:18080}}{reverse\_proxy backend:18080}.
El resto utiliza el frontend estático y compresión gzip. Esta separación
evita que la compresión de una respuesta JSON modifique un ETag fuerte
de borrador y haga inválido el
\texorpdfstring{\softcode{If-Match}}{If-Match} enviado por el cliente.
Es un ejemplo de cómo una optimización de infraestructura puede afectar
una regla de concurrencia HTTP.

El navegador trabaja con el mismo origen para páginas y API. Por ello no
se recurre a CORS universal con cookies. La configuración local usa HTTP
loopback, mientras que la preparación futura de producción requiere
HTTPS y cookies Secure. Un archivo de ejemplo con dominio
\texorpdfstring{\softcode{.invalid}}{.invalid} no implica que exista un
certificado emitido o un servidor público
\texorpdfstring{\softtext{contratado.}}{contratado.}

\chapter{Pruebas y evidencia que explican la confianza
disponible}\label{pruebas-y-evidencia-que-explican-la-confianza-disponible}

\section{Qué significa una prueba y qué
no}\label{quuxe9-significa-una-prueba-y-quuxe9-no}

Una prueba unitaria ejecuta una función o un grupo pequeño de objetos.
Es rápida y útil para transformación, validación o
\texorpdfstring{\softtext{serialización.}}{serialización.} Una prueba de
SQL utiliza PostgreSQL real y puede proteger constraints, funciones,
locks y efectos \texorpdfstring{\softtext{persistentes.}}{persistentes.}
Una prueba HTTP atraviesa autenticación, filtros, controlador y acceso a
BD. Una prueba E2E maneja la interfaz en un navegador y observa el
comportamiento para el usuario.

Estas capas son
\texorpdfstring{\softtext{complementarias.}}{complementarias.} Un test
de React que verifica un formato decimal no demuestra que dos
aplicaciones simultáneas respeten el saldo. Una captura visual no
acredita \texorpdfstring{\softtext{persistencia.}}{persistencia.} Una
comprobación de IDs
\texorpdfstring{\softtext{RN/TEST/TASK}}{RN/TEST/TASK} demuestra
cobertura documental, no ejecución del dominio. El manual utiliza la
matriz de pruebas para vincular cada afirmación con su nivel y límites.

El backend final registra24 pruebas:15 unitarias,8HTTP contra Caddy15180
y una integración Testcontainers
\texorpdfstring{\softtext{PostgreSQL18.6}}{PostgreSQL18.6} con
\texorpdfstring{\softtext{migracionesV1-V13.}}{migracionesV1-V13.} El
frontend registra13 unitarias y16 recorridos UI, separados en15
integrados y una prueba adicional de caída física del proveedor. La SQL
registró157 intentos GOLD,22 guardas de versión y66 comandos de
\texorpdfstring{\softtext{pesajes/reportes.}}{pesajes/reportes.} Estas
cuentas no deben sumarse para presentar un porcentaje de cobertura sin
\texorpdfstring{\softtext{denominador.}}{denominador.}

\section{Oráculos económicos y
concurrencia}\label{oruxe1culos-econuxf3micos-y-concurrencia}

Un esperado fijo evita que la prueba comparta el mismo error de cálculo
que el código. En GOLD-02 se escribe explícitamente que cuatro
posiciones de1100 producen costo4400 y margen1800; no se obtiene el
esperado llamando a la misma función que confirma la venta. GOLD-03
\texorpdfstring{\softtext{fija0.34/0.33/0.33}}{fija0.34/0.33/0.33} y
suma1.00, incluyendo el orden estable.

Para una carrera se utilizan dos conexiones independientes y barreras
que coordinan el momento de
\texorpdfstring{\softtext{competencia.}}{competencia.} El estado final
debe demostrar un solo ganador o un rechazo sin efecto. Dormir unos
milisegundos y observar un mensaje no es una prueba fiable de exclusión
mutua. También se comprueba el reintento tras pérdida de respuesta: la
clave previa recupera el resultado y los registros no aumentan.

La matriz \texorpdfstring{\softtext{TEST-01..35}}{TEST-01..35} conserva
los escenarios diseñados y registra evidencia SQL/API/UI por caso. Donde
un caso quedó acotado, la limitación no se transforma en una aprobación
universal. La parte humana del modelo y la rúbrica sigue pendiente
aunque las verificaciones automáticas pasen.

\section{Rendimiento medido}\label{rendimiento-medido}

El benchmark QA se preparó con1901 comandos canónicos y repitió1901 sin
\texorpdfstring{\softtext{duplicación.}}{duplicación.} El dataset propio
contiene200 animales,200 compras,100 ventas,300 gastos,100 pagos y100
entregas; la base QA completa conserva además fixtures
\texorpdfstring{\softtext{anteriores.}}{anteriores.} La medida
utiliza800 lecturas y100 calentamientos, con conexión única y cuatro
conexiones \texorpdfstring{\softtext{persistentes.}}{persistentes.} Los
tiempos incluyen consulta, transferencia y decodificación psycopg, pero
excluyen HTTP, Java y React.

{\def\LTcaptype{none} % do not increment counter
\begin{longtable}[]{@{}
  >{\raggedright\arraybackslash}p{(\linewidth - 4\tabcolsep) * \real{0.3333}}
  >{\raggedright\arraybackslash}p{(\linewidth - 4\tabcolsep) * \real{0.3333}}
  >{\raggedright\arraybackslash}p{(\linewidth - 4\tabcolsep) * \real{0.3333}}@{}}
\toprule\noalign{}
\begin{minipage}[b]{\linewidth}\raggedright
Consulta
\end{minipage} & \begin{minipage}[b]{\linewidth}\raggedright
p50/p95 una conexión en ms
\end{minipage} & \begin{minipage}[b]{\linewidth}\raggedright
p50/p95 cuatro conexiones en ms
\end{minipage} \\
\midrule\noalign{}
\endhead
\bottomrule\noalign{}
\endlastfoot
Inventario & 9.580 /14.088 & 9.478 /13.322 \\
Cartera & 6.866 /10.825 & 10.025 /12.372 \\
Liquidación & 1.534 /3.718 & 1.616 /2.813 \\
Expediente & 0.244 /0.970 & 0.444 /1.075 \\
\end{longtable}
}

La evidencia pertenece a un i7-13620H con16CPU lógicos y
\texorpdfstring{\softtext{aproximadamente15724.6MiB}}{aproximadamente15724.6MiB}
de memoria. Las consultas se ejecutaron con caché caliente y sin
\texorpdfstring{\softtext{escritores.}}{escritores.} Los EXPLAIN del
wrapper PL/pgSQL no desglosan todos los planes internos. Es una
observación útil del MVP; no permite inferir el número de clientes de
producción ni garantizar un p95 HTTP.

\section{Cómo reproducir una medida sin alterar la
demo}\label{cuxf3mo-reproducir-una-medida-sin-alterar-la-demo}

\texorpdfstring{\softcode{db/tests/benchmark.py --prepare}}{db/tests/benchmark.py --prepare}
y \texorpdfstring{\softcode{--measure}}{--measure} están fuera del
runner ordinario porque conservan datos QA y exigen una ventana
\texorpdfstring{\softtext{coordinada.}}{coordinada.} La base autorizada
es
\texorpdfstring{\softcode{bovicierre\_release\_replay}}{bovicierre\_release\_replay},
no la principal. El propio script valida perfil y contexto. Antes de
medir, se coordina una pausa de
\texorpdfstring{\softtext{API/escritores;}}{API/escritores;} después se
reanudan. La reproducción no se incluye como un paso automático del
arranque de usuarios.

La prueba visual se realizó con Chrome visible y GPU Intel integrada,
\texorpdfstring{\softtext{incluyendo390/768/1440}}{incluyendo390/768/1440}
px y un flujo sólo con teclado. No se ha certificado WCAG ni se ha
probado teléfono físico. Para una nueva animación costosa o efecto GPU,
deben repetirse medidas específicas sobre el hardware más débil del
usuario; las capturas actuales no acreditan los FPS de un efecto que aún
no existe.

\chapter{Mantenimiento y diagnóstico con
criterio}\label{mantenimiento-y-diagnuxf3stico-con-criterio}

\section{Si no abre la interfaz}\label{si-no-abre-la-interfaz}

Verifica el estado de los servicios propios, el puerto15180 y la URL. Si
Caddy no inició, consulta sus logs sin imprimir secretos. Si el HTML
aparece pero los datos fallan, comprueba
\texorpdfstring{\softcode{/api/v1/salud}}{/api/v1/salud}: un problema de
backend o autenticación puede ser distinto de un problema de frontend.
Si \texorpdfstring{\softcode{make academic}}{make academic}
rechaza18080, identifica la API manual y termina su proceso
gradualmente; no mates procesos de otros proyectos.

\section{Si un comando devuelve un
conflicto}\label{si-un-comando-devuelve-un-conflicto}

Un409 puede indicar una reserva incompatible, una regla o una clave
reutilizada con contenido diferente. Lee
\texorpdfstring{\softcode{codigo}}{codigo},
\texorpdfstring{\softcode{mensaje}}{mensaje} y los campos públicos de la
respuesta. Un412 indica versión revisada obsoleta; un428 indica falta de
la precondición requerida. No soluciones ambos generando automáticamente
una versión o una clave nueva. Primero actualiza el detalle y comprende
si ya existe un resultado para la intención anterior.

Un422 expresa una regla de dominio, como composición de pesaje
incorrecta o asignaciones que no concilian. El usuario debe corregir las
entradas \texorpdfstring{\softtext{permitidas.}}{permitidas.} Un503
expresa una dependencia necesaria no disponible; no prueba por sí solo
que un comando ya enviado hizo rollback. Si el resultado puede haberse
comprometido, se conserva la intención y se consulta su estado.

\section{Si la réplica no converge}\label{si-la-ruxe9plica-no-converge}

Consulta suscripciones, apply, slots y WAL retenido. Un error de apply
no se resuelve saltando transacciones al azar. Verifica las migraciones,
constraints, propietarios y las referencias de origen. Recuerda que DDL,
roles y secuencias no se replican como filas. La ausencia de actividad
no permite afirmar retraso cero: puede significar que no hubo nada nuevo
o que falta una medida.

Antes de volver a configurar
\texorpdfstring{\softcode{copy\_data=false}}{copy\_data=false}, debes
probar una base compatible común. Usarlo porque ``parece que ya están
iguales'' puede omitir datos. Si faltó un slot o los datos divergieron,
el runbook de recuperación explica la investigación y
\texorpdfstring{\softtext{reconstrucción.}}{reconstrucción.} No cambies
el propietario para que B reciba escrituras que corresponden a A.

\section{Si falta un documento}\label{si-falta-un-documento}

El archivo privado y PostgreSQL no comparten una transacción ACID. Por
eso existe staging, metadata y limpieza
\texorpdfstring{\softtext{controlada.}}{controlada.} Comprueba
existencia, hash, estado y
\texorpdfstring{\softtext{autorización.}}{autorización.} Un archivo
huérfano no debe borrarse sólo por no verlo en la pantalla: puede estar
en staging o corresponder a una versión antigua todavía
\texorpdfstring{\softtext{referenciada.}}{referenciada.} La limpieza
respeta \texorpdfstring{\softtext{edad/estado}}{edad/estado} y falla de
forma conservadora ante dudas.

Una \texorpdfstring{\softtext{descarga403/404}}{descarga403/404} puede
ser una protección de acceso, no un archivo perdido. El UUID no es una
contraseña ni concede derecho a descargar. El QR lleva a una ruta
\texorpdfstring{\softtext{autenticada.}}{autenticada.} Nunca conviertas
un directorio privado en una ruta estática de nginx para ``hacer que
abra''.

\section{Cómo extender el sistema sin romper sus
fronteras}\label{cuxf3mo-extender-el-sistema-sin-romper-sus-fronteras}

Para una nueva acción, primero identifica la regla, el dueño, el periodo
y el efecto \texorpdfstring{\softtext{persistente.}}{persistente.}
Decide si es una consulta o un comando
\texorpdfstring{\softtext{económico.}}{económico.} Actualiza el contrato
con
\texorpdfstring{\softtext{entradas/respuestas/errores}}{entradas/respuestas/errores}
y luego el SQL canónico y el adaptador Java. El frontend debe conservar
las claves y versiones
\texorpdfstring{\softtext{pertinentes.}}{pertinentes.} Añade una prueba
de comportamiento para el riesgo concreto; no copies cada línea de
implementación como si eso garantizara
\texorpdfstring{\softtext{corrección.}}{corrección.}

Una nueva columna exige migración numerada centralmente y sincronización
con el diccionario y modelo lógico. No edites retrospectivamente una
migración aplicada para esconder un upgrade: crea la siguiente versión.
Coordina A/B antes de publicar cambios
\texorpdfstring{\softtext{dependientes.}}{dependientes.} Si la columna
es técnica y local, documenta por qué no se replica. Si agrega datos
sensibles, revisa permisos, logs y
\texorpdfstring{\softtext{exportaciones.}}{exportaciones.}

Una nueva integración externa comienza por la operación requerida y una
API oficial mantenida. API Anything puede ayudar en lecturas públicas
autorizadas cuando no existe una opción suficiente, pero un MCP es una
herramienta del agente, no el runtime automático de
\texorpdfstring{\softtext{BoviCierre.}}{BoviCierre.} Para que el
producto consuma un servicio se necesita su propio adaptador backend,
validación, dominios autorizados, timeout, errores y credenciales
privadas. Mejorar la UI no obliga a copiar el backend de una referencia
visual.

\chapter{Herramientas de desarrollo, trazabilidad y
responsabilidades}\label{herramientas-de-desarrollo-trazabilidad-y-responsabilidades}

Las herramientas auxiliares ayudan a construir y comprobar el producto,
pero no forman parte de la experiencia diaria del operador. Chrome
DevTools MCP sirve para inspeccionar DOM/CSS, consola, red y capturas.
UI/UX Pro Max fundamenta jerarquía, espaciado, componentes y estados.
API Anything ofrece operaciones concretas de lectura cuando
\texorpdfstring{\softtext{corresponden.}}{corresponden.} Playwright es
el runner reproducible de pruebas frontend; no debe confundirse con un
servidor MCP que controla un navegador en una sesión de trabajo.

Graphify conecta reglas, tablas, comandos, endpoints, pantallas y
pruebas. El corte final registra2509 nodos,4881 aristas y248
comunidades, con53 documentos reinterpretados y102 fuentes
\texorpdfstring{\softtext{contrastadas.}}{contrastadas.} Los documentos
históricos sin cambios mantienen extracción anterior. Una arista
conceptual no demuestra una FK ni una prueba aprobada: se comprueba
contra SQL, catálogo o evidencia. Sus salidas permanecen excluidas de
Git.

Ponytail orienta una implementación pequeña pero completa: reutilizar
piezas y evitar abstracciones sin frontera real, sin recortar controles
monetarios, seguridad o
\texorpdfstring{\softtext{accesibilidad.}}{accesibilidad.} ECC aporta
skills y mecanismos de trabajo según lo instalado; no existe una
obligación de invocar un comando \texorpdfstring{\softcode{/ecc}}{/ecc}
universal. Las instalaciones locales se documentan por versión y los
\texorpdfstring{\softtext{perfiles/caches}}{perfiles/caches} no se
incluyen en el bundle de React.

Aldo conserva responsabilidad humana de backend, BD e infraestructura;
Jared de frontend. Los agentes ejecutaron trabajo técnico, pero no
suplantan la revisión de ambos. Las55 tareas originales tienen52
verificaciones técnicas locales, una pendiente externa y dos diferidas.
Las55 revisiones humanas siguen
\texorpdfstring{\softtext{pendientes.}}{pendientes.} El commit MVP9282
b75 y el ajuste de wrapper92 d3 d41 están publicados en main; las notas
antiguas de ``sin publicar'' describen el corte anterior, no el estado
GitHub actual.

\chapter{Preparación comercial y límites del
alcance}\label{preparaciuxf3n-comercial-y-luxedmites-del-alcance}

La hipótesis inicial es3000MXN por sede/mes, hasta cinco usuarios,5GB de
documentos y dos horas mensuales de soporte. Implantación y capacitación
se cotizarían aparte; un piloto posterior tendría30 días de alcance
acordado. Es una oferta a validar, no precio
\texorpdfstring{\softtext{demostrado.}}{demostrado.} El valor del ganado
pertenece al negocio del cliente: no se contabiliza como ingreso SaaS.

Las entrevistas todavía no se
\texorpdfstring{\softtext{realizaron.}}{realizaron.} La cifra pública
histórica de bovinos en Chiapas contextualiza la actividad regional; no
determina el inventario actual ni el número de negocios
\texorpdfstring{\softtext{compradores.}}{compradores.} El snapshot
geográfico de INEGI conserva claves y procedencia; los datos sintéticos
permanecen separados. Una referencia SNIIM exige categoría, unidad,
presentación, mercado y periodo, y no sustituye un precio pactado.

OCR, captura offline avanzada, marketplace, custodia de dinero, CFDI y
autorizaciones sanitarias quedan fuera. La apertura de inventario real
previo al piloto exige un procedimiento conciliado; no se inventan
compras para \texorpdfstring{\softtext{sustituirlo.}}{sustituirlo.}
Traspasos que cambien dueño escritor y devoluciones físicas después de
entrega también requieren evolución del alcance.

La preparación para una VM Linux con Compose y Caddy conserva archivos
de ejemplo, redes, secretos y runbooks. No se contrató servidor, no se
cambió DNS y no se emitió TLS público. Dos PostgreSQL en el mismo host
son una demostración académica con escritores disjuntos; no redundancia
física contra pérdida de la máquina. El siguiente paso comercial
requiere revisión, usuarios reales, datos iniciales, recuperación
externa y medición en el servidor elegido.

\chapter{Cómo interpretar variables, cifras y
fórmulas}\label{cuxf3mo-interpretar-variables-cifras-y-fuxf3rmulas}

\section{Tres clases de variables}\label{tres-clases-de-variables}

En este software ``variable'' puede significar tres cosas
\texorpdfstring{\softtext{diferentes.}}{diferentes.} Una
\textbf{variable del negocio} representa una medida o importe: bruto,
tara, precio, saldo o costo. Una \textbf{variable del programa} es un
nombre que referencia un objeto durante la ejecución, como el resultado
de una consulta o el
\texorpdfstring{\softtext{AbortController}}{AbortController} de una
petición. Una \textbf{variable de configuración} selecciona una
dirección, puerto o recurso antes de arrancar, como
\texorpdfstring{\softcode{BOVI\_DB\_A\_URL}}{BOVI\_DB\_A\_URL}. Ninguna
de esas tres es intercambiable con las otras.

La columna \texorpdfstring{\softcode{importe}}{importe} de un asiento
conserva un hecho \texorpdfstring{\softtext{persistente.}}{persistente.}
La variable local \texorpdfstring{\softcode{result}}{result} de Java
puede contener el JSON de ese hecho. El estado React con ese JSON es una
copia para mostrarlo. Si el operador cambia la presentación a formato
monetario con separadores, no está modificando la columna. La autoridad
económica sigue en \texorpdfstring{\softtext{PostgreSQL.}}{PostgreSQL.}

\section{Notación y unidades}\label{notaciuxf3n-y-unidades}

En las fórmulas de este manual, \(B\) representa peso bruto, \(T\) tara
y \(N\) peso neto, todos en
\texorpdfstring{\softtext{kilogramos.}}{kilogramos.} \(p\) es precio por
kilogramo, con unidad MXN/kg. \(n\) es el número entero de animales y
\(p_h\) precio por cabeza. \(V\) es venta neta reconocida, \(C\) costo
de posiciones vendidas y \(G_v\) gasto directo de venta. \(M\) es
margen. Nunca debe sumarse una cantidad de kg a un importe en MXN.

\[
N=B-T,\qquad 0\leq T\leq B.
\]

Por ejemplo, \texorpdfstring{\softtext{bruto5200.000}}{bruto5200.000} kg
y \texorpdfstring{\softtext{tara200.000}}{tara200.000} kg producen
\texorpdfstring{\softtext{neto5000.000}}{neto5000.000} kg. Una tara
mayor que el bruto es inválida, no un neto negativo que se arregla
tomando su valor absoluto. Un peso individual desconocido se representa
como ausencia, no como cero medido.

\[
\Delta W=W_i-W_f,\qquad
\Delta W_{\%}=100\frac{W_i-W_f}{W_i},\quad W_i>0.
\]

Los símbolos \(W_i\) y \(W_f\) representan pesos de conjuntos
comparables, inicial y final. Si el peso final es mayor, la variación es
negativa y se conserva: no se convierte en merma cero. Si los animales
de los dos tickets difieren, la aritmética puede producir un número pero
el resultado carece de significado para esa
\texorpdfstring{\softtext{comparación.}}{comparación.} Primero se valida
la composición exacta.

\section{Precio y reconocimiento}\label{precio-y-reconocimiento}

\[
\mathrm{Base}_{kg}=N\,p,\qquad
\mathrm{Base}_{cabeza}=n\,p_h.
\]

En la primera expresión kg multiplicado por MXN/kg produce MXN; en la
segunda, cabezas multiplicadas por MXN/cabeza también produce MXN. Esa
comprobación dimensional ayuda a detectar una condición aplicada al lote
\texorpdfstring{\softtext{equivocado.}}{equivocado.} Las condiciones
comerciales aceptadas determinan la base, no una referencia externa de
mercado.

\[
L=\operatorname{Redondear}_{2}
\left(\mathrm{Base}+\mathrm{Cargos}-\mathrm{Bonificaciones}\right).
\]

\(L\) representa el total de la obligación según las reglas de
\texorpdfstring{\softtext{materialización.}}{materialización.} Los
conceptos pueden materializarse y redondearse individualmente conforme
al SQL canónico; esta expresión resume su significado, no autoriza una
segunda implementación en Java o React. En particular, mostrar una
previsualización no registra
\texorpdfstring{\softtext{reconocimiento.}}{reconocimiento.} La
confirmación persistente produce el hecho económico una sola vez.

\section{Saldo, disponible y margen}\label{saldo-disponible-y-margen}

\[
S=L+\sum_j A^{cargo}_j-\sum_k A^{abono}_k-\sum_m P^{neto}_m.
\]

\(S\) es saldo de una
\texorpdfstring{\softtext{liquidación.}}{liquidación.} \(A^{cargo}\)
incrementa la obligación y \(A^{abono}\) la reduce. \(P^{neto}\) incluye
aplicaciones y sus reversos pertinentes; sumar una aplicación histórica
revertida como si todavía existiera reduciría el saldo dos veces. Un
pago registrado pero no validado no se puede aplicar para eludir esta
regla.

\[
D=I_{validado}-\sum_m P^{neto}_m-\sum_r R_r.
\]

\(D\) es dinero disponible de un pago validado; \(I_{validado}\) es su
importe y \(R\) son reembolsos
\texorpdfstring{\softtext{vinculados.}}{vinculados.} Disponible y saldo
son dos límites distintos. Una aplicación sólo procede cuando respeta
ambos y los ámbitos de empresa, sede, contraparte, moneda y sentido. Un
disponible de200 no significa otra venta de200.

\[
M=V-C-G_v.
\]

En GOLD-02, \(V=6400\), \(C=4400\) y \(G_v=200\), por lo que \(M=1800\).
Cobrar únicamente1200 no convierte el margen en un resultado de caja
de1200 menos costos: margen y dinero cobrado son ejes
\texorpdfstring{\softtext{diferentes.}}{diferentes.} Un cierre o reporte
debe explicar qué hechos y periodo utiliza.

\section{Precisión decimal y
redondeo}\label{precisiuxf3n-decimal-y-redondeo}

PostgreSQL utiliza NUMERIC, que almacena valores decimales de precisión
fijada. \texorpdfstring{\softcode{NUMERIC(18,2)}}{NUMERIC(18,2)}
permite18 dígitos en total y2 decimales; tiene16 dígitos enteros. El
límite finito positivo es \(10^{16}-10^{-2}\).
\texorpdfstring{\softcode{NUMERIC(14,3)}}{NUMERIC(14,3)} para peso
permite11 dígitos enteros y tres decimales. El dominio incorpora cotas y
validación de no finitos: disponer de un tipo numérico no reemplaza las
reglas de entrada.

Java interpreta entradas monetarias con herramientas decimales cuando
valida, mientras los importes confirmados provienen del SQL. El JSON
utiliza cadenas para no introducir una representación binaria
\texorpdfstring{\softtext{aproximada.}}{aproximada.} Una cadena
\texorpdfstring{\softcode{"1000.00"}}{"1000.00"} conserva significado
monetario; \texorpdfstring{\softcode{"400.000"}}{"400.000"} conserva
escala de kg. TypeScript puede dar formato visual, pero no recomputa la
liquidación con \texorpdfstring{\softcode{parseFloat}}{parseFloat}.

Para \(x\geq0\), una descripción matemática del redondeo HALF\_UP a
centavos es:

\[
\operatorname{Redondear}_{2}(x)=\frac{\lfloor100x+0.5\rfloor}{100}.
\]

\texorpdfstring{\softtext{Así,12.345}}{Así,12.345} se
\texorpdfstring{\softtext{materializa12.35}}{materializa12.35} y12.344
se \texorpdfstring{\softtext{materializa12.34.}}{materializa12.34.} Esta
fórmula ilustrativa se limita a no negativos. Los efectos contrarios se
representan con clases de asiento o registros de reversión según el
dominio, no introduciendo libremente un importe negativo para saltar
\texorpdfstring{\softtext{restricciones.}}{restricciones.}

\section{Distribución y conservación del
total}\label{distribuciuxf3n-y-conservaciuxf3n-del-total}

Un gasto total \(G\) se asigna a \(n\) posiciones por pesos relativos
\(w_i\) válidos. Por cabeza se toma \(w_i=1\). Por peso individual se
usa una medición individual existente, nunca el promedio de un ticket
grupal. Para importe individual la entrada declara cada importe y debe
conciliar \texorpdfstring{\softtext{directamente.}}{directamente.}

\[
a_i^{exacto}=G\frac{w_i}{\sum_{j=1}^{n}w_j},\qquad
r=G-\sum_{i=1}^{n}\operatorname{Redondear}_{2}(a_i^{exacto}).
\]

\(r\) es el residuo después de redondear. La distribución corrige
centavos en orden estable de posición, sin inventar o perder dinero.
GOLD-03 exige que1.00 entre tres posiciones iguales
\texorpdfstring{\softtext{termine0.34,0.33}}{termine0.34,0.33} y0.33. Si
tres importes redondeados inicialmente dan0.99, queda un centavo que se
asigna a la primera posición del orden. Si la suma inicial excede el
total, la política corrige el residuo negativo respetando la validez de
cada \texorpdfstring{\softtext{asignación.}}{asignación.}

\[
\sum_{i=1}^{n}a_i^{final}=G.
\]

Esta igualdad es una invariante: debe permanecer verdadera
independientemente de cómo se presente la tabla. La implementación
concreta y las validaciones de costo no negativo se explican en el
capítulo SQL y se contrastan con esperados fijos.

\section{Tiempo y resultado
histórico}\label{tiempo-y-resultado-histuxf3rico}

La fecha del documento identifica cuándo ocurrió o se fechó la
evidencia. La fecha de reconocimiento determina el periodo gerencial del
hecho. El timestamp de registro indica cuándo el servidor lo grabó. Un
flete de septiembre registrado y reconocido en octubre conserva las tres
nociones; no se modifica silenciosamente la fecha para reabrir un cierre
pasado.

Si un costo compite con el cierre, ambos bloquean la misma fila del
periodo. El costo entra antes del corte o el cierre gana y el costo se
rechaza sin efecto en ese periodo. Reconocerlo en otro periodo exige una
intención nueva \texorpdfstring{\softtext{explícita.}}{explícita.} El
reintento de red conserva los parámetros anteriores; no puede
convertirse de forma automática en un costo del mes siguiente.

\chapter{Diagramas de hechos y
estados}\label{diagramas-de-hechos-y-estados}

\section{Identidad y ciclos}\label{identidad-y-ciclos}

\begin{figure}[htbp]
\centering
\begin{tikzpicture}[node distance=12mm and 10mm,every node/.style={font=\small},box/.style={draw=brand,rounded corners,fill=surface,align=center,minimum width=32mm,minimum height=12mm},>=stealth]
\node[box] (a) {Animal\\identidad permanente};
\node[box,below left=of a] (c1) {Ciclo1\\costo propio};
\node[box,below right=of a] (c2) {Ciclo2\\costo propio};
\node[box,below=of c1] (v1) {Vendido y entregado\\CERRADA\_VENTA};
\node[box,below=of c2] (v2) {Recompra\\nueva adquisición};
\draw[->] (a)--(c1);\draw[->] (a)--(c2);\draw[->] (c1)--(v1);\draw[->] (c2)--(v2);
\draw[->,dashed] (v1)--node[below,font=\scriptsize]{habilita recompra válida}(v2);
\end{tikzpicture}
\caption{Un nuevo ciclo conserva el anterior. Nunca se habilitan dos posiciones activas para la misma identidad en el ámbito definido.}
\end{figure}

\section{Venta, pago y entrega son ejes
separados}\label{venta-pago-y-entrega-son-ejes-separados}

\begin{figure}[htbp]
\centering
\begin{tikzpicture}[node distance=11mm,every node/.style={font=\small},box/.style={draw=brand,fill=surface,align=center,minimum width=45mm,minimum height=10mm},>=stealth]
\node[box] (o) {Confirmar venta};
\node[box,below left=of o] (l) {Liquidación\\obligación};
\node[box,below right=of o] (c) {Costo de venta\\sale de inventario};
\node[box,below=of l] (p) {Cobro validado\\aplicación reduce saldo};
\node[box,below=of c] (e) {Entrega física\\cierra posición};
\draw[->] (o)--(l);\draw[->] (o)--(c);\draw[->] (l)--(p);\draw[->] (c)--(e);
\end{tikzpicture}
\caption{Pagar y entregar no vuelven a reconocer la venta. Cada hecho conserva su propia fecha y evidencia.}
\end{figure}

\section{Mensaje de muerte entre
dueños}\label{mensaje-de-muerte-entre-dueuxf1os}

\begin{figure}[htbp]
\centering
\begin{tikzpicture}[node distance=7mm,every node/.style={font=\small},box/.style={draw=brand,fill=surface,align=center,text width=110mm,minimum height=9mm},>=stealth]
\node[box] (b) {B: baja, pérdida, auditoría y outbox en una transacción};
\node[box,below=of b] (lease) {Worker B: reclama lease en transacción corta y termina};
\node[box,below=of lease] (net) {Llamada a A fuera de locks económicos de B};
\node[box,below=of net] (a) {A: inbox idempotente y actualización del maestro};
\node[box,below=of a] (ack) {B: registra acuse; la réplica distribuye identidad};
\draw[->] (b)--(lease);\draw[->] (lease)--(net);\draw[->] (net)--(a);\draw[->] (a)--(ack);
\end{tikzpicture}
\caption{Acuse perdido se recupera con inbox. No existe una transacción distribuida A/B ni una segunda pérdida.}
\end{figure}

\chapter{\texorpdfstring{La base de datos de
\texorpdfstring{\softtext{BoviCierre:}}{BoviCierre:} modelo, comandos y
reconocimiento}{La base de datos de  modelo, comandos y reconocimiento}}\label{la-base-de-datos-de-bovicierre-modelo-comandos-y-reconocimiento}

\section{1. Propósito y forma de estudiar el
sistema}\label{propuxf3sito-y-forma-de-estudiar-el-sistema}

Este capítulo explica la implementación de base de datos del checkout
real de \texorpdfstring{\softtext{BoviCierre.}}{BoviCierre.} Se estudian
las migraciones V1-V13, las funciones vigentes, el modelo lógico y el
catálogo físico observado. Su finalidad es didáctica: seguir una
identidad, una medición, una obligación y un importe a través de sus
transformaciones, entendiendo qué asegura cada restricción y qué
necesita una \texorpdfstring{\softtext{transacción.}}{transacción.} No
se presenta una lista de sentencias para ejecutar indiscriminadamente
sobre la base principal. Las consultas ilustrativas de catálogo son de
lectura; los ejemplos de escritura representan escenarios sintéticos que
deben ensayarse únicamente en QA propia.

Hay cinco conceptos que deben distinguirse desde el primer ejercicio:
identidad, intención, evidencia, reconocimiento y dinero. El animal
tiene una identidad que permanece; seleccionar animales para una compra
expresa una intención; aprobar un ticket aporta evidencia; confirmar la
compra reconoce una obligación y un ciclo de inventario; pagar y aplicar
dinero modifica caja y saldo. Recepción y entrega agregan hechos
físicos. Cada paso explica una parte del trato, pero no debe reconocer
otra compra o venta simplemente porque el usuario avanzó a otra
pantalla.

El modelo evita una tabla única de animales que mezcle su nombre, el
precio actual, la deuda pendiente, el pago y la fecha de entrega. Esa
mezcla parecería sencilla en una demostración pequeña, pero impediría
conservar un ciclo anterior después de recompra y explicaría mal una
venta pagada que todavía no se entrega. En
\texorpdfstring{\softtext{BoviCierre,}}{BoviCierre,} cada eje tiene
relaciones propias. Las consultas construyen una explicación conjunta,
mientras las escrituras mantienen la identidad y la historia de cada
hecho.

Las fuentes se citan con ruta y línea de inicio. Una función puede
continuar muchas líneas después, por lo que la referencia orienta al
bloque y no sustituye su lectura completa. Los valores GOLD son
ficticios y fijos. No se incluyen contraseñas, cookies, conninfo, hashes
de usuarios reales, contenido de sesiones ni datos comerciales privados.
Los tiempos de pruebas se atribuyen a sus artefactos; no se vuelven a
ejecutar aquí para producir una narración más cómoda.

\section{2. Qué DDL manda y por qué existen dos
recuentos}\label{quuxe9-ddl-manda-y-por-quuxe9-existen-dos-recuentos}

La autoridad física reside en
\texorpdfstring{\softtext{db/migrations.}}{db/migrations.} V1 crea
tablas, claves y parte del soporte; las migraciones siguientes añaden o
sustituyen objetos de manera
\texorpdfstring{\softtext{controlada.}}{controlada.}
\texorpdfstring{\softtext{tools/esquema/esquema.py}}{tools/esquema/esquema.py}
es la representación lógica para diagramas y
\texorpdfstring{\softtext{comparación.}}{comparación.} No debe
transformarse automáticamente en DDL sin revisar
\texorpdfstring{\softtext{semántica.}}{semántica.} Un diagrama o una
lista de columnas no contienen necesariamente todos los CHECK, los
índices parciales, los triggers, las columnas generadas o los
\texorpdfstring{\softtext{privilegios.}}{privilegios.}

Las consultas de sólo lectura realizadas en A y B el 12 de octubre de
2026 confirmaron39 tablas ordinarias en bovi y7 en infra:46 tablas
ordinarias y86FK en esos dos esquemas. La representación lógica
documenta42 tablas y83FK. No se trata de una contradicción ni de tres
relaciones artificiales para llegar a un número esperado. Las cuatro
tablas físicas adicionales son
\texorpdfstring{\softtext{infra.tarifa\_snapshot,}}{infra.tarifa\_snapshot,}
\texorpdfstring{\softtext{infra.spring\_session,}}{infra.spring\_session,}
\texorpdfstring{\softtext{infra.spring\_session\_attributes}}{infra.spring\_session\_attributes}
e
\texorpdfstring{\softtext{infra.flyway\_schema\_history.}}{infra.flyway\_schema\_history.}
La captura agrega dos FK a usuario y sucursal; los atributos agregan una
FK hacia su sesión. Así,42 más4 da46 y83 más3 da86.

Ese alcance excluye vistas, catálogos internos de
\texorpdfstring{\softtext{PostgreSQL,}}{PostgreSQL,} las tablas
extranjeras del esquema
\texorpdfstring{\softtext{ext\_proveedor}}{ext\_proveedor} y las tablas
ordinarias del servidor proveedor. Tampoco incorpora proyectos ajenos
del equipo. El modelo lógico de42 ya incluye soportes técnicos como
\texorpdfstring{\softtext{folio\_contador,}}{folio\_contador,} inbox e
idempotencia; por tanto,42 no significa42 entidades exclusivamente
ganaderas. El MER conceptual y el catálogo relacional persiguen
objetivos distintos y deben defenderse como tales.

El conteo no demuestra que todas las45RN estén probadas en todas las
\texorpdfstring{\softtext{combinaciones.}}{combinaciones.} Demuestra que
un catálogo con un alcance explícito corresponde a la implementación
observada. La comparación lógica contra A/B verifica columnas, tipos,
nulabilidad y claves declaradas; las pruebas de permisos, triggers,
carreras y oráculos económicos verifican mecanismos que el conteo no
expresa.

\textbf{Fuentes:}
\texorpdfstring{\softtext{tools/esquema/esquema.py:1;}}{tools/esquema/esquema.py:1;}
\texorpdfstring{\softtext{db/migrations/V1\_\_esquema.sql:1;}}{db/migrations/V1\_\_esquema.sql:1;}
\texorpdfstring{\softtext{db/README.md:3;}}{db/README.md:3;}
\texorpdfstring{\softtext{docs/database/relational-model.md:1.}}{docs/database/relational-model.md:1.}

\subsection{2.1 La historia de migración no debe
inventarse}\label{la-historia-de-migraciuxf3n-no-debe-inventarse}

En la instancia principal existente, V1-V3 fueron adoptadas mediante
baseline3 después de comparar su catálogo con una reproducción limpia.
El historial declara adopción auditada de catálogo, no ejecución
original de esos archivos por Flyway. V4-V13 sí aparecen como
migraciones \texorpdfstring{\softtext{ejecutadas.}}{ejecutadas.} Una
base completamente vacía aplica la secuencia V1-V13. Ambas historias
pueden terminar en una estructura equivalente, pero no son la misma
historia operativa.

El backend no lleva Flyway para crear objetos al iniciar una petición o
al arrancar el servidor. Un trabajo
\texorpdfstring{\softtext{Java21/Flyway}}{Java21/Flyway} separado aplica
y valida migraciones con una identidad
\texorpdfstring{\softtext{administrativa.}}{administrativa.} Las cuentas
de negocio no necesitan DDL. Esta separación impide que una credencial
usada por una solicitud HTTP modifique triggers o cree una tabla para
eludir una \texorpdfstring{\softtext{validación.}}{validación.}

CREATE OR REPLACE FUNCTION reemplaza la definición de una función
existente con la misma firma. No significa que V2, V4, V8 y V13 sean
cuatro algoritmos monetarios activos
\texorpdfstring{\softtext{simultáneamente.}}{simultáneamente.} La última
definición de
\texorpdfstring{\softtext{bovi.cmd\_ejecutar}}{bovi.cmd\_ejecutar} es
V13. La consulta
\texorpdfstring{\softtext{bovi.consultar}}{bovi.consultar} y los
auxiliares de pesaje y reporte vigentes son V12. Cálculo, reparto,
normalización y validación de periodos conservan auxiliares definidos en
V2. Lectura FDW y captura local quedan según V9.

{\def\LTcaptype{none} % do not increment counter
\begin{longtable}[]{@{}
  >{\raggedright\arraybackslash}p{(\linewidth - 4\tabcolsep) * \real{0.3333}}
  >{\raggedright\arraybackslash}p{(\linewidth - 4\tabcolsep) * \real{0.3333}}
  >{\raggedright\arraybackslash}p{(\linewidth - 4\tabcolsep) * \real{0.3333}}@{}}
\toprule\noalign{}
\begin{minipage}[b]{\linewidth}\raggedright
Migración
\end{minipage} & \begin{minipage}[b]{\linewidth}\raggedright
Aporte
\end{minipage} & \begin{minipage}[b]{\linewidth}\raggedright
Razón que debe explicarse
\end{minipage} \\
\midrule\noalign{}
\endhead
\bottomrule\noalign{}
\endlastfoot
V1 & Estructura, claves, restricciones e índices & Establece integridad
de filas y \texorpdfstring{\softtext{relaciones.}}{relaciones.} \\
V2 & Comandos, cálculo, reparto, periodos y triggers & Define
reconocimiento atómico y
\texorpdfstring{\softtext{auxiliares.}}{auxiliares.} \\
V3 & Consultas, sesiones, grants y trabajador & Separa interfaces
técnicas y de negocio. \\
V4 & Captura de proveedor y DTO completos & Materializa procedencia y
respuestas \texorpdfstring{\softtext{persistibles.}}{persistibles.} \\
V5 & Corte por periodo de reconocimiento & Evita que documento posterior
altere un cierre. \\
V6 & Dinero a centavos y paginación & Conserva precisión original sin
romper DTO monetario. \\
V7 & Filtros de consulta & Evita búsqueda ignorada y pagina con orden
estable. \\
V8 & Borrador KG incompleto y reloj real & Capturar pendiente no
autoriza confirmar sin medición. \\
V9 & Lector A, captura dueña y fecha PATCH & Saca red de locks y
preserva fecha omitida. \\
V10 & Posiciones disponibles para venta & Estado solo no basta si hay
reserva o baja. \\
V11 & Guarda de versión al
\texorpdfstring{\softtext{cancelar/anular}}{cancelar/anular} & Un lock
no detecta por sí solo intención obsoleta. \\
V12 & Tickets y reportes canónicos & Conserva composición propia y
margen \texorpdfstring{\softtext{histórico.}}{histórico.} \\
V13 & Igualdad de conjunto en corrección & Rechaza
\texorpdfstring{\softtext{omisiones/adiciones}}{omisiones/adiciones}
antes de tocar original. \\
\end{longtable}
}

\section{3. Leer PK, FK, NN y UK como mecanismos
distintos}\label{leer-pk-fk-nn-y-uk-como-mecanismos-distintos}

PK identifica la fila y obliga a que sus columnas no sean NULL. NN, NOT
NULL, exige un valor. FK obliga a que la combinación referenciada
exista. UK, UNIQUE, impide duplicados de una
\texorpdfstring{\softtext{combinación.}}{combinación.} Una UK que
contiene la PK y un atributo redundante puede servir como destino de una
FK que también comprueba tipo, pero no es necesariamente una clave
candidata mínima. La minimalidad exige que no sobre ninguna columna.

En los hechos locales se utiliza con frecuencia la identidad compuesta
E,S,id: \texorpdfstring{\softtext{empresa\_id,}}{empresa\_id,}
\texorpdfstring{\softtext{sucursal\_id}}{sucursal\_id} e
\texorpdfstring{\softtext{identificador.}}{identificador.} Una FK local
repite empresa y sede además del UUID. Esto evita que un pago de una
sede encuentre una liquidación de otra simplemente porque su UUID
existe. Un UUID difícil de adivinar tampoco concede autorización:
estructura y permisos son garantías distintas.

Las entidades dependientes usan claves parciales.
\texorpdfstring{\softtext{numero\_lote}}{numero\_lote} identifica un
lote dentro de una operación;
\texorpdfstring{\softtext{numero\_version}}{numero\_version} identifica
una propuesta dentro de un contrato o una versión dentro de un periodo.
El mismo número puede existir en otro
\texorpdfstring{\softtext{propietario.}}{propietario.} Presentar sólo
lote1 sin operación es perder una parte esencial de la identidad.

Una FK opcional normalmente utiliza MATCH SIMPLE. E y S pueden ser
obligatorios mientras
\texorpdfstring{\softtext{documento\_id}}{documento\_id} es NULL; al
faltar el identificador opcional, la referencia se omite. Imponer MATCH
FULL indiscriminadamente exigiría todos nulos o todos completos y
rompería ese caso normal. Cuando el negocio exige exactamente una
fuente, se aplica un CHECK XOR, como en
\texorpdfstring{\softtext{asignacion\_costo,}}{asignacion\_costo,} no
una nulabilidad \texorpdfstring{\softtext{improvisada.}}{improvisada.}

Los índices únicos parciales aplican a un subconjunto, por ejemplo
posiciones activas. Los CHECK examinan valores de una fila. Los triggers
controlan estados de padres o
\texorpdfstring{\softtext{transiciones.}}{transiciones.} Los comandos
con FOR UPDATE protegen condiciones agregadas
\texorpdfstring{\softtext{concurrentes.}}{concurrentes.} No se
intercambian sin perder garantías: un CHECK de la aplicación nueva no
puede verificar por sí solo la suma de otras aplicaciones
\texorpdfstring{\softtext{concurrentes.}}{concurrentes.}

Las columnas obligatorias no vuelven obligatorio un dato distinto del
que \texorpdfstring{\softtext{representan.}}{representan.} Una fecha de
creación NN no sustituye una fecha de reconocimiento elegida. Un actor
NN no autoriza a utilizar un UUID inventado. Un documento opcional puede
seguir siendo privado. En el diccionario íntegro de anexos deben leerse
las claves completas y acciones referenciales, además de las marcas por
columna.

\section{4. Mapa de las42 tablas lógicas y las cuatro técnicas
adicionales}\label{mapa-de-las42-tablas-luxf3gicas-y-las-cuatro-tuxe9cnicas-adicionales}

Esta sección explica por qué existe cada tabla y qué columnas deben
reconocerse al recorrer el sistema. El anexo de diccionario conserva
todas las columnas y declaraciones; aquí se sigue su función. Se utiliza
E para empresa\_id y S para
\texorpdfstring{\softtext{sucursal\_id.}}{sucursal\_id.} Los nombres de
módulos del modelo lógico no son esquemas físicos separados: las tablas
de negocio están en bovi y el soporte explícito indicado en infra.

\subsection{4.1 Organización, acceso y maestros
centrales}\label{organizaciuxf3n-acceso-y-maestros-centrales}

\textbf{empresa.} Su PK es id;
\texorpdfstring{\softtext{razon\_social,}}{razon\_social,} estado,
\texorpdfstring{\softtext{nodo\_origen}}{nodo\_origen} y creado\_en
completan organización y autoría, mientras rfc admite ausencia
declarada. Es raíz del scope, no una etiqueta que pueda cambiarse para
habilitar datos de otra empresa. La instalación configura una empresa;
disponer de claves empresariales no acredita RLS multiempresa comercial.

\textbf{sucursal.} PK E,id y UK E,codigo distinguen identidad técnica de
código humano. nodo\_dueno fija A o B;
\texorpdfstring{\softtext{zona\_horaria}}{zona\_horaria} fija contexto
de \texorpdfstring{\softtext{calendario.}}{calendario.} La FK empresa
conecta sede con
\texorpdfstring{\softtext{organización.}}{organización.} Es maestro
escrito en A aunque describa una sede cuyo dueño financiero sea B. Dueño
no significa custodio físico del ganado.

\textbf{\texorpdfstring{\softtext{rol\_aplicacion.}}{rol\_aplicacion.}}
PK codigo y descripcion forman catálogo de roles humanos. ADMIN,
OPERADOR, CAJA y AUDITOR no equivalen a cuentas
\texorpdfstring{\softtext{PostgreSQL.}}{PostgreSQL.} La descripción
permite comunicar responsabilidades; los permisos se materializan cuando
una membresía vigente conecta usuario y sede con ese código.

\textbf{usuario.} PK E,id y UK
\texorpdfstring{\softtext{E,login\_normalizado}}{E,login\_normalizado}
mantienen identidad de acceso. nombre y estado explican presentación y
vigencia. \texorpdfstring{\softtext{hash\_contrasena}}{hash\_contrasena}
es privado y contiene bcrypt, no texto legible. La autenticación tiene
cuenta técnica separada para leerlo. Desactivar usuario limita
solicitudes nuevas, pero no cambia actor de auditorías
\texorpdfstring{\softtext{históricas.}}{históricas.}

\textbf{\texorpdfstring{\softtext{membresia\_usuario.}}{membresia\_usuario.}}
PK
\texorpdfstring{\softtext{E,S,usuario\_id,rol\_codigo}}{E,S,usuario\_id,rol\_codigo}
expresa una relación ternaria. Sus tres FK conectan sede, persona y
catálogo de roles. vigente permite retirar autorización sin borrar
usuario. La misma persona puede tener varios roles o permisos diferentes
entre sedes. El comando de asignación exige ADMIN y no permite elegir
otra S en el cuerpo.

\textbf{parte.} PK E,id identifica
\texorpdfstring{\softtext{contraparte.}}{contraparte.}
\texorpdfstring{\softtext{tipo\_persona,}}{tipo\_persona,} nombre,
telefono y \texorpdfstring{\softtext{es\_sintetico}}{es\_sintetico}
ayudan a explicar declaración y DEMO. El RFC tiene unicidad parcial
cuando existe. Parte no se divide en proveedor y cliente excluyentes:
puede vender y comprar. Una liquidación conserva snapshots para que un
renombre posterior no reescriba el
\texorpdfstring{\softtext{comprobante.}}{comprobante.}

\textbf{predio.} PK E,id identifica lugar declarado. nombre es
obligatorio;
\texorpdfstring{\softtext{clave\_upp\_declarada,}}{clave\_upp\_declarada,}
municipio,
\texorpdfstring{\softtext{entidad\_federativa}}{entidad\_federativa} y
\texorpdfstring{\softtext{parte\_responsable\_id}}{parte\_responsable\_id}
aportan información opcional. La FK al responsable conecta declaración
comercial. Capturar una UPP no demuestra consulta o certificación
oficial; el término declarada evita atribuir una validación que no
existe.

\textbf{animal.} PK E,id mantiene identidad permanente y UK
\texorpdfstring{\softtext{E,sede\_operadora\_id,id}}{E,sede\_operadora\_id,id}
permite comprobar pertenencia operativa. sexo y situacion son datos
obligatorios; raza y fecha de nacimiento son
\texorpdfstring{\softtext{declaradas.}}{declaradas.} La FK hacia
sucursal fija sede \texorpdfstring{\softtext{autorizada.}}{autorizada.}
VIGENTE,
\texorpdfstring{\softtext{IDENTIDAD\_EN\_REVISION}}{IDENTIDAD\_EN\_REVISION}
y MUERTO tienen efectos distintos sobre nuevos
\texorpdfstring{\softtext{compromisos.}}{compromisos.}

\textbf{\texorpdfstring{\softtext{animal\_identificacion.}}{animal\_identificacion.}}
PK E,id y UK
\texorpdfstring{\softtext{E,sistema,codigo\_normalizado}}{E,sistema,codigo\_normalizado}
resuelven \texorpdfstring{\softtext{colisiones.}}{colisiones.}
\texorpdfstring{\softtext{codigo\_declarado}}{codigo\_declarado}
conserva lo escrito; el normalizado facilita
\texorpdfstring{\softtext{comparación.}}{comparación.} Las FK enlazan
animal y evidencia documental opcional con su sede. La normalización no
acepta símbolos arbitrarios: primero valida y luego elimina
\texorpdfstring{\softtext{separadores.}}{separadores.} Escritura central
A hace competir las solicitudes en una misma unicidad.

\subsection{4.2 Intenciones comerciales y condiciones
versionadas}\label{intenciones-comerciales-y-condiciones-versionadas}

\textbf{operacion.} PK E,S,id, UK de tipo/folio y una UK adicional
id/tipo describen el agregado. parte\_id y referencias opcionales a
predios conectan contexto. tipo, moneda,
\texorpdfstring{\softtext{fecha\_operacion,}}{fecha\_operacion,} estado
y \texorpdfstring{\softtext{estado\_entrega}}{estado\_entrega} no
mezclan dinero con entrega.
\texorpdfstring{\softtext{version\_lock}}{version\_lock} controla
concurrencia;
\texorpdfstring{\softtext{numero\_version\_aceptada}}{numero\_version\_aceptada}
enlaza la versión contractual mediante FK diferible.

\textbf{\texorpdfstring{\softtext{folio\_contador.}}{folio\_contador.}}
PK E,S,tipo,anio identifica un contador interno. ultimo cambia mediante
UPSERT \texorpdfstring{\softtext{transaccional.}}{transaccional.} Su FK
hacia sede mantiene contexto. Este soporte aparece en el modelo lógico
aunque no sea entidad ganadera. Un folio puede tener huecos y no
constituye folio fiscal
\texorpdfstring{\softtext{autorizado.}}{autorizado.}

\textbf{\texorpdfstring{\softtext{lote\_operacion.}}{lote\_operacion.}}
PK
\texorpdfstring{\softtext{E,S,operacion\_id,numero\_lote}}{E,S,operacion\_id,numero\_lote}
expresa \texorpdfstring{\softtext{dependencia.}}{dependencia.}
descripcion es una etiqueta opcional; FK identifica la
\texorpdfstring{\softtext{operación.}}{operación.} El lote agrupa
selección y \texorpdfstring{\softtext{condiciones.}}{condiciones.}
Cambiar su ordinalidad cuando existe historia puede aplicar el precio de
un grupo a otro, por lo que identidad de lote y composición deben
conservarse al editar.

\textbf{\texorpdfstring{\softtext{contrato\_version.}}{contrato\_version.}}
PK
\texorpdfstring{\softtext{E,S,operacion\_id,numero\_version}}{E,S,operacion\_id,numero\_version}
distingue \texorpdfstring{\softtext{propuestas.}}{propuestas.} estado,
\texorpdfstring{\softtext{politica\_redondeo}}{politica\_redondeo} y
\texorpdfstring{\softtext{requiere\_recepcion}}{requiere\_recepcion}
definen decisión contractual;
\texorpdfstring{\softtext{documento\_id}}{documento\_id} es evidencia
opcional. \texorpdfstring{\softtext{aceptada\_por}}{aceptada\_por} y
\texorpdfstring{\softtext{aceptada\_en}}{aceptada\_en} son nulos antes
de aceptar. La FK a operación y la documental conservan contexto.
ACEPTADA se protege contra edición o borrado.

\textbf{\texorpdfstring{\softtext{condicion\_lote.}}{condicion\_lote.}}
Su PK combina operación, versión y lote. FK hacia contrato, lote y
pesaje unen precio y evidencia.
\texorpdfstring{\softtext{unidad\_precio,}}{unidad\_precio,}
\texorpdfstring{\softtext{precio\_unitario,}}{precio\_unitario,}
\texorpdfstring{\softtext{deduccion\_pct,}}{deduccion\_pct,}
\texorpdfstring{\softtext{cargos\_pactados}}{cargos\_pactados} y
\texorpdfstring{\softtext{bonificaciones\_pactadas}}{bonificaciones\_pactadas}
explican fórmula.
\texorpdfstring{\softtext{pesaje\_base\_id}}{pesaje\_base\_id} puede
faltar en borrador KG desde V8; las guardas de acuerdo y confirmación
exigen completarlo, no simulan un total cero.

\subsection{4.3 Selección, ciclos y compromisos
físicos}\label{selecciuxf3n-ciclos-y-compromisos-fuxedsicos}

\textbf{\texorpdfstring{\softtext{seleccion\_compra.}}{seleccion\_compra.}}
PK incorpora operación, lote y animal; UK adicional impide repetir
animal entre lotes del mismo trato. La FK con
\texorpdfstring{\softtext{tipo\_operacion}}{tipo\_operacion} comprueba
COMPRA y la FK hacia animal incluye sede operadora. Seleccionar aún no
crea inventario: representa intención cuya adquisición se reconoce al
confirmar.

\textbf{\texorpdfstring{\softtext{posicion\_inventario.}}{posicion\_inventario.}}
PK E,S,id y UK de \texorpdfstring{\softtext{animal/ciclo}}{animal/ciclo}
mantienen historia de
\texorpdfstring{\softtext{adquisición.}}{adquisición.} Otra clave
vincula compra, lote y animal. El índice parcial impide dos posiciones
activas por animal en empresa. estado distingue pendiente de recepción,
disponible, vendida pendiente de entrega, cerrada y baja. El costo es
derivado, no un saldo manual editable en la fila.

\textbf{\texorpdfstring{\softtext{seleccion\_venta.}}{seleccion\_venta.}}
PK incluye operación, lote y
\texorpdfstring{\softtext{posicion\_id.}}{posicion\_id.} La UK impide
repetir posición dentro de una venta y la FK de tipo exige VENTA. Se
vende un ciclo concreto, no una identidad animal ambigua. Una venta
puede reunir posiciones de varias compras sin perder el costo de origen
de cada una.

\textbf{\texorpdfstring{\softtext{reserva\_venta.}}{reserva\_venta.}} PK
E,S,id y FK a selección exacta vinculan compromiso con
\texorpdfstring{\softtext{venta/lote/posición.}}{venta/lote/posición.}
creada\_en, vence\_en, cerrada\_en y estado explican duración. Un índice
único parcial limita ACTIVA por posición. El índice no utiliza now:
vencer requiere
\texorpdfstring{\softtext{transición/comprobación}}{transición/comprobación}
\texorpdfstring{\softtext{explícita.}}{explícita.} Reservada se deriva,
no es otro estado persistido de posición.

\textbf{\texorpdfstring{\softtext{movimiento\_inventario.}}{movimiento\_inventario.}}
PK E,S,id y FK hacia posición, operación, predio, custodio, documento y
original de reverso conservan hecho físico.
\texorpdfstring{\softtext{ocurrido\_en}}{ocurrido\_en} expresa cuándo
ocurrió, no cuándo se reconoció una venta. Recepción y entrega no
generan deuda ni costo adicional. Baja crea pérdida mediante comando y
evidencia, conservando identidad y ciclo.

\subsection{4.4 Mediciones y expediente
documental}\label{mediciones-y-expediente-documental}

\textbf{pesaje.} PK E,S,id y FK hacia
\texorpdfstring{\softtext{lote/documento/original}}{lote/documento/original}
conservan ticket. bruto\_kg, tara\_kg y neto\_kg tienen precisión de
tres decimales; CHECK impide valores
\texorpdfstring{\softtext{incoherentes.}}{incoherentes.} modo, fuente,
bascula, \texorpdfstring{\softtext{ticket\_folio}}{ticket\_folio} y
pesado\_en explican
\texorpdfstring{\softtext{procedencia.}}{procedencia.} APROBADO congela;
CORREGIDO conserva original y otro ticket apunta a
\texorpdfstring{\softtext{corrige\_a\_id.}}{corrige\_a\_id.}

\textbf{\texorpdfstring{\softtext{pesaje\_animal.}}{pesaje\_animal.}} PK
\texorpdfstring{\softtext{E,S,pesaje\_id,animal\_id}}{E,S,pesaje\_id,animal\_id}
y dos FK materializan el conjunto medido.
\texorpdfstring{\softtext{peso\_individual\_kg}}{peso\_individual\_kg}
admite NULL porque una báscula grupal no midió necesariamente cada
individuo. Ese NULL no se rellena dividiendo el neto. La aprobación
congela composición y la corrección conserva conjunto exacto.

\textbf{documento.} PK E,S,id, unicidad de clave de almacenamiento y
enlace al documento previo conservan versiones. SHA256, tamaño y MIME
describen bytes privados almacenados fuera de BD. Fechas, emisor, folio,
origen y destino son datos
\texorpdfstring{\softtext{declarados.}}{declarados.} Una revisión de v1
no se transforma en aprobación de v2.

\textbf{\texorpdfstring{\softtext{operacion\_documento.}}{operacion\_documento.}}
PK
\texorpdfstring{\softtext{E,S,operacion\_id,documento\_id}}{E,S,operacion\_id,documento\_id}
evita relación duplicada. rol explica función de evidencia y ambas FK
mantienen scope. Asociar archivo no confirma importe, deuda ni
\texorpdfstring{\softtext{movimiento.}}{movimiento.} El servidor
controla descarga; conocer una clave interna no autoriza obtener los
bytes.

\textbf{\texorpdfstring{\softtext{revision\_administrativa.}}{revision\_administrativa.}}
PK E,S,id y CHECK XOR exigen operación o identificación, exactamente
una. criterio, resultado,
\texorpdfstring{\softtext{responsable\_id}}{responsable\_id} y
\texorpdfstring{\softtext{revisado\_en}}{revisado\_en} explican
revisión.
\texorpdfstring{\softtext{revision\_anterior\_id}}{revision\_anterior\_id}
conserva \texorpdfstring{\softtext{continuidad.}}{continuidad.} El
encabezado no acredita certificación sanitaria ni permite editar una
liquidación \texorpdfstring{\softtext{confirmada.}}{confirmada.}

\textbf{\texorpdfstring{\softtext{revision\_documento.}}{revision\_documento.}}
PK
\texorpdfstring{\softtext{E,S,revision\_id,documento\_id}}{E,S,revision\_id,documento\_id}
enlaza versión exacta. resultado y observacion guardan juicio sobre lo
revisado. Las dos FK sostienen asociación; una versión documental nueva
deja intacta esa \texorpdfstring{\softtext{referencia.}}{referencia.} Es
historia administrativa protegida, no una bandera mutable que siempre
apunta al último archivo.

\subsection{4.5 Obligaciones, costos y
calendario}\label{obligaciones-costos-y-calendario}

\textbf{\texorpdfstring{\softtext{periodo\_cierre.}}{periodo\_cierre.}}
PK E,S,id, FK a sede, fechas inclusivas, estado y
\texorpdfstring{\softtext{version\_lock}}{version\_lock} definen mes.
EXCLUDE GiST impide superposición; política complementaria mantiene un
abierto y meses \texorpdfstring{\softtext{consecutivos.}}{consecutivos.}
La fila es mutex común para cierre y reconocimiento, evitando que un
costo se cuele después del corte bajo fecha documental antigua.

\textbf{\texorpdfstring{\softtext{cierre\_periodo\_version.}}{cierre\_periodo\_version.}}
PK
\texorpdfstring{\softtext{E,S,periodo\_id,numero\_version}}{E,S,periodo\_id,numero\_version}
y FK al periodo conservan CIERRE o
\texorpdfstring{\softtext{REAPERTURA.}}{REAPERTURA.} actor\_id, motivo,
corte\_en,
\texorpdfstring{\softtext{hash\_conciliacion}}{hash\_conciliacion} y
totales explican decisión. V12 añade snapshot de
\texorpdfstring{\softtext{reporte\_margen}}{reporte\_margen} al JSON de
cierres nuevos. Reabrir no borra versiones previas ni inventa un
desglose inexistente de un cierre antiguo.

\textbf{\texorpdfstring{\softtext{liquidacion.}}{liquidacion.}} PK
E,S,id y UK por operación permiten una obligación principal. FK al
contrato aceptado y periodo conservan
\texorpdfstring{\softtext{reconocimiento.}}{reconocimiento.} total,
direccion y snapshots son inmutables; vence\_el puede ser NULL. El saldo
se deriva de ajustes y aplicaciones; generar un PDF o recibir dinero no
modifica el total original.

\textbf{\texorpdfstring{\softtext{concepto\_liquidacion.}}{concepto\_liquidacion.}}
PK
\texorpdfstring{\softtext{E,S,liquidacion\_id,numero\_concepto}}{E,S,liquidacion\_id,numero\_concepto}
identifica desglose
\texorpdfstring{\softtext{dependiente.}}{dependiente.} cantidad, unidad,
\texorpdfstring{\softtext{precio\_unitario}}{precio\_unitario} e importe
explican BASE, CARGO o
\texorpdfstring{\softtext{BONIFICACION.}}{BONIFICACION.} Lote y pesaje
aportan evidencia. Una bonificación pactada puede materializar importe
negativo según tipo; los signos de valoración continúan derivados, no
introducidos
\texorpdfstring{\softtext{arbitrariamente.}}{arbitrariamente.}

\textbf{\texorpdfstring{\softtext{ajuste\_liquidacion.}}{ajuste\_liquidacion.}}
PK E,S,id enlaza liquidación, documento, original de ajuste y periodo.
naturaleza \texorpdfstring{\softtext{CARGO/ABONO}}{CARGO/ABONO} da signo
a monto positivo. motivo, fuente y actor justifican
\texorpdfstring{\softtext{corrección.}}{corrección.} No reescribe
\texorpdfstring{\softtext{liquidacion.total.}}{liquidacion.total.} Una
reducción de deuda y una bonificación de costo son intenciones distintas
aunque sus cifras coincidan.

\textbf{\texorpdfstring{\softtext{gasto\_operacion.}}{gasto\_operacion.}}
PK E,S,id y FK a operación, proveedor, documento, gasto bonificado y
periodo mantienen causa. naturaleza y clasificacion separan
\texorpdfstring{\softtext{costo/bonificación}}{costo/bonificación} y
\texorpdfstring{\softtext{capitalizable/directo/otro.}}{capitalizable/directo/otro.}
La referencia de proveedor permite detectar duplicados
\texorpdfstring{\softtext{declarados.}}{declarados.}
\texorpdfstring{\softtext{snapshot\_fuente\_externa}}{snapshot\_fuente\_externa}
conserva tarifa aceptada; su importe no se vuelve a consultar en FDW
dentro de locks.

\textbf{\texorpdfstring{\softtext{asignacion\_costo.}}{asignacion\_costo.}}
PK E,S,id y FK hacia posición y fuentes permiten seguir cuotas. Un CHECK
XOR exige concepto, gasto o ajuste como única fuente. metodo,
base\_valor y
\texorpdfstring{\softtext{lote\_reparto\_id}}{lote\_reparto\_id}
explican \texorpdfstring{\softtext{distribución.}}{distribución.}
periodo\_id ancla efecto. La suma del reparto debe ser exactamente igual
al importe, conservando residuo en centavos
\texorpdfstring{\softtext{determinista.}}{determinista.}

\textbf{\texorpdfstring{\softtext{asiento\_valoracion.}}{asiento\_valoracion.}}
PK E,S,id enlaza posición, periodo, asignación y
\texorpdfstring{\softtext{operación.}}{operación.} clase, sentido e
importe generan
\texorpdfstring{\softtext{delta\_inventario,}}{delta\_inventario,}
\texorpdfstring{\softtext{delta\_costo\_venta}}{delta\_costo\_venta} y
\texorpdfstring{\softtext{delta\_perdida.}}{delta\_perdida.} Los deltas
no son campos libres del
\texorpdfstring{\softtext{formulario.}}{formulario.} Un costo tardío de
compra puede atribuirse a la venta original de la posición ya vendida;
una baja atribuye pérdida, no venta.

\subsection{4.6 Dinero, historia e
integración}\label{dinero-historia-e-integraciuxf3n}

\textbf{pago.} PK E,S,id conecta parte, origen de reembolso, documento y
periodo. tipo es \texorpdfstring{\softtext{COBRO/PAGO;}}{COBRO/PAGO;}
naturaleza
\texorpdfstring{\softtext{NORMAL/REEMBOLSO,}}{NORMAL/REEMBOLSO,} un eje
distinto. REGISTRADO mantiene periodo y
\texorpdfstring{\softtext{fecha\_reconocimiento}}{fecha\_reconocimiento}
nulos hasta validar. importe, método, referencia y
\texorpdfstring{\softtext{recibido\_el}}{recibido\_el} describen caja.
Disponible es derivado y puede quedar como anticipo sin fabricar ingreso
adicional.

\textbf{\texorpdfstring{\softtext{aplicacion\_pago.}}{aplicacion\_pago.}}
PK E,S,id y FK hacia pago, liquidación, original de reverso y periodo
modelan N:M con identidad. naturaleza
\texorpdfstring{\softtext{APLICACION/REVERSO}}{APLICACION/REVERSO} da
signo a importe positivo.
\texorpdfstring{\softtext{reverso\_de\_id}}{reverso\_de\_id} único
impide segundo contramovimiento del mismo original. FK no basta: el
comando verifica contraparte, moneda, dirección, saldo y disponible bajo
locks.

\textbf{\texorpdfstring{\softtext{auditoria\_evento.}}{auditoria\_evento.}}
PK E,S,id conserva actor, comando, correlación, entidad y cambio JSON.
El resultado se inserta junto al efecto, no después en otra
\texorpdfstring{\softtext{transacción.}}{transacción.} Desde V13 el
motivo de corrección puede recuperarse del resultado auditado del ticket
nuevo. Los logs técnicos no deben copiar secretos por el hecho de
existir auditoría interna.

\textbf{\texorpdfstring{\softtext{evento\_outbox.}}{evento\_outbox.}} PK
E,S,id guarda tipo, versión de payload, datos y propietario de despacho.
creado\_en no demuestra entrega. El evento nace con el hecho económico o
físico y es inmutable. La tabla de entrega separa reintentos locales. La
muerte en B puede actualizar A después, sin retener red dentro de la
transacción de pérdida.

\textbf{\texorpdfstring{\softtext{mensaje\_inbox.}}{mensaje\_inbox.}} PK
\texorpdfstring{\softtext{consumidor,fuente\_id,event\_id}}{consumidor,fuente\_id,event\_id}
deduplica consumo local. resultado y
\texorpdfstring{\softtext{recibido\_en}}{recibido\_en} conservan
respuesta. No se publica como negocio replicado. Consumir muerte central
no reconoce otra pérdida: devuelve resultado si el evento ya fue
aplicado y permite responder a un acuse perdido.

\textbf{\texorpdfstring{\softtext{clave\_idempotencia.}}{clave\_idempotencia.}}
PK
\texorpdfstring{\softtext{E,S,principal\_id,comando,clave}}{E,S,principal\_id,comando,clave}
delimita \texorpdfstring{\softtext{intención.}}{intención.}
\texorpdfstring{\softtext{hash\_peticion,}}{hash\_peticion,} estado y
resultado distinguen replay de conflicto. expira\_en es infinity en el
comando actual. La persistencia es permanente para deduplicación
financiera, no caché de sesión. Se consulta sólo por el principal
correspondiente y después de autorización vigente.

\textbf{\texorpdfstring{\softtext{infra.nodo\_config.}}{infra.nodo\_config.}}
PK nodo y empresa\_id definen identidad técnica local. No se replica: A
y B deben seguir reconociéndose como sí mismos. La app no modifica este
soporte para adjudicarse otra sede. Una restauración aislada debe
reconstruir conexiones, no reutilizar suscripciones hacia el proyecto
original.

\textbf{\texorpdfstring{\softtext{infra.entrega\_outbox.}}{infra.entrega\_outbox.}}
PK \texorpdfstring{\softtext{evento\_id,destino}}{evento\_id,destino}
conserva intentos, próximo intento, acuse, error y lease. Los nombres
físicos son \texorpdfstring{\softtext{lease\_token}}{lease\_token} y
\texorpdfstring{\softtext{lease\_vence\_en.}}{lease\_vence\_en.} No
tiene FK declarada al evento en este catálogo: funciones validan su
relación. La lease durable termina antes de red y un acuse antiguo no
puede confirmar trabajo reclamado por otro.

\textbf{\texorpdfstring{\softtext{infra.importacion\_fuente.}}{infra.importacion\_fuente.}}
PK id y datos de
\texorpdfstring{\softtext{empresa/sede/fuente/hash/corte/conteos}}{empresa/sede/fuente/hash/corte/conteos}
documentan \texorpdfstring{\softtext{procedencia.}}{procedencia.}
\texorpdfstring{\softtext{CSV/OCR/FDW}}{CSV/OCR/FDW} son vocabulario, no
prueba de OCR \texorpdfstring{\softtext{implementado.}}{implementado.}
Es soporte local; no atribuirle FK ausentes. Importar una propuesta no
autoriza confirmar automáticamente importes.

\subsection{4.7 Las cuatro tablas físicas
adicionales}\label{las-cuatro-tablas-fuxedsicas-adicionales}

\textbf{\texorpdfstring{\softtext{infra.tarifa\_snapshot.}}{infra.tarifa\_snapshot.}}
PK id y FK
\texorpdfstring{\softtext{E,actor\_id→usuario}}{E,actor\_id→usuario} y
\texorpdfstring{\softtext{E,S→sucursal}}{E,S→sucursal} añaden dos
\texorpdfstring{\softtext{referencias.}}{referencias.} ruta, precio
numeric(14,2), \texorpdfstring{\softtext{hash\_sha256,}}{hash\_sha256,}
snapshot, \texorpdfstring{\softtext{capturado\_en}}{capturado\_en} y
\texorpdfstring{\softtext{aceptada\_en}}{aceptada\_en} representan
captura privada. Precio fuente de cuatro decimales queda en JSON; el
monto canónico a centavos se compara con gasto. El token sólo vale para
\texorpdfstring{\softtext{actor/sede/empresa}}{actor/sede/empresa} de
captura.

\textbf{\texorpdfstring{\softtext{infra.spring\_session.}}{infra.spring\_session.}}
PK primary\_id y unicidad de session\_id sostienen sesión JDBC.
\texorpdfstring{\softtext{creation\_time,}}{creation\_time,}
\texorpdfstring{\softtext{last\_access\_time,}}{last\_access\_time,}
\texorpdfstring{\softtext{max\_inactive\_interval,}}{max\_inactive\_interval,}
\texorpdfstring{\softtext{expiry\_time}}{expiry\_time} y
\texorpdfstring{\softtext{principal\_name}}{principal\_name} son
técnicos. No se replican ni sustituyen membresía vigente. El rol auth
puede modificar sesiones; la app de negocio no. Restore invalida estado
efímero.

\textbf{\texorpdfstring{\softtext{infra.spring\_session\_attributes.}}{infra.spring\_session\_attributes.}}
PK
\texorpdfstring{\softtext{session\_primary\_id,attribute\_name}}{session\_primary\_id,attribute\_name}
y FK a \texorpdfstring{\softtext{spring\_session}}{spring\_session} con
CASCADE preservan atributos
\texorpdfstring{\softtext{serializados.}}{serializados.}
\texorpdfstring{\softtext{attribute\_bytes}}{attribute\_bytes} contiene
bytes privados que no se inspeccionan en el manual. CASCADE es adecuado
para sesiones, no una autorización de borrar hechos económicos
\texorpdfstring{\softtext{confirmados.}}{confirmados.}

\textbf{\texorpdfstring{\softtext{infra.flyway\_schema\_history.}}{infra.flyway\_schema\_history.}}
PK \texorpdfstring{\softtext{installed\_rank}}{installed\_rank} y
columnas version, description, type, script, checksum,
\texorpdfstring{\softtext{installed\_by/on,}}{installed\_by/on,}
\texorpdfstring{\softtext{execution\_time}}{execution\_time} y success
explican \texorpdfstring{\softtext{migración.}}{migración.} No tiene FK
\texorpdfstring{\softtext{empresarial.}}{empresarial.} Su éxito registra
\texorpdfstring{\softtext{DDL/adopción,}}{DDL/adopción,} no todasRN
aprobadas. La entrada baseline3 conserva la verdad de instalación
existente.

\section{5. El protocolo de un comando
SQL}\label{el-protocolo-de-un-comando-sql}

La firma vigente recibe ocho argumentos y retorna JSONB:

\begin{Shaded}
\begin{Highlighting}[]
\NormalTok{bovi.cmd\_ejecutar(}
\NormalTok{    p\_empresa uuid, p\_sucursal uuid, p\_actor uuid,}
\NormalTok{    p\_clave text, p\_hash text, p\_correlation uuid,}
\NormalTok{    p\_comando text, p\_payload jsonb}
\NormalTok{) returns jsonb}
\end{Highlighting}
\end{Shaded}

No son ocho datos comerciales elegidos libremente por el
\texorpdfstring{\softtext{formulario.}}{formulario.} Empresa y sede se
resuelven en contexto autorizado; actor proviene de sesión; clave
identifica intención; hash lo calcula el servidor sobre cuerpo canónico;
correlation relaciona seguimiento; comando selecciona rama; payload
aporta datos tipados de esa
\texorpdfstring{\softtext{intención.}}{intención.} El backend no acepta
una página de proveedor fabricada por el navegador como procedencia
\texorpdfstring{\softtext{auténtica.}}{auténtica.}

La función es SECURITY DEFINER y se ejecuta con derechos del rol técnico
dueño de comandos. Eso permite insertar hechos sin conceder DML directo
a la app. También obliga a ser extremadamente explícito con scope y
objetos \texorpdfstring{\softtext{cualificados.}}{cualificados.}
\texorpdfstring{\softtext{search\_path=pg\_catalog}}{search\_path=pg\_catalog}
y referencias bovi/infra evitan que un objeto de otro esquema sustituya
una función esperada. PUBLIC no tiene EXECUTE.

El control fn\_scope comprueba empresa configurada, usuario activo,
membresía vigente y rol
\texorpdfstring{\softtext{compatible.}}{compatible.} ADMIN satisface
roles operativos \texorpdfstring{\softtext{específicos.}}{específicos.}
Si la intención modifica maestros, exige A. En hechos locales exige que
la sede tenga como dueño el nodo actual. No utiliza
\texorpdfstring{\softtext{current\_user}}{current\_user} como persona:
\texorpdfstring{\softtext{current\_user}}{current\_user} es la identidad
técnica de ejecución, especialmente en SECURITY DEFINER.

Después se valida clave de16-80 caracteres y hash
\texorpdfstring{\softtext{hexadecimal64.}}{hexadecimal64.} INSERT ON
CONFLICT reclama el espacio de intención; SELECT FOR UPDATE bloquea la
fila \texorpdfstring{\softtext{idempotente.}}{idempotente.} Si el hash
difiere se rechaza BV405. Si estado es COMPLETADA se devuelve resultado
\texorpdfstring{\softtext{persistido.}}{persistido.} Sólo una intención
nueva alcanza la lógica del agregado y sus comprobaciones de versión.

Ese orden resuelve un caso central: confirmar con versión0 incrementa el
agregado a1. Si se perdió la respuesta, repetir misma clave y cuerpo con
versión0 debe devolver la liquidación original, no412. En cambio, una
nueva clave intentando modificar versión0 de un agregado que ya está en1
debe fallar. Replay y nueva escritura obsoleta son intenciones
distintas.

Cada rama compone su resultado antes del commit. Al terminar se inserta
\texorpdfstring{\softtext{auditoria\_evento,}}{auditoria\_evento,}
\texorpdfstring{\softtext{evento\_outbox}}{evento\_outbox} y se marca
idempotencia COMPLETADA con ese mismo resultado. Todo pertenece a la
transacción local del comando. Si una FK o una validación falla,
rollback elimina también reclamo nuevo, auditoría y outbox de esa
\texorpdfstring{\softtext{intención.}}{intención.} No se anuncia un
éxito parcial ni se consulta después del commit para construir otra
respuesta.

\textbf{Fuente:}
\texorpdfstring{\softtext{db/migrations/V13\_\_correccion\_composicion\_exacta.sql:2.}}{db/migrations/V13\_\_correccion\_composicion\_exacta.sql:2.}

\subsection{5.1 Leer las variables locales sin confundirlas con
columnas}\label{leer-las-variables-locales-sin-confundirlas-con-columnas}

e,s,a son alias de empresa, sede y actor. j es el payload normalizado; n
es la identidad del nodo. result guarda el JSON de respuesta y prev es
una fila tipada de
\texorpdfstring{\softtext{clave\_idempotencia.}}{clave\_idempotencia.}
op,original,pos,pay,liq,app y gasto son filas \%ROWTYPE de operación,
ticket original, posición, pago, liquidación, aplicación y gasto. Usar
\%ROWTYPE permite que el registro local corresponda al tipo de la tabla,
pero no implica que se haya leído o bloqueado: eso depende del SELECT
concreto.

r y c son record flexibles para resultados de consultas; lot y x son
elementos JSON. idn comienza como UUID nuevo y puede reasignarse al
recurso retornado. opid identifica operación, pid identifica
\texorpdfstring{\softtext{pago/posición/periodo}}{pago/posición/periodo}
según la rama, lid liquidación, wid ticket, aid asignación y per
periodo. La reutilización de pid exige leer su contexto: no se debe
seguir una variable sólo por su nombre suponiendo que siempre significa
pago.

amount es monto de la acción; totaln puede ser total comercial o importe
remanente de reaplicación; cost representa costo o cuota temporal según
rama. cnt cuenta posiciones; k e idx recorren distribución o conceptos;
seqn forma folio o versión de cierre; contractn identifica la versión
\texorpdfstring{\softtext{contractual.}}{contractual.} ref agrupa
reparto; ids es arreglo de identidades; clas y nature seleccionan
vocabularios de clase y naturaleza; f conserva fecha contextual;
preserve decide si existe historia que obliga a mantener
\texorpdfstring{\softtext{composición.}}{composición.}

La directiva
\texorpdfstring{\softtext{variable\_conflict}}{variable\_conflict}
use\_column evita una ambigüedad al coincidir nombres de variables y
columnas. No sustituye escribir consultas cualificadas ni entender qué
valor se compara. Un SELECT INTO no es un UPDATE y tampoco un lock salvo
que incluya FOR UPDATE. Una variable numérica exacta no prueba redondeo
correcto si después se materializa otro dato con una política distinta.

El bloque de DECLARE sirve como mapa de lectura. Conviene recorrerlo y
después señalar dónde se asigna cada variable
\texorpdfstring{\softtext{importante.}}{importante.} Para confirmar
compra, seguir
\texorpdfstring{\softtext{per→op→contractn→totaln→lid→pos→cost.}}{per→op→contractn→totaln→lid→pos→cost.}
Para aplicar dinero, seguir
\texorpdfstring{\softtext{per→pid→pay→lid→liq→amount.}}{per→pid→pay→lid→liq→amount.}
Para corregir ticket, seguir opid derivado del
\texorpdfstring{\softtext{original→op→original→j→idn→result.}}{original→op→original→j→idn→result.}
Ese ejercicio revela dependencias reales y evita una explicación
genérica de ``la función valida todo''.

\section{6. Aritmética decimal, unidades y
materialización}\label{aritmuxe9tica-decimal-unidades-y-materializaciuxf3n}

Una cantidad exacta no empieza en el reporte: empieza cuando se recibe
el dato. El contrato Dinero utiliza un objeto con importe como cadena y
moneda MXN. Peso utiliza kg como cadena. Una tarifa puede tener cuatro
decimales; un importe confirmado se materializa a dos. PostgreSQL
trabaja con NUMERIC y Java transporta
\texorpdfstring{\softtext{BigDecimal;}}{BigDecimal;} el navegador no
envía un total calculado para que la BD lo acepte como verdad.

Ejemplos de representación:

\begin{Shaded}
\begin{Highlighting}[]
\FunctionTok{\{}\DataTypeTok{"importe"}\FunctionTok{:}\StringTok{"12000.01"}\FunctionTok{,}\DataTypeTok{"moneda"}\FunctionTok{:}\StringTok{"MXN"}\FunctionTok{\}}
\FunctionTok{\{}\DataTypeTok{"kg"}\FunctionTok{:}\StringTok{"4900.000"}\FunctionTok{\}}
\end{Highlighting}
\end{Shaded}

La cadena conserva escala y evita pasar por una aproximación binaria
flotante. No basta con llamar decimal a una variable si después una
conversión usa double para sumar. La autoridad que confirma importes es
SQL. Java y UI presentan los resultados canónicos y mantienen el
contexto de intención; no implementan otra fórmula que pueda divergir
del resultado \texorpdfstring{\softtext{persistido.}}{persistido.}

numeric(18,2) expresa precisión total18 y escala2, no dieciocho enteros
más dos decimales. numeric(14,3) conserva masa con tres decimales;
numeric(14,4) admite tarifas con cuatro. Una restricción de precisión
limita magnitud
\texorpdfstring{\softtext{representable.}}{representable.} Los CHECK
evitan montos inválidos, NaN o infinitos según el campo. Las funciones
comprueban también condiciones semánticas como importe de aplicación
positivo y ya redondeado a centavos.

\subsection{6.1 Neto, variación física y deducción
contractual}\label{neto-variaciuxf3n-fuxedsica-y-deducciuxf3n-contractual}

Para un ticket \texorpdfstring{\softtext{BRUTO\_TARA,}}{BRUTO\_TARA,}
sean B la masa bruta en kilogramos, T la tara y N el neto. La ecuación y
sus condiciones son:

\[
N=B-T,\qquad 0\leq T\leq B,\qquad N>0.
\]

Bruto incluye el elemento que se descontará; tara es esa parte y neto
representa masa liquidable
\texorpdfstring{\softtext{registrada.}}{registrada.} Un ticket
\texorpdfstring{\softtext{NETO\_REPORTADO}}{NETO\_REPORTADO} conserva un
neto declarado sin inventar bruto y tara. La precisión y procedencia
siguen siendo \texorpdfstring{\softtext{explícitas.}}{explícitas.} No
debe llenarse bruto con neto sólo para evitar NULL, porque se estaría
fabricando una medición distinta de la real.

Para comparar dos tickets, sean W\_i el peso inicial y W\_f el final. La
variación y porcentaje son:

\[
\Delta W=W_i-W_f,\qquad
v=100\,\frac{W_i-W_f}{W_i},\qquad W_i>0.
\]

W\_i y W\_f se expresan en kg; v se expresa en
\texorpdfstring{\softtext{porcentaje.}}{porcentaje.} Si el peso final es
mayor, la variación es negativa y se conserva ese signo. Una ganancia de
peso no debe convertirse a cero por comodidad visual. Antes de calcular
se exige igualdad del conjunto de animales; dos totales de lotes
distintos no producen una merma
\texorpdfstring{\softtext{comparable.}}{comparable.}

En GOLD-01, W\_i=5000 kg y W\_f=4900 kg. La diferencia es100 kg y el
\texorpdfstring{\softtext{porcentaje2.}}{porcentaje2.} Ese2\% describe
evidencia física. No autoriza descontar nuevamente2\% del precio. La
deducción contractual d existe en
\texorpdfstring{\softtext{condicion\_lote}}{condicion\_lote} y se aplica
solamente si se pactó. Mezclar ambas variables produciría doble
descuento.

Cuando la unidad comercial es KG:

\[
Q_{\mathrm{kg}}=N\left(1-\frac{d}{100}\right).
\]

Q es cantidad liquidable; N es el neto del ticket seleccionado y d es
porcentaje pactado. Si N=1000 kg y d=1.5, Q=985 kg. Si no se pactó
deducción, d=0 y Q=N. Este ejemplo enseña la fórmula, no agrega un
descuento inexistente a GOLD-01. El porcentaje de variación sigue siendo
información \texorpdfstring{\softtext{independiente.}}{independiente.}

Para CABEZA, Q=n es el conteo de animales o posiciones seleccionadas en
ese lote. No se usa el conteo total de la sede ni el número de filas
después de una unión duplicada. La selección tiene unicidad por
operación y las FK de tipo impiden contabilizar una selección de venta
como compra.

\textbf{Fuentes:}
\texorpdfstring{\softtext{db/migrations/V2\_\_comandos.sql:70;}}{db/migrations/V2\_\_comandos.sql:70;}
\texorpdfstring{\softtext{docs/domain/states-and-calculations.md:47;}}{docs/domain/states-and-calculations.md:47;}
\texorpdfstring{\softtext{db/migrations/V12\_\_pesajes\_reportes\_canonicos.sql:430.}}{db/migrations/V12\_\_pesajes\_reportes\_canonicos.sql:430.}

\subsection{6.2 Base por lote y total
comercial}\label{base-por-lote-y-total-comercial}

Sean P\_l la tarifa del lote l, Q\_l su cantidad liquidable, C\_l los
cargos y B\_l las bonificaciones pactadas. Se materializan base y total:

\[
A_l=\operatorname{round}_{2}(Q_lP_l),
\qquad
L_l=A_l+C_l-B_l,
\qquad
L=\sum_{l=1}^{m}L_l.
\]

A\_l es base monetaria a centavos; L\_l total del lote y L total de
\texorpdfstring{\softtext{operación.}}{operación.} m es número de lotes
con condiciones de la versión relevante. Se redondea al materializar el
concepto, no después de sumar una tarifa ya truncada. La política
HALF\_UP a dos decimales responde al reconocimiento de cada importe
\texorpdfstring{\softtext{persistido.}}{persistido.} Cargos y
bonificaciones están expresados como importes y se materializan como
conceptos distintos.

El SQL selecciona la versión máxima de contrato cuando la operación
sigue BORRADOR; para una operación acordada o confirmada utiliza la
versión aceptada. Esta elección evita simular siempre versión1 cuando
existe \texorpdfstring{\softtext{reapertura.}}{reapertura.} Una
propuesta nueva no cambia una liquidación confirmada: la liquidación
conserva \texorpdfstring{\softtext{numero\_version}}{numero\_version} y
\texorpdfstring{\softtext{snapshot\_condiciones.}}{snapshot\_condiciones.}

Para una compra de5000 kg a48MXN/kg, la base
\texorpdfstring{\softtext{es240000.00.}}{es240000.00.} Con gastos de
adquisición posteriores
\texorpdfstring{\softtext{de12000.00,}}{de12000.00,} el costo de
posiciones llega \texorpdfstring{\softtext{a252000.00,}}{a252000.00,}
pero la obligación comercial de compra no se modifica automáticamente
por esos gastos. Esa separación es esencial: costo gerencial y deuda con
la contraparte no son la misma cifra.

En una venta de4900 kg a60MXN/kg, la base
\texorpdfstring{\softtext{es294000.00.}}{es294000.00.} Si hay un gasto
directo de venta de8000.00, reduce margen, no necesariamente total
adeudado por comprador. Los cargos que sí se pactaron con comprador
deben estar en condiciones o ajustes de obligación, no camuflados como
gastos propios.

\subsection{6.3 Reparto exacto y residuo
determinista}\label{reparto-exacto-y-residuo-determinista}

El reparto no utiliza pesos que no fueron medidos. Se admiten
\texorpdfstring{\softtext{POR\_CABEZA,}}{POR\_CABEZA,}
\texorpdfstring{\softtext{POR\_PESO\_MEDIDO}}{POR\_PESO\_MEDIDO} e
\texorpdfstring{\softtext{IMPORTE\_INDIVIDUAL.}}{IMPORTE\_INDIVIDUAL.}
En los dos primeros, el gasto G se distribuye con base w\_i para cada
posición i. Por cabeza w\_i=1. Por peso w\_i es medición individual
aprobada de la evidencia
\texorpdfstring{\softtext{compatible.}}{compatible.}

\[
a_i^{*}=G\,\frac{w_i}{\sum_{j=1}^{n}w_j},
\qquad
r_i=\operatorname{round}_{2}(a_i^{*}),
\qquad
R=G-\sum_{i=1}^{n}r_i.
\]

a\_i estrella es cuota exacta antes de redondear; r\_i cuota a centavos
y R residuo. R no desaparece y tampoco se asigna a una posición
\texorpdfstring{\softtext{arbitraria.}}{arbitraria.} El SQL convierte a
centavos y corrige las primeras filas ordenadas por
\texorpdfstring{\softtext{posicion\_id}}{posicion\_id}
\texorpdfstring{\softtext{ascendente.}}{ascendente.} El orden estable
hace reproducible el resultado al consultar o reintentar la misma
\texorpdfstring{\softtext{intención.}}{intención.}

Para G=1.00 y tres posiciones iguales, las cuotas exactas son un tercio.
Después de redondear
\texorpdfstring{\softtext{resultan0.33,0.33,0.33;}}{resultan0.33,0.33,0.33;}
queda0.01. La primera posición recibe el centavo y las cuotas finales
\texorpdfstring{\softtext{son0.34,0.33,0.33.}}{son0.34,0.33,0.33.} Ese
es GOLD-03, un esperado fijo. No se obtiene llamando a la misma función
que se está probando para fabricar el oráculo.

La corrección también puede ser negativa: si el redondeo provisional
excedió G, el algoritmo resta centavos según el mismo orden. Para una
bonificación, el signo nace de naturaleza validada; el usuario no
introduce deltas monetarios libres. El importe acumulado de una posición
no puede volverse negativo para ocultar un reparto excesivo.

En \texorpdfstring{\softtext{POR\_PESO\_MEDIDO}}{POR\_PESO\_MEDIDO} se
consultan pesos individuales del ticket aprobado ligado a condiciones de
la liquidación de compra. Si una posición no tiene medición individual
positiva, se rechaza BV018. Un ticket grupal de300 kg para dos animales
no demuestra150 kg para cada uno. La ausencia de medición es información
significativa y debe permanecer NULL.

En
\texorpdfstring{\softtext{IMPORTE\_INDIVIDUAL,}}{IMPORTE\_INDIVIDUAL,}
se envían todas las posiciones una sola vez con cantidades explícitas
positivas a centavos. Las identidades deben pertenecer al conjunto y la
suma debe igualar el gasto. No se permite omitir una posición para
absorber residuo sin
\texorpdfstring{\softtext{explicación.}}{explicación.} Una
\texorpdfstring{\softtext{distribución0.40/0.60}}{distribución0.40/0.60}
de1.00 es válida si corresponde a todo el
\texorpdfstring{\softtext{conjunto;0.40/0.59}}{conjunto;0.40/0.59} no
cuadra aunque ``casi'' llegue al total.

\textbf{Fuente:}
\texorpdfstring{\softtext{db/migrations/V2\_\_comandos.sql:64;}}{db/migrations/V2\_\_comandos.sql:64;}
\texorpdfstring{\softtext{db/migrations/V2\_\_comandos.sql:99.}}{db/migrations/V2\_\_comandos.sql:99.}

\subsection{6.4 Saldo, disponible y margen son tres resultados
diferentes}\label{saldo-disponible-y-margen-son-tres-resultados-diferentes}

Para una liquidación, sean L\_0 total confirmado original, C ajustes
CARGO, B ajustes ABONO y A aplicaciones netas. El total ajustado y saldo
son:

\[
L_{\mathrm{ajustado}}=L_0+\sum C-\sum B,
\qquad
S=L_{\mathrm{ajustado}}-A.
\]

Aplicaciones netas significa suma de APLICACION menos suma de REVERSO.
L\_0 permanece en la fila
\texorpdfstring{\softtext{histórica.}}{histórica.} S no se introduce en
un UPDATE manual para cuadrar. Si se modifica una aplicación, cambia la
vista del saldo al añadir un contramovimiento, no al borrar evidencia.

Para un pago normal con importe P\_0, aplicaciones netas A\_p y
reembolsos F:

\[
D=P_0-A_p-\sum F.
\]

D es \texorpdfstring{\softtext{disponible.}}{disponible.} Puede ser
positivo aunque una obligación ya tenga saldo cero: es un anticipo o
saldo a favor. Un reembolso disminuye D y representa flujo real de caja.
No cambia la venta reconocida ni el costo del
\texorpdfstring{\softtext{inventario.}}{inventario.} El nuevo pago de
naturaleza REEMBOLSO conserva referencia al original y toma dirección
opuesta: COBRO original produce PAGO de devolución, y viceversa.

Para margen de venta, sean V la venta neta reconocida, K el costo de
posiciones vendidas y G\_d los gastos directos:

\[
M=V-K-G_d.
\]

M no usa cobros como ingreso. Un comprador que aún no pagó puede haber
generado venta y margen reconocidos; dinero recibido como anticipo sin
venta no genera margen. Separar devengo gerencial y caja evita que el
reporte mejore o empeore sólo porque se aplicó un pago.

Los reportes calculan sobre todo el conjunto filtrado, no sobre la
primera página visible. Si la cartera muestra25 filas de200, la UI no
suma esas25 como si fueran toda la sede. La consulta y el contrato
indican paginación, corte y fuente. El margen agrega el filtro completo
y conserva faltantes, evitando etiquetar una cifra provisional como
cierre \texorpdfstring{\softtext{definitivo.}}{definitivo.}

\subsection{6.5 Asientos de valoración y signos
derivados}\label{asientos-de-valoraciuxf3n-y-signos-derivados}

C es costo acumulado de una posición. Una adquisición genera
\texorpdfstring{\softtext{inventario+C.}}{inventario+C.} Venta genera
inventario−C y costo vendido+C. Baja sin venta genera inventario−C y
\texorpdfstring{\softtext{pérdida+C.}}{pérdida+C.} Un costo tardío δ se
dirige al eje correspondiente al estado actual del ciclo. Una
bonificación produce efecto contrario bajo límites que evitan un costo
acumulado negativo.

{\def\LTcaptype{none} % do not increment counter
\begin{longtable}[]{@{}lll@{}}
\toprule\noalign{}
Hecho & Delta de inventario & Otro efecto \\
\midrule\noalign{}
\endhead
\bottomrule\noalign{}
\endlastfoot
Adquisición C & +C & Sin costo vendido ni pérdida. \\
Venta C & −C & Costo vendido+C. \\
Costo tardío capitalizado δ & +δ & Sin nueva deuda por sí solo. \\
Costo tardío ya vendido δ & 0 & Costo
\texorpdfstring{\softtext{vendido+δ.}}{vendido+δ.} \\
Baja C & −C & Pérdida+C, sin venta. \\
Costo tardío tras baja δ & 0 &
\texorpdfstring{\softtext{Pérdida+δ.}}{Pérdida+δ.} \\
Anulación antes de entrega C neto & +C neto & Costo vendido−C neto. \\
\end{longtable}
}

Las columnas
\texorpdfstring{\softtext{delta\_inventario,}}{delta\_inventario,}
\texorpdfstring{\softtext{delta\_costo\_venta}}{delta\_costo\_venta} y
\texorpdfstring{\softtext{delta\_perdida}}{delta\_perdida} son
generadas. Clase y sentido determinan su fórmula a partir de importe. Un
usuario no puede enviar tres números contradictorios para que un reporte
aparente cuadrar. Los asientos preservan operación y posición para
atribuir costo tardío a la venta original, incluso si su fuente
administrativa fue la compra.

\section{7. Recorridos reales de
comandos}\label{recorridos-reales-de-comandos}

\subsection{7.1 Maestros: crear sin duplicar
identidad}\label{maestros-crear-sin-duplicar-identidad}

\texorpdfstring{\softtext{crear\_contraparte}}{crear\_contraparte} y
\texorpdfstring{\softtext{crear\_predio}}{crear\_predio} escriben en A
después de fn\_scope con modo maestro.
\texorpdfstring{\softtext{Registrar\_animal}}{Registrar\_animal} crea
identidad y recorre todo el arreglo de identificaciones, no sólo la
primera. Cada código se valida y normaliza; la UK de
\texorpdfstring{\softtext{sistema/código}}{sistema/código} resuelve
colisiones en el escritor central. Si una identificación posterior del
mismo comando colisiona, rollback retira también animal e
identificaciones anteriores de esa
\texorpdfstring{\softtext{intención.}}{intención.}

\texorpdfstring{\softtext{crear\_usuario}}{crear\_usuario} recibe un
hash bcrypt generado por el backend, nunca acepta del cliente una
contraseña ya presentada como hash para eludir políticas del servidor.
\texorpdfstring{\softtext{actualizar\_usuario}}{actualizar\_usuario}
conserva empresa y comprueba
\texorpdfstring{\softtext{existencia.}}{existencia.}
\texorpdfstring{\softtext{asignar\_membresia}}{asignar\_membresia}
enlaza usuario y rol con la sede
\texorpdfstring{\softtext{autorizada.}}{autorizada.} La cuenta ADMIN de
una sede no puede introducir
\texorpdfstring{\softtext{sucursal\_id}}{sucursal\_id} de otra en el
cuerpo y aprovechar que el maestro vive en A.

Las altas maestras se replican hacia B. Si B todavía no observa una
identidad, no debe inventar una copia para continuar compra. El servidor
espera o informa disponibilidad según la
\texorpdfstring{\softtext{operación.}}{operación.} La replicación
asíncrona tiene visibilidad eventual; la unicidad canónica no se
consigue permitiendo que ambos nodos inventen filas y después las
comparen por nombre.

\subsection{7.2 Borrador, edición y
acuerdo}\label{borrador-ediciuxf3n-y-acuerdo}

\texorpdfstring{\softtext{crear\_operacion}}{crear\_operacion} obtiene
folio mediante contador, inserta cabecera BORRADOR y contrato PROPUESTA
versión1. Los lotes se numeran en orden del arreglo. La rama coloca
selección de animales en compra o posiciones en venta, y condiciones por
lote. Si se aporta neto para KG, se crea ticket propuesto con
composición correspondiente, todavía sin aprobación
\texorpdfstring{\softtext{implícita.}}{implícita.}

\texorpdfstring{\softtext{editar\_operacion}}{editar\_operacion} exige
estado BORRADOR y
\texorpdfstring{\softtext{version\_lock}}{version\_lock} actual bajo
lock de \texorpdfstring{\softtext{operación.}}{operación.} Una ACORDADA
requiere reabrir, no PATCH directo. Si existe historia de pesajes o
aceptación previa, preserve exige que el conjunto
\texorpdfstring{\softtext{lote/elemento}}{lote/elemento} entrante
coincida con el anterior. Se pueden cambiar condiciones de la propuesta
vigente sin borrar condiciones aceptadas. No se elimina historia para
facilitar el \texorpdfstring{\softtext{formulario.}}{formulario.}

Cuando no existe historia, la edición puede reconstruir selecciones y
lotes del borrador. Cuando sí existe, conserva composición y
\texorpdfstring{\softtext{ordinalidad.}}{ordinalidad.} La política
\texorpdfstring{\softtext{requiere\_recepcion}}{requiere\_recepcion} se
cambia sólo si llega
\texorpdfstring{\softtext{explícitamente.}}{explícitamente.}
\texorpdfstring{\softtext{fecha\_operacion}}{fecha\_operacion} también
se conserva cuando PATCH omite ese campo: usar
\texorpdfstring{\softtext{current\_date}}{current\_date} como reemplazo
automático habría alterado una fecha documental sin intención del
usuario.

\texorpdfstring{\softtext{acordar\_operacion}}{acordar\_operacion}
bloquea operación, verifica
\texorpdfstring{\softtext{BORRADOR/version}}{BORRADOR/version} y
condiciones, acepta contrato y cambia a ACORDADA. El trigger de V8
completa validación de acuerdo: no admite KG sin ticket aprobado y
conjunto \texorpdfstring{\softtext{compatible.}}{compatible.}
\texorpdfstring{\softtext{reabrir\_operacion}}{reabrir\_operacion}
verifica \texorpdfstring{\softtext{ACORDADA/version}}{ACORDADA/version}
y clona
\texorpdfstring{\softtext{contrato/condiciones}}{contrato/condiciones}
hacia la siguiente propuesta, incrementando
\texorpdfstring{\softtext{version\_lock}}{version\_lock} y devolviendo
BORRADOR. El contrato anterior queda ACEPTADA e inmutable.

\subsection{7.3 Ticket propuesto, aprobación y corrección
V13}\label{ticket-propuesto-aprobaciuxf3n-y-correcciuxf3n-v13}

\texorpdfstring{\softtext{registrar\_pesaje}}{registrar\_pesaje} bloquea
operación y admite únicamente contexto editable
\texorpdfstring{\softtext{BORRADOR/ACORDADA.}}{BORRADOR/ACORDADA.}
Calcula neto de bruto/tara cuando corresponde o conserva neto reportado.
Inserta el conjunto indicado con pesos individuales reales o NULL.
Aprobar exige composición y cambia PROPUESTA a APROBADO. Esa aprobación
congela medición y asociaciones mediante triggers.

\texorpdfstring{\softtext{corregir\_pesaje}}{corregir\_pesaje} deriva
opid del ticket original indicado por
\texorpdfstring{\softtext{ruta/contexto.}}{ruta/contexto.} Después
bloquea operación y original en ese orden. El original debe estar
APROBADO. Si se envía
\texorpdfstring{\softtext{operacion\_id}}{operacion\_id} o
\texorpdfstring{\softtext{numero\_lote}}{numero\_lote} incompatible con
el original, rechaza; no permite usar una ruta para modificar el estado
de un ticket y crear el nuevo bajo otra
\texorpdfstring{\softtext{operación.}}{operación.}

El motivo es obligatorio, recortado y de1-200
\texorpdfstring{\softtext{caracteres.}}{caracteres.} V13 calcula
diferencia de conjuntos original−nuevo y
\texorpdfstring{\softtext{nuevo−original.}}{nuevo−original.} Si
cualquiera contiene elementos, rechaza BV013 antes de UPDATE del
original. Un subconjunto de dos animales reducido a uno no pasa; un
superconjunto que añade un tercero existente tampoco. Validar sólo FK
habría permitido un tercero válido, pero ajeno al ticket.

Después de las guardas se cambia únicamente estado del original a
CORREGIDO. Se crea ticket nuevo PROPUESTA con
\texorpdfstring{\softtext{corrige\_a\_id,}}{corrige\_a\_id,} la misma
composición y mediciones nuevas. En BORRADOR puede actualizarse la
referencia de condición vigente al nuevo ticket KG; una condición
aceptada no se reescribe por esa vía. El ticket corregido requiere su
propia aprobación antes de
\texorpdfstring{\softtext{utilizarlo.}}{utilizarlo.}

El motivo forma parte del resultado del nuevo ticket y ese resultado se
guarda en
\texorpdfstring{\softtext{auditoria\_evento.cambio}}{auditoria\_evento.cambio}
e \texorpdfstring{\softtext{idempotencia.}}{idempotencia.}
\texorpdfstring{\softtext{fn\_pesaje\_json}}{fn\_pesaje\_json} lo
recupera por entidad\_id y comando
\texorpdfstring{\softtext{corregir\_pesaje.}}{corregir\_pesaje.} Un
ticket antiguo sin motivo auditado devuelve
\texorpdfstring{\softtext{motivo\_correccion}}{motivo\_correccion} NULL;
no se fabrica una justificación
\texorpdfstring{\softtext{histórica.}}{histórica.} Replay ocurre antes
de volver a exigir original APROBADO, pues el original ya está CORREGIDO
después del primer éxito.

\textbf{Fuentes:}
\texorpdfstring{\softtext{db/migrations/V13\_\_correccion\_composicion\_exacta.sql:99;}}{db/migrations/V13\_\_correccion\_composicion\_exacta.sql:99;}
\texorpdfstring{\softtext{db/migrations/V12\_\_pesajes\_reportes\_canonicos.sql:2.}}{db/migrations/V12\_\_pesajes\_reportes\_canonicos.sql:2.}

\subsection{7.4 Reserva y reloj de
confirmación}\label{reserva-y-reloj-de-confirmaciuxf3n}

\texorpdfstring{\softtext{reservar\_venta}}{reservar\_venta} bloquea
operación y posiciones por id
\texorpdfstring{\softtext{ascendente.}}{ascendente.} Primero expira
reservas activas cuyo vencimiento ya pasó. Comprueba estado DISPONIBLE,
animal VIGENTE y ausencia de baja local. La nueva reserva tiene
vencimiento explícito o valor por defecto de15 minutos desde reloj real.
Un índice parcial impide dos ACTIVA para la misma posición.

\texorpdfstring{\softtext{liberar\_reserva}}{liberar\_reserva} puede
derivar operación a partir de reserva\_id y procesa posiciones de la
venta. Libera las ACTIVA relacionadas; no vende ni reconoce ingreso.
Cancelar una operación libera reservas y cambia estado, sin crear
\texorpdfstring{\softtext{liquidación.}}{liquidación.} Un usuario que
conserva pantalla antigua no puede interpretar una reserva liberada o
vencida como compromiso vigente.

now permanece fijo al comienzo de
\texorpdfstring{\softtext{transacción.}}{transacción.} Si una sesión
empieza a las10:00 y espera un lock hasta10:16, una reserva que
vencía10:15 podría parecer vigente al compararla con now. Por eso la
confirmación usa
\texorpdfstring{\softtext{clock\_timestamp:}}{clock\_timestamp:} vuelve
a consultar reloj real durante el
\texorpdfstring{\softtext{procedimiento.}}{procedimiento.} También hay
una comprobación al final de consumo para detectar vencimiento ocurrido
durante una confirmación larga. Si falla, toda la confirmación se
revierte.

\subsection{7.5 Confirmar compra y
venta}\label{confirmar-compra-y-venta}

La confirmación reconoce periodo explícito y bloquea su fila antes de
los recursos \texorpdfstring{\softtext{financieros.}}{financieros.} Lee
operación bajo lock, exige versión actual y estado
\texorpdfstring{\softtext{BORRADOR/ACORDADA,}}{BORRADOR/ACORDADA,}
comprueba lotes, cantidades y evidencia. Los tickets KG deben estar
aprobados y tener el conjunto
\texorpdfstring{\softtext{seleccionado.}}{seleccionado.} El contrato
actual se acepta si todavía era propuesta, y se calcula total por
función canónica.

Se inserta liquidación principal con dirección según tipo: compra
POR\_PAGAR, venta \texorpdfstring{\softtext{POR\_COBRAR.}}{POR\_COBRAR.}
Se copian nombre/RFC de contraparte y
\texorpdfstring{\softtext{condiciones.}}{condiciones.} Se materializan
conceptos base, cargos y
\texorpdfstring{\softtext{bonificaciones.}}{bonificaciones.} Las filas y
snapshots corresponden a la versión aceptada, de modo que una tarifa o
nombre posterior no recalcula la
\texorpdfstring{\softtext{obligación.}}{obligación.}

En compra, se bloquea identidad de cada animal en orden estable, exige
VIGENTE y no habilita otro ciclo frente a una baja pendiente de
revisión. Se crean posiciones con
\texorpdfstring{\softtext{numero\_ciclo}}{numero\_ciclo} siguiente.
\texorpdfstring{\softtext{requiere\_recepcion}}{requiere\_recepcion}
decide
\texorpdfstring{\softtext{PENDIENTE\_RECEPCION}}{PENDIENTE\_RECEPCION} o
\texorpdfstring{\softtext{DISPONIBLE.}}{DISPONIBLE.} Los conceptos se
distribuyen a posiciones con residuo exacto, y los asientos ADQUISICION
reconocen costo inicial positivo.

En venta, las posiciones se bloquean por id y cada una exige DISPONIBLE
y reserva ACTIVA vigente por
\texorpdfstring{\softtext{clock\_timestamp.}}{clock\_timestamp.} Se toma
su inventario acumulado como costo, se crea asiento VENTA, se cambia a
\texorpdfstring{\softtext{VENDIDA\_PENDIENTE\_ENTREGA}}{VENDIDA\_PENDIENTE\_ENTREGA}
y se consume reserva. La entrega posterior cierra físicamente el ciclo,
sin crear otro asiento de venta. El resultado canónico es la liquidación
completa y se persiste tal cual para replay.

\section{8. Dinero aplicado, reversos y correcciones sin borrar
originales}\label{dinero-aplicado-reversos-y-correcciones-sin-borrar-originales}

\subsection{8.1 Registrar y validar un pago son actos
distintos}\label{registrar-y-validar-un-pago-son-actos-distintos}

\texorpdfstring{\softtext{registrar\_pago}}{registrar\_pago} crea una
fila REGISTRADO de naturaleza NORMAL. La rama vigente deja periodo\_id y
\texorpdfstring{\softtext{fecha\_reconocimiento}}{fecha\_reconocimiento}
NULL. Su importe y referencia describen dinero declarado; no descuenta
una obligación ni reconoce venta. Exigir ya un cierre abierto para una
simple captura de dinero habría confundido registro con
\texorpdfstring{\softtext{validación.}}{validación.} El estado permite
comprobar esa separación en la consulta.

\texorpdfstring{\softtext{validar\_pago}}{validar\_pago} exige CAJA y
llama
\texorpdfstring{\softtext{fn\_periodo\_payload}}{fn\_periodo\_payload}
con periodo y fecha elegidos. Cambia un REGISTRADO a VALIDADO y fija
\texorpdfstring{\softtext{reconocimiento.}}{reconocimiento.} Si no
encuentra el pago en ese estado, rechaza. La validación no crea una
aplicación implícita: el usuario puede tener un pago validado
completamente disponible, por ejemplo un anticipo. La venta y su costo
continúan independientes de esta caja.

El resultado de registrar, validar o reembolsar se construye desde
\texorpdfstring{\softtext{v\_disponible\_pago,}}{v\_disponible\_pago,}
se normaliza a JSON decimal y conserva pago\_id. Se guarda como
respuesta del comando. No se debe presentar un disponible inventado por
la UI sumando únicamente las aplicaciones que aparecen en una página.

\subsection{8.2 Aplicar con saldo y disponible bajo
locks}\label{aplicar-con-saldo-y-disponible-bajo-locks}

\texorpdfstring{\softtext{aplicar\_pago}}{aplicar\_pago} bloquea primero
periodo, luego pago y las obligaciones en orden estable de identidad.
pay es la fila de dinero; liq la obligación; amount el monto de la
\texorpdfstring{\softtext{aplicación.}}{aplicación.} La rama admite un
arreglo de aplicaciones y ordena por
\texorpdfstring{\softtext{liquidacion\_id.}}{liquidacion\_id.} Esa
jerarquía reduce el riesgo de que dos comandos bloqueen los mismos
recursos en órdenes opuestos.

Antes de cada inserción comprueba pago VALIDADO y NORMAL. Pago y
obligación deben pertenecer al mismo E/S, a la misma contraparte y
moneda, y tener dirección compatible: un COBRO satisface
\texorpdfstring{\softtext{POR\_COBRAR;}}{POR\_COBRAR;} un PAGO satisface
POR\_PAGAR. La FK conserva scope estructural, pero no prueba esa
compatibilidad semántica de dirección y
\texorpdfstring{\softtext{contraparte.}}{contraparte.}

La cantidad debe ser positiva, estar a centavos y no superar ni
disponible del pago ni saldo de la
\texorpdfstring{\softtext{obligación.}}{obligación.} Estas sumatorias se
leen mientras se conservan los locks. Si dos conexiones compiten por
pago600 y cada una pretende aplicar400, una puede aplicar400; la otra
observa disponible200 y falla. Revisar disponible sólo antes del lock
permitiría que ambas vieran600 y juntas
\texorpdfstring{\softtext{aplicaran800.}}{aplicaran800.}

El loop de aplicaciones pertenece a una sola
\texorpdfstring{\softtext{transacción.}}{transacción.} Si la segunda
obligación del arreglo es incompatible, no se conserva una primera
aplicación parcial como éxito del comando. Se revierten todas las
inserciones de esa \texorpdfstring{\softtext{intención.}}{intención.}
Esto facilita que un reintento de la misma clave tenga un significado
definido.

\subsection{8.3 Reversar no equivale a
eliminar}\label{reversar-no-equivale-a-eliminar}

\texorpdfstring{\softtext{revertir\_aplicacion}}{revertir\_aplicacion}
busca el original de naturaleza APLICACION y bloquea pago y liquidación
asociados. Inserta una fila REVERSO con el mismo importe y
\texorpdfstring{\softtext{reverso\_de\_id.}}{reverso\_de\_id.} La UK
sobre referencia impide repetir reverso del mismo original. El CHECK
exige que naturaleza REVERSO corresponda a un original enlazado.

Las vistas usan signo por
\texorpdfstring{\softtext{naturaleza.}}{naturaleza.} Si existía
aplicación1000, el reverso1000 lleva la suma neta a0 y recupera
\texorpdfstring{\softtext{disponible.}}{disponible.} Las dos filas
siguen explicando qué se hizo y qué se corrigió. DELETE habría borrado
esa explicación y perjudicado reconciliación entre saldo histórico,
auditoría y \texorpdfstring{\softtext{comprobantes.}}{comprobantes.}

No existe un permiso general de ``cambiar importes
\texorpdfstring{\softtext{históricos”.}}{históricos”.} El comando de
reverso utiliza actor, periodo y contexto actuales. Un reverso
reconocido este mes puede afectar el saldo vivo de una obligación
anterior sin modificar el corte ya guardado del periodo anterior. La
cronología de reconocimientos continúa
\texorpdfstring{\softtext{verificable.}}{verificable.}

\subsection{8.4 Bonificación de obligación completamente pagada:
GOLD-07}\label{bonificaciuxf3n-de-obligaciuxf3n-completamente-pagada-gold-07}

Considérese obligación1000 y pago1000 aplicado. Se solicita ABONO100. Si
simplemente se insertara el ajuste, la obligación ajustada quedaría900
frente a aplicaciones netas1000. El modelo rechaza esa
\texorpdfstring{\softtext{sobreaplicación.}}{sobreaplicación.} El
comando guiado permite
\texorpdfstring{\softtext{reaplicar\_pagos:true}}{reaplicar\_pagos:true}
para ejecutar la conciliación completa de manera atómica.

La rama
\texorpdfstring{\softtext{ajustar\_liquidacion}}{ajustar\_liquidacion}
obtiene periodo, bloquea los pagos involucrados por id y después
\texorpdfstring{\softtext{liquidación.}}{liquidación.} Comprueba que el
abono no exceda total ajustado. Si excede saldo actualmente disponible
para abonar, requiere la opción explícita de reverso y
\texorpdfstring{\softtext{reaplicación.}}{reaplicación.} totaln se
convierte en la obligación remanente después del abono, no en un nuevo
total original.

Se recorren aplicaciones vigentes sin reverso por
\texorpdfstring{\softtext{aplicada\_en}}{aplicada\_en} e id. Cada una se
\texorpdfstring{\softtext{contrarresta.}}{contrarresta.} cost se usa
aquí como cantidad que se reaplica:
\texorpdfstring{\softtext{least(app.importe,totaln).}}{least(app.importe,totaln).}
Su nombre no significa costo ganadero en esta rama; leer la asignación
revela el \texorpdfstring{\softtext{significado.}}{significado.} El
remanente disminuye después de cada
\texorpdfstring{\softtext{reaplicación.}}{reaplicación.}

En el esperado fijo, reverso1000 recupera dinero y aplicación900 lo
vuelve a vincular; ajuste ABONO100 deja obligación900 y saldo0.
Disponible queda100. Reembolsar100 crea flujo contrario de caja y
devuelve \texorpdfstring{\softtext{disponible0.}}{disponible0.} Todos
los pasos de
\texorpdfstring{\softtext{ajuste/reconciliación}}{ajuste/reconciliación}
se confirman juntos; el reembolso es una intención posterior de dinero
real, no una eliminación de
\texorpdfstring{\softtext{aplicación.}}{aplicación.}

Si
\texorpdfstring{\softtext{reaplicar\_pagos:false}}{reaplicar\_pagos:false}
o ausente no autoriza la conciliación necesaria, la función devuelve
BV020. Un fallo intermedio revierte reversos, nuevas aplicaciones y
ajuste. El usuario no debe quedar frente a un estado temporal confirmado
donde el pago desapareció o la obligación quedó negativa.

\textbf{Fuente:}
\texorpdfstring{\softtext{db/migrations/V13\_\_correccion\_composicion\_exacta.sql:214;}}{db/migrations/V13\_\_correccion\_composicion\_exacta.sql:214;}
\texorpdfstring{\softtext{db/tests/gold.py:96.}}{db/tests/gold.py:96.}

\section{9. Gastos, costo tardío, baja y
anulación}\label{gastos-costo-tarduxedo-baja-y-anulaciuxf3n}

\texorpdfstring{\softtext{registrar\_gasto}}{registrar\_gasto} exige
periodo abierto para su reconocimiento contextual y crea un gasto
BORRADOR. Si utiliza proveedor externo, recibe
\texorpdfstring{\softtext{snapshot\_id}}{snapshot\_id} de una captura
aceptada del mismo E/S/actor; compara el importe con el precio canónico
local. No permite que el navegador envíe
\texorpdfstring{\softtext{snapshot\_fuente\_externa}}{snapshot\_fuente\_externa}
arbitrario sin token. El JSON copiado al gasto conserva procedencia y
\texorpdfstring{\softtext{aceptación.}}{aceptación.}

\texorpdfstring{\softtext{repartir\_gasto}}{repartir\_gasto} o
\texorpdfstring{\softtext{registrar\_costo\_tardio}}{registrar\_costo\_tardio}
vuelve a validar periodo, bloquea gasto y exige BORRADOR. Una
BONIFICACION debe enlazar costo aprobado de misma operación y
clasificación, sin superar saldo de bonificación permitido. Así se evita
reducir tres veces un flete que sólo se reconoció una vez. El gasto
original sigue \texorpdfstring{\softtext{consultable.}}{consultable.}

Cuando clasificación es CAPITALIZABLE, se obtienen posiciones de esa
compra en orden y se bloquean. fn\_reparto determina cuotas por método.
Cada asignación conserva fuente y grupo. Si una cuota es cero no se
inventa asiento monetario positivo; si haría negativo el costo
acumulado, falla. La clase del asiento depende del estado de posición:
inventario, venta reconocida o baja.

Un detalle importante es opid en costo tardío de venta. Aunque el gasto
se registró sobre la compra, la posición ya vendida requiere atribuir
\texorpdfstring{\softtext{delta\_costo\_venta}}{delta\_costo\_venta} a
su venta original. La rama obtiene esa operación desde asiento VENTA.
Sin esa reasignación, un reporte filtrado por venta omitiría un costo
tardío real porque su origen administrativo era compra.

GOLD-02 comienza con10 animales a1000, adquisición10000 y gasto1000
distribuido:1100 por posición. Vender4 a1600 produce6400, costo4400 y
gasto directo200; margen1800,
\texorpdfstring{\softtext{inventario6600.}}{inventario6600.} GOLD-06
añade500 a la compra,50 por posición: costo vendido4600, inventario6900
y \texorpdfstring{\softtext{margen1600.}}{margen1600.} Bonificar200,20
por posición, reduce costo vendido a4520, inventario6780 y aumenta
margen a1680. Las deudas comerciales originales no cambian
automáticamente por esos gastos propios.

Una baja exige ADMIN, posición libre DISPONIBLE o
\texorpdfstring{\softtext{PENDIENTE\_RECEPCION,}}{PENDIENTE\_RECEPCION,}
motivo permitido y documento existente del scope. Se rechaza si hay
reserva vigente. Se toma costo capitalizado, se inserta asiento BAJA y
se cambia posición a BAJA con
\texorpdfstring{\softtext{motivo/cierre.}}{motivo/cierre.} GOLD-08
conserva costo1100 convertido en pérdida1100, inventario0 y, después de
un costo tardío50,
\texorpdfstring{\softtext{pérdida1150.}}{pérdida1150.} No hay venta ni
cuenta por cobrar inventada.

Si motivo es MUERTE y el nodo es A, se actualiza identidad central
dentro del contexto \texorpdfstring{\softtext{autorizado.}}{autorizado.}
Si es B, se crea
\texorpdfstring{\softtext{ANIMAL\_MUERTO}}{ANIMAL\_MUERTO} junto a
pérdida local y se entrega después. Una réplica que todavía no ve animal
MUERTO no permite eludir la baja local mediante nueva adquisición: el
comando y los estados conservan esa protección operativa.

\texorpdfstring{\softtext{anular\_venta}}{anular\_venta} no es cancelar
borrador. Exige venta CONFIRMADA, versión actual y ninguna entrega
asociada. Bloquea periodo, operación, posiciones y
\texorpdfstring{\softtext{liquidación.}}{liquidación.} Las aplicaciones
deben estar revertidas previamente: saldo debe coincidir con total
ajustado. El comando restaura a inventario el costo neto de cada
posición, crea contrasiento de venta y abona obligación para llevarla al
efecto correspondiente, conservando liquidación original.

En ANUL-01, costo inicial vendido1100 más tardío100 se anula antes de
entrega. El contrasiento restituye1200 a inventario y deja costo neto de
venta0. Tomar únicamente1100 habría ignorado el ajuste tardío; borrar la
venta habría perdido explicación de la
\texorpdfstring{\softtext{restitución.}}{restitución.} La posición
vuelve a DISPONIBLE y el ciclo no se transforma en otro ciclo de
\texorpdfstring{\softtext{adquisición.}}{adquisición.}

\textbf{Fuente:}
\texorpdfstring{\softtext{db/migrations/V13\_\_correccion\_composicion\_exacta.sql:235;}}{db/migrations/V13\_\_correccion\_composicion\_exacta.sql:235;}
\texorpdfstring{\softtext{db/migrations/V13\_\_correccion\_composicion\_exacta.sql:271;}}{db/migrations/V13\_\_correccion\_composicion\_exacta.sql:271;}
\texorpdfstring{\softtext{db/migrations/V13\_\_correccion\_composicion\_exacta.sql:331.}}{db/migrations/V13\_\_correccion\_composicion\_exacta.sql:331.}

\section{10. Periodos: elegir reconocimiento y conservar
historia}\label{periodos-elegir-reconocimiento-y-conservar-historia}

Una fecha documental no determina automáticamente el mes de
\texorpdfstring{\softtext{reconocimiento.}}{reconocimiento.} El sistema
distingue \texorpdfstring{\softtext{fecha\_operacion}}{fecha\_operacion}
o \texorpdfstring{\softtext{fecha\_economica,}}{fecha\_economica,}
\texorpdfstring{\softtext{fecha\_reconocimiento}}{fecha\_reconocimiento}
y tiempo técnico
\texorpdfstring{\softtext{registrado\_en.}}{registrado\_en.} Un
documento fechado en otro mes puede conocerse hoy y reconocerse
explícitamente en un periodo abierto. Cambiar un campo documental no
debe insertar retroactivamente costo en un periodo que ya se cerró.

\texorpdfstring{\softtext{crear\_periodo}}{crear\_periodo} exige ADMIN y
toma un advisory lock por
\texorpdfstring{\softtext{empresa/sede}}{empresa/sede} para serializar
la creación del \texorpdfstring{\softtext{calendario.}}{calendario.} El
primer mes se declara
\texorpdfstring{\softtext{explícitamente.}}{explícitamente.} Debe
comenzar en primer día del mes y terminar en último; un siguiente
periodo sólo es consecutivo y no puede coexistir con otro abierto.
EXCLUDE agrega protección frente a rangos
\texorpdfstring{\softtext{superpuestos.}}{superpuestos.} La aplicación
no crea automáticamente un mes al descubrir que falta.

fn\_periodo recibe e,s,f,pid: empresa, sede, fecha de reconocimiento e
identificador de periodo. Ambos f y pid son
\texorpdfstring{\softtext{obligatorios.}}{obligatorios.} SELECT FOR
UPDATE toma la fila exacta de E/S/id; debe existir, estar ABIERTO y
contener f entre \texorpdfstring{\softtext{desde/hasta.}}{desde/hasta.}
Devuelve pid para asignarlo al hecho.
\texorpdfstring{\softtext{fn\_periodo\_payload}}{fn\_periodo\_payload}
sólo adapta las propiedades JSON a esos argumentos, no elige una fecha
silenciosa para resolver un error.

\texorpdfstring{\softtext{cerrar\_periodo}}{cerrar\_periodo} bloquea esa
misma fila. Si costo tardío la bloqueó primero, el cierre espera y
después observa el costo confirmado; si cierre ganó, el costo espera y
al retomar encuentra CERRADO, falla BV025 y no deja asignación, asiento,
auditoría o intención nueva parcial. Una nueva intención puede elegir el
siguiente mes abierto. No se modifica la fecha de un intento antiguo
para hacerlo pasar como replay igual.

El inventario de corte suma asientos asociados a periodos cuyo hasta es
anterior o igual al periodo objetivo. No filtra simplemente
\texorpdfstring{\softtext{asiento.fecha\_economica}}{asiento.fecha\_economica}
contra la fecha final, porque en datos históricos ese nombre puede haber
transportado significados diferentes y un documento futuro no debe
excluir reconocimiento actual. La relación periodo\_id es la autoridad
del calendario de cada hecho.

Costo vendido y pérdida se suman por periodo
\texorpdfstring{\softtext{reconocido.}}{reconocido.} Cartera se lee del
saldo vivo mientras el lock y la política de único periodo abierto
impiden que otro reconocimiento posterior cambie el mismo corte
\texorpdfstring{\softtext{concurrentemente.}}{concurrentemente.} Esa
coherencia es local de sede; no promete un snapshot distribuido
simultáneo de A y B.

El cierre inserta
\texorpdfstring{\softtext{cierre\_periodo\_version}}{cierre\_periodo\_version}
con número siguiente, actor, motivo, corte, hash y totales. El cambio de
periodo a CERRADO pertenece a la misma
\texorpdfstring{\softtext{transacción.}}{transacción.} Reabrir sólo el
último cerrado sin periodo
\texorpdfstring{\softtext{posterior/abierto}}{posterior/abierto} crea
versión REAPERTURA y conserva las
\texorpdfstring{\softtext{anteriores.}}{anteriores.} No hace UPDATE de
un cierre aceptado para volver a presentar otras cifras.

\subsection{10.1 El reporte histórico desde
V12}\label{el-reporte-histuxf3rico-desde-v12}

Al cerrar se materializa
\texorpdfstring{\softtext{reporte\_margen}}{reporte\_margen} dentro de
totales. Una consulta de ese periodo cerrado sin búsqueda ni filtro de
operación puede devolver el snapshot y su
\texorpdfstring{\softtext{numero\_version\_cierre.}}{numero\_version\_cierre.}
La etiqueta CERRADO se utiliza sólo cuando no contiene faltantes. Un
cierre que aún tuvo revisión o gasto pendiente conserva PROVISIONAL,
aunque su calendario esté cerrado.

Un caso QA fijó margen90 al cerrar. Se reconoció después costo20 en el
siguiente periodo; el margen vivo bajó70, pero la consulta del cierre
anterior siguió devolviendo el mismo
\texorpdfstring{\softtext{snapshot90.}}{snapshot90.} No se cambia de
etiqueta para ocultar diferencia: son dos preguntas legítimas, el corte
anterior y el estado actual.

Los cierres anteriores a V12 que no contienen
\texorpdfstring{\softtext{reporte\_margen}}{reporte\_margen} no se
reconstruyen como supuesta historia aceptada. Devuelven PROVISIONAL y
\texorpdfstring{\softtext{CIERRE\_SIN\_REPORTE\_MARGEN\_HISTORICO.}}{CIERRE\_SIN\_REPORTE\_MARGEN\_HISTORICO.}
Esta rama se revisó en código; no se encontró un caso principal A
actualmente cerrado que permitiera atribuir una prueba física específica
de ese legado. El límite se mantiene
\texorpdfstring{\softtext{explícito.}}{explícito.}

Si se filtra por operación o q dentro de periodo cerrado, no se presenta
el total global guardado como si fuera el desglose filtrado. La consulta
explica
\texorpdfstring{\softtext{REPORTE\_FILTRADO\_SIN\_DESGLOSE\_HISTORICO}}{REPORTE\_FILTRADO\_SIN\_DESGLOSE\_HISTORICO}
y \texorpdfstring{\softtext{provisionalidad.}}{provisionalidad.} El
contrato de cartera actual no admite periodo\_id histórico: el backend
lo rechaza, en vez de ignorarlo o inventar filas históricas no
\texorpdfstring{\softtext{almacenadas.}}{almacenadas.}

\textbf{Fuentes:}
\texorpdfstring{\softtext{db/migrations/V2\_\_comandos.sql:55;}}{db/migrations/V2\_\_comandos.sql:55;}
\texorpdfstring{\softtext{db/migrations/V5\_\_corte\_por\_periodo\_reconocimiento.sql:1;}}{db/migrations/V5\_\_corte\_por\_periodo\_reconocimiento.sql:1;}
\texorpdfstring{\softtext{db/migrations/V12\_\_pesajes\_reportes\_canonicos.sql:12;}}{db/migrations/V12\_\_pesajes\_reportes\_canonicos.sql:12;}
\texorpdfstring{\softtext{db/migrations/V13\_\_correccion\_composicion\_exacta.sql:300.}}{db/migrations/V13\_\_correccion\_composicion\_exacta.sql:300.}

\section{11. Privilegios, SECURITY DEFINER y frontera de
confianza}\label{privilegios-security-definer-y-frontera-de-confianza}

Los roles PostgreSQL son cuentas técnicas, distintas de los roles
humanos de \texorpdfstring{\softtext{aplicación.}}{aplicación.}
bovi\_owner administra estructura y
\texorpdfstring{\softtext{bovi\_command}}{bovi\_command} posee funciones
privilegiadas; son NOLOGIN.
\texorpdfstring{\softtext{bovi\_migrator}}{bovi\_migrator} puede asumir
propiedad con privilegios administrativos separados. bovi\_app\_a y
bovi\_app\_b conectan a negocio con SELECT permitido y EXECUTE de
interfaces, sin
\texorpdfstring{\softtext{INSERT/UPDATE/DELETE/DDL}}{INSERT/UPDATE/DELETE/DDL}
directo. bovi\_auth lee
\texorpdfstring{\softtext{identidad/membresía}}{identidad/membresía} y
tiene DML específicamente sobre sesiones.
\texorpdfstring{\softtext{bovi\_worker}}{bovi\_worker} ejecuta funciones
de entrega y referencia de archivos. Monitor y backup tienen propósitos
distintos.

El rol técnico no inventa la persona
\texorpdfstring{\softtext{autorizada.}}{autorizada.} fn\_scope comprueba
p\_actor contra usuario y membresía vigentes. Una sesión válida en el
navegador no permite a un cuerpo elegir otra persona. Las cuentas SQL
compartidas pertenecen al servidor confiable; no se distribuyen al
navegador. El diseño no utiliza JWT directo de clientes sobre tablas ni
declara RLS multiempresa completa.

SECURITY DEFINER permite que app ejecute una función que inserta hechos
bajo derechos \texorpdfstring{\softtext{controlados.}}{controlados.} Su
seguridad depende de no exponer funciones generales de ejecución SQL,
comprobar scope y dueño, cualificar objetos y fijar
\texorpdfstring{\softtext{search\_path}}{search\_path} seguro. Si una
función simplemente aceptara tabla y sentencia arbitrarias, el NOLOGIN
del dueño no impediría abuso mediante EXECUTE.

PUBLIC se revoca para funciones de negocio. App recibe las firmas
\texorpdfstring{\softtext{necesarias.}}{necesarias.} El SELECT de
usuario se limita a columnas no sensibles; auth tiene lectura del hash
para su función legítima. La función de consulta elimina secretos del
JSON. Controles de API y restricciones SQL se complementan; conocer un
UUID no reemplaza membresía ni propiedad.

Los triggers protegen hechos confirmados frente a
\texorpdfstring{\softtext{UPDATE/DELETE,}}{UPDATE/DELETE,} incluso en
una simulación controlada con rol dueño de funciones. Esto demuestra
algo distinto de app sin DML. La guardia de contrato aceptado preserva
historia; la de pesaje permite únicamente transición
\texorpdfstring{\softtext{APROBADO→CORREGIDO}}{APROBADO→CORREGIDO} sin
cambiar sus otros valores. La de composición impide editar animales de
un ticket que ya no es PROPUESTA.

No se afirma que una cuenta superusuario carezca de capacidad
administrativa para alterar el sistema. Las pruebas se hacen con roles
definidos y contextos
\texorpdfstring{\softtext{controlados.}}{controlados.} PGAudit y JIT son
posibles conceptos de seguridad, pero este capítulo no declara que estén
instalados ni \texorpdfstring{\softtext{certificados.}}{certificados.}
La auditoría de negocio implementada es explícita en tablas y comandos.

\textbf{Fuentes:}
\texorpdfstring{\softtext{db/bootstrap/01\_roles.sql:1;}}{db/bootstrap/01\_roles.sql:1;}
\texorpdfstring{\softtext{db/migrations/V2\_\_comandos.sql:3;}}{db/migrations/V2\_\_comandos.sql:3;}
\texorpdfstring{\softtext{db/migrations/V2\_\_comandos.sql:29;}}{db/migrations/V2\_\_comandos.sql:29;}
\texorpdfstring{\softtext{db/migrations/V3\_\_consultas\_sesiones\_grants.sql:151.}}{db/migrations/V3\_\_consultas\_sesiones\_grants.sql:151.}

\section{12. Propietarios fijos, replicación y
fallos}\label{propietarios-fijos-replicaciuxf3n-y-fallos}

A escribe maestros y centraliza
\texorpdfstring{\softtext{autenticación.}}{autenticación.} Los hechos
locales se escriben en el nodo dueño de la sede. B no se convierte en
escritor de A por indisponibilidad, ni A confirma una venta de B para
continuar servicio. Pgpool trabaja raw con un backend fijo por puerto y
sin promoción \texorpdfstring{\softtext{automática.}}{automática.}

Las publicaciones bidireccionales replican tablas comerciales y filtran
hechos por \texorpdfstring{\softtext{nodo\_origen.}}{nodo\_origen.} Las
suscripciones utilizan origin=none para evitar reenviar cambios que
provienen de otro origen. Esa opción no resuelve conflictos si se
permite a ambos nodos escribir el mismo agregado. La disciplina de
propietario evita crear el conflicto; los constraints y conciliación
detectan inconsistencias
\texorpdfstring{\softtext{observables.}}{observables.}

El apply de suscripción no representa una nueva intención comercial
humana. Por eso no debe generar otra auditoría económica o volver a
despachar cada fila recibida como orden. Outbox conserva
\texorpdfstring{\softtext{propietario\_despacho}}{propietario\_despacho}
y el worker procesa hechos de su nodo. Inbox, claves idempotentes y
estados de entrega son locales y no se publican como hechos
\texorpdfstring{\softtext{comerciales.}}{comerciales.}

Una partición puede retrasar visibilidad de maestros o hechos. Se
informa fuente y antigüedad cuando se conoce; un timestamp agregado de
sincronización no demuestra que toda identidad o tarifa esté
\texorpdfstring{\softtext{actualizada.}}{actualizada.} La lectura de
réplica nunca transfiere propiedad de escritura. La conciliación revisa
huérfanos y relaciones balanceadas después de
\texorpdfstring{\softtext{convergencia.}}{convergencia.}

Si A cae, nuevas solicitudes autenticadas de A y B fallan503 porque la
autorización vigente es central. Una transacción B ya autorizada y en
ejecución puede concluir en su propio nodo. Son situaciones diferentes:
terminar un trabajo validado no equivale a admitir otro login o nueva
petición sin A. La caída de B no cambia su dueño ni reinicia
arbitrariamente el proceso para ocultarla.

Las pruebas registraron paradas físicas de contenedores propios
\texorpdfstring{\softtext{A/B/proveedor}}{A/B/proveedor} y reanudación
en finally. La partición de réplica se simuló deshabilitando stream
apply en QA, no cortando una red física. Es una evidencia útil con
alcance \texorpdfstring{\softtext{específico.}}{específico.} Dos
PostgreSQL en la misma laptop no son alta disponibilidad frente a
pérdida del host.

\section{13. Muerte B→A: outbox, inbox y lease
durable}\label{muerte-ba-outbox-inbox-y-lease-durable}

B reconoce baja, pérdida y evento
\texorpdfstring{\softtext{ANIMAL\_MUERTO}}{ANIMAL\_MUERTO} en una
transacción local. No necesita mantener una transacción financiera
abierta esperando que A acepte red. El hecho económico ya queda íntegro;
el evento durable expresa trabajo pendiente para actualizar identidad
central.

\texorpdfstring{\softtext{reclamar\_entregas}}{reclamar\_entregas}
recibe p\_worker UUID y límite de1-100. Selecciona pendientes de
despacho local, sin acuse, con lease
\texorpdfstring{\softtext{vencida/ausente}}{vencida/ausente} y próximo
intento \texorpdfstring{\softtext{habilitado.}}{habilitado.} Usa FOR
UPDATE SKIP LOCKED para que otro trabajador pueda procesar diferentes
\texorpdfstring{\softtext{pendientes.}}{pendientes.} Guarda
\texorpdfstring{\softtext{lease\_token=p\_worker}}{lease\_token=p\_worker}
y \texorpdfstring{\softtext{lease\_vence\_en}}{lease\_vence\_en} treinta
segundos después, e incrementa intentos.

La transacción de reclamo termina antes de conectar con A. El UUID
identifica al trabajador de ese reclamo; no es actor comercial de una
venta. A consume mediante
\texorpdfstring{\softtext{consumir\_muerte,}}{consumir\_muerte,} que
exige nodo A y empresa configurada, evento remoto convergente
\texorpdfstring{\softtext{ANIMAL\_MUERTO}}{ANIMAL\_MUERTO} de B y
payload idéntico al outbox observado. No acepta una orden arbitraria que
declare otro animal muerto.

La PK del inbox
\texorpdfstring{\softtext{IDENTIDAD\_CENTRAL/B/eventid}}{IDENTIDAD\_CENTRAL/B/eventid}
deduplica. Sólo el primer consumo actualiza animal de la sede
correspondiente y registra auditoría central. No crea otro asiento de
pérdida: la pérdida pertenece a B. Repetir el consumo devuelve estado
APLICADO. La convergencia del evento y su validación evitan reconocer
una identidad a partir de información remota que aún no se observó.

\texorpdfstring{\softtext{confirmar\_entrega}}{confirmar\_entrega}
verifica evento, destino,
\texorpdfstring{\softtext{lease\_token}}{lease\_token} y vigencia. Un
acuse del trabajador antiguo no confirma una entrega reclamada después
por otro. \texorpdfstring{\softtext{fallar\_entrega}}{fallar\_entrega}
guarda error saneado, próximo intento con demora lineal cinco segundos
por intento limitada a cinco minutos y limpia lease. No guarda conninfo
ni contraseña dentro del
\texorpdfstring{\softtext{diagnóstico.}}{diagnóstico.}

Si se pierde el acuse después de consumo A, B mantiene pendiente y
entrega otra vez. El inbox devuelve resultado y luego se confirma
entrega local. Si un worker abandona lease, otro puede reclamar al
vencer. Las pruebas simularon abandono dejando vencer la lease y acuse
perdido omitiendo la primera llamada; no se atribuye a esas pruebas
haber matado un proceso Java.

\texorpdfstring{\softtext{despachar\_outbox}}{despachar\_outbox} tiene
una finalidad DEMO local para otros eventos, con inbox y acuse locales.
No debe confundirse con garantizar entrega a una empresa externa real.
\texorpdfstring{\softtext{archivo\_referenciado}}{archivo\_referenciado}
permite al worker consultar si una huella documental sigue
\texorpdfstring{\softtext{comprometida.}}{comprometida.} La limpieza de
bytes exige respuesta de ambos propietarios; si uno falla, no se elimina
contenido basándose sólo en el otro.

\textbf{Fuente:}
\texorpdfstring{\softtext{db/migrations/V3\_\_consultas\_sesiones\_grants.sql:131.}}{db/migrations/V3\_\_consultas\_sesiones\_grants.sql:131.}

\section{14. FDW sin transacción financiera
remota}\label{fdw-sin-transacciuxf3n-financiera-remota}

El proveedor simulado conserva
\texorpdfstring{\softtext{tarifa\_transporte}}{tarifa\_transporte} con
precio numeric(14,4) y
\texorpdfstring{\softtext{ticket\_bascula}}{ticket\_bascula} con neto
\texorpdfstring{\softtext{numeric(14,3).}}{numeric(14,3).} El rol remoto
\texorpdfstring{\softtext{bovi\_fdw\_read}}{bovi\_fdw\_read} tiene
SELECT. En A, \texorpdfstring{\softtext{postgres\_fdw}}{postgres\_fdw}
proporciona \texorpdfstring{\softtext{ext\_proveedor}}{ext\_proveedor} y
el servidor prov\_sim. B no conserva vínculo FDW de proveedor desde V9.

El backend solicita
\texorpdfstring{\softtext{leer\_tarifas\_proveedor}}{leer\_tarifas\_proveedor}
en A mediante autocommit, incluso si la sede interesada es B. La función
comprueba actor y scope maestro, pagina con orden de ruta y construye
procedencia, unidad, versión, fecha, precisión original y Dinero a
centavos. Después llama
\texorpdfstring{\softtext{capturar\_tarifas\_proveedor}}{capturar\_tarifas\_proveedor}
en el dueño de S, también fuera de locks
\texorpdfstring{\softtext{financieros.}}{financieros.}

Capturar valida contexto E/S/actor, arreglo acotado, hash, fuente DEMO,
moneda, política y coherencia entre cuatro decimales originales y dos
\texorpdfstring{\softtext{canónicos.}}{canónicos.} Guarda precio a
centavos por el tipo físico14,2 y JSON de precisión original. Hash
comprueba integridad del contenido transportado; no es una firma que
convierta una entrada de navegador en fuente confiable. La API no
permite suministrar esa página interna como body público.

\texorpdfstring{\softtext{aceptar\_tarifa}}{aceptar\_tarifa} marca
aceptación en captura local sólo para ese actor y scope.
\texorpdfstring{\softtext{registrar\_gasto}}{registrar\_gasto} copia
snapshot aceptado y compara el monto local. Una fuente futura cambiada o
caída no altera gasto aceptado anterior. Si no se leyó una captura, no
se inventa una tarifa para que la compra continúe.

La frontera \texorpdfstring{\softtext{real12000.0050}}{real12000.0050}
del proveedor se redondea a
\texorpdfstring{\softtext{Dinero12000.01.}}{Dinero12000.01.} El JSON
\texorpdfstring{\softtext{conserva12000.0050}}{conserva12000.0050} y
política HALF\_UP\_2; el catálogo de captura14,2
\texorpdfstring{\softtext{almacena12000.01.}}{almacena12000.01.} Por eso
comparar \texorpdfstring{\softtext{gasto12000.01}}{gasto12000.01} con
captura funciona. Una sospecha inicial de comparación contra cuatro
decimales se descartó al verificar catálogo y ejecución, sin añadir una
migración \texorpdfstring{\softtext{redundante.}}{redundante.}

No se usa FDW para atomicidad monetaria o two phase commit de proveedor.
Si la red se mantiene dentro de un lock de pago o periodo, una fuente
lenta bloquearía negocio local y complicaría cierre. La lectura previa y
token aceptado reducen esa dependencia, con procedencia
\texorpdfstring{\softtext{explícita.}}{explícita.} El proveedor es DEMO,
no SINIIGA, SENASICA ni empresa comercial real.

\textbf{Fuente:}
\texorpdfstring{\softtext{db/migrations/V9\_\_lector\_a\_captura\_owner\_fecha.sql:2;}}{db/migrations/V9\_\_lector\_a\_captura\_owner\_fecha.sql:2;}
\texorpdfstring{\softtext{db/bootstrap/bootstrap.py:60.}}{db/bootstrap/bootstrap.py:60.}

\section{15. Consultas canónicas: qué responde cada
pregunta}\label{consultas-canuxf3nicas-quuxe9-responde-cada-pregunta}

\texorpdfstring{\softtext{bovi.consultar}}{bovi.consultar} recibe
e,s,a,recurso,params y retorna JSON. recurso no es una sentencia SQL del
navegador: el CASE decide tablas y ramas
\texorpdfstring{\softtext{permitidas.}}{permitidas.} La función vuelve a
validar scope. Las partes dinámicas provienen de una lista de recursos y
filtros definidos, mientras valores se pasan con USING. No debe
confundirse SQL dinámico interno controlado con permitir un WHERE
arbitrario enviado por cliente.

tab selecciona tabla o vista; filtro reúne predicados permitidos; sql
conserva sentencia formada. rows y detail guardan JSON; totaln es el
conteo completo, size tamaño y page número de página. idn identifica un
detalle; searchq es búsqueda recortada y pattern su versión escapada
para ILIKE literal. statefilter y typefilter no aceptan cualquier
palabra: el recurso define vocabulario admisible.

La paginación es base1. El desplazamiento se obtiene mediante:

\[
\mathrm{offset}=(\mathrm{pagina}-1)\,\mathrm{tamano}.
\]

Para página1 y tamaño25, offset0; para página2, offset25. ORDER BY
estable evita que un mismo corte distribuya filas aleatoriamente entre
páginas. El conteo total no es longitud de rows: una página de25 puede
corresponder a200 \texorpdfstring{\softtext{resultados.}}{resultados.}
Una solicitud de tamaño mayor que100 o página menor que1 se rechaza, no
se normaliza
\texorpdfstring{\softtext{silenciosamente.}}{silenciosamente.}

En q, caracteres porcentual y guion bajo se escapan para tratarse
\texorpdfstring{\softtext{literalmente.}}{literalmente.} El sistema no
convierte una búsqueda de usuario en comodín libre ni en SQL. Estado y
tipo se comprueban según operaciones, posiciones o pagos.
\texorpdfstring{\softtext{disponibles\_para\_venta}}{disponibles\_para\_venta}
no es una comparación de texto: el parámetro booleano verdadero filtra
DISPONIBLE, animal VIGENTE, sin reserva vigente y sin baja local. False
o ausencia incluye todas las posiciones del scope.

El detalle de operación une contratos, versión actual, lotes, selección
y tickets. Los tickets provienen de
\texorpdfstring{\softtext{fn\_pesaje\_json,}}{fn\_pesaje\_json,} por lo
que la composición histórica no se reconstruye con animales actualmente
presentes en un lote editado. El detalle de liquidación aporta conceptos
y aplicaciones canónicas y la contraparte necesaria para aplicar dinero
\texorpdfstring{\softtext{compatible.}}{compatible.} Documento retira
clave interna y añade revisiones exactas.

simulador consulta
\texorpdfstring{\softtext{fn\_calcular\_liquidacion}}{fn\_calcular\_liquidacion}
sin reconocer compra o venta. No es un comando monetario independiente
que el navegador pueda presentar como
\texorpdfstring{\softtext{confirmación.}}{confirmación.} Si falta peso
KG, no se fabrica una cantidad cero ni una liquidación principal. Un
resultado de cálculo es una previsualización; el comando vuelve a
validar evidencia, periodo, versión y reserva antes de confirmar.

\subsection{15.1 Variables del reporte y corte
observable}\label{variables-del-reporte-y-corte-observable}

fn\_reporte recibe recurso y params y construye metadata con corte\_en,
\texorpdfstring{\softtext{fuente\_consulta}}{fuente\_consulta} y
\texorpdfstring{\softtext{antiguedad\_replica\_segundos.}}{antiguedad\_replica\_segundos.}
ownernode procede de sucursal; today utiliza su
\texorpdfstring{\softtext{zona\_horaria.}}{zona\_horaria.}
\texorpdfstring{\softtext{statement\_timestamp}}{statement\_timestamp}
representa corte de esa sentencia. En el propietario, antigüedad se
expresa0; en una réplica donde no se dispone de medida individual fiable
queda NULL, no un número inventado a partir de una estadística global.

Cartera selecciona sólo saldos positivos. direction admite
\texorpdfstring{\softtext{POR\_COBRAR/POR\_PAGAR}}{POR\_COBRAR/POR\_PAGAR}
y overdue es booleano verdadero o falso. La propiedad vencida del JSON
debe ser booleana; una cadena ``true'' no equivale al tipo esperado.
Backend convierte el parámetro de query tipado antes de llegar a SQL.
vence\_el NULL produce
\texorpdfstring{\softtext{dias\_vencida}}{dias\_vencida} NULL: no se
declara mora de una obligación sin fecha.

La CTE filtered define el conjunto, paginated aplica
\texorpdfstring{\softtext{límite/offset}}{límite/offset} y el resultado
conserva filas, pagina, tamano y total. etiqueta ACTUALIZADO expresa
consulta actual. No hay una suma de primera página que parezca una
cartera completa. Historial de corte se conserva en cierre; un parámetro
nuevo que la consulta no implementa debe rechazarse en la interfaz, no
aceptarse \texorpdfstring{\softtext{silenciosamente.}}{silenciosamente.}

Margen utiliza CTE scoped, income, valuation, expense, relevant y flags.
scoped delimita operaciones; income suma ventas y ajustes con el periodo
propio de cada hecho; valuation suma costo y pérdida; expense suma
directos aprobados. relevant evita evaluar faltantes de objetos ajenos
al filtro. flags identifica gasto pendiente y revisión administrativa
pendiente.

Es importante que una venta de un mes anterior pueda recibir costo
reconocido hoy: el periodo del costo no se reemplaza por el periodo de
la \texorpdfstring{\softtext{liquidación.}}{liquidación.} A la vez,
consulta de un cierre histórico sin filtros recupera snapshot; no suma
datos tardíos de otro mes. Esta combinación distingue margen actual de
corte aceptado y evita una cifra financiera que cambia de significado
según la pantalla.

\textbf{Fuente:}
\texorpdfstring{\softtext{db/migrations/V12\_\_pesajes\_reportes\_canonicos.sql:12;}}{db/migrations/V12\_\_pesajes\_reportes\_canonicos.sql:12;}
\texorpdfstring{\softtext{db/migrations/V12\_\_pesajes\_reportes\_canonicos.sql:430.}}{db/migrations/V12\_\_pesajes\_reportes\_canonicos.sql:430.}

\section{16. Ejemplos seguros para estudiar sin modificar
hechos}\label{ejemplos-seguros-para-estudiar-sin-modificar-hechos}

Para verificar versión del servidor y contar tablas ordinarias, una
conexión de lectura autorizada puede utilizar consultas de catálogo. No
se incluye una contraseña en la línea de comandos ni se imprime
configuración de user mapping.

\begin{Shaded}
\begin{Highlighting}[]
\KeywordTok{SELECT}\NormalTok{ current\_database(), current\_setting(}\StringTok{\textquotesingle{}server\_version\textquotesingle{}}\NormalTok{);}

\KeywordTok{SELECT}\NormalTok{ n.nspname, }\FunctionTok{count}\NormalTok{(}\OperatorTok{*}\NormalTok{) }\KeywordTok{AS}\NormalTok{ tablas\_ordinarias}
\KeywordTok{FROM}\NormalTok{ pg\_class c}
\KeywordTok{JOIN}\NormalTok{ pg\_namespace n }\KeywordTok{ON}\NormalTok{ n.}\KeywordTok{oid}\OperatorTok{=}\NormalTok{c.relnamespace}
\KeywordTok{WHERE}\NormalTok{ c.relkind}\OperatorTok{=}\StringTok{\textquotesingle{}r\textquotesingle{}} \KeywordTok{AND}\NormalTok{ n.nspname }\KeywordTok{IN}\NormalTok{ (}\StringTok{\textquotesingle{}bovi\textquotesingle{}}\NormalTok{,}\StringTok{\textquotesingle{}infra\textquotesingle{}}\NormalTok{)}
\KeywordTok{GROUP} \KeywordTok{BY}\NormalTok{ n.nspname}
\KeywordTok{ORDER} \KeywordTok{BY}\NormalTok{ n.nspname;}

\KeywordTok{SELECT} \FunctionTok{count}\NormalTok{(}\OperatorTok{*}\NormalTok{) }\KeywordTok{AS}\NormalTok{ claves\_externas}
\KeywordTok{FROM}\NormalTok{ pg\_constraint k}
\KeywordTok{JOIN}\NormalTok{ pg\_namespace n }\KeywordTok{ON}\NormalTok{ n.}\KeywordTok{oid}\OperatorTok{=}\NormalTok{k.connamespace}
\KeywordTok{WHERE}\NormalTok{ k.contype}\OperatorTok{=}\StringTok{\textquotesingle{}f\textquotesingle{}} \KeywordTok{AND}\NormalTok{ n.nspname }\KeywordTok{IN}\NormalTok{ (}\StringTok{\textquotesingle{}bovi\textquotesingle{}}\NormalTok{,}\StringTok{\textquotesingle{}infra\textquotesingle{}}\NormalTok{);}
\end{Highlighting}
\end{Shaded}

Estas consultas no leen filas
\texorpdfstring{\softtext{empresariales.}}{empresariales.} Su ejecución
concreta en el capítulo
\texorpdfstring{\softtext{obtuvo46/86}}{obtuvo46/86} en ambos nodos. Si
otra instalación devuelve cantidades distintas, se identifica base,
nivel de migración y esquemas incluidos antes de concluir que existe
error. No se borra una tabla para hacer coincidir una cifra
\texorpdfstring{\softtext{documental.}}{documental.}

Para estudiar una definición privilegiada, se inspeccionan propiedades
de función, sin copiar
\texorpdfstring{\softtext{credenciales.}}{credenciales.} Una consulta a
pg\_proc puede revisar prosecdef y configuración del
\texorpdfstring{\softtext{search\_path.}}{search\_path.} La definición
debe contrastarse con su última
\texorpdfstring{\softtext{migración.}}{migración.} Enumerar nombres por
regex estático no descubre todos los triggers creados en bloques DO
dinámicos; por eso el atlas del manual advierte ese límite y el catálogo
completa la \texorpdfstring{\softtext{observación.}}{observación.}

Para ensayar un comando de negocio, la alternativa apropiada es un
escenario QA propio con identidad y periodo conocidos. Se comienza
BEGIN, se llama la firma con parámetros de servidor autenticado y se
comprueban cifras fijas; después ROLLBACK. No se ejecuta ese ejemplo
sobre main, ni se utiliza un hash inventado como si fuera la
canonicalización de \texorpdfstring{\softtext{producción.}}{producción.}
El harness gold.py muestra una práctica concreta de rollback y
savepoints para negativas.

Un SAVEPOINT permite probar una excepción sin perder toda la transacción
de fixture. Después del error se hace ROLLBACK TO SAVEPOINT y se
verifica que el hecho no apareció. Eso no significa que la aplicación
deba atrapar cualquier error y confirmar el resto: el comando económico
conserva atomicidad de su
\texorpdfstring{\softtext{intención.}}{intención.} La técnica pertenece
a una prueba \texorpdfstring{\softtext{controlada.}}{controlada.}

\section{17. Bootstrap, despliegue local y evidencia de
rendimiento}\label{bootstrap-despliegue-local-y-evidencia-de-rendimiento}

bootstrap prepara roles, identidad local, maestros DEMO,
\texorpdfstring{\softtext{publicaciones/suscripciones}}{publicaciones/suscripciones}
y proveedor simulado. Los datos financieros de la semilla posterior se
construyen por comandos, no insertando totales deseados directamente en
liquidacion o
\texorpdfstring{\softtext{asiento\_valoracion.}}{asiento\_valoracion.}
La inicialización organizativa no es una venta demostrada ni un
bootstrap comercial de
\texorpdfstring{\softtext{producción.}}{producción.}

Las guardas iniciales rechazan
\texorpdfstring{\softtext{BOVI\_ENV=production}}{BOVI\_ENV=production} o
perfil production antes de cualquier acceso o mutación de
\texorpdfstring{\softtext{bootstrap/seed}}{bootstrap/seed} DEMO.
Configuración preparada para una VM no ejecuta ese aprovisionamiento ni
activa TLS real. Roles comerciales, empresa real, inventario inicial
conciliado y secretos de producción requieren otra preparación
\texorpdfstring{\softtext{autorizada.}}{autorizada.}

La topología académica usa
\texorpdfstring{\softtext{PostgreSQL18.6}}{PostgreSQL18.6} y
\texorpdfstring{\softtext{Pgpool4.7.3}}{Pgpool4.7.3} raw, en loopback
con puertos propios. A/B tienen límites de memoria de1GiB y
\texorpdfstring{\softtext{proveedor512MiB;}}{proveedor512MiB;}
\texorpdfstring{\softtext{shared\_buffers128MiB,}}{shared\_buffers128MiB,}
\texorpdfstring{\softtext{work\_mem4MiB,}}{work\_mem4MiB,}
\texorpdfstring{\softtext{maintenance\_work\_mem64MiB}}{maintenance\_work\_mem64MiB}
y \texorpdfstring{\softtext{max\_connections50}}{max\_connections50}
acotan consumo. Cada Pgpool tiene16 procesos con max\_pool1. Conexiones
de autenticación A, replicación, migración y administración consumen
presupuesto adicional al
\texorpdfstring{\softtext{negocio/worker.}}{negocio/worker.}

Estos límites no prueban capacidad comercial. El benchmark propio
utilizó QA separada,200 animales,300 operaciones,300 gastos,100 pagos
y100 entregas creados mediante1901 comandos. Repetir1901 conservó
identidades, resultados y conteos. Después se midieron800 lecturas
canónicas por Pgpool con1 y4 clientes,100 warmups y4EXPLAIN ANALYZE
BUFFERS.

La máquina era Intel i7-13620H con16 lógicos y
\texorpdfstring{\softtext{unos15.7GiB}}{unos15.7GiB} reportados en el
artefacto. Caché caliente, dataset sintético moderado y ausencia de
HTTP/carga escritora delimitan la medición. Los planes del wrapper
PL/pgSQL no muestran todos los planes internos. Se reportan percentiles
reales, pero no se extrapola un SLO del piloto, un número de clientes
contratados o escalabilidad frente a miles de escrituras
\texorpdfstring{\softtext{simultáneas.}}{simultáneas.}

\textbf{Fuentes:}
\texorpdfstring{\softtext{db/bootstrap/bootstrap.py:1;}}{db/bootstrap/bootstrap.py:1;}
\texorpdfstring{\softtext{db/bootstrap/01\_roles.sql:1;}}{db/bootstrap/01\_roles.sql:1;}
\texorpdfstring{\softtext{db/tests/benchmark-report.md:1;}}{db/tests/benchmark-report.md:1;}
\texorpdfstring{\softtext{infra/compose/compose.academic.yml:1.}}{infra/compose/compose.academic.yml:1.}

\section{18. Glosario SQL orientado a esta
implementación}\label{glosario-sql-orientado-a-esta-implementaciuxf3n}

\textbf{Transacción local.} Unidad que confirma o revierte todos los
efectos del comando en un
\texorpdfstring{\softtext{propietario.}}{propietario.} No incluye
mágicamente otro PostgreSQL o un archivo remoto. Outbox y staging
documental resuelven límites externos sin afirmar atomicidad distribuida
\texorpdfstring{\softtext{inexistente.}}{inexistente.}

\textbf{FOR UPDATE.} Bloquea filas seleccionadas frente a cambios
\texorpdfstring{\softtext{incompatibles.}}{incompatibles.} Debe
utilizarse antes de leer condiciones agregadas que luego se
\texorpdfstring{\softtext{materializan.}}{materializan.} No comprueba
versión del cliente por sí mismo:
\texorpdfstring{\softtext{version\_lock}}{version\_lock} exige
comparación explícita bajo el lock.

\textbf{SKIP LOCKED.} Permite que un trabajador omita pendientes que
otro mantiene \texorpdfstring{\softtext{bloqueados.}}{bloqueados.} Es
útil para colas, no para omitir silenciosamente una posición que debía
formar parte de una venta.

\textbf{Advisory transaction lock.} Lock por clave de aplicación que
dura la \texorpdfstring{\softtext{transacción.}}{transacción.} En
\texorpdfstring{\softtext{crear\_periodo}}{crear\_periodo} serializa
calendario por \texorpdfstring{\softtext{empresa/sede.}}{empresa/sede.}
No agrega una FK ni vuelve correctos datos que no cumplen
\texorpdfstring{\softtext{predicados.}}{predicados.}

\textbf{\%ROWTYPE y record.} El primero utiliza estructura de una tabla;
el segundo adopta resultados variables. Un registro cargado puede estar
obsoleto si se leyó sin lock. El tipo no equivale a autorización o
vigencia.

\textbf{JSONB.} Representación consultable de contexto, snapshot y
respuesta. No convierte payload de navegador en fuente confiable. Se
valida forma, alcance y valores; un hash no demuestra autenticidad
externa si quien lo envía también puede
\texorpdfstring{\softtext{fabricarlo.}}{fabricarlo.}

\textbf{RETURN NEXT.} Emite una fila de una función que retorna
conjunto, como reparto. Permite que cada posición reciba base e importe
\texorpdfstring{\softtext{explícitos.}}{explícitos.} El consumidor
todavía necesita garantizar que todas las cuotas se materializan en una
misma \texorpdfstring{\softtext{transacción.}}{transacción.}

\textbf{ON CONFLICT.} Define respuesta al conflicto de unicidad. En
idempotencia reclama o encuentra intención; en contador incrementa
secuencia; en inbox evita segundo consumo. No es una instrucción
universal de ignorar errores
\texorpdfstring{\softtext{económicos.}}{económicos.}

\textbf{Columna generada.} Valor derivado por el motor de otras
columnas, como delta de
\texorpdfstring{\softtext{valoración.}}{valoración.} Protege contra tres
signos contradictorios enviados por cliente, pero la clase y el sentido
todavía requieren vocabulario y autorizaciones correctos.

\textbf{Snapshot.} Información conocida y materializada en un acto. Un
snapshot de contraparte no sigue futuros renombres; un cierre no
incorpora automáticamente costos de meses
\texorpdfstring{\softtext{posteriores.}}{posteriores.} Esa estabilidad
explica por qué no basta consultar maestros actuales al imprimir
historia.

\textbf{P0001 y DETAIL.} Las reglas de negocio utilizan SQLSTATE P0001
con JSON público que contiene código BV, regla\_id y campos. BV412 o
BV405 no son SQLSTATE
\texorpdfstring{\softtext{inventados.}}{inventados.} Errores
\texorpdfstring{\softtext{nativos23503/23505/42501}}{nativos23503/23505/42501}
conservan su mecanismo y se traducen controladamente en API, sin exponer
SQL ni secretos.

\textbf{\texorpdfstring{\softtext{Idempotencia.}}{Idempotencia.}} Misma
intención autorizada y mismo contenido produce respuesta anterior;
contenido distinto bajo la misma clave genera conflicto. No significa
ejecutar dos veces y esperar casualmente que el estado quede igual.

\textbf{Reverso.} Nuevo hecho que contrarresta uno anterior y conserva
\texorpdfstring{\softtext{referencia.}}{referencia.} No borra el
original ni oculta el periodo en que se decidió corregir. Su unicidad
impide anular repetidamente el mismo efecto.

\section{19. Qué quedó comprobado y qué no debe
prometerse}\label{quuxe9-queduxf3-comprobado-y-quuxe9-no-debe-prometerse}

La evidencia local incluye157 llamadas de GOLD fijo,22 guardas de
versiones,66 comandos QA de
\texorpdfstring{\softtext{tickets/reportes/cierre,}}{tickets/reportes/cierre,}
carreras reales con dos conexiones y barreras,
\texorpdfstring{\softtext{roles/constraints/triggers,}}{roles/constraints/triggers,}
replicación, proveedor y fallos. Backend registró24 pruebas:15
unitarias,8HTTP reales por Caddy
\texorpdfstring{\softtext{y1PostgreSQL18.6}}{y1PostgreSQL18.6} limpio
con \texorpdfstring{\softtext{migraciones1-13.}}{migraciones1-13.} UI
registró16 casos integrados con límites propios. Estas cifras no son
porcentajes de cobertura universal.

El restore V13 se realizó en proyecto aislado y preservó firmas,
archivos, deduplicación, inbox y conexiones reconstruidas; invalidó
sesiones y probó compra nueva1000, replay con ETag original, lookup y
descarga con hash. Pausa8.038 s y
\texorpdfstring{\softtext{recuperación28.616}}{recuperación28.616} s son
medidas de ese ensayo en el mismo host. No prueban PITR, copia externa o
recuperación frente a pérdida física de la laptop.

La lectura del catálogo hecha para este capítulo confirma estado
estructural actual de A/B sin nuevas
\texorpdfstring{\softtext{mutaciones.}}{mutaciones.} No repite fallos o
carreras ni atribuye una ejecución reciente a logs
\texorpdfstring{\softtext{anteriores.}}{anteriores.} El propósito
pedagógico es enseñar cómo verificar una afirmación: la tabla dice dónde
vive el hecho, la restricción qué combinación protege, el comando cuándo
materializa, el test qué caso ejercitó y el artefacto quién lo ejecutó.

Una defensa sólida distingue garantía estructural, garantía
transaccional, permiso técnico, autorización humana y resultado
observado. ``Existe FK'' no demuestra mismo comprador; ``hay lock'' no
demuestra If-Match; ``replica'' no demuestra alta disponibilidad; ``pasó
un GOLD'' no demuestra todas las combinaciones; ``se preparó
production'' no demuestra
\texorpdfstring{\softtext{despliegue.}}{despliegue.} Mantener esas
diferencias hace el manual útil para estudiar el producto real y para
decidir responsablemente su siguiente
\texorpdfstring{\softtext{incremento.}}{incremento.}

\section{20. Cuaderno de oráculos fijos para
repasar}\label{cuaderno-de-oruxe1culos-fijos-para-repasar}

Un oráculo independiente dice qué resultado se esperaba antes de
ejecutar la \texorpdfstring{\softtext{implementación.}}{implementación.}
No se recalcula la respuesta llamando a la misma función bajo prueba.
Los siguientes valores proceden de los casos GOLD documentados y
permiten comprobar si se está confundiendo obligación, costo, dinero o
\texorpdfstring{\softtext{movimiento.}}{movimiento.} Cada cifra debe
explicarse también por sus filas fuente; coincidir accidentalmente en el
total no prueba un historial correcto.

{\def\LTcaptype{none} % do not increment counter
\begin{longtable}[]{@{}
  >{\raggedright\arraybackslash}p{(\linewidth - 4\tabcolsep) * \real{0.3333}}
  >{\raggedright\arraybackslash}p{(\linewidth - 4\tabcolsep) * \real{0.3333}}
  >{\raggedright\arraybackslash}p{(\linewidth - 4\tabcolsep) * \real{0.3333}}@{}}
\toprule\noalign{}
\begin{minipage}[b]{\linewidth}\raggedright
Caso
\end{minipage} & \begin{minipage}[b]{\linewidth}\raggedright
Esperado fijo
\end{minipage} & \begin{minipage}[b]{\linewidth}\raggedright
Qué distinción verifica
\end{minipage} \\
\midrule\noalign{}
\endhead
\bottomrule\noalign{}
\endlastfoot
GOLD-01 & Costo252000; margen34000; saldo244000;
\texorpdfstring{\softtext{variación100kg/2\%}}{variación100kg/2\%} &
Peso y deducción contractual
\texorpdfstring{\softtext{independientes.}}{independientes.} \\
GOLD-02 & Costo4400; margen1800; inventario6600 & Venta parcial conserva
ciclos restantes. \\
GOLD-03 &
\texorpdfstring{\softtext{0.34+0.33+0.33=1.00}}{0.34+0.33+0.33=1.00} &
Residuo de centavos por identidad estable. \\
GOLD-04A &
\texorpdfstring{\softtext{Obligación1000,pago600,aplicaciones400+200→saldo400,disponible0}}{Obligación1000,pago600,aplicaciones400+200→saldo400,disponible0}
& Parciales y replay sin
\texorpdfstring{\softtext{sobreaplicación.}}{sobreaplicación.} \\
GOLD-04B &
\texorpdfstring{\softtext{Obligación1000,pago1200,aplica1000→saldo0,disponible200}}{Obligación1000,pago1200,aplica1000→saldo0,disponible200}
& Anticipo no es otra venta. \\
GOLD-05 & Recompra genera ciclo2, nunca dos posiciones activas &
Identidad animal separada de ciclo. \\
GOLD-06 &
\texorpdfstring{\softtext{Tardío4600/6900/1600;}}{Tardío4600/6900/1600;}
\texorpdfstring{\softtext{bonificación4520/6780/1680}}{bonificación4520/6780/1680}
& Costo
\texorpdfstring{\softtext{vendido/inventario/margen}}{vendido/inventario/margen}
sin cambiar deuda. \\
GOLD-07 &
\texorpdfstring{\softtext{Reverso1000,reaplica900,abono100→saldo0,disponible100;}}{Reverso1000,reaplica900,abono100→saldo0,disponible100;}
\texorpdfstring{\softtext{reembolso→0}}{reembolso→0} & Corrección pagada
atómica y caja separada. \\
GOLD-08 &
\texorpdfstring{\softtext{Inventario0,pérdida1100→1150}}{Inventario0,pérdida1100→1150}
& Baja no genera ingreso ni cartera. \\
GOLD-09 & Ciclo anterior conserva costo y nuevo ciclo obtiene propio &
Historia no se recicla. \\
GOLD-10 & Costo entra antes del corte o
\texorpdfstring{\softtext{fallaBV025sin}}{fallaBV025sin} efecto &
Periodo como lock común de
\texorpdfstring{\softtext{reconocimiento.}}{reconocimiento.} \\
ANUL-01 & Inventario restituido1200,costo neto de venta0 & Anulación
incorpora costo tardío previo. \\
\end{longtable}
}

En GOLD-04A, la primera aplicación400 deja saldo600 y
\texorpdfstring{\softtext{disponible200.}}{disponible200.} La segunda200
deja saldo400 y \texorpdfstring{\softtext{disponible0.}}{disponible0.}
Repetir la segunda con la misma clave devuelve el resultado anterior sin
reducir nuevamente saldo. Una nueva intención por0.01 no puede pasar,
porque no hay \texorpdfstring{\softtext{disponible.}}{disponible.} No
sería suficiente comprobar únicamente que la obligación todavía tiene
saldo400: deben cumplirse ambos límites.

En GOLD-04B sobra200 después de pagar totalmente la
\texorpdfstring{\softtext{obligación.}}{obligación.} Ese importe se
conserva como disponible del pago, no como margen ni ingreso de otra
\texorpdfstring{\softtext{operación.}}{operación.} Puede aplicarse a
otra obligación compatible o reembolsarse con su propio hecho de caja y
periodo. Un reporte que sumara dinero recibido como venta mostraría1200
en vez de1000 y mezclaría anticipo con
\texorpdfstring{\softtext{reconocimiento.}}{reconocimiento.}

En GOLD-05 y09, entregar cierra el ciclo anterior y permite una nueva
adquisición del mismo animal vigente. La historia del primer ciclo
conserva costo y \texorpdfstring{\softtext{operación.}}{operación.} Las
condiciones actuales del segundo trato no recalculan el primero. La
unicidad parcial impide dos ciclos activos; no impide conservar dos
filas históricas cuando uno está cerrado.

En GOLD-10 importa ensayar ambos órdenes de la carrera, no dormir unos
segundos y suponer
\texorpdfstring{\softtext{simultaneidad.}}{simultaneidad.} Dos
conexiones y barrera observan qué lock bloquea al otro. Si costo ganó,
el cierre refleja el importe; si cierre ganó, el reconocimiento
posterior falla y rollback conserva ausencia de efectos parciales. El
esperado es una alternativa coherente por orden, no una cifra fija única
independiente de quién obtuvo el lock.

\textbf{Fuentes:}
\texorpdfstring{\softtext{docs/domain/states-and-calculations.md:61;}}{docs/domain/states-and-calculations.md:61;}
\texorpdfstring{\softtext{db/tests/gold.py:1;}}{db/tests/gold.py:1;}
\texorpdfstring{\softtext{db/tests/closure\_order.py:1;}}{db/tests/closure\_order.py:1;}
\texorpdfstring{\softtext{db/tests/weighing\_reports.py:1.}}{db/tests/weighing\_reports.py:1.}

\chapter{Backend, comunicación y seguridad: cómo funciona el programa
Java
real}\label{backend-comunicaciuxf3n-y-seguridad-cuxf3mo-funciona-el-programa-java-real}

\section{Alcance de este capítulo y forma de
leerlo}\label{alcance-de-este-capuxedtulo-y-forma-de-leerlo}

Este capítulo explica el backend que existe en el repositorio de
\texorpdfstring{\softtext{BoviCierre,}}{BoviCierre,} a partir de una
lectura completa de sus 24 archivos Java de producción, sus dos archivos
de configuración de aplicación, el POM de Maven y el contrato OpenAPI
0.2.0. La lectura se realizó el 12 de octubre de 2026. Las pruebas que
se citan al final son ejecuciones registradas el día anterior; escribir
este capítulo no significa que se hayan vuelto a ejecutar servicios,
comandos económicos o pruebas.

El lector puede seguir tres niveles. Primero, comprender la circulación
de una solicitud: navegador, filtros, controlador, autorización,
conexión propietaria, función SQL y respuesta. Segundo, estudiar cada
clase y sus métodos. Tercero, reconstruir las secuencias completas de
sesión, confirmación y archivo privado. Los nombres conservan la
ortografía del código para que sea posible buscarlos con el editor. Las
referencias tienen el formato ruta:línea; indican el comienzo del método
o bloque pertinente, no una afirmación de que todo el comportamiento
cabe en una sola línea.

Aquí, \textbf{campo} significa dato conservado dentro de un objeto Java;
\textbf{variable local}, dato temporal de una llamada;
\textbf{argumento}, valor que recibe un método; \textbf{salida}, su
resultado o \texorpdfstring{\softtext{excepción.}}{excepción.} Un objeto
de Spring registrado como servicio o controlador suele compartirse entre
solicitudes, mientras los argumentos y variables locales pertenecen a
cada \texorpdfstring{\softtext{ejecución.}}{ejecución.} Por ello importa
distinguir las dependencias finales de los contadores concurrentes y de
los objetos \texorpdfstring{\softtext{temporales.}}{temporales.}

Este texto también diferencia cuatro objetos que se confunden
fácilmente: usuario, empresa, sede y nodo. El usuario es una identidad
\texorpdfstring{\softtext{autenticada.}}{autenticada.} La empresa es el
ámbito empresarial configurado en el
\texorpdfstring{\softtext{despliegue.}}{despliegue.} La sede es la
unidad autorizada donde ocurre una
\texorpdfstring{\softtext{operación.}}{operación.} El nodo A o B es el
PostgreSQL que posee la capacidad de escribir un agregado. Solicitar una
sede B no concede acceso a B: primero deben comprobarse usuario y
\texorpdfstring{\softtext{membresía.}}{membresía.}

\section{Qué está implementado y qué corresponde al
diseño}\label{quuxe9-estuxe1-implementado-y-quuxe9-corresponde-al-diseuxf1o}

La aplicación ejecutable usa Java 21, Spring Boot 4.1.1, Spring MVC,
Spring Security, JDBC y Spring Session JDBC. Es un proceso Java con
componentes distribuidos en paquetes. No hay un microservicio
independiente por venta, pago o documento. Dentro del mismo proceso, los
controladores reciben HTTP,
\texorpdfstring{\softtext{AccessService}}{AccessService} consulta
autorización, \texorpdfstring{\softtext{CommandService}}{CommandService}
selecciona el propietario y PostgreSQL ejecuta las reglas
\texorpdfstring{\softtext{económicas.}}{económicas.}

El documento de arquitectura plantea módulos como identidad, inventario,
pagos y liquidaciones, además de carpetas domain, application y
adapters. En los archivos inspeccionados no existen esas capas completas
para cada módulo. Los paquetes reales son acceso, operaciones,
evidencia,
\texorpdfstring{\softtext{evidencia/application,}}{evidencia/application,}
integraciones, reportes y shared.
\texorpdfstring{\softtext{CoreController}}{CoreController} concentra
varias familias de rutas y
\texorpdfstring{\softtext{CommandService}}{CommandService} concentra la
comunicación canónica JDBC. Presentar el diagrama planeado como una
lista de clases ya programadas sería
\texorpdfstring{\softtext{incorrecto.}}{incorrecto.}

En particular,
\textbf{\texorpdfstring{\softtext{OwnerDatabaseRouter}}{OwnerDatabaseRouter}}
y
\textbf{\texorpdfstring{\softtext{BranchTransactionRunner}}{BranchTransactionRunner}}
aparecen en el diseño, pero no como clases del backend actual. Su
función se realiza de manera compacta en
\texorpdfstring{\softtext{CommandService}}{CommandService} mediante
mapas de \texorpdfstring{\softtext{JdbcTemplate}}{JdbcTemplate} y
\texorpdfstring{\softtext{TransactionTemplate.}}{TransactionTemplate.}
Tampoco existen clases Java Money o PesoKg en
\texorpdfstring{\softtext{producción.}}{producción.} La representación
monetaria que circula por los controladores es JSON con cadenas; los
importes confirmados se calculan en
\texorpdfstring{\softtext{PostgreSQL.}}{PostgreSQL.} La norma del
proyecto pide NUMERIC y BigDecimal para aritmética exacta, pero este
código Java no contiene una segunda calculadora financiera basada en
\texorpdfstring{\softtext{BigDecimal.}}{BigDecimal.}

Sí existe un puerto explícito:
\texorpdfstring{\softtext{DocumentStorage.}}{DocumentStorage.}
\texorpdfstring{\softtext{DocumentController}}{DocumentController}
depende de esa interfaz y
\texorpdfstring{\softtext{PrivateDocuments}}{PrivateDocuments}
implementa el almacenamiento local. Es una separación concreta,
comprobable en el constructor y en los métodos invocados; no debe
extrapolarse a puertos inexistentes para todas las demás
\texorpdfstring{\softtext{operaciones.}}{operaciones.}

La consecuencia para quien exponga el proyecto es sencilla: se puede
explicar la arquitectura objetivo y después mostrar la implementación
compacta que la materializa
\texorpdfstring{\softtext{parcialmente.}}{parcialmente.} La
autorización, elección del propietario, transacción local, idempotencia,
almacenamiento privado y trabajador existen. Una separación exhaustiva
de dominio, casos de uso y adaptadores por cada módulo es trabajo
adicional.

\section{Dependencias, configuración y entrada del
proceso}\label{dependencias-configuraciuxf3n-y-entrada-del-proceso}

\subsection{Qué aporta el POM}\label{quuxe9-aporta-el-pom}

El archivo \texorpdfstring{\softcode{backend/pom.xml}}{backend/pom.xml}
fija spring-boot-starter-parent 4.1.1 y
\texorpdfstring{\softtext{java.version}}{java.version} 21. El parent
administra un conjunto compatible de dependencias; no aparece una
versión independiente escrita junto a cada starter. La versión del
proyecto es \texorpdfstring{\softtext{0.1.0-SNAPSHOT}}{0.1.0-SNAPSHOT} y
el grupo es \texorpdfstring{\softtext{mx.bovicierre.}}{mx.bovicierre.}
El plugin spring-boot-maven-plugin prepara el empaquetado de la
\texorpdfstring{\softtext{aplicación.}}{aplicación.}

{\def\LTcaptype{none} % do not increment counter
\begin{longtable}[]{@{}
  >{\raggedright\arraybackslash}p{(\linewidth - 4\tabcolsep) * \real{0.3333}}
  >{\raggedright\arraybackslash}p{(\linewidth - 4\tabcolsep) * \real{0.3333}}
  >{\raggedright\arraybackslash}p{(\linewidth - 4\tabcolsep) * \real{0.3333}}@{}}
\toprule\noalign{}
\begin{minipage}[b]{\linewidth}\raggedright
Dependencia
\end{minipage} & \begin{minipage}[b]{\linewidth}\raggedright
Función real en este backend
\end{minipage} & \begin{minipage}[b]{\linewidth}\raggedright
Alcance
\end{minipage} \\
\midrule\noalign{}
\endhead
\bottomrule\noalign{}
\endlastfoot
spring-boot-starter-webmvc & Anotaciones HTTP, servidor y
\texorpdfstring{\softtext{ResponseEntity}}{ResponseEntity} &
Producción \\
spring-boot-starter-validation & Validación del record Login y errores
de argumentos & Producción \\
spring-boot-starter-security & Filtros, contexto de sesión, CSRF y
BCrypt & Producción \\
spring-boot-starter-jdbc &
\texorpdfstring{\softtext{JdbcTemplate,}}{JdbcTemplate,}
\texorpdfstring{\softtext{DataSourceTransactionManager}}{DataSourceTransactionManager}
y acceso JDBC & Producción \\
spring-boot-starter-session-jdbc & Persistencia de sesiones en
PostgreSQL & Producción \\
postgresql & Driver JDBC y lectura estructurada de errores PostgreSQL &
Producción \\
\texorpdfstring{\softtext{com.google.zxing:core}}{com.google.zxing:core}
3.5.4 & Construcción local de la matriz QR & Producción \\
starters de pruebas, security-test y testcontainers-postgresql & Pruebas
Java, HTTP y PostgreSQL real & Test \\
\end{longtable}
}

El POM de runtime no añade Flyway. Las migraciones son una actividad
separada. Tampoco incorpora una biblioteca de renderizado PDF:
\texorpdfstring{\softtext{SettlementPdf}}{SettlementPdf} construye el
comprobante \texorpdfstring{\softtext{directamente.}}{directamente.}
Testcontainers no se convierte en un servicio del producto; su
dependencia tiene alcance test.

\subsection{\texorpdfstring{\texorpdfstring{\softtext{application.yml}}{application.yml}
y el perfil de
producción}{ y el perfil de producción}}\label{application.yml-y-el-perfil-de-producciuxf3n}

\texorpdfstring{\softcode{backend/src/main/resources/application.yml:1}}{backend/src/main/resources/application.yml:1}
establece apagado graceful, dirección predeterminada 127.0.0.1 y puerto
18080.
\texorpdfstring{\softtext{BOVI\_BIND\_ADDRESS}}{BOVI\_BIND\_ADDRESS} y
\texorpdfstring{\softtext{BOVI\_API\_PORT}}{BOVI\_API\_PORT} permiten
cambiar esos valores. En el Compose integrado el backend escucha
internamente para que Caddy lo alcance; el acceso del usuario se realiza
por el origen local 15180. El puerto interno de Java y el puerto
publicado del proxy son cosas distintas.

La cookie se llama
\texorpdfstring{\softtext{BOVI\_SESSION,}}{BOVI\_SESSION,} tiene
HttpOnly y \texorpdfstring{\softtext{SameSite=lax,}}{SameSite=lax,} y
Secure=false en la configuración local. HttpOnly impide que el
JavaScript ordinario lea la cookie; no reemplaza la autorización del
servidor. SameSite reduce algunos escenarios de envío entre sitios; no
sustituye el token CSRF. El archivo
\texorpdfstring{\softtext{application-production.yml:1}}{application-production.yml:1}
cambia Secure a true. Es una condición necesaria para el perfil HTTPS,
pero ese pequeño archivo no configura por sí solo DNS, certificados,
publicación segura o una plataforma productiva completa.

Spring Session usa un timeout de 30 minutos, la tabla
\texorpdfstring{\softtext{infra.spring\_session}}{infra.spring\_session}
y initialize-schema: never. La aplicación no crea esa tabla al arrancar.
Su existencia y permisos deben proceder de la preparación de la base.
Las propiedades de apagado permiten esperar hasta 30 segundos por fase y
por finalización de tareas
\texorpdfstring{\softtext{programadas.}}{programadas.} Esto sirve al
procedimiento de respaldo con parada ordenada; no hace que una copia
arbitraria de archivos y base sea automáticamente coherente.

Los multipart tienen máximo por archivo de 20971520 bytes, exactamente
20 MiB, y máximo por solicitud de 22020096 bytes, 21 MiB. El margen de
la solicitud cubre el envoltorio y campos
\texorpdfstring{\softtext{adicionales.}}{adicionales.}
\texorpdfstring{\softtext{PrivateDocuments}}{PrivateDocuments} también
comprueba el límite de contenido. La primera barrera puede producir 413;
la comprobación del servicio produce un error de dominio cuando el
archivo está vacío o excede su límite.

La configuración separa cinco conexiones: a, b, auth, worker-a y
worker-b. Los valores predeterminados de negocio pasan por Pgpool local
19999 y 19998; auth apunta al PostgreSQL A técnico en 15432. Los
trabajadores toman las URL de negocio y un usuario distinto,
\texorpdfstring{\softtext{bovi\_worker.}}{bovi\_worker.} Las contraseñas
se reciben mediante propiedades de entorno; sus valores predeterminados
están vacíos. Este capítulo no incluye valores
\texorpdfstring{\softtext{funcionales.}}{funcionales.}

\texorpdfstring{\softtext{BOVI\_DOCUMENTS\_DIR}}{BOVI\_DOCUMENTS\_DIR}
determina el directorio privado. Sin esa variable, la propiedad
\texorpdfstring{\softtext{bovi.documents}}{bovi.documents} usa
\texorpdfstring{\softtext{./var/private-documents;}}{./var/private-documents;}
el lanzamiento local documentado configura un directorio externo.
\texorpdfstring{\softtext{bovi.outbox-enabled,}}{bovi.outbox-enabled,}
alimentada por
\texorpdfstring{\softtext{BOVI\_OUTBOX\_ENABLED,}}{BOVI\_OUTBOX\_ENABLED,}
habilita o deshabilita el
\texorpdfstring{\softtext{trabajador.}}{trabajador.} El nombre de
aplicación es \texorpdfstring{\softtext{bovicierre.}}{bovicierre.}

\subsection{\texorpdfstring{\texorpdfstring{\softtext{BoviCierreApplication:}}{BoviCierreApplication:}
un método de
arranque}{ un método de arranque}}\label{bovicierreapplication-un-muxe9todo-de-arranque}

Referencia:
\texorpdfstring{\softcode{backend/src/main/java/mx/bovicierre/BoviCierreApplication.java:8}}{backend/src/main/java/mx/bovicierre/BoviCierreApplication.java:8}.
La clase tiene
\texorpdfstring{\softtext{SpringBootApplication,}}{SpringBootApplication,}
que define el punto de configuración del proceso, y
\texorpdfstring{\softtext{EnableScheduling,}}{EnableScheduling,} que
permite las tareas
\texorpdfstring{\softtext{programadas.}}{programadas.} No tiene campos
propios \texorpdfstring{\softtext{permanentes.}}{permanentes.}

\textbf{main(String{[}{]} args)} crea la variable application con una
instancia de
\texorpdfstring{\softtext{SpringApplication.}}{SpringApplication.}
Calcula directory a partir de
\texorpdfstring{\softtext{BOVI\_RUNTIME\_DIR}}{BOVI\_RUNTIME\_DIR} o del
directorio privado de ejecución bajo user.home. Registra
\texorpdfstring{\softtext{ApplicationPidFileWriter}}{ApplicationPidFileWriter}
para escribir api.pid y finalmente llama
\texorpdfstring{\softtext{application.run(args).}}{application.run(args).}
El argumento args pasa a Spring; la variable directory no selecciona una
sede ni una empresa, únicamente la ubicación del PID.

El llamador habitual es el lanzador Java del JAR. La salida no es una
respuesta HTTP: es un contexto Spring en
\texorpdfstring{\softtext{funcionamiento.}}{funcionamiento.} El archivo
PID permite que un procedimiento local identifique su proceso propio. No
debe confundirse ese PID con el identificador de una transacción o con
el de un \texorpdfstring{\softtext{contenedor.}}{contenedor.}

\subsection{\texorpdfstring{\texorpdfstring{\softtext{DatabaseConfiguration:}}{DatabaseConfiguration:}
cinco pools y una transacción de
sesión}{ cinco pools y una transacción de sesión}}\label{databaseconfiguration-cinco-pools-y-una-transacciuxf3n-de-sesiuxf3n}

Referencia:
\texorpdfstring{\softcode{backend/src/main/java/mx/bovicierre/shared/DatabaseConfiguration.java:15}}{backend/src/main/java/mx/bovicierre/shared/DatabaseConfiguration.java:15}.
La clase Configuration no guarda un router mutable global. Crea beans
con nombres explícitos que después se inyectan con Qualifier.

\textbf{pool(Environment env, String name)} es la fábrica privada.
prefix se forma como
\texorpdfstring{\softtext{bovi.database}}{bovi.database} más el nombre.
config es un \texorpdfstring{\softtext{HikariConfig}}{HikariConfig}
donde se cargan URL, usuario y contraseña
\texorpdfstring{\softtext{obligatorios.}}{obligatorios.} Los pools
worker reciben máximo dos conexiones; los demás, cuatro.
\texorpdfstring{\softtext{MinimumIdle}}{MinimumIdle} es cero, el tiempo
de adquisición es 3000 ms, el de validación 1000 ms y
\texorpdfstring{\softtext{InitializationFailTimeout}}{InitializationFailTimeout}
es -1. Este último valor permite iniciar el pool sin exigir una conexión
exitosa inmediata; las solicitudes posteriores todavía pueden fallar.

Los métodos
\textbf{\texorpdfstring{\softtext{ownerA(env)}}{ownerA(env)}} y
\textbf{\texorpdfstring{\softtext{ownerB(env)}}{ownerB(env)}} llaman
pool con a y b.
\textbf{\texorpdfstring{\softtext{workerA(env)}}{workerA(env)}} y
\textbf{\texorpdfstring{\softtext{workerB(env)}}{workerB(env)}} usan
worker-a y worker-b.
\textbf{\texorpdfstring{\softtext{authDataSource(env)}}{authDataSource(env)}}
usa auth y está marcado
\texorpdfstring{\softtext{SpringSessionDataSource.}}{SpringSessionDataSource.}
\textbf{\texorpdfstring{\softtext{sessionTransactionManager(ds)}}{sessionTransactionManager(ds)}}
construye
\texorpdfstring{\softtext{DataSourceTransactionManager}}{DataSourceTransactionManager}
sobre el bean \texorpdfstring{\softtext{authDataSource}}{authDataSource}
y lleva
\texorpdfstring{\softtext{SpringSessionTransactionManager.}}{SpringSessionTransactionManager.}
Son seis métodos de publicación de beans más la fábrica privada; ninguno
decide dinámicamente mover una transacción ya iniciada a otro servidor.

La entrada Environment representa configuración resuelta por Spring, no
un objeto enviado por el navegador. La salida de las primeras cinco
fábricas es un \texorpdfstring{\softtext{DataSource.}}{DataSource.} La
última devuelve el administrador dedicado a transacciones de sesión. Los
llamadores son el contenedor Spring,
\texorpdfstring{\softtext{CommandService,}}{CommandService,}
\texorpdfstring{\softtext{AccessService,}}{AccessService,} NodeHealth y
\texorpdfstring{\softtext{OutboxWorker,}}{OutboxWorker,} según sus
Qualifier.

Que existan cinco pools no significa cinco bases distintas: autorización
y varios pools pueden llegar al mismo PostgreSQL A con roles distintos.
La distinción separa permisos, capacidad de conexiones y tareas, aunque
compartan servidor físico.

\section{Comunicación HTTP y forma de los
datos}\label{comunicaciuxf3n-http-y-forma-de-los-datos}

\subsection{Qué viaja desde el
navegador}\label{quuxe9-viaja-desde-el-navegador}

Los clientes usan rutas bajo /api/v1. La sede suele venir en
\texorpdfstring{\softtext{/sedes/\{sedeId\};}}{/sedes/\{sedeId\};} los
recursos usan UUID en la ruta. Los controladores reciben Authentication
proporcionada por Spring. Un cuerpo que contenga actor\_id no puede
sustituir la identidad del contexto de seguridad. En
\texorpdfstring{\softtext{CoreController}}{CoreController} se rechazan
explícitamente \texorpdfstring{\softtext{empresa\_id,}}{empresa\_id,}
\texorpdfstring{\softtext{sucursal\_id,}}{sucursal\_id,} actor\_id y
otros campos que el servidor debe
\texorpdfstring{\softtext{determinar.}}{determinar.}

Una mutación ordinaria lleva un cuerpo JSON, X-XSRF-TOKEN y, para
comandos de negocio,
\texorpdfstring{\softtext{Idempotency-Key.}}{Idempotency-Key.} Algunas
mutaciones de operación también requieren If-Match. La cookie mantiene
la sesión. Son controles con funciones diferentes: cookie para localizar
la sesión, CSRF para proteger el envío, idempotencia para identificar un
intento y ETag para comprobar la versión.

\begin{Shaded}
\begin{Highlighting}[]
\FunctionTok{\{}
  \DataTypeTok{"importe"}\FunctionTok{:} \StringTok{"1200.00"}\FunctionTok{,}
  \DataTypeTok{"moneda"}\FunctionTok{:} \StringTok{"MXN"}
\FunctionTok{\}}
\end{Highlighting}
\end{Shaded}

Este ejemplo demuestra la estructura Dinero. El valor monetario es una
cadena, no 1200.00 como número JSON. El peso usa otra estructura:

\begin{Shaded}
\begin{Highlighting}[]
\FunctionTok{\{}
  \DataTypeTok{"kg"}\FunctionTok{:} \StringTok{"300.000"}
\FunctionTok{\}}
\end{Highlighting}
\end{Shaded}

Un precio unitario puede ser la cadena ``10.0000''. Los dos, tres o
cuatro decimales dependen del tipo, no de un formato visual elegido por
el navegador. Un peso individual desconocido puede omitirse en una
propuesta o aparecer nulo en un resultado; el sistema no debe repartir
el peso grupal y fingir que fue medido para cada animal.

\subsection{Lecturas, comandos y documentos
derivados}\label{lecturas-comandos-y-documentos-derivados}

Una consulta recibe filtros y devuelve un JSON canónico o una
representación derivada. Una mutación llama una función de comando.
Exportar un CSV, dibujar un QR o generar el PDF de una liquidación
consulta información existente; esas salidas no crean otra venta ni otra
deuda.

Una excepción importante es consultar tarifas: el GET consulta el
proveedor mediante A y guarda una captura técnica en el propietario de
la sede. Por tanto, no es una lectura completamente libre de escritura
local. La captura se mantiene separada de la transacción económica y de
la aceptación comercial.

OpenAPI describe 70
\texorpdfstring{\softtext{operationId.}}{operationId.} Los nombres de
parámetros animalId,
\texorpdfstring{\softtext{operacionId}}{operacionId} o
\texorpdfstring{\softtext{documentoId}}{documentoId} del contrato pueden
aparecer como id en una anotación Java. Para comprender la equivalencia
importa la posición y el tipo del segmento de ruta, no que ambos
archivos utilicen el mismo nombre local.

Las respuestas de comando no se fabrican con una fórmula en el
\texorpdfstring{\softtext{controlador.}}{controlador.} En el camino
ordinario, PostgreSQL devuelve jsonb, JDBC lo convierte a texto y
\texorpdfstring{\softtext{ResponseEntity}}{ResponseEntity} lo publica
como \texorpdfstring{\softtext{application/json.}}{application/json.}
Ese transporte conserva las cadenas monetarias
\texorpdfstring{\softtext{canónicas.}}{canónicas.}

\section{Autenticación y autorización explicadas clase por
clase}\label{autenticaciuxf3n-y-autorizaciuxf3n-explicadas-clase-por-clase}

\subsection{\texorpdfstring{\texorpdfstring{\softtext{AccessService:}}{AccessService:}
la identidad y las membresías
vigentes}{ la identidad y las membresías vigentes}}\label{accessservice-la-identidad-y-las-membresuxedas-vigentes}

Referencia:
\texorpdfstring{\softcode{backend/src/main/java/mx/bovicierre/acceso/AccessService.java:14}}{backend/src/main/java/mx/bovicierre/acceso/AccessService.java:14}.
El único campo es auth, un
\texorpdfstring{\softtext{JdbcTemplate}}{JdbcTemplate} construido sobre
\texorpdfstring{\softtext{authDataSource.}}{authDataSource.} El
constructor
\textbf{\texorpdfstring{\softtext{AccessService(source)}}{AccessService(source)}}
fija esa dependencia técnica. Este servicio no utiliza ownerB para
decidir si alguien puede entrar.

\textbf{login(String login)} normaliza el texto con trim y minúsculas de
\texorpdfstring{\softtext{Locale.ROOT.}}{Locale.ROOT.} La consulta
parametrizada busca id,
\texorpdfstring{\softtext{empresa\_id,}}{empresa\_id,} nombre,
\texorpdfstring{\softtext{hash\_contrasena}}{hash\_contrasena} y estado.
El filtro empresarial procede de
\texorpdfstring{\softtext{infra.nodo\_config,}}{infra.nodo\_config,} no
del cuerpo de login. rows es la lista de coincidencias; se devuelve su
primera fila únicamente cuando tiene tamaño uno. Cero o varias filas
producen null. Este método localiza el registro; la comparación de
contraseña ocurre en
\texorpdfstring{\softtext{SessionController.}}{SessionController.}

\textbf{user(Authentication authentication)} revisa ausencia, indicador
de autenticación y tipo
\texorpdfstring{\softtext{AnonymousAuthenticationToken.}}{AnonymousAuthenticationToken.}
Si el contexto no representa una sesión activa, lanza AUTH401. Su salida
es un UUID convertido desde
\texorpdfstring{\softtext{authentication.getName().}}{authentication.getName().}
No toma el UUID de una cabecera libre. Lo llaman prácticamente todos los
controladores protegidos y el filtro de autorización vigente.

\textbf{current(UUID user)} consulta id y nombre para un usuario ACTIVO,
nuevamente limitado a la empresa
\texorpdfstring{\softtext{configurada.}}{configurada.} Una cantidad
distinta de una fila produce AUTH401 con el mensaje de sesión no
\texorpdfstring{\softtext{autorizada.}}{autorizada.} La cookie puede
existir y, aun así, este método rechazar al usuario bloqueado. Lo llaman
el filtro de seguridad, authorize y la consulta de sesión actual.

\textbf{sites(UUID user)} une
\texorpdfstring{\softtext{membresia\_usuario}}{membresia\_usuario} con
sucursal, exige membresía vigente y ordena por código de sede y rol.
Devuelve una lista de filas. Si el usuario tiene dos roles en una sede,
inicialmente hay dos filas para esa sede. Cada fila contiene identidad
de sede, empresa, código, nombre, nodo dueño y
\texorpdfstring{\softtext{rol\_codigo.}}{rol\_codigo.} La agrupación
para la interfaz ocurre después.

\textbf{authorize(UUID user, UUID site, boolean mutation)} primero
ejecuta current. Filtra las filas de sites para la sede solicitada; rows
guarda ese \texorpdfstring{\softtext{subconjunto.}}{subconjunto.} Si
está vacío, AUTH403. roles reúne códigos
\texorpdfstring{\softtext{diferentes.}}{diferentes.} Para una mutación
exige al menos ADMIN, OPERADOR o CAJA, de modo que un AUDITOR solo no
puede escribir. first aporta la empresa y el
\texorpdfstring{\softtext{propietario.}}{propietario.} El resultado es
el record Scope.

\textbf{Scope} conserva company, site, owner y roles, con accesores
company(), site(), owner() y roles(). Java genera también constructor,
igualdad, hashCode y toString propios de un record; no son comandos
económicos escritos a mano. Un Scope es la decisión obtenida para esa
llamada. No significa que los roles se copien permanentemente dentro de
la cookie.

authorize solo aplica la primera clasificación de escritura.
\texorpdfstring{\softtext{CommandService}}{CommandService} exige después
el rol pertinente al comando y PostgreSQL valida su propio contexto.
Esto evita atribuir a mutation=true una precisión que no tiene: permitir
una escritura general no equivale a autorizar cualquier ajuste
\texorpdfstring{\softtext{financiero.}}{financiero.}

\subsection{\texorpdfstring{\texorpdfstring{\softtext{AvailabilityFilter:}}{AvailabilityFilter:}
comprobar A antes de procesar la
API}{ comprobar A antes de procesar la API}}\label{availabilityfilter-comprobar-a-antes-de-procesar-la-api}

Referencia:
\texorpdfstring{\softcode{backend/src/main/java/mx/bovicierre/acceso/AvailabilityFilter.java:19}}{backend/src/main/java/mx/bovicierre/acceso/AvailabilityFilter.java:19}.
Es un
\texorpdfstring{\softtext{OncePerRequestFilter}}{OncePerRequestFilter}
registrado con la prioridad más alta. Su campo auth apunta al
\texorpdfstring{\softtext{JdbcTemplate}}{JdbcTemplate} técnico central,
creado por el constructor
\textbf{\texorpdfstring{\softtext{AvailabilityFilter(ds)}}{AvailabilityFilter(ds)}}.

\textbf{\texorpdfstring{\softtext{shouldNotFilter(request)}}{shouldNotFilter(request)}}
devuelve true cuando la URI no empieza con /api/ o cuando es exactamente
\texorpdfstring{\softtext{/api/v1/salud.}}{/api/v1/salud.} Es decir,
omite el chequeo en esos casos. La salud necesita permanecer consultable
incluso si A falla. El GET anónimo de CSRF sí queda dentro del chequeo
de \texorpdfstring{\softtext{disponibilidad.}}{disponibilidad.}

\textbf{\texorpdfstring{\softtext{doFilterInternal(request,}}{doFilterInternal(request,}
response, chain)} ejecuta select 1 sobre auth y después
\texorpdfstring{\softtext{chain.doFilter.}}{chain.doFilter.} Si aparece
\texorpdfstring{\softtext{DataAccessException,}}{DataAccessException,}
utiliza
\texorpdfstring{\softtext{ApiErrors.database}}{ApiErrors.database} y
\texorpdfstring{\softtext{ApiErrors.write.}}{ApiErrors.write.} No envía
la solicitud de B a A como sustituto ni la de A a B: A es aquí
dependencia de autorización y sesión.

Las variables request y response son objetos HTTP; chain representa el
resto de filtros y despacho. La comprobación de conectividad no lee
roles. Su papel es convertir la indisponibilidad central en una
respuesta coherente antes de continuar. El bloque try también envuelve
la cadena, por lo que puede recoger una
\texorpdfstring{\softtext{DataAccessException}}{DataAccessException} que
se propague desde ella.

\subsection{\texorpdfstring{\texorpdfstring{\softtext{SecurityConfiguration:}}{SecurityConfiguration:}
seguridad de
infraestructura}{ seguridad de infraestructura}}\label{securityconfiguration-seguridad-de-infraestructura}

Referencia:
\texorpdfstring{\softcode{backend/src/main/java/mx/bovicierre/acceso/SecurityConfiguration.java:21}}{backend/src/main/java/mx/bovicierre/acceso/SecurityConfiguration.java:21}.
Esta clase declara beans; no conserva campos propios. Conviene
distinguir sus métodos porque no todos realizan la misma
\texorpdfstring{\softtext{protección.}}{protección.}

\textbf{\texorpdfstring{\softtext{noGeneratedUsers()}}{noGeneratedUsers()}}
entrega un
\texorpdfstring{\softtext{UserDetailsService}}{UserDetailsService} cuya
función de búsqueda siempre lanza
\texorpdfstring{\softtext{UsernameNotFoundException.}}{UsernameNotFoundException.}
Evita que el producto se apoye en un usuario automático de demostración
generado por configuración
\texorpdfstring{\softtext{predeterminada.}}{predeterminada.} No contiene
una contraseña fija.

\textbf{\texorpdfstring{\softtext{passwordEncoder()}}{passwordEncoder()}}
crea
\texorpdfstring{\softtext{BCryptPasswordEncoder}}{BCryptPasswordEncoder}
con coste 12.
\textbf{\texorpdfstring{\softtext{securityContextRepository()}}{securityContextRepository()}}
crea
\texorpdfstring{\softtext{HttpSessionSecurityContextRepository,}}{HttpSessionSecurityContextRepository,}
que guarda el contexto autenticado en sesión.
\textbf{\texorpdfstring{\softtext{csrfRepository()}}{csrfRepository()}}
crea
\texorpdfstring{\softtext{HttpSessionCsrfTokenRepository}}{HttpSessionCsrfTokenRepository}
y fija la cabecera
\texorpdfstring{\softtext{X-XSRF-TOKEN.}}{X-XSRF-TOKEN.} Estos objetos
se inyectan también en
\texorpdfstring{\softtext{SessionController.}}{SessionController.}

\textbf{security(http, access, context, csrf)} construye
\texorpdfstring{\softtext{SecurityFilterChain.}}{SecurityFilterChain.}
Declara públicas las rutas de CSRF, login y salud; las demás requieren
\texorpdfstring{\softtext{autenticación.}}{autenticación.} Configura
CSRF con su repositorio y
\texorpdfstring{\softtext{CsrfTokenRequestAttributeHandler,}}{CsrfTokenRequestAttributeHandler,}
así como el repositorio del contexto de seguridad. Deshabilita request
cache, formulario de login, HTTP Basic y logout automático porque las
rutas JSON propias se ocupan de esas acciones.

Una ruta pública no queda exenta automáticamente de CSRF. Login es
pública para que un usuario todavía anónimo pueda enviar credenciales,
pero el POST sigue protegido. Esa distinción explica el GET de
bootstrap: primero se obtiene el token y después se realiza la mutación.

El
\texorpdfstring{\softtext{authenticationEntryPoint}}{authenticationEntryPoint}
escribe AUTH401 en formato JSON. El
\texorpdfstring{\softtext{accessDeniedHandler}}{accessDeniedHandler}
escribe AUTH403 por falta de permiso o CSRF inválido. No se redirige a
una página HTML de \texorpdfstring{\softtext{contraseña.}}{contraseña.}

Dentro de security se crea además un filtro anónimo con su propio
\textbf{\texorpdfstring{\softtext{doFilterInternal(request,}}{doFilterInternal(request,}
response, chain)}. Después de
\texorpdfstring{\softtext{AnonymousAuthenticationFilter,}}{AnonymousAuthenticationFilter,}
obtiene authentication desde
\texorpdfstring{\softtext{SecurityContextHolder.}}{SecurityContextHolder.}
Si no es anónima, verifica
\texorpdfstring{\softtext{access.current(access.user(authentication)).}}{access.current(access.user(authentication)).}
Agrega X-Correlation-ID y continúa la cadena. Captura
\texorpdfstring{\softtext{ApiException}}{ApiException} y
\texorpdfstring{\softtext{DataAccessException}}{DataAccessException}
para que una revocación o un problema de base conserve la respuesta
pública prevista.

Las autoridades de Spring de la sesión están vacías en el login actual;
las reglas de negocio no dependen de anotaciones hasRole sobre una copia
antigua de permisos. El control vigente se consulta en
\texorpdfstring{\softtext{AccessService.}}{AccessService.} La conexión
central y las comprobaciones SQL son parte de ese contrato.

\subsection{\texorpdfstring{\texorpdfstring{\softtext{SessionController:}}{SessionController:}
los cinco endpoints de sesión y
salud}{ los cinco endpoints de sesión y salud}}\label{sessioncontroller-los-cinco-endpoints-de-sesiuxf3n-y-salud}

Referencia:
\texorpdfstring{\softcode{backend/src/main/java/mx/bovicierre/acceso/SessionController.java:26}}{backend/src/main/java/mx/bovicierre/acceso/SessionController.java:26}.
Sus campos son access, encoder, contexts, csrf, dummyHash y health. El
constructor
\textbf{\texorpdfstring{\softtext{SessionController(…)}}{SessionController(…)}}
recibe las dependencias y genera dummyHash a partir de un UUID
aleatorio. Ese hash permite realizar una comparación BCrypt también
cuando no se encontró usuario; reduce una diferencia evidente de camino,
pero no debe describirse como una garantía formal de tiempo constante.

\textbf{\texorpdfstring{\softtext{csrf(CsrfToken}}{csrf(CsrfToken}
token)} responde al GET
\texorpdfstring{\softtext{/api/v1/sesion/csrf.}}{/api/v1/sesion/csrf.}
Devuelve un mapa con token y headerName y establece Cache-Control:
no-store. El valor de token lo suministra la infraestructura de CSRF. El
cliente debe utilizar el nombre de cabecera recibido, no asumir otra
variante.

\textbf{login(Login input, request, response)} verifica que la
contraseña ocupe como máximo 72 bytes UTF-8. El record Login tiene
NotBlank y Size para login y contrasena, pero una longitud de caracteres
no es idéntica a una longitud de bytes. Una contraseña con caracteres
multibyte puede alcanzar el límite de BCrypt con menos caracteres
visibles.

user contiene el resultado de
\texorpdfstring{\softtext{access.login.}}{access.login.} hash usa
dummyHash si no hay usuario o elimina el prefijo de almacenamiento
\{bcrypt\} si lo hay. Se compara la contraseña y se comprueba el estado
ACTIVO. El mensaje de fallo no revela si el usuario existía o cuál de
las condiciones falló.

Al tener éxito, login obtiene o crea la sesión, cambia su identificador,
crea un contexto vacío y le agrega un
\texorpdfstring{\softtext{UsernamePasswordAuthenticationToken}}{UsernamePasswordAuthenticationToken}
autenticado cuyo nombre es el UUID del usuario. Las credenciales del
token son null. Guarda el contexto y borra el token CSRF anterior con
\texorpdfstring{\softtext{csrf.saveToken(null,…).}}{csrf.saveToken(null,…).}
Devuelve 204. El cliente necesita pedir un token nuevo después de esa
rotación.

\textbf{logout(request, response)} obtiene la sesión sin crearla. Si
existe, la invalida, limpia
\texorpdfstring{\softtext{SecurityContextHolder}}{SecurityContextHolder}
y devuelve 204. El método no contiene una instrucción manual para borrar
la cookie por nombre; la identidad deja de ser válida porque se invalidó
la sesión del lado servidor y se limpió el contexto.

\textbf{current(authentication)} devuelve GET
\texorpdfstring{\softtext{/api/v1/sesion/actual.}}{/api/v1/sesion/actual.}
id identifica al usuario; user contiene su nombre vigente. sites es un
\texorpdfstring{\softtext{LinkedHashMap}}{LinkedHashMap} por UUID. Cada
primera fila de sede se copia en entry, se elimina rol\_codigo y se
crean listas roles y
\texorpdfstring{\softtext{capacidades.}}{capacidades.} Las siguientes
filas agregan roles a la sede ya existente. Toda sede recibe CONSULTAR;
una fila que no sea AUDITOR puede incorporar ESCRIBIR.

Estas capacidades simplifican la interfaz, pero no representan permisos
específicos para todos los comandos. CAJA puede mostrar ESCRIBIR y no
estar autorizada para un ajuste reservado a ADMIN. La decisión final
sigue en los servicios y el SQL.

\textbf{health()} publica GET
\texorpdfstring{\softtext{/api/v1/salud}}{/api/v1/salud} delegando a
\texorpdfstring{\softtext{NodeHealth.status.}}{NodeHealth.status.}
\textbf{Login} es un record de entrada con accesores login() y
contrasena(): login no vacío y máximo 80 caracteres, contraseña no vacía
y máximo 72 caracteres, además del control de bytes dentro del método.

\subsection{\texorpdfstring{\texorpdfstring{\softtext{UserController:}}{UserController:}
crear usuarios sin enviar hashes desde el
cliente}{ crear usuarios sin enviar hashes desde el cliente}}\label{usercontroller-crear-usuarios-sin-enviar-hashes-desde-el-cliente}

Referencia:
\texorpdfstring{\softcode{backend/src/main/java/mx/bovicierre/acceso/UserController.java:14}}{backend/src/main/java/mx/bovicierre/acceso/UserController.java:14}.
Tiene access, commands y encoder; el constructor fija las tres
\texorpdfstring{\softtext{dependencias.}}{dependencias.}

El método privado \textbf{write(site, command, body, key, target, auth,
request, status)} obtiene el usuario, autoriza mutación y exige ADMIN.
Parsear el JSON ocurre después de esa
\texorpdfstring{\softtext{comprobación.}}{comprobación.} No admite
contrasena fuera de
\texorpdfstring{\softtext{crear\_usuario}}{crear\_usuario} ni permite
\texorpdfstring{\softtext{hash\_contrasena,}}{hash\_contrasena,}
\texorpdfstring{\softtext{empresa\_id,}}{empresa\_id,}
\texorpdfstring{\softtext{sucursal\_id}}{sucursal\_id} o actor\_id
suministrados por el navegador. target, cuando existe, fija usuario\_id
desde la ruta.

fingerprint se calcula antes de reemplazar la contraseña por BCrypt.
Esto importa porque BCrypt genera una sal aleatoria: dos codificaciones
de la misma contraseña no producen el mismo texto. Si la intención
idempotente usara el hash salado nuevo, un reintento legítimo podría
parecer un contenido diferente.

password es el nodo retirado de payload. Debe ser textual, no blanco y
ocupar hasta 72 bytes UTF-8. Luego se agrega
\texorpdfstring{\softtext{hash\_contrasena}}{hash\_contrasena} con
prefijo \{bcrypt\}. La contraseña en claro no se manda como campo del
comando SQL. En
\texorpdfstring{\softtext{crear\_usuario}}{crear\_usuario} es
\texorpdfstring{\softtext{obligatoria.}}{obligatoria.} La llamada usa la
sobrecarga de \texorpdfstring{\softtext{CommandService}}{CommandService}
que recibe fingerprint
\texorpdfstring{\softtext{explícito.}}{explícito.} La huella de
intención tampoco se publica en las respuestas del controlador ni en sus
logs \texorpdfstring{\softtext{ordinarios.}}{ordinarios.}

\textbf{create(\ldots)} atiende POST a usuarios, selecciona
\texorpdfstring{\softtext{crear\_usuario}}{crear\_usuario} y responde
201. \textbf{update(\ldots)} atiende PATCH
\texorpdfstring{\softtext{/usuarios/\{id\},}}{/usuarios/\{id\},}
selecciona
\texorpdfstring{\softtext{actualizar\_usuario}}{actualizar\_usuario} y
responde 200. \textbf{membership(\ldots)} atiende POST
\texorpdfstring{\softtext{/usuarios/\{id\}/membresias,}}{/usuarios/\{id\}/membresias,}
selecciona
\texorpdfstring{\softtext{asignar\_membresia}}{asignar\_membresia} y
responde 200. Sus parámetros proporcionan sede, cuerpo, clave, usuario
objetivo y contexto HTTP; el método común aplica el resto.

Estos comandos son maestros y se ejecutan en A, incluso cuando la sede
cuya membresía se administra es B. La creación de un usuario no implica
por sí sola una membresía para todas las sedes. La asignación es un
comando separado.

\section{\texorpdfstring{\texorpdfstring{\softtext{JsonSupport:}}{JsonSupport:}
frontera de datos, ETag y huella de
intención}{ frontera de datos, ETag y huella de intención}}\label{jsonsupport-frontera-de-datos-etag-y-huella-de-intenciuxf3n}

Referencia:
\texorpdfstring{\softcode{backend/src/main/java/mx/bovicierre/shared/JsonSupport.java:11}}{backend/src/main/java/mx/bovicierre/shared/JsonSupport.java:11}.
Es una clase final de utilidades con constructor privado vacío. Su campo
público estático JSON es un JsonMapper
\texorpdfstring{\softtext{compartido.}}{compartido.} No conserva el
cuerpo de la última solicitud.

\textbf{object(String body)} lee el árbol JSON. node debe existir y ser
objeto; un array raíz, una cadena raíz o un texto inválido generan
BV400. La salida es \texorpdfstring{\softtext{ObjectNode.}}{ObjectNode.}
Los controladores ordinarios que reciben String utilizan este método,
mientras Login utiliza deserialización y validación del framework.

\textbf{\texorpdfstring{\softtext{sorted(JsonNode}}{sorted(JsonNode}
value)} es recursivo y privado. Para un objeto crea result y un TreeMap
fields, ordena las claves y procesa cada hijo. Para un array crea otro
array y procesa los hijos en su mismo orden. Para un escalar devuelve el
valor. Ordenar las propiedades no significa ordenar los animales de una
lista: cambiar el orden de un array sigue siendo un cambio de intención
para esta huella.

\textbf{hash(String command, JsonNode body)} serializa el cuerpo
ordenado, lo precede con el nombre de comando y un salto de línea, lo
codifica en UTF-8 y calcula SHA-256
\texorpdfstring{\softtext{hexadecimal.}}{hexadecimal.} La variable
content es el arreglo de bytes firmado por esa huella; el nombre del
comando evita que idéntico JSON de acciones distintas sea la misma
\texorpdfstring{\softtext{intención.}}{intención.}

\[
H = \operatorname{SHA256}\bigl(\operatorname{UTF8}(comando \,\Vert\, \operatorname{LF} \,\Vert\, JSON_{ordenado}(cuerpo))\bigr)
\]

En la ecuación, LF representa el byte de salto de línea, 0 x0A, que el
código inserta entre el comando y el JSON. La huella no es una firma
digital ni prueba la identidad de quien envió el cuerpo. La
autenticación, el ámbito y la clave completan el control. Un cuerpo
inválido puede tener una huella calculable y después ser rechazado:
generar hash no equivale a validar el dominio.

\textbf{version(String etag)} lanza 428 si If-Match está ausente. Exige
un entero no negativo entre comillas, sin ceros iniciales salvo cero, y
con hasta 18 dígitos. Extrae el texto interior y lo convierte a long. No
acepta un ETag débil con W/, una lista de versiones ni un sufijo de
\texorpdfstring{\softtext{compresión.}}{compresión.} El ETag sirve como
precondición de versión, no como clave de sesión.

\textbf{\texorpdfstring{\softtext{animalIdentifications(ObjectNode}}{animalIdentifications(ObjectNode}
payload)} toma identificaciones o construye una forma compatible a
partir de sistema y codigo. Exige array no vacío. Cada identification
debe ser objeto, sistema y codigo textuales, sistema en SINIIGA, SINIDA,
INTERNO u OTRO, y código de 1 a 40 caracteres ASCII
\texorpdfstring{\softtext{permitidos.}}{permitidos.} Debe contener al
menos una letra o dígito; espacios o guiones solos no bastan.

\textbf{\texorpdfstring{\softtext{decimals(JsonNode}}{decimals(JsonNode}
node)} recorre objetos y arrays. Rechaza entradas llamadas total, saldo,
disponible, margen y
\texorpdfstring{\softtext{costo\_vigente}}{costo\_vigente} porque el
servidor las deriva. Para importe, monto,
\texorpdfstring{\softtext{cargos\_pactados}}{cargos\_pactados} y
\texorpdfstring{\softtext{bonificaciones\_pactadas}}{bonificaciones\_pactadas}
admite un objeto Dinero con importe y moneda MXN o una cadena decimal.
Los objetos se recorren nuevamente y se valida su importe interior.

{\def\LTcaptype{none} % do not increment counter
\begin{longtable}[]{@{}
  >{\raggedright\arraybackslash}p{(\linewidth - 4\tabcolsep) * \real{0.3333}}
  >{\raggedright\arraybackslash}p{(\linewidth - 4\tabcolsep) * \real{0.3333}}
  >{\raggedright\arraybackslash}p{(\linewidth - 4\tabcolsep) * \real{0.3333}}@{}}
\toprule\noalign{}
\begin{minipage}[b]{\linewidth}\raggedright
Familia
\end{minipage} & \begin{minipage}[b]{\linewidth}\raggedright
Forma textual comprobada
\end{minipage} & \begin{minipage}[b]{\linewidth}\raggedright
Responsabilidad restante
\end{minipage} \\
\midrule\noalign{}
\endhead
\bottomrule\noalign{}
\endlastfoot
Importe, monto, cargos y bonificaciones & Dos decimales; signo opcional
y hasta 16 dígitos enteros & SQL decide signo válido y efecto
económico \\
Precio unitario & Cuatro decimales; sin signo negativo y hasta 10
dígitos enteros & SQL comprueba condiciones de precio \\
Kg y pesos nombrados, valor & Tres decimales; sin signo negativo y hasta
11 dígitos enteros & SQL valida tara, neto y magnitud \\
Cantidad, \texorpdfstring{\softtext{deducción\_pct}}{deducción\_pct} y
base\_valor & Cadena con hasta seis decimales según el patrón general &
SQL valida unidad y rango \\
\end{longtable}
}

Este validador no es un validador integral de todo OpenAPI. Reconoce una
lista de nombres de campo y formas decimales; no convierte cada petición
en un DTO Java \texorpdfstring{\softtext{específico.}}{específico.} Para
explicar la seguridad correctamente, deben separarse forma JSON, campos
protegidos, autorización y restricciones canónicas SQL.

\section{\texorpdfstring{\texorpdfstring{\softtext{CommandService:}}{CommandService:}
enrutamiento, transacción local y
observabilidad}{ enrutamiento, transacción local y observabilidad}}\label{commandservice-enrutamiento-transacciuxf3n-local-y-observabilidad}

Referencia:
\texorpdfstring{\softcode{backend/src/main/java/mx/bovicierre/shared/CommandService.java:16}}{backend/src/main/java/mx/bovicierre/shared/CommandService.java:16}.
Esta es la pieza central de comunicación
\texorpdfstring{\softtext{económica.}}{económica.} Todos los
controladores de comandos la llaman, incluyendo usuarios, documentos y
aceptación de tarifas.

\subsection{Campos y constructor}\label{campos-y-constructor}

\textbf{access} conserva
\texorpdfstring{\softtext{AccessService.}}{AccessService.} \textbf{jdbc}
es el mapa A/B de
\texorpdfstring{\softtext{JdbcTemplate;}}{JdbcTemplate;}
\textbf{transactions}, el mapa A/B de
\texorpdfstring{\softtext{TransactionTemplate.}}{TransactionTemplate.}
El constructor
\textbf{\texorpdfstring{\softtext{CommandService(access,}}{CommandService(access,}
a, b)} crea un administrador
\texorpdfstring{\softtext{DataSourceTransactionManager}}{DataSourceTransactionManager}
independiente por fuente. Las claves A y B fijan la relación entre
plantilla JDBC y transacción del mismo
\texorpdfstring{\softtext{propietario.}}{propietario.} No hay un
DataSource por defecto que concentre accidentalmente todas las
escrituras en A.

\textbf{MASTER} contiene
\texorpdfstring{\softtext{crear\_contraparte,}}{crear\_contraparte,}
\texorpdfstring{\softtext{registrar\_animal,}}{registrar\_animal,}
\texorpdfstring{\softtext{crear\_predio,}}{crear\_predio,}
\texorpdfstring{\softtext{crear\_usuario,}}{crear\_usuario,}
\texorpdfstring{\softtext{actualizar\_usuario}}{actualizar\_usuario} y
\texorpdfstring{\softtext{asignar\_membresia.}}{asignar\_membresia.} Es
una clasificación por nombre de comando controlado por el servidor.
\textbf{LOG} escribe mensajes
\texorpdfstring{\softtext{estructurados.}}{estructurados.}
\textbf{\texorpdfstring{\softtext{processStarted}}{processStarted}}
conserva el instante de creación del servicio.

Los LongAdder \textbf{successes},
\textbf{\texorpdfstring{\softtext{commandErrors}}{commandErrors}} y
\textbf{\texorpdfstring{\softtext{commandNanos}}{commandNanos}} acumulan
resultados y duración de comandos. \textbf{queries},
\textbf{\texorpdfstring{\softtext{queryErrors}}{queryErrors}},
\textbf{\texorpdfstring{\softtext{slowQueries}}{slowQueries}} y
\textbf{queryNanos} acumulan datos de consultas medidas. LongAdder
permite incrementos concurrentes; estos números no son asientos
financieros ni están protegidos por la transacción SQL. Reiniciar Java
reinicia los \texorpdfstring{\softtext{contadores.}}{contadores.}

\subsection{Las dos sobrecargas
command}\label{las-dos-sobrecargas-command}

\textbf{command(user, site, command, payload, key, correlation)} calcula
\texorpdfstring{\softtext{JsonSupport.hash}}{JsonSupport.hash} y llama
la sobrecarga completa. \textbf{command(user, site, command, payload,
key, correlation, hash)} acepta una huella previamente
\texorpdfstring{\softtext{construida.}}{construida.}
\texorpdfstring{\softtext{UserController}}{UserController} la necesita
para la contraseña y
\texorpdfstring{\softtext{DocumentController}}{DocumentController} para
su intención \texorpdfstring{\softtext{documental.}}{documental.}

La ejecución completa comienza con started, owner inicialmente
\texorpdfstring{\softtext{NO\_RESUELTO,}}{NO\_RESUELTO,} outcome
inicialmente FALLO y errorCode nulo. scope resulta de authorize con
\texorpdfstring{\softtext{mutation=true.}}{mutation=true.} owner se
vuelve A para un maestro o el propietario de scope para un agregado
local.

required selecciona los roles que pueden realizar la acción:

{\def\LTcaptype{none} % do not increment counter
\begin{longtable}[]{@{}
  >{\raggedright\arraybackslash}p{(\linewidth - 4\tabcolsep) * \real{0.3333}}
  >{\raggedright\arraybackslash}p{(\linewidth - 4\tabcolsep) * \real{0.3333}}
  >{\raggedright\arraybackslash}p{(\linewidth - 4\tabcolsep) * \real{0.3333}}@{}}
\toprule\noalign{}
\begin{minipage}[b]{\linewidth}\raggedright
Clasificación escrita en Java
\end{minipage} & \begin{minipage}[b]{\linewidth}\raggedright
Roles admitidos
\end{minipage} & \begin{minipage}[b]{\linewidth}\raggedright
Ejemplos por nombre
\end{minipage} \\
\midrule\noalign{}
\endhead
\bottomrule\noalign{}
\endlastfoot
Administración & ADMIN & Nombres que contienen periodo;
\texorpdfstring{\softtext{anular\_venta,}}{anular\_venta,}
\texorpdfstring{\softtext{registrar\_baja,}}{registrar\_baja,}
\texorpdfstring{\softtext{registrar\_revision,}}{registrar\_revision,}
\texorpdfstring{\softtext{ajustar\_liquidacion}}{ajustar\_liquidacion} y
administración de usuarios \\
Caja & ADMIN o CAJA &
\texorpdfstring{\softtext{registrar\_pago,}}{registrar\_pago,}
\texorpdfstring{\softtext{validar\_pago,}}{validar\_pago,}
\texorpdfstring{\softtext{aplicar\_pago,}}{aplicar\_pago,}
\texorpdfstring{\softtext{revertir\_aplicacion}}{revertir\_aplicacion} y
\texorpdfstring{\softtext{reembolsar\_pago}}{reembolsar\_pago} \\
Operación general & ADMIN u OPERADOR & Otros comandos, como borradores,
pesajes, reservas, documentos y aceptación de tarifa \\
\end{longtable}
}

Esta tabla reproduce la clasificación del código, no una inferencia
sobre cualquier combinación de payload. Por ejemplo, el controlador de
movimientos usa el nombre genérico
\texorpdfstring{\softtext{registrar\_movimiento;}}{registrar\_movimiento;}
no debe afirmarse que Java aplica automáticamente toda regla de un alias
\texorpdfstring{\softtext{registrar\_baja}}{registrar\_baja} por el mero
texto del \texorpdfstring{\softtext{formulario.}}{formulario.} Las
reglas específicas del comando y sus datos se comprueban también en el
SQL.

Después se valida Idempotency-Key: no nula, no blanca y longitud de 16 a
80. Se ejecutan decimals y, para
\texorpdfstring{\softtext{registrar\_animal,}}{registrar\_animal,}
\texorpdfstring{\softtext{animalIdentifications.}}{animalIdentifications.}
node repite la elección A o
\texorpdfstring{\softtext{scope.owner.}}{scope.owner.} Si no existe una
plantilla para ese nodo, devuelve BV501/503. No intenta enviar al nodo
contrario.

Cada attempt abre una transacción local mediante
\texorpdfstring{\softtext{transactions.get(node).execute.}}{transactions.get(node).execute.}
Dentro del callback usa exactamente
\texorpdfstring{\softtext{jdbc.get(node)}}{jdbc.get(node)} y una
sentencia parametrizada:

\begin{Shaded}
\begin{Highlighting}[]
\KeywordTok{select}\NormalTok{ bovi.cmd\_ejecutar(?,?,?,?,?,?,?,?:}\CharTok{:jsonb}\NormalTok{):}\CharTok{:text}
\end{Highlighting}
\end{Shaded}

Los ocho argumentos son empresa autorizada, sede autorizada, usuario,
clave, huella, correlación, comando y payload
\texorpdfstring{\softtext{serializado.}}{serializado.} Los signos de
interrogación son parámetros JDBC; el cuerpo no se concatena como SQL.
Las funciones y tablas internas pertenecen a las migraciones, no a una
instrucción arbitraria recibida del navegador.

response es el JSON devuelto.
\texorpdfstring{\softtext{TransactionTemplate}}{TransactionTemplate}
cierra la transacción antes de que execute entregue el resultado al
llamador. Si todo concluye, outcome pasa a
\texorpdfstring{\softtext{COMPLETADA.}}{COMPLETADA.} La API no ejecuta
después una consulta financiera para reconstruir el resultado de un
commit conocido; recibe el resultado persistible producido por el
comando canónico.

\subsection{Qué se reintenta y qué
no}\label{quuxe9-se-reintenta-y-quuxe9-no}

Si \texorpdfstring{\softtext{DataAccessException}}{DataAccessException}
tiene como causa más específica una SQLException 40001 o 40P01, y
attempt es menor que dos, se espera un tiempo acotado y se repite la
misma \texorpdfstring{\softtext{intención.}}{intención.} El programa
permite como máximo tres ejecuciones de transacción, con dos esperas. La
espera aleatoria base va de 20 a 80 ms y se multiplica por attempt+1:
primero 20-80 ms; después 40-160 ms.

Si el hilo se interrumpe, se restablece su indicador y se lanza 503.
Otros errores no activan este bucle. Una desconexión de respuesta con
resultado incierto no se transforma en un reintento ciego ilimitado. El
cliente dispone de la consulta autenticada de resultado.

\begin{Shaded}
\begin{Highlighting}[]
\NormalTok{autorizar y elegir propietario}
\NormalTok{para intento = 0, 1, 2:}
\NormalTok{    ejecutar una transacción en ese propietario}
\NormalTok{    si éxito: devolver JSON canónico}
\NormalTok{    si SQLSTATE es 40001 o 40P01 y aún hay intento:}
\NormalTok{        esperar un intervalo acotado}
\NormalTok{        repetir con la misma clave, huella y correlación}
\NormalTok{    en otro caso: propagar error}
\NormalTok{registrar siempre duración y resultado público}
\end{Highlighting}
\end{Shaded}

El catch de \texorpdfstring{\softtext{ApiException}}{ApiException} copia
su código y clasifica RECHAZADA para estados menores de 500 o FALLO para
los demás. El catch de
\texorpdfstring{\softtext{RuntimeException}}{RuntimeException} obtiene
un mapeo para registrar un código seguro, pero vuelve a lanzar la
excepción original. El advice HTTP será quien produzca la respuesta
\texorpdfstring{\softtext{correspondiente.}}{correspondiente.}

El finally suma elapsed en nanosegundos, incrementa un contador y crea
event con evento, comando,
\texorpdfstring{\softtext{correlation\_id,}}{correlation\_id,}
\texorpdfstring{\softtext{nodo\_dueno,}}{nodo\_dueno,}
\texorpdfstring{\softtext{duracion\_ms}}{duracion\_ms} y resultado.
Puede incluir codigo. No agrega payload, clave, hash de petición o
\texorpdfstring{\softtext{contraseña.}}{contraseña.} Este registro ayuda
a localizar una ejecución; la auditoría económica permanece en
\texorpdfstring{\softtext{PostgreSQL.}}{PostgreSQL.}

\subsection{\texorpdfstring{query,
\texorpdfstring{\softtext{providerTariffs,}}{providerTariffs,}
timedQuery y
diagnostics}{query,  timedQuery y diagnostics}}\label{query-providertariffs-timedquery-y-diagnostics}

\textbf{query(user, site, resource, params)} autoriza lectura. usuarios
exige ADMIN. Un lookup de comando maestro también usa A. Contrapartes,
animales, animal, predios y usuarios se consultan en A; los demás
recursos van al propietario de la sede. Ejecuta
\texorpdfstring{\softtext{bovi.consultar}}{bovi.consultar} con empresa,
sede, usuario, recurso y params JSON mediante
\texorpdfstring{\softtext{timedQuery.}}{timedQuery.} No abre la
\texorpdfstring{\softtext{TransactionTemplate}}{TransactionTemplate} de
un comando para una consulta ordinaria.

\textbf{\texorpdfstring{\softtext{providerTariffs(user,}}{providerTariffs(user,}
site, params)} autoriza lectura y hace dos llamadas separadas. Primero
timedQuery en A invoca
\texorpdfstring{\softtext{leer\_tarifas\_proveedor.}}{leer\_tarifas\_proveedor.}
quotes contiene la página de fuente y contexto interno. Luego otra
llamada, en el dueño de la sede, invoca
\texorpdfstring{\softtext{capturar\_tarifas\_proveedor}}{capturar\_tarifas\_proveedor}
con ese resultado interno. Las conexiones trabajan en autocommit y la
primera llamada termina antes de iniciar la segunda. El cliente no
entrega una página supuestamente leída del proveedor.

\textbf{\texorpdfstring{\softtext{timedQuery(node,}}{timedQuery(node,}
resource, query)} recibe un Supplier. started mide el comienzo; outcome
y errorCode describen el resultado. Un fallo incrementa
\texorpdfstring{\softtext{queryErrors}}{queryErrors} y se propaga. El
finally incrementa queries y
\texorpdfstring{\softtext{queryNanos.}}{queryNanos.} Si dura al menos un
segundo, incrementa \texorpdfstring{\softtext{slowQueries}}{slowQueries}
y escribe \texorpdfstring{\softtext{CONSULTA\_LENTA}}{CONSULTA\_LENTA}
con recurso, nodo, duración y resultado. No imprime la sentencia ni los
\texorpdfstring{\softtext{parámetros.}}{parámetros.}

\textbf{diagnostics()} devuelve mapas con instante de proceso,
advertencia de reinicio de contadores, comandos completados y errores,
duración total, consultas totales, errores, lentas, umbral de 1000 ms y
duración total. Las conversiones de duración son unidades temporales, no
operaciones \texorpdfstring{\softtext{monetarias.}}{monetarias.}

Las denegaciones que ocurren antes de timedQuery no forman parte de sus
contadores de consulta. Tampoco se cuenta cada select de
\texorpdfstring{\softtext{AccessService}}{AccessService} como consulta
económica medida. Esta distinción evita llamar al diagnóstico un medidor
exhaustivo de toda actividad JDBC.

\section{\texorpdfstring{\texorpdfstring{\softtext{CoreController:}}{CoreController:}
un adaptador HTTP
compartido}{ un adaptador HTTP compartido}}\label{corecontroller-un-adaptador-http-compartido}

Referencia:
\texorpdfstring{\softcode{backend/src/main/java/mx/bovicierre/operaciones/CoreController.java:16}}{backend/src/main/java/mx/bovicierre/operaciones/CoreController.java:16}.
Los campos access y commands se fijan en el constructor
\textbf{\texorpdfstring{\softtext{CoreController(access,}}{CoreController(access,}
commands)}. Esta clase concentra las familias de recursos; no contiene
el cálculo de costo, saldo o margen.

\subsection{Los métodos auxiliares}\label{los-muxe9todos-auxiliares}

\textbf{relative(request)} toma uri y elimina el prefijo hasta la sede.
Por ejemplo, de una ruta completa extrae
\texorpdfstring{\softtext{operaciones/\{id\}/confirmar.}}{operaciones/\{id\}/confirmar.}
Se utiliza para determinar qué recurso o identificador corresponde a la
ruta \texorpdfstring{\softtext{registrada.}}{registrada.} Está apoyado
en la estructura de esas rutas, no es un router de
\texorpdfstring{\softtext{PostgreSQL.}}{PostgreSQL.}

\textbf{result(json, status)} prepara
\texorpdfstring{\softtext{ResponseEntity}}{ResponseEntity} con
\texorpdfstring{\softtext{application/json.}}{application/json.} Lee
node para saber si hay
\texorpdfstring{\softtext{version\_lock}}{version\_lock} y, en ese caso,
agrega un ETag con ese número entre comillas. Devuelve el texto JSON
recibido. No agrega una versión inventada a objetos que no la contienen.

\textbf{write(site, command, body, key, match, idKey, id, version,
status, authentication, request)} concentra el camino de escritura.
payload se construye como objeto JSON; cuerpo ausente o blanco se
convierte en objeto vacío. Rechaza campos de ámbito, origen y
\texorpdfstring{\softtext{version\_lock}}{version\_lock} enviados en el
cuerpo. Si hay id de ruta, lo inserta usando idKey.

Para \texorpdfstring{\softtext{corregir\_pesaje}}{corregir\_pesaje}
comprueba que \texorpdfstring{\softtext{corrige\_a\_id}}{corrige\_a\_id}
opcional sea textual, interpretable como UUID y coincidente con la ruta.
reason debe ser motivo textual no blanco de hasta 200 unidades de
longitud Java. Se fija después
\texorpdfstring{\softtext{corrige\_a\_id}}{corrige\_a\_id} al original
de la ruta. El SQL V13 comprueba el conjunto de animales; el controlador
no pretende reconstruir esa composición desde el lote actual.

Cuando version es true, write obtiene
\texorpdfstring{\softtext{version\_lock}}{version\_lock} de
\texorpdfstring{\softtext{JsonSupport.version(match).}}{JsonSupport.version(match).}
La clave match representa If-Match, no un entero libre del cuerpo. La
llamada a \texorpdfstring{\softtext{commands.command}}{commands.command}
proporciona usuario del contexto y correlación de request. La variable
json recibe el resultado y response lo
\texorpdfstring{\softtext{representa.}}{representa.} Para
\texorpdfstring{\softtext{registrar\_gasto}}{registrar\_gasto} se
obtiene gasto\_id o id y se agrega Location de su recurso, con el mismo
cuerpo \texorpdfstring{\softtext{confirmado.}}{confirmado.}

\textbf{integer(value)} convierte un texto a int; si falla, BV400 de
\texorpdfstring{\softtext{paginación.}}{paginación.}
\textbf{\texorpdfstring{\softtext{validateIdFilter(name,}}{validateIdFilter(name,}
value)} valida únicamente periodo\_id y
\texorpdfstring{\softtext{operacion\_id.}}{operacion\_id.} Además de
\texorpdfstring{\softtext{UUID.fromString,}}{UUID.fromString,} exige que
la representación canónica coincida sin distinguir
\texorpdfstring{\softtext{mayúsculas.}}{mayúsculas.} Esto rechaza formas
abreviadas que un parser permisivo podría aceptar.

\subsection{Lecturas públicas de esta
clase}\label{lecturas-puxfablicas-de-esta-clase}

\textbf{\texorpdfstring{\softtext{compare(sedeId,}}{compare(sedeId,}
\texorpdfstring{\softtext{inicial\_id,}}{inicial\_id,} final\_id, auth)}
construye params de dos UUID y consulta
\texorpdfstring{\softtext{comparacion\_pesajes.}}{comparacion\_pesajes.}
El cálculo de variación o la respuesta
\texorpdfstring{\softtext{NO\_COMPARABLE}}{NO\_COMPARABLE} se realiza en
SQL; Java no divide
\texorpdfstring{\softtext{porcentajes.}}{porcentajes.}

\textbf{\texorpdfstring{\softtext{commandResult(sedeId,}}{commandResult(sedeId,}
comando, clave, auth)} limita comando a letras minúsculas y guion bajo,
de 3 a 60 caracteres, y la clave a 16-80. Consulta el recurso comando
con esos filtros. La base limita por principal y sede, y
\texorpdfstring{\softtext{CommandService}}{CommandService} dirige los
maestros a A. Una ausencia devuelve un error de recurso; no prueba que
el intento original jamás haya
\texorpdfstring{\softtext{confirmado.}}{confirmado.}

\textbf{\texorpdfstring{\softtext{list(sedeId,}}{list(sedeId,} query,
auth, request)} convierte la ruta relativa a resource reemplazando
barras por guiones bajos. Rechaza periodo\_id para
\texorpdfstring{\softtext{reportes\_cartera}}{reportes\_cartera} porque
esa cartera es actual; la historia del cierre se consulta por su
versión. params recibe UUID verificados y booleanos reales para
\texorpdfstring{\softtext{disponibles\_para\_venta}}{disponibles\_para\_venta}
y vencida. Otros filtros siguen siendo textos. page y size tienen
valores predeterminados 1 y 25, y límites página positiva y tamaño
1-100.

\textbf{\texorpdfstring{\softtext{one(sedeId,}}{one(sedeId,} id, query,
auth, request)} identifica el recurso específico por terminación o
prefijo:
\texorpdfstring{\softtext{periodo\_versiones,}}{periodo\_versiones,}
\texorpdfstring{\softtext{operacion\_pesajes,}}{operacion\_pesajes,}
\texorpdfstring{\softtext{operacion\_expediente,}}{operacion\_expediente,}
simulador, animal, operacion, liquidacion, pago,
\texorpdfstring{\softtext{reparto\_gasto}}{reparto\_gasto} o gasto.
Revisa los filtros UUID, agrega el id de ruta y consulta. La variable
relative aquí es un texto local distinto del método que lo produjo.

\subsection{Escrituras y selección de
acciones}\label{escrituras-y-selecciuxf3n-de-acciones}

\textbf{create(\ldots)} obtiene suffix, el último segmento, y busca un
comando en un mapa fijo. Contrapartes, predios y animales son maestros;
operaciones, pagos, gastos, periodos y movimientos utilizan su selector
\texorpdfstring{\softtext{correspondiente.}}{correspondiente.} Pasa
status 201 a write.

\textbf{edit(\ldots)} atiende PATCH de
\texorpdfstring{\softtext{operación.}}{operación.} Llama write con
\texorpdfstring{\softtext{editar\_operacion,}}{editar\_operacion,}
\texorpdfstring{\softtext{operacion\_id,}}{operacion\_id,} el UUID de
ruta, comprobación de versión y respuesta 200.

\textbf{action(\ldots)} obtiene suffix y selecciona el comando con
switch. Para reabrir distingue una operación de un periodo examinando la
ruta. Calcula idKey según la primera familia: operación, pesaje,
reserva, liquidación, gasto, pago, aplicación o periodo. También calcula
status: 201 para nuevas entidades o movimientos de ciertas acciones; 200
para otras.

Las acciones confirmar, cancelar, anular y acordar requieren versión, al
igual que
\texorpdfstring{\softtext{reabrir\_operacion.}}{reabrir\_operacion.}
\texorpdfstring{\softtext{Reabrir\_periodo}}{Reabrir\_periodo} no
utiliza en este controlador la misma bandera. La protección de periodo y
su estado se realiza en el comando canónico; no debe atribuirse al
método una comprobación If-Match que no ejecuta.

{\def\LTcaptype{none} % do not increment counter
\begin{longtable}[]{@{}
  >{\raggedright\arraybackslash}p{(\linewidth - 4\tabcolsep) * \real{0.3333}}
  >{\raggedright\arraybackslash}p{(\linewidth - 4\tabcolsep) * \real{0.3333}}
  >{\raggedright\arraybackslash}p{(\linewidth - 4\tabcolsep) * \real{0.3333}}@{}}
\toprule\noalign{}
\begin{minipage}[b]{\linewidth}\raggedright
Familia de acción
\end{minipage} & \begin{minipage}[b]{\linewidth}\raggedright
Selector SQL elegido
\end{minipage} & \begin{minipage}[b]{\linewidth}\raggedright
Respuesta de este adaptador
\end{minipage} \\
\midrule\noalign{}
\endhead
\bottomrule\noalign{}
\endlastfoot
acordar / reabrir operación / confirmar &
\texorpdfstring{\softtext{acordar\_operacion}}{acordar\_operacion} /
\texorpdfstring{\softtext{reabrir\_operacion}}{reabrir\_operacion} /
\texorpdfstring{\softtext{confirmar\_operacion}}{confirmar\_operacion} &
200; versión requerida \\
cancelar / anular &
\texorpdfstring{\softtext{cancelar\_operacion}}{cancelar\_operacion} /
\texorpdfstring{\softtext{anular\_venta}}{anular\_venta} & 200; versión
requerida \\
pesajes / aprobar / correcciones &
\texorpdfstring{\softtext{registrar\_pesaje}}{registrar\_pesaje} /
\texorpdfstring{\softtext{aprobar\_pesaje}}{aprobar\_pesaje} /
\texorpdfstring{\softtext{corregir\_pesaje}}{corregir\_pesaje} & 201 /
200 / 201 \\
reservas / liberar &
\texorpdfstring{\softtext{reservar\_venta}}{reservar\_venta} /
\texorpdfstring{\softtext{liberar\_reserva}}{liberar\_reserva} & 201 /
200 \\
ajustes / reparto &
\texorpdfstring{\softtext{ajustar\_liquidacion}}{ajustar\_liquidacion} /
\texorpdfstring{\softtext{repartir\_gasto}}{repartir\_gasto} & 201 /
201 \\
validar / aplicaciones / reembolsos &
\texorpdfstring{\softtext{validar\_pago}}{validar\_pago} /
\texorpdfstring{\softtext{aplicar\_pago}}{aplicar\_pago} /
\texorpdfstring{\softtext{reembolsar\_pago}}{reembolsar\_pago} & 200 /
201 / 201 \\
reverso &
\texorpdfstring{\softtext{revertir\_aplicacion}}{revertir\_aplicacion} &
201 \\
cierre o cerrar / reabrir periodo &
\texorpdfstring{\softtext{cerrar\_periodo}}{cerrar\_periodo} /
\texorpdfstring{\softtext{reabrir\_periodo}}{reabrir\_periodo} & 201 /
200 \\
revisiones de operación &
\texorpdfstring{\softtext{registrar\_revision}}{registrar\_revision} &
200 \\
\end{longtable}
}

El controlador conserva alias como cerrar junto a cierre y una ruta de
revisiones de \texorpdfstring{\softtext{operación.}}{operación.} El
inventario de OpenAPI más adelante enumera las rutas contratadas; las
anotaciones pueden incluir compatibilidades
\texorpdfstring{\softtext{adicionales.}}{adicionales.}

\section{\texorpdfstring{\texorpdfstring{\softtext{ApiException}}{ApiException}
y ApiErrors: del fallo interno al error
público}{ y ApiErrors: del fallo interno al error público}}\label{apiexception-y-apierrors-del-fallo-interno-al-error-puxfablico}

\subsection{La excepción
transportable}\label{la-excepciuxf3n-transportable}

Referencia:
\texorpdfstring{\softcode{backend/src/main/java/mx/bovicierre/shared/ApiException.java:3}}{backend/src/main/java/mx/bovicierre/shared/ApiException.java:3}.
\texorpdfstring{\softtext{ApiException}}{ApiException} extiende
\texorpdfstring{\softtext{RuntimeException}}{RuntimeException} y tiene
status, code, rule y fields. status y code son finales. rule puede
quedar nula y fields inicia como lista vacía.
\textbf{\texorpdfstring{\softtext{ApiException(status,}}{ApiException(status,}
code, message)} guarda el estado HTTP y código, y entrega message al
constructor de
\texorpdfstring{\softtext{RuntimeException.}}{RuntimeException.}

No incluye una conexión ni una sentencia SQL. Su mensaje debe ser apto
para el usuario. El detalle de campos es una lista de mapas con nombres
públicos; el advice puede publicar esas piezas.

\subsection{Métodos de ApiErrors}\label{muxe9todos-de-apierrors}

Referencia:
\texorpdfstring{\softcode{backend/src/main/java/mx/bovicierre/shared/ApiErrors.java:15}}{backend/src/main/java/mx/bovicierre/shared/ApiErrors.java:15}.
La clase
\texorpdfstring{\softtext{RestControllerAdvice}}{RestControllerAdvice}
no tiene campos \texorpdfstring{\softtext{persistentes.}}{persistentes.}
\textbf{correlation(request)} reutiliza un UUID del atributo
\texorpdfstring{\softtext{correlation\_id}}{correlation\_id} si existe;
de lo contrario crea id, lo guarda y lo devuelve. El filtro y los
controladores usan el mismo atributo para mantener una correlación por
solicitud.

\textbf{body(ex, id)} construye un
\texorpdfstring{\softtext{LinkedHashMap}}{LinkedHashMap} con codigo,
mensaje, \texorpdfstring{\softtext{correlation\_id}}{correlation\_id} y
campos; agrega regla\_id cuando no es nula. \textbf{write(response, ex,
id)} se usa desde filtros, fuera del camino habitual del
\texorpdfstring{\softtext{controlador.}}{controlador.} Fija estado,
codificación UTF-8,
\texorpdfstring{\softtext{application/problem+json}}{application/problem+json}
y X-Correlation-ID antes de escribir el JSON.

\textbf{api(ex, request)} es el handler de
\texorpdfstring{\softtext{ApiException:}}{ApiException:} devuelve
\texorpdfstring{\softtext{ResponseEntity}}{ResponseEntity} con estado,
media type de problema y body. \textbf{oversized(ex, request)} traduce
\texorpdfstring{\softtext{MaxUploadSizeExceededException}}{MaxUploadSizeExceededException}
a 413/BV023. \textbf{invalid(ex, request)} convierte cabeceras
faltantes, incompatibilidad de tipos, mensajes HTTP ilegibles y
validación del argumento en 400/BV400.

Hay dos métodos llamados \textbf{database}. La sobrecarga estática
\textbf{database(Throwable ex)} examina el fallo. La sobrecarga handler
\textbf{\texorpdfstring{\softtext{database(RuntimeException}}{database(RuntimeException}
ex, request)} atiende
\texorpdfstring{\softtext{DataAccessException}}{DataAccessException} y
\texorpdfstring{\softtext{TransactionException,}}{TransactionException,}
llama a la anterior y después a api. \textbf{internal()} devuelve
500/BV500 con un mensaje genérico.

\subsection{Cómo se valida el DETAIL de
PostgreSQL}\label{cuxf3mo-se-valida-el-detail-de-postgresql}

database(Throwable) avanza por cause hasta hallar
\texorpdfstring{\softtext{SQLException.}}{SQLException.} Si no existe,
una imposibilidad de conexión JDBC reconocida produce 503; los demás
fallos genéricos producen internal. Cuando SQLSTATE es P0001 y hay
\texorpdfstring{\softtext{ServerErrorMessage,}}{ServerErrorMessage,}
interpreta DETAIL como JSON.

error es el árbol del detalle. code debe cumplir BV seguido de una
categoría 0, 4 o 5 y dos dígitos; rule debe cumplir RN y tres dígitos.
fields debe ser un array de hasta 50 elementos. Cada elemento necesita
campo y problema no blancos, de hasta 160 caracteres y sin controles.
\texorpdfstring{\softtext{publicFields}}{publicFields} copia únicamente
esas dos \texorpdfstring{\softtext{propiedades.}}{propiedades.} Un
detalle mal formado no se repite en la respuesta: se transforma en error
interno.

El mensaje que escribió PostgreSQL no se publica sin control. El mapper
elige mensajes españoles conocidos para clave repetida, versión,
recurso, disponibilidad, reserva, sobreaplicación y periodo; usa
mensajes generales para los demás. Esto conserva trazabilidad mediante
códigos sin filtrar nombres internos, fragmentos SQL o datos privados de
una \texorpdfstring{\softtext{excepción.}}{excepción.}

{\def\LTcaptype{none} % do not increment counter
\begin{longtable}[]{@{}
  >{\raggedright\arraybackslash}p{(\linewidth - 4\tabcolsep) * \real{0.3333}}
  >{\raggedright\arraybackslash}p{(\linewidth - 4\tabcolsep) * \real{0.3333}}
  >{\raggedright\arraybackslash}p{(\linewidth - 4\tabcolsep) * \real{0.3333}}@{}}
\toprule\noalign{}
\begin{minipage}[b]{\linewidth}\raggedright
Error reconocido
\end{minipage} & \begin{minipage}[b]{\linewidth}\raggedright
Resultado HTTP
\end{minipage} & \begin{minipage}[b]{\linewidth}\raggedright
Lectura práctica
\end{minipage} \\
\midrule\noalign{}
\endhead
\bottomrule\noalign{}
\endlastfoot
BV412, BV428, BV404 & 412, 428, 404 & Versión obsoleta, precondición
ausente o recurso no disponible \\
Otros BV4xx del detalle & 409 & Conflicto; BV403 es reserva vencida, no
permiso HTTP403 \\
BV0xx del detalle & 422 & Regla de dominio incumplida \\
BV5xx del detalle o conexión reconocida & 503 & Dependencia requerida no
disponible \\
23505 con restricción uk\_ & 409/BV401 & Conflicto de unicidad
conocido \\
23001 con fk\_ & 409/BV401 & Referencias impiden modificar \\
23503 con fk\_ o CHECK reconocido & 422/BV015 & Referencia de ámbito o
condición declarativa \\
23502, 22P02, 22001, 22003, 22007 & 422/BV015 & Dato requerido, tipo,
longitud, rango o fecha inválidos \\
42501 & \texorpdfstring{\softtext{403/AUTH403}}{403/AUTH403} & Permiso
rechazado \\
40001 o 40P01 tras los reintentos & 409/BV409 & Conflicto concurrente,
conservar la misma clave \\
\end{longtable}
}

Los nombres de restricción se revisan para las categorías que lo
necesitan; no toda excepción se presenta como
\texorpdfstring{\softtext{indisponibilidad.}}{indisponibilidad.} Una SQL
no reconocida termina como 500, lo cual permite distinguir un defecto de
programación de un dueño caído.

Un ejemplo público, sin valores reales de sesión, es:

\begin{Shaded}
\begin{Highlighting}[]
\FunctionTok{\{}
  \DataTypeTok{"codigo"}\FunctionTok{:} \StringTok{"BV412"}\FunctionTok{,}
  \DataTypeTok{"mensaje"}\FunctionTok{:} \StringTok{"La versión cambió; recargue el recurso."}\FunctionTok{,}
  \DataTypeTok{"correlation\_id"}\FunctionTok{:} \StringTok{"00000000{-}0000{-}4000{-}8000{-}000000000001"}\FunctionTok{,}
  \DataTypeTok{"campos"}\FunctionTok{:} \OtherTok{[]}\FunctionTok{,}
  \DataTypeTok{"regla\_id"}\FunctionTok{:} \StringTok{"RN{-}017"}
\FunctionTok{\}}
\end{Highlighting}
\end{Shaded}

La correlación del ejemplo es
\texorpdfstring{\softtext{ilustrativa.}}{ilustrativa.} Para soporte se
conserva el UUID de la respuesta real; no se necesita enviar al operador
una contraseña o el texto de una sentencia para localizar el intento.

\section{Archivos privados: puerto, adaptador, controlador y
limpieza}\label{archivos-privados-puerto-adaptador-controlador-y-limpieza}

\subsection{\texorpdfstring{\texorpdfstring{\softtext{DocumentStorage:}}{DocumentStorage:}
un contrato pequeño y
real}{ un contrato pequeño y real}}\label{documentstorage-un-contrato-pequeuxf1o-y-real}

Referencia:
\texorpdfstring{\softcode{backend/src/main/java/mx/bovicierre/evidencia/application/DocumentStorage.java:7}}{backend/src/main/java/mx/bovicierre/evidencia/application/DocumentStorage.java:7}.
La interfaz expone
\textbf{\texorpdfstring{\softtext{stage(MultipartFile}}{stage(MultipartFile}
file)} y \textbf{path(String sha256)}. stage devuelve un Upload cuyo
ciclo de vida debe cerrarse; path devuelve una ruta local válida para
contenido autorizado previamente por el llamador.

El record \textbf{Stored} contiene sha256, size, mime y name, con
accesores sha256(), size(), mime() y name(). La interfaz anidada
\textbf{Upload} extiende
\texorpdfstring{\softtext{AutoCloseable}}{AutoCloseable} y expone
stored() y close(). El controlador puede usar try-with-resources sin
conocer FileLock o Hikari. Los metadatos de negocio siguen perteneciendo
al comando SQL.

\subsection{\texorpdfstring{\texorpdfstring{\softtext{PrivateDocuments:}}{PrivateDocuments:}
almacenamiento por
huella}{ almacenamiento por huella}}\label{privatedocuments-almacenamiento-por-huella}

Referencia:
\texorpdfstring{\softcode{backend/src/main/java/mx/bovicierre/evidencia/PrivateDocuments.java:27}}{backend/src/main/java/mx/bovicierre/evidencia/PrivateDocuments.java:27}.
\textbf{root} es el directorio absoluto normalizado; \textbf{LOCKS} es
un \texorpdfstring{\softtext{ConcurrentHashMap}}{ConcurrentHashMap}
estático de bloqueos locales por directorio y huella. El constructor
\textbf{\texorpdfstring{\softtext{PrivateDocuments(root)}}{PrivateDocuments(root)}}
normaliza la ruta \texorpdfstring{\softtext{configurada.}}{configurada.}
No acepta una ruta física como parámetro del usuario HTTP.

\textbf{store(file)} es un método de conveniencia: llama stage y
devuelve stored dentro de un try que inmediatamente cierra el Upload.
\texorpdfstring{\softtext{DocumentController}}{DocumentController}
utiliza directamente stage para conservar el bloqueo durante el comando
SQL. No debe atribuirse a store el mismo alcance de bloqueo alrededor de
una transacción que sucede después de
\texorpdfstring{\softtext{devolverlo.}}{devolverlo.}

La clase interna \textbf{Upload} conserva stored y lock. Su constructor
privado une el resultado de escritura y el HeldLock. \textbf{stored()}
devuelve los metadatos del contenido; \textbf{close()} delega la
\texorpdfstring{\softtext{liberación.}}{liberación.} El record privado
\textbf{HeldLock} contiene local, channel y file: respectivamente
\texorpdfstring{\softtext{ReentrantLock,}}{ReentrantLock,}
\texorpdfstring{\softtext{FileChannel}}{FileChannel} y FileLock. Su
\textbf{close()} libera primero el bloqueo de archivo, cierra el canal y
desbloquea el \texorpdfstring{\softtext{ReentrantLock}}{ReentrantLock}
en finally; si aparece IOException, registra un aviso genérico. Sus
accesores generados no agregan reglas nuevas.

\textbf{directories()} crea root y .locks y aplica permisos POSIX
rwx------ a ambos. \textbf{lock(hash, wait)} obtiene el
\texorpdfstring{\softtext{ReentrantLock}}{ReentrantLock} con
\texorpdfstring{\softtext{computeIfAbsent.}}{computeIfAbsent.} En modo
wait bloquea; en modo sin espera usa tryLock y puede devolver null. Abre
el archivo
\texorpdfstring{\softtext{.locks/\{hash\}.lock}}{.locks/\{hash\}.lock}
con permisos rw------- y obtiene un FileLock entre procesos. local
impide colisiones dentro de la JVM; el bloqueo de archivo coordina
procesos que comparten
\texorpdfstring{\softtext{directorio.}}{directorio.} Si falla, cierra
channel y libera local.

\textbf{stage(file)} comprueba archivo no vacío y máximo de 20 MiB.
bytes contiene el contenido completo en memoria. mime proviene de type y
hash de SHA-256. Crea directorios y obtiene held esperando el bloqueo.
temporary tiene una ruta de huella, UUID aleatorio y sufijo .part.

El método crea y escribe temporary, fija sus permisos y la mueve
atómicamente a
\texorpdfstring{\softtext{root/\{hash\},}}{root/\{hash\},} con reemplazo
permitido. name toma el nombre original o documento; reemplaza
caracteres no permitidos y lo limita a 160. Devuelve Upload con Stored y
el bloqueo aún retenido. El finally elimina la parte temporal si
persiste. Un error de E/S produce BV503; no se realiza un movimiento no
atómico \texorpdfstring{\softtext{alternativo.}}{alternativo.}

\textbf{type(bytes)} examina firmas iniciales: \%PDF-, los ocho bytes
PNG o los tres primeros JPEG. Devuelve
\texorpdfstring{\softtext{application/pdf,}}{application/pdf,} image/png
o \texorpdfstring{\softtext{image/jpeg.}}{image/jpeg.} Otro contenido
produce 422/BV023. Es detección de firma, no un análisis estructural
integral ni un escáner
\texorpdfstring{\softtext{antimalware.}}{antimalware.}

\textbf{path(sha256)} exige exactamente 64 caracteres hexadecimales
minúsculos, resuelve bajo root y exige archivo regular sin seguir enlace
simbólico final. Lee bytes y recalcula hash. Una diferencia produce 409
de integridad; indisponibilidad de lectura, 503. Devuelve Path
únicamente después de esa
\texorpdfstring{\softtext{comprobación.}}{comprobación.} Quien invoca
path debe haber autorizado el documento; conocer su huella no sustituye
ese paso.

\subsection{\texorpdfstring{cleanup y
\texorpdfstring{\softtext{CleanupMain:}}{CleanupMain:} seguridad por
candidato}{cleanup y  seguridad por candidato}}\label{cleanup-y-cleanupmain-seguridad-por-candidato}

\textbf{cleanup(execute, referenced)} devuelve el record
\textbf{Cleanup}, cuyos campos son modo, candidatos, eliminados y
omitidos y cuyos accesores tienen esos nombres. Las listas de resultado
se copian antes de
\texorpdfstring{\softtext{devolverlas.}}{devolverlas.}

paths recorre el \texorpdfstring{\softtext{directorio.}}{directorio.}
Solo considera archivos regulares, sin seguir enlaces, que parezcan una
huella o un .part asociado y sean anteriores a 24 horas. hash extrae los
primeros 64 \texorpdfstring{\softtext{caracteres.}}{caracteres.} Con un
bloqueo sin espera, omite contenido ocupado. Después vuelve a comprobar
existencia, antigüedad y
\texorpdfstring{\softtext{referenced.test(hash)}}{referenced.test(hash)}
bajo bloqueo. Solo entonces agrega un candidato y, si execute es true,
lo elimina.

La limpieza predeterminada
\texorpdfstring{\softtext{diagnostica.}}{diagnostica.} No elimina un
archivo reciente ni uno que ambos chequeos no hayan permitido considerar
huérfano. Tampoco es una transacción global de sistema de archivos: si
falla una comprobación después de borrar candidatos previos, aquellos
borrados ya realizados no se revierten.

Referencia del lanzador:
\texorpdfstring{\softcode{backend/src/main/java/mx/bovicierre/evidencia/CleanupMain.java:9}}{backend/src/main/java/mx/bovicierre/evidencia/CleanupMain.java:9}.
El constructor privado impide uso ordinario como objeto.
\textbf{main(args)} acepta cero argumentos o únicamente
\texorpdfstring{\softtext{–execute.}}{–execute.} env contiene el
entorno; secrets selecciona directorio externo por
\texorpdfstring{\softtext{BOVI\_SECRETS\_DIR,}}{BOVI\_SECRETS\_DIR,}
alternativa
\texorpdfstring{\softtext{BOVI\_SECRET\_DIR}}{BOVI\_SECRET\_DIR} o
ubicación \texorpdfstring{\softtext{predeterminada.}}{predeterminada.}
password se lee del archivo
\texorpdfstring{\softtext{worker\_password}}{worker\_password} sin
\texorpdfstring{\softtext{imprimirse.}}{imprimirse.} directory determina
el \texorpdfstring{\softtext{almacenamiento.}}{almacenamiento.}

Abre conexiones a y b con rol
\texorpdfstring{\softtext{bovi\_worker}}{bovi\_worker} y prepara qa y qb
para
\texorpdfstring{\softtext{archivo\_referenciado(?),}}{archivo\_referenciado(?),}
cada una con timeout de cinco segundos. El Predicate usado por cleanup
ejecuta ambos chequeos y devuelve inA o inB. Cualquier fallo impide
eliminar ese candidato y corta el procedimiento; si ni siquiera se
pueden abrir ambas conexiones, no inicia el escaneo. result se imprime
como JSON de \texorpdfstring{\softtext{diagnóstico.}}{diagnóstico.} Este
main no inicia Spring, HTTP ni los trabajadores
\texorpdfstring{\softtext{programados.}}{programados.}

\subsection{\texorpdfstring{\texorpdfstring{\softtext{DocumentController}}{DocumentController}
y sus cuatro
operaciones}{ y sus cuatro operaciones}}\label{documentcontroller-y-sus-cuatro-operaciones}

Referencia:
\texorpdfstring{\softcode{backend/src/main/java/mx/bovicierre/evidencia/DocumentController.java:16}}{backend/src/main/java/mx/bovicierre/evidencia/DocumentController.java:16}.
Tiene files de tipo
\texorpdfstring{\softtext{DocumentStorage,}}{DocumentStorage,} access y
commands. El constructor
\textbf{\texorpdfstring{\softtext{DocumentController(files,}}{DocumentController(files,}
access, commands)} permite inyectar el adaptador local mediante el
puerto.

\textbf{upload(\ldots)} recibe
\texorpdfstring{\softtext{MultipartFile}}{MultipartFile} archivo,
\texorpdfstring{\softtext{tipo\_documental,}}{tipo\_documental,}
operación opcional, documento previo opcional, clave y contexto. user y
scope se obtienen del servidor. Rechaza tipo documental blanco o mayor
de 60 \texorpdfstring{\softtext{caracteres.}}{caracteres.} El try de
\texorpdfstring{\softtext{files.stage}}{files.stage} conserva upload
hasta que termine el comando.

stored aporta huella, tamaño, MIME y nombre saneado. payload incorpora
esos datos y \texorpdfstring{\softtext{storage\_key:}}{storage\_key:}
una huella de almacenamiento que combina empresa, sede, usuario, clave y
huella del contenido. Esta clave técnica no es la ruta del archivo
físico. Si hay operación, se agrega su ID. Si hay documento previo,
previous se consulta con autorización, se agrega la referencia y se
calcula version como previa más uno.

El comando
\texorpdfstring{\softtext{registrar\_documento}}{registrar\_documento}
usa una huella explícita de intención que incluye el nombre original
recibido. No solo compara los bytes: cambiar metadatos o nombre con la
misma clave puede entrar en conflicto. Devuelve 201 con el JSON
canónico. Al salir del bloque se libera el bloqueo físico.

Publicar el contenido ocurre antes del SQL. Si el SQL rechaza o la
respuesta queda incierta, puede quedar un objeto sin
\texorpdfstring{\softtext{referencia.}}{referencia.} No se borra
precipitadamente: podría existir un commit cuyo resultado no llegó. La
antigüedad, el chequeo de referencias A/B y el mismo bloqueo sustentan
la limpieza posterior.

\textbf{\texorpdfstring{\softtext{metadata(sedeId,}}{metadata(sedeId,}
id, auth)} consulta documento y devuelve JSON. \textbf{download(\ldots)}
consulta esos metadatos con ámbito autorizado, obtiene metadata,
verifica files.path de sha256 y entrega
\texorpdfstring{\softtext{FileSystemResource.}}{FileSystemResource.} Usa
\texorpdfstring{\softtext{application/octet-stream,}}{application/octet-stream,}
nombre de descarga basado en el UUID y X-Content-Type-Options: nosniff.
El navegador no recibe una ruta libre a abrir en disco.

\textbf{review(\ldots)} parsea el cuerpo. Si falta documentos, construye
una lista con el documento de la ruta y resultado. Si el cliente ya
proporcionó la lista, no la reemplaza. Delega
\texorpdfstring{\softtext{registrar\_revision,}}{registrar\_revision,}
cuya autorización ADMIN y validación de
\texorpdfstring{\softtext{vínculos/versiones}}{vínculos/versiones}
corresponden al servicio y SQL. Devuelve 200. Es importante no confundir
esta construcción condicional con la comprobación estricta
\texorpdfstring{\softtext{corrige\_a\_id}}{corrige\_a\_id} del
controlador de pesajes.

\section{Integraciones: tarifas, outbox y estado de
operación}\label{integraciones-tarifas-outbox-y-estado-de-operaciuxf3n}

\subsection{\texorpdfstring{\texorpdfstring{\softtext{ProviderController}}{ProviderController}}{}}\label{providercontroller}

Referencia:
\texorpdfstring{\softcode{backend/src/main/java/mx/bovicierre/integraciones/ProviderController.java:15}}{backend/src/main/java/mx/bovicierre/integraciones/ProviderController.java:15}.
El constructor recibe y conserva access y commands.
\textbf{\texorpdfstring{\softtext{tariffs(sedeId,}}{tariffs(sedeId,}
query, auth)} convierte pagina y tamano, por defecto 1 y 25, limita
tamaño a 100 y ruta a 160
\texorpdfstring{\softtext{caracteres.}}{caracteres.} params contiene
números de paginación y, si existe, ruta. Devuelve
\texorpdfstring{\softtext{application/json}}{application/json} con
no-store mediante
\texorpdfstring{\softtext{providerTariffs.}}{providerTariffs.}

La consulta se lee exclusivamente en A. Después se guarda la captura en
el dueño local. B no necesita consultar directamente el servidor externo
para aceptar una captura ya recibida. El importe posterior se toma de
esa captura local, no de una tarifa reenviada libremente por el
navegador.

\textbf{\texorpdfstring{\softtext{accept(sedeId,}}{accept(sedeId,} body,
key, auth, request)} exige un objeto con exactamente
\texorpdfstring{\softtext{snapshot\_id}}{snapshot\_id} textual y
comprobable como UUID. No acepta precio, hash o contexto empresarial
adicionales del cliente. Invoca
\texorpdfstring{\softtext{aceptar\_tarifa}}{aceptar\_tarifa} con clave y
correlación y devuelve 201. La aceptación y el gasto son comandos
locales distintos; el gasto canónico verifica que la captura corresponda
a actor y ámbito.

\subsection{\texorpdfstring{\texorpdfstring{\softtext{OutboxWorker:}}{OutboxWorker:}
transacciones separadas, sin atomicidad
distribuida}{ transacciones separadas, sin atomicidad distribuida}}\label{outboxworker-transacciones-separadas-sin-atomicidad-distribuida}

Referencia:
\texorpdfstring{\softcode{backend/src/main/java/mx/bovicierre/integraciones/OutboxWorker.java:16}}{backend/src/main/java/mx/bovicierre/integraciones/OutboxWorker.java:16}.
Sus campos jdbc y transactions tienen mapas A/B propios, construidos con
workerA y workerB; enabled decide si actúa. worker es un UUID nuevo para
esta instancia del proceso. unavailable es un HashSet que evita repetir
el mismo aviso en cada ciclo; LOG publica mensajes técnicos
\texorpdfstring{\softtext{genéricos.}}{genéricos.}

El constructor
\textbf{\texorpdfstring{\softtext{OutboxWorker(a,}}{OutboxWorker(a,} b,
enabled)} crea las dos plantillas y transacciones con roles de
\texorpdfstring{\softtext{trabajador.}}{trabajador.} El único método de
trabajo \textbf{dispatch()} está programado con espera inicial de 10000
ms y demora fija de 5000 ms entre
\texorpdfstring{\softtext{ejecuciones.}}{ejecuciones.} No es una
solicitud HTTP.

Primero, si enabled es false, sale. Para node A y B, llama
\texorpdfstring{\softtext{despachar\_outbox(50)}}{despachar\_outbox(50)}
en una transacción local. Un fallo de un nodo se registra una vez y
conserva trabajo pendiente; una recuperación puede registrarse cuando
unavailable pierde esa marca.

Después reclama en B hasta 20 entregas con
\texorpdfstring{\softtext{reclamar\_entregas(worker).}}{reclamar\_entregas(worker).}
claimed es JSON; cada event contiene id, empresa, sede y payload. La
transacción de reclamo termina antes de contactar A. Para cada evento,
una transacción A ejecuta
\texorpdfstring{\softtext{consumir\_muerte}}{consumir\_muerte} y una
transacción B posterior confirma entrega con id y worker.

Si consume o confirma falla, otra transacción B llama
\texorpdfstring{\softtext{fallar\_entrega}}{fallar\_entrega} y registra
IDENTIDAD como pendiente. Si tampoco puede completar ese manejo, el
catch exterior conserva la reserva temporal hasta vencer y avisa que el
propietario \texorpdfstring{\softtext{reintentará.}}{reintentará.}

\begin{Shaded}
\begin{Highlighting}[]
\NormalTok{B: reclamar entregas y confirmar la transacción local}
\NormalTok{por cada evento reclamado:}
\NormalTok{    A: consumir identidad de forma idempotente y confirmar}
\NormalTok{    B: confirmar acuse para este trabajador}
\NormalTok{si falla:}
\NormalTok{    B: registrar fallo/backoff, cuando sea posible}
\NormalTok{    conservar evento pendiente o esperar vencimiento de lease}
\end{Highlighting}
\end{Shaded}

La lease y el backoff persistentes se definen en SQL; Java utiliza sus
funciones, no un temporizador como prueba de posesión del evento. Si A
confirmó y el acuse B se perdió, el consumo se puede repetir. La
deduplicación de inbox evita repetir el efecto central. El diseño admite
reintentos y no debe anunciar entrega externa exactamente una vez.

\subsection{\texorpdfstring{NodeHealth y
\texorpdfstring{\softtext{DiagnosticsController}}{DiagnosticsController}}{NodeHealth y }}\label{nodehealth-y-diagnosticscontroller}

Referencia:
\texorpdfstring{\softcode{backend/src/main/java/mx/bovicierre/shared/NodeHealth.java:11}}{backend/src/main/java/mx/bovicierre/shared/NodeHealth.java:11}.
El constructor
\textbf{\texorpdfstring{\softtext{NodeHealth(a,}}{NodeHealth(a,} b,
auth, workerA, workerB)} crea nodes para probar A, B y
\texorpdfstring{\softtext{AUTORIZACION\_A,}}{AUTORIZACION\_A,} y pools
para observar esas tres fuentes más WORKER\_A y WORKER\_B.

\textbf{status()} ejecuta select 1 en cada nodes y construye statuses
con DISPONIBLE o
\texorpdfstring{\softtext{NO\_DISPONIBLE.}}{NO\_DISPONIBLE.} Devuelve
estado ACTIVO del proceso y los estados separados. El endpoint responde
200; estado ACTIVO no significa que las tres dependencias estén
\texorpdfstring{\softtext{disponibles.}}{disponibles.}

\textbf{pools()} examina cada fuente Hikari. bean es
\texorpdfstring{\softtext{HikariPoolMXBean:}}{HikariPoolMXBean:} result
recibe activos, inactivos, esperando y máximo. Si el pool todavía no
dispone del bean, informa
\texorpdfstring{\softtext{NO\_DISPONIBLE.}}{NO\_DISPONIBLE.} No publica
URL ni \texorpdfstring{\softtext{contraseña.}}{contraseña.}

Referencia:
\texorpdfstring{\softcode{backend/src/main/java/mx/bovicierre/integraciones/DiagnosticsController.java:15}}{backend/src/main/java/mx/bovicierre/integraciones/DiagnosticsController.java:15}.
Conserva access, commands y health por su
\texorpdfstring{\softtext{constructor.}}{constructor.}
\textbf{\texorpdfstring{\softtext{snapshot(sedeId,}}{snapshot(sedeId,}
auth)} exige ADMIN vigente. sync intenta consultar sincronizacion; si
hay
\texorpdfstring{\softtext{DataAccessException,}}{DataAccessException,}
conserva un estado público de indisponibilidad en esa sección. result
combina diagnostics, pools, sincronizacion y
\texorpdfstring{\softtext{capturado\_en.}}{capturado\_en.} Devuelve
no-store.

Un operador no obtiene estos detalles por conocer la ruta. La respuesta
tampoco es un certificado de retraso cero de cada flujo de réplica:
informa el corte que proporciona la consulta canónica y la
disponibilidad instantánea de
\texorpdfstring{\softtext{conexiones.}}{conexiones.}

\section{CSV, QR y PDF: producir representaciones sin reconocer otra
operación}\label{csv-qr-y-pdf-producir-representaciones-sin-reconocer-otra-operaciuxf3n}

\subsection{\texorpdfstring{\texorpdfstring{\softtext{CsvController:}}{CsvController:}
tres formas de
reporte}{ tres formas de reporte}}\label{csvcontroller-tres-formas-de-reporte}

Referencia:
\texorpdfstring{\softcode{backend/src/main/java/mx/bovicierre/reportes/CsvController.java:15}}{backend/src/main/java/mx/bovicierre/reportes/CsvController.java:15}.
Tiene access y commands, fijados por
\textbf{\texorpdfstring{\softtext{CsvController(access,}}{CsvController(access,}
commands)}.
\textbf{\texorpdfstring{\softtext{export(sedeId,}}{export(sedeId,} tipo,
auth)} admite cartera, margen e inventario; otro tipo produce 404. user
se obtiene de la sesión y rows acumula los nodos que serán
\texorpdfstring{\softtext{exportados.}}{exportados.}

Para page desde uno, consulta reportes\_ más tipo con tamaño 100.
response es el JSON SQL. En margen agrega el objeto completo como una
fila y termina: el margen canónico es un agregado, no una página que
Java deba sumar. Para cartera, items sale de filas; para inventario
utiliza elementos y conserva una compatibilidad con items. Agrega esas
filas y termina al alcanzar total o cuando la página está vacía. Si
todavía faltan filas y ya alcanzó 10000, rechaza el reporte grande.

La salida es text/csv UTF-8 con Content-Disposition y nombre por tipo.
La implementación acumula las filas en memoria; no es un streaming de
millones de registros ni toma filtros libres del endpoint JSON para la
descarga.

\textbf{cell(value)} examina trimmed, el texto sin espacios iniciales.
Si comienza por un signo de fórmula =, +, - o @, o por
\texorpdfstring{\softtext{tabulación/salto}}{tabulación/salto} de línea
de control al inicio original, antepone una comilla simple. Después
duplica comillas dobles y encierra toda celda entre comillas. Esta
protección evita que un nombre de contraparte se interprete
ordinariamente como fórmula al abrir la hoja.

\textbf{csv(rows)} devuelve cadena vacía si no hay filas. keys toma las
propiedades de la primera fila y define las columnas; csv es el
\texorpdfstring{\softtext{StringBuilder.}}{StringBuilder.} Por cada row
construye cells: escalares se vuelven texto y objetos o arrays se
serializan como JSON. Usa CRLF entre registros. Un Dinero permanece como
objeto JSON escapado en la celda; no se convierte en double ni se
recalcula.

\subsection{\texorpdfstring{\texorpdfstring{\softtext{AnimalQrController:}}{AnimalQrController:}
matriz local para una ficha
autorizada}{ matriz local para una ficha autorizada}}\label{animalqrcontroller-matriz-local-para-una-ficha-autorizada}

Referencia:
\texorpdfstring{\softcode{backend/src/main/java/mx/bovicierre/evidencia/AnimalQrController.java:15}}{backend/src/main/java/mx/bovicierre/evidencia/AnimalQrController.java:15}.
El constructor
\textbf{\texorpdfstring{\softtext{AnimalQrController(access,}}{AnimalQrController(access,}
commands)} conserva las
\texorpdfstring{\softtext{dependencias.}}{dependencias.}
\textbf{qr(sedeId, id, auth)} consulta animal con ámbito autorizado
antes de crear el gráfico. url contiene
\texorpdfstring{\softtext{/?sede=\{UUID\}\#animales/\{UUID\};}}{/?sede=\{UUID\}\#animales/\{UUID\};}
no incluye cookie ni token. Devuelve SVG con no-store.

\textbf{svg(url)} llama
\texorpdfstring{\softtext{QRCodeWriter,}}{QRCodeWriter,} con formato
QR\_CODE y dimensiones
\texorpdfstring{\softtext{automáticas.}}{automáticas.} matrix contiene
los módulos. Los índices y y x recorren la matriz; cada módulo negro
agrega a path un cuadrado SVG mediante M, h y v. El resultado tiene
viewBox, fondo blanco, un path negro y etiqueta accesible. Ninguna API
externa recibe los
\texorpdfstring{\softtext{identificadores.}}{identificadores.}

El QR contiene una referencia de navegación, no una
\texorpdfstring{\softtext{credencial.}}{credencial.} Al abrir la ficha
se aplican los controles de sesión. La generación local tampoco
demuestra haber probado una cámara de teléfono físico; eso corresponde a
evidencia de dispositivo aparte.

\subsection{\texorpdfstring{\texorpdfstring{\softtext{SettlementController}}{SettlementController}
y
\texorpdfstring{\softtext{SettlementPdf}}{SettlementPdf}}{ y }}\label{settlementcontroller-y-settlementpdf}

Referencia:
\texorpdfstring{\softcode{backend/src/main/java/mx/bovicierre/evidencia/SettlementController.java:13}}{backend/src/main/java/mx/bovicierre/evidencia/SettlementController.java:13}.
El constructor conserva access y commands.
\textbf{\texorpdfstring{\softtext{receipt(sedeId,}}{receipt(sedeId,} id,
auth)} consulta liquidacion, convierte su JSON a liquidation y llama
\texorpdfstring{\softtext{SettlementPdf.render.}}{SettlementPdf.render.}
La salida son bytes
\texorpdfstring{\softtext{application/pdf}}{application/pdf} y un nombre
\texorpdfstring{\softtext{liquidacion-\{id\}.pdf.}}{liquidacion-\{id\}.pdf.}
No hay llamada a
\texorpdfstring{\softtext{confirmar\_operacion}}{confirmar\_operacion}
ni a \texorpdfstring{\softtext{registrar\_pago.}}{registrar\_pago.}

Referencia del generador:
\texorpdfstring{\softcode{backend/src/main/java/mx/bovicierre/evidencia/SettlementPdf.java:9}}{backend/src/main/java/mx/bovicierre/evidencia/SettlementPdf.java:9}.
Su campo ENCODING es windows-1252 y su constructor privado evita
instanciar la utilidad. \textbf{money(value)} concatena importe y moneda
si recibió un objeto o toma su texto si recibió una forma
\texorpdfstring{\softtext{compatible.}}{compatible.}
\textbf{bytes(text)} aplica ENCODING. Ninguno calcula un saldo.

\textbf{render(liquidation)} arma lines: encabezado de comprobante
interno no fiscal, IDs, contraparte, dirección, total, cargos, abonos,
aplicaciones, saldo y
\texorpdfstring{\softtext{vencimiento.}}{vencimiento.} concept recorre
conceptos y agrega cantidad, unidad, precio e importe como textos
\texorpdfstring{\softtext{canónicos.}}{canónicos.} Se agrega una
aclaración de que el documento deriva de hechos
\texorpdfstring{\softtext{confirmados.}}{confirmados.}

objects reúne objetos PDF en bytes. Primero se crean catálogo, colección
de páginas y fuente Helvetica con
\texorpdfstring{\softtext{WinAnsiEncoding.}}{WinAnsiEncoding.} pages
resulta de dividir las líneas en grupos de 40; kids enumera sus
\texorpdfstring{\softtext{referencias.}}{referencias.} Cada page genera
un objeto de página y otro de contenido; contentId señala este último.
MediaBox usa 612 × 792 puntos, tamaño carta, no A4.

content escribe comandos de texto: tamaño 10, avance 14 y punto inicial.
Se escapan barra inversa y paréntesis y se eliminan controles que
podrían romper el texto. stream contiene los bytes del contenido; object
agrega su longitud y
\texorpdfstring{\softtext{delimitadores.}}{delimitadores.} pdf empieza
con PDF-1.4, offsets guarda la posición de cada objeto y xref la
posición de la tabla de
\texorpdfstring{\softtext{referencias.}}{referencias.} El trailer señala
el catálogo y startxref permite encontrar la tabla.

Esta estructura explica por qué el archivo puede abrirse sin una
biblioteca PDF adicional. No incrusta una fuente Unicode completa, no
genera firma digital y no es un documento fiscal. El PDF de este
comprobante del producto tampoco es el compilador del manual que el
lector está \texorpdfstring{\softtext{consultando.}}{consultando.}

\section{Tres secuencias completas para reconstruir la
comunicación}\label{tres-secuencias-completas-para-reconstruir-la-comunicaciuxf3n}

\subsection{Secuencia 1: iniciar sesión y obtener una sede
autorizada}\label{secuencia-1-iniciar-sesiuxf3n-y-obtener-una-sede-autorizada}

\begin{enumerate}
\def\labelenumi{\arabic{enumi}.}
\tightlist
\item
  El navegador solicita GET
  \texorpdfstring{\softtext{/api/v1/sesion/csrf.}}{/api/v1/sesion/csrf.}
  \texorpdfstring{\softtext{AvailabilityFilter}}{AvailabilityFilter}
  prueba A técnico en
  \texorpdfstring{\softcode{AvailabilityFilter.java:23}}{AvailabilityFilter.java:23}
  y la configuración permite acceso anónimo en
  \texorpdfstring{\softcode{SecurityConfiguration.java:27}}{SecurityConfiguration.java:27}.
\item
  \texorpdfstring{\softtext{SessionController.csrf,}}{SessionController.csrf,}
  \texorpdfstring{\softcode{SessionController.java:37}}{SessionController.java:37},
  devuelve token y \texorpdfstring{\softtext{headerName,}}{headerName,}
  con no-store. El navegador conserva la cookie de sesión y utiliza la
  cabecera indicada. Un token ilustrativo no sirve en otro contexto
  real.
\item
  POST
  \texorpdfstring{\softtext{/api/v1/sesion/login}}{/api/v1/sesion/login}
  envía login y contrasena en JSON y el token CSRF. La validación del
  record y el control de 72 bytes están en
  \texorpdfstring{\softcode{SessionController.java:38}}{SessionController.java:38}.
\item
  \texorpdfstring{\softtext{AccessService.login,}}{AccessService.login,}
  \texorpdfstring{\softcode{AccessService.java:17}}{AccessService.java:17},
  busca el usuario en la empresa
  \texorpdfstring{\softtext{configurada.}}{configurada.}
  \texorpdfstring{\softtext{SessionController}}{SessionController}
  compara BCrypt, exige ACTIVO y usa dummyHash si falta una identidad.
\item
  \texorpdfstring{\softcode{SessionController.java:46}}{SessionController.java:46}
  crea o utiliza la sesión y cambia su
  \texorpdfstring{\softtext{identificador.}}{identificador.} Guarda el
  contexto autenticado, invalida el CSRF anterior y devuelve 204.
\item
  El navegador vuelve a solicitar CSRF y consulta
  \texorpdfstring{\softtext{/sesion/actual.}}{/sesion/actual.}
  \texorpdfstring{\softcode{SessionController.java:59}}{SessionController.java:59}
  agrupa sedes, roles y
  \texorpdfstring{\softtext{capacidades.}}{capacidades.} El usuario
  puede elegir una sede de esa lista.
\item
  En solicitudes posteriores el filtro de
  \texorpdfstring{\softcode{SecurityConfiguration.java:38}}{SecurityConfiguration.java:38}
  vuelve a consultar current. Bloquear al usuario afecta una nueva
  petición aunque el navegador conserve su cookie.
\end{enumerate}

\begin{Shaded}
\begin{Highlighting}[]
\NormalTok{GET CSRF {-}\textgreater{} cookie de sesión + token inicial}
\NormalTok{POST login + token {-}\textgreater{} usuario ACTIVO y contraseña comprobados}
\NormalTok{rotación de sesión {-}\textgreater{} 204 y token antiguo invalidado}
\NormalTok{GET CSRF nuevo {-}\textgreater{} token de la sesión autenticada}
\NormalTok{GET sesión actual {-}\textgreater{} sedes y roles vigentes}
\end{Highlighting}
\end{Shaded}

Las referencias abreviadas de esta sección tienen el prefijo
\texorpdfstring{\softcode{backend/src/main/java/mx/bovicierre/acceso/}}{backend/src/main/java/mx/bovicierre/acceso/}.
El flujo no usa JWT, HTTP Basic ni una contraseña fija introducida por
el ejemplo.

\subsection{Secuencia 2: confirmar una operación y recuperar una
respuesta
perdida}\label{secuencia-2-confirmar-una-operaciuxf3n-y-recuperar-una-respuesta-perdida}

\begin{enumerate}
\def\labelenumi{\arabic{enumi}.}
\tightlist
\item
  GET
  \texorpdfstring{\softtext{/api/v1/sedes/\{sedeId\}/operaciones/\{id\}}}{/api/v1/sedes/\{sedeId\}/operaciones/\{id\}}
  usa
  \texorpdfstring{\softtext{CoreController.one,}}{CoreController.one,}
  \texorpdfstring{\softcode{operaciones/CoreController.java:96}}{operaciones/CoreController.java:96}.
  El JSON que contiene
  \texorpdfstring{\softtext{version\_lock}}{version\_lock} produce ETag
  por result, línea24.
\item
  El operador prepara un intento con la versión original,
  Idempotency-Key nueva y cuerpo de
  \texorpdfstring{\softtext{reconocimiento.}}{reconocimiento.} El precio
  y lotes pertenecen al borrador; el navegador no envía un total
  \texorpdfstring{\softtext{autoritativo.}}{autoritativo.}
\end{enumerate}

\begin{Shaded}
\begin{Highlighting}[]
\FunctionTok{\{}
  \DataTypeTok{"periodo\_id"}\FunctionTok{:} \StringTok{"00000000{-}0000{-}4000{-}8000{-}000000000002"}\FunctionTok{,}
  \DataTypeTok{"fecha\_reconocimiento"}\FunctionTok{:} \StringTok{"2026{-}10{-}12"}
\FunctionTok{\}}
\end{Highlighting}
\end{Shaded}

\begin{enumerate}
\def\labelenumi{\arabic{enumi}.}
\setcounter{enumi}{2}
\tightlist
\item
  POST
  \texorpdfstring{\softtext{/operaciones/\{id\}/confirmar}}{/operaciones/\{id\}/confirmar}
  lleva If-Match, por ejemplo ``3'', y
  \texorpdfstring{\softtext{X-XSRF-TOKEN.}}{X-XSRF-TOKEN.} action,
  \texorpdfstring{\softcode{CoreController.java:113}}{CoreController.java:113},
  elige
  \texorpdfstring{\softtext{confirmar\_operacion;}}{confirmar\_operacion;}
  write, línea30, fija
  \texorpdfstring{\softtext{operacion\_id}}{operacion\_id} y
  \texorpdfstring{\softtext{version\_lock}}{version\_lock} del servidor.
\item
  \texorpdfstring{\softcode{CommandService.java:32}}{CommandService.java:32}
  autoriza nuevamente, escoge el dueño y abre su
  \texorpdfstring{\softtext{TransactionTemplate.}}{TransactionTemplate.}
  No abre otra transacción en el nodo opuesto.
\item
  PostgreSQL V13 reclama la intención bajo el ámbito y principal. La
  fila idempotente y el retorno previo aparecen en
  \texorpdfstring{\softcode{db/migrations/V13\_\_correccion\_composicion\_exacta.sql:12}}{db/migrations/V13\_\_correccion\_composicion\_exacta.sql:12}.
  Si la intención está completada, devuelve el resultado original antes
  de la nueva comprobación de versión.
\item
  Para una ejecución nueva, el SQL comprueba periodo de reconocimiento y
  versión bajo bloqueo, condiciones, pesajes y recursos. La entrada a
  confirmación se observa en V13:129. Los cálculos y escrituras son
  responsabilidad del capítulo de base de datos.
\item
  El SQL produce hechos económicos, auditoría, outbox y resultado del
  intento. La transacción local confirma; Java devuelve ese JSON. Una
  excepción revierte las escrituras de esa
  \texorpdfstring{\softtext{ejecución.}}{ejecución.}
\item
  Si el navegador perdió la respuesta, conserva clave, cuerpo y ETag
  \texorpdfstring{\softtext{originales.}}{originales.} Consulta GET
  \texorpdfstring{\softtext{/comandos/confirmar\_operacion/\{clave\},}}{/comandos/confirmar\_operacion/\{clave\},}
  \texorpdfstring{\softcode{CoreController.java:64}}{CoreController.java:64}.
  Un resultado completado permite recuperar el objeto sin otro POST.
\item
  Si el resultado aún no es visible, una ausencia no demuestra rollback:
  la transacción puede seguir en curso o la consulta puede fallar.
  Cuando corresponda reintentar, se conserva exactamente la
  \texorpdfstring{\softtext{intención.}}{intención.} Cambiar fecha,
  versión o cuerpo con la misma clave provoca conflicto.
\end{enumerate}

\[
\text{Misma clave y misma intención} \Rightarrow \text{resultado persistido anterior}
\]

\[
\text{Misma clave y distinta intención} \Rightarrow 409
\]

Este retorno no implica volver a crear liquidación, deuda o asientos.
Los bloqueos y la clave persistida coordinan intentos simultáneos; la
cabecera del navegador, por sí sola, no garantiza
\texorpdfstring{\softtext{idempotencia.}}{idempotencia.}

\subsection{Secuencia 3: publicar una versión documental y
descargarla}\label{secuencia-3-publicar-una-versiuxf3n-documental-y-descargarla}

\begin{enumerate}
\def\labelenumi{\arabic{enumi}.}
\tightlist
\item
  El cliente envía multipart a
  \texorpdfstring{\softtext{/api/v1/sedes/\{sedeId\}/documentos}}{/api/v1/sedes/\{sedeId\}/documentos}
  con archivo y
  \texorpdfstring{\softtext{tipo\_documental;}}{tipo\_documental;}
  opcionalmente operación y documento previo. Conserva su clave de
  intento y token CSRF.
\item
  \texorpdfstring{\softtext{DocumentController.upload,}}{DocumentController.upload,}
  \texorpdfstring{\softcode{evidencia/DocumentController.java:22}}{evidencia/DocumentController.java:22},
  obtiene usuario y Scope y valida el tipo
  \texorpdfstring{\softtext{documental.}}{documental.} No utiliza una
  ruta de disco enviada por el cliente.
\item
  \texorpdfstring{\softtext{PrivateDocuments.stage,}}{PrivateDocuments.stage,}
  \texorpdfstring{\softcode{evidencia/PrivateDocuments.java:64}}{evidencia/PrivateDocuments.java:64},
  valida tamaño y firma, calcula SHA-256, adquiere bloqueo local y de
  archivo, escribe .part y publica el contenido mediante movimiento
  atómico.
\item
  Upload devuelve Stored y mantiene el bloqueo.
  \texorpdfstring{\softtext{DocumentController}}{DocumentController}
  construye metadata y
  \texorpdfstring{\softtext{storage\_key.}}{storage\_key.} Si es una
  nueva versión, consulta el previo y obtiene su número para
  \texorpdfstring{\softtext{incrementarlo.}}{incrementarlo.}
\item
  \texorpdfstring{\softtext{commands.command}}{commands.command} ejecuta
  \texorpdfstring{\softtext{registrar\_documento}}{registrar\_documento}
  con su intención y transacción local. Un commit deja metadata y
  resultado idempotente; el cierre del bloque libera el bloqueo. Un
  rechazo puede dejar contenido físico huérfano, sin inventar metadata
  aprobada.
\item
  Un retry igual puede recuperar el metadata anterior. La limpieza no
  intenta resolver esa incertidumbre borrando inmediatamente el archivo.
\item
  GET
  \texorpdfstring{\softtext{/documentos/\{id\}/contenido}}{/documentos/\{id\}/contenido}
  consulta primero metadata
  \texorpdfstring{\softtext{autorizado.}}{autorizado.}
  \texorpdfstring{\softtext{PrivateDocuments.path,}}{PrivateDocuments.path,}
  línea116, verifica ruta, archivo regular y hash.
  \texorpdfstring{\softtext{DocumentController.download,}}{DocumentController.download,}
  línea42, publica el recurso como adjunto, sin interpretación de tipo
  en el navegador.
\item
  Una revisión administrativa referencia la versión documental
  correspondiente; crear otra versión no convierte retroactivamente la
  primera revisión en revisión de todos los contenidos futuros.
\item
  Pasadas 24 horas, el procedimiento de limpieza puede diagnosticar
  objetos sin \texorpdfstring{\softtext{referencia.}}{referencia.}
  \texorpdfstring{\softtext{CleanupMain}}{CleanupMain} consulta A y B
  bajo rol técnico. Solo el modo explícito de ejecución elimina
  candidatos comprobados y no
  \texorpdfstring{\softtext{bloqueados.}}{bloqueados.}
\end{enumerate}

Aquí no hay una transacción distribuida que englobe simultáneamente
PostgreSQL y el disco. La protección procede del protocolo de staging,
huella, bloqueo durante metadata, idempotencia y limpieza
\texorpdfstring{\softtext{conservadora.}}{conservadora.}

\section{Inventario del contrato HTTP y sus llamadores
Java}\label{inventario-del-contrato-http-y-sus-llamadores-java}

Todas las rutas siguientes llevan el prefijo /api/v1. Para compactar la
lectura, \textbf{S} representa
\texorpdfstring{\softtext{/sedes/\{sedeId\},}}{/sedes/\{sedeId\},}
\textbf{O} representa
\texorpdfstring{\softtext{operaciones/\{operacionId\},}}{operaciones/\{operacionId\},}
\textbf{P}
\texorpdfstring{\softtext{pagos/\{pagoId\},}}{pagos/\{pagoId\},}
\textbf{D}
\texorpdfstring{\softtext{documentos/\{documentoId\},}}{documentos/\{documentoId\},}
\textbf{L}
\texorpdfstring{\softtext{liquidaciones/\{liquidacionId\},}}{liquidaciones/\{liquidacionId\},}
\textbf{G}
\texorpdfstring{\softtext{gastos/\{gastoId\}}}{gastos/\{gastoId\}} y
\textbf{T}
\texorpdfstring{\softtext{periodos/\{periodoId\}.}}{periodos/\{periodoId\}.}
La tabla enumera los 70
\texorpdfstring{\softtext{operationId,}}{operationId,} no 70 clases.

{\def\LTcaptype{none} % do not increment counter
\begin{longtable}[]{@{}
  >{\raggedright\arraybackslash}p{(\linewidth - 4\tabcolsep) * \real{0.3333}}
  >{\raggedright\arraybackslash}p{(\linewidth - 4\tabcolsep) * \real{0.3333}}
  >{\raggedright\arraybackslash}p{(\linewidth - 4\tabcolsep) * \real{0.3333}}@{}}
\toprule\noalign{}
\begin{minipage}[b]{\linewidth}\raggedright
Método y ruta relativa
\end{minipage} & \begin{minipage}[b]{\linewidth}\raggedright
\texorpdfstring{\softtext{operationId}}{operationId} del contrato
\end{minipage} & \begin{minipage}[b]{\linewidth}\raggedright
Método Java principal
\end{minipage} \\
\midrule\noalign{}
\endhead
\bottomrule\noalign{}
\endlastfoot
POST \texorpdfstring{\softtext{/sesion/login}}{/sesion/login} & login &
\texorpdfstring{\softtext{SessionController.login}}{SessionController.login} \\
POST \texorpdfstring{\softtext{/sesion/logout}}{/sesion/logout} & logout
&
\texorpdfstring{\softtext{SessionController.logout}}{SessionController.logout} \\
GET \texorpdfstring{\softtext{/sesion/actual}}{/sesion/actual} &
\texorpdfstring{\softtext{sesionActual}}{sesionActual} &
\texorpdfstring{\softtext{SessionController.current}}{SessionController.current} \\
GET \texorpdfstring{\softtext{/sesion/csrf}}{/sesion/csrf} &
\texorpdfstring{\softtext{obtenerCsrf}}{obtenerCsrf} &
\texorpdfstring{\softtext{SessionController.csrf}}{SessionController.csrf} \\
GET /salud &
\texorpdfstring{\softtext{consultarSaludProceso}}{consultarSaludProceso}
&
\texorpdfstring{\softtext{SessionController.health}}{SessionController.health} \\
GET \texorpdfstring{\softtext{S/contrapartes}}{S/contrapartes} &
\texorpdfstring{\softtext{listarContrapartes}}{listarContrapartes} &
\texorpdfstring{\softtext{CoreController.list}}{CoreController.list} \\
POST \texorpdfstring{\softtext{S/contrapartes}}{S/contrapartes} &
\texorpdfstring{\softtext{crearContraparte}}{crearContraparte} &
\texorpdfstring{\softtext{CoreController.create}}{CoreController.create} \\
POST S/animales &
\texorpdfstring{\softtext{registrarIdentidadAnimal}}{registrarIdentidadAnimal}
&
\texorpdfstring{\softtext{CoreController.create}}{CoreController.create} \\
GET S/animales &
\texorpdfstring{\softtext{listarAnimales}}{listarAnimales} &
\texorpdfstring{\softtext{CoreController.list}}{CoreController.list} \\
GET
\texorpdfstring{\softtext{S/animales/\{animalId\}}}{S/animales/\{animalId\}}
& \texorpdfstring{\softtext{fichaAnimal}}{fichaAnimal} &
\texorpdfstring{\softtext{CoreController.one}}{CoreController.one} \\
GET
\texorpdfstring{\softtext{S/animales/\{animalId\}/qr}}{S/animales/\{animalId\}/qr}
& qrAnimal &
\texorpdfstring{\softtext{AnimalQrController.qr}}{AnimalQrController.qr} \\
GET \texorpdfstring{\softtext{S/operaciones}}{S/operaciones} &
\texorpdfstring{\softtext{bandejaOperaciones}}{bandejaOperaciones} &
\texorpdfstring{\softtext{CoreController.list}}{CoreController.list} \\
POST \texorpdfstring{\softtext{S/operaciones}}{S/operaciones} &
\texorpdfstring{\softtext{crearBorrador}}{crearBorrador} &
\texorpdfstring{\softtext{CoreController.create}}{CoreController.create} \\
GET S/O & \texorpdfstring{\softtext{obtenerOperacion}}{obtenerOperacion}
& \texorpdfstring{\softtext{CoreController.one}}{CoreController.one} \\
PATCH S/O & \texorpdfstring{\softtext{editarBorrador}}{editarBorrador} &
\texorpdfstring{\softtext{CoreController.edit}}{CoreController.edit} \\
POST \texorpdfstring{\softtext{S/O/confirmar}}{S/O/confirmar} &
\texorpdfstring{\softtext{confirmarOperacion}}{confirmarOperacion} &
\texorpdfstring{\softtext{CoreController.action}}{CoreController.action} \\
POST \texorpdfstring{\softtext{S/O/pesajes}}{S/O/pesajes} &
\texorpdfstring{\softtext{proponerPesaje}}{proponerPesaje} &
\texorpdfstring{\softtext{CoreController.action}}{CoreController.action} \\
POST
\texorpdfstring{\softtext{S/pesajes/\{pesajeId\}/aprobar}}{S/pesajes/\{pesajeId\}/aprobar}
& \texorpdfstring{\softtext{aprobarPesaje}}{aprobarPesaje} &
\texorpdfstring{\softtext{CoreController.action}}{CoreController.action} \\
POST
\texorpdfstring{\softtext{S/pesajes/\{pesajeId\}/correcciones}}{S/pesajes/\{pesajeId\}/correcciones}
& \texorpdfstring{\softtext{corregirPesaje}}{corregirPesaje} &
\texorpdfstring{\softtext{CoreController.action}}{CoreController.action} \\
POST \texorpdfstring{\softtext{S/O/reservas}}{S/O/reservas} &
\texorpdfstring{\softtext{reservarPosiciones}}{reservarPosiciones} &
\texorpdfstring{\softtext{CoreController.action}}{CoreController.action} \\
POST
\texorpdfstring{\softtext{S/reservas/\{reservaId\}/liberar}}{S/reservas/\{reservaId\}/liberar}
& \texorpdfstring{\softtext{liberarReserva}}{liberarReserva} &
\texorpdfstring{\softtext{CoreController.action}}{CoreController.action} \\
GET
\texorpdfstring{\softtext{S/inventario/posiciones}}{S/inventario/posiciones}
& \texorpdfstring{\softtext{listarPosiciones}}{listarPosiciones} &
\texorpdfstring{\softtext{CoreController.list}}{CoreController.list} \\
POST
\texorpdfstring{\softtext{S/inventario/movimientos}}{S/inventario/movimientos}
& \texorpdfstring{\softtext{registrarMovimiento}}{registrarMovimiento} &
\texorpdfstring{\softtext{CoreController.create}}{CoreController.create} \\
GET S/L &
\texorpdfstring{\softtext{detalleLiquidacion}}{detalleLiquidacion} &
\texorpdfstring{\softtext{CoreController.one}}{CoreController.one} \\
POST \texorpdfstring{\softtext{S/L/ajustes}}{S/L/ajustes} &
\texorpdfstring{\softtext{ajustarLiquidacion}}{ajustarLiquidacion} &
\texorpdfstring{\softtext{CoreController.action}}{CoreController.action} \\
GET \texorpdfstring{\softtext{S/L/comprobante}}{S/L/comprobante} &
\texorpdfstring{\softtext{comprobanteInterno}}{comprobanteInterno} &
\texorpdfstring{\softtext{SettlementController.receipt}}{SettlementController.receipt} \\
POST S/gastos &
\texorpdfstring{\softtext{registrarGasto}}{registrarGasto} &
\texorpdfstring{\softtext{CoreController.create}}{CoreController.create} \\
GET S/gastos & \texorpdfstring{\softtext{listarGastos}}{listarGastos} &
\texorpdfstring{\softtext{CoreController.list}}{CoreController.list} \\
GET \texorpdfstring{\softtext{S/G/reparto}}{S/G/reparto} &
\texorpdfstring{\softtext{previsualizarReparto}}{previsualizarReparto} &
\texorpdfstring{\softtext{CoreController.one}}{CoreController.one} \\
POST \texorpdfstring{\softtext{S/G/reparto}}{S/G/reparto} &
\texorpdfstring{\softtext{aprobarReparto}}{aprobarReparto} &
\texorpdfstring{\softtext{CoreController.action}}{CoreController.action} \\
POST S/pagos & \texorpdfstring{\softtext{registrarPago}}{registrarPago}
&
\texorpdfstring{\softtext{CoreController.create}}{CoreController.create} \\
GET S/pagos & \texorpdfstring{\softtext{listarPagos}}{listarPagos} &
\texorpdfstring{\softtext{CoreController.list}}{CoreController.list} \\
POST \texorpdfstring{\softtext{S/P/aplicaciones}}{S/P/aplicaciones} &
\texorpdfstring{\softtext{aplicarPago}}{aplicarPago} &
\texorpdfstring{\softtext{CoreController.action}}{CoreController.action} \\
POST
\texorpdfstring{\softtext{S/aplicaciones/\{aplicacionId\}/reverso}}{S/aplicaciones/\{aplicacionId\}/reverso}
& \texorpdfstring{\softtext{revertirAplicacion}}{revertirAplicacion} &
\texorpdfstring{\softtext{CoreController.action}}{CoreController.action} \\
POST \texorpdfstring{\softtext{S/P/reembolsos}}{S/P/reembolsos} &
\texorpdfstring{\softtext{registrarReembolso}}{registrarReembolso} &
\texorpdfstring{\softtext{CoreController.action}}{CoreController.action} \\
GET \texorpdfstring{\softtext{S/reportes/cartera}}{S/reportes/cartera} &
\texorpdfstring{\softtext{reporteCartera}}{reporteCartera} &
\texorpdfstring{\softtext{CoreController.list}}{CoreController.list} \\
GET \texorpdfstring{\softtext{S/reportes/margen}}{S/reportes/margen} &
\texorpdfstring{\softtext{reporteMargen}}{reporteMargen} &
\texorpdfstring{\softtext{CoreController.list}}{CoreController.list} \\
POST S/T/cierre &
\texorpdfstring{\softtext{cerrarPeriodo}}{cerrarPeriodo} &
\texorpdfstring{\softtext{CoreController.action}}{CoreController.action} \\
POST \texorpdfstring{\softtext{S/documentos}}{S/documentos} &
\texorpdfstring{\softtext{subirDocumento}}{subirDocumento} &
\texorpdfstring{\softtext{DocumentController.upload}}{DocumentController.upload} \\
GET \texorpdfstring{\softtext{S/documentos}}{S/documentos} &
\texorpdfstring{\softtext{listarDocumentos}}{listarDocumentos} &
\texorpdfstring{\softtext{CoreController.list}}{CoreController.list} \\
GET \texorpdfstring{\softtext{S/D/contenido}}{S/D/contenido} &
\texorpdfstring{\softtext{descargarDocumento}}{descargarDocumento} &
\texorpdfstring{\softtext{DocumentController.download}}{DocumentController.download} \\
GET
\texorpdfstring{\softtext{S/integraciones/proveedor/tarifas}}{S/integraciones/proveedor/tarifas}
&
\texorpdfstring{\softtext{consultarTarifasProveedor}}{consultarTarifasProveedor}
&
\texorpdfstring{\softtext{ProviderController.tariffs}}{ProviderController.tariffs} \\
GET \texorpdfstring{\softtext{S/liquidaciones}}{S/liquidaciones} &
\texorpdfstring{\softtext{listarLiquidaciones}}{listarLiquidaciones} &
\texorpdfstring{\softtext{CoreController.list}}{CoreController.list} \\
POST \texorpdfstring{\softtext{S/P/validar}}{S/P/validar} &
\texorpdfstring{\softtext{validarPago}}{validarPago} &
\texorpdfstring{\softtext{CoreController.action}}{CoreController.action} \\
GET
\texorpdfstring{\softtext{S/comandos/\{comando\}/\{clave\}}}{S/comandos/\{comando\}/\{clave\}}
&
\texorpdfstring{\softtext{consultarResultadoComando}}{consultarResultadoComando}
&
\texorpdfstring{\softtext{CoreController.commandResult}}{CoreController.commandResult} \\
GET S/usuarios &
\texorpdfstring{\softtext{listarUsuarios}}{listarUsuarios} &
\texorpdfstring{\softtext{CoreController.list}}{CoreController.list} \\
POST S/usuarios & \texorpdfstring{\softtext{crearUsuario}}{crearUsuario}
&
\texorpdfstring{\softtext{UserController.create}}{UserController.create} \\
GET S/periodos &
\texorpdfstring{\softtext{listarPeriodos}}{listarPeriodos} &
\texorpdfstring{\softtext{CoreController.list}}{CoreController.list} \\
POST S/periodos & \texorpdfstring{\softtext{crearPeriodo}}{crearPeriodo}
&
\texorpdfstring{\softtext{CoreController.create}}{CoreController.create} \\
GET S/predios & \texorpdfstring{\softtext{listarPredios}}{listarPredios}
&
\texorpdfstring{\softtext{CoreController.list}}{CoreController.list} \\
POST S/predios & \texorpdfstring{\softtext{crearPredio}}{crearPredio} &
\texorpdfstring{\softtext{CoreController.create}}{CoreController.create} \\
PATCH
\texorpdfstring{\softtext{S/usuarios/\{usuarioId\}}}{S/usuarios/\{usuarioId\}}
& \texorpdfstring{\softtext{editarUsuario}}{editarUsuario} &
\texorpdfstring{\softtext{UserController.update}}{UserController.update} \\
POST
\texorpdfstring{\softtext{S/usuarios/\{usuarioId\}/membresias}}{S/usuarios/\{usuarioId\}/membresias}
& \texorpdfstring{\softtext{asignarMembresia}}{asignarMembresia} &
\texorpdfstring{\softtext{UserController.membership}}{UserController.membership} \\
POST \texorpdfstring{\softtext{S/T/reabrir}}{S/T/reabrir} &
\texorpdfstring{\softtext{reabrirPeriodo}}{reabrirPeriodo} &
\texorpdfstring{\softtext{CoreController.action}}{CoreController.action} \\
GET \texorpdfstring{\softtext{S/T/versiones}}{S/T/versiones} &
\texorpdfstring{\softtext{consultarVersionesCierre}}{consultarVersionesCierre}
& \texorpdfstring{\softtext{CoreController.one}}{CoreController.one} \\
POST \texorpdfstring{\softtext{S/O/acordar}}{S/O/acordar} &
\texorpdfstring{\softtext{acordarOperacion}}{acordarOperacion} &
\texorpdfstring{\softtext{CoreController.action}}{CoreController.action} \\
POST \texorpdfstring{\softtext{S/O/reabrir}}{S/O/reabrir} &
\texorpdfstring{\softtext{reabrirOperacion}}{reabrirOperacion} &
\texorpdfstring{\softtext{CoreController.action}}{CoreController.action} \\
POST \texorpdfstring{\softtext{S/O/cancelar}}{S/O/cancelar} &
\texorpdfstring{\softtext{cancelarOperacion}}{cancelarOperacion} &
\texorpdfstring{\softtext{CoreController.action}}{CoreController.action} \\
POST S/O/anular &
\texorpdfstring{\softtext{anularOperacion}}{anularOperacion} &
\texorpdfstring{\softtext{CoreController.action}}{CoreController.action} \\
GET \texorpdfstring{\softtext{S/O/simulador}}{S/O/simulador} &
\texorpdfstring{\softtext{simularOperacion}}{simularOperacion} &
\texorpdfstring{\softtext{CoreController.one}}{CoreController.one} \\
GET \texorpdfstring{\softtext{S/O/expediente}}{S/O/expediente} &
\texorpdfstring{\softtext{consultarExpediente}}{consultarExpediente} &
\texorpdfstring{\softtext{CoreController.one}}{CoreController.one} \\
GET S/P & \texorpdfstring{\softtext{consultarPago}}{consultarPago} &
\texorpdfstring{\softtext{CoreController.one}}{CoreController.one} \\
GET S/G & \texorpdfstring{\softtext{consultarGasto}}{consultarGasto} &
\texorpdfstring{\softtext{CoreController.one}}{CoreController.one} \\
GET S/D &
\texorpdfstring{\softtext{consultarDocumento}}{consultarDocumento} &
\texorpdfstring{\softtext{DocumentController.metadata}}{DocumentController.metadata} \\
POST \texorpdfstring{\softtext{S/D/revisiones}}{S/D/revisiones} &
\texorpdfstring{\softtext{revisarDocumento}}{revisarDocumento} &
\texorpdfstring{\softtext{DocumentController.review}}{DocumentController.review} \\
GET
\texorpdfstring{\softtext{S/pesajes/comparacion}}{S/pesajes/comparacion}
& \texorpdfstring{\softtext{compararPesajes}}{compararPesajes} &
\texorpdfstring{\softtext{CoreController.compare}}{CoreController.compare} \\
GET
\texorpdfstring{\softtext{S/reportes/\{tipo\}/csv}}{S/reportes/\{tipo\}/csv}
& \texorpdfstring{\softtext{exportarCsv}}{exportarCsv} &
\texorpdfstring{\softtext{CsvController.export}}{CsvController.export} \\
GET \texorpdfstring{\softtext{S/sincronizacion}}{S/sincronizacion} &
\texorpdfstring{\softtext{consultarSincronizacion}}{consultarSincronizacion}
&
\texorpdfstring{\softtext{CoreController.list}}{CoreController.list} \\
POST
\texorpdfstring{\softtext{S/integraciones/proveedor/tarifas/aceptaciones}}{S/integraciones/proveedor/tarifas/aceptaciones}
& \texorpdfstring{\softtext{aceptarTarifa}}{aceptarTarifa} &
\texorpdfstring{\softtext{ProviderController.accept}}{ProviderController.accept} \\
GET
\texorpdfstring{\softtext{S/integraciones/diagnostico}}{S/integraciones/diagnostico}
& \texorpdfstring{\softtext{consultarDiagnostico}}{consultarDiagnostico}
&
\texorpdfstring{\softtext{DiagnosticsController.snapshot}}{DiagnosticsController.snapshot} \\
\end{longtable}
}

El catálogo muestra reutilización de métodos. Para conocer un contrato
completo hay que consultar su esquema, cabeceras y errores OpenAPI;
compartir
\texorpdfstring{\softtext{CoreController.action}}{CoreController.action}
no hace iguales un reverso, una confirmación y una reserva.

\section{Cobertura, evidencia y límites de este
capítulo}\label{cobertura-evidencia-y-luxedmites-de-este-capuxedtulo}

Se explicaron las 24 unidades de fuente de producción, sus constructores
explícitos, todos sus métodos declarados y los campos
\texorpdfstring{\softtext{relevantes.}}{relevantes.} También se
describieron los records anidados y sus accesores generados, el filtro
anónimo de seguridad, argumentos HTTP y variables que deciden ámbito,
transacción, hash, metadatos, representación y
\texorpdfstring{\softtext{reintentos.}}{reintentos.} Los imports y
métodos heredados del framework no se presentan como lógica nueva de
\texorpdfstring{\softtext{BoviCierre.}}{BoviCierre.}

La evidencia
\texorpdfstring{\softcode{backend/verification-result-v13.json}}{backend/verification-result-v13.json}
registra 24 pruebas aprobadas: 15 unitarias, ocho HTTP por Caddy15180 y
un PostgreSQL 18.6 fresco con V1-V13. Incluye contenido de CSV,
composición de pesajes, motivos, idempotencia, sesiones, archivos y
controles de \texorpdfstring{\softtext{autorización.}}{autorización.}
Este conjunto no prueba exhaustivamente las 70 operaciones y todas las
combinaciones de entrada.

Hay límites concretos que deben conservarse al explicar el código. Los
permisos detallados se verifican en servicios y SQL, no en la capacidad
genérica ESCRIBIR. La sesión y autorización dependen de A. No hay
failover automático de propietario, 2PC ni transacción conjunta entre
disco y base. Los uploads y exports usan memoria y tamaño acotado; LOCKS
no tiene una política explícita de expulsión en esta clase. El
almacenamiento requiere permisos POSIX y movimiento atómico del sistema
de archivos utilizado.

La firma inicial de un archivo no es
\texorpdfstring{\softtext{antimalware.}}{antimalware.} El PDF no
incrusta Unicode completo ni firma fiscalmente un documento. El QR no
acredita pruebas de cámara móvil. La aplicación inspeccionada no
implementa en estas clases un mecanismo de rate limiting de login;
BCrypt y mensajes uniformes no sustituyen ese control. Las métricas no
son historial persistente, p95 de HTTP ni certificación de capacidad.
Estos límites identifican el alcance ejecutable; no autorizan nuevas
modificaciones del producto por escribir el manual.

Finalmente, una clase compacta puede cumplir una responsabilidad sin
estar separada en diez archivos. Eso explica por qué el router real es
una decisión en
\texorpdfstring{\softtext{CommandService}}{CommandService} y por qué hay
un puerto documental concreto. Para evaluar mantenibilidad futura
conviene comparar responsabilidades observadas con la arquitectura
planeada, conservando primero las reglas SQL, el contrato y la
evidencia.

\section{Fuentes primarias y localización en el
repositorio}\label{fuentes-primarias-y-localizaciuxf3n-en-el-repositorio}

Prefijo \textbf{J}:
\texorpdfstring{\softcode{backend/src/main/java/mx/bovicierre/}}{backend/src/main/java/mx/bovicierre/}.
Las referencias J de esta lista deben expandirse con ese prefijo. Cada
archivo Java fue leído completo, no únicamente la línea citada.

{\def\LTcaptype{none} % do not increment counter
\begin{longtable}[]{@{}
  >{\raggedright\arraybackslash}p{(\linewidth - 4\tabcolsep) * \real{0.3333}}
  >{\raggedright\arraybackslash}p{(\linewidth - 4\tabcolsep) * \real{0.3333}}
  >{\raggedright\arraybackslash}p{(\linewidth - 4\tabcolsep) * \real{0.3333}}@{}}
\toprule\noalign{}
\begin{minipage}[b]{\linewidth}\raggedright
Fuente
\end{minipage} & \begin{minipage}[b]{\linewidth}\raggedright
Punto de entrada para estudiar
\end{minipage} & \begin{minipage}[b]{\linewidth}\raggedright
Contenido utilizado
\end{minipage} \\
\midrule\noalign{}
\endhead
\bottomrule\noalign{}
\endlastfoot
\texorpdfstring{\softcode{J/BoviCierreApplication.java:8}}{J/BoviCierreApplication.java:8}
& main & Inicio, scheduling y PID \\
\texorpdfstring{\softcode{J/shared/DatabaseConfiguration.java:15}}{J/shared/DatabaseConfiguration.java:15}
& pool y beans & Cinco conexiones y transacción de sesión \\
\texorpdfstring{\softcode{J/shared/CommandService.java:16}}{J/shared/CommandService.java:16}
& command y query & Router compacto, transacciones, reintentos y
diagnóstico \\
\texorpdfstring{\softcode{J/shared/JsonSupport.java:11}}{J/shared/JsonSupport.java:11}
& utilidades & JSON, decimales, hash, ETag e identificaciones \\
\texorpdfstring{\softcode{J/shared/ApiException.java:3}}{J/shared/ApiException.java:3}
& constructor & Excepción pública \\
\texorpdfstring{\softcode{J/shared/ApiErrors.java:15}}{J/shared/ApiErrors.java:15}
& handlers & Traducción de SQLSTATE y UTF-8 \\
\texorpdfstring{\softcode{J/shared/NodeHealth.java:11}}{J/shared/NodeHealth.java:11}
& status y pools & Disponibilidad y Hikari \\
\texorpdfstring{\softcode{J/acceso/AccessService.java:14}}{J/acceso/AccessService.java:14}
& authorize & Usuario, empresa, sedes y roles \\
\texorpdfstring{\softcode{J/acceso/AvailabilityFilter.java:19}}{J/acceso/AvailabilityFilter.java:19}
& \texorpdfstring{\softtext{doFilterInternal}}{doFilterInternal} &
Disponibilidad central \\
\texorpdfstring{\softcode{J/acceso/SecurityConfiguration.java:21}}{J/acceso/SecurityConfiguration.java:21}
& security & Contexto, CSRF y filtros \\
\texorpdfstring{\softcode{J/acceso/SessionController.java:26}}{J/acceso/SessionController.java:26}
& \texorpdfstring{\softtext{login/current}}{login/current} & Sesión y
salud \\
\texorpdfstring{\softcode{J/acceso/UserController.java:14}}{J/acceso/UserController.java:14}
& write & Usuarios, BCrypt y membresías \\
\texorpdfstring{\softcode{J/operaciones/CoreController.java:16}}{J/operaciones/CoreController.java:16}
& \texorpdfstring{\softtext{write/action}}{write/action} & Adaptador de
recursos económicos \\
\texorpdfstring{\softcode{J/evidencia/application/DocumentStorage.java:7}}{J/evidencia/application/DocumentStorage.java:7}
& interfaces & Puerto y ciclos de contenido \\
\texorpdfstring{\softcode{J/evidencia/PrivateDocuments.java:27}}{J/evidencia/PrivateDocuments.java:27}
& \texorpdfstring{\softtext{stage/cleanup/path}}{stage/cleanup/path} &
Adaptador local y locks \\
\texorpdfstring{\softcode{J/evidencia/DocumentController.java:16}}{J/evidencia/DocumentController.java:16}
& upload & Metadata, contenido y revisión \\
\texorpdfstring{\softcode{J/evidencia/CleanupMain.java:9}}{J/evidencia/CleanupMain.java:9}
& main &
\texorpdfstring{\softtext{Diagnóstico/limpieza}}{Diagnóstico/limpieza}
técnica \\
\texorpdfstring{\softcode{J/evidencia/AnimalQrController.java:15}}{J/evidencia/AnimalQrController.java:15}
& qr/svg & QR local \\
\texorpdfstring{\softcode{J/evidencia/SettlementController.java:13}}{J/evidencia/SettlementController.java:13}
& receipt & Comprobante autorizado \\
\texorpdfstring{\softcode{J/evidencia/SettlementPdf.java:9}}{J/evidencia/SettlementPdf.java:9}
& render & Estructura PDF \\
\texorpdfstring{\softcode{J/integraciones/ProviderController.java:15}}{J/integraciones/ProviderController.java:15}
& \texorpdfstring{\softtext{tariffs/accept}}{tariffs/accept} & Fuente y
captura aceptada \\
\texorpdfstring{\softcode{J/integraciones/OutboxWorker.java:16}}{J/integraciones/OutboxWorker.java:16}
& dispatch & Reclamo, consumo y acuse \\
\texorpdfstring{\softcode{J/integraciones/DiagnosticsController.java:15}}{J/integraciones/DiagnosticsController.java:15}
& snapshot & Diagnóstico ADMIN \\
\texorpdfstring{\softcode{J/reportes/CsvController.java:15}}{J/reportes/CsvController.java:15}
& \texorpdfstring{\softtext{export/csv/cell}}{export/csv/cell} &
Representaciones y protección de celdas \\
\texorpdfstring{\softcode{backend/pom.xml:1}}{backend/pom.xml:1} &
Dependencias & Boot, Java, JDBC, seguridad, sesión y ZXing \\
\texorpdfstring{\softcode{backend/src/main/resources/application.yml:1}}{backend/src/main/resources/application.yml:1}
& Propiedades & Límites, conexiones, sesión y apagado \\
\texorpdfstring{\softcode{backend/src/main/resources/application-production.yml:1}}{backend/src/main/resources/application-production.yml:1}
& Perfil & Cookie Secure \\
\texorpdfstring{\softcode{docs/api/openapi.yaml:1}}{docs/api/openapi.yaml:1}
& Contrato completo & 70
\texorpdfstring{\softtext{operationId;}}{operationId;} tipos, cabeceras
y estados \\
\texorpdfstring{\softcode{docs/development/backend-context.md:1}}{docs/development/backend-context.md:1}
& Contexto vigente & V13/V12, evidencia y despliegue integrado \\
\texorpdfstring{\softcode{docs/architecture/overview.md:1}}{docs/architecture/overview.md:1}
& Diseño & Módulos y clases planeados, diferenciados del código \\
\texorpdfstring{\softcode{db/migrations/V13\_\_correccion\_composicion\_exacta.sql:12}}{db/migrations/V13\_\_correccion\_composicion\_exacta.sql:12}
& Frontera SQL & Resultado idempotente y confirmación canónica \\
\texorpdfstring{\softcode{backend/verification-result-v13.json:1}}{backend/verification-result-v13.json:1}
& Registro histórico & 24 pruebas ejecutadas el 11-oct-2026 \\
\end{longtable}
}

Para verificar afirmaciones de cálculo y restricciones físicas,
continuar con el capítulo de PostgreSQL y las migraciones
\texorpdfstring{\softtext{canónicas.}}{canónicas.} Este capítulo muestra
cómo el programa Java llega a esas funciones, cómo protege y presenta la
comunicación, y dónde termina su
\texorpdfstring{\softtext{responsabilidad.}}{responsabilidad.}

\chapter{\texorpdfstring{Frontend de
\texorpdfstring{\softtext{BoviCierre:}}{BoviCierre:} manual didáctico
del software
implementado}{Frontend de  manual didáctico del software implementado}}\label{frontend-de-bovicierre-manual-diduxe1ctico-del-software-implementado}

\section{1. Qué hace este capítulo y cómo
leerlo}\label{quuxe9-hace-este-capuxedtulo-y-cuxf3mo-leerlo}

Este capítulo explica el frontend real de
\texorpdfstring{\softtext{BoviCierre,}}{BoviCierre,} una aplicación para
registrar y cerrar compras y ventas de ganado en pie, conservar su
evidencia y conocer las obligaciones por cobrar o pagar. Se basa en la
lectura del código vigente el 12 de octubre de 2026. Los recorridos y
resultados citados corresponden a la evidencia local del 11 de octubre;
preparar este capítulo no ejecutó nuevamente la aplicación ni alteró
registros. Todos los ejemplos son DEMO, sintéticos y
\texorpdfstring{\softtext{didácticos.}}{didácticos.} Sus identificadores
se deben sustituir por los devueltos por una API
\texorpdfstring{\softtext{autorizada.}}{autorizada.} No contienen
contraseñas, cookies ni
\texorpdfstring{\softtext{credenciales.}}{credenciales.}

El lector puede empezar por el recorrido operativo y regresar al
catálogo de componentes cuando encuentre un nombre
\texorpdfstring{\softtext{desconocido.}}{desconocido.} También puede
estudiar primero las capas de datos, porque explican por qué ciertos
botones permanecen deshabilitados después de un problema de red. La
interfaz no consiste sólo en una colección de formularios: conserva la
relación entre una intención humana, una versión de un expediente y el
resultado que confirmó el servidor. Comprender esa relación permite
diagnosticar un error sin repetir accidentalmente una operación
\texorpdfstring{\softtext{económica.}}{económica.}

La separación principal es la siguiente. React organiza controles,
presenta resultados y conserva el estado de una
\texorpdfstring{\softtext{interacción.}}{interacción.} La API autentica,
autoriza y conduce comandos hacia las reglas del dominio. PostgreSQL
calcula importes y valida las transiciones
\texorpdfstring{\softtext{económicas.}}{económicas.} Por tanto, una
liquidación visible no se obtiene sumando celdas en el navegador. Es una
representación de un resultado canónico recibido. Aunque el código use
\texorpdfstring{\softcode{Number}}{Number} para páginas, números de lote
o contadores, esos usos no constituyen aritmética monetaria. Importes y
pesos viajan como cadenas decimales.

Hay quince pantallas conceptuales en la documentación de producto, pero
el menú implementado contiene diecinueve destinos. Algunas pantallas
conceptuales abarcan una lista y un detalle; otras acciones adicionales,
como documentos, predios o periodos, aparecen como destinos propios. Esa
diferencia no significa que existan diecinueve aplicaciones ni
diecinueve diseños
\texorpdfstring{\softtext{independientes.}}{independientes.} Todos los
destinos reutilizan la misma estructura, campos, tablas, avisos y reglas
de \texorpdfstring{\softtext{recuperación.}}{recuperación.}

Fuentes principales:
\texorpdfstring{\softcode{frontend/src/App.tsx:7}}{frontend/src/App.tsx:7},
\texorpdfstring{\softcode{frontend/src/screens.tsx:8}}{frontend/src/screens.tsx:8},
\texorpdfstring{\softcode{frontend/src/ui.tsx:4}}{frontend/src/ui.tsx:4},
\texorpdfstring{\softcode{frontend/src/api.ts:13}}{frontend/src/api.ts:13},
\texorpdfstring{\softcode{docs/ui/flows-and-screens.md:18}}{docs/ui/flows-and-screens.md:18},
\texorpdfstring{\softcode{docs/development/frontend-context.md:1}}{docs/development/frontend-context.md:1}.

\section{2. Mapa de archivos, dependencias y
ejecución}\label{mapa-de-archivos-dependencias-y-ejecuciuxf3n}

La aplicación tiene una entrada pequeña, un contenedor de navegación,
una colección de pantallas y cuatro módulos
\texorpdfstring{\softtext{compartidos.}}{compartidos.} Esta organización
permite seguir un dato desde el campo donde se captura hasta la
solicitud HTTP que lo envía. No hay que buscar una segunda biblioteca de
componentes ni un estado global basado en Redux: las dependencias de
ejecución declaradas son React y React DOM. El resto corresponde a
compilación, tipos, generación del contrato y pruebas.

{\def\LTcaptype{none} % do not increment counter
\begin{longtable}[]{@{}
  >{\raggedright\arraybackslash}p{(\linewidth - 4\tabcolsep) * \real{0.3333}}
  >{\raggedright\arraybackslash}p{(\linewidth - 4\tabcolsep) * \real{0.3333}}
  >{\raggedright\arraybackslash}p{(\linewidth - 4\tabcolsep) * \real{0.3333}}@{}}
\toprule\noalign{}
\begin{minipage}[b]{\linewidth}\raggedright
Archivo o grupo
\end{minipage} & \begin{minipage}[b]{\linewidth}\raggedright
Responsabilidad
\end{minipage} & \begin{minipage}[b]{\linewidth}\raggedright
Punto de entrada
\end{minipage} \\
\midrule\noalign{}
\endhead
\bottomrule\noalign{}
\endlastfoot
\texorpdfstring{\softcode{main.tsx}}{main.tsx} y
\texorpdfstring{\softcode{index.html}}{index.html} & Montar React y
cargar estilos &
\texorpdfstring{\softcode{frontend/src/main.tsx:5}}{frontend/src/main.tsx:5} \\
\texorpdfstring{\softcode{App.tsx}}{App.tsx} & Acceso, sede, navegación,
catálogos y despacho de vistas &
\texorpdfstring{\softcode{frontend/src/App.tsx:19}}{frontend/src/App.tsx:19} \\
\texorpdfstring{\softcode{screens.tsx}}{screens.tsx} & Operaciones,
liquidaciones, dinero y controles específicos &
\texorpdfstring{\softcode{frontend/src/screens.tsx:11}}{frontend/src/screens.tsx:11} \\
\texorpdfstring{\softcode{ui.tsx}}{ui.tsx} & Recursos, avisos, campos,
comandos, tablas y listas &
\texorpdfstring{\softcode{frontend/src/ui.tsx:4}}{frontend/src/ui.tsx:4} \\
\texorpdfstring{\softcode{api.ts}}{api.ts} & HTTP, CSRF, errores e
intentos persistidos &
\texorpdfstring{\softcode{frontend/src/api.ts:1}}{frontend/src/api.ts:1} \\
\texorpdfstring{\softcode{domain.ts}}{domain.ts} y
\texorpdfstring{\softcode{workflow.ts}}{workflow.ts} & Representación
decimal, tipos y composición de lotes &
\texorpdfstring{\softcode{frontend/src/domain.ts:1}}{frontend/src/domain.ts:1} \\
\texorpdfstring{\softcode{styles.css}}{styles.css} & Identidad,
estructura y adaptación de tamaños &
\texorpdfstring{\softcode{frontend/src/styles.css:1}}{frontend/src/styles.css:1} \\
\texorpdfstring{\softcode{generated/api-schema.ts}}{generated/api-schema.ts}
& Tipos producidos desde OpenAPI &
\texorpdfstring{\softcode{frontend/package.json:16}}{frontend/package.json:16} \\
\texorpdfstring{\softcode{tests}}{tests} y
\texorpdfstring{\softcode{tests\_e2e}}{tests\_e2e} & Verificación
unitaria y recorridos reales &
\texorpdfstring{\softcode{frontend/playwright.config.ts:2}}{frontend/playwright.config.ts:2} \\
\end{longtable}
}

Las versiones declaradas son React 19.2.4, TypeScript 5.9.3, Vite 8.0.3
y Playwright 1.58.2. El archivo de paquetes fija las versiones y el
archivo de bloqueo permite una instalación reproducible mediante
\texorpdfstring{\softcode{npm ci}}{npm ci}. Los entornos documentados
son Node 20.20.2 o 24.21.0; el rango técnico de
\texorpdfstring{\softcode{engines}}{engines} es más amplio y no debe
confundirse con una prueba realizada en cada versión que admite ese
rango. Saber qué está declarado y qué se ejecutó evita convertir
compatibilidad teórica en evidencia práctica.

La entrada integrada es
\texorpdfstring{\softcode{http://localhost:15180}}{http://localhost:15180},
servida por Caddy y los contenedores locales. Vite ofrece una
alternativa de desarrollo en
\texorpdfstring{\softcode{http://127.0.0.1:15173}}{http://127.0.0.1:15173},
con proxy de \texorpdfstring{\softcode{/api}}{/api} hacia
\texorpdfstring{\softcode{http://127.0.0.1:18080}}{http://127.0.0.1:18080}.
El navegador sigue utilizando rutas relativas del mismo origen. No
descubre nodos PostgreSQL ni selecciona al propietario de una sede. La
elección de base de datos pertenece al backend. Tampoco deben iniciarse
dos procesos de API sobre el mismo puerto para intentar resolver una
falla del frontend.

{\def\LTcaptype{none} % do not increment counter
\begin{longtable}[]{@{}
  >{\raggedright\arraybackslash}p{(\linewidth - 4\tabcolsep) * \real{0.3333}}
  >{\raggedright\arraybackslash}p{(\linewidth - 4\tabcolsep) * \real{0.3333}}
  >{\raggedright\arraybackslash}p{(\linewidth - 4\tabcolsep) * \real{0.3333}}@{}}
\toprule\noalign{}
\begin{minipage}[b]{\linewidth}\raggedright
Script
\end{minipage} & \begin{minipage}[b]{\linewidth}\raggedright
Qué hace
\end{minipage} & \begin{minipage}[b]{\linewidth}\raggedright
Qué no demuestra por sí solo
\end{minipage} \\
\midrule\noalign{}
\endhead
\bottomrule\noalign{}
\endlastfoot
\texorpdfstring{\softcode{dev}}{dev} & Inicia Vite en loopback y puerto
15173 & Persistencia correcta o recorrido aprobado \\
\texorpdfstring{\softcode{generate:api}}{generate:api} & Produce tipos
desde OpenAPI & Que el servidor entregue todos los DTO correctamente \\
\texorpdfstring{\softcode{typecheck}}{typecheck} & Comprueba TypeScript
con \texorpdfstring{\softcode{tsc -b}}{tsc -b} & Resultado funcional en
un navegador \\
\texorpdfstring{\softcode{test}}{test} & Ejecuta pruebas unitarias con
\texorpdfstring{\softcode{tsx}}{tsx} & Comandos reales de PostgreSQL \\
\texorpdfstring{\softcode{build}}{build} & Comprueba tipos y genera
archivos estáticos & Rendimiento GPU ni aceptación humana \\
\texorpdfstring{\softcode{preview}}{preview} & Sirve la compilación
local & Despliegue comercial \\
\texorpdfstring{\softcode{test:e2e}}{test:e2e} & Ejecuta Playwright
según su configuración & Cobertura universal o prueba en teléfono
físico \\
\end{longtable}
}

Estos comandos se explican para que el lector los entienda; no son
instrucciones para ejecutarlos indiscriminadamente sobre un entorno
\texorpdfstring{\softtext{compartido.}}{compartido.} Los recorridos E2E
crean datos DEMO y necesitan una ventana
\texorpdfstring{\softtext{coordinada.}}{coordinada.} La prueba física
del proveedor además requiere que otro responsable detenga y reanude
únicamente ese servicio. El runner no controla Docker por su cuenta.

Fuentes:
\texorpdfstring{\softcode{frontend/package.json:1}}{frontend/package.json:1},
\texorpdfstring{\softcode{frontend/vite.config.ts:3}}{frontend/vite.config.ts:3},
\texorpdfstring{\softcode{frontend/tsconfig.json:1}}{frontend/tsconfig.json:1},
\texorpdfstring{\softcode{frontend/README.md:15}}{frontend/README.md:15},
\texorpdfstring{\softcode{frontend/playwright.config.ts:2}}{frontend/playwright.config.ts:2}.

\section{3. HTML inicial y montaje de
React}\label{html-inicial-y-montaje-de-react}

\texorpdfstring{\softcode{index.html}}{index.html} declara idioma
español, UTF-8, viewport adaptable, color de tema y el elemento
\texorpdfstring{\softcode{root}}{root}. Ese elemento es el punto donde
React administra el árbol de interfaz. El archivo HTML no contiene una
lista precargada de ganado ni una sesión de ejemplo. Su tarea es ofrecer
una base mínima y cargar el módulo de entrada. El título de la pestaña
identifica el producto, mientras que los encabezados internos cambian
con la pantalla activa.

\texorpdfstring{\softcode{main.tsx}}{main.tsx} importa React,
\texorpdfstring{\softcode{createRoot}}{createRoot},
\texorpdfstring{\softcode{App}}{App} y la hoja de estilos. La llamada
\texorpdfstring{\softcode{createRoot(document.getElementById('root')!)}}{createRoot(document.getElementById('root')!)}
establece una única raíz. El operador de afirmación no nula comunica a
TypeScript que el HTML esperado contiene ese elemento; no constituye una
recuperación si alguien elimina el
\texorpdfstring{\softtext{contenedor.}}{contenedor.} Después se
renderiza \texorpdfstring{\softcode{App}}{App} dentro de
\texorpdfstring{\softcode{React.StrictMode}}{React.StrictMode}. Este
modo ayuda a detectar comportamientos problemáticos durante el
desarrollo, pero puede ejecutar ciclos adicionales de efectos. Por eso
una consulta anónima inicial puede repetirse y debe distinguirse de un
comando económico duplicado.

Los eventos de envío requieren una acción del usuario y las claves de
idempotencia se conservan en el módulo de comandos. El montaje de React
no confirma una compra. Esta distinción es útil al interpretar la
consola: ver una consulta repetida en desarrollo no permite concluir que
se hayan creado dos
\texorpdfstring{\softtext{obligaciones.}}{obligaciones.} Hay que
identificar método, ruta y clave del comando antes de diagnosticar
\texorpdfstring{\softtext{duplicación.}}{duplicación.}

El CSS se carga desde la entrada, de modo que componentes de acceso,
catálogo y dinero reciben los mismos tokens. No existe una hoja distinta
para cada sede. La sede modifica el ámbito de datos, no la identidad
gráfica. Si la aplicación no muestra nada, la revisión comienza por el
HTML, la carga del módulo y errores
\texorpdfstring{\softtext{JavaScript;}}{JavaScript;} si muestra una
pantalla pero falla una solicitud, el diagnóstico se mueve a
\texorpdfstring{\softcode{request}}{request}, sus parámetros y el
servidor.

Fuentes:
\texorpdfstring{\softcode{frontend/index.html:2}}{frontend/index.html:2},
\texorpdfstring{\softcode{frontend/src/main.tsx:1}}{frontend/src/main.tsx:1},
\texorpdfstring{\softcode{frontend/src/main.tsx:5}}{frontend/src/main.tsx:5}.

\section{4. Tipos del dominio y representación
exacta}\label{tipos-del-dominio-y-representaciuxf3n-exacta}

\texorpdfstring{\softcode{domain.ts}}{domain.ts} exporta cinco tipos
básicos. \texorpdfstring{\softcode{Money}}{Money} reutiliza el esquema
\texorpdfstring{\softcode{Dinero}}{Dinero} generado desde OpenAPI.
\texorpdfstring{\softcode{Sede}}{Sede} contiene identidad, empresa,
código, nombre, roles y algunos datos opcionales de capacidad o
\texorpdfstring{\softtext{disponibilidad.}}{disponibilidad.}
\texorpdfstring{\softcode{Session}}{Session} conserva las sedes
autorizadas y admite tanto la identidad superior como una estructura de
usuario usada como
\texorpdfstring{\softtext{compatibilidad.}}{compatibilidad.}
\texorpdfstring{\softcode{Row}}{Row} representa objetos de consulta con
campos desconocidos e identificador opcional.
\texorpdfstring{\softcode{Page}}{Page} contiene
\texorpdfstring{\softcode{elementos}}{elementos}, página, tamaño y
total.

Usar \texorpdfstring{\softcode{Row}}{Row} permite que una tabla genérica
reciba columnas distintas, pero también significa que TypeScript no
garantiza automáticamente cada campo leído. Por eso el código verifica
tipos al mostrar dinero, peso o listas, y ofrece textos de ausencia. Los
esquemas generados se emplean con mayor precisión en entradas
específicas, como \texorpdfstring{\softcode{PredioNuevo}}{PredioNuevo} y
\texorpdfstring{\softcode{PesajeNuevo}}{PesajeNuevo}. Es importante
reconocer esta combinación: el frontend tiene tipos de contrato, aunque
no todas sus lecturas sean objetos completamente
\texorpdfstring{\softtext{especializados.}}{especializados.}

\subsection{4.1 Función decimal}\label{funciuxf3n-decimal}

Firma:
\texorpdfstring{\softcode{decimal(value: unknown): string}}{decimal(value: unknown): string}.
Su responsabilidad es representar una cadena numérica recibida, sin
convertirla a coma flotante. Si el valor no es una cadena decimal
reconocida, devuelve
\texorpdfstring{\softcode{No disponible}}{No disponible}. Cuando es
válido, separa parte entera y fracción, añade comas de agrupación a la
parte entera y conserva exactamente la fracción. Por ejemplo,
\texorpdfstring{\softcode{9007199254740993.01}}{9007199254740993.01} se
muestra como
\texorpdfstring{\softcode{9,007,199,254,740,993.01}}{9,007,199,254,740,993.01},
pese a superar el entero que JavaScript representa con precisión como
\texorpdfstring{\softcode{number}}{number}.

La función acepta un signo negativo para presentación, porque un valor
calculado puede \texorpdfstring{\softtext{necesitarlo.}}{necesitarlo.}
No normaliza la escala ni decide si un importe está permitido en una
captura. También conserva \texorpdfstring{\softcode{0.000}}{0.000} como
\texorpdfstring{\softcode{0.000}}{0.000}: tres decimales de peso no se
reducen silenciosamente a dos. La llaman
\texorpdfstring{\softcode{money}}{money},
\texorpdfstring{\softcode{text}}{text} y la comparación de tickets. Su
error visible es un texto de ausencia, no una excepción que derribe la
pantalla.

\subsection{4.2 Función money}\label{funciuxf3n-money}

Firma:
\texorpdfstring{\softcode{money(value: unknown): string}}{money(value: unknown): string}.
Comprueba que exista un objeto con
\texorpdfstring{\softcode{importe}}{importe}, pasa ese importe a
\texorpdfstring{\softcode{decimal}}{decimal} y añade la moneda. Cuando
falta moneda utiliza la presentación MXN; cuando falta el objeto
esperado muestra
\texorpdfstring{\softcode{No disponible}}{No disponible}. Se emplea en
resumen de liquidación, disponible del pago, tarifas, reportes e
\texorpdfstring{\softtext{inventario.}}{inventario.} No calcula el saldo
y no valida que una moneda arbitraria sea comercialmente
\texorpdfstring{\softtext{compatible.}}{compatible.} Esas guardas siguen
perteneciendo al servidor.

La diferencia entre formato y cálculo debe quedar clara. Presentar
\texorpdfstring{\softcode{1000.00}}{1000.00} con una coma facilita la
lectura, pero no cambia su valor ni introduce un nuevo hecho. El texto
mostrado nunca se vuelve a analizar con
\texorpdfstring{\softcode{parseFloat}}{parseFloat} para confirmar una
\texorpdfstring{\softtext{operación.}}{operación.} Si la API omite
saldo, la función no inventa \texorpdfstring{\softcode{0.00}}{0.00}; la
ausencia permanece visible y permite investigar el DTO.

\subsection{\texorpdfstring{4.3 Función
\texorpdfstring{\softtext{canonicalDecimal}}{canonicalDecimal}}{4.3 Función }}\label{funciuxf3n-canonicaldecimal}

Firma:
\texorpdfstring{\softcode{canonicalDecimal(value: string, scale: number): string}}{canonicalDecimal(value: string, scale: number): string}.
Valida una captura positiva mediante una expresión regular, elimina
ceros iniciales redundantes y completa ceros hasta la escala
\texorpdfstring{\softtext{solicitada.}}{solicitada.} Para escala dos
permite hasta dieciséis dígitos enteros; para tres, once; para cuatro,
diez; las otras escalas usadas por porcentajes tienen otra cota. La
salida de \texorpdfstring{\softcode{048.5}}{048.5} con escala cuatro es
\texorpdfstring{\softcode{48.5000}}{48.5000}. La salida de
\texorpdfstring{\softcode{0}}{0} con escala dos es
\texorpdfstring{\softcode{0.00}}{0.00}.

Rechaza exponentes, comas de miles, negativos, infinitos, texto vacío y
exceso de \texorpdfstring{\softtext{precisión.}}{precisión.} Rechazar
\texorpdfstring{\softcode{1.001}}{1.001} al pedir dos decimales evita
que la interfaz decida un redondeo
\texorpdfstring{\softtext{económico.}}{económico.} Su responsabilidad es
validar y normalizar
\texorpdfstring{\softtext{representación.}}{representación.} Si hay que
redondear un cálculo confirmado, lo hace SQL conforme a la política
canónica. \texorpdfstring{\softcode{canonicalDecimal}}{canonicalDecimal}
lanza una excepción con mensaje de captura;
\texorpdfstring{\softcode{CommandForm}}{CommandForm} la captura, enfoca
el resumen y deja al usuario corregir antes de enviar.

\subsection{4.4 Función text}\label{funciuxf3n-text}

Firma:
\texorpdfstring{\softcode{text(value: unknown): string}}{text(value: unknown): string}.
Es el adaptador de presentación más utilizado. Traduce ausencia a
\texorpdfstring{\softcode{No disponible}}{No disponible}, booleanos a
\texorpdfstring{\softcode{Sí}}{Sí} o \texorpdfstring{\softcode{No}}{No},
dinero mediante \texorpdfstring{\softcode{money}}{money}, peso mediante
\texorpdfstring{\softcode{decimal}}{decimal} y la unidad kg, arreglos
mediante la unión de sus representaciones e identificaciones mediante
sistema y código declarado. Un objeto genérico que no reconoce se
representa como \texorpdfstring{\softcode{Ver detalle}}{Ver detalle}.
Finalmente sustituye guiones bajos de estados por espacios.

Esta función explica por qué el código SQL puede usar
\texorpdfstring{\softcode{PENDIENTE\_RECEPCION}}{PENDIENTE\_RECEPCION} y
el usuario leer
\texorpdfstring{\softcode{PENDIENTE RECEPCION}}{PENDIENTE RECEPCION}.
También explica que un objeto de JSON no se vuelque íntegramente en una
celda. La representación compacta facilita trabajar con datos, aunque no
sustituye una vista de detalle cuando el objeto posee información que la
tabla no expone.

\subsection{4.5 Función today}\label{funciuxf3n-today}

Firma:
\texorpdfstring{\softcode{today(date = new Date()): string}}{today(date = new Date()): string}.
Obtiene año, mes y día con
\texorpdfstring{\softcode{Intl.DateTimeFormat}}{Intl.DateTimeFormat} en
\texorpdfstring{\softcode{America/Mexico\_City}}{America/Mexico\_City} y
construye una fecha ISO. No utiliza simplemente la fecha UTC de
\texorpdfstring{\softcode{toISOString}}{toISOString}, porque una tarde
mexicana puede corresponder al día siguiente en UTC. La usan fechas
predeterminadas del trato, pagos, gastos y
\texorpdfstring{\softtext{reconocimiento.}}{reconocimiento.} Su
parámetro opcional permite comprobar instantes fijos en pruebas.

Una fecha predeterminada es una ayuda de captura, no una decisión
contable \texorpdfstring{\softtext{automática.}}{automática.} El usuario
sigue seleccionando periodo y fecha, y el servidor comprueba que la
combinación sea admisible. El ejemplo unitario
\texorpdfstring{\softcode{2026-10-12T01:00:00Z}}{2026-10-12T01:00:00Z}
produce \texorpdfstring{\softcode{2026-10-11}}{2026-10-11} en México,
mostrando el problema que se evitó.

Fuentes:
\texorpdfstring{\softcode{frontend/src/domain.ts:2}}{frontend/src/domain.ts:2},
\texorpdfstring{\softcode{frontend/src/domain.ts:7}}{frontend/src/domain.ts:7},
\texorpdfstring{\softcode{frontend/src/domain.ts:12}}{frontend/src/domain.ts:12},
\texorpdfstring{\softcode{frontend/src/domain.ts:17}}{frontend/src/domain.ts:17},
\texorpdfstring{\softcode{frontend/src/domain.ts:22}}{frontend/src/domain.ts:22},
\texorpdfstring{\softcode{frontend/src/domain.ts:32}}{frontend/src/domain.ts:32},
\texorpdfstring{\softcode{frontend/tests/domain.test.ts:8}}{frontend/tests/domain.test.ts:8}.

\section{5. Dos auxiliares que preservan la identidad
comercial}\label{dos-auxiliares-que-preservan-la-identidad-comercial}

\texorpdfstring{\softcode{replaceLot(lots: Row[], replacement: Row): Row[]}}{replaceLot(lots: Row[], replacement: Row): Row[]}
reemplaza el lote con el mismo
\texorpdfstring{\softcode{numero\_lote}}{numero\_lote}, conserva los
demás y ordena el resultado por ese número. No modifica el arreglo
original. La ordenación usa números de lote, no precios.
\texorpdfstring{\softcode{LotEditor}}{LotEditor} lo llama para construir
el PATCH completo de
\texorpdfstring{\softtext{condiciones.}}{condiciones.} Es relevante
porque el servidor puede relacionar la ordinalidad del arreglo con la
identidad del lote: editar el primero no debe enviarlo al final ni
intercambiar animales con el segundo.

La prueba unitaria conserva la referencia del lote intacto y comprueba
que el precio original del primer objeto no cambie. Además, el recorrido
real editó dos lotes antes y después de un ticket aprobado. Esto no
convierte al auxiliar en un motor contractual: su alcance es construir
una composición ordenada. La autorización para editar, el estado de
borrador y la conservación de versiones se validan en el servidor.

\texorpdfstring{\softcode{operationAnimals(operation: Row, positions: Row[] = []): string[]}}{operationAnimals(operation: Row, positions: Row[] = []): string[]}
obtiene los animales de cada lote. En una compra puede leer
\texorpdfstring{\softcode{lote.animales}}{lote.animales}. En una venta
tiene que resolver
\texorpdfstring{\softcode{lote.posiciones}}{lote.posiciones} contra las
posiciones consultadas y obtener sus
\texorpdfstring{\softcode{animal\_id}}{animal\_id}. Si una posición no
tiene identidad resuelta, lanza un error que exige consultar toda la
\texorpdfstring{\softtext{composición.}}{composición.} Nunca utiliza un
arreglo vacío de animales como sustituto silencioso de una venta
conocida.

El error se captura antes de enviar un ticket. Esta conducta protege un
dato fundamental: el pesaje corresponde a un conjunto concreto de
animales. Tener trescientos kilogramos sin saber qué identidades
integraron el grupo no equivale a un ticket válido.
\texorpdfstring{\softcode{OperationEditor}}{OperationEditor} usa el
auxiliar al presentar pesos opcionales y al construir la propuesta;
\texorpdfstring{\softcode{useSalePositions}}{useSalePositions} lo emplea
para verificar que terminó de resolver las posiciones
\texorpdfstring{\softtext{requeridas.}}{requeridas.}

Fuentes:
\texorpdfstring{\softcode{frontend/src/workflow.ts:3}}{frontend/src/workflow.ts:3},
\texorpdfstring{\softcode{frontend/src/workflow.ts:8}}{frontend/src/workflow.ts:8},
\texorpdfstring{\softcode{frontend/src/screens.tsx:14}}{frontend/src/screens.tsx:14},
\texorpdfstring{\softcode{frontend/src/screens.tsx:20}}{frontend/src/screens.tsx:20},
\texorpdfstring{\softcode{frontend/tests/workflow.test.ts:6}}{frontend/tests/workflow.test.ts:6}.

\section{6. La frontera HTTP: solicitar, reconocer errores y recuperar
comandos}\label{la-frontera-http-solicitar-reconocer-errores-y-recuperar-comandos}

El módulo \texorpdfstring{\softcode{api.ts}}{api.ts} concentra reglas
que sería peligroso copiar de manera distinta en cada pantalla. Sus
variables de módulo son \texorpdfstring{\softcode{csrf}}{csrf},
\texorpdfstring{\softcode{csrfHeader}}{csrfHeader} y
\texorpdfstring{\softcode{principal}}{principal}. No son tokens
almacenados en disco por el frontend.
\texorpdfstring{\softcode{csrf}}{csrf} conserva el valor actual en
memoria; \texorpdfstring{\softcode{csrfHeader}}{csrfHeader} comienza con
un nombre conocido y se ajusta a la cabecera autorizada que entrega el
servidor. \texorpdfstring{\softcode{principal}}{principal} identifica a
quién pertenece el espacio de intentos
\texorpdfstring{\softtext{persistidos.}}{persistidos.} El navegador
maneja la cookie de sesión; el código no la extrae para guardarla en
\texorpdfstring{\softcode{localStorage}}{localStorage}.

\subsection{6.1 ApiError y clasificación del
resultado}\label{apierror-y-clasificaciuxf3n-del-resultado}

\texorpdfstring{\softcode{ApiError}}{ApiError} extiende
\texorpdfstring{\softcode{Error}}{Error}. Su constructor recibe
\texorpdfstring{\softcode{status}}{status} y detalles opcionales:
mensaje, regla, referencia de correlación y lista de campos
\texorpdfstring{\softtext{problemáticos.}}{problemáticos.} Para 401
explica que terminó la sesión; para 403 y 404 utiliza un mismo mensaje
de falta de acceso, evitando diferencias innecesarias de
\texorpdfstring{\softtext{presentación.}}{presentación.} Para 412 avisa
que otra persona cambió el registro y para 428 que falta su versión.
Para 503 ofrece conservar la captura. El estado cero representa un
resultado que no se pudo conocer por red o por interpretación de la
respuesta.

Una excepción no describe automáticamente el estado económico del
servidor. Un POST pudo confirmar y perder su respuesta. Por eso
\texorpdfstring{\softcode{keepAttempt(error: Error): boolean}}{keepAttempt(error: Error): boolean}
clasifica si debe conservarse la
\texorpdfstring{\softtext{intención.}}{intención.} Retiene errores que
no sean \texorpdfstring{\softcode{ApiError}}{ApiError}, estado cero,
401, 408, 409, 425, 429 y cualquier 5 xx. Descarta el intento pendiente
en rechazos conocidos como 400, 403, 404, 412, 413, 422 o 428,
permitiendo corregir una captura resuelta. El 409 tiene tratamiento
conservador y ofrece una acción explícita para resolver y volver a
capturar.

{\def\LTcaptype{none} % do not increment counter
\begin{longtable}[]{@{}
  >{\raggedright\arraybackslash}p{(\linewidth - 4\tabcolsep) * \real{0.3333}}
  >{\raggedright\arraybackslash}p{(\linewidth - 4\tabcolsep) * \real{0.3333}}
  >{\raggedright\arraybackslash}p{(\linewidth - 4\tabcolsep) * \real{0.3333}}@{}}
\toprule\noalign{}
\begin{minipage}[b]{\linewidth}\raggedright
Situación observada
\end{minipage} & \begin{minipage}[b]{\linewidth}\raggedright
Tratamiento del intento
\end{minipage} & \begin{minipage}[b]{\linewidth}\raggedright
Siguiente conducta
\end{minipage} \\
\midrule\noalign{}
\endhead
\bottomrule\noalign{}
\endlastfoot
Respuesta perdida o 5xx & Mantener clave, cuerpo y versión & Consultar
resultado antes de reenviar \\
Sesión expirada & Mantener almacenamiento por principal & Reautenticar
al mismo usuario \\
Comando en curso 425 & Mantener intención & Esperar y volver a
consultar \\
Conflicto 409 & Mantener hasta resolución explícita & Revisar conflicto,
no cambiar clave a escondidas \\
Regla 422 o captura 400 & Resolver el intento rechazado & Corregir
campos y formular una intención nueva \\
Versión 412 o 428 & Resolver el envío inválido & Obtener versión vigente
antes de continuar \\
\end{longtable}
}

\subsection{\texorpdfstring{6.2
\texorpdfstring{\softtext{setPrincipal,}}{setPrincipal,}
\texorpdfstring{\softtext{clearSession}}{clearSession} y
\texorpdfstring{\softtext{refreshCsrf}}{refreshCsrf}}{6.2   y }}\label{setprincipal-clearsession-y-refreshcsrf}

\texorpdfstring{\softcode{setPrincipal(id: string | undefined)}}{setPrincipal(id: string | undefined)}
actualiza la identidad que forma el prefijo del
\texorpdfstring{\softtext{almacenamiento.}}{almacenamiento.} La llaman
\texorpdfstring{\softcode{App}}{App} al recuperar una sesión o completar
login y las pruebas de
\texorpdfstring{\softtext{aislamiento.}}{aislamiento.}
\texorpdfstring{\softcode{clearSession()}}{clearSession()} borra CSRF e
identidad en memoria y emite
\texorpdfstring{\softcode{session-cleared}}{session-cleared} cuando
existe ventana. No recorre ni elimina las entradas anteriores de
\texorpdfstring{\softcode{sessionStorage}}{sessionStorage}. Por ello el
mismo principal puede recuperar posteriormente un intento, mientras que
otro principal no obtiene la misma clave de
\texorpdfstring{\softtext{almacenamiento.}}{almacenamiento.}

\texorpdfstring{\softcode{refreshCsrf()}}{refreshCsrf()} solicita
\texorpdfstring{\softcode{/sesion/csrf}}{/sesion/csrf} mediante
\texorpdfstring{\softcode{request}}{request}, lee token y nombre de
cabecera y admite únicamente
\texorpdfstring{\softcode{X-CSRF-TOKEN}}{X-CSRF-TOKEN} o
\texorpdfstring{\softcode{X-XSRF-TOKEN}}{X-XSRF-TOKEN}. Un nombre
distinto provoca error. Después actualiza las variables en memoria. El
acceso lo llama antes del POST de login y otra vez después, para
utilizar el contexto que corresponde a la sesión recién
\texorpdfstring{\softtext{establecida.}}{establecida.} Una prueba
unitaria comprueba la cabecera recibida, sin presentar ese ensayo como
una prueba de cookie real.

Estas funciones separan tres conceptos que suelen confundirse: identidad
de usuario, cookie de sesión y defensa CSRF. La identidad organiza el
almacenamiento local; la cookie autentica ante el servidor; el token
CSRF acompaña escrituras del mismo origen. Ninguno de esos mecanismos
sustituye la autorización de una sede o un recurso en el backend.

\subsection{6.3 request}\label{request}

Firma:
\texorpdfstring{\softcode{request<T>(path: string, options: RequestInit = \{\}): Promise<\{data: T; etag: string | null\}>}}{request<T>(path: string, options: RequestInit = \{\}): Promise<\{data: T; etag: string | null\}>}.
Construye \texorpdfstring{\softcode{Headers}}{Headers} a partir de las
opciones. Si el cuerpo es una cadena y falta tipo de contenido, declara
JSON. Si la petición tiene método explícito distinto de GET y existe
CSRF, añade la cabecera
\texorpdfstring{\softtext{correspondiente.}}{correspondiente.} Ejecuta
\texorpdfstring{\softcode{fetch('/api/v1' + path, ...)}}{fetch('/api/v1' + path, ...)}
con
\texorpdfstring{\softcode{credentials: 'same-origin'}}{credentials: 'same-origin'}.
Las rutas de negocio llegan normalmente como
\texorpdfstring{\softcode{/sedes/<sede>/...}}{/sedes/<sede>/...}.

Un fallo de \texorpdfstring{\softcode{fetch}}{fetch} se convierte en
\texorpdfstring{\softcode{ApiError(0)}}{ApiError(0)}. Una respuesta no
exitosa intenta leer sus detalles JSON; si es 401, emite
\texorpdfstring{\softcode{session-expired}}{session-expired} antes de
lanzar el error. Una respuesta 204 devuelve datos
\texorpdfstring{\softtext{indefinidos.}}{indefinidos.} Las demás
respuestas exitosas se interpretan como JSON; si esa interpretación
falla, vuelve a tratar el resultado como
\texorpdfstring{\softtext{desconocido.}}{desconocido.} Finalmente
conserva la cabecera \texorpdfstring{\softcode{ETag}}{ETag}, que no es
una cantidad económica, sino una versión para
\texorpdfstring{\softtext{concurrencia.}}{concurrencia.}

Los recursos cancelables pasan
\texorpdfstring{\softcode{signal}}{signal} en
\texorpdfstring{\softcode{options}}{options}. La función no conserva una
caché global de respuestas ni reintenta
\texorpdfstring{\softtext{automáticamente.}}{automáticamente.} Tampoco
incorpora un selector de base de datos. Si una vista muestra datos
obsoletos, hay que investigar su ciclo de montaje, revisión y petición;
no existe un almacén universal oculto que garantice frescura. Las
descargas y QR utilizan enlaces o imágenes del mismo origen, pues
\texorpdfstring{\softcode{request}}{request} espera JSON y no es el
adaptador de todos los tipos binarios.

\subsection{\texorpdfstring{6.4 Attempt,
\texorpdfstring{\softtext{createAttempt}}{createAttempt} y
\texorpdfstring{\softtext{executeAttempt}}{executeAttempt}}{6.4 Attempt,  y }}\label{attempt-createattempt-y-executeattempt}

\texorpdfstring{\softcode{Attempt}}{Attempt} contiene
\texorpdfstring{\softcode{path}}{path},
\texorpdfstring{\softcode{method}}{method},
\texorpdfstring{\softcode{body}}{body},
\texorpdfstring{\softcode{key}}{key} y una versión
\texorpdfstring{\softcode{etag}}{etag} opcional.
\texorpdfstring{\softcode{createAttempt(path, body, method = 'POST', etag?)}}{createAttempt(path, body, method = 'POST', etag?)}
serializa el cuerpo inmediatamente y crea una UUID mediante
\texorpdfstring{\softcode{crypto.randomUUID()}}{crypto.randomUUID()}.
Serializar en ese momento evita que una mutación posterior del objeto de
captura cambie el contenido de un reintento. La clave se crea una vez
por intención; no una vez por cada clic posterior.

\texorpdfstring{\softcode{executeAttempt<T>(attempt, checkResult = false)}}{executeAttempt<T>(attempt, checkResult = false)}
permite consultar previamente el comando. Si
\texorpdfstring{\softcode{checkResult}}{checkResult} está activo y el
resultado está completado, devuelve el resultado persistido sin nuevo
POST. Si no está completado y la consulta permite continuar, utiliza
\texorpdfstring{\softcode{request}}{request} con el método y cuerpo
originales, \texorpdfstring{\softcode{Idempotency-Key}}{Idempotency-Key}
y \texorpdfstring{\softcode{If-Match}}{If-Match} cuando el intento posee
ETag. Un envío de edición queda ligado, por tanto, tanto al contenido
como a la versión que el usuario revisó.

La ecuación conceptual es una identidad de intención, no una operación
numérica:

\[
I=(\text{principal},\text{sede},\text{ruta},\text{método},\text{clave},\text{cuerpo},\text{versión})
\]

Cambiar cualquiera de sus elementos relevantes después de un resultado
desconocido no constituye necesariamente un reintento seguro. La
interfaz protege la captura para que primero se resuelva el comando
anterior. Este principio resulta especialmente importante para pagos,
ajustes, gastos y archivos, donde un registro duplicado podría
representar una decisión humana que nunca ocurrió.

\subsection{\texorpdfstring{6.5
\texorpdfstring{\softtext{storageKey,}}{storageKey,}
\texorpdfstring{\softtext{readPending,}}{readPending,}
\texorpdfstring{\softtext{rememberPending}}{rememberPending} y
\texorpdfstring{\softtext{forgetPending}}{forgetPending}}{6.5    y }}\label{storagekey-readpending-rememberpending-y-forgetpending}

\texorpdfstring{\softcode{storageKey(path)}}{storageKey(path)} es
interno. Devuelve el prefijo
\texorpdfstring{\softcode{bovi:pending:}}{bovi:pending:} seguido del
principal y la ruta, o no devuelve clave si falta identidad. La ruta
incluye la sede.
\texorpdfstring{\softcode{readPending(path)}}{readPending(path)} lee y
analiza la entrada correspondiente; si el almacenamiento o JSON falla,
devuelve ausencia.
\texorpdfstring{\softcode{rememberPending(attempt)}}{rememberPending(attempt)}
guarda el intento serializado cuando existe principal.
\texorpdfstring{\softcode{forgetPending(path)}}{forgetPending(path)}
elimina solamente esa entrada. Los errores de almacenamiento se
absorben, porque el formulario montado aún puede conservar su intento en
memoria.

Este almacenamiento es por pestaña. No es una cola offline ni una
sincronización de borradores entre
\texorpdfstring{\softtext{dispositivos.}}{dispositivos.} Tampoco cifra
el contenido comercial que una intención ordinaria guarda. La aplicación
evita persistir contraseña o bytes de un archivo en sus formularios
especiales, pero un cuerpo ordinario puede contener concepto,
identificadores e importe. El acceso al mismo origen y la seguridad del
entorno del navegador siguen siendo
\texorpdfstring{\softtext{relevantes.}}{relevantes.}

Después de recargar, un formulario ordinario puede encontrar su
intención y ofrecer
\texorpdfstring{\softtext{recuperación.}}{recuperación.} Un campo que
alguien estaba escribiendo pero que nunca envió no adquiere esa
persistencia
\texorpdfstring{\softtext{automáticamente.}}{automáticamente.} Es
diferente decir «el borrador guardado existe en SQL», «el intento
enviado está pendiente» y «un texto todavía vive sólo en un
\texorpdfstring{\softcode{useState}}{useState}». El capítulo usa esas
expresiones por separado.

\subsection{\texorpdfstring{6.6
\texorpdfstring{\softtext{commandName}}{commandName} y
\texorpdfstring{\softtext{commandResult}}{commandResult}}{6.6  y }}\label{commandname-y-commandresult}

\texorpdfstring{\softcode{commandName(attempt)}}{commandName(attempt)}
es interno y traduce ruta y método al nombre canónico de comando.
Reconoce creación de contraparte, predio, animal, operación, gasto,
pago, documento, periodo y usuario; reconoce PATCH de operación o
usuario; y distingue acciones como confirmar, acordar, reabrir, aprobar,
corregir, repartir, aplicar, revertir, reembolsar o asignar
\texorpdfstring{\softtext{membresía.}}{membresía.} La aceptación de
tarifa posee un caso específico por su ruta más profunda. Reabrir una
operación se diferencia de reabrir un periodo.

\texorpdfstring{\softcode{commandResult<T>(attempt)}}{commandResult<T>(attempt)}
obtiene ese nombre y la sede desde la ruta, y consulta
\texorpdfstring{\softcode{/sedes/<sede>/comandos/<comando>/<clave>}}{/sedes/<sede>/comandos/<comando>/<clave>}.
\texorpdfstring{\softcode{COMPLETADA}}{COMPLETADA} devuelve
\texorpdfstring{\softcode{resultado}}{resultado}.
\texorpdfstring{\softcode{EN\_CURSO}}{EN\_CURSO} lanza 425 y evita otro
envío. Un estado desconocido produce resultado incierto. Un 404 del
lookup significa que no encontró ese comando para continuar con el mismo
intento; otros errores se propagan. Si la ruta no tiene nombre
reconocido, devuelve
\texorpdfstring{\softcode{completed: false}}{completed: false}, una
razón para mantener sincronizada esa traducción al añadir comandos.

El lookup no revela resultados de cualquier usuario. El servidor
autentica y verifica el ámbito. El frontend organiza las claves para
ayudar al usuario correcto, pero la garantía final no puede depender de
que alguien no altere
\texorpdfstring{\softcode{sessionStorage}}{sessionStorage}. Las pruebas
de reintento verifican que después de una respuesta entregada como 502
el navegador recupera el resultado con un único POST real.

Fuentes:
\texorpdfstring{\softcode{frontend/src/api.ts:1}}{frontend/src/api.ts:1},
\texorpdfstring{\softcode{frontend/src/api.ts:6}}{frontend/src/api.ts:6},
\texorpdfstring{\softcode{frontend/src/api.ts:11}}{frontend/src/api.ts:11},
\texorpdfstring{\softcode{frontend/src/api.ts:13}}{frontend/src/api.ts:13},
\texorpdfstring{\softcode{frontend/src/api.ts:23}}{frontend/src/api.ts:23},
\texorpdfstring{\softcode{frontend/src/api.ts:27}}{frontend/src/api.ts:27},
\texorpdfstring{\softcode{frontend/src/api.ts:33}}{frontend/src/api.ts:33},
\texorpdfstring{\softcode{frontend/src/api.ts:42}}{frontend/src/api.ts:42},
\texorpdfstring{\softcode{frontend/tests/domain.test.ts:19}}{frontend/tests/domain.test.ts:19},
\texorpdfstring{\softcode{frontend/tests/workflow.test.ts:25}}{frontend/tests/workflow.test.ts:25}.

\section{7. La biblioteca pequeña de interfaz
compartida}\label{la-biblioteca-pequeuxf1a-de-interfaz-compartida}

\subsection{\texorpdfstring{7.1
\texorpdfstring{\softtext{useResource}}{useResource} y la consulta
local}{7.1  y la consulta local}}\label{useresource-y-la-consulta-local}

Firma:
\texorpdfstring{\softcode{useResource<T>(path: string, revision = 0)}}{useResource<T>(path: string, revision = 0)}.
Es un hook de consulta, usado por listas y detalles. Conserva un estado
con datos opcionales, ETag opcional,
\texorpdfstring{\softcode{loading}}{loading} y error; además conserva un
contador local de recarga. Un efecto depende de ruta, revisión y
contador. Cada ejecución marca carga, crea
\texorpdfstring{\softcode{AbortController}}{AbortController}, solicita
datos y actualiza estado sólo mientras permanece activa. Al desmontarse
o cambiar dependencias, desactiva el efecto y aborta la solicitud.

El retorno añade \texorpdfstring{\softcode{reload}}{reload}, una función
que incrementa el contador. Pulsar «Actualizar consulta» no ejecuta un
cálculo en memoria; produce otra consulta. El parámetro
\texorpdfstring{\softcode{revision}}{revision} permite que una escritura
confirmada invalide consultas asociadas. Es habitual que
\texorpdfstring{\softcode{onDone}}{onDone} incremente una revisión de
pantalla y que las listas vuelvan a cargar sus datos
\texorpdfstring{\softtext{canónicos.}}{canónicos.}

Este hook no mantiene resultados antiguos durante la carga: comienza una
nueva consulta con
\texorpdfstring{\softcode{\{loading: true\}}}{\{loading: true\}}.
Tampoco es una caché compartida entre pantallas. El aislamiento depende
de rutas con sede, de la clave de montaje del contenedor y de la
\texorpdfstring{\softtext{cancelación.}}{cancelación.} Al cambiar de B a
A, un resultado tardío de B no debe volver a poblar la vista. El
recorrido real retuvo precisamente una consulta de B y comprobó su
aborto al cambiar de sede.

\subsection{7.2 Notice y State}\label{notice-y-state}

\texorpdfstring{\softcode{Notice(\{error\})}}{Notice(\{error\})} produce
un aviso con \texorpdfstring{\softcode{role="alert"}}{role="alert"}
cuando existe error. Expone el mensaje y, si es
\texorpdfstring{\softcode{ApiError}}{ApiError} con regla, añade regla y
referencia de \texorpdfstring{\softtext{correlación.}}{correlación.} No
transforma una negativa en éxito. Si no existe error, no renderiza
contenido. Lo llaman formularios, recursos y pantallas que realizan
alguna actualización directa.

\texorpdfstring{\softcode{State(\{loading, error, empty, reload, children\})}}{State(\{loading, error, empty, reload, children\})}
organiza estados comunes. Cargando muestra un texto con
\texorpdfstring{\softcode{role="status"}}{role="status"}. Error muestra
\texorpdfstring{\softcode{Notice}}{Notice} y «Volver a
\texorpdfstring{\softtext{consultar».}}{consultar».} Vacío muestra una
explicación inicial. Si ninguno aplica, devuelve los hijos. El orden
importa: una carga no se presenta como lista vacía y un error no se
confunde con ausencia de registros. Algunas vistas proporcionan mensajes
vacíos más específicos dentro de sus hijos.

\texorpdfstring{\softcode{reload}}{reload} sólo pertenece a consultas.
El reintento de un comando tiene otro botón y otra lógica. Mezclar ambos
sería peligroso: repetir un GET para conocer datos actuales no equivale
a repetir un POST \texorpdfstring{\softtext{económico.}}{económico.} La
biblioteca hace visible esa diferencia mediante funciones y controles
distintos.

\subsection{\texorpdfstring{7.3
\texorpdfstring{\softtext{FieldErrors,}}{FieldErrors,} FormScope, Field
y
FormField}{7.3  FormScope, Field y FormField}}\label{fielderrors-formscope-field-y-formfield}

\texorpdfstring{\softcode{FieldErrors}}{FieldErrors} es un contexto con
mensajes por nombre. \texorpdfstring{\softcode{FormScope}}{FormScope}
aporta un prefijo de identificadores por
\texorpdfstring{\softtext{formulario.}}{formulario.}
\texorpdfstring{\softcode{Field}}{Field} recibe etiqueta, nombre, tipo,
obligatoriedad, valor controlado,
\texorpdfstring{\softcode{onChange}}{onChange}, opciones, ayuda, error y
\texorpdfstring{\softtext{deshabilitación.}}{deshabilitación.} Construye
un identificador usando el contexto; conecta etiqueta y control; y
asocia ayuda o error mediante
\texorpdfstring{\softcode{aria-describedby}}{aria-describedby}. Añade
\texorpdfstring{\softcode{aria-invalid}}{aria-invalid} cuando
\texorpdfstring{\softtext{corresponde.}}{corresponde.}

Si recibe opciones genera un select con una opción inicial
\texorpdfstring{\softtext{«Selecciona…».}}{«Selecciona…».} En caso
contrario genera input. Contraseñas de acceso permiten gestores con
\texorpdfstring{\softcode{current-password}}{current-password}; login
usa \texorpdfstring{\softcode{username}}{username}. Los campos de
importe, precio y neto utilizan
\texorpdfstring{\softcode{inputMode="decimal"}}{inputMode="decimal"},
que facilita el teclado apropiado pero no valida por sí mismo una regla
monetaria. La obligatoriedad nativa y la normalización decimal son
defensas \texorpdfstring{\softtext{complementarias.}}{complementarias.}

\texorpdfstring{\softcode{FormField}}{FormField}, exportado desde
\texorpdfstring{\softcode{App.tsx}}{App.tsx}, envuelve
\texorpdfstring{\softcode{Field}}{Field} con un
\texorpdfstring{\softcode{useState}}{useState} inicial. Firma resumida:
\texorpdfstring{\softcode{FormField(\{name, label, initial, options, required, type, help\})}}{FormField(\{name, label, initial, options, required, type, help\})}.
Permite escribir campos de captura sin repetir el par
\texorpdfstring{\softtext{valor/setter.}}{valor/setter.} El valor
inicial se toma al montar; por eso varias pantallas utilizan
\texorpdfstring{\softcode{key}}{key} para remontar un formulario cuando
cambia de destino o versión. Modificar la prop
\texorpdfstring{\softcode{initial}}{initial} de un componente ya montado
no equivale a sincronizar su estado
\texorpdfstring{\softtext{automáticamente.}}{automáticamente.}

\subsection{\texorpdfstring{7.4
\texorpdfstring{\softtext{CommandForm}}{CommandForm} y sus funciones
internas}{7.4  y sus funciones internas}}\label{commandform-y-sus-funciones-internas}

Firma:
\texorpdfstring{\softcode{CommandForm(\{path, build, method = 'POST', etag, children, label, onDone\})}}{CommandForm(\{path, build, method = 'POST', etag, children, label, onDone\})}.
\texorpdfstring{\softcode{build}}{build} convierte
\texorpdfstring{\softcode{FormData}}{FormData} en el cuerpo que exige la
acción. \texorpdfstring{\softcode{onDone}}{onDone} recibe el resultado
del servidor y puede navegar o aumentar una revisión. El componente
posee un identificador \texorpdfstring{\softcode{useId}}{useId}, intento
pendiente, banderas \texorpdfstring{\softcode{busy}}{busy} y
\texorpdfstring{\softcode{success}}{success}, error y una referencia al
resumen de mensajes. Inicializa el intento mediante
\texorpdfstring{\softcode{readPending(path)}}{readPending(path)}.

La función interna \texorpdfstring{\softcode{submit}}{submit} cancela el
envío nativo. Si ya hay un intento, ejecuta exactamente ese intento con
lookup previo. Si no existe, lee el formulario, llama a
\texorpdfstring{\softcode{build}}{build}, crea y guarda el intento y lo
ejecuta. Un error de construcción, como decimal inválido, se muestra sin
enviar un cuerpo \texorpdfstring{\softtext{defectuoso.}}{defectuoso.} La
función interna \texorpdfstring{\softcode{execute}}{execute} administra
carga, error y éxito; elimina el pendiente después de un resultado
confirmado; llama a \texorpdfstring{\softcode{onDone}}{onDone}; y
consulta \texorpdfstring{\softcode{keepAttempt}}{keepAttempt} al fallar.
Finalmente vuelve a habilitar la acción de espera, pero puede mantener
protegidos los campos por una intención no resuelta.

El \texorpdfstring{\softcode{fieldset}}{fieldset} queda deshabilitado
mientras se espera o existe un pendiente. La acción principal cambia
entre el rótulo original, «Esperando respuesta\ldots» y «Reintentar el
mismo comando». Para 409 aparece «Resolver conflicto y volver a
capturar», que elimina expresamente el pendiente. Esa acción exige
revisión humana del conflicto; no genera automáticamente otro POST. Los
errores de campos producen enlaces hacia los controles identificados con
el prefijo del formulario, y el resumen se enfoca después del envío
fallido.

Un error de prueba ilustró una consecuencia de este diseño: buscar un
formulario por el texto «Confirmar venta» deja de encontrarlo después de
un resultado incierto, porque el botón cambia de nombre. La prueba final
utiliza un campo estable del
\texorpdfstring{\softtext{formulario.}}{formulario.} Entender la máquina
de estados evita corregir el producto para satisfacer un selector de
prueba \texorpdfstring{\softtext{equivocado.}}{equivocado.}

\subsection{7.5 DataTable y List}\label{datatable-y-list}

\texorpdfstring{\softcode{DataTable(\{rows, columns, onSelect\})}}{DataTable(\{rows, columns, onSelect\})}
presenta una tabla dentro de una región enfocable y
\texorpdfstring{\softtext{desplazable.}}{desplazable.} Cada columna
tiene clave, etiqueta y una marca opcional de dinero. Encabezados
utilizan \texorpdfstring{\softcode{scope="col"}}{scope="col"}. Las
celdas pasan por \texorpdfstring{\softcode{text}}{text}; el estado
aparece como badge. Las columnas monetarias tienen alineación y cifras
tabulares. Si existe \texorpdfstring{\softcode{onSelect}}{onSelect},
añade «Abrir» con texto de contexto accesible, de modo que dos botones
visualmente iguales no sean indistinguibles para quien navega por
controles.

\texorpdfstring{\softcode{List(\{path, columns, revision = 0, onSelect, filters = []\})}}{List(\{path, columns, revision = 0, onSelect, filters = []\})}
combina consulta paginada, búsqueda y filtros. Conserva página, texto
escrito, búsqueda aplicada y valores de filtros. Envía página desde uno
y tamaño veinte. El texto escrito no produce necesariamente una
solicitud por tecla: enviar «Buscar» copia la consulta recortada y
reinicia página. Cambiar un filtro también regresa a la primera página.
«Anterior» se deshabilita en página uno y «Siguiente» cuando ya se
alcanza el total.

La búsqueda y los filtros consultan el servidor. La tabla no descarga
toda la sede para filtrarla en
\texorpdfstring{\softtext{JavaScript.}}{JavaScript.} Un detalle puede
seleccionarse mediante \texorpdfstring{\softcode{onSelect}}{onSelect},
pero una lista sin ese callback permanece de consulta. El número de
resultados pertenece al servidor; no se deduce contando sólo las filas
visibles. Algunas pantallas usan consultas especiales fuera de
\texorpdfstring{\softcode{List}}{List}, como cartera, pero conservan el
mismo criterio de paginación y mensajes.

Fuentes:
\texorpdfstring{\softcode{frontend/src/ui.tsx:4}}{frontend/src/ui.tsx:4},
\texorpdfstring{\softcode{frontend/src/ui.tsx:10}}{frontend/src/ui.tsx:10},
\texorpdfstring{\softcode{frontend/src/ui.tsx:11}}{frontend/src/ui.tsx:11},
\texorpdfstring{\softcode{frontend/src/ui.tsx:17}}{frontend/src/ui.tsx:17},
\texorpdfstring{\softcode{frontend/src/ui.tsx:19}}{frontend/src/ui.tsx:19},
\texorpdfstring{\softcode{frontend/src/ui.tsx:23}}{frontend/src/ui.tsx:23},
\texorpdfstring{\softcode{frontend/src/ui.tsx:30}}{frontend/src/ui.tsx:30},
\texorpdfstring{\softcode{frontend/src/ui.tsx:33}}{frontend/src/ui.tsx:33},
\texorpdfstring{\softcode{frontend/src/App.tsx:15}}{frontend/src/App.tsx:15}.

\section{8. App: acceso, sede y
navegación}\label{app-acceso-sede-y-navegaciuxf3n}

\texorpdfstring{\softcode{App()}}{App()} es el componente exportado por
defecto. Sus estados principales son sesión, carga inicial, error
global, sede, pantalla y fila activa. La fila activa aporta un
identificador para abrir detalles. Al montar, registra listeners de
expiración de sesión y cambio de hash, intenta recuperar
\texorpdfstring{\softcode{/sesion/actual}}{/sesion/actual}, refresca
CSRF y establece principal. Un 401 anónimo inicial se maneja como
ausencia de sesión; otros errores pueden aparecer en acceso. La limpieza
del efecto elimina los listeners.

Un segundo efecto observa la sesión e intenta seleccionar el parámetro
\texorpdfstring{\softcode{sede}}{sede} de la URL sólo si coincide con
una sede autorizada y
\texorpdfstring{\softtext{disponible.}}{disponible.} Conocer una UUID no
concede acceso. El selector inicial muestra las sedes de la sesión, sus
roles y disponibilidad declarada. Si no hay sedes, explica que se debe
pedir acceso. La interfaz no utiliza una sede de ejemplo escondida para
que el recorrido parezca conectado.

\texorpdfstring{\softcode{readRoute()}}{readRoute()} interpreta el hash
como pantalla e
\texorpdfstring{\softtext{identificador.}}{identificador.} Usa
operaciones como destino predeterminado y decodifica el
\texorpdfstring{\softtext{identificador.}}{identificador.}
\texorpdfstring{\softcode{chooseSede(site)}}{chooseSede(site)} escribe
la sede en el query string mediante
\texorpdfstring{\softcode{replaceState}}{replaceState} y actualiza el
estado. \texorpdfstring{\softcode{go(next, row?)}}{go(next, row?)}
cambia pantalla y fila activa, conserva la sede, construye el hash,
añade historial con \texorpdfstring{\softcode{pushState}}{pushState} e
intenta enfocar el título. No usa una biblioteca de router: el
enrutamiento local está contenido en estas funciones y el listener de
hash.

\texorpdfstring{\softcode{Brand()}}{Brand()} representa el nombre y
símbolo del producto; su SVG decorativo está oculto al árbol accesible.
\texorpdfstring{\softcode{description(screen)}}{description(screen)}
selecciona una explicación breve de la pantalla o utiliza una
descripción general. Son funciones de presentación sin
\texorpdfstring{\softtext{solicitudes.}}{solicitudes.} El nombre de la
sede y los roles del encabezado provienen de la sesión, no del diseño
gráfico.

\texorpdfstring{\softcode{Login(\{onLogin, initialError\})}}{Login(\{onLogin, initialError\})}
conserva login, contraseña, error y estado de espera. Enviar cancela la
navegación del formulario, obtiene CSRF anónimo, envía login, elimina la
contraseña del estado después de éxito, refresca CSRF y recupera la
sesión. \texorpdfstring{\softcode{onLogin}}{onLogin} comunica el
resultado al \texorpdfstring{\softtext{contenedor.}}{contenedor.} Una
negativa permanece como mensaje; un envío pendiente cambia la etiqueta
del botón. El capítulo no reproduce el contenido de contraseña ni ofrece
una credencial DEMO fija.

\texorpdfstring{\softcode{Workspace(\{sede, screen, active, go\})}}{Workspace(\{sede, screen, active, go\})}
construye
\texorpdfstring{\softcode{base = '/sedes/' + sede.id}}{base = '/sedes/' + sede.id},
conserva una revisión local y define
\texorpdfstring{\softcode{bump}}{bump} para aumentarla después de
comandos. Calcula \texorpdfstring{\softcode{operator}}{operator} a
partir de ADMIN u OPERADOR. Despacha cada destino hacia listas,
formularios o componentes
\texorpdfstring{\softtext{especializados.}}{especializados.} Su clave de
montaje combina sede, pantalla e identificador activo. Cambiar ámbito
desmonta el estado anterior, evitando conservar una selección de animal
o un formulario de otra
\texorpdfstring{\softtext{operación.}}{operación.}

\texorpdfstring{\softcode{str(f, name)}}{str(f, name)} obtiene texto de
\texorpdfstring{\softcode{FormData}}{FormData}, usando vacío cuando
falta.
\texorpdfstring{\softcode{cash(f, name = 'importe')}}{cash(f, name = 'importe')}
construye
\texorpdfstring{\softcode{\{importe: canonicalDecimal(..., 2), moneda: 'MXN'\}}}{\{importe: canonicalDecimal(..., 2), moneda: 'MXN'\}}.
No suma dinero.
\texorpdfstring{\softcode{PartyField(\{base, initial, name, label, required\})}}{PartyField(\{base, initial, name, label, required\})}
combina búsqueda explícita y select de contrapartes: conserva query y
search, solicita hasta cien resultados y puede buscar por texto o
identidad. Se reutiliza para
\texorpdfstring{\softtext{comprador/vendedor,}}{comprador/vendedor,}
responsable de predio y proveedor de gasto. Si la consulta falla muestra
aviso. La búsqueda permite encontrar un registro que no estaba entre las
primeras cien opciones; no es un paginador completo del catálogo.

\subsection{8.1 Las quince pantallas conceptuales y sus
destinos}\label{las-quince-pantallas-conceptuales-y-sus-destinos}

{\def\LTcaptype{none} % do not increment counter
\begin{longtable}[]{@{}
  >{\raggedright\arraybackslash}p{(\linewidth - 4\tabcolsep) * \real{0.3333}}
  >{\raggedright\arraybackslash}p{(\linewidth - 4\tabcolsep) * \real{0.3333}}
  >{\raggedright\arraybackslash}p{(\linewidth - 4\tabcolsep) * \real{0.3333}}@{}}
\toprule\noalign{}
\begin{minipage}[b]{\linewidth}\raggedright
Pantalla conceptual
\end{minipage} & \begin{minipage}[b]{\linewidth}\raggedright
Destino o composición real
\end{minipage} & \begin{minipage}[b]{\linewidth}\raggedright
Fuente
\end{minipage} \\
\midrule\noalign{}
\endhead
\bottomrule\noalign{}
\endlastfoot
1. Acceso y selección & \texorpdfstring{\softcode{Login}}{Login} y
selector de sede de \texorpdfstring{\softcode{App}}{App} &
\texorpdfstring{\softcode{frontend/src/App.tsx:26}}{frontend/src/App.tsx:26} \\
2. Bandeja & \texorpdfstring{\softcode{operaciones}}{operaciones}, lista
con tipo y estado &
\texorpdfstring{\softcode{frontend/src/App.tsx:37}}{frontend/src/App.tsx:37} \\
3. Contrapartes & Directorio y alta; predios como destino adicional &
\texorpdfstring{\softcode{frontend/src/App.tsx:38}}{frontend/src/App.tsx:38} \\
4. Ficha animal & Lista,
\texorpdfstring{\softcode{AnimalDetail}}{AnimalDetail} y
\texorpdfstring{\softcode{AnimalCycles}}{AnimalCycles} &
\texorpdfstring{\softcode{frontend/src/App.tsx:40}}{frontend/src/App.tsx:40} \\
5. Inventario & Posiciones, costo, estado y CSV &
\texorpdfstring{\softcode{frontend/src/App.tsx:41}}{frontend/src/App.tsx:41} \\
6. Compra y venta & Alta de borrador y
\texorpdfstring{\softcode{OperationEditor}}{OperationEditor} &
\texorpdfstring{\softcode{frontend/src/App.tsx:35}}{frontend/src/App.tsx:35} \\
7. Tickets & Lista de operaciones y
\texorpdfstring{\softcode{TicketHistory}}{TicketHistory} &
\texorpdfstring{\softcode{frontend/src/App.tsx:52}}{frontend/src/App.tsx:52} \\
8. Liquidación & Lista y
\texorpdfstring{\softcode{Settlement}}{Settlement} &
\texorpdfstring{\softcode{frontend/src/App.tsx:44}}{frontend/src/App.tsx:44} \\
9. Dinero & \texorpdfstring{\softcode{Payments}}{Payments}, validación y
aplicación &
\texorpdfstring{\softcode{frontend/src/App.tsx:45}}{frontend/src/App.tsx:45} \\
10. Expediente & Lista de operaciones y
\texorpdfstring{\softcode{CaseFile}}{CaseFile} &
\texorpdfstring{\softcode{frontend/src/App.tsx:36}}{frontend/src/App.tsx:36} \\
11. Cartera & \texorpdfstring{\softcode{Report}}{Report} con búsqueda y
vencimiento &
\texorpdfstring{\softcode{frontend/src/App.tsx:46}}{frontend/src/App.tsx:46} \\
12. Margen & \texorpdfstring{\softcode{Report}}{Report} con agregados
canónicos &
\texorpdfstring{\softcode{frontend/src/App.tsx:46}}{frontend/src/App.tsx:46} \\
13. Usuarios & \texorpdfstring{\softcode{Users}}{Users} y controles
ADMIN &
\texorpdfstring{\softcode{frontend/src/App.tsx:54}}{frontend/src/App.tsx:54} \\
14. Sincronización &
\texorpdfstring{\softcode{ControlView}}{ControlView} de consulta &
\texorpdfstring{\softcode{frontend/src/App.tsx:55}}{frontend/src/App.tsx:55} \\
15. Simulador & Lista y \texorpdfstring{\softcode{Simulator}}{Simulator}
&
\texorpdfstring{\softcode{frontend/src/App.tsx:56}}{frontend/src/App.tsx:56} \\
\end{longtable}
}

Los destinos adicionales son predios, movimientos, gastos, periodos y
documentos; parte de su función se integra en las pantallas
\texorpdfstring{\softtext{conceptuales.}}{conceptuales.} El menú agrupa
Operaciones, Inventario, Dinero y Control. En escritorio aparece una
navegación lateral; en móvil, un selector nativo agrupado. Un destino
desconocido produce un mensaje de función sin endpoint autorizado, en
lugar de inventar datos. Ese fallback no demuestra que una función
pendiente esté conectada.

\subsection{8.2 Catálogos y ficha
animal}\label{catuxe1logos-y-ficha-animal}

Las contrapartes tienen nombre, tipo de persona, RFC y teléfono; el alta
usa \texorpdfstring{\softcode{CommandForm}}{CommandForm} y luego aumenta
revisión. Los predios agregan nombre obligatorio, municipio, entidad,
UPP declarada y responsable opcional. El constructor recorta el nombre y
rechaza un texto vacío. Su aviso aclara que la declaración comercial no
acredita propiedad ni autorización sanitaria. Los animales se registran
con sexo y una identificación de sistema INTERNO, SINIIGA, SINIDA u
OTRO.

\texorpdfstring{\softcode{AnimalDetail(\{base, id\})}}{AnimalDetail(\{base, id\})}
solicita la ficha, muestra situación y convergencia, presenta
identificaciones y solicita el QR mediante una imagen autorizada del
mismo origen. La imagen no concede una sesión. Su texto alternativo
explica que abre una ficha con acceso
\texorpdfstring{\softtext{autorizado.}}{autorizado.}
\texorpdfstring{\softcode{AnimalCycles(\{base, id\})}}{AnimalCycles(\{base, id\})}
conserva página y consulta posiciones con
\texorpdfstring{\softcode{q}}{q} igual al animal, tamaño veinte y sede
en la ruta. Muestra número de ciclo, estado y costo vigente. La clave
adicional \texorpdfstring{\softcode{base + id}}{base + id} reinicia su
paginación al cambiar animal o sede.

Una recompra no debe sobrescribir el costo de un ciclo cerrado anterior.
La tabla expone precisamente esa historia: un animal permanente puede
tener un ciclo vendido y entregado y otro
\texorpdfstring{\softtext{disponible.}}{disponible.} La disponibilidad
se consulta en inventario; no se obtiene suponiendo que toda identidad
registrada está lista para vender. Las altas de catálogo no confirman
por sí mismas operaciones
\texorpdfstring{\softtext{económicas.}}{económicas.}

Fuentes:
\texorpdfstring{\softcode{frontend/src/App.tsx:13}}{frontend/src/App.tsx:13},
\texorpdfstring{\softcode{frontend/src/App.tsx:14}}{frontend/src/App.tsx:14},
\texorpdfstring{\softcode{frontend/src/App.tsx:16}}{frontend/src/App.tsx:16},
\texorpdfstring{\softcode{frontend/src/App.tsx:18}}{frontend/src/App.tsx:18},
\texorpdfstring{\softcode{frontend/src/App.tsx:19}}{frontend/src/App.tsx:19},
\texorpdfstring{\softcode{frontend/src/App.tsx:23}}{frontend/src/App.tsx:23},
\texorpdfstring{\softcode{frontend/src/App.tsx:31}}{frontend/src/App.tsx:31},
\texorpdfstring{\softcode{frontend/src/App.tsx:33}}{frontend/src/App.tsx:33},
\texorpdfstring{\softcode{frontend/src/App.tsx:59}}{frontend/src/App.tsx:59}.

\section{9. Operaciones, lotes, pesajes y
acuerdos}\label{operaciones-lotes-pesajes-y-acuerdos}

\texorpdfstring{\softcode{PeriodFields(\{base\})}}{PeriodFields(\{base\})}
consulta hasta cien periodos y presenta sólo los que llegan como
ABIERTO; ofrece fecha de reconocimiento con
\texorpdfstring{\softcode{today()}}{today()}. No abre un periodo
\texorpdfstring{\softtext{automáticamente.}}{automáticamente.}
\texorpdfstring{\softcode{recognition(f)}}{recognition(f)} exporta un
constructor con \texorpdfstring{\softcode{periodo\_id}}{periodo\_id} y
\texorpdfstring{\softcode{fecha\_reconocimiento}}{fecha\_reconocimiento};
lo reutilizan confirmación, validación de pago, aplicación, gastos,
ajustes, reversos, reembolsos y bajas. Cada comando conserva su contexto
económico \texorpdfstring{\softtext{explícito.}}{explícito.}

\texorpdfstring{\softcode{inputDecimal(value)}}{inputDecimal(value)} es
interno y adapta valores iniciales de condiciones: lee
\texorpdfstring{\softcode{importe}}{importe} si recibe dinero, devuelve
vacío para \texorpdfstring{\softcode{undefined}}{undefined} y convierte
otros valores a texto. No calcula ni redondea. Lo llama
\texorpdfstring{\softcode{LotEditor}}{LotEditor} para campos de cargos y
\texorpdfstring{\softtext{bonificaciones.}}{bonificaciones.} Este
auxiliar de representación no sustituye la validación posterior de
\texorpdfstring{\softcode{canonicalDecimal}}{canonicalDecimal} o del
contrato.

\texorpdfstring{\softcode{OperationEditor(\{base, sede, id, go\})}}{OperationEditor(\{base, sede, id, go\})}
conserva revisión, paso y último ticket propuesto. Consulta la operación
y su ETag, resuelve posiciones de venta y comprueba permiso de
\texorpdfstring{\softtext{operación.}}{operación.} Los pasos son
condiciones, pesaje y revisión. Muestra folio, estado operativo, estado
físico, contraparte y versión. Un bloque plegable presenta contratos
\texorpdfstring{\softtext{anteriores.}}{anteriores.} Si la operación
está confirmada, comunica que cierre económico, recepción, entrega y
pago son hechos separados; permite consultar liquidación sin ofrecer
reescribir el cierre.

El constructor de ticket selecciona un lote, conserva el neto como
\texorpdfstring{\softcode{Peso\{kg\}}}{Peso\{kg\}} de escala tres,
convierte fecha/hora a ISO y añade folio y documento opcional. Los
animales salen de
\texorpdfstring{\softcode{operationAnimals}}{operationAnimals}, no de un
conteo. Los pesos individuales se incluyen únicamente cuando el usuario
los escribió. El total del lote no se divide entre animales para
inventar mediciones
\texorpdfstring{\softtext{individuales.}}{individuales.} Después del
POST, el componente conserva el ticket y puede ofrecer su aprobación
\texorpdfstring{\softtext{explícita.}}{explícita.}

\texorpdfstring{\softcode{useSalePositions(base, operation)}}{useSalePositions(base, operation)}
es un hook interno distinto de
\texorpdfstring{\softcode{useResource}}{useResource}. Su identidad de
consulta es la composición de posiciones serializada para una venta.
Mantiene datos, carga y error; depende de base e identidad, por lo que
no descarta identidades resueltas sólo porque una recarga transitoria
deje momentáneamente ausente el detalle. Su función interna
\texorpdfstring{\softcode{load}}{load} consulta inventario por páginas
de cien hasta resolver el conjunto requerido o alcanzar el total.
Después usa
\texorpdfstring{\softcode{operationAnimals}}{operationAnimals} para
comprobar que no faltó alguna posición. El efecto también cancela su
solicitud al cambiar ámbito o
\texorpdfstring{\softtext{desmontarse.}}{desmontarse.}

\texorpdfstring{\softcode{LotForm(props)}}{LotForm(props)} conserva el
número de lote \texorpdfstring{\softtext{seleccionado.}}{seleccionado.}
Si el estado es ACORDADA muestra condiciones aceptadas y un comando de
reapertura con ETag. Si es editable, permite elegir lotes, añadir el
siguiente número, montar
\texorpdfstring{\softcode{LotEditor}}{LotEditor} por clave y aceptar
condiciones mediante \texorpdfstring{\softcode{/acordar}}{/acordar}.
\texorpdfstring{\softcode{LotEditor}}{LotEditor} conserva búsqueda y
selección, consulta animales para compra o posiciones para venta y
deshabilita posiciones reservadas o no
\texorpdfstring{\softtext{disponibles.}}{disponibles.} Guardar envía
todas las condiciones a través de
\texorpdfstring{\softcode{replaceLot}}{replaceLot}, conservando el resto
de lotes.

Los campos incluyen unidad KG o CABEZA, precio con cuatro decimales,
pesaje base, deducción con seis decimales y
\texorpdfstring{\softtext{cargos/bonificaciones}}{cargos/bonificaciones}
con dos. La preferencia de recepción se obtiene del acuerdo actual y se
conserva al editar. Guardar un borrador KG sin ticket puede ser útil
para preparar el expediente, pero no autoriza aceptar ni confirmar un
acuerdo \texorpdfstring{\softtext{incompleto.}}{incompleto.} El servidor
exige evidencia aprobada en las transiciones
\texorpdfstring{\softtext{correspondientes.}}{correspondientes.} Reabrir
conserva versiones aceptadas y produce una propuesta nueva; no edita la
historia anterior.

La reserva de venta envía posiciones guardadas y vencimiento ISO.
Confirmar utiliza reconocimiento y ETag. La interfaz no manda el total
calculado. Una reserva vencida o ajena provoca rechazo y exige revisión
de \texorpdfstring{\softtext{disponibilidad.}}{disponibilidad.} Un botón
deshabilitado en la lista es una ayuda; el servidor debe seguir
rechazando dos confirmaciones competidoras aunque ambos usuarios hayan
visto una posición disponible antes de reservar.

\subsection{9.1 Historial, corrección y
comparación}\label{historial-correcciuxf3n-y-comparaciuxf3n}

\texorpdfstring{\softcode{TicketHistory(\{base, id, sede\})}}{TicketHistory(\{base, id, sede\})}
conserva revisión y ticket
\texorpdfstring{\softtext{seleccionado.}}{seleccionado.} Acepta una
respuesta paginada o un arreglo de tickets y construye una columna de
composición con animal e indicación de peso
\texorpdfstring{\softtext{individual.}}{individual.} Cuando no hubo
medición muestra «Sin medición
\texorpdfstring{\softtext{individual».}}{individual».} Presenta versión
corregida, motivo y fecha. Sólo ADMIN u OPERADOR obtiene acciones de
selección para corregir; otros perfiles siguen
\texorpdfstring{\softtext{consultando.}}{consultando.}

\texorpdfstring{\softcode{TicketCorrection(\{base, ticket, onDone\})}}{TicketCorrection(\{base, ticket, onDone\})}
tiene tres ramas. Una PROPUESTA puede aprobarse después de revisión. Un
ticket distinto de APROBADO conserva historia y no ofrece corrección
editable. Un APROBADO sin composición original completa muestra un aviso
y exige recargar; no toma animales del lote actual como reemplazo. Para
una corrección válida envía neto nuevo, motivo,
\texorpdfstring{\softcode{corrige\_a\_id}}{corrige\_a\_id}, número de
lote y exactamente los animales
\texorpdfstring{\softtext{originales.}}{originales.} La aprobación
posterior enlaza la propuesta y conserva el original como CORREGIDO.

\texorpdfstring{\softcode{TicketComparison(\{base, tickets\})}}{TicketComparison(\{base, tickets\})}
conserva primer y último identificador y ofrece dos
\texorpdfstring{\softtext{selectores.}}{selectores.} Cuando ambos
existen monta
\texorpdfstring{\softcode{ComparisonResult(\{base, first, last\})}}{ComparisonResult(\{base, first, last\})}.
Este componente consulta el endpoint canónico y muestra
\texorpdfstring{\softtext{NO\_COMPARABLE}}{NO\_COMPARABLE} cuando los
grupos difieren. Si recibe COMPARABLE presenta variación kg y porcentaje
devueltos como cadenas. No calcula la diferencia ni la convierte
automáticamente en descuento. Una tercera rama avisa si el servidor no
entregó comparación \texorpdfstring{\softtext{disponible.}}{disponible.}

Fuentes:
\texorpdfstring{\softcode{frontend/src/screens.tsx:8}}{frontend/src/screens.tsx:8},
\texorpdfstring{\softcode{frontend/src/screens.tsx:9}}{frontend/src/screens.tsx:9},
\texorpdfstring{\softcode{frontend/src/screens.tsx:10}}{frontend/src/screens.tsx:10},
\texorpdfstring{\softcode{frontend/src/screens.tsx:11}}{frontend/src/screens.tsx:11},
\texorpdfstring{\softcode{frontend/src/screens.tsx:16}}{frontend/src/screens.tsx:16},
\texorpdfstring{\softcode{frontend/src/screens.tsx:24}}{frontend/src/screens.tsx:24},
\texorpdfstring{\softcode{frontend/src/screens.tsx:29}}{frontend/src/screens.tsx:29},
\texorpdfstring{\softcode{frontend/src/screens.tsx:64}}{frontend/src/screens.tsx:64},
\texorpdfstring{\softcode{frontend/src/screens.tsx:68}}{frontend/src/screens.tsx:68},
\texorpdfstring{\softcode{frontend/src/screens.tsx:109}}{frontend/src/screens.tsx:109}.

\section{10. Liquidaciones, dinero y correcciones
económicas}\label{liquidaciones-dinero-y-correcciones-econuxf3micas}

\texorpdfstring{\softcode{Settlement(\{base, sede, id, go\})}}{Settlement(\{base, sede, id, go\})}
consulta una liquidación y conserva revisión. ADMIN o CAJA puede pasar a
registrar cobro o pago. El resumen muestra saldo, vencimiento y
dirección; el detalle separa operación económica, entrega física y
\texorpdfstring{\softtext{contraparte.}}{contraparte.} Conceptos, total
original, cargos, abonos, aplicaciones netas y saldo provienen del
servidor. El comprobante se abre como recurso autorizado y se identifica
como interno, no CFDI.

El formulario de ajuste sólo se presenta para ADMIN. Envía naturaleza,
monto canónico, motivo, fuente, reconocimiento y
\texorpdfstring{\softcode{reaplicar\_pagos}}{reaplicar\_pagos}. Esa
bandera solicita conciliación atómica cuando un abono reduce una
obligación pagada. No significa que React ejecute primero reversos y
después nuevas aplicaciones en varias llamadas. Una sola intención llega
al servidor, que conserva originales y registra los hechos necesarios
dentro de su \texorpdfstring{\softtext{transacción.}}{transacción.}

\texorpdfstring{\softcode{FinancialHistory(\{base, data, caja, onDone\})}}{FinancialHistory(\{base, data, caja, onDone\})}
muestra ajustes y
\texorpdfstring{\softtext{aplicaciones.}}{aplicaciones.} Conserva una
aplicación \texorpdfstring{\softtext{seleccionada.}}{seleccionada.}
Cuando el perfil puede operar caja y la fila no es ya un REVERSO, ofrece
\texorpdfstring{\softcode{/aplicaciones/<id>/reverso}}{/aplicaciones/<id>/reverso}
con motivo y
\texorpdfstring{\softtext{reconocimiento.}}{reconocimiento.} El original
sigue visible; el callback recarga datos. Esta acción independiente
sirve para otros casos y no reemplaza la conciliación atómica solicitada
por el ajuste guiado.

\texorpdfstring{\softcode{Payments(\{base, sede, context\})}}{Payments(\{base, sede, context\})}
conserva revisión, pago seleccionado, error de recarga, destino,
búsqueda aplicada y página de
\texorpdfstring{\softtext{obligaciones.}}{obligaciones.} El contexto de
una liquidación puede preseleccionar dirección, contraparte y UUID
buscada. La función interna
\texorpdfstring{\softcode{refreshPayment(id)}}{refreshPayment(id)}
recupera el pago completo mediante
\texorpdfstring{\softcode{request}}{request}, actualiza la selección y
aumenta revisión; \texorpdfstring{\softcode{bump}}{bump} la llama
después de validaciones o
\texorpdfstring{\softtext{aplicaciones.}}{aplicaciones.} Estas consultas
directas no añaden un
\texorpdfstring{\softcode{AbortController}}{AbortController} propio como
\texorpdfstring{\softcode{useResource}}{useResource}; el ámbito de la
pantalla sigue estando separado por el montaje del
\texorpdfstring{\softtext{contenedor.}}{contenedor.}

La lista de pagos filtra estado REGISTRADO o VALIDADO y tipo COBRO o
PAGO. Registrar dinero envía importe, contraparte, método, referencia y
fecha. No exige en ese formulario el mismo reconocimiento que validar.
Un pago REGISTRADO ofrece validación con periodo y fecha; uno VALIDADO
permite aplicar. La lista de obligaciones consulta veinte por página,
incorpora búsqueda y muestra candidatos con saldo positivo, misma
contraparte y dirección
\texorpdfstring{\softtext{compatible.}}{compatible.} Ese filtrado ayuda
a seleccionar, pero el servidor vuelve a validar empresa, sede, moneda,
dirección y \texorpdfstring{\softtext{disponibilidad.}}{disponibilidad.}

Una aplicación envía un arreglo con liquidación e importe, además de
\texorpdfstring{\softtext{reconocimiento.}}{reconocimiento.} El destino
sólo aparece si todavía está dentro del conjunto compatible visible. No
se descuenta saldo localmente al pulsar el botón: después de éxito se
recupera el pago y las
\texorpdfstring{\softtext{obligaciones.}}{obligaciones.}
\texorpdfstring{\softcode{RefundForm(\{base, payment, onDone\})}}{RefundForm(\{base, payment, onDone\})}
permite documentar salida del disponible con importe, método, fecha,
referencia y
\texorpdfstring{\softtext{reconocimiento.}}{reconocimiento.} Conserva el
pago original y utiliza otra intención y otra clave.

Tres magnitudes deben leerse por separado: obligación pendiente,
disponible de un pago y resultado económico de una venta. Un anticipo
puede existir sin liquidación saldada; una entrega física puede estar
pendiente con saldo cero; un margen puede incorporar gastos que no
aumentan la deuda del comprador. Las fórmulas siguientes describen
categorías canónicas, no operaciones ejecutadas por la interfaz:

\[
\text{Saldo de obligación}=\text{Total original}+\text{Cargos}-\text{Abonos}-\text{Aplicaciones netas}
\]

\[
\text{Disponible del pago}=\text{Importe validado}-\text{Aplicaciones netas}-\text{Reembolsos}
\]

La UI recibe cada resultado de la consulta
\texorpdfstring{\softtext{autorizada.}}{autorizada.} Si una cifra falta,
muestra ausencia. No deriva otra cifra para hacer que el desglose
parezca completo.

Fuentes:
\texorpdfstring{\softcode{frontend/src/screens.tsx:33}}{frontend/src/screens.tsx:33},
\texorpdfstring{\softcode{frontend/src/screens.tsx:36}}{frontend/src/screens.tsx:36},
\texorpdfstring{\softcode{frontend/src/screens.tsx:38}}{frontend/src/screens.tsx:38},
\texorpdfstring{\softcode{frontend/src/screens.tsx:97}}{frontend/src/screens.tsx:97},
\texorpdfstring{\softcode{frontend/src/screens.tsx:98}}{frontend/src/screens.tsx:98}.

\section{11. Gastos, tarifas aceptadas y
periodos}\label{gastos-tarifas-aceptadas-y-periodos}

\texorpdfstring{\softcode{ExpenseCapture(\{base, onDone\})}}{ExpenseCapture(\{base, onDone\})}
conserva naturaleza COSTO o BONIFICACION y snapshot aceptado. Renderiza
primero \texorpdfstring{\softcode{ProviderTariffs}}{ProviderTariffs} y
después el \texorpdfstring{\softtext{formulario.}}{formulario.} Sin
snapshot, el usuario captura importe mediante
\texorpdfstring{\softcode{cash}}{cash}; con snapshot, se presenta la
tarifa aceptada y el cuerpo utiliza el
\texorpdfstring{\softcode{precio}}{precio} recibido y su
\texorpdfstring{\softcode{snapshot\_id}}{snapshot\_id}. La UI no permite
escribir una procedencia externa arbitraria para hacerla pasar por una
aceptación válida.

El gasto incluye operación, clasificación CAPITALIZABLE,
\texorpdfstring{\softtext{DIRECTO\_VENTA}}{DIRECTO\_VENTA} u OTRO,
concepto, fecha económica declarada y reconocimiento
\texorpdfstring{\softtext{explícito.}}{explícito.} Proveedor, referencia
de origen y documento son
\texorpdfstring{\softtext{opcionales.}}{opcionales.} Si se elige
BONIFICACION, aparece el gasto original que se bonifica. Cambiar la
naturaleza cambia campos visibles, pero no elimina el hecho original. La
composición del cuerpo y la autorización se vuelven a comprobar en SQL.
Un gasto directo de venta puede afectar margen sin incrementar la
obligación con la \texorpdfstring{\softtext{contraparte.}}{contraparte.}

\texorpdfstring{\softcode{ProviderTariffs(\{base, onAccepted\})}}{ProviderTariffs(\{base, onAccepted\})}
es interno. Conserva snapshot seleccionado y puede inicializarlo desde
un intento de aceptación pendiente. Consulta cien tarifas, encuentra la
fila seleccionada y adapta errores de red o 503 a un mensaje contextual
del proveedor. No afirma que la sede o ambos nodos estén caídos. El
select expone ruta, precio y versión; una ficha muestra fuente, unidad y
vigencia. La aceptación envía sólo
\texorpdfstring{\softcode{snapshot\_id}}{snapshot\_id}. Al confirmar,
comunica al padre la fila consultada
\texorpdfstring{\softtext{correspondiente.}}{correspondiente.} Cambiar
selección elimina el snapshot aceptado del padre hasta aceptar la nueva
versión.

La aceptación y el gasto son comandos distintos. Consultar una tarifa
tampoco equivale a aceptarla. El backend lee la fuente mediante el
lector permitido y captura el snapshot en el propietario de la sede; esa
topología no se controla desde un select. La prueba de UI identifica la
fuente como DEMO simulada. El lector no debe concluir que existe un
convenio con un proveedor comercial ni una consulta sanitaria oficial.

\texorpdfstring{\softcode{ExpenseDetail(\{base, sede, id\})}}{ExpenseDetail(\{base, sede, id\})}
conserva revisión y método. Solicita una previsualización de reparto.
Para POR\_CABEZA o
\texorpdfstring{\softtext{POR\_PESO\_MEDIDO}}{POR\_PESO\_MEDIDO} muestra
bases e importes \texorpdfstring{\softtext{canónicos.}}{canónicos.} Para
\texorpdfstring{\softtext{IMPORTE\_INDIVIDUAL}}{IMPORTE\_INDIVIDUAL}
utiliza la consulta por cabeza para conocer destinos, pero oculta sus
importes calculados y pide una captura individual por posición. No
ofrece esa previsualización como aprobación de los valores manuales. El
servidor exige composición y suma exacta del gasto.

El POST de reparto incluye método y reconocimiento; en modo individual
agrega \texorpdfstring{\softcode{asignaciones}}{asignaciones} con
posición y dinero. Un constructor rechaza intentar repartir
individualmente sin destinos
\texorpdfstring{\softtext{consultados.}}{consultados.} Cambiar método
remonta el formulario mediante clave, evitando reutilizar campos de otro
método. \texorpdfstring{\softcode{onDone}}{onDone} aumenta revisión. La
UI no ajusta la última posición con una diferencia calculada localmente
para «cuadrar» centavos; esa política corresponde al cálculo canónico.

\texorpdfstring{\softcode{Periods(\{base, sede\})}}{Periods(\{base, sede\})}
conserva revisión y fila
\texorpdfstring{\softtext{seleccionada.}}{seleccionada.} Todos consultan
la lista; ADMIN puede seleccionar para
\texorpdfstring{\softtext{administrar.}}{administrar.}
\texorpdfstring{\softcode{PeriodVersions(\{base, id\})}}{PeriodVersions(\{base, id\})}
consulta versiones inmutables y admite respuesta paginada o arreglo.
Muestra número, tipo, motivo y corte. El formulario elige
\texorpdfstring{\softcode{/cierre}}{/cierre} cuando el periodo está
ABIERTO y \texorpdfstring{\softcode{/reabrir}}{/reabrir} en el otro
caso, siempre con motivo. Después de éxito elimina selección y recarga.

El alta solicita primer y último día de un mes. El texto explica
continuidad y un único periodo abierto, pero la UI no reemplaza la
validación de orden, límites y conciliación del servidor. La prueba real
cerró y reabrió el periodo corriente y conservó versiones. El siguiente
mes se comprobó en QA aislado; no se adelantó por iniciativa de la
prueba el calendario de la demostración principal.

Fuentes:
\texorpdfstring{\softcode{frontend/src/screens.tsx:83}}{frontend/src/screens.tsx:83},
\texorpdfstring{\softcode{frontend/src/screens.tsx:87}}{frontend/src/screens.tsx:87},
\texorpdfstring{\softcode{frontend/src/screens.tsx:88}}{frontend/src/screens.tsx:88},
\texorpdfstring{\softcode{frontend/src/screens.tsx:92}}{frontend/src/screens.tsx:92},
\texorpdfstring{\softcode{frontend/src/screens.tsx:96}}{frontend/src/screens.tsx:96}.

\section{12. Archivos privados y usuarios
sensibles}\label{archivos-privados-y-usuarios-sensibles}

\texorpdfstring{\softcode{Documents(\{base, sede\})}}{Documents(\{base, sede\})}
conserva revisión y documento
\texorpdfstring{\softtext{seleccionado.}}{seleccionado.} Presenta lista
con tipo, formato y versión; ADMIN u OPERADOR obtiene carga. Elegir una
fila monta
\texorpdfstring{\softcode{DocumentDetail(\{base, sede, id\})}}{DocumentDetail(\{base, sede, id\})}
por \texorpdfstring{\softtext{identificador.}}{identificador.} Ese
detalle consulta metadatos, muestra si el contenido es sintético, ofrece
descarga autorizada y muestra
\texorpdfstring{\softtext{revisiones.}}{revisiones.} ADMIN puede
registrar una revisión con operación, criterio, resultado y documento
exacto. El callback vuelve a consultar el detalle.

\texorpdfstring{\softcode{DocumentUpload(\{base, onDone\})}}{DocumentUpload(\{base, onDone\})}
necesita un tratamiento especial porque
\texorpdfstring{\softcode{File}}{File} y bytes no deben guardarse como
un JSON ordinario. Conserva intento, carga, error, éxito y una
referencia \texorpdfstring{\softcode{upload}}{upload} al
\texorpdfstring{\softcode{FormData}}{FormData} en memoria. Su función
interna \texorpdfstring{\softcode{submit(form)}}{submit(form)} consulta
primero el resultado si ya existe comando pendiente. Si encuentra
COMPLETADA, termina sin subir otra vez. Si necesita enviar, verifica que
haya archivo, tamaño máximo veinte MiB y MIME PDF, JPEG o PNG.

Calcula SHA-256 con
\texorpdfstring{\softcode{crypto.subtle.digest}}{crypto.subtle.digest}
sobre los bytes y forma una huella con nombre, tamaño, hash, tipo
documental y operación vinculada. La huella se guarda como cuerpo del
\texorpdfstring{\softcode{Attempt}}{Attempt}; el archivo real permanece
en memoria. Después de recargar, seleccionar un archivo diferente o
cambiar \texorpdfstring{\softtext{tipo/vínculo}}{tipo/vínculo} no
coincide con la intención antigua y produce error. La nueva selección
debe reproducir contenido y metadatos del mismo envío. Cuando la
consulta de resultado confirma la subida, no se necesita recapturar
bytes.

El POST real usa \texorpdfstring{\softcode{FormData}}{FormData}, no la
huella JSON. \texorpdfstring{\softcode{request}}{request} no impone
\texorpdfstring{\softcode{Content-Type: application/json}}{Content-Type: application/json}
porque el cuerpo no es cadena; el navegador construye el límite
multipart. Se envía la misma clave de
\texorpdfstring{\softtext{idempotencia.}}{idempotencia.} Después de
éxito borra pendiente y referencia de memoria. En un rechazo conocido
permite corregir; ante resultado incierto conserva lo necesario para
recuperar. El atributo \texorpdfstring{\softcode{accept}}{accept} ayuda
a elegir archivo, pero las verificaciones del cliente y del servidor
siguen siendo \texorpdfstring{\softtext{necesarias.}}{necesarias.}

\texorpdfstring{\softcode{Users(\{base, sede\})}}{Users(\{base, sede\})}
conserva revisión y usuario
\texorpdfstring{\softtext{seleccionado.}}{seleccionado.} Una lista
muestra nombre, login, roles y estado. ADMIN obtiene alta sensible,
PATCH de \texorpdfstring{\softtext{nombre/estado}}{nombre/estado} y
membresía por sede. El PATCH permite ACTIVO o BLOQUEADO. La membresía
envía un código ADMIN, OPERADOR, CAJA o AUDITOR y un booleano
\texorpdfstring{\softcode{vigente}}{vigente}; retirar un rol es un
comando explícito, no borrar visualmente una etiqueta. La autorización
vigente se comprueba también en el backend.

\texorpdfstring{\softcode{SensitiveUserCreate(\{base, onDone\})}}{SensitiveUserCreate(\{base, onDone\})}
conserva intento, bandera de espera, error, éxito y una referencia de
texto \texorpdfstring{\softcode{body}}{body}. Su
\texorpdfstring{\softcode{submit}}{submit} consulta el resultado
pendiente antes de intentar crear de nuevo. Construye JSON con login,
nombre y contraseña sólo en memoria; exige contraseña no vacía; obtiene
SHA-256 del contenido completo y persiste exclusivamente la huella con
la clave. Si se recarga y el comando no terminó, el mismo contenido debe
recapturarse para coincidir con la huella. La contraseña nunca se copia
a \texorpdfstring{\softcode{sessionStorage}}{sessionStorage}.

El POST utiliza el JSON original en memoria. Después de éxito elimina
referencia y pendiente, informa al contenedor y solicita reset del
\texorpdfstring{\softtext{formulario.}}{formulario.} El callback
incorpora nombre y estado ACTIVO para presentar inmediatamente el
usuario recién creado, mientras que la lista puede consultarse otra vez.
Crear un usuario no concede automáticamente todos los roles de todas las
sedes. Por ello la prueba de AUDITOR realiza una membresía separada
antes de iniciar su nueva sesión.

La protección mediante huella no convierte este mecanismo en un almacén
de contraseñas ni en un servicio de autenticación offline. La sesión y
el hash de contraseña del servidor siguen perteneciendo a otra capa. La
prueba unitaria y la prueba UI comprueban ausencia de contraseña
persistida; no se presenta ese resultado como una auditoría completa de
\texorpdfstring{\softtext{penetración.}}{penetración.}

Fuentes:
\texorpdfstring{\softcode{frontend/src/screens.tsx:46}}{frontend/src/screens.tsx:46},
\texorpdfstring{\softcode{frontend/src/screens.tsx:47}}{frontend/src/screens.tsx:47},
\texorpdfstring{\softcode{frontend/src/screens.tsx:49}}{frontend/src/screens.tsx:49},
\texorpdfstring{\softcode{frontend/src/screens.tsx:51}}{frontend/src/screens.tsx:51},
\texorpdfstring{\softcode{frontend/src/screens.tsx:99}}{frontend/src/screens.tsx:99},
\texorpdfstring{\softcode{frontend/src/screens.tsx:100}}{frontend/src/screens.tsx:100},
\texorpdfstring{\softcode{frontend/src/screens.tsx:102}}{frontend/src/screens.tsx:102}.

\section{13. Expediente, movimientos, reportes y
simulador}\label{expediente-movimientos-reportes-y-simulador}

\texorpdfstring{\softcode{CaseFile(\{base, id\})}}{CaseFile(\{base, id\})}
solicita
\texorpdfstring{\softcode{/operaciones/<id>/expediente}}{/operaciones/<id>/expediente}.
Admite eventos o hechos en la estructura consultada y extrae documentos
\texorpdfstring{\softtext{vinculados.}}{vinculados.} Representa fechas,
tipo o acción y descripción o estado en una línea temporal. Cuando
faltan eventos lo declara; cuando faltan documentos no inventa una
evidencia. Sus enlaces descargan la versión autorizada del mismo origen.
El botón de actualizar utiliza
\texorpdfstring{\softcode{reload}}{reload}, de modo que la historia
permanece sustentada por consultas, no por una narración construida de
clics locales.

\texorpdfstring{\softcode{MovementCapture(\{base, sede\})}}{MovementCapture(\{base, sede\})}
conserva tipo RECEPCION
\texorpdfstring{\softtext{inicialmente.}}{inicialmente.} El usuario
captura posiciones separadas por espacios o comas; el constructor
produce un arreglo no vacío de textos y convierte fecha/hora local a
ISO. ENTREGA exige identificador de la venta
\texorpdfstring{\softtext{confirmada.}}{confirmada.} BAJA aparece sólo
para ADMIN y agrega motivo PERDIDA o MUERTE, documento y
\texorpdfstring{\softtext{reconocimiento.}}{reconocimiento.} RECEPCION y
ENTREGA no incluyen reconocimiento económico nuevo en ese cuerpo: son
hechos físicos distintos del cierre ya
\texorpdfstring{\softtext{confirmado.}}{confirmado.}

\texorpdfstring{\softcode{ControlView(\{base, kind\})}}{ControlView(\{base, kind\})}
consulta usuarios o sincronización, pero la navegación actual emplea
\texorpdfstring{\softcode{Users}}{Users} para administración y
\texorpdfstring{\softcode{ControlView}}{ControlView} para
\texorpdfstring{\softtext{sincronización.}}{sincronización.} Su rama de
usuarios es una representación de consulta reutilizable, no un segundo
CRUD completo. La rama de sincronización muestra estado, última
actualización y avisos cuando llegan. No presenta automáticamente todos
los campos internos del diagnóstico ni nombres de hosts. Un valor
ausente se mantiene como ausencia y una etiqueta ASINCRONA no se
reemplaza por «todo sincronizado» por
\texorpdfstring{\softtext{conveniencia.}}{conveniencia.}

\texorpdfstring{\softcode{Simulator(\{base, id\})}}{Simulator(\{base, id\})}
consulta el cálculo de una operación, presenta total y conceptos y
etiqueta «No registra hechos». Si hay faltantes recibidos los muestra.
Un borrador KG incompleto puede producir error y no un total artificial
de cero. Consultar el simulador no genera inventario ni obligación, no
acepta condiciones y no confirma. Su resultado sirve para revisar una
propuesta, no para saltarse la acción económica
\texorpdfstring{\softtext{explícita.}}{explícita.}

\texorpdfstring{\softcode{Report(\{base, kind, go\})}}{Report(\{base, kind, go\})}
conserva página, texto escrito, búsqueda aplicada y filtro de
\texorpdfstring{\softtext{vencimiento.}}{vencimiento.} Para cartera
construye query con tamaño veinte y booleano textual
\texorpdfstring{\softcode{true}}{true} o
\texorpdfstring{\softcode{false}}{false} en
\texorpdfstring{\softcode{vencida}}{vencida}; el backend normaliza la
entrada. Muestra \texorpdfstring{\softcode{filas}}{filas}, total y
controles de página. Elegir una fila navega hacia la
\texorpdfstring{\softtext{liquidación.}}{liquidación.} Para margen no
agrega paginación local, porque recibe un agregado canónico de todo el
ámbito \texorpdfstring{\softtext{consultado.}}{consultado.}

Ambos reportes muestran corte, etiqueta de consolidación, fuente y
antigüedad cuando la API los informa. Una antigüedad nula se explica
expresamente; no se transforma en cero. Las métricas de margen son venta
neta, costo vendido, gastos directos, margen y pérdidas. El componente
llama a \texorpdfstring{\softcode{money}}{money} para mostrarlas, pero
no suma o resta esos campos. El corte identifica la consulta, no promete
una consistencia global instantánea entre todos los streams de
\texorpdfstring{\softtext{replicación.}}{replicación.}

Los enlaces CSV de cartera, margen e inventario apuntan a
\texorpdfstring{\softcode{/api/v1/sedes/<sede>/reportes/<tipo>/csv}}{/api/v1/sedes/<sede>/reportes/<tipo>/csv}.
Dicen «Exportar CSV de toda la sede». No llevan la búsqueda o filtros de
la tabla. Un lector que buscó una liquidación y luego descarga debe
esperar toda la sede autorizada, no únicamente la fila visible. La
descarga usa cookie y autorización del servidor. No se llama a un
servicio externo para producir el archivo ni se convierte dinero a
números binarios para escribirlo en el cliente.

Fuentes:
\texorpdfstring{\softcode{frontend/src/screens.tsx:41}}{frontend/src/screens.tsx:41},
\texorpdfstring{\softcode{frontend/src/screens.tsx:75}}{frontend/src/screens.tsx:75},
\texorpdfstring{\softcode{frontend/src/screens.tsx:79}}{frontend/src/screens.tsx:79},
\texorpdfstring{\softcode{frontend/src/screens.tsx:80}}{frontend/src/screens.tsx:80},
\texorpdfstring{\softcode{frontend/src/screens.tsx:82}}{frontend/src/screens.tsx:82},
\texorpdfstring{\softcode{frontend/src/App.tsx:41}}{frontend/src/App.tsx:41}.

\section{14. Recorrido completo con cuerpos de solicitud
didácticos}\label{recorrido-completo-con-cuerpos-de-solicitud-diduxe1cticos}

Este recorrido combina las acciones de formularios con la forma de sus
mensajes. Los UUID siguientes son marcadores sintácticamente válidos
para estudiar el contrato; no son instrucciones para crear registros
sobre una sede \texorpdfstring{\softtext{desconocida.}}{desconocida.} En
una ejecución real se toman identificadores de las respuestas y se
obtiene la versión mediante GET. Se supone una sede autorizada, una
contraparte DEMO, dos animales, un periodo abierto y roles suficientes
para operación y caja. La ilustración conserva el caso de dos cabezas
a500.0000 MXN: el servidor
\texorpdfstring{\softtext{devolvió1000.00}}{devolvió1000.00} en el
recorrido \texorpdfstring{\softtext{evidenciado.}}{evidenciado.}

\subsection{14.1 Preparar identidad y crear
borrador}\label{preparar-identidad-y-crear-borrador}

Primero se inicia sesión sin copiar sus secretos al capítulo, se
selecciona sede y se crea o consulta una
\texorpdfstring{\softtext{contraparte.}}{contraparte.} El cuerpo
comercial de alta puede ser:

\begin{Shaded}
\begin{Highlighting}[]
\FunctionTok{\{}\DataTypeTok{"nombre"}\FunctionTok{:}\StringTok{"DEMO Contraparte didáctica"}\FunctionTok{,}\DataTypeTok{"tipo\_persona"}\FunctionTok{:}\StringTok{"FISICA"}\FunctionTok{\}}
\end{Highlighting}
\end{Shaded}

Cada identidad animal necesita su propia alta; un ejemplo de entrada es:

\begin{Shaded}
\begin{Highlighting}[]
\FunctionTok{\{}\DataTypeTok{"sexo"}\FunctionTok{:}\StringTok{"M"}\FunctionTok{,}\DataTypeTok{"identificaciones"}\FunctionTok{:}\OtherTok{[}\FunctionTok{\{}\DataTypeTok{"sistema"}\FunctionTok{:}\StringTok{"INTERNO"}\FunctionTok{,}\DataTypeTok{"codigo"}\FunctionTok{:}\StringTok{"DEMO{-}MANUAL{-}001"}\FunctionTok{\}}\OtherTok{]}\FunctionTok{\}}
\end{Highlighting}
\end{Shaded}

Después, el formulario de Nueva operación envía a
\texorpdfstring{\softcode{/sedes/<S>/operaciones}}{/sedes/<S>/operaciones}:

\begin{Shaded}
\begin{Highlighting}[]
\FunctionTok{\{}
  \DataTypeTok{"tipo"}\FunctionTok{:}\StringTok{"COMPRA"}\FunctionTok{,}
  \DataTypeTok{"parte\_id"}\FunctionTok{:}\StringTok{"10000000{-}0000{-}4000{-}8000{-}000000000001"}\FunctionTok{,}
  \DataTypeTok{"fecha\_operacion"}\FunctionTok{:}\StringTok{"2026{-}10{-}11"}\FunctionTok{,}
  \DataTypeTok{"requiere\_recepcion"}\FunctionTok{:}\KeywordTok{true}
\FunctionTok{\}}
\end{Highlighting}
\end{Shaded}

Crear el borrador devuelve su identificador, pero no significa que haya
posiciones ni deuda. Abrir ese identificador produce GET del detalle y
su ETag. La edición usa la versión recibida, no presupone que un POST de
creación haya entregado la cabecera. Los tests de fixtures también
consultan GET entre transiciones por esta razón.

\subsection{14.2 Guardar condiciones, aceptar y
confirmar}\label{guardar-condiciones-aceptar-y-confirmar}

El PATCH a
\texorpdfstring{\softcode{/sedes/<S>/operaciones/<OC>}}{/sedes/<S>/operaciones/<OC>}
conserva un arreglo ordenado de lotes. Para dos animales, su forma
simplificada es:

\begin{Shaded}
\begin{Highlighting}[]
\FunctionTok{\{}
  \DataTypeTok{"lotes"}\FunctionTok{:}\OtherTok{[}\FunctionTok{\{}
    \DataTypeTok{"numero\_lote"}\FunctionTok{:}\DecValTok{1}\FunctionTok{,}
    \DataTypeTok{"animales"}\FunctionTok{:}\OtherTok{[}
      \StringTok{"20000000{-}0000{-}4000{-}8000{-}000000000001"}\OtherTok{,}
      \StringTok{"20000000{-}0000{-}4000{-}8000{-}000000000002"}
    \OtherTok{]}\FunctionTok{,}
    \DataTypeTok{"unidad\_precio"}\FunctionTok{:}\StringTok{"CABEZA"}\FunctionTok{,}
    \DataTypeTok{"precio\_unitario"}\FunctionTok{:}\StringTok{"500.0000"}
  \FunctionTok{\}}\OtherTok{]}\FunctionTok{,}
  \DataTypeTok{"requiere\_recepcion"}\FunctionTok{:}\KeywordTok{true}
\FunctionTok{\}}
\end{Highlighting}
\end{Shaded}

Lo acompañan CSRF,
\texorpdfstring{\softcode{Idempotency-Key}}{Idempotency-Key} del intento
e \texorpdfstring{\softcode{If-Match}}{If-Match} con el ETag obtenido.
El servidor guarda la propuesta. Luego «Aceptar condiciones» envía
\texorpdfstring{\softcode{\{\}}}{\{\}} a
\texorpdfstring{\softcode{/acordar}}{/acordar} con versión vigente. Si
el acuerdo está aceptado y se desea modificarlo, se utiliza
\texorpdfstring{\softcode{/reabrir}}{/reabrir}, se guarda otra propuesta
y se acepta otra versión; no se edita directamente la versión anterior.

Confirmar compra envía únicamente su reconocimiento económico:

\begin{Shaded}
\begin{Highlighting}[]
\FunctionTok{\{}
  \DataTypeTok{"periodo\_id"}\FunctionTok{:}\StringTok{"30000000{-}0000{-}4000{-}8000{-}000000000001"}\FunctionTok{,}
  \DataTypeTok{"fecha\_reconocimiento"}\FunctionTok{:}\StringTok{"2026{-}10{-}11"}
\FunctionTok{\}}
\end{Highlighting}
\end{Shaded}

El total esperado de esta ilustración es1000.00, pero ese campo no está
en el cuerpo de \texorpdfstring{\softtext{confirmación.}}{confirmación.}
El servidor devuelve una liquidación cuyo
\texorpdfstring{\softcode{total}}{total} y
\texorpdfstring{\softcode{saldo}}{saldo} usan dinero canónico. La
navegación puede tomar \texorpdfstring{\softcode{data.id}}{data.id} como
identificador de la liquidación; no debe conservarlo erróneamente como
UUID de \texorpdfstring{\softtext{operación.}}{operación.} El inventario
resultante será pendiente de recepción porque así se pactó.

\subsection{14.3 Recibir físicamente y preparar
venta}\label{recibir-fuxedsicamente-y-preparar-venta}

En Inventario se buscan los animales y se copian los identificadores de
posición devueltos. Recibir utiliza otro comando:

\begin{Shaded}
\begin{Highlighting}[]
\FunctionTok{\{}
  \DataTypeTok{"tipo"}\FunctionTok{:}\StringTok{"RECEPCION"}\FunctionTok{,}
  \DataTypeTok{"posiciones"}\FunctionTok{:}\OtherTok{[}
    \StringTok{"40000000{-}0000{-}4000{-}8000{-}000000000001"}\OtherTok{,}
    \StringTok{"40000000{-}0000{-}4000{-}8000{-}000000000002"}
  \OtherTok{]}\FunctionTok{,}
  \DataTypeTok{"ocurrido\_en"}\FunctionTok{:}\StringTok{"2026{-}10{-}11T16:00:00.000Z"}
\FunctionTok{\}}
\end{Highlighting}
\end{Shaded}

No incluye pago ni un segundo total de compra. Después de confirmarlo se
consulta la \texorpdfstring{\softtext{disponibilidad.}}{disponibilidad.}
Para vender se crea otro borrador con tipo VENTA y contraparte
\texorpdfstring{\softtext{autorizada.}}{autorizada.} Su PATCH selecciona
posiciones, no identidades sin ciclo:

\begin{Shaded}
\begin{Highlighting}[]
\FunctionTok{\{}
  \DataTypeTok{"lotes"}\FunctionTok{:}\OtherTok{[}\FunctionTok{\{}
    \DataTypeTok{"numero\_lote"}\FunctionTok{:}\DecValTok{1}\FunctionTok{,}
    \DataTypeTok{"posiciones"}\FunctionTok{:}\OtherTok{[}
      \StringTok{"40000000{-}0000{-}4000{-}8000{-}000000000001"}\OtherTok{,}
      \StringTok{"40000000{-}0000{-}4000{-}8000{-}000000000002"}
    \OtherTok{]}\FunctionTok{,}
    \DataTypeTok{"unidad\_precio"}\FunctionTok{:}\StringTok{"CABEZA"}\FunctionTok{,}
    \DataTypeTok{"precio\_unitario"}\FunctionTok{:}\StringTok{"500.0000"}
  \FunctionTok{\}}\OtherTok{]}\FunctionTok{,}
  \DataTypeTok{"requiere\_recepcion"}\FunctionTok{:}\KeywordTok{false}
\FunctionTok{\}}
\end{Highlighting}
\end{Shaded}

Se aceptan condiciones y se reserva el conjunto antes de confirmar. El
vencimiento siempre debe ser futuro para la intención actual:

\begin{Shaded}
\begin{Highlighting}[]
\FunctionTok{\{}
  \DataTypeTok{"posiciones"}\FunctionTok{:}\OtherTok{[}
    \StringTok{"40000000{-}0000{-}4000{-}8000{-}000000000001"}\OtherTok{,}
    \StringTok{"40000000{-}0000{-}4000{-}8000{-}000000000002"}
  \OtherTok{]}\FunctionTok{,}
  \DataTypeTok{"vence\_en"}\FunctionTok{:}\StringTok{"2026{-}10{-}11T20:00:00.000Z"}
\FunctionTok{\}}
\end{Highlighting}
\end{Shaded}

La fecha es didáctica, no válida para ejecutar el ejemplo en otro día.
Si venció, no se inventa una extensión en el navegador. El usuario
revisa disponibilidad y captura una nueva reserva conforme a la
respuesta. Confirmar venta vuelve a usar periodo y fecha explícitos,
versión vigente y clave propia. El servidor decide inventario y
obligación en la misma
\texorpdfstring{\softtext{transacción.}}{transacción.}

\subsection{14.4 Variante por kg}\label{variante-por-kg}

Cuando el precio es KG, guardar puede preparar la composición antes de
tener ticket, pero aceptar exige base aprobada. El ticket enviado
conserva identidades animales, aunque la venta se haya seleccionado por
posiciones:

\begin{Shaded}
\begin{Highlighting}[]
\FunctionTok{\{}
  \DataTypeTok{"numero\_lote"}\FunctionTok{:}\DecValTok{1}\FunctionTok{,}
  \DataTypeTok{"modo"}\FunctionTok{:}\StringTok{"NETO\_REPORTADO"}\FunctionTok{,}
  \DataTypeTok{"neto"}\FunctionTok{:\{}\DataTypeTok{"kg"}\FunctionTok{:}\StringTok{"300.000"}\FunctionTok{\},}
  \DataTypeTok{"ticket\_folio"}\FunctionTok{:}\StringTok{"DEMO{-}MANUAL{-}PESO"}\FunctionTok{,}
  \DataTypeTok{"pesado\_en"}\FunctionTok{:}\StringTok{"2026{-}10{-}11T17:00:00.000Z"}\FunctionTok{,}
  \DataTypeTok{"animales"}\FunctionTok{:}\OtherTok{[}
    \FunctionTok{\{}\DataTypeTok{"animal\_id"}\FunctionTok{:}\StringTok{"20000000{-}0000{-}4000{-}8000{-}000000000001"}\FunctionTok{,}\DataTypeTok{"peso\_individual"}\FunctionTok{:\{}\DataTypeTok{"kg"}\FunctionTok{:}\StringTok{"100.000"}\FunctionTok{\}\}}\OtherTok{,}
    \FunctionTok{\{}\DataTypeTok{"animal\_id"}\FunctionTok{:}\StringTok{"20000000{-}0000{-}4000{-}8000{-}000000000002"}\FunctionTok{,}\DataTypeTok{"peso\_individual"}\FunctionTok{:\{}\DataTypeTok{"kg"}\FunctionTok{:}\StringTok{"200.000"}\FunctionTok{\}\}}
  \OtherTok{]}
\FunctionTok{\}}
\end{Highlighting}
\end{Shaded}

Los pesos individuales sólo se incluyen si existen
\texorpdfstring{\softtext{mediciones.}}{mediciones.} Una propuesta se
aprueba mediante su endpoint y se vincula a la
\texorpdfstring{\softtext{condición.}}{condición.} Con
\texorpdfstring{\softtext{precio10.0000}}{precio10.0000} y
\texorpdfstring{\softtext{neto300.000,}}{neto300.000,} el servidor
\texorpdfstring{\softtext{devolvió3000.00}}{devolvió3000.00} en la
regresión real. La igualdad didáctica describe ese ejemplo, no una
implementación \texorpdfstring{\softtext{JavaScript:}}{JavaScript:}

\[
300.000\ \mathrm{kg}\times10.0000\ \mathrm{MXN/kg}=3000.00\ \mathrm{MXN}
\]

\subsection{14.5 Registrar, validar y aplicar
dinero}\label{registrar-validar-y-aplicar-dinero}

Desde liquidación se abre Cobros y pagos con contexto de dirección y
\texorpdfstring{\softtext{contraparte.}}{contraparte.} Registrar
utiliza:

\begin{Shaded}
\begin{Highlighting}[]
\FunctionTok{\{}
  \DataTypeTok{"parte\_id"}\FunctionTok{:}\StringTok{"10000000{-}0000{-}4000{-}8000{-}000000000001"}\FunctionTok{,}
  \DataTypeTok{"tipo"}\FunctionTok{:}\StringTok{"COBRO"}\FunctionTok{,}
  \DataTypeTok{"importe"}\FunctionTok{:\{}\DataTypeTok{"importe"}\FunctionTok{:}\StringTok{"1000.00"}\FunctionTok{,}\DataTypeTok{"moneda"}\FunctionTok{:}\StringTok{"MXN"}\FunctionTok{\},}
  \DataTypeTok{"metodo"}\FunctionTok{:}\StringTok{"EFECTIVO"}\FunctionTok{,}
  \DataTypeTok{"referencia"}\FunctionTok{:}\StringTok{"DEMO Manual"}\FunctionTok{,}
  \DataTypeTok{"recibido\_el"}\FunctionTok{:}\StringTok{"2026{-}10{-}11"}
\FunctionTok{\}}
\end{Highlighting}
\end{Shaded}

La creación de pago no salda automáticamente la
\texorpdfstring{\softtext{liquidación.}}{liquidación.} Se valida con
reconocimiento y se elige una obligación
\texorpdfstring{\softtext{compatible.}}{compatible.} Aplicar usa un
cuerpo diferente:

\begin{Shaded}
\begin{Highlighting}[]
\FunctionTok{\{}
  \DataTypeTok{"periodo\_id"}\FunctionTok{:}\StringTok{"30000000{-}0000{-}4000{-}8000{-}000000000001"}\FunctionTok{,}
  \DataTypeTok{"fecha\_reconocimiento"}\FunctionTok{:}\StringTok{"2026{-}10{-}11"}\FunctionTok{,}
  \DataTypeTok{"aplicaciones"}\FunctionTok{:}\OtherTok{[}\FunctionTok{\{}
    \DataTypeTok{"liquidacion\_id"}\FunctionTok{:}\StringTok{"50000000{-}0000{-}4000{-}8000{-}000000000001"}\FunctionTok{,}
    \DataTypeTok{"importe"}\FunctionTok{:\{}\DataTypeTok{"importe"}\FunctionTok{:}\StringTok{"1000.00"}\FunctionTok{,}\DataTypeTok{"moneda"}\FunctionTok{:}\StringTok{"MXN"}\FunctionTok{\}}
  \FunctionTok{\}}\OtherTok{]}
\FunctionTok{\}}
\end{Highlighting}
\end{Shaded}

Después se recuperan saldo y
\texorpdfstring{\softtext{disponible.}}{disponible.} Un saldo cero
recibido confirma la aplicación económica, no la entrega física. Si el
comprador aún no recibió ganado, ENTREGA se registra por separado con
posiciones, venta confirmada y fecha. Una bonificación posterior
de100.00 puede solicitar
\texorpdfstring{\softcode{reaplicar\_pagos:true}}{reaplicar\_pagos:true}
en un ABONO ADMIN; el servidor conserva originales, deja aplicaciones
\texorpdfstring{\softtext{netas900.00}}{netas900.00} y
\texorpdfstring{\softtext{disponible100.00}}{disponible100.00} para
reembolso. La interfaz no ejecuta esa conciliación sumando renglones ni
enviando tres comandos ocultos.

Fuentes:
\texorpdfstring{\softcode{frontend/src/App.tsx:35}}{frontend/src/App.tsx:35},
\texorpdfstring{\softcode{frontend/src/screens.tsx:31}}{frontend/src/screens.tsx:31},
\texorpdfstring{\softcode{frontend/src/screens.tsx:14}}{frontend/src/screens.tsx:14},
\texorpdfstring{\softcode{frontend/src/screens.tsx:39}}{frontend/src/screens.tsx:39},
\texorpdfstring{\softcode{frontend/src/screens.tsx:82}}{frontend/src/screens.tsx:82},
\texorpdfstring{\softcode{frontend/tests\_e2e/closure.spec.ts:33}}{frontend/tests\_e2e/closure.spec.ts:33},
\texorpdfstring{\softcode{frontend/tests\_e2e/closure.spec.ts:41}}{frontend/tests\_e2e/closure.spec.ts:41},
\texorpdfstring{\softcode{frontend/tests\_e2e/closure.spec.ts:88}}{frontend/tests\_e2e/closure.spec.ts:88}.

\section{15. Estados, borrador y resultado incierto: cómo pensar una
falla}\label{estados-borrador-y-resultado-incierto-cuxf3mo-pensar-una-falla}

Un borrador existe en SQL sólo después de que el servidor confirma su
alta o edición. Pasar de un paso visual a otro no guarda automáticamente
los campos. «Guardar condiciones del lote» es una acción explícita; el
callback actualiza la revisión y consulta el detalle. Dejar una página
con texto no enviado puede perder ese texto. La aplicación no ofrece
captura offline ni autosave universal, y esa ausencia debe explicarse al
usuario antes de asumir que una propuesta quedó
\texorpdfstring{\softtext{persistida.}}{persistida.}

Para un comando ordinario, el instante de envío congela cuerpo, ruta,
método y versión en \texorpdfstring{\softcode{Attempt}}{Attempt}.
Mientras no se conozca el resultado, el formulario bloquea sus campos.
Recargar puede recuperar esa intención por principal y ruta, pero no
restaura necesariamente toda la presentación previa. Por ejemplo, puede
ser necesario volver al detalle correcto; el contenido original que se
reenvía proviene del intento, no de escribir valores aproximados en
controles nuevos.

Un 401 provoca un evento que lleva de regreso al acceso. La identidad y
CSRF de memoria se eliminan; el almacenamiento queda separado por el
usuario anterior. Al entrar como otro usuario no aparece su intento. Al
entrar como el mismo usuario puede recuperarse y consultar el resultado.
Este mecanismo no permite continuar con una cookie expirada ni sustituir
la autorización del lookup. Cada respuesta económica recuperada debe
proceder del servidor
\texorpdfstring{\softtext{autenticado.}}{autenticado.}

Un 412 significa que la versión revisada ya no es la actual. El usuario
obtiene el registro vigente y vuelve a valorar sus cambios. Un 428
significa que no acompañó versión cuando se
\texorpdfstring{\softtext{necesitaba.}}{necesitaba.} No se arreglan
enviando una versión inventada. La interfaz utiliza ETag del GET y, en
el editor, dispone de una representación entre comillas de
\texorpdfstring{\softcode{version\_lock}}{version\_lock} recibido como
fallback. Esa versión sigue viniendo del servidor. El proxy integrado
debe conservar el ETag fuerte; transformarlo con compresión de API puede
romper la \texorpdfstring{\softtext{comparación.}}{comparación.}

Un 422 conocido habilita corregir la captura porque el intento anterior
fue rechazado. Un 5 xx o una respuesta que no puede interpretarse no
tiene esa certeza. Lo correcto es consultar primero. Tampoco se afirma
que cada 500 o504 haya sido provocado realmente en E2E: su clasificación
está cubierta por unitarias. La prueba502 alteró sólo la entrega de una
respuesta201 real, lo que reproduce precisamente el riesgo de desconocer
una confirmación ya realizada.

En cargas documentales y creación sensible la recuperación es más
cuidadosa. Se guarda clave y huella, no el contenido protegido. Si el
lookup confirma, termina sin pedir otra subida o
\texorpdfstring{\softtext{contraseña.}}{contraseña.} Si no confirma y la
memoria desapareció al recargar, se recaptura exactamente el contenido
original. Cambiar el archivo o contraseña mientras la intención está
pendiente produciría otra decisión, por lo que la comparación de huella
rechaza el cambio.

La UX de conflicto mantiene una referencia de regla para soporte y un
paso de resolución visible. No anuncia «todo está bien» después de un
timeout ni muestra un saldo cero para ocultar una consulta ausente. Para
diagnosticar, conviene separar: error de captura antes de enviar,
rechazo conocido del servidor, versión antigua y resultado
\texorpdfstring{\softtext{desconocido.}}{desconocido.} Esa clasificación
permite actuar con propósito y evita tratar todos los avisos como una
invitación a pulsar Guardar
\texorpdfstring{\softtext{repetidamente.}}{repetidamente.}

Fuentes:
\texorpdfstring{\softcode{frontend/src/ui.tsx:26}}{frontend/src/ui.tsx:26},
\texorpdfstring{\softcode{frontend/src/api.ts:11}}{frontend/src/api.ts:11},
\texorpdfstring{\softcode{frontend/src/api.ts:19}}{frontend/src/api.ts:19},
\texorpdfstring{\softcode{frontend/src/screens.tsx:53}}{frontend/src/screens.tsx:53},
\texorpdfstring{\softcode{frontend/src/screens.tsx:102}}{frontend/src/screens.tsx:102},
\texorpdfstring{\softcode{frontend/tests/workflow.test.ts:21}}{frontend/tests/workflow.test.ts:21},
\texorpdfstring{\softcode{frontend/tests\_e2e/closure.spec.ts:90}}{frontend/tests\_e2e/closure.spec.ts:90}.

\section{16. CSS, accesibilidad y adaptación de
tamaños}\label{css-accesibilidad-y-adaptaciuxf3n-de-tamauxf1os}

La hoja de estilos define fondo
\texorpdfstring{\softcode{\#F6F7F4}}{\#F6F7F4}, texto
\texorpdfstring{\softcode{\#17211B}}{\#17211B}, acción
\texorpdfstring{\softcode{\#1F5A43}}{\#1F5A43}, secundario
\texorpdfstring{\softcode{\#5A685F}}{\#5A685F}, borde
\texorpdfstring{\softcode{\#D8E0D8}}{\#D8E0D8}, superficie suave
\texorpdfstring{\softcode{\#E9EFE9}}{\#E9EFE9} y error
\texorpdfstring{\softcode{\#9F2922}}{\#9F2922}. La fuente comienza por
Inter como preferencia y continúa con familias del sistema; no hay una
descarga de fuente declarada en esa hoja. Los tokens reducen decisiones
aisladas y permiten que acceso, inventario y dinero conserven la misma
identidad. El foco visible usa tres píxeles de color cálido y separación
del control.

El layout de escritorio usa
\texorpdfstring{\softcode{app-shell}}{app-shell} flexible, lateral fijo
y un \texorpdfstring{\softcode{workspace}}{workspace} con
\texorpdfstring{\softcode{min-width:0}}{min-width:0}. Este último
detalle permite que el contenido flexible encoja sin imponer su ancho
intrínseco al documento. \texorpdfstring{\softcode{main}}{main} posee
ancho completo y máximo1600 px. Las tarjetas
\texorpdfstring{\softcode{panel}}{panel} utilizan borde, radio y
espaciado \texorpdfstring{\softtext{consistente.}}{consistente.}
\texorpdfstring{\softcode{split}}{split} organiza dos columnas mediante
\texorpdfstring{\softcode{minmax}}{minmax}, y
\texorpdfstring{\softcode{form-grid}}{form-grid} distribuye campos.
\texorpdfstring{\softcode{facts}}{facts} y
\texorpdfstring{\softcode{metric-grid}}{metric-grid} utilizan columnas
adaptables para información y métricas.

Los botones tienen objetivo mínimo44 px y los controles de campo45 px.
Las cifras monetarias usan
\texorpdfstring{\softcode{font-variant-numeric:tabular-nums}}{font-variant-numeric:tabular-nums},
alineación derecha y una representación completa. Un identificador largo
puede envolver mediante
\texorpdfstring{\softcode{overflow-wrap:anywhere}}{overflow-wrap:anywhere};
no se utiliza como motivo para abreviar una UUID y hacer ambiguo el
registro \texorpdfstring{\softtext{seleccionado.}}{seleccionado.} Los
estados combinan texto y badge; el color por sí solo no contiene toda la
\texorpdfstring{\softtext{información.}}{información.}

Las tablas desplazan dentro de
\texorpdfstring{\softcode{table-region}}{table-region}, que tiene
\texorpdfstring{\softcode{overflow-x:auto}}{overflow-x:auto} y ancho
máximo \texorpdfstring{\softtext{disponible.}}{disponible.} En móvil la
tabla conserva un ancho mínimo570 px: el usuario recorre la región, no
toda la página. La regla
\texorpdfstring{\softcode{.text-button\{position:relative\}}}{.text-button\{position:relative\}}
y el anclaje del texto \texorpdfstring{\softcode{sr-only}}{sr-only}
evitan que un texto invisible posicionado fuera de su control extienda
el documento. Una región accesible con tabulación permite llegar a la
tabla mediante teclado y comprender qué conjunto contiene.

A1100 px el lateral y espaciados se reducen y el layout dividido pasa a
una columna. A767 px se apilan bloques, resumen económico y
\texorpdfstring{\softtext{formulario.}}{formulario.} Las reglas
posteriores ocultan la navegación lateral y muestran
\texorpdfstring{\softcode{mobile-navigation}}{mobile-navigation}, un
select nativo con grupos. No deben interpretarse reglas anteriores del
archivo aisladas de la cascada: al final existe una corrección que
reemplaza el menú horizontal móvil por el selector compacto. Esa
secuencia explica el estilo aplicado realmente.

El acceso móvil reduce la banda gráfica a la marca y coloca el
formulario antes de una región promocional extensa. La tipografía de
inputs pasa a16 px. \texorpdfstring{\softcode{skip}}{skip} aparece al
obtener foco y salta hacia contenido. Los campos poseen etiqueta,
errores asociados y ayudas. El resumen de error se puede enfocar; no se
mueve el foco arbitrariamente al cambiar cada carácter. La preferencia
\texorpdfstring{\softcode{prefers-reduced-motion}}{prefers-reduced-motion}
elimina animaciones, incluido el indicador giratorio.

La evidencia de390,768 y1440 px comprueba ausencia de desplazamiento
horizontal global en los recorridos
\texorpdfstring{\softtext{registrados.}}{registrados.} No es equivalente
a probar un teclado virtual de teléfono, cada combinación de texto
largo, un lector de pantalla o conformidad WCAG global. La suite de
teclado sí completa un flujo principal real sin ratón ni
\texorpdfstring{\softcode{locator.focus}}{locator.focus}; el uso de un
\texorpdfstring{\softcode{focus}}{focus} interno de la aplicación para
anunciar navegación o error es parte del comportamiento del producto y
no una trampa del runner.

Fuentes:
\texorpdfstring{\softcode{frontend/src/styles.css:1}}{frontend/src/styles.css:1},
\texorpdfstring{\softcode{frontend/src/styles.css:3}}{frontend/src/styles.css:3},
\texorpdfstring{\softcode{frontend/src/styles.css:5}}{frontend/src/styles.css:5},
\texorpdfstring{\softcode{frontend/src/styles.css:6}}{frontend/src/styles.css:6},
\texorpdfstring{\softcode{frontend/src/styles.css:8}}{frontend/src/styles.css:8},
\texorpdfstring{\softcode{frontend/src/styles.css:11}}{frontend/src/styles.css:11},
\texorpdfstring{\softcode{frontend/src/ui.tsx:21}}{frontend/src/ui.tsx:21},
\texorpdfstring{\softcode{frontend/src/App.tsx:29}}{frontend/src/App.tsx:29}.

\section{17. Pruebas: qué verifican y qué límites
conservan}\label{pruebas-quuxe9-verifican-y-quuxe9-luxedmites-conservan}

Hay trece pruebas unitarias entre
\texorpdfstring{\softcode{domain.test.ts}}{domain.test.ts} y
\texorpdfstring{\softcode{workflow.test.ts}}{workflow.test.ts}.
Verifican formato exacto de un importe mayor que el entero seguro de
\texorpdfstring{\softtext{JavaScript,}}{JavaScript,} captura decimal y
rechazo de
\texorpdfstring{\softtext{exponentes/negativos/exceso}}{exponentes/negativos/exceso}
de escala, conservación de cuerpo y ETag al mutar el objeto original,
estructura accesible de campos, cabecera CSRF, fecha mexicana,
recuperación de lookup sin POST adicional, estado EN\_CURSO,
identificadores únicos entre formularios, orden de lotes, composición
animal de venta y aislamiento del intento al expirar sesión.

Varias unitarias sustituyen \texorpdfstring{\softcode{fetch}}{fetch} o
almacenamiento por objetos
\texorpdfstring{\softtext{controlados.}}{controlados.} Esa técnica
permite observar exactamente qué se envía y restaurar el objeto original
al terminar. No es evidencia de PostgreSQL real. El render a markup
comprueba atributos, pero no certifica navegación con lector de
pantalla. El capítulo mantiene esas distinciones porque una prueba
pequeña es valiosa precisamente cuando se conoce el comportamiento
acotado que demuestra.

La configuración Playwright usa un trabajador, sin paralelismo ni
retries, timeout de caso240 s y expectativas15 s. Abre Chrome instalado
con ventana visible y perfil separado. Puede recibir
\texorpdfstring{\softcode{BOVI\_UI\_TEST\_URL}}{BOVI\_UI\_TEST\_URL}; el
valor de desarrollo predeterminado es15173 y la tanda final usó
\texorpdfstring{\softtext{Caddy15180.}}{Caddy15180.} No utiliza trazas
de navegación para versionar cookies o
\texorpdfstring{\softtext{contraseñas.}}{contraseñas.} Los secretos
necesarios permanecen en el directorio privado local y no se reproducen
aquí.

\texorpdfstring{\softcode{runtime.ts}}{runtime.ts} exporta un
\texorpdfstring{\softcode{test}}{test} extendido con fixture automática
\texorpdfstring{\softcode{runtimeLog}}{runtimeLog}. Escucha errores de
consola, errores JavaScript y respuestas HTTP de error; conserva título,
estado, duración y URL base en JSON temporal. Los contextos secundarios
necesitan observación propia: el caso de reservas registra por separado
su segundo operador. Por ello un total de errores del contexto principal
no debe presentarse como un inventario exhaustivo de cualquier ventana
\texorpdfstring{\softtext{imaginable.}}{imaginable.}

{\def\LTcaptype{none} % do not increment counter
\begin{longtable}[]{@{}
  >{\raggedright\arraybackslash}p{(\linewidth - 4\tabcolsep) * \real{0.3333}}
  >{\raggedright\arraybackslash}p{(\linewidth - 4\tabcolsep) * \real{0.3333}}
  >{\raggedright\arraybackslash}p{(\linewidth - 4\tabcolsep) * \real{0.3333}}@{}}
\toprule\noalign{}
\begin{minipage}[b]{\linewidth}\raggedright
Grupo de prueba
\end{minipage} & \begin{minipage}[b]{\linewidth}\raggedright
Comportamiento concreto
\end{minipage} & \begin{minipage}[b]{\linewidth}\raggedright
Fuente
\end{minipage} \\
\midrule\noalign{}
\endhead
\bottomrule\noalign{}
\endlastfoot
Flujo principal & Compra, recepción, venta, cobro, abono atómico,
reembolso, archivo y QR &
\texorpdfstring{\softcode{frontend/tests\_e2e/closure.spec.ts:27}}{frontend/tests\_e2e/closure.spec.ts:27} \\
Respuesta perdida & POST real confirmado, entrega abortada y lookup sin
segundo POST &
\texorpdfstring{\softcode{frontend/tests\_e2e/closure.spec.ts:48}}{frontend/tests\_e2e/closure.spec.ts:48} \\
Auditor & \texorpdfstring{\softtext{Alta/membresía}}{Alta/membresía}
real, controles ocultos y escritura directa403 &
\texorpdfstring{\softcode{frontend/tests\_e2e/closure.spec.ts:57}}{frontend/tests\_e2e/closure.spec.ts:57} \\
Gastos & CompraKG, negativa422, pesos medidos, tres repartos y tarifa
aceptada &
\texorpdfstring{\softcode{frontend/tests\_e2e/closure.spec.ts:63}}{frontend/tests\_e2e/closure.spec.ts:63} \\
Acuerdos & Reapertura, edición y dos versiones aceptadas conservadas &
\texorpdfstring{\softcode{frontend/tests\_e2e/closure.spec.ts:73}}{frontend/tests\_e2e/closure.spec.ts:73} \\
Periodos &
\texorpdfstring{\softtext{Cierre/reapertura}}{Cierre/reapertura}
corriente e historia de versiones &
\texorpdfstring{\softcode{frontend/tests\_e2e/closure.spec.ts:76}}{frontend/tests\_e2e/closure.spec.ts:76} \\
Regresión KG & Dos lotes preservados y
\texorpdfstring{\softtext{venta300.000kg/3000.00}}{venta300.000kg/3000.00}
&
\texorpdfstring{\softcode{frontend/tests\_e2e/closure.spec.ts:81}}{frontend/tests\_e2e/closure.spec.ts:81} \\
HTTP502 & Entrega simulada502 después de POST201; clave conservada &
\texorpdfstring{\softcode{frontend/tests\_e2e/closure.spec.ts:90}}{frontend/tests\_e2e/closure.spec.ts:90} \\
Cambio de sede & GET retenido B abortado al pasar a A &
\texorpdfstring{\softcode{frontend/tests\_e2e/closure.spec.ts:94}}{frontend/tests\_e2e/closure.spec.ts:94} \\
Predios & Validación, alta, búsqueda, páginas20+1 y refresh &
\texorpdfstring{\softcode{frontend/tests\_e2e/additional.spec.ts:21}}{frontend/tests\_e2e/additional.spec.ts:21} \\
Ciclos & Entrega física y recompra con dos ciclos históricos &
\texorpdfstring{\softcode{frontend/tests\_e2e/additional.spec.ts:27}}{frontend/tests\_e2e/additional.spec.ts:27} \\
Corrección & Original100kg, versión110kg, motivo y grupo incompatible &
\texorpdfstring{\softcode{frontend/tests\_e2e/additional.spec.ts:32}}{frontend/tests\_e2e/additional.spec.ts:32} \\
Reportes & \texorpdfstring{\softtext{Corte/fuente}}{Corte/fuente} y
descarga CSV real de toda la sede &
\texorpdfstring{\softcode{frontend/tests\_e2e/additional.spec.ts:37}}{frontend/tests\_e2e/additional.spec.ts:37} \\
Reservas & Dos usuarios, conflicto409 y vencimiento real409 &
\texorpdfstring{\softcode{frontend/tests\_e2e/additional.spec.ts:42}}{frontend/tests\_e2e/additional.spec.ts:42} \\
Teclado390 &
\texorpdfstring{\softtext{Compra→recepción→venta→pago}}{Compra→recepción→venta→pago}
sin ratón ni \texorpdfstring{\softtext{locator.focus}}{locator.focus} &
\texorpdfstring{\softcode{frontend/tests\_e2e/keyboard.spec.ts:21}}{frontend/tests\_e2e/keyboard.spec.ts:21} \\
Proveedor físico & GET503, captura preservada y recuperación200 &
\texorpdfstring{\softcode{frontend/tests\_e2e/provider.spec.ts:6}}{frontend/tests\_e2e/provider.spec.ts:6} \\
\end{longtable}
}

Los helpers de
\texorpdfstring{\softcode{closure.spec.ts}}{closure.spec.ts} buscan
contrapartes y animales mediante controles reales, seleccionan periodos
y esperan una respuesta por
\texorpdfstring{\softtext{método/ruta.}}{método/ruta.} Los de
\texorpdfstring{\softcode{additional.spec.ts}}{additional.spec.ts}
pueden crear precondiciones DEMO mediante API explícita; su función
\texorpdfstring{\softcode{operation}}{operation} obtiene GET y ETag
entre \texorpdfstring{\softtext{transiciones.}}{transiciones.} Por eso
no se atribuye a clicks la creación de cada
\texorpdfstring{\softtext{precondición.}}{precondición.} El helper
\texorpdfstring{\softcode{responsive}}{responsive} cambia viewport y
comprueba el ancho del documento, además de capturar.

En \texorpdfstring{\softcode{keyboard.spec.ts}}{keyboard.spec.ts},
\texorpdfstring{\softcode{tabTo}}{tabTo} llega al destino usando Tab y
sólo lee si obtuvo foco; \texorpdfstring{\softcode{enter}}{enter},
\texorpdfstring{\softcode{input}}{input},
\texorpdfstring{\softcode{select}}{select} y
\texorpdfstring{\softcode{datetime}}{datetime} utilizan eventos de
teclado. El datetime se captura por segmentos del control nativo de
Chrome. Sólo identidad y contraparte son fixtures API; el flujo
financiero utiliza
\texorpdfstring{\softtext{formularios.}}{formularios.} La evidencia
registra534 eventos,364Tab, ancho390 y
\texorpdfstring{\softtext{scrollWidth375.}}{scrollWidth375.} Estos
números describen esa corrida, no una cantidad universal de pasos para
cualquier conjunto de datos.

La prueba física del proveedor requiere
\texorpdfstring{\softcode{BOVI\_PROVIDER\_OUTAGE=true}}{BOVI\_PROVIDER\_OUTAGE=true}.
Si no existe esa autorización coordinada, se omite con razón
\texorpdfstring{\softtext{explícita.}}{explícita.} En la tanda final se
ejecutó aparte, sin omisiones. Después del503 crea un archivo
\texorpdfstring{\softcode{ready}}{ready}; espera que el responsable
reanude, compruebe salud y escriba literalmente
\texorpdfstring{\softcode{resumed}}{resumed}; entonces pulsa volver a
consultar. No acepta precio ni crea gasto durante la caída. El login sí
registra sesión, por lo que el resultado declara cero escrituras
comerciales, no ausencia total de POST.

El cierre registrado fue15/15 integrados y1/1 proveedor físico, cero
\texorpdfstring{\softtext{retries/skips.}}{retries/skips.} El reportero
de la tanda15 mostró1,3 min redondeados; las duraciones individuales
\texorpdfstring{\softtext{sumaron73.520}}{sumaron73.520} ms, sin
confundir esa suma con tiempo global. El proveedor
\texorpdfstring{\softtext{tuvo47.431,536}}{tuvo47.431,536} ms de tanda
y46.406 ms de caso. Se registraron24 entradas esperadas de consola en
contextos principales y cero errores
\texorpdfstring{\softtext{JavaScript;401}}{JavaScript;401} anónimo,400
de motivo excesivo,409,403,422,502 inducido,503 físico y entrega
abortada se \texorpdfstring{\softtext{clasificaron.}}{clasificaron.} No
se anuncia consola vacía.

La GPU registrada fue Intel integrada mediante
\texorpdfstring{\softtext{ANGLE/OpenGL,}}{ANGLE/OpenGL,} sin
\texorpdfstring{\softtext{SwiftShader.}}{SwiftShader.} Eso demuestra el
dispositivo de esas capturas, no una medición FPS. La restauración de
datos, las pruebas de nodos y el benchmark SQL tienen sus propias
evidencias del coordinador; no se atribuyen al runner frontend. Tampoco
hay entrevistas, prueba con clientes reales, aprobación humana o
despliegue público derivados de estos
\texorpdfstring{\softtext{resultados.}}{resultados.}

Fuentes:
\texorpdfstring{\softcode{frontend/tests/domain.test.ts:8}}{frontend/tests/domain.test.ts:8},
\texorpdfstring{\softcode{frontend/tests/workflow.test.ts:6}}{frontend/tests/workflow.test.ts:6},
\texorpdfstring{\softcode{frontend/playwright.config.ts:2}}{frontend/playwright.config.ts:2},
\texorpdfstring{\softcode{frontend/tests\_e2e/runtime.ts:3}}{frontend/tests\_e2e/runtime.ts:3},
\texorpdfstring{\softcode{frontend/tests\_e2e/keyboard.spec.ts:8}}{frontend/tests\_e2e/keyboard.spec.ts:8},
\texorpdfstring{\softcode{frontend/tests\_e2e/provider.spec.ts:7}}{frontend/tests\_e2e/provider.spec.ts:7},
\texorpdfstring{\softcode{docs/testing/evidence/ui/bovicierre-ui-final-run.json:1}}{docs/testing/evidence/ui/bovicierre-ui-final-run.json:1},
\texorpdfstring{\softcode{docs/development/frontend-context.md:25}}{docs/development/frontend-context.md:25}.

\section{18. Guía de estudio y lectura del código sin alterar
datos}\label{guuxeda-de-estudio-y-lectura-del-cuxf3digo-sin-alterar-datos}

Para estudiar el sistema conviene seguir primero una lectura sin
\texorpdfstring{\softtext{ejecución.}}{ejecución.} Abra
\texorpdfstring{\softcode{main.tsx}}{main.tsx}, identifique
\texorpdfstring{\softcode{App}}{App}, siga
\texorpdfstring{\softcode{Workspace}}{Workspace} hacia una pantalla y
anote su ruta. Después busque el
\texorpdfstring{\softcode{build}}{build} del formulario y el helper que
normaliza sus campos. Compare el cuerpo con OpenAPI y observe qué campos
no envía: total y saldo de confirmación no se calculan en GUI. Continúe
hacia \texorpdfstring{\softcode{CommandForm}}{CommandForm},
\texorpdfstring{\softcode{createAttempt}}{createAttempt} y
\texorpdfstring{\softcode{request}}{request}, y termine identificando el
\texorpdfstring{\softcode{onDone}}{onDone} que consulta o navega después
del resultado.

Una segunda lectura puede seguir el camino de un error. Pregunte qué
ocurre antes del envío, qué ocurre ante422 y qué cambia ante502. Observe
cuándo se borra el intento, dónde queda el ETag y cuál es la identidad
que forma la clave de
\texorpdfstring{\softtext{almacenamiento.}}{almacenamiento.} En archivos
o usuarios, examine la diferencia entre huella persistida y contenido de
memoria. Este ejercicio revela más del diseño que memorizar colores o
endpoints aislados.

Una tercera lectura puede estudiar estado y alcance. Localice
\texorpdfstring{\softcode{key}}{key} en
\texorpdfstring{\softcode{Workspace}}{Workspace},
\texorpdfstring{\softcode{AnimalCycles}}{AnimalCycles},
\texorpdfstring{\softcode{LotEditor}}{LotEditor}, correcciones y
destinos de pago. Identifique efectos que abortan consultas y callbacks
directos que recargan recursos. No suponga que todo prop inicial
actualiza un \texorpdfstring{\softcode{useState}}{useState} ya montado.
La clave de React, las dependencias del efecto y la ruta con sede
trabajan juntas, pero cada mecanismo resuelve una parte distinta del
problema.

Por último, relacione funciones con evidencia. Una función de formato se
estudia con unitarias; una recuperación de POST confirmado requiere la
prueba de red correspondiente; una afirmación responsive requiere sus
métricas y capturas; una incidencia física requiere coordinación de
servicios. La documentación puede explicar cómo repetir un caso, pero no
sustituye su \texorpdfstring{\softtext{ejecución.}}{ejecución.} Durante
una sesión compartida no se lanzan pruebas que escriben mientras otro
responsable realiza backup o altera funciones SQL.

El frontend implementado permite trabajar con evidencia y estados
concretos. Su claridad depende de conservar las distinciones que
estructuran el dominio: identidad frente a ciclo, borrador frente a
cierre, consulta frente a aceptación de tarifa, dinero registrado frente
a aplicación, reverso frente a borrado y resultado incierto frente a
rechazo conocido. Cada una tiene componentes, mensajes y comandos
visibles en el código descrito. La revisión humana del producto sigue
siendo un paso separado del hecho de que esas funciones existan y sus
casos locales hayan pasado.

\chapter{Atlas de scripts Python y sus
variables}\label{atlas-de-scripts-python-y-sus-variables}

El atlas registra firmas y nombres de variables desde AST de Python, con
línea inicial y final. Un nombre local sólo existe durante su función;
una constante de módulo puede compartirse dentro del script. No
confundir estos nombres con variables de entorno. La lista incluye
asignaciones directas y targets de bucles; argumentos, bindings de
excepciones y closures deben leerse con el fragmento de código para
conocer su alcance exacto.

\section{\texorpdfstring{\texorpdfstring{\softtext{scripts/doctor.py}}{scripts/doctor.py}}{}}\label{scriptsdoctor.py}

Diagnóstico local sin modificar servicios ni mostrar configuración
privada.

\textbf{Variables de módulo:}
\texorpdfstring{\softcode{COMMANDS}}{COMMANDS}.

\subsection{\texorpdfstring{\texorpdfstring{\softtext{inspect\_command}}{inspect\_command}}{}}\label{inspect_command}

\textbf{Firma:}
\texorpdfstring{\softcode{inspect\_command(command)}}{inspect\_command(command)}.
Fuente
\texorpdfstring{\softcode{scripts/doctor.py:21-29}}{scripts/doctor.py:21-29}.

Consulta una versión con timeout y registra ausencia o exit code sin
alterar servicios.

\textbf{Variables locales asignadas:}
\texorpdfstring{\softcode{lines}}{lines},
\texorpdfstring{\softcode{result}}{result}.

\subsection{listening}\label{listening}

\textbf{Firma:}
\texorpdfstring{\softcode{listening(port)}}{listening(port)}. Fuente
\texorpdfstring{\softcode{scripts/doctor.py:32-35}}{scripts/doctor.py:32-35}.

Comprueba si hay un listener TCP
\texorpdfstring{\softtext{en127.0.0.1}}{en127.0.0.1} para el puerto
indicado.

\subsection{resources}\label{resources}

\textbf{Firma:} \texorpdfstring{\softcode{resources()}}{resources()}.
Fuente
\texorpdfstring{\softcode{scripts/doctor.py:38-41}}{scripts/doctor.py:38-41}.

Obtiene sistema, shell, CPU y memoria desde información del host.

\textbf{Variables locales asignadas:}
\texorpdfstring{\softcode{memory}}{memory},
\texorpdfstring{\softcode{total}}{total}.

\subsection{main}\label{main}

\textbf{Firma:} \texorpdfstring{\softcode{main()}}{main()}. Fuente
\texorpdfstring{\softcode{scripts/doctor.py:44-57}}{scripts/doctor.py:44-57}.

Entrada CLI que interpreta las opciones y coordina las funciones del
archivo; su firma no concede permisos
\texorpdfstring{\softtext{comerciales.}}{comerciales.}

\textbf{Variables locales asignadas:}
\texorpdfstring{\softcode{result}}{result}.

\section{\texorpdfstring{\texorpdfstring{\softtext{scripts/academic.py}}{scripts/academic.py}}{}}\label{scriptsacademic.py}

Arranque académico integrado; rechaza un backend manual
\texorpdfstring{\softtext{concurrente.}}{concurrente.}

\textbf{Variables de módulo:} \texorpdfstring{\softcode{ROOT}}{ROOT}.

\subsection{busy}\label{busy}

\textbf{Firma:} \texorpdfstring{\softcode{busy(port)}}{busy(port)}.
Fuente
\texorpdfstring{\softcode{scripts/academic.py:11-14}}{scripts/academic.py:11-14}.

Detecta el puerto de API manual ocupado antes del arranque integrado.

\subsection{main}\label{main-1}

\textbf{Firma:} \texorpdfstring{\softcode{main()}}{main()}. Fuente
\texorpdfstring{\softcode{scripts/academic.py:17-28}}{scripts/academic.py:17-28}.

Entrada CLI que interpreta las opciones y coordina las funciones del
archivo; su firma no concede permisos
\texorpdfstring{\softtext{comerciales.}}{comerciales.}

\textbf{Variables locales asignadas:}
\texorpdfstring{\softcode{docs}}{docs},
\texorpdfstring{\softcode{env}}{env},
\texorpdfstring{\softcode{private}}{private}.

\section{\texorpdfstring{\texorpdfstring{\softtext{scripts/dev.py}}{scripts/dev.py}}{}}\label{scriptsdev.py}

Ejecuta API y Vite en primer plano; Ctrl+C detiene únicamente sus hijos.

\textbf{Variables de módulo:} \texorpdfstring{\softcode{ROOT}}{ROOT},
\texorpdfstring{\softcode{COMMANDS}}{COMMANDS}.

\subsection{stop}\label{stop}

\textbf{Firma:}
\texorpdfstring{\softcode{stop(children)}}{stop(children)}. Fuente
\texorpdfstring{\softcode{scripts/dev.py:16-25}}{scripts/dev.py:16-25}.

Termina únicamente los grupos de procesos hijos creados por make dev;
espera y limita la escalada.

\textbf{Variables locales asignadas:}
\texorpdfstring{\softcode{child}}{child}.

\subsection{main}\label{main-2}

\textbf{Firma:} \texorpdfstring{\softcode{main()}}{main()}. Fuente
\texorpdfstring{\softcode{scripts/dev.py:28-50}}{scripts/dev.py:28-50}.

Entrada CLI que interpreta las opciones y coordina las funciones del
archivo; su firma no concede permisos
\texorpdfstring{\softtext{comerciales.}}{comerciales.}

\textbf{Variables locales asignadas:}
\texorpdfstring{\softcode{child}}{child},
\texorpdfstring{\softcode{children}}{children},
\texorpdfstring{\softcode{command}}{command},
\texorpdfstring{\softcode{directory}}{directory}.

\section{\texorpdfstring{\texorpdfstring{\softtext{scripts/recovery.py}}{scripts/recovery.py}}{}}\label{scriptsrecovery.py}

Respaldo privado y ensayo aislado
\texorpdfstring{\softtext{BoviCierre;}}{BoviCierre;} nunca elimina
\texorpdfstring{\softtext{volúmenes.}}{volúmenes.}

\textbf{Variables de módulo:} \texorpdfstring{\softcode{ROOT}}{ROOT},
\texorpdfstring{\softcode{PRIVATE}}{PRIVATE},
\texorpdfstring{\softcode{SECRETS}}{SECRETS},
\texorpdfstring{\softcode{DOCS}}{DOCS},
\texorpdfstring{\softcode{SOURCE}}{SOURCE}.

\subsection{run}\label{run}

\textbf{Firma:}
\texorpdfstring{\softcode{run(args, **kwargs)}}{run(args, **kwargs)}.
Fuente
\texorpdfstring{\softcode{scripts/recovery.py:24-29}}{scripts/recovery.py:24-29}.

Ejecuta una utilidad y conserva su salida; en recuperación filtra el
error para no divulgar cadenas de conexión.

\textbf{Variables locales asignadas:}
\texorpdfstring{\softcode{result}}{result}.

\subsection{sql}\label{sql}

\textbf{Firma:}
\texorpdfstring{\softcode{sql(project, node, query)}}{sql(project, node, query)}.
Fuente
\texorpdfstring{\softcode{scripts/recovery.py:32-37}}{scripts/recovery.py:32-37}.

Ejecuta SQL sobre un nodo identificado; en recuperación verifica la
etiqueta del contenedor antes de
\texorpdfstring{\softtext{utilizarlo.}}{utilizarlo.}

\textbf{Variables locales asignadas:}
\texorpdfstring{\softcode{label}}{label},
\texorpdfstring{\softcode{name}}{name}.

\subsection{manifest}\label{manifest}

\textbf{Firma:}
\texorpdfstring{\softcode{manifest(directory)}}{manifest(directory)}.
Fuente
\texorpdfstring{\softcode{scripts/recovery.py:40-49}}{scripts/recovery.py:40-49}.

Recorre archivos privados, rechaza enlaces simbólicos y obtiene SHA-256
por ruta relativa.

\textbf{Variables locales asignadas:}
\texorpdfstring{\softcode{directory}}{directory},
\texorpdfstring{\softcode{path}}{path},
\texorpdfstring{\softcode{records}}{records},
\texorpdfstring{\softcode{str}}{str}.

\subsection{signature}\label{signature}

\textbf{Firma:}
\texorpdfstring{\softcode{signature(project, node)}}{signature(project, node)}.
Fuente
\texorpdfstring{\softcode{scripts/recovery.py:52-54}}{scripts/recovery.py:52-54}.

Obtiene conteos y huellas de filas de negocio replicables para comprobar
\texorpdfstring{\softtext{convergencia.}}{convergencia.}

\textbf{Variables locales asignadas:}
\texorpdfstring{\softcode{tables}}{tables}.

\subsection{\texorpdfstring{\texorpdfstring{\softtext{local\_signature}}{local\_signature}}{}}\label{local_signature}

\textbf{Firma:}
\texorpdfstring{\softcode{local\_signature(project, node)}}{local\_signature(project, node)}.
Fuente
\texorpdfstring{\softcode{scripts/recovery.py:57-58}}{scripts/recovery.py:57-58}.

Comprueba tablas locales de deduplicación, inbox, entrega, captura y
Flyway que no se comparan como réplicas.

\subsection{\texorpdfstring{\texorpdfstring{\softtext{private\_dir}}{private\_dir}}{}}\label{private_dir}

\textbf{Firma:}
\texorpdfstring{\softcode{private\_dir(path)}}{private\_dir(path)}.
Fuente
\texorpdfstring{\softcode{scripts/recovery.py:61-63}}{scripts/recovery.py:61-63}.

Crea un directorio nuevo con permisos privados; la existencia previa es
un error.

\subsection{\texorpdfstring{\texorpdfstring{\softtext{backend\_mode}}{backend\_mode}}{}}\label{backend_mode}

\textbf{Firma:}
\texorpdfstring{\softcode{backend\_mode(pid)}}{backend\_mode(pid)}.
Fuente
\texorpdfstring{\softcode{scripts/recovery.py:66-84}}{scripts/recovery.py:66-84}.

Distingue la API Docker propia de la Java nativa comprobando proceso,
cwd, puerto y runtime.

\textbf{Variables locales asignadas:}
\texorpdfstring{\softcode{container}}{container},
\texorpdfstring{\softcode{environment}}{environment},
\texorpdfstring{\softcode{native}}{native},
\texorpdfstring{\softcode{pid}}{pid},
\texorpdfstring{\softcode{pid\_file}}{pid\_file},
\texorpdfstring{\softcode{process}}{process},
\texorpdfstring{\softcode{running}}{running},
\texorpdfstring{\softcode{state}}{state}.

\subsection{backup}\label{backup}

\textbf{Firma:}
\texorpdfstring{\softcode{backup(destination, pid)}}{backup(destination, pid)}.
Fuente
\texorpdfstring{\softcode{scripts/recovery.py:87-153}}{scripts/recovery.py:87-153}.

Pausa la API propia, espera coherencia, copia
\texorpdfstring{\softtext{documentos/secretos}}{documentos/secretos} y
dumps, escribe manifiesto y reanuda en finally.

\textbf{Variables locales asignadas:}
\texorpdfstring{\softcode{\_}}{\_}, \texorpdfstring{\softcode{a}}{a},
\texorpdfstring{\softcode{active}}{active},
\texorpdfstring{\softcode{api\_container}}{api\_container},
\texorpdfstring{\softcode{artifacts}}{artifacts},
\texorpdfstring{\softcode{b}}{b},
\texorpdfstring{\softcode{container}}{container},
\texorpdfstring{\softcode{data}}{data},
\texorpdfstring{\softcode{env}}{env},
\texorpdfstring{\softcode{local\_signatures}}{local\_signatures},
\texorpdfstring{\softcode{log}}{log},
\texorpdfstring{\softcode{log\_path}}{log\_path},
\texorpdfstring{\softcode{node}}{node},
\texorpdfstring{\softcode{path}}{path},
\texorpdfstring{\softcode{pending}}{pending},
\texorpdfstring{\softcode{process}}{process},
\texorpdfstring{\softcode{report}}{report},
\texorpdfstring{\softcode{restart}}{restart},
\texorpdfstring{\softcode{roles}}{roles},
\texorpdfstring{\softcode{source\_documents}}{source\_documents},
\texorpdfstring{\softcode{source\_secrets}}{source\_secrets},
\texorpdfstring{\softcode{started}}{started}.

\subsection{restore}\label{restore}

\textbf{Firma:}
\texorpdfstring{\softcode{restore(destination)}}{restore(destination)}.
Fuente
\texorpdfstring{\softcode{scripts/recovery.py:157-211}}{scripts/recovery.py:157-211}.

Crea un destino aislado, restaura
\texorpdfstring{\softtext{hashes/datos/roles}}{hashes/datos/roles} y
reconstruye FDW y suscripciones sólo hacia el destino.

\textbf{Variables locales asignadas:}
\texorpdfstring{\softcode{command}}{command},
\texorpdfstring{\softcode{compose}}{compose},
\texorpdfstring{\softcode{config}}{config},
\texorpdfstring{\softcode{container}}{container},
\texorpdfstring{\softcode{digest}}{digest},
\texorpdfstring{\softcode{env}}{env},
\texorpdfstring{\softcode{expected}}{expected},
\texorpdfstring{\softcode{name}}{name},
\texorpdfstring{\softcode{node}}{node},
\texorpdfstring{\softcode{port}}{port},
\texorpdfstring{\softcode{project}}{project},
\texorpdfstring{\softcode{records}}{records},
\texorpdfstring{\softcode{report}}{report},
\texorpdfstring{\softcode{roles}}{roles},
\texorpdfstring{\softcode{service}}{service},
\texorpdfstring{\softcode{started}}{started},
\texorpdfstring{\softcode{work}}{work}.

\subsection{\texorpdfstring{\texorpdfstring{\softtext{verify\_http}}{verify\_http}}{}}\label{verify_http}

\textbf{Firma:}
\texorpdfstring{\softcode{verify\_http(destination)}}{verify\_http(destination)}.
Fuente
\texorpdfstring{\softcode{scripts/recovery.py:214-271}}{scripts/recovery.py:214-271}.

API restaurada en el puerto exclusivo18180; jamás usa el puerto
\texorpdfstring{\softtext{fuente18080.}}{fuente18080.}

\textbf{Variables locales asignadas:}
\texorpdfstring{\softcode{animal}}{animal},
\texorpdfstring{\softcode{base}}{base},
\texorpdfstring{\softcode{confirmed}}{confirmed},
\texorpdfstring{\softcode{csrf}}{csrf},
\texorpdfstring{\softcode{data}}{data},
\texorpdfstring{\softcode{detail}}{detail},
\texorpdfstring{\softcode{document}}{document},
\texorpdfstring{\softcode{documents}}{documents},
\texorpdfstring{\softcode{etag}}{etag},
\texorpdfstring{\softcode{headers}}{headers},
\texorpdfstring{\softcode{inventory}}{inventory},
\texorpdfstring{\softcode{key}}{key},
\texorpdfstring{\softcode{lookup}}{lookup},
\texorpdfstring{\softcode{opener}}{opener},
\texorpdfstring{\softcode{operation}}{operation},
\texorpdfstring{\softcode{party}}{party},
\texorpdfstring{\softcode{party\_body}}{party\_body},
\texorpdfstring{\softcode{path}}{path},
\texorpdfstring{\softcode{period}}{period},
\texorpdfstring{\softcode{periods}}{periods},
\texorpdfstring{\softcode{raw}}{raw},
\texorpdfstring{\softcode{recognition}}{recognition},
\texorpdfstring{\softcode{report}}{report},
\texorpdfstring{\softcode{session}}{session},
\texorpdfstring{\softcode{site}}{site},
\texorpdfstring{\softcode{today}}{today}.

\subsection{main}\label{main-3}

\textbf{Firma:} \texorpdfstring{\softcode{main()}}{main()}. Fuente
\texorpdfstring{\softcode{scripts/recovery.py:274-301}}{scripts/recovery.py:274-301}.

Entrada CLI que interpreta las opciones y coordina las funciones del
archivo; su firma no concede permisos
\texorpdfstring{\softtext{comerciales.}}{comerciales.}

\textbf{Variables locales asignadas:}
\texorpdfstring{\softcode{args}}{args},
\texorpdfstring{\softcode{parser}}{parser},
\texorpdfstring{\softcode{path}}{path}.

\subsection{call}\label{call}

\textbf{Firma:}
\texorpdfstring{\softcode{call(method, path, payload=None, key=None, match=None, binary=False)}}{call(method, path, payload=None, key=None, match=None, binary=False)}.
Fuente
\texorpdfstring{\softcode{scripts/recovery.py:221-232}}{scripts/recovery.py:221-232}.

Helper HTTP local que conserva
\texorpdfstring{\softtext{cookies/CSRF}}{cookies/CSRF} y comprueba el
resultado o el contrato según el script.

\textbf{Variables locales asignadas:}
\texorpdfstring{\softcode{csrf}}{csrf},
\texorpdfstring{\softcode{data}}{data},
\texorpdfstring{\softcode{headers}}{headers},
\texorpdfstring{\softcode{raw}}{raw}.

\section{\texorpdfstring{\texorpdfstring{\softtext{scripts/sync-task-tables.py}}{scripts/sync-task-tables.py}}{}}\label{scriptssync-task-tables.py}

Genera las \texorpdfstring{\softtext{tablas/listas}}{tablas/listas} del
README desde el backlog canónico, sin cambiar estados.

\textbf{Variables de módulo:} \texorpdfstring{\softcode{ROOT}}{ROOT},
\texorpdfstring{\softcode{backlog}}{backlog},
\texorpdfstring{\softcode{sections}}{sections},
\texorpdfstring{\softcode{readme}}{readme},
\texorpdfstring{\softcode{text}}{text},
\texorpdfstring{\softcode{start}}{start},
\texorpdfstring{\softcode{end}}{end},
\texorpdfstring{\softcode{text}}{text}.

\section{\texorpdfstring{\texorpdfstring{\softtext{tools/esquema/comparar\_catalogo.py}}{tools/esquema/comparar\_catalogo.py}}{}}\label{toolsesquemacomparar_catalogo.py}

Contrasta columnas, tipos, nullabilidad, PK, UK declaradas y FK lógicas
contra PostgreSQL propio.

\textbf{Variables de módulo:} \texorpdfstring{\softcode{SQL}}{SQL}.

\subsection{qualified}\label{qualified}

\textbf{Firma:}
\texorpdfstring{\softcode{qualified(name)}}{qualified(name)}. Fuente
\texorpdfstring{\softcode{tools/esquema/comparar\_catalogo.py:23-24}}{tools/esquema/comparar\_catalogo.py:23-24}.

Resuelve el nombre del objeto con su esquema para compararlo sin
\texorpdfstring{\softtext{ambigüedad.}}{ambigüedad.}

\subsection{\texorpdfstring{\texorpdfstring{\softtext{normalize\_type}}{normalize\_type}}{}}\label{normalize_type}

\textbf{Firma:}
\texorpdfstring{\softcode{normalize\_type(value)}}{normalize\_type(value)}.
Fuente
\texorpdfstring{\softcode{tools/esquema/comparar\_catalogo.py:27-29}}{tools/esquema/comparar\_catalogo.py:27-29}.

Normaliza representaciones del catálogo para contrastarlas con el modelo
lógico.

\textbf{Variables locales asignadas:}
\texorpdfstring{\softcode{value}}{value}.

\subsection{compare}\label{compare}

\textbf{Firma:}
\texorpdfstring{\softcode{compare(tables, catalog)}}{compare(tables, catalog)}.
Fuente
\texorpdfstring{\softcode{tools/esquema/comparar\_catalogo.py:32-61}}{tools/esquema/comparar\_catalogo.py:32-61}.

Compara
\texorpdfstring{\softtext{columnas/tipos/nullabilidad}}{columnas/tipos/nullabilidad}
y claves con el catálogo; no demuestra todas las RN.

\textbf{Variables locales asignadas:}
\texorpdfstring{\softcode{\_}}{\_},
\texorpdfstring{\softcode{actual}}{actual},
\texorpdfstring{\softcode{actual\_columns}}{actual\_columns},
\texorpdfstring{\softcode{actual\_fks}}{actual\_fks},
\texorpdfstring{\softcode{actual\_keys}}{actual\_keys},
\texorpdfstring{\softcode{column}}{column},
\texorpdfstring{\softcode{columns}}{columns},
\texorpdfstring{\softcode{errors}}{errors},
\texorpdfstring{\softcode{expected}}{expected},
\texorpdfstring{\softcode{expected\_fks}}{expected\_fks},
\texorpdfstring{\softcode{foreign\_keys}}{foreign\_keys},
\texorpdfstring{\softcode{key}}{key},
\texorpdfstring{\softcode{keys}}{keys},
\texorpdfstring{\softcode{name}}{name},
\texorpdfstring{\softcode{primary}}{primary},
\texorpdfstring{\softcode{table}}{table},
\texorpdfstring{\softcode{unique}}{unique},
\texorpdfstring{\softcode{unique\_keys}}{unique\_keys}.

\subsection{main}\label{main-4}

\textbf{Firma:} \texorpdfstring{\softcode{main()}}{main()}. Fuente
\texorpdfstring{\softcode{tools/esquema/comparar\_catalogo.py:64-86}}{tools/esquema/comparar\_catalogo.py:64-86}.

Entrada CLI que interpreta las opciones y coordina las funciones del
archivo; su firma no concede permisos
\texorpdfstring{\softtext{comerciales.}}{comerciales.}

\textbf{Variables locales asignadas:}
\texorpdfstring{\softcode{args}}{args},
\texorpdfstring{\softcode{catalog}}{catalog},
\texorpdfstring{\softcode{errors}}{errors},
\texorpdfstring{\softcode{labels}}{labels},
\texorpdfstring{\softcode{parser}}{parser},
\texorpdfstring{\softcode{raw}}{raw},
\texorpdfstring{\softcode{table}}{table}.

\section{\texorpdfstring{\texorpdfstring{\softtext{tools/dev/check\_api\_contract.py}}{tools/dev/check\_api\_contract.py}}{}}\label{toolsdevcheck_api_contract.py}

Valida tráfico DEMO real contra OpenAPI; no guarda payloads ni cookies.

\textbf{Variables de módulo:} \texorpdfstring{\softcode{ROOT}}{ROOT}.

\subsection{main}\label{main-5}

\textbf{Firma:} \texorpdfstring{\softcode{main()}}{main()}. Fuente
\texorpdfstring{\softcode{tools/dev/check\_api\_contract.py:17-72}}{tools/dev/check\_api\_contract.py:17-72}.

Entrada CLI que interpreta las opciones y coordina las funciones del
archivo; su firma no concede permisos
\texorpdfstring{\softtext{comerciales.}}{comerciales.}

\textbf{Variables locales asignadas:}
\texorpdfstring{\softcode{args}}{args},
\texorpdfstring{\softcode{content}}{content},
\texorpdfstring{\softcode{contract}}{contract},
\texorpdfstring{\softcode{csrf}}{csrf},
\texorpdfstring{\softcode{errors}}{errors},
\texorpdfstring{\softcode{headers}}{headers},
\texorpdfstring{\softcode{media}}{media},
\texorpdfstring{\softcode{op}}{op},
\texorpdfstring{\softcode{opener}}{opener},
\texorpdfstring{\softcode{parser}}{parser},
\texorpdfstring{\softcode{payload}}{payload},
\texorpdfstring{\softcode{raw}}{raw},
\texorpdfstring{\softcode{request}}{request},
\texorpdfstring{\softcode{resolver}}{resolver},
\texorpdfstring{\softcode{response}}{response},
\texorpdfstring{\softcode{results}}{results},
\texorpdfstring{\softcode{schema}}{schema},
\texorpdfstring{\softcode{secret}}{secret},
\texorpdfstring{\softcode{site}}{site},
\texorpdfstring{\softcode{spec}}{spec},
\texorpdfstring{\softcode{suffix}}{suffix},
\texorpdfstring{\softcode{url}}{url},
\texorpdfstring{\softcode{validator}}{validator}.

\subsection{call}\label{call-1}

\textbf{Firma:}
\texorpdfstring{\softcode{call(method, path, template=None, body=None, expected=200)}}{call(method, path, template=None, body=None, expected=200)}.
Fuente
\texorpdfstring{\softcode{tools/dev/check\_api\_contract.py:30-57}}{tools/dev/check\_api\_contract.py:30-57}.

Helper HTTP local que conserva
\texorpdfstring{\softtext{cookies/CSRF}}{cookies/CSRF} y comprueba el
resultado o el contrato según el script.

\textbf{Variables locales asignadas:}
\texorpdfstring{\softcode{content}}{content},
\texorpdfstring{\softcode{errors}}{errors},
\texorpdfstring{\softcode{headers}}{headers},
\texorpdfstring{\softcode{media}}{media},
\texorpdfstring{\softcode{op}}{op},
\texorpdfstring{\softcode{payload}}{payload},
\texorpdfstring{\softcode{raw}}{raw},
\texorpdfstring{\softcode{request}}{request},
\texorpdfstring{\softcode{response}}{response},
\texorpdfstring{\softcode{schema}}{schema},
\texorpdfstring{\softcode{spec}}{spec},
\texorpdfstring{\softcode{validator}}{validator}.

\section{\texorpdfstring{\texorpdfstring{\softtext{db/bootstrap/bootstrap.py}}{db/bootstrap/bootstrap.py}}{}}\label{dbbootstrapbootstrap.py}

Bootstrap aislado, repetible; nunca opera contenedores fuera de
\texorpdfstring{\softtext{bovicierre-academic.}}{bovicierre-academic.}

\textbf{Variables de módulo:} \texorpdfstring{\softcode{ROOT}}{ROOT},
\texorpdfstring{\softcode{SECRET}}{SECRET},
\texorpdfstring{\softcode{PROJECT}}{PROJECT},
\texorpdfstring{\softcode{DATABASE}}{DATABASE},
\texorpdfstring{\softcode{E}}{E}, \texorpdfstring{\softcode{A}}{A},
\texorpdfstring{\softcode{B}}{B}, \texorpdfstring{\softcode{U}}{U},
\texorpdfstring{\softcode{masters}}{masters},
\texorpdfstring{\softcode{local}}{local},
\texorpdfstring{\softcode{restore\_converged}}{restore\_converged},
\texorpdfstring{\softcode{seed}}{seed},
\texorpdfstring{\softcode{provider}}{provider},
\texorpdfstring{\softcode{fdw}}{fdw}.

\subsection{lit}\label{lit}

\textbf{Firma:} \texorpdfstring{\softcode{lit(v)}}{lit(v)}. Fuente
\texorpdfstring{\softcode{db/bootstrap/bootstrap.py:13-13}}{db/bootstrap/bootstrap.py:13-13}.

Escapa literal SQL usado en bootstrap controlado; no es el mecanismo
HTTP de filtros públicos.

\subsection{sql}\label{sql-1}

\textbf{Firma:}
\texorpdfstring{\softcode{sql(node, query)}}{sql(node, query)}. Fuente
\texorpdfstring{\softcode{db/bootstrap/bootstrap.py:14-21}}{db/bootstrap/bootstrap.py:14-21}.

Ejecuta SQL sobre un nodo identificado; en recuperación verifica la
etiqueta del contenedor antes de
\texorpdfstring{\softtext{utilizarlo.}}{utilizarlo.}

\textbf{Variables locales asignadas:}
\texorpdfstring{\softcode{error}}{error},
\texorpdfstring{\softcode{file}}{file},
\texorpdfstring{\softcode{p}}{p}.

\subsection{secret}\label{secret}

\textbf{Firma:} \texorpdfstring{\softcode{secret(n)}}{secret(n)}. Fuente
\texorpdfstring{\softcode{db/bootstrap/bootstrap.py:22-22}}{db/bootstrap/bootstrap.py:22-22}.

Lee el archivo privado requerido por bootstrap; su valor se excluye de
este manual.

\section{\texorpdfstring{\texorpdfstring{\softtext{db/seed/demo.py}}{db/seed/demo.py}}{}}\label{dbseeddemo.py}

Semilla pequeña repetible por comandos y claves permanentes; no borra ni
inserta hechos \texorpdfstring{\softtext{financieros.}}{financieros.}

\textbf{Variables de módulo:} \texorpdfstring{\softcode{ROOT}}{ROOT},
\texorpdfstring{\softcode{SECRETS}}{SECRETS},
\texorpdfstring{\softcode{fixture}}{fixture},
\texorpdfstring{\softcode{dataset}}{dataset},
\texorpdfstring{\softcode{day}}{day},
\texorpdfstring{\softcode{label}}{label},
\texorpdfstring{\softcode{e}}{e}, \texorpdfstring{\softcode{s}}{s},
\texorpdfstring{\softcode{actor}}{actor}.

\subsection{query}\label{query}

\textbf{Firma:}
\texorpdfstring{\softcode{query(sql, args=())}}{query(sql, args=())}.
Fuente
\texorpdfstring{\softcode{db/seed/demo.py:21-22}}{db/seed/demo.py:21-22}.

Ejecuta una consulta parametrizada de la semilla y recupera su resultado
para comprobar el estado
\texorpdfstring{\softtext{persistido.}}{persistido.}

\textbf{Variables locales asignadas:}
\texorpdfstring{\softcode{row}}{row}.

\subsection{money}\label{money}

\textbf{Firma:} \texorpdfstring{\softcode{money(value)}}{money(value)}.
Fuente
\texorpdfstring{\softcode{db/seed/demo.py:23-23}}{db/seed/demo.py:23-23}.

Construye el DTO
\texorpdfstring{\softtext{importe/moneda}}{importe/moneda} con cadena
decimal MXN, sin aritmética flotante.

\subsection{command}\label{command}

\textbf{Firma:}
\texorpdfstring{\softcode{command(name, step, body)}}{command(name, step, body)}.
Fuente
\texorpdfstring{\softcode{db/seed/demo.py:24-26}}{db/seed/demo.py:24-26}.

Construye clave y hash del payload y llama
\texorpdfstring{\softtext{cmd\_ejecutar;}}{cmd\_ejecutar;} no inserta
hechos financieros
\texorpdfstring{\softtext{directamente.}}{directamente.}

\textbf{Variables locales asignadas:}
\texorpdfstring{\softcode{digest}}{digest}.

\section{\texorpdfstring{\texorpdfstring{\softtext{db/tests/benchmark.py}}{db/tests/benchmark.py}}{}}\label{dbtestsbenchmark.py}

Comandos DEMO idempotentes y lecturas canónicas QA; jamás modifica main.

\textbf{Variables de módulo:} \texorpdfstring{\softcode{ROOT}}{ROOT},
\texorpdfstring{\softcode{FIXTURE}}{FIXTURE},
\texorpdfstring{\softcode{DB}}{DB},
\texorpdfstring{\softcode{DATASET}}{DATASET},
\texorpdfstring{\softcode{E}}{E}, \texorpdfstring{\softcode{S}}{S},
\texorpdfstring{\softcode{U}}{U},
\texorpdfstring{\softcode{STATE}}{STATE},
\texorpdfstring{\softcode{CALLS}}{CALLS},
\texorpdfstring{\softcode{FINANCIAL}}{FINANCIAL}.

\subsection{connect}\label{connect}

\textbf{Firma:}
\texorpdfstring{\softcode{connect(pool=False, readonly=False)}}{connect(pool=False, readonly=False)}.
Fuente
\texorpdfstring{\softcode{db/tests/benchmark.py:32-43}}{db/tests/benchmark.py:32-43}.

Abre una conexión psycopg A con timeouts; puede pasar por Pgpool y
limitarse a lectura para medir.

\textbf{Variables locales asignadas:}
\texorpdfstring{\softcode{options}}{options},
\texorpdfstring{\softcode{secret}}{secret}.

\subsection{guard}\label{guard}

\textbf{Firma:}
\texorpdfstring{\softcode{guard(connection)}}{guard(connection)}. Fuente
\texorpdfstring{\softcode{db/tests/benchmark.py:46-60}}{db/tests/benchmark.py:46-60}.

Verifica base QA, servidor18,
\texorpdfstring{\softtext{empresaDEMO,}}{empresaDEMO,} nodoA y las
guardas V13 antes del benchmark.

\textbf{Variables locales asignadas:}
\texorpdfstring{\softcode{company}}{company},
\texorpdfstring{\softcode{definition}}{definition}.

\subsection{command}\label{command-1}

\textbf{Firma:}
\texorpdfstring{\softcode{command(connection, name, step, body, financial)}}{command(connection, name, step, body, financial)}.
Fuente
\texorpdfstring{\softcode{db/tests/benchmark.py:63-73}}{db/tests/benchmark.py:63-73}.

Construye clave y hash del payload y llama
\texorpdfstring{\softtext{cmd\_ejecutar;}}{cmd\_ejecutar;} no inserta
hechos financieros
\texorpdfstring{\softtext{directamente.}}{directamente.}

\textbf{Variables locales asignadas:}
\texorpdfstring{\softcode{CALLS}}{CALLS},
\texorpdfstring{\softcode{digest}}{digest},
\texorpdfstring{\softcode{key}}{key},
\texorpdfstring{\softcode{name}}{name},
\texorpdfstring{\softcode{payload}}{payload}.

\subsection{money}\label{money-1}

\textbf{Firma:}
\texorpdfstring{\softcode{money(amount)}}{money(amount)}. Fuente
\texorpdfstring{\softcode{db/tests/benchmark.py:76-77}}{db/tests/benchmark.py:76-77}.

Construye el DTO
\texorpdfstring{\softtext{importe/moneda}}{importe/moneda} con cadena
decimal MXN, sin aritmética flotante.

\subsection{counts}\label{counts}

\textbf{Firma:}
\texorpdfstring{\softcode{counts(connection)}}{counts(connection)}.
Fuente
\texorpdfstring{\softcode{db/tests/benchmark.py:80-84}}{db/tests/benchmark.py:80-84}.

Obtiene conteos de tablas relevantes para comprobar que el replay no
duplicó hechos.

\subsection{prepare}\label{prepare}

\textbf{Firma:} \texorpdfstring{\softcode{prepare()}}{prepare()}. Fuente
\texorpdfstring{\softcode{db/tests/benchmark.py:87-190}}{db/tests/benchmark.py:87-190}.

Crea el dataset moderado QA por comandos canónicos y contrasta esperados
fijos; conserva su \texorpdfstring{\softtext{manifiesto.}}{manifiesto.}

\textbf{Variables locales asignadas:}
\texorpdfstring{\softcode{actual}}{actual},
\texorpdfstring{\softcode{animal}}{animal},
\texorpdfstring{\softcode{before}}{before},
\texorpdfstring{\softcode{expected}}{expected},
\texorpdfstring{\softcode{expense}}{expense},
\texorpdfstring{\softcode{financial}}{financial},
\texorpdfstring{\softcode{index}}{index},
\texorpdfstring{\softcode{key}}{key},
\texorpdfstring{\softcode{party}}{party},
\texorpdfstring{\softcode{payment}}{payment},
\texorpdfstring{\softcode{payments}}{payments},
\texorpdfstring{\softcode{period}}{period},
\texorpdfstring{\softcode{position}}{position},
\texorpdfstring{\softcode{purchase}}{purchase},
\texorpdfstring{\softcode{purchases}}{purchases},
\texorpdfstring{\softcode{result}}{result},
\texorpdfstring{\softcode{sale}}{sale},
\texorpdfstring{\softcode{sales}}{sales},
\texorpdfstring{\softcode{schema}}{schema},
\texorpdfstring{\softcode{settlement}}{settlement},
\texorpdfstring{\softcode{settlements}}{settlements},
\texorpdfstring{\softcode{started}}{started},
\texorpdfstring{\softcode{state}}{state},
\texorpdfstring{\softcode{tag}}{tag}.

\subsection{utc}\label{utc}

\textbf{Firma:} \texorpdfstring{\softcode{utc()}}{utc()}. Fuente
\texorpdfstring{\softcode{db/tests/benchmark.py:193-194}}{db/tests/benchmark.py:193-194}.

Devuelve el instante actual UTC en formato ISO8601 para fechar
evidencia.

\subsection{revision}\label{revision}

\textbf{Firma:} \texorpdfstring{\softcode{revision()}}{revision()}.
Fuente
\texorpdfstring{\softcode{db/tests/benchmark.py:197-201}}{db/tests/benchmark.py:197-201}.

Registra commit y estado dirty del checkout en el momento de la
\texorpdfstring{\softtext{ejecución.}}{ejecución.}

\subsection{machine}\label{machine}

\textbf{Firma:} \texorpdfstring{\softcode{machine()}}{machine()}. Fuente
\texorpdfstring{\softcode{db/tests/benchmark.py:204-212}}{db/tests/benchmark.py:204-212}.

Lee CPU, memoria y carga del host desde procfs; no estima consumo
exclusivo de \texorpdfstring{\softtext{PostgreSQL.}}{PostgreSQL.}

\textbf{Variables locales asignadas:}
\texorpdfstring{\softcode{cpu}}{cpu},
\texorpdfstring{\softcode{mem}}{mem}.

\subsection{percentiles}\label{percentiles}

\textbf{Firma:}
\texorpdfstring{\softcode{percentiles(samples)}}{percentiles(samples)}.
Fuente
\texorpdfstring{\softcode{db/tests/benchmark.py:215-218}}{db/tests/benchmark.py:215-218}.

Obtiene mediana p50 y p95 por rango
\texorpdfstring{\softtext{ceil(0.95*n),}}{ceil(0.95*n),} además de
\texorpdfstring{\softtext{mínimos/máximos}}{mínimos/máximos} y tamaño de
muestra.

\subsection{params}\label{params}

\textbf{Firma:}
\texorpdfstring{\softcode{params(resource, index, state)}}{params(resource, index, state)}.
Fuente
\texorpdfstring{\softcode{db/tests/benchmark.py:221-227}}{db/tests/benchmark.py:221-227}.

Selecciona filtros, páginas e IDs conocidos para rotar consultas sobre
el dataset \texorpdfstring{\softtext{autorizado.}}{autorizado.}

\subsection{measure}\label{measure}

\textbf{Firma:} \texorpdfstring{\softcode{measure()}}{measure()}. Fuente
\texorpdfstring{\softcode{db/tests/benchmark.py:230-285}}{db/tests/benchmark.py:230-285}.

Calienta y mide consultas canónicas con una/cuatro conexiones, conserva
muestras y planes con límites
\texorpdfstring{\softtext{declarados.}}{declarados.}

\textbf{Variables locales asignadas:}
\texorpdfstring{\softcode{barrier}}{barrier},
\texorpdfstring{\softcode{begin}}{begin},
\texorpdfstring{\softcode{body}}{body},
\texorpdfstring{\softcode{elapsed}}{elapsed},
\texorpdfstring{\softcode{index}}{index},
\texorpdfstring{\softcode{name}}{name},
\texorpdfstring{\softcode{output}}{output},
\texorpdfstring{\softcode{plan}}{plan},
\texorpdfstring{\softcode{query}}{query},
\texorpdfstring{\softcode{report}}{report},
\texorpdfstring{\softcode{resources}}{resources},
\texorpdfstring{\softcode{samples}}{samples},
\texorpdfstring{\softcode{started}}{started},
\texorpdfstring{\softcode{state}}{state},
\texorpdfstring{\softcode{value}}{value},
\texorpdfstring{\softcode{workers}}{workers}.

\subsection{main}\label{main-6}

\textbf{Firma:} \texorpdfstring{\softcode{main()}}{main()}. Fuente
\texorpdfstring{\softcode{db/tests/benchmark.py:288-315}}{db/tests/benchmark.py:288-315}.

Entrada CLI que interpreta las opciones y coordina las funciones del
archivo; su firma no concede permisos
\texorpdfstring{\softtext{comerciales.}}{comerciales.}

\textbf{Variables locales asignadas:}
\texorpdfstring{\softcode{args}}{args},
\texorpdfstring{\softcode{mock}}{mock},
\texorpdfstring{\softcode{mode}}{mode},
\texorpdfstring{\softcode{parser}}{parser},
\texorpdfstring{\softcode{report}}{report},
\texorpdfstring{\softcode{summary}}{summary}.

\subsection{git}\label{git}

\textbf{Firma:} \texorpdfstring{\softcode{git(*args)}}{git(*args)}.
Fuente
\texorpdfstring{\softcode{db/tests/benchmark.py:198-200}}{db/tests/benchmark.py:198-200}.

Consulta \texorpdfstring{\softtext{commit/estado}}{commit/estado}
mediante subprocess sin modificar Git.

\subsection{execute}\label{execute}

\textbf{Firma:}
\texorpdfstring{\softcode{execute(connection, name, index)}}{execute(connection, name, index)}.
Fuente
\texorpdfstring{\softcode{db/tests/benchmark.py:243-249}}{db/tests/benchmark.py:243-249}.

Cronometra una consulta canónica del benchmark y valida su resultado
antes de conservar la muestra.

\textbf{Variables locales asignadas:}
\texorpdfstring{\softcode{begin}}{begin},
\texorpdfstring{\softcode{body}}{body},
\texorpdfstring{\softcode{elapsed}}{elapsed},
\texorpdfstring{\softcode{value}}{value}.

\subsection{worker}\label{worker}

\textbf{Firma:}
\texorpdfstring{\softcode{worker(worker\_id)}}{worker(worker\_id)}.
Fuente
\texorpdfstring{\softcode{db/tests/benchmark.py:262-272}}{db/tests/benchmark.py:262-272}.

Mide con una conexión persistente por participante y una barrera por
recurso; respeta el máximo de cuatro.

\textbf{Variables locales asignadas:}
\texorpdfstring{\softcode{begin}}{begin},
\texorpdfstring{\softcode{index}}{index},
\texorpdfstring{\softcode{name}}{name},
\texorpdfstring{\softcode{output}}{output},
\texorpdfstring{\softcode{samples}}{samples}.

\chapter{Atlas de funciones SQL y declaraciones
locales}\label{atlas-de-funciones-sql-y-declaraciones-locales}

Este atlas conserva la última definición textual de cada función
encontrada en las migraciones ordenadas V1-V13. Es una referencia
estática al código; las migraciones que alteran objetos con bloques
dinámicos deben interpretarse con su explicación y la evidencia de
catálogo. La última definición de
\texorpdfstring{\softcode{cmd\_ejecutar}}{cmd\_ejecutar} está enV13, la
consulta principal y el reporte canónico enV12. Las declaraciones
muestran tipos y valores iniciales públicos del código, nunca una
configuración privada ejecutada.

\section{\texorpdfstring{\texorpdfstring{\softtext{bovi.archivo\_referenciado}}{bovi.archivo\_referenciado}}{}}\label{bovi.archivo_referenciado}

\textbf{Fuente de la definición:}
\texorpdfstring{\softcode{db/migrations/V3\_\_consultas\_sesiones\_grants.sql:150}}{db/migrations/V3\_\_consultas\_sesiones\_grants.sql:150}.

\textbf{Argumentos:}
\texorpdfstring{\softcode{p\_sha256 text}}{p\_sha256 text}.

\textbf{Retorno y propiedades:}
\texorpdfstring{\softcode{boolean LANGUAGE sql STABLE SECURITY DEFINER SET search\_path = pg\_catalog}}{boolean LANGUAGE sql STABLE SECURITY DEFINER SET search\_path = pg\_catalog}.

Para entender sus efectos, consulta la explicación SQL del cuerpo
principal y las RN. El prefijo \texorpdfstring{\softcode{fn\_}}{fn\_}
agrupa auxiliares; una función SECURITY DEFINER requiere los grants
explícitos y el contexto validado. Una llamada no autoriza por sí sola
al usuario humano.

\section{\texorpdfstring{\texorpdfstring{\softtext{bovi.capturar\_tarifas\_proveedor}}{bovi.capturar\_tarifas\_proveedor}}{}}\label{bovi.capturar_tarifas_proveedor}

\textbf{Fuente de la definición:}
\texorpdfstring{\softcode{db/migrations/V9\_\_lector\_a\_captura\_owner\_fecha.sql:17}}{db/migrations/V9\_\_lector\_a\_captura\_owner\_fecha.sql:17}.

\textbf{Argumentos:}
\texorpdfstring{\softcode{e uuid,s uuid,a uuid,pagina\_fuente jsonb}}{e uuid,s uuid,a uuid,pagina\_fuente jsonb}.

\textbf{Retorno y propiedades:}
\texorpdfstring{\softcode{jsonb LANGUAGE plpgsql SECURITY DEFINER SET search\_path=pg\_catalog}}{jsonb LANGUAGE plpgsql SECURITY DEFINER SET search\_path=pg\_catalog}.

\textbf{Variables locales declaradas:}

\begin{itemize}
\tightlist
\item
  \texorpdfstring{\softcode{item jsonb}}{item jsonb}.
\item
  \texorpdfstring{\softcode{snap jsonb}}{snap jsonb}.
\item
  \texorpdfstring{\softcode{sid uuid}}{sid uuid}.
\item
  \texorpdfstring{\softcode{result jsonb:='[]'}}{result jsonb:='[]'}.
\item
  \texorpdfstring{\softcode{source\_price numeric}}{source\_price numeric}.
\end{itemize}

Para entender sus efectos, consulta la explicación SQL del cuerpo
principal y las RN. El prefijo \texorpdfstring{\softcode{fn\_}}{fn\_}
agrupa auxiliares; una función SECURITY DEFINER requiere los grants
explícitos y el contexto validado. Una llamada no autoriza por sí sola
al usuario humano.

\section{\texorpdfstring{\texorpdfstring{\softtext{bovi.cmd\_ejecutar}}{bovi.cmd\_ejecutar}}{}}\label{bovi.cmd_ejecutar}

\textbf{Fuente de la definición:}
\texorpdfstring{\softcode{db/migrations/V13\_\_correccion\_composicion\_exacta.sql:2}}{db/migrations/V13\_\_correccion\_composicion\_exacta.sql:2}.

\textbf{Argumentos:}
\texorpdfstring{\softcode{p\_empresa uuid,p\_sucursal uuid,p\_actor uuid,p\_clave text,p\_hash text,p\_correlation uuid,p\_comando text,p\_payload jsonb}}{p\_empresa uuid,p\_sucursal uuid,p\_actor uuid,p\_clave text,p\_hash text,p\_correlation uuid,p\_comando text,p\_payload jsonb}.

\textbf{Retorno y propiedades:}
\texorpdfstring{\softcode{jsonb LANGUAGE plpgsql SECURITY DEFINER SET search\_path = pg\_catalog}}{jsonb LANGUAGE plpgsql SECURITY DEFINER SET search\_path = pg\_catalog}.

\textbf{Variables locales declaradas:}

\begin{itemize}
\tightlist
\item
  \texorpdfstring{\softcode{e uuid:=p\_empresa}}{e uuid:=p\_empresa}.
\item
  \texorpdfstring{\softcode{s uuid:=p\_sucursal}}{s uuid:=p\_sucursal}.
\item
  \texorpdfstring{\softcode{a uuid:=p\_actor}}{a uuid:=p\_actor}.
\item
  \texorpdfstring{\softcode{j jsonb:=p\_payload}}{j jsonb:=p\_payload}.
\item
  \texorpdfstring{\softcode{n text:=bovi.fn\_nodo()}}{n text:=bovi.fn\_nodo()}.
\item
  \texorpdfstring{\softcode{result jsonb}}{result jsonb}.
\item
  \texorpdfstring{\softcode{prev bovi.clave\_idempotencia\%ROWTYPE}}{prev bovi.clave\_idempotencia\%ROWTYPE}.
\item
  \texorpdfstring{\softcode{op bovi.operacion\%ROWTYPE}}{op bovi.operacion\%ROWTYPE}.
\item
  \texorpdfstring{\softcode{original bovi.pesaje\%ROWTYPE}}{original bovi.pesaje\%ROWTYPE}.
\item
  \texorpdfstring{\softcode{pos bovi.posicion\_inventario\%ROWTYPE}}{pos bovi.posicion\_inventario\%ROWTYPE}.
\item
  \texorpdfstring{\softcode{pay bovi.pago\%ROWTYPE}}{pay bovi.pago\%ROWTYPE}.
\item
  \texorpdfstring{\softcode{liq bovi.liquidacion\%ROWTYPE}}{liq bovi.liquidacion\%ROWTYPE}.
\item
  \texorpdfstring{\softcode{app bovi.aplicacion\_pago\%ROWTYPE}}{app bovi.aplicacion\_pago\%ROWTYPE}.
\item
  \texorpdfstring{\softcode{gasto bovi.gasto\_operacion\%ROWTYPE}}{gasto bovi.gasto\_operacion\%ROWTYPE}.
\item
  \texorpdfstring{\softcode{r record}}{r record}.
\item
  \texorpdfstring{\softcode{c record}}{c record}.
\item
  \texorpdfstring{\softcode{lot jsonb}}{lot jsonb}.
\item
  \texorpdfstring{\softcode{x jsonb}}{x jsonb}.
\item
  \texorpdfstring{\softcode{idn uuid:=uuidv7()}}{idn uuid:=uuidv7()}.
\item
  \texorpdfstring{\softcode{opid uuid:=coalesce((j->>'operacion\_id')::uuid,(j->>'id')::uuid)}}{opid uuid:=coalesce((j->>'operacion\_id')::uuid,(j->>'id')::uuid)}.
\item
  \texorpdfstring{\softcode{pid uuid}}{pid uuid}.
\item
  \texorpdfstring{\softcode{lid uuid}}{lid uuid}.
\item
  \texorpdfstring{\softcode{wid uuid}}{wid uuid}.
\item
  \texorpdfstring{\softcode{aid uuid}}{aid uuid}.
\item
  \texorpdfstring{\softcode{per uuid}}{per uuid}.
\item
  \texorpdfstring{\softcode{amount numeric}}{amount numeric}.
\item
  \texorpdfstring{\softcode{totaln numeric}}{totaln numeric}.
\item
  \texorpdfstring{\softcode{cost numeric}}{cost numeric}.
\item
  \texorpdfstring{\softcode{cnt integer}}{cnt integer}.
\item
  \texorpdfstring{\softcode{k integer}}{k integer}.
\item
  \texorpdfstring{\softcode{idx integer:=0}}{idx integer:=0}.
\item
  \texorpdfstring{\softcode{seqn integer}}{seqn integer}.
\item
  \texorpdfstring{\softcode{contractn smallint}}{contractn smallint}.
\item
  \texorpdfstring{\softcode{preserve boolean:=false}}{preserve boolean:=false}.
\item
  \texorpdfstring{\softcode{clas text}}{clas text}.
\item
  \texorpdfstring{\softcode{nature text}}{nature text}.
\item
  \texorpdfstring{\softcode{f date:=coalesce((j->>'fecha\_economica')::date,(j->>'fecha\_operacion')::date,current\_date)}}{f date:=coalesce((j->>'fecha\_economica')::date,(j->>'fecha\_operacion')::date,current\_date)}.
\item
  \texorpdfstring{\softcode{ref uuid}}{ref uuid}.
\item
  \texorpdfstring{\softcode{ids uuid[]}}{ids uuid[]}.
\item
  \texorpdfstring{\softcode{state text}}{state text}.
\end{itemize}

\textbf{Ramas de despacho presentes en esta definición:}
\texorpdfstring{\softcode{aceptar\_tarifa}}{aceptar\_tarifa},
\texorpdfstring{\softcode{crear\_usuario}}{crear\_usuario},
\texorpdfstring{\softcode{actualizar\_usuario}}{actualizar\_usuario},
\texorpdfstring{\softcode{asignar\_membresia}}{asignar\_membresia},
\texorpdfstring{\softcode{crear\_contraparte}}{crear\_contraparte},
\texorpdfstring{\softcode{crear\_predio}}{crear\_predio},
\texorpdfstring{\softcode{registrar\_animal}}{registrar\_animal},
\texorpdfstring{\softcode{COMPRA}}{COMPRA},
\texorpdfstring{\softcode{acordar\_operacion}}{acordar\_operacion},
\texorpdfstring{\softcode{aprobar\_pesaje}}{aprobar\_pesaje},
\texorpdfstring{\softcode{COMPRA}}{COMPRA},
\texorpdfstring{\softcode{registrar\_pago}}{registrar\_pago},
\texorpdfstring{\softcode{validar\_pago}}{validar\_pago},
\texorpdfstring{\softcode{aplicar\_pago}}{aplicar\_pago},
\texorpdfstring{\softcode{revertir\_aplicacion}}{revertir\_aplicacion},
\texorpdfstring{\softcode{reembolsar\_pago}}{reembolsar\_pago},
\texorpdfstring{\softcode{COBRO}}{COBRO},
\texorpdfstring{\softcode{ajustar\_liquidacion}}{ajustar\_liquidacion},
\texorpdfstring{\softcode{registrar\_gasto}}{registrar\_gasto},
\texorpdfstring{\softcode{crear\_periodo}}{crear\_periodo},
\texorpdfstring{\softcode{registrar\_documento}}{registrar\_documento},
\texorpdfstring{\softcode{registrar\_revision}}{registrar\_revision},
\texorpdfstring{\softcode{cancelar\_operacion}}{cancelar\_operacion},
\texorpdfstring{\softcode{anular\_venta}}{anular\_venta}.

Para entender sus efectos, consulta la explicación SQL del cuerpo
principal y las RN. El prefijo \texorpdfstring{\softcode{fn\_}}{fn\_}
agrupa auxiliares; una función SECURITY DEFINER requiere los grants
explícitos y el contexto validado. Una llamada no autoriza por sí sola
al usuario humano.

\section{\texorpdfstring{\texorpdfstring{\softtext{bovi.confirmar\_entrega}}{bovi.confirmar\_entrega}}{}}\label{bovi.confirmar_entrega}

\textbf{Fuente de la definición:}
\texorpdfstring{\softcode{db/migrations/V3\_\_consultas\_sesiones\_grants.sql:140}}{db/migrations/V3\_\_consultas\_sesiones\_grants.sql:140}.

\textbf{Argumentos:}
\texorpdfstring{\softcode{eventid uuid,worker uuid}}{eventid uuid,worker uuid}.

\textbf{Retorno y propiedades:}
\texorpdfstring{\softcode{void LANGUAGE plpgsql SECURITY DEFINER SET search\_path = pg\_catalog}}{void LANGUAGE plpgsql SECURITY DEFINER SET search\_path = pg\_catalog}.

Para entender sus efectos, consulta la explicación SQL del cuerpo
principal y las RN. El prefijo \texorpdfstring{\softcode{fn\_}}{fn\_}
agrupa auxiliares; una función SECURITY DEFINER requiere los grants
explícitos y el contexto validado. Una llamada no autoriza por sí sola
al usuario humano.

\section{\texorpdfstring{\texorpdfstring{\softtext{bovi.consultar}}{bovi.consultar}}{}}\label{bovi.consultar}

\textbf{Fuente de la definición:}
\texorpdfstring{\softcode{db/migrations/V12\_\_pesajes\_reportes\_canonicos.sql:430}}{db/migrations/V12\_\_pesajes\_reportes\_canonicos.sql:430}.

\textbf{Argumentos:}
\texorpdfstring{\softcode{e uuid,s uuid,a uuid,recurso text,params jsonb}}{e uuid,s uuid,a uuid,recurso text,params jsonb}.

\textbf{Retorno y propiedades:}
\texorpdfstring{\softcode{jsonb LANGUAGE plpgsql SECURITY DEFINER SET search\_path = pg\_catalog}}{jsonb LANGUAGE plpgsql SECURITY DEFINER SET search\_path = pg\_catalog}.

\textbf{Variables locales declaradas:}

\begin{itemize}
\tightlist
\item
  \texorpdfstring{\softcode{tab text}}{tab text}.
\item
  \texorpdfstring{\softcode{filtro text:=''}}{filtro text:=''}.
\item
  \texorpdfstring{\softcode{sql text}}{sql text}.
\item
  \texorpdfstring{\softcode{rows jsonb}}{rows jsonb}.
\item
  \texorpdfstring{\softcode{detail jsonb}}{detail jsonb}.
\item
  \texorpdfstring{\softcode{totaln integer}}{totaln integer}.
\item
  \texorpdfstring{\softcode{size integer:=coalesce((params->>'tamano')::integer,25)}}{size integer:=coalesce((params->>'tamano')::integer,25)}.
\item
  \texorpdfstring{\softcode{page integer:=coalesce((params->>'pagina')::integer,1)}}{page integer:=coalesce((params->>'pagina')::integer,1)}.
\item
  \texorpdfstring{\softcode{idn uuid:=(params->>'id')::uuid}}{idn uuid:=(params->>'id')::uuid}.
\item
  \texorpdfstring{\softcode{searchq text:=nullif(btrim(params->>'q'),'')}}{searchq text:=nullif(btrim(params->>'q'),'')}.
\item
  \texorpdfstring{\softcode{pattern text}}{pattern text}.
\item
  \texorpdfstring{\softcode{statefilter text:=nullif(params->>'estado','')}}{statefilter text:=nullif(params->>'estado','')}.
\item
  \texorpdfstring{\softcode{typefilter text:=nullif(params->>'tipo','')}}{typefilter text:=nullif(params->>'tipo','')}.
\end{itemize}

Para entender sus efectos, consulta la explicación SQL del cuerpo
principal y las RN. El prefijo \texorpdfstring{\softcode{fn\_}}{fn\_}
agrupa auxiliares; una función SECURITY DEFINER requiere los grants
explícitos y el contexto validado. Una llamada no autoriza por sí sola
al usuario humano.

\section{\texorpdfstring{\texorpdfstring{\softtext{bovi.consultar\_tarifas\_proveedor}}{bovi.consultar\_tarifas\_proveedor}}{}}\label{bovi.consultar_tarifas_proveedor}

\textbf{Fuente de la definición:}
\texorpdfstring{\softcode{db/migrations/V9\_\_lector\_a\_captura\_owner\_fecha.sql:35}}{db/migrations/V9\_\_lector\_a\_captura\_owner\_fecha.sql:35}.

\textbf{Argumentos:}
\texorpdfstring{\softcode{e uuid,s uuid,a uuid,params jsonb}}{e uuid,s uuid,a uuid,params jsonb}.

\textbf{Retorno y propiedades:}
\texorpdfstring{\softcode{jsonb LANGUAGE plpgsql SECURITY DEFINER SET search\_path=pg\_catalog}}{jsonb LANGUAGE plpgsql SECURITY DEFINER SET search\_path=pg\_catalog}.

Para entender sus efectos, consulta la explicación SQL del cuerpo
principal y las RN. El prefijo \texorpdfstring{\softcode{fn\_}}{fn\_}
agrupa auxiliares; una función SECURITY DEFINER requiere los grants
explícitos y el contexto validado. Una llamada no autoriza por sí sola
al usuario humano.

\section{\texorpdfstring{\texorpdfstring{\softtext{bovi.consumir\_muerte}}{bovi.consumir\_muerte}}{}}\label{bovi.consumir_muerte}

\textbf{Fuente de la definición:}
\texorpdfstring{\softcode{db/migrations/V3\_\_consultas\_sesiones\_grants.sql:121}}{db/migrations/V3\_\_consultas\_sesiones\_grants.sql:121}.

\textbf{Argumentos:}
\texorpdfstring{\softcode{e uuid,s uuid,eventid uuid,payload jsonb}}{e uuid,s uuid,eventid uuid,payload jsonb}.

\textbf{Retorno y propiedades:}
\texorpdfstring{\softcode{jsonb LANGUAGE plpgsql SECURITY DEFINER SET search\_path = pg\_catalog}}{jsonb LANGUAGE plpgsql SECURITY DEFINER SET search\_path = pg\_catalog}.

Para entender sus efectos, consulta la explicación SQL del cuerpo
principal y las RN. El prefijo \texorpdfstring{\softcode{fn\_}}{fn\_}
agrupa auxiliares; una función SECURITY DEFINER requiere los grants
explícitos y el contexto validado. Una llamada no autoriza por sí sola
al usuario humano.

\section{\texorpdfstring{\texorpdfstring{\softtext{bovi.despachar\_outbox}}{bovi.despachar\_outbox}}{}}\label{bovi.despachar_outbox}

\textbf{Fuente de la definición:}
\texorpdfstring{\softcode{db/migrations/V3\_\_consultas\_sesiones\_grants.sql:113}}{db/migrations/V3\_\_consultas\_sesiones\_grants.sql:113}.

\textbf{Argumentos:}
\texorpdfstring{\softcode{p\_limite integer DEFAULT 50}}{p\_limite integer DEFAULT 50}.

\textbf{Retorno y propiedades:}
\texorpdfstring{\softcode{jsonb LANGUAGE plpgsql SECURITY DEFINER SET search\_path = pg\_catalog}}{jsonb LANGUAGE plpgsql SECURITY DEFINER SET search\_path = pg\_catalog}.

\textbf{Variables locales declaradas:}

\begin{itemize}
\tightlist
\item
  \texorpdfstring{\softcode{r record}}{r record}.
\item
  \texorpdfstring{\softcode{cnt integer:=0}}{cnt integer:=0}.
\end{itemize}

Para entender sus efectos, consulta la explicación SQL del cuerpo
principal y las RN. El prefijo \texorpdfstring{\softcode{fn\_}}{fn\_}
agrupa auxiliares; una función SECURITY DEFINER requiere los grants
explícitos y el contexto validado. Una llamada no autoriza por sí sola
al usuario humano.

\section{\texorpdfstring{\texorpdfstring{\softtext{bovi.fallar\_entrega}}{bovi.fallar\_entrega}}{}}\label{bovi.fallar_entrega}

\textbf{Fuente de la definición:}
\texorpdfstring{\softcode{db/migrations/V3\_\_consultas\_sesiones\_grants.sql:145}}{db/migrations/V3\_\_consultas\_sesiones\_grants.sql:145}.

\textbf{Argumentos:}
\texorpdfstring{\softcode{eventid uuid,worker uuid}}{eventid uuid,worker uuid}.

\textbf{Retorno y propiedades:}
\texorpdfstring{\softcode{void LANGUAGE plpgsql SECURITY DEFINER SET search\_path = pg\_catalog}}{void LANGUAGE plpgsql SECURITY DEFINER SET search\_path = pg\_catalog}.

Para entender sus efectos, consulta la explicación SQL del cuerpo
principal y las RN. El prefijo \texorpdfstring{\softcode{fn\_}}{fn\_}
agrupa auxiliares; una función SECURITY DEFINER requiere los grants
explícitos y el contexto validado. Una llamada no autoriza por sí sola
al usuario humano.

\section{\texorpdfstring{\texorpdfstring{\softtext{bovi.fn\_calcular\_liquidacion}}{bovi.fn\_calcular\_liquidacion}}{}}\label{bovi.fn_calcular_liquidacion}

\textbf{Fuente de la definición:}
\texorpdfstring{\softcode{db/migrations/V2\_\_comandos.sql:70}}{db/migrations/V2\_\_comandos.sql:70}.

\textbf{Argumentos:}
\texorpdfstring{\softcode{e uuid,s uuid,op uuid}}{e uuid,s uuid,op uuid}.

\textbf{Retorno y propiedades:}
\texorpdfstring{\softcode{TABLE(numero\_lote smallint,cantidad numeric,unidad text,precio numeric,base numeric,cargos numeric,bonificaciones numeric,total numeric,pesaje\_id uuid) LANGUAGE sql STABLE SET search\_path = pg\_catalog}}{TABLE(numero\_lote smallint,cantidad numeric,unidad text,precio numeric,base numeric,cargos numeric,bonificaciones numeric,total numeric,pesaje\_id uuid) LANGUAGE sql STABLE SET search\_path = pg\_catalog}.

Para entender sus efectos, consulta la explicación SQL del cuerpo
principal y las RN. El prefijo \texorpdfstring{\softcode{fn\_}}{fn\_}
agrupa auxiliares; una función SECURITY DEFINER requiere los grants
explícitos y el contexto validado. Una llamada no autoriza por sí sola
al usuario humano.

\section{\texorpdfstring{\texorpdfstring{\softtext{bovi.fn\_codigo}}{bovi.fn\_codigo}}{}}\label{bovi.fn_codigo}

\textbf{Fuente de la definición:}
\texorpdfstring{\softcode{db/migrations/V2\_\_comandos.sql:78}}{db/migrations/V2\_\_comandos.sql:78}.

\textbf{Argumentos:} \texorpdfstring{\softcode{c text}}{c text}.

\textbf{Retorno y propiedades:}
\texorpdfstring{\softcode{text LANGUAGE plpgsql IMMUTABLE SET search\_path = pg\_catalog}}{text LANGUAGE plpgsql IMMUTABLE SET search\_path = pg\_catalog}.

\textbf{Variables locales declaradas:}

\begin{itemize}
\tightlist
\item
  \texorpdfstring{\softcode{norm text:=upper(c)}}{norm text:=upper(c)}.
\end{itemize}

Para entender sus efectos, consulta la explicación SQL del cuerpo
principal y las RN. El prefijo \texorpdfstring{\softcode{fn\_}}{fn\_}
agrupa auxiliares; una función SECURITY DEFINER requiere los grants
explícitos y el contexto validado. Una llamada no autoriza por sí sola
al usuario humano.

\section{\texorpdfstring{\texorpdfstring{\softtext{bovi.fn\_json\_decimal}}{bovi.fn\_json\_decimal}}{}}\label{bovi.fn_json_decimal}

\textbf{Fuente de la definición:}
\texorpdfstring{\softcode{db/migrations/V3\_\_consultas\_sesiones\_grants.sql:24}}{db/migrations/V3\_\_consultas\_sesiones\_grants.sql:24}.

\textbf{Argumentos:} \texorpdfstring{\softcode{j jsonb}}{j jsonb}.

\textbf{Retorno y propiedades:}
\texorpdfstring{\softcode{jsonb LANGUAGE plpgsql IMMUTABLE SET search\_path = pg\_catalog}}{jsonb LANGUAGE plpgsql IMMUTABLE SET search\_path = pg\_catalog}.

\textbf{Variables locales declaradas:}

\begin{itemize}
\tightlist
\item
  \texorpdfstring{\softcode{k text}}{k text}.
\item
  \texorpdfstring{\softcode{v jsonb}}{v jsonb}.
\item
  \texorpdfstring{\softcode{result jsonb:='\{\}'}}{result jsonb:='\{\}'}.
\end{itemize}

Para entender sus efectos, consulta la explicación SQL del cuerpo
principal y las RN. El prefijo \texorpdfstring{\softcode{fn\_}}{fn\_}
agrupa auxiliares; una función SECURITY DEFINER requiere los grants
explícitos y el contexto validado. Una llamada no autoriza por sí sola
al usuario humano.

\section{\texorpdfstring{\texorpdfstring{\softtext{bovi.fn\_liquidacion\_json}}{bovi.fn\_liquidacion\_json}}{}}\label{bovi.fn_liquidacion_json}

\textbf{Fuente de la definición:}
\texorpdfstring{\softcode{db/migrations/V3\_\_consultas\_sesiones\_grants.sql:35}}{db/migrations/V3\_\_consultas\_sesiones\_grants.sql:35}.

\textbf{Argumentos:}
\texorpdfstring{\softcode{e uuid,s uuid,lid uuid}}{e uuid,s uuid,lid uuid}.

\textbf{Retorno y propiedades:}
\texorpdfstring{\softcode{jsonb LANGUAGE sql STABLE SET search\_path = pg\_catalog}}{jsonb LANGUAGE sql STABLE SET search\_path = pg\_catalog}.

Para entender sus efectos, consulta la explicación SQL del cuerpo
principal y las RN. El prefijo \texorpdfstring{\softcode{fn\_}}{fn\_}
agrupa auxiliares; una función SECURITY DEFINER requiere los grants
explícitos y el contexto validado. Una llamada no autoriza por sí sola
al usuario humano.

\section{\texorpdfstring{\texorpdfstring{\softtext{bovi.fn\_nodo}}{bovi.fn\_nodo}}{}}\label{bovi.fn_nodo}

\textbf{Fuente de la definición:}
\texorpdfstring{\softcode{db/migrations/V2\_\_comandos.sql:2}}{db/migrations/V2\_\_comandos.sql:2}.

\textbf{Argumentos:} ``.

\textbf{Retorno y propiedades:}
\texorpdfstring{\softcode{text LANGUAGE sql STABLE SET search\_path = pg\_catalog}}{text LANGUAGE sql STABLE SET search\_path = pg\_catalog}.

Para entender sus efectos, consulta la explicación SQL del cuerpo
principal y las RN. El prefijo \texorpdfstring{\softcode{fn\_}}{fn\_}
agrupa auxiliares; una función SECURITY DEFINER requiere los grants
explícitos y el contexto validado. Una llamada no autoriza por sí sola
al usuario humano.

\section{\texorpdfstring{\texorpdfstring{\softtext{bovi.fn\_normalizar\_payload}}{bovi.fn\_normalizar\_payload}}{}}\label{bovi.fn_normalizar_payload}

\textbf{Fuente de la definición:}
\texorpdfstring{\softcode{db/migrations/V2\_\_comandos.sql:81}}{db/migrations/V2\_\_comandos.sql:81}.

\textbf{Argumentos:}
\texorpdfstring{\softcode{cmd text,j jsonb}}{cmd text,j jsonb}.

\textbf{Retorno y propiedades:}
\texorpdfstring{\softcode{jsonb LANGUAGE plpgsql IMMUTABLE SET search\_path = pg\_catalog}}{jsonb LANGUAGE plpgsql IMMUTABLE SET search\_path = pg\_catalog}.

\textbf{Variables locales declaradas:}

\begin{itemize}
\tightlist
\item
  \texorpdfstring{\softcode{outj jsonb:=j}}{outj jsonb:=j}.
\item
  \texorpdfstring{\softcode{lot jsonb}}{lot jsonb}.
\item
  \texorpdfstring{\softcode{lots jsonb:='[]'}}{lots jsonb:='[]'}.
\item
  \texorpdfstring{\softcode{v text}}{v text}.
\end{itemize}

Para entender sus efectos, consulta la explicación SQL del cuerpo
principal y las RN. El prefijo \texorpdfstring{\softcode{fn\_}}{fn\_}
agrupa auxiliares; una función SECURITY DEFINER requiere los grants
explícitos y el contexto validado. Una llamada no autoriza por sí sola
al usuario humano.

\section{\texorpdfstring{\texorpdfstring{\softtext{bovi.fn\_periodo}}{bovi.fn\_periodo}}{}}\label{bovi.fn_periodo}

\textbf{Fuente de la definición:}
\texorpdfstring{\softcode{db/migrations/V2\_\_comandos.sql:55}}{db/migrations/V2\_\_comandos.sql:55}.

\textbf{Argumentos:}
\texorpdfstring{\softcode{e uuid,s uuid,f date,pid uuid}}{e uuid,s uuid,f date,pid uuid}.

\textbf{Retorno y propiedades:}
\texorpdfstring{\softcode{uuid LANGUAGE plpgsql SET search\_path = pg\_catalog}}{uuid LANGUAGE plpgsql SET search\_path = pg\_catalog}.

\textbf{Variables locales declaradas:}

\begin{itemize}
\tightlist
\item
  \texorpdfstring{\softcode{r bovi.periodo\_cierre\%ROWTYPE}}{r bovi.periodo\_cierre\%ROWTYPE}.
\end{itemize}

Para entender sus efectos, consulta la explicación SQL del cuerpo
principal y las RN. El prefijo \texorpdfstring{\softcode{fn\_}}{fn\_}
agrupa auxiliares; una función SECURITY DEFINER requiere los grants
explícitos y el contexto validado. Una llamada no autoriza por sí sola
al usuario humano.

\section{\texorpdfstring{\texorpdfstring{\softtext{bovi.fn\_periodo\_payload}}{bovi.fn\_periodo\_payload}}{}}\label{bovi.fn_periodo_payload}

\textbf{Fuente de la definición:}
\texorpdfstring{\softcode{db/migrations/V2\_\_comandos.sql:63}}{db/migrations/V2\_\_comandos.sql:63}.

\textbf{Argumentos:}
\texorpdfstring{\softcode{e uuid,s uuid,j jsonb}}{e uuid,s uuid,j jsonb}.

\textbf{Retorno y propiedades:}
\texorpdfstring{\softcode{uuid LANGUAGE sql SET search\_path = pg\_catalog}}{uuid LANGUAGE sql SET search\_path = pg\_catalog}.

Para entender sus efectos, consulta la explicación SQL del cuerpo
principal y las RN. El prefijo \texorpdfstring{\softcode{fn\_}}{fn\_}
agrupa auxiliares; una función SECURITY DEFINER requiere los grants
explícitos y el contexto validado. Una llamada no autoriza por sí sola
al usuario humano.

\section{\texorpdfstring{\texorpdfstring{\softtext{bovi.fn\_pesaje\_json}}{bovi.fn\_pesaje\_json}}{}}\label{bovi.fn_pesaje_json}

\textbf{Fuente de la definición:}
\texorpdfstring{\softcode{db/migrations/V12\_\_pesajes\_reportes\_canonicos.sql:2}}{db/migrations/V12\_\_pesajes\_reportes\_canonicos.sql:2}.

\textbf{Argumentos:}
\texorpdfstring{\softcode{e uuid,s uuid,pid uuid}}{e uuid,s uuid,pid uuid}.

\textbf{Retorno y propiedades:}
\texorpdfstring{\softcode{jsonb LANGUAGE sql STABLE SET search\_path=pg\_catalog}}{jsonb LANGUAGE sql STABLE SET search\_path=pg\_catalog}.

Para entender sus efectos, consulta la explicación SQL del cuerpo
principal y las RN. El prefijo \texorpdfstring{\softcode{fn\_}}{fn\_}
agrupa auxiliares; una función SECURITY DEFINER requiere los grants
explícitos y el contexto validado. Una llamada no autoriza por sí sola
al usuario humano.

\section{\texorpdfstring{\texorpdfstring{\softtext{bovi.fn\_repartir}}{bovi.fn\_repartir}}{}}\label{bovi.fn_repartir}

\textbf{Fuente de la definición:}
\texorpdfstring{\softcode{db/migrations/V2\_\_comandos.sql:64}}{db/migrations/V2\_\_comandos.sql:64}.

\textbf{Argumentos:}
\texorpdfstring{\softcode{m numeric,n integer}}{m numeric,n integer}.

\textbf{Retorno y propiedades:}
\texorpdfstring{\softcode{TABLE(numero integer,importe numeric) LANGUAGE sql IMMUTABLE SET search\_path = pg\_catalog}}{TABLE(numero integer,importe numeric) LANGUAGE sql IMMUTABLE SET search\_path = pg\_catalog}.

Para entender sus efectos, consulta la explicación SQL del cuerpo
principal y las RN. El prefijo \texorpdfstring{\softcode{fn\_}}{fn\_}
agrupa auxiliares; una función SECURITY DEFINER requiere los grants
explícitos y el contexto validado. Una llamada no autoriza por sí sola
al usuario humano.

\section{\texorpdfstring{\texorpdfstring{\softtext{bovi.fn\_reparto}}{bovi.fn\_reparto}}{}}\label{bovi.fn_reparto}

\textbf{Fuente de la definición:}
\texorpdfstring{\softcode{db/migrations/V2\_\_comandos.sql:99}}{db/migrations/V2\_\_comandos.sql:99}.

\textbf{Argumentos:}
\texorpdfstring{\softcode{e uuid,s uuid,gid uuid,metodo text,asignaciones jsonb DEFAULT NULL}}{e uuid,s uuid,gid uuid,metodo text,asignaciones jsonb DEFAULT NULL}.

\textbf{Retorno y propiedades:}
\texorpdfstring{\softcode{TABLE(posicion\_id uuid,base\_valor numeric,importe numeric) LANGUAGE plpgsql STABLE SET search\_path = pg\_catalog}}{TABLE(posicion\_id uuid,base\_valor numeric,importe numeric) LANGUAGE plpgsql STABLE SET search\_path = pg\_catalog}.

\textbf{Variables locales declaradas:}

\begin{itemize}
\tightlist
\item
  \texorpdfstring{\softcode{g bovi.gasto\_operacion\%ROWTYPE}}{g bovi.gasto\_operacion\%ROWTYPE}.
\item
  \texorpdfstring{\softcode{ids uuid[]}}{ids uuid[]}.
\item
  \texorpdfstring{\softcode{weights numeric[]}}{weights numeric[]}.
\item
  \texorpdfstring{\softcode{valuesn numeric[]}}{valuesn numeric[]}.
\item
  \texorpdfstring{\softcode{sourceweight numeric}}{sourceweight numeric}.
\item
  \texorpdfstring{\softcode{entry jsonb}}{entry jsonb}.
\item
  \texorpdfstring{\softcode{amount numeric}}{amount numeric}.
\item
  \texorpdfstring{\softcode{totalweight numeric}}{totalweight numeric}.
\item
  \texorpdfstring{\softcode{roundedtotal numeric}}{roundedtotal numeric}.
\item
  \texorpdfstring{\softcode{residual integer}}{residual integer}.
\item
  \texorpdfstring{\softcode{idx integer}}{idx integer}.
\item
  \texorpdfstring{\softcode{signn integer}}{signn integer}.
\end{itemize}

Para entender sus efectos, consulta la explicación SQL del cuerpo
principal y las RN. El prefijo \texorpdfstring{\softcode{fn\_}}{fn\_}
agrupa auxiliares; una función SECURITY DEFINER requiere los grants
explícitos y el contexto validado. Una llamada no autoriza por sí sola
al usuario humano.

\section{\texorpdfstring{\texorpdfstring{\softtext{bovi.fn\_reporte}}{bovi.fn\_reporte}}{}}\label{bovi.fn_reporte}

\textbf{Fuente de la definición:}
\texorpdfstring{\softcode{db/migrations/V12\_\_pesajes\_reportes\_canonicos.sql:12}}{db/migrations/V12\_\_pesajes\_reportes\_canonicos.sql:12}.

\textbf{Argumentos:}
\texorpdfstring{\softcode{e uuid,s uuid,recurso text,params jsonb}}{e uuid,s uuid,recurso text,params jsonb}.

\textbf{Retorno y propiedades:}
\texorpdfstring{\softcode{jsonb LANGUAGE plpgsql STABLE SET search\_path=pg\_catalog}}{jsonb LANGUAGE plpgsql STABLE SET search\_path=pg\_catalog}.

\textbf{Variables locales declaradas:}

\begin{itemize}
\tightlist
\item
  \texorpdfstring{\softcode{metadata jsonb}}{metadata jsonb}.
\item
  \texorpdfstring{\softcode{result jsonb}}{result jsonb}.
\item
  \texorpdfstring{\softcode{pattern text}}{pattern text}.
\item
  \texorpdfstring{\softcode{q text:=nullif(btrim(params->>'q'),'')}}{q text:=nullif(btrim(params->>'q'),'')}.
\item
  \texorpdfstring{\softcode{size integer:=coalesce((params->>'tamano')::integer,25)}}{size integer:=coalesce((params->>'tamano')::integer,25)}.
\item
  \texorpdfstring{\softcode{page integer:=coalesce((params->>'pagina')::integer,1)}}{page integer:=coalesce((params->>'pagina')::integer,1)}.
\item
  \texorpdfstring{\softcode{periodid uuid:=(params->>'periodo\_id')::uuid}}{periodid uuid:=(params->>'periodo\_id')::uuid}.
\item
  \texorpdfstring{\softcode{operationid uuid:=(params->>'operacion\_id')::uuid}}{operationid uuid:=(params->>'operacion\_id')::uuid}.
\item
  \texorpdfstring{\softcode{periodstate text}}{periodstate text}.
\item
  \texorpdfstring{\softcode{faltantes jsonb:='[]'}}{faltantes jsonb:='[]'}.
\item
  \texorpdfstring{\softcode{ownernode text}}{ownernode text}.
\item
  \texorpdfstring{\softcode{closure record}}{closure record}.
\item
  \texorpdfstring{\softcode{today date}}{today date}.
\item
  \texorpdfstring{\softcode{direction text:=params->>'direccion'}}{direction text:=params->>'direccion'}.
\item
  \texorpdfstring{\softcode{overdue boolean}}{overdue boolean}.
\end{itemize}

Para entender sus efectos, consulta la explicación SQL del cuerpo
principal y las RN. El prefijo \texorpdfstring{\softcode{fn\_}}{fn\_}
agrupa auxiliares; una función SECURITY DEFINER requiere los grants
explícitos y el contexto validado. Una llamada no autoriza por sí sola
al usuario humano.

\section{\texorpdfstring{\texorpdfstring{\softtext{bovi.fn\_scope}}{bovi.fn\_scope}}{}}\label{bovi.fn_scope}

\textbf{Fuente de la definición:}
\texorpdfstring{\softcode{db/migrations/V2\_\_comandos.sql:3}}{db/migrations/V2\_\_comandos.sql:3}.

\textbf{Argumentos:}
\texorpdfstring{\softcode{e uuid,s uuid,a uuid,rol text DEFAULT NULL,maestro boolean DEFAULT false}}{e uuid,s uuid,a uuid,rol text DEFAULT NULL,maestro boolean DEFAULT false}.

\textbf{Retorno y propiedades:}
\texorpdfstring{\softcode{void LANGUAGE plpgsql SET search\_path = pg\_catalog}}{void LANGUAGE plpgsql SET search\_path = pg\_catalog}.

Para entender sus efectos, consulta la explicación SQL del cuerpo
principal y las RN. El prefijo \texorpdfstring{\softcode{fn\_}}{fn\_}
agrupa auxiliares; una función SECURITY DEFINER requiere los grants
explícitos y el contexto validado. Una llamada no autoriza por sí sola
al usuario humano.

\section{\texorpdfstring{\texorpdfstring{\softtext{bovi.fn\_validar\_conciliacion\_propietario}}{bovi.fn\_validar\_conciliacion\_propietario}}{}}\label{bovi.fn_validar_conciliacion_propietario}

\textbf{Fuente de la definición:}
\texorpdfstring{\softcode{db/migrations/V3\_\_consultas\_sesiones\_grants.sql:38}}{db/migrations/V3\_\_consultas\_sesiones\_grants.sql:38}.

\textbf{Argumentos:}
\texorpdfstring{\softcode{e uuid,s uuid,a uuid}}{e uuid,s uuid,a uuid}.

\textbf{Retorno y propiedades:}
\texorpdfstring{\softcode{jsonb LANGUAGE plpgsql SECURITY DEFINER SET search\_path = pg\_catalog}}{jsonb LANGUAGE plpgsql SECURITY DEFINER SET search\_path = pg\_catalog}.

\textbf{Variables locales declaradas:}

\begin{itemize}
\tightlist
\item
  \texorpdfstring{\softcode{fk record}}{fk record}.
\item
  \texorpdfstring{\softcode{localcols text}}{localcols text}.
\item
  \texorpdfstring{\softcode{joincols text}}{joincols text}.
\item
  \texorpdfstring{\softcode{nonnull text}}{nonnull text}.
\item
  \texorpdfstring{\softcode{missing bigint}}{missing bigint}.
\item
  \texorpdfstring{\softcode{totalmissing bigint:=0}}{totalmissing bigint:=0}.
\item
  \texorpdfstring{\softcode{errors jsonb:='[]'}}{errors jsonb:='[]'}.
\item
  \texorpdfstring{\softcode{settlements bigint}}{settlements bigint}.
\end{itemize}

Para entender sus efectos, consulta la explicación SQL del cuerpo
principal y las RN. El prefijo \texorpdfstring{\softcode{fn\_}}{fn\_}
agrupa auxiliares; una función SECURITY DEFINER requiere los grants
explícitos y el contexto validado. Una llamada no autoriza por sí sola
al usuario humano.

\section{\texorpdfstring{\texorpdfstring{\softtext{bovi.leer\_tarifas\_proveedor}}{bovi.leer\_tarifas\_proveedor}}{}}\label{bovi.leer_tarifas_proveedor}

\textbf{Fuente de la definición:}
\texorpdfstring{\softcode{db/migrations/V9\_\_lector\_a\_captura\_owner\_fecha.sql:2}}{db/migrations/V9\_\_lector\_a\_captura\_owner\_fecha.sql:2}.

\textbf{Argumentos:}
\texorpdfstring{\softcode{e uuid,s uuid,a uuid,params jsonb}}{e uuid,s uuid,a uuid,params jsonb}.

\textbf{Retorno y propiedades:}
\texorpdfstring{\softcode{jsonb LANGUAGE plpgsql SECURITY DEFINER SET search\_path=pg\_catalog}}{jsonb LANGUAGE plpgsql SECURITY DEFINER SET search\_path=pg\_catalog}.

\textbf{Variables locales declaradas:}

\begin{itemize}
\tightlist
\item
  \texorpdfstring{\softcode{r record}}{r record}.
\item
  \texorpdfstring{\softcode{sid uuid}}{sid uuid}.
\item
  \texorpdfstring{\softcode{snap jsonb}}{snap jsonb}.
\item
  \texorpdfstring{\softcode{result jsonb:='[]'}}{result jsonb:='[]'}.
\item
  \texorpdfstring{\softcode{countn integer:=0}}{countn integer:=0}.
\item
  \texorpdfstring{\softcode{totaln integer}}{totaln integer}.
\item
  \texorpdfstring{\softcode{page integer:=coalesce((params->>'pagina')::integer,1)}}{page integer:=coalesce((params->>'pagina')::integer,1)}.
\item
  \texorpdfstring{\softcode{size integer:=coalesce((params->>'tamano')::integer,25)}}{size integer:=coalesce((params->>'tamano')::integer,25)}.
\end{itemize}

Para entender sus efectos, consulta la explicación SQL del cuerpo
principal y las RN. El prefijo \texorpdfstring{\softcode{fn\_}}{fn\_}
agrupa auxiliares; una función SECURITY DEFINER requiere los grants
explícitos y el contexto validado. Una llamada no autoriza por sí sola
al usuario humano.

\section{\texorpdfstring{\texorpdfstring{\softtext{bovi.reclamar\_entregas}}{bovi.reclamar\_entregas}}{}}\label{bovi.reclamar_entregas}

\textbf{Fuente de la definición:}
\texorpdfstring{\softcode{db/migrations/V3\_\_consultas\_sesiones\_grants.sql:131}}{db/migrations/V3\_\_consultas\_sesiones\_grants.sql:131}.

\textbf{Argumentos:}
\texorpdfstring{\softcode{p\_worker uuid,p\_limite integer DEFAULT 50}}{p\_worker uuid,p\_limite integer DEFAULT 50}.

\textbf{Retorno y propiedades:}
\texorpdfstring{\softcode{jsonb LANGUAGE plpgsql SECURITY DEFINER SET search\_path = pg\_catalog}}{jsonb LANGUAGE plpgsql SECURITY DEFINER SET search\_path = pg\_catalog}.

\textbf{Variables locales declaradas:}

\begin{itemize}
\tightlist
\item
  \texorpdfstring{\softcode{r record}}{r record}.
\item
  \texorpdfstring{\softcode{claimed jsonb:='[]'}}{claimed jsonb:='[]'}.
\end{itemize}

Para entender sus efectos, consulta la explicación SQL del cuerpo
principal y las RN. El prefijo \texorpdfstring{\softcode{fn\_}}{fn\_}
agrupa auxiliares; una función SECURITY DEFINER requiere los grants
explícitos y el contexto validado. Una llamada no autoriza por sí sola
al usuario humano.

\section{\texorpdfstring{\texorpdfstring{\softtext{bovi.trg\_acuerdo\_completo}}{bovi.trg\_acuerdo\_completo}}{}}\label{bovi.trg_acuerdo_completo}

\textbf{Fuente de la definición:}
\texorpdfstring{\softcode{db/migrations/V8\_\_borrador\_kg\_y\_expiracion.sql:5}}{db/migrations/V8\_\_borrador\_kg\_y\_expiracion.sql:5}.

\textbf{Argumentos:} ``.

\textbf{Retorno y propiedades:}
\texorpdfstring{\softcode{trigger LANGUAGE plpgsql SET search\_path=pg\_catalog}}{trigger LANGUAGE plpgsql SET search\_path=pg\_catalog}.

\textbf{Variables locales declaradas:}

\begin{itemize}
\tightlist
\item
  \texorpdfstring{\softcode{ver smallint}}{ver smallint}.
\item
  \texorpdfstring{\softcode{r record}}{r record}.
\item
  \texorpdfstring{\softcode{w bovi.pesaje\%ROWTYPE}}{w bovi.pesaje\%ROWTYPE}.
\end{itemize}

Para entender sus efectos, consulta la explicación SQL del cuerpo
principal y las RN. El prefijo \texorpdfstring{\softcode{fn\_}}{fn\_}
agrupa auxiliares; una función SECURITY DEFINER requiere los grants
explícitos y el contexto validado. Una llamada no autoriza por sí sola
al usuario humano.

\section{\texorpdfstring{\texorpdfstring{\softtext{bovi.trg\_aprobado}}{bovi.trg\_aprobado}}{}}\label{bovi.trg_aprobado}

\textbf{Fuente de la definición:}
\texorpdfstring{\softcode{db/migrations/V2\_\_comandos.sql:29}}{db/migrations/V2\_\_comandos.sql:29}.

\textbf{Argumentos:} ``.

\textbf{Retorno y propiedades:}
\texorpdfstring{\softcode{trigger LANGUAGE plpgsql SET search\_path = pg\_catalog}}{trigger LANGUAGE plpgsql SET search\_path = pg\_catalog}.

Para entender sus efectos, consulta la explicación SQL del cuerpo
principal y las RN. El prefijo \texorpdfstring{\softcode{fn\_}}{fn\_}
agrupa auxiliares; una función SECURITY DEFINER requiere los grants
explícitos y el contexto validado. Una llamada no autoriza por sí sola
al usuario humano.

\section{\texorpdfstring{\texorpdfstring{\softtext{bovi.trg\_borrador}}{bovi.trg\_borrador}}{}}\label{bovi.trg_borrador}

\textbf{Fuente de la definición:}
\texorpdfstring{\softcode{db/migrations/V2\_\_comandos.sql:22}}{db/migrations/V2\_\_comandos.sql:22}.

\textbf{Argumentos:} ``.

\textbf{Retorno y propiedades:}
\texorpdfstring{\softcode{trigger LANGUAGE plpgsql SET search\_path = pg\_catalog}}{trigger LANGUAGE plpgsql SET search\_path = pg\_catalog}.

\textbf{Variables locales declaradas:}

\begin{itemize}
\tightlist
\item
  \texorpdfstring{\softcode{j jsonb:=CASE WHEN TG\_OP='DELETE' THEN to\_jsonb(OLD) ELSE to\_jsonb(NEW) END}}{j jsonb:=CASE WHEN TG\_OP='DELETE' THEN to\_jsonb(OLD) ELSE to\_jsonb(NEW) END}.
\end{itemize}

Para entender sus efectos, consulta la explicación SQL del cuerpo
principal y las RN. El prefijo \texorpdfstring{\softcode{fn\_}}{fn\_}
agrupa auxiliares; una función SECURITY DEFINER requiere los grants
explícitos y el contexto validado. Una llamada no autoriza por sí sola
al usuario humano.

\section{\texorpdfstring{\texorpdfstring{\softtext{bovi.trg\_inmutable}}{bovi.trg\_inmutable}}{}}\label{bovi.trg_inmutable}

\textbf{Fuente de la definición:}
\texorpdfstring{\softcode{db/migrations/V2\_\_comandos.sql:20}}{db/migrations/V2\_\_comandos.sql:20}.

\textbf{Argumentos:} ``.

\textbf{Retorno y propiedades:}
\texorpdfstring{\softcode{trigger LANGUAGE plpgsql SET search\_path = pg\_catalog}}{trigger LANGUAGE plpgsql SET search\_path = pg\_catalog}.

Para entender sus efectos, consulta la explicación SQL del cuerpo
principal y las RN. El prefijo \texorpdfstring{\softcode{fn\_}}{fn\_}
agrupa auxiliares; una función SECURITY DEFINER requiere los grants
explícitos y el contexto validado. Una llamada no autoriza por sí sola
al usuario humano.

\section{\texorpdfstring{\texorpdfstring{\softtext{bovi.trg\_pesaje\_composicion}}{bovi.trg\_pesaje\_composicion}}{}}\label{bovi.trg_pesaje_composicion}

\textbf{Fuente de la definición:}
\texorpdfstring{\softcode{db/migrations/V2\_\_comandos.sql:35}}{db/migrations/V2\_\_comandos.sql:35}.

\textbf{Argumentos:} ``.

\textbf{Retorno y propiedades:}
\texorpdfstring{\softcode{trigger LANGUAGE plpgsql SET search\_path = pg\_catalog}}{trigger LANGUAGE plpgsql SET search\_path = pg\_catalog}.

\textbf{Variables locales declaradas:}

\begin{itemize}
\tightlist
\item
  \texorpdfstring{\softcode{j jsonb:=CASE WHEN TG\_OP='DELETE' THEN to\_jsonb(OLD) ELSE to\_jsonb(NEW) END}}{j jsonb:=CASE WHEN TG\_OP='DELETE' THEN to\_jsonb(OLD) ELSE to\_jsonb(NEW) END}.
\end{itemize}

Para entender sus efectos, consulta la explicación SQL del cuerpo
principal y las RN. El prefijo \texorpdfstring{\softcode{fn\_}}{fn\_}
agrupa auxiliares; una función SECURITY DEFINER requiere los grants
explícitos y el contexto validado. Una llamada no autoriza por sí sola
al usuario humano.

\section{\texorpdfstring{\texorpdfstring{\softtext{bovi.trg\_propietario}}{bovi.trg\_propietario}}{}}\label{bovi.trg_propietario}

\textbf{Fuente de la definición:}
\texorpdfstring{\softcode{db/migrations/V2\_\_comandos.sql:9}}{db/migrations/V2\_\_comandos.sql:9}.

\textbf{Argumentos:} ``.

\textbf{Retorno y propiedades:}
\texorpdfstring{\softcode{trigger LANGUAGE plpgsql SET search\_path = pg\_catalog}}{trigger LANGUAGE plpgsql SET search\_path = pg\_catalog}.

\textbf{Variables locales declaradas:}

\begin{itemize}
\tightlist
\item
  \texorpdfstring{\softcode{j jsonb:=to\_jsonb(NEW)}}{j jsonb:=to\_jsonb(NEW)}.
\item
  \texorpdfstring{\softcode{oldj jsonb:=to\_jsonb(OLD)}}{oldj jsonb:=to\_jsonb(OLD)}.
\item
  \texorpdfstring{\softcode{n text:=bovi.fn\_nodo()}}{n text:=bovi.fn\_nodo()}.
\end{itemize}

Para entender sus efectos, consulta la explicación SQL del cuerpo
principal y las RN. El prefijo \texorpdfstring{\softcode{fn\_}}{fn\_}
agrupa auxiliares; una función SECURITY DEFINER requiere los grants
explícitos y el contexto validado. Una llamada no autoriza por sí sola
al usuario humano.

\section{\texorpdfstring{\texorpdfstring{\softtext{bovi.trg\_snapshot\_tarifa}}{bovi.trg\_snapshot\_tarifa}}{}}\label{bovi.trg_snapshot_tarifa}

\textbf{Fuente de la definición:}
\texorpdfstring{\softcode{db/migrations/V4\_\_capturas\_proveedor\_dto.sql:11}}{db/migrations/V4\_\_capturas\_proveedor\_dto.sql:11}.

\textbf{Argumentos:} ``.

\textbf{Retorno y propiedades:}
\texorpdfstring{\softcode{trigger LANGUAGE plpgsql SET search\_path=pg\_catalog}}{trigger LANGUAGE plpgsql SET search\_path=pg\_catalog}.

Para entender sus efectos, consulta la explicación SQL del cuerpo
principal y las RN. El prefijo \texorpdfstring{\softcode{fn\_}}{fn\_}
agrupa auxiliares; una función SECURITY DEFINER requiere los grants
explícitos y el contexto validado. Una llamada no autoriza por sí sola
al usuario humano.

\chapter{Referencia completa del contrato
HTTP}\label{referencia-completa-del-contrato-http}

Esta referencia se extrae de
\texorpdfstring{\softtext{OpenAPI0.2.0}}{OpenAPI0.2.0} en la revisión92
d3 d41. Describe el acuerdo versionado; sus respuestas declaradas no
equivalen a una ejecución nueva de cada
\texorpdfstring{\softtext{combinación.}}{combinación.} Las clases y
casos de prueba explicados en el cuerpo identifican la implementación y
las evidencias reales. Los nombres de parámetros y esquemas conservan su
ortografía para buscar directamente en el código.

\section{\texorpdfstring{POST
\texorpdfstring{\softtext{/sesion/login}}{/sesion/login}}{POST }}\label{post-sesionlogin}

Inicia sesión (cookie HttpOnly)

\textbf{Identificador:} \texorpdfstring{\softcode{login}}{login}.

{\def\LTcaptype{none} % do not increment counter
\begin{longtable}[]{@{}lll@{}}
\toprule\noalign{}
Parámetro y ubicación & Tipo & Requerido y propósito \\
\midrule\noalign{}
\endhead
\bottomrule\noalign{}
\endlastfoot
\texorpdfstring{\softcode{X-XSRF-TOKEN}}{X-XSRF-TOKEN} (header) & string
& Sí. \\
\end{longtable}
}

\textbf{Cuerpo
\texorpdfstring{\softtext{application/json:}}{application/json:}}
object. \texorpdfstring{\softtext{Obligatorio.}}{Obligatorio.}

{\def\LTcaptype{none} % do not increment counter
\begin{longtable}[]{@{}
  >{\raggedright\arraybackslash}p{(\linewidth - 4\tabcolsep) * \real{0.3333}}
  >{\raggedright\arraybackslash}p{(\linewidth - 4\tabcolsep) * \real{0.3333}}
  >{\raggedright\arraybackslash}p{(\linewidth - 4\tabcolsep) * \real{0.3333}}@{}}
\toprule\noalign{}
\begin{minipage}[b]{\linewidth}\raggedright
HTTP
\end{minipage} & \begin{minipage}[b]{\linewidth}\raggedright
Significado declarado
\end{minipage} & \begin{minipage}[b]{\linewidth}\raggedright
Cuerpo y media type
\end{minipage} \\
\midrule\noalign{}
\endhead
\bottomrule\noalign{}
\endlastfoot
204 & Sesión creada (Set-Cookie) & Sin cuerpo descrito \\
401 & Sin sesión &
\texorpdfstring{\softtext{application/problem+json:}}{application/problem+json:}
Error \\
503 & Dependencia no disponible; una respuesta perdida puede requerir
comprobar resultado con la misma clave. &
\texorpdfstring{\softtext{application/problem+json:}}{application/problem+json:}
Error \\
422 & Viola una regla de negocio &
\texorpdfstring{\softtext{application/problem+json:}}{application/problem+json:}
Error \\
\end{longtable}
}

\section{\texorpdfstring{POST
\texorpdfstring{\softtext{/sesion/logout}}{/sesion/logout}}{POST }}\label{post-sesionlogout}

Logout

\textbf{Identificador:} \texorpdfstring{\softcode{logout}}{logout}.

{\def\LTcaptype{none} % do not increment counter
\begin{longtable}[]{@{}lll@{}}
\toprule\noalign{}
Parámetro y ubicación & Tipo & Requerido y propósito \\
\midrule\noalign{}
\endhead
\bottomrule\noalign{}
\endlastfoot
\texorpdfstring{\softcode{X-XSRF-TOKEN}}{X-XSRF-TOKEN} (header) & string
& Sí. \\
\end{longtable}
}

{\def\LTcaptype{none} % do not increment counter
\begin{longtable}[]{@{}
  >{\raggedright\arraybackslash}p{(\linewidth - 4\tabcolsep) * \real{0.3333}}
  >{\raggedright\arraybackslash}p{(\linewidth - 4\tabcolsep) * \real{0.3333}}
  >{\raggedright\arraybackslash}p{(\linewidth - 4\tabcolsep) * \real{0.3333}}@{}}
\toprule\noalign{}
\begin{minipage}[b]{\linewidth}\raggedright
HTTP
\end{minipage} & \begin{minipage}[b]{\linewidth}\raggedright
Significado declarado
\end{minipage} & \begin{minipage}[b]{\linewidth}\raggedright
Cuerpo y media type
\end{minipage} \\
\midrule\noalign{}
\endhead
\bottomrule\noalign{}
\endlastfoot
204 & Sesión cerrada & Sin cuerpo descrito \\
401 & Sin sesión &
\texorpdfstring{\softtext{application/problem+json:}}{application/problem+json:}
Error \\
403 & Sin permiso para la sede o la acción &
\texorpdfstring{\softtext{application/problem+json:}}{application/problem+json:}
Error \\
503 & Dependencia no disponible; una respuesta perdida puede requerir
comprobar resultado con la misma clave. &
\texorpdfstring{\softtext{application/problem+json:}}{application/problem+json:}
Error \\
422 & Viola una regla de negocio &
\texorpdfstring{\softtext{application/problem+json:}}{application/problem+json:}
Error \\
\end{longtable}
}

\section{\texorpdfstring{GET
\texorpdfstring{\softtext{/sesion/actual}}{/sesion/actual}}{GET }}\label{get-sesionactual}

Usuario actual, sedes y capacidades autorizadas

\textbf{Identificador:}
\texorpdfstring{\softcode{sesionActual}}{sesionActual}.

{\def\LTcaptype{none} % do not increment counter
\begin{longtable}[]{@{}
  >{\raggedright\arraybackslash}p{(\linewidth - 4\tabcolsep) * \real{0.3333}}
  >{\raggedright\arraybackslash}p{(\linewidth - 4\tabcolsep) * \real{0.3333}}
  >{\raggedright\arraybackslash}p{(\linewidth - 4\tabcolsep) * \real{0.3333}}@{}}
\toprule\noalign{}
\begin{minipage}[b]{\linewidth}\raggedright
HTTP
\end{minipage} & \begin{minipage}[b]{\linewidth}\raggedright
Significado declarado
\end{minipage} & \begin{minipage}[b]{\linewidth}\raggedright
Cuerpo y media type
\end{minipage} \\
\midrule\noalign{}
\endhead
\bottomrule\noalign{}
\endlastfoot
200 & OK &
\texorpdfstring{\softtext{application/json:}}{application/json:}
\texorpdfstring{\softtext{SesionActual}}{SesionActual} \\
401 & Sin sesión &
\texorpdfstring{\softtext{application/problem+json:}}{application/problem+json:}
Error \\
403 & Sin permiso para la sede o la acción &
\texorpdfstring{\softtext{application/problem+json:}}{application/problem+json:}
Error \\
503 & Dependencia no disponible; una respuesta perdida puede requerir
comprobar resultado con la misma clave. &
\texorpdfstring{\softtext{application/problem+json:}}{application/problem+json:}
Error \\
\end{longtable}
}

\section{\texorpdfstring{GET
\texorpdfstring{\softtext{/sedes/\{sedeId\}/contrapartes}}{/sedes/\{sedeId\}/contrapartes}}{GET }}\label{get-sedessedeidcontrapartes}

Listar contrapartes

\textbf{Identificador:}
\texorpdfstring{\softcode{listarContrapartes}}{listarContrapartes}.

{\def\LTcaptype{none} % do not increment counter
\begin{longtable}[]{@{}lll@{}}
\toprule\noalign{}
Parámetro y ubicación & Tipo & Requerido y propósito \\
\midrule\noalign{}
\endhead
\bottomrule\noalign{}
\endlastfoot
\texorpdfstring{\softcode{sedeId}}{sedeId} (path) & string (uuid) &
Sí. \\
\texorpdfstring{\softcode{pagina}}{pagina} (query) & integer & No. \\
\texorpdfstring{\softcode{tamano}}{tamano} (query) & integer & No. \\
\texorpdfstring{\softcode{q}}{q} (query) & string & No. \\
\end{longtable}
}

{\def\LTcaptype{none} % do not increment counter
\begin{longtable}[]{@{}
  >{\raggedright\arraybackslash}p{(\linewidth - 4\tabcolsep) * \real{0.3333}}
  >{\raggedright\arraybackslash}p{(\linewidth - 4\tabcolsep) * \real{0.3333}}
  >{\raggedright\arraybackslash}p{(\linewidth - 4\tabcolsep) * \real{0.3333}}@{}}
\toprule\noalign{}
\begin{minipage}[b]{\linewidth}\raggedright
HTTP
\end{minipage} & \begin{minipage}[b]{\linewidth}\raggedright
Significado declarado
\end{minipage} & \begin{minipage}[b]{\linewidth}\raggedright
Cuerpo y media type
\end{minipage} \\
\midrule\noalign{}
\endhead
\bottomrule\noalign{}
\endlastfoot
200 & OK &
\texorpdfstring{\softtext{application/json:}}{application/json:}
objeto \\
403 & Sin permiso para la sede o la acción &
\texorpdfstring{\softtext{application/problem+json:}}{application/problem+json:}
Error \\
401 & Sin sesión &
\texorpdfstring{\softtext{application/problem+json:}}{application/problem+json:}
Error \\
503 & Dependencia no disponible; una respuesta perdida puede requerir
comprobar resultado con la misma clave. &
\texorpdfstring{\softtext{application/problem+json:}}{application/problem+json:}
Error \\
400 & Forma o parámetro inválido &
\texorpdfstring{\softtext{application/problem+json:}}{application/problem+json:}
Error \\
\end{longtable}
}

\section{\texorpdfstring{POST
\texorpdfstring{\softtext{/sedes/\{sedeId\}/contrapartes}}{/sedes/\{sedeId\}/contrapartes}}{POST }}\label{post-sedessedeidcontrapartes}

Crear contraparte

\textbf{Identificador:}
\texorpdfstring{\softcode{crearContraparte}}{crearContraparte}.
Escritura maestra; se ejecuta en el nodo A.

{\def\LTcaptype{none} % do not increment counter
\begin{longtable}[]{@{}lll@{}}
\toprule\noalign{}
Parámetro y ubicación & Tipo & Requerido y propósito \\
\midrule\noalign{}
\endhead
\bottomrule\noalign{}
\endlastfoot
\texorpdfstring{\softcode{sedeId}}{sedeId} (path) & string (uuid) &
Sí. \\
\texorpdfstring{\softcode{Idempotency-Key}}{Idempotency-Key} (header) &
string & Sí. \\
\texorpdfstring{\softcode{X-XSRF-TOKEN}}{X-XSRF-TOKEN} (header) & string
& Sí. \\
\end{longtable}
}

\textbf{Cuerpo
\texorpdfstring{\softtext{application/json:}}{application/json:}}
\texorpdfstring{\softtext{ContraparteNueva.}}{ContraparteNueva.}
\texorpdfstring{\softtext{Obligatorio.}}{Obligatorio.}

{\def\LTcaptype{none} % do not increment counter
\begin{longtable}[]{@{}
  >{\raggedright\arraybackslash}p{(\linewidth - 4\tabcolsep) * \real{0.3333}}
  >{\raggedright\arraybackslash}p{(\linewidth - 4\tabcolsep) * \real{0.3333}}
  >{\raggedright\arraybackslash}p{(\linewidth - 4\tabcolsep) * \real{0.3333}}@{}}
\toprule\noalign{}
\begin{minipage}[b]{\linewidth}\raggedright
HTTP
\end{minipage} & \begin{minipage}[b]{\linewidth}\raggedright
Significado declarado
\end{minipage} & \begin{minipage}[b]{\linewidth}\raggedright
Cuerpo y media type
\end{minipage} \\
\midrule\noalign{}
\endhead
\bottomrule\noalign{}
\endlastfoot
201 & Creada &
\texorpdfstring{\softtext{application/json:}}{application/json:}
Contraparte \\
409 & Conflicto de concurrencia,
\texorpdfstring{\softtext{version\_lock}}{version\_lock} o idempotencia
(misma clave, payload distinto) &
\texorpdfstring{\softtext{application/problem+json:}}{application/problem+json:}
Error \\
422 & Viola una regla de negocio &
\texorpdfstring{\softtext{application/problem+json:}}{application/problem+json:}
Error \\
503 & Dependencia no disponible; una respuesta perdida puede requerir
comprobar resultado con la misma clave. &
\texorpdfstring{\softtext{application/problem+json:}}{application/problem+json:}
Error \\
401 & Sin sesión &
\texorpdfstring{\softtext{application/problem+json:}}{application/problem+json:}
Error \\
403 & Sin permiso para la sede o la acción &
\texorpdfstring{\softtext{application/problem+json:}}{application/problem+json:}
Error \\
400 & Forma o parámetro inválido &
\texorpdfstring{\softtext{application/problem+json:}}{application/problem+json:}
Error \\
\end{longtable}
}

\section{\texorpdfstring{POST
\texorpdfstring{\softtext{/sedes/\{sedeId\}/animales}}{/sedes/\{sedeId\}/animales}}{POST }}\label{post-sedessedeidanimales}

Registrar identidad animal

\textbf{Identificador:}
\texorpdfstring{\softcode{registrarIdentidadAnimal}}{registrarIdentidadAnimal}.
Ejecutado en el maestro A. Si la sede es B, la respuesta indica
\texorpdfstring{\softtext{PENDIENTE\_CONVERGENCIA}}{PENDIENTE\_CONVERGENCIA}
hasta que la identidad sea visible en B.

{\def\LTcaptype{none} % do not increment counter
\begin{longtable}[]{@{}lll@{}}
\toprule\noalign{}
Parámetro y ubicación & Tipo & Requerido y propósito \\
\midrule\noalign{}
\endhead
\bottomrule\noalign{}
\endlastfoot
\texorpdfstring{\softcode{sedeId}}{sedeId} (path) & string (uuid) &
Sí. \\
\texorpdfstring{\softcode{Idempotency-Key}}{Idempotency-Key} (header) &
string & Sí. \\
\texorpdfstring{\softcode{X-XSRF-TOKEN}}{X-XSRF-TOKEN} (header) & string
& Sí. \\
\end{longtable}
}

\textbf{Cuerpo
\texorpdfstring{\softtext{application/json:}}{application/json:}}
object. \texorpdfstring{\softtext{Obligatorio.}}{Obligatorio.}

{\def\LTcaptype{none} % do not increment counter
\begin{longtable}[]{@{}
  >{\raggedright\arraybackslash}p{(\linewidth - 4\tabcolsep) * \real{0.3333}}
  >{\raggedright\arraybackslash}p{(\linewidth - 4\tabcolsep) * \real{0.3333}}
  >{\raggedright\arraybackslash}p{(\linewidth - 4\tabcolsep) * \real{0.3333}}@{}}
\toprule\noalign{}
\begin{minipage}[b]{\linewidth}\raggedright
HTTP
\end{minipage} & \begin{minipage}[b]{\linewidth}\raggedright
Significado declarado
\end{minipage} & \begin{minipage}[b]{\linewidth}\raggedright
Cuerpo y media type
\end{minipage} \\
\midrule\noalign{}
\endhead
\bottomrule\noalign{}
\endlastfoot
201 & Registrado &
\texorpdfstring{\softtext{application/json:}}{application/json:}
Animal \\
409 & Conflicto de concurrencia,
\texorpdfstring{\softtext{version\_lock}}{version\_lock} o idempotencia
(misma clave, payload distinto) &
\texorpdfstring{\softtext{application/problem+json:}}{application/problem+json:}
Error \\
503 & Dependencia no disponible; una respuesta perdida puede requerir
comprobar resultado con la misma clave. &
\texorpdfstring{\softtext{application/problem+json:}}{application/problem+json:}
Error \\
401 & Sin sesión &
\texorpdfstring{\softtext{application/problem+json:}}{application/problem+json:}
Error \\
403 & Sin permiso para la sede o la acción &
\texorpdfstring{\softtext{application/problem+json:}}{application/problem+json:}
Error \\
422 & Viola una regla de negocio &
\texorpdfstring{\softtext{application/problem+json:}}{application/problem+json:}
Error \\
400 & Forma o parámetro inválido &
\texorpdfstring{\softtext{application/problem+json:}}{application/problem+json:}
Error \\
\end{longtable}
}

\section{\texorpdfstring{GET
\texorpdfstring{\softtext{/sedes/\{sedeId\}/animales}}{/sedes/\{sedeId\}/animales}}{GET }}\label{get-sedessedeidanimales}

\texorpdfstring{\softtext{listarAnimales}}{listarAnimales}

\textbf{Identificador:}
\texorpdfstring{\softcode{listarAnimales}}{listarAnimales}.

{\def\LTcaptype{none} % do not increment counter
\begin{longtable}[]{@{}
  >{\raggedright\arraybackslash}p{(\linewidth - 4\tabcolsep) * \real{0.3333}}
  >{\raggedright\arraybackslash}p{(\linewidth - 4\tabcolsep) * \real{0.3333}}
  >{\raggedright\arraybackslash}p{(\linewidth - 4\tabcolsep) * \real{0.3333}}@{}}
\toprule\noalign{}
\begin{minipage}[b]{\linewidth}\raggedright
Parámetro y ubicación
\end{minipage} & \begin{minipage}[b]{\linewidth}\raggedright
Tipo
\end{minipage} & \begin{minipage}[b]{\linewidth}\raggedright
Requerido y propósito
\end{minipage} \\
\midrule\noalign{}
\endhead
\bottomrule\noalign{}
\endlastfoot
\texorpdfstring{\softcode{sedeId}}{sedeId} (path) & string (uuid) &
Sí. \\
\texorpdfstring{\softcode{pagina}}{pagina} (query) & integer & No. \\
\texorpdfstring{\softcode{tamano}}{tamano} (query) & integer & No. \\
\texorpdfstring{\softcode{q}}{q} (query) & string &
\texorpdfstring{\softtext{No. Búsqueda}}{No. Búsqueda} literal acotada;
\% y \_ no actúan como comodines. \\
\end{longtable}
}

{\def\LTcaptype{none} % do not increment counter
\begin{longtable}[]{@{}
  >{\raggedright\arraybackslash}p{(\linewidth - 4\tabcolsep) * \real{0.3333}}
  >{\raggedright\arraybackslash}p{(\linewidth - 4\tabcolsep) * \real{0.3333}}
  >{\raggedright\arraybackslash}p{(\linewidth - 4\tabcolsep) * \real{0.3333}}@{}}
\toprule\noalign{}
\begin{minipage}[b]{\linewidth}\raggedright
HTTP
\end{minipage} & \begin{minipage}[b]{\linewidth}\raggedright
Significado declarado
\end{minipage} & \begin{minipage}[b]{\linewidth}\raggedright
Cuerpo y media type
\end{minipage} \\
\midrule\noalign{}
\endhead
\bottomrule\noalign{}
\endlastfoot
200 & OK &
\texorpdfstring{\softtext{application/json:}}{application/json:}
objeto \\
401 & Sin sesión &
\texorpdfstring{\softtext{application/problem+json:}}{application/problem+json:}
Error \\
403 & Sin permiso para la sede o la acción &
\texorpdfstring{\softtext{application/problem+json:}}{application/problem+json:}
Error \\
503 & Dependencia no disponible; una respuesta perdida puede requerir
comprobar resultado con la misma clave. &
\texorpdfstring{\softtext{application/problem+json:}}{application/problem+json:}
Error \\
400 & Forma o parámetro inválido &
\texorpdfstring{\softtext{application/problem+json:}}{application/problem+json:}
Error \\
\end{longtable}
}

\section{\texorpdfstring{GET
\texorpdfstring{\softtext{/sedes/\{sedeId\}/animales/\{animalId\}}}{/sedes/\{sedeId\}/animales/\{animalId\}}}{GET }}\label{get-sedessedeidanimalesanimalid}

Ficha animal

\textbf{Identificador:}
\texorpdfstring{\softcode{fichaAnimal}}{fichaAnimal}.

{\def\LTcaptype{none} % do not increment counter
\begin{longtable}[]{@{}lll@{}}
\toprule\noalign{}
Parámetro y ubicación & Tipo & Requerido y propósito \\
\midrule\noalign{}
\endhead
\bottomrule\noalign{}
\endlastfoot
\texorpdfstring{\softcode{sedeId}}{sedeId} (path) & string (uuid) &
Sí. \\
\texorpdfstring{\softcode{animalId}}{animalId} (path) & string (uuid) &
Sí. \\
\end{longtable}
}

{\def\LTcaptype{none} % do not increment counter
\begin{longtable}[]{@{}
  >{\raggedright\arraybackslash}p{(\linewidth - 4\tabcolsep) * \real{0.3333}}
  >{\raggedright\arraybackslash}p{(\linewidth - 4\tabcolsep) * \real{0.3333}}
  >{\raggedright\arraybackslash}p{(\linewidth - 4\tabcolsep) * \real{0.3333}}@{}}
\toprule\noalign{}
\begin{minipage}[b]{\linewidth}\raggedright
HTTP
\end{minipage} & \begin{minipage}[b]{\linewidth}\raggedright
Significado declarado
\end{minipage} & \begin{minipage}[b]{\linewidth}\raggedright
Cuerpo y media type
\end{minipage} \\
\midrule\noalign{}
\endhead
\bottomrule\noalign{}
\endlastfoot
200 & Ficha con ciclos históricos &
\texorpdfstring{\softtext{application/json:}}{application/json:}
\texorpdfstring{\softtext{FichaAnimal}}{FichaAnimal} \\
404 & No existe o no es accesible (política uniforme) &
\texorpdfstring{\softtext{application/problem+json:}}{application/problem+json:}
Error \\
401 & Sin sesión &
\texorpdfstring{\softtext{application/problem+json:}}{application/problem+json:}
Error \\
403 & Sin permiso para la sede o la acción &
\texorpdfstring{\softtext{application/problem+json:}}{application/problem+json:}
Error \\
503 & Dependencia no disponible; una respuesta perdida puede requerir
comprobar resultado con la misma clave. &
\texorpdfstring{\softtext{application/problem+json:}}{application/problem+json:}
Error \\
400 & Forma o parámetro inválido &
\texorpdfstring{\softtext{application/problem+json:}}{application/problem+json:}
Error \\
\end{longtable}
}

\section{\texorpdfstring{GET
\texorpdfstring{\softtext{/sedes/\{sedeId\}/operaciones}}{/sedes/\{sedeId\}/operaciones}}{GET }}\label{get-sedessedeidoperaciones}

Bandeja operaciones

\textbf{Identificador:}
\texorpdfstring{\softcode{bandejaOperaciones}}{bandejaOperaciones}.

{\def\LTcaptype{none} % do not increment counter
\begin{longtable}[]{@{}
  >{\raggedright\arraybackslash}p{(\linewidth - 4\tabcolsep) * \real{0.3333}}
  >{\raggedright\arraybackslash}p{(\linewidth - 4\tabcolsep) * \real{0.3333}}
  >{\raggedright\arraybackslash}p{(\linewidth - 4\tabcolsep) * \real{0.3333}}@{}}
\toprule\noalign{}
\begin{minipage}[b]{\linewidth}\raggedright
Parámetro y ubicación
\end{minipage} & \begin{minipage}[b]{\linewidth}\raggedright
Tipo
\end{minipage} & \begin{minipage}[b]{\linewidth}\raggedright
Requerido y propósito
\end{minipage} \\
\midrule\noalign{}
\endhead
\bottomrule\noalign{}
\endlastfoot
\texorpdfstring{\softcode{sedeId}}{sedeId} (path) & string (uuid) &
Sí. \\
\texorpdfstring{\softcode{pagina}}{pagina} (query) & integer & No. \\
\texorpdfstring{\softcode{tamano}}{tamano} (query) & integer & No. \\
\texorpdfstring{\softcode{estado}}{estado} (query) &
\texorpdfstring{\softtext{EstadoOperacion}}{EstadoOperacion} & No. \\
\texorpdfstring{\softcode{tipo}}{tipo} (query) & string & No. \\
\texorpdfstring{\softcode{q}}{q} (query) & string &
\texorpdfstring{\softtext{No. Búsqueda}}{No. Búsqueda} literal acotada;
\% y \_ no actúan como comodines. \\
\end{longtable}
}

{\def\LTcaptype{none} % do not increment counter
\begin{longtable}[]{@{}
  >{\raggedright\arraybackslash}p{(\linewidth - 4\tabcolsep) * \real{0.3333}}
  >{\raggedright\arraybackslash}p{(\linewidth - 4\tabcolsep) * \real{0.3333}}
  >{\raggedright\arraybackslash}p{(\linewidth - 4\tabcolsep) * \real{0.3333}}@{}}
\toprule\noalign{}
\begin{minipage}[b]{\linewidth}\raggedright
HTTP
\end{minipage} & \begin{minipage}[b]{\linewidth}\raggedright
Significado declarado
\end{minipage} & \begin{minipage}[b]{\linewidth}\raggedright
Cuerpo y media type
\end{minipage} \\
\midrule\noalign{}
\endhead
\bottomrule\noalign{}
\endlastfoot
200 & OK &
\texorpdfstring{\softtext{application/json:}}{application/json:}
objeto \\
401 & Sin sesión &
\texorpdfstring{\softtext{application/problem+json:}}{application/problem+json:}
Error \\
403 & Sin permiso para la sede o la acción &
\texorpdfstring{\softtext{application/problem+json:}}{application/problem+json:}
Error \\
503 & Dependencia no disponible; una respuesta perdida puede requerir
comprobar resultado con la misma clave. &
\texorpdfstring{\softtext{application/problem+json:}}{application/problem+json:}
Error \\
400 & Forma o parámetro inválido &
\texorpdfstring{\softtext{application/problem+json:}}{application/problem+json:}
Error \\
\end{longtable}
}

\section{\texorpdfstring{POST
\texorpdfstring{\softtext{/sedes/\{sedeId\}/operaciones}}{/sedes/\{sedeId\}/operaciones}}{POST }}\label{post-sedessedeidoperaciones}

Crear borrador

\textbf{Identificador:}
\texorpdfstring{\softcode{crearBorrador}}{crearBorrador}.

{\def\LTcaptype{none} % do not increment counter
\begin{longtable}[]{@{}lll@{}}
\toprule\noalign{}
Parámetro y ubicación & Tipo & Requerido y propósito \\
\midrule\noalign{}
\endhead
\bottomrule\noalign{}
\endlastfoot
\texorpdfstring{\softcode{sedeId}}{sedeId} (path) & string (uuid) &
Sí. \\
\texorpdfstring{\softcode{Idempotency-Key}}{Idempotency-Key} (header) &
string & Sí. \\
\texorpdfstring{\softcode{X-XSRF-TOKEN}}{X-XSRF-TOKEN} (header) & string
& Sí. \\
\end{longtable}
}

\textbf{Cuerpo
\texorpdfstring{\softtext{application/json:}}{application/json:}}
\texorpdfstring{\softtext{OperacionNueva.}}{OperacionNueva.}
\texorpdfstring{\softtext{Obligatorio.}}{Obligatorio.}

{\def\LTcaptype{none} % do not increment counter
\begin{longtable}[]{@{}
  >{\raggedright\arraybackslash}p{(\linewidth - 4\tabcolsep) * \real{0.3333}}
  >{\raggedright\arraybackslash}p{(\linewidth - 4\tabcolsep) * \real{0.3333}}
  >{\raggedright\arraybackslash}p{(\linewidth - 4\tabcolsep) * \real{0.3333}}@{}}
\toprule\noalign{}
\begin{minipage}[b]{\linewidth}\raggedright
HTTP
\end{minipage} & \begin{minipage}[b]{\linewidth}\raggedright
Significado declarado
\end{minipage} & \begin{minipage}[b]{\linewidth}\raggedright
Cuerpo y media type
\end{minipage} \\
\midrule\noalign{}
\endhead
\bottomrule\noalign{}
\endlastfoot
201 & Borrador creado &
\texorpdfstring{\softtext{application/json:}}{application/json:}
Operacion \\
401 & Sin sesión &
\texorpdfstring{\softtext{application/problem+json:}}{application/problem+json:}
Error \\
403 & Sin permiso para la sede o la acción &
\texorpdfstring{\softtext{application/problem+json:}}{application/problem+json:}
Error \\
503 & Dependencia no disponible; una respuesta perdida puede requerir
comprobar resultado con la misma clave. &
\texorpdfstring{\softtext{application/problem+json:}}{application/problem+json:}
Error \\
422 & Viola una regla de negocio &
\texorpdfstring{\softtext{application/problem+json:}}{application/problem+json:}
Error \\
400 & Forma o parámetro inválido &
\texorpdfstring{\softtext{application/problem+json:}}{application/problem+json:}
Error \\
\end{longtable}
}

\section{\texorpdfstring{GET
\texorpdfstring{\softtext{/sedes/\{sedeId\}/operaciones/\{operacionId\}}}{/sedes/\{sedeId\}/operaciones/\{operacionId\}}}{GET }}\label{get-sedessedeidoperacionesoperacionid}

Obtener operacion

\textbf{Identificador:}
\texorpdfstring{\softcode{obtenerOperacion}}{obtenerOperacion}.

{\def\LTcaptype{none} % do not increment counter
\begin{longtable}[]{@{}lll@{}}
\toprule\noalign{}
Parámetro y ubicación & Tipo & Requerido y propósito \\
\midrule\noalign{}
\endhead
\bottomrule\noalign{}
\endlastfoot
\texorpdfstring{\softcode{sedeId}}{sedeId} (path) & string (uuid) &
Sí. \\
\texorpdfstring{\softcode{operacionId}}{operacionId} (path) & string
(uuid) & Sí. \\
\end{longtable}
}

{\def\LTcaptype{none} % do not increment counter
\begin{longtable}[]{@{}
  >{\raggedright\arraybackslash}p{(\linewidth - 4\tabcolsep) * \real{0.3333}}
  >{\raggedright\arraybackslash}p{(\linewidth - 4\tabcolsep) * \real{0.3333}}
  >{\raggedright\arraybackslash}p{(\linewidth - 4\tabcolsep) * \real{0.3333}}@{}}
\toprule\noalign{}
\begin{minipage}[b]{\linewidth}\raggedright
HTTP
\end{minipage} & \begin{minipage}[b]{\linewidth}\raggedright
Significado declarado
\end{minipage} & \begin{minipage}[b]{\linewidth}\raggedright
Cuerpo y media type
\end{minipage} \\
\midrule\noalign{}
\endhead
\bottomrule\noalign{}
\endlastfoot
200 & OK &
\texorpdfstring{\softtext{application/json:}}{application/json:}
Operacion \\
404 & No existe o no es accesible (política uniforme) &
\texorpdfstring{\softtext{application/problem+json:}}{application/problem+json:}
Error \\
401 & Sin sesión &
\texorpdfstring{\softtext{application/problem+json:}}{application/problem+json:}
Error \\
403 & Sin permiso para la sede o la acción &
\texorpdfstring{\softtext{application/problem+json:}}{application/problem+json:}
Error \\
503 & Dependencia no disponible; una respuesta perdida puede requerir
comprobar resultado con la misma clave. &
\texorpdfstring{\softtext{application/problem+json:}}{application/problem+json:}
Error \\
400 & Forma o parámetro inválido &
\texorpdfstring{\softtext{application/problem+json:}}{application/problem+json:}
Error \\
\end{longtable}
}

\section{\texorpdfstring{PATCH
\texorpdfstring{\softtext{/sedes/\{sedeId\}/operaciones/\{operacionId\}}}{/sedes/\{sedeId\}/operaciones/\{operacionId\}}}{PATCH }}\label{patch-sedessedeidoperacionesoperacionid}

Editar borrador

\textbf{Identificador:}
\texorpdfstring{\softcode{editarBorrador}}{editarBorrador}. Edita lotes,
selecciones y condiciones de la versión en curso. Sólo
\texorpdfstring{\softtext{BORRADOR/ACORDADA.}}{BORRADOR/ACORDADA.}

{\def\LTcaptype{none} % do not increment counter
\begin{longtable}[]{@{}lll@{}}
\toprule\noalign{}
Parámetro y ubicación & Tipo & Requerido y propósito \\
\midrule\noalign{}
\endhead
\bottomrule\noalign{}
\endlastfoot
\texorpdfstring{\softcode{sedeId}}{sedeId} (path) & string (uuid) &
Sí. \\
\texorpdfstring{\softcode{operacionId}}{operacionId} (path) & string
(uuid) & Sí. \\
\texorpdfstring{\softcode{If-Match}}{If-Match} (header) & string & Sí.
\texorpdfstring{\softtext{version\_lock}}{version\_lock} recibido en
ETag \\
\texorpdfstring{\softcode{X-XSRF-TOKEN}}{X-XSRF-TOKEN} (header) & string
& Sí. \\
\texorpdfstring{\softcode{Idempotency-Key}}{Idempotency-Key} (header) &
string & Sí. \\
\end{longtable}
}

\textbf{Cuerpo
\texorpdfstring{\softtext{application/json:}}{application/json:}}
\texorpdfstring{\softtext{EdicionBorrador.}}{EdicionBorrador.}
\texorpdfstring{\softtext{Obligatorio.}}{Obligatorio.}

{\def\LTcaptype{none} % do not increment counter
\begin{longtable}[]{@{}
  >{\raggedright\arraybackslash}p{(\linewidth - 4\tabcolsep) * \real{0.3333}}
  >{\raggedright\arraybackslash}p{(\linewidth - 4\tabcolsep) * \real{0.3333}}
  >{\raggedright\arraybackslash}p{(\linewidth - 4\tabcolsep) * \real{0.3333}}@{}}
\toprule\noalign{}
\begin{minipage}[b]{\linewidth}\raggedright
HTTP
\end{minipage} & \begin{minipage}[b]{\linewidth}\raggedright
Significado declarado
\end{minipage} & \begin{minipage}[b]{\linewidth}\raggedright
Cuerpo y media type
\end{minipage} \\
\midrule\noalign{}
\endhead
\bottomrule\noalign{}
\endlastfoot
200 & Actualizado &
\texorpdfstring{\softtext{application/json:}}{application/json:}
Operacion \\
409 & Conflicto de concurrencia,
\texorpdfstring{\softtext{version\_lock}}{version\_lock} o idempotencia
(misma clave, payload distinto) &
\texorpdfstring{\softtext{application/problem+json:}}{application/problem+json:}
Error \\
422 & Viola una regla de negocio &
\texorpdfstring{\softtext{application/problem+json:}}{application/problem+json:}
Error \\
412 & If-Match no corresponde a la versión revisada &
\texorpdfstring{\softtext{application/problem+json:}}{application/problem+json:}
Error \\
428 & Falta If-Match obligatorio &
\texorpdfstring{\softtext{application/problem+json:}}{application/problem+json:}
Error \\
401 & Sin sesión &
\texorpdfstring{\softtext{application/problem+json:}}{application/problem+json:}
Error \\
403 & Sin permiso para la sede o la acción &
\texorpdfstring{\softtext{application/problem+json:}}{application/problem+json:}
Error \\
503 & Dependencia no disponible; una respuesta perdida puede requerir
comprobar resultado con la misma clave. &
\texorpdfstring{\softtext{application/problem+json:}}{application/problem+json:}
Error \\
400 & Forma o parámetro inválido &
\texorpdfstring{\softtext{application/problem+json:}}{application/problem+json:}
Error \\
\end{longtable}
}

\section{\texorpdfstring{POST
\texorpdfstring{\softtext{/sedes/\{sedeId\}/operaciones/\{operacionId\}/confirmar}}{/sedes/\{sedeId\}/operaciones/\{operacionId\}/confirmar}}{POST }}\label{post-sedessedeidoperacionesoperacionidconfirmar}

Confirmar operacion

\textbf{Identificador:}
\texorpdfstring{\softcode{confirmarOperacion}}{confirmarOperacion}.
\texorpdfstring{\softtext{ConfirmarCompra}}{ConfirmarCompra} o
\texorpdfstring{\softtext{ConfirmarVenta}}{ConfirmarVenta} según el tipo
(RN-021, RN-022). Atómico.

{\def\LTcaptype{none} % do not increment counter
\begin{longtable}[]{@{}lll@{}}
\toprule\noalign{}
Parámetro y ubicación & Tipo & Requerido y propósito \\
\midrule\noalign{}
\endhead
\bottomrule\noalign{}
\endlastfoot
\texorpdfstring{\softcode{sedeId}}{sedeId} (path) & string (uuid) &
Sí. \\
\texorpdfstring{\softcode{operacionId}}{operacionId} (path) & string
(uuid) & Sí. \\
\texorpdfstring{\softcode{Idempotency-Key}}{Idempotency-Key} (header) &
string & Sí. \\
\texorpdfstring{\softcode{If-Match}}{If-Match} (header) & string & Sí.
\texorpdfstring{\softtext{version\_lock}}{version\_lock} recibido en
ETag \\
\texorpdfstring{\softcode{X-XSRF-TOKEN}}{X-XSRF-TOKEN} (header) & string
& Sí. \\
\end{longtable}
}

\textbf{Cuerpo
\texorpdfstring{\softtext{application/json:}}{application/json:}}
\texorpdfstring{\softtext{Reconocimiento.}}{Reconocimiento.}
\texorpdfstring{\softtext{Obligatorio.}}{Obligatorio.}

{\def\LTcaptype{none} % do not increment counter
\begin{longtable}[]{@{}
  >{\raggedright\arraybackslash}p{(\linewidth - 4\tabcolsep) * \real{0.3333}}
  >{\raggedright\arraybackslash}p{(\linewidth - 4\tabcolsep) * \real{0.3333}}
  >{\raggedright\arraybackslash}p{(\linewidth - 4\tabcolsep) * \real{0.3333}}@{}}
\toprule\noalign{}
\begin{minipage}[b]{\linewidth}\raggedright
HTTP
\end{minipage} & \begin{minipage}[b]{\linewidth}\raggedright
Significado declarado
\end{minipage} & \begin{minipage}[b]{\linewidth}\raggedright
Cuerpo y media type
\end{minipage} \\
\midrule\noalign{}
\endhead
\bottomrule\noalign{}
\endlastfoot
200 & Confirmada (o resultado previo si es reintento con la misma clave
y payload) &
\texorpdfstring{\softtext{application/json:}}{application/json:}
Liquidacion \\
409 & Conflicto de concurrencia,
\texorpdfstring{\softtext{version\_lock}}{version\_lock} o idempotencia
(misma clave, payload distinto) &
\texorpdfstring{\softtext{application/problem+json:}}{application/problem+json:}
Error \\
422 & Viola una regla de negocio &
\texorpdfstring{\softtext{application/problem+json:}}{application/problem+json:}
Error \\
503 & Dependencia no disponible; una respuesta perdida puede requerir
comprobar resultado con la misma clave. &
\texorpdfstring{\softtext{application/problem+json:}}{application/problem+json:}
Error \\
412 & If-Match no corresponde a la versión revisada &
\texorpdfstring{\softtext{application/problem+json:}}{application/problem+json:}
Error \\
428 & Falta If-Match obligatorio &
\texorpdfstring{\softtext{application/problem+json:}}{application/problem+json:}
Error \\
401 & Sin sesión &
\texorpdfstring{\softtext{application/problem+json:}}{application/problem+json:}
Error \\
403 & Sin permiso para la sede o la acción &
\texorpdfstring{\softtext{application/problem+json:}}{application/problem+json:}
Error \\
400 & Forma o parámetro inválido &
\texorpdfstring{\softtext{application/problem+json:}}{application/problem+json:}
Error \\
\end{longtable}
}

\section{\texorpdfstring{POST
\texorpdfstring{\softtext{/sedes/\{sedeId\}/operaciones/\{operacionId\}/pesajes}}{/sedes/\{sedeId\}/operaciones/\{operacionId\}/pesajes}}{POST }}\label{post-sedessedeidoperacionesoperacionidpesajes}

Proponer pesaje

\textbf{Identificador:}
\texorpdfstring{\softcode{proponerPesaje}}{proponerPesaje}.

{\def\LTcaptype{none} % do not increment counter
\begin{longtable}[]{@{}lll@{}}
\toprule\noalign{}
Parámetro y ubicación & Tipo & Requerido y propósito \\
\midrule\noalign{}
\endhead
\bottomrule\noalign{}
\endlastfoot
\texorpdfstring{\softcode{sedeId}}{sedeId} (path) & string (uuid) &
Sí. \\
\texorpdfstring{\softcode{operacionId}}{operacionId} (path) & string
(uuid) & Sí. \\
\texorpdfstring{\softcode{Idempotency-Key}}{Idempotency-Key} (header) &
string & Sí. \\
\texorpdfstring{\softcode{X-XSRF-TOKEN}}{X-XSRF-TOKEN} (header) & string
& Sí. \\
\end{longtable}
}

\textbf{Cuerpo
\texorpdfstring{\softtext{application/json:}}{application/json:}}
\texorpdfstring{\softtext{PesajeNuevo.}}{PesajeNuevo.}
\texorpdfstring{\softtext{Obligatorio.}}{Obligatorio.}

{\def\LTcaptype{none} % do not increment counter
\begin{longtable}[]{@{}
  >{\raggedright\arraybackslash}p{(\linewidth - 4\tabcolsep) * \real{0.3333}}
  >{\raggedright\arraybackslash}p{(\linewidth - 4\tabcolsep) * \real{0.3333}}
  >{\raggedright\arraybackslash}p{(\linewidth - 4\tabcolsep) * \real{0.3333}}@{}}
\toprule\noalign{}
\begin{minipage}[b]{\linewidth}\raggedright
HTTP
\end{minipage} & \begin{minipage}[b]{\linewidth}\raggedright
Significado declarado
\end{minipage} & \begin{minipage}[b]{\linewidth}\raggedright
Cuerpo y media type
\end{minipage} \\
\midrule\noalign{}
\endhead
\bottomrule\noalign{}
\endlastfoot
201 & Propuesto &
\texorpdfstring{\softtext{application/json:}}{application/json:}
Pesaje \\
422 & Viola una regla de negocio &
\texorpdfstring{\softtext{application/problem+json:}}{application/problem+json:}
Error \\
401 & Sin sesión &
\texorpdfstring{\softtext{application/problem+json:}}{application/problem+json:}
Error \\
403 & Sin permiso para la sede o la acción &
\texorpdfstring{\softtext{application/problem+json:}}{application/problem+json:}
Error \\
503 & Dependencia no disponible; una respuesta perdida puede requerir
comprobar resultado con la misma clave. &
\texorpdfstring{\softtext{application/problem+json:}}{application/problem+json:}
Error \\
400 & Forma o parámetro inválido &
\texorpdfstring{\softtext{application/problem+json:}}{application/problem+json:}
Error \\
\end{longtable}
}

\section{\texorpdfstring{POST
\texorpdfstring{\softtext{/sedes/\{sedeId\}/pesajes/\{pesajeId\}/aprobar}}{/sedes/\{sedeId\}/pesajes/\{pesajeId\}/aprobar}}{POST }}\label{post-sedessedeidpesajespesajeidaprobar}

Aprobar pesaje

\textbf{Identificador:}
\texorpdfstring{\softcode{aprobarPesaje}}{aprobarPesaje}.

{\def\LTcaptype{none} % do not increment counter
\begin{longtable}[]{@{}lll@{}}
\toprule\noalign{}
Parámetro y ubicación & Tipo & Requerido y propósito \\
\midrule\noalign{}
\endhead
\bottomrule\noalign{}
\endlastfoot
\texorpdfstring{\softcode{sedeId}}{sedeId} (path) & string (uuid) &
Sí. \\
\texorpdfstring{\softcode{pesajeId}}{pesajeId} (path) & string (uuid) &
Sí. \\
\texorpdfstring{\softcode{Idempotency-Key}}{Idempotency-Key} (header) &
string & Sí. \\
\texorpdfstring{\softcode{X-XSRF-TOKEN}}{X-XSRF-TOKEN} (header) & string
& Sí. \\
\end{longtable}
}

{\def\LTcaptype{none} % do not increment counter
\begin{longtable}[]{@{}
  >{\raggedright\arraybackslash}p{(\linewidth - 4\tabcolsep) * \real{0.3333}}
  >{\raggedright\arraybackslash}p{(\linewidth - 4\tabcolsep) * \real{0.3333}}
  >{\raggedright\arraybackslash}p{(\linewidth - 4\tabcolsep) * \real{0.3333}}@{}}
\toprule\noalign{}
\begin{minipage}[b]{\linewidth}\raggedright
HTTP
\end{minipage} & \begin{minipage}[b]{\linewidth}\raggedright
Significado declarado
\end{minipage} & \begin{minipage}[b]{\linewidth}\raggedright
Cuerpo y media type
\end{minipage} \\
\midrule\noalign{}
\endhead
\bottomrule\noalign{}
\endlastfoot
200 & Aprobado y composición congelada &
\texorpdfstring{\softtext{application/json:}}{application/json:}
Pesaje \\
422 & Viola una regla de negocio &
\texorpdfstring{\softtext{application/problem+json:}}{application/problem+json:}
Error \\
401 & Sin sesión &
\texorpdfstring{\softtext{application/problem+json:}}{application/problem+json:}
Error \\
403 & Sin permiso para la sede o la acción &
\texorpdfstring{\softtext{application/problem+json:}}{application/problem+json:}
Error \\
503 & Dependencia no disponible; una respuesta perdida puede requerir
comprobar resultado con la misma clave. &
\texorpdfstring{\softtext{application/problem+json:}}{application/problem+json:}
Error \\
400 & Forma o parámetro inválido &
\texorpdfstring{\softtext{application/problem+json:}}{application/problem+json:}
Error \\
\end{longtable}
}

\section{\texorpdfstring{POST
\texorpdfstring{\softtext{/sedes/\{sedeId\}/pesajes/\{pesajeId\}/correcciones}}{/sedes/\{sedeId\}/pesajes/\{pesajeId\}/correcciones}}{POST }}\label{post-sedessedeidpesajespesajeidcorrecciones}

Corregir pesaje

\textbf{Identificador:}
\texorpdfstring{\softcode{corregirPesaje}}{corregirPesaje}. Crea un
nuevo pesaje que corrige al original (no lo edita).

{\def\LTcaptype{none} % do not increment counter
\begin{longtable}[]{@{}lll@{}}
\toprule\noalign{}
Parámetro y ubicación & Tipo & Requerido y propósito \\
\midrule\noalign{}
\endhead
\bottomrule\noalign{}
\endlastfoot
\texorpdfstring{\softcode{sedeId}}{sedeId} (path) & string (uuid) &
Sí. \\
\texorpdfstring{\softcode{pesajeId}}{pesajeId} (path) & string (uuid) &
Sí. \\
\texorpdfstring{\softcode{Idempotency-Key}}{Idempotency-Key} (header) &
string & Sí. \\
\texorpdfstring{\softcode{X-XSRF-TOKEN}}{X-XSRF-TOKEN} (header) & string
& Sí. \\
\end{longtable}
}

\textbf{Cuerpo
\texorpdfstring{\softtext{application/json:}}{application/json:}}
\texorpdfstring{\softtext{PesajeCorreccion.}}{PesajeCorreccion.}
\texorpdfstring{\softtext{Obligatorio.}}{Obligatorio.}

{\def\LTcaptype{none} % do not increment counter
\begin{longtable}[]{@{}
  >{\raggedright\arraybackslash}p{(\linewidth - 4\tabcolsep) * \real{0.3333}}
  >{\raggedright\arraybackslash}p{(\linewidth - 4\tabcolsep) * \real{0.3333}}
  >{\raggedright\arraybackslash}p{(\linewidth - 4\tabcolsep) * \real{0.3333}}@{}}
\toprule\noalign{}
\begin{minipage}[b]{\linewidth}\raggedright
HTTP
\end{minipage} & \begin{minipage}[b]{\linewidth}\raggedright
Significado declarado
\end{minipage} & \begin{minipage}[b]{\linewidth}\raggedright
Cuerpo y media type
\end{minipage} \\
\midrule\noalign{}
\endhead
\bottomrule\noalign{}
\endlastfoot
201 & Corrección registrada &
\texorpdfstring{\softtext{application/json:}}{application/json:}
Pesaje \\
401 & Sin sesión &
\texorpdfstring{\softtext{application/problem+json:}}{application/problem+json:}
Error \\
403 & Sin permiso para la sede o la acción &
\texorpdfstring{\softtext{application/problem+json:}}{application/problem+json:}
Error \\
503 & Dependencia no disponible; una respuesta perdida puede requerir
comprobar resultado con la misma clave. &
\texorpdfstring{\softtext{application/problem+json:}}{application/problem+json:}
Error \\
422 & Viola una regla de negocio &
\texorpdfstring{\softtext{application/problem+json:}}{application/problem+json:}
Error \\
400 & Forma o parámetro inválido &
\texorpdfstring{\softtext{application/problem+json:}}{application/problem+json:}
Error \\
\end{longtable}
}

\section{\texorpdfstring{POST
\texorpdfstring{\softtext{/sedes/\{sedeId\}/operaciones/\{operacionId\}/reservas}}{/sedes/\{sedeId\}/operaciones/\{operacionId\}/reservas}}{POST }}\label{post-sedessedeidoperacionesoperacionidreservas}

Reservar posiciones

\textbf{Identificador:}
\texorpdfstring{\softcode{reservarPosiciones}}{reservarPosiciones}.

{\def\LTcaptype{none} % do not increment counter
\begin{longtable}[]{@{}lll@{}}
\toprule\noalign{}
Parámetro y ubicación & Tipo & Requerido y propósito \\
\midrule\noalign{}
\endhead
\bottomrule\noalign{}
\endlastfoot
\texorpdfstring{\softcode{sedeId}}{sedeId} (path) & string (uuid) &
Sí. \\
\texorpdfstring{\softcode{operacionId}}{operacionId} (path) & string
(uuid) & Sí. \\
\texorpdfstring{\softcode{Idempotency-Key}}{Idempotency-Key} (header) &
string & Sí. \\
\texorpdfstring{\softcode{X-XSRF-TOKEN}}{X-XSRF-TOKEN} (header) & string
& Sí. \\
\end{longtable}
}

\textbf{Cuerpo
\texorpdfstring{\softtext{application/json:}}{application/json:}}
object. \texorpdfstring{\softtext{Obligatorio.}}{Obligatorio.}

{\def\LTcaptype{none} % do not increment counter
\begin{longtable}[]{@{}
  >{\raggedright\arraybackslash}p{(\linewidth - 4\tabcolsep) * \real{0.3333}}
  >{\raggedright\arraybackslash}p{(\linewidth - 4\tabcolsep) * \real{0.3333}}
  >{\raggedright\arraybackslash}p{(\linewidth - 4\tabcolsep) * \real{0.3333}}@{}}
\toprule\noalign{}
\begin{minipage}[b]{\linewidth}\raggedright
HTTP
\end{minipage} & \begin{minipage}[b]{\linewidth}\raggedright
Significado declarado
\end{minipage} & \begin{minipage}[b]{\linewidth}\raggedright
Cuerpo y media type
\end{minipage} \\
\midrule\noalign{}
\endhead
\bottomrule\noalign{}
\endlastfoot
201 & Reservas activas &
\texorpdfstring{\softtext{application/json:}}{application/json:}
colección de Reserva \\
409 & Posición no disponible (otra reserva activa o vendida) &
\texorpdfstring{\softtext{application/problem+json:}}{application/problem+json:}
Error \\
401 & Sin sesión &
\texorpdfstring{\softtext{application/problem+json:}}{application/problem+json:}
Error \\
403 & Sin permiso para la sede o la acción &
\texorpdfstring{\softtext{application/problem+json:}}{application/problem+json:}
Error \\
503 & Dependencia no disponible; una respuesta perdida puede requerir
comprobar resultado con la misma clave. &
\texorpdfstring{\softtext{application/problem+json:}}{application/problem+json:}
Error \\
422 & Viola una regla de negocio &
\texorpdfstring{\softtext{application/problem+json:}}{application/problem+json:}
Error \\
400 & Forma o parámetro inválido &
\texorpdfstring{\softtext{application/problem+json:}}{application/problem+json:}
Error \\
\end{longtable}
}

\section{\texorpdfstring{POST
\texorpdfstring{\softtext{/sedes/\{sedeId\}/reservas/\{reservaId\}/liberar}}{/sedes/\{sedeId\}/reservas/\{reservaId\}/liberar}}{POST }}\label{post-sedessedeidreservasreservaidliberar}

Liberar reserva

\textbf{Identificador:}
\texorpdfstring{\softcode{liberarReserva}}{liberarReserva}.

{\def\LTcaptype{none} % do not increment counter
\begin{longtable}[]{@{}lll@{}}
\toprule\noalign{}
Parámetro y ubicación & Tipo & Requerido y propósito \\
\midrule\noalign{}
\endhead
\bottomrule\noalign{}
\endlastfoot
\texorpdfstring{\softcode{sedeId}}{sedeId} (path) & string (uuid) &
Sí. \\
\texorpdfstring{\softcode{reservaId}}{reservaId} (path) & string (uuid)
& Sí. \\
\texorpdfstring{\softcode{Idempotency-Key}}{Idempotency-Key} (header) &
string & Sí. \\
\texorpdfstring{\softcode{X-XSRF-TOKEN}}{X-XSRF-TOKEN} (header) & string
& Sí. \\
\end{longtable}
}

{\def\LTcaptype{none} % do not increment counter
\begin{longtable}[]{@{}
  >{\raggedright\arraybackslash}p{(\linewidth - 4\tabcolsep) * \real{0.3333}}
  >{\raggedright\arraybackslash}p{(\linewidth - 4\tabcolsep) * \real{0.3333}}
  >{\raggedright\arraybackslash}p{(\linewidth - 4\tabcolsep) * \real{0.3333}}@{}}
\toprule\noalign{}
\begin{minipage}[b]{\linewidth}\raggedright
HTTP
\end{minipage} & \begin{minipage}[b]{\linewidth}\raggedright
Significado declarado
\end{minipage} & \begin{minipage}[b]{\linewidth}\raggedright
Cuerpo y media type
\end{minipage} \\
\midrule\noalign{}
\endhead
\bottomrule\noalign{}
\endlastfoot
200 & Liberada &
\texorpdfstring{\softtext{application/json:}}{application/json:}
Reserva \\
401 & Sin sesión &
\texorpdfstring{\softtext{application/problem+json:}}{application/problem+json:}
Error \\
403 & Sin permiso para la sede o la acción &
\texorpdfstring{\softtext{application/problem+json:}}{application/problem+json:}
Error \\
503 & Dependencia no disponible; una respuesta perdida puede requerir
comprobar resultado con la misma clave. &
\texorpdfstring{\softtext{application/problem+json:}}{application/problem+json:}
Error \\
422 & Viola una regla de negocio &
\texorpdfstring{\softtext{application/problem+json:}}{application/problem+json:}
Error \\
400 & Forma o parámetro inválido &
\texorpdfstring{\softtext{application/problem+json:}}{application/problem+json:}
Error \\
\end{longtable}
}

\section{\texorpdfstring{GET
\texorpdfstring{\softtext{/sedes/\{sedeId\}/inventario/posiciones}}{/sedes/\{sedeId\}/inventario/posiciones}}{GET }}\label{get-sedessedeidinventarioposiciones}

Listar posiciones

\textbf{Identificador:}
\texorpdfstring{\softcode{listarPosiciones}}{listarPosiciones}.

{\def\LTcaptype{none} % do not increment counter
\begin{longtable}[]{@{}
  >{\raggedright\arraybackslash}p{(\linewidth - 4\tabcolsep) * \real{0.3333}}
  >{\raggedright\arraybackslash}p{(\linewidth - 4\tabcolsep) * \real{0.3333}}
  >{\raggedright\arraybackslash}p{(\linewidth - 4\tabcolsep) * \real{0.3333}}@{}}
\toprule\noalign{}
\begin{minipage}[b]{\linewidth}\raggedright
Parámetro y ubicación
\end{minipage} & \begin{minipage}[b]{\linewidth}\raggedright
Tipo
\end{minipage} & \begin{minipage}[b]{\linewidth}\raggedright
Requerido y propósito
\end{minipage} \\
\midrule\noalign{}
\endhead
\bottomrule\noalign{}
\endlastfoot
\texorpdfstring{\softcode{sedeId}}{sedeId} (path) & string (uuid) &
Sí. \\
\texorpdfstring{\softcode{pagina}}{pagina} (query) & integer & No. \\
\texorpdfstring{\softcode{tamano}}{tamano} (query) & integer & No. \\
\texorpdfstring{\softcode{estado}}{estado} (query) &
\texorpdfstring{\softtext{EstadoPosicion}}{EstadoPosicion} & No. \\
\texorpdfstring{\softcode{disponibles\_para\_venta}}{disponibles\_para\_venta}
(query) & boolean & No.~true: sólo posiciones disponibles, sin reserva
activa vigente ni identidad bloqueada; false: no aplicar ese filtro. \\
\texorpdfstring{\softcode{q}}{q} (query) & string &
\texorpdfstring{\softtext{No. Filtro}}{No. Filtro} de búsqueda
acotado. \\
\end{longtable}
}

{\def\LTcaptype{none} % do not increment counter
\begin{longtable}[]{@{}
  >{\raggedright\arraybackslash}p{(\linewidth - 4\tabcolsep) * \real{0.3333}}
  >{\raggedright\arraybackslash}p{(\linewidth - 4\tabcolsep) * \real{0.3333}}
  >{\raggedright\arraybackslash}p{(\linewidth - 4\tabcolsep) * \real{0.3333}}@{}}
\toprule\noalign{}
\begin{minipage}[b]{\linewidth}\raggedright
HTTP
\end{minipage} & \begin{minipage}[b]{\linewidth}\raggedright
Significado declarado
\end{minipage} & \begin{minipage}[b]{\linewidth}\raggedright
Cuerpo y media type
\end{minipage} \\
\midrule\noalign{}
\endhead
\bottomrule\noalign{}
\endlastfoot
200 & OK &
\texorpdfstring{\softtext{application/json:}}{application/json:}
objeto \\
401 & Sin sesión &
\texorpdfstring{\softtext{application/problem+json:}}{application/problem+json:}
Error \\
403 & Sin permiso para la sede o la acción &
\texorpdfstring{\softtext{application/problem+json:}}{application/problem+json:}
Error \\
503 & Dependencia no disponible; una respuesta perdida puede requerir
comprobar resultado con la misma clave. &
\texorpdfstring{\softtext{application/problem+json:}}{application/problem+json:}
Error \\
400 & Forma o parámetro inválido &
\texorpdfstring{\softtext{application/problem+json:}}{application/problem+json:}
Error \\
\end{longtable}
}

\section{\texorpdfstring{POST
\texorpdfstring{\softtext{/sedes/\{sedeId\}/inventario/movimientos}}{/sedes/\{sedeId\}/inventario/movimientos}}{POST }}\label{post-sedessedeidinventariomovimientos}

Registrar movimiento

\textbf{Identificador:}
\texorpdfstring{\softcode{registrarMovimiento}}{registrarMovimiento}.
Recepción, entrega o baja. No genera venta, costo ni obligación
(RN-020).

{\def\LTcaptype{none} % do not increment counter
\begin{longtable}[]{@{}lll@{}}
\toprule\noalign{}
Parámetro y ubicación & Tipo & Requerido y propósito \\
\midrule\noalign{}
\endhead
\bottomrule\noalign{}
\endlastfoot
\texorpdfstring{\softcode{sedeId}}{sedeId} (path) & string (uuid) &
Sí. \\
\texorpdfstring{\softcode{Idempotency-Key}}{Idempotency-Key} (header) &
string & Sí. \\
\texorpdfstring{\softcode{X-XSRF-TOKEN}}{X-XSRF-TOKEN} (header) & string
& Sí. \\
\end{longtable}
}

\textbf{Cuerpo
\texorpdfstring{\softtext{application/json:}}{application/json:}}
object. \texorpdfstring{\softtext{Obligatorio.}}{Obligatorio.}

{\def\LTcaptype{none} % do not increment counter
\begin{longtable}[]{@{}
  >{\raggedright\arraybackslash}p{(\linewidth - 4\tabcolsep) * \real{0.3333}}
  >{\raggedright\arraybackslash}p{(\linewidth - 4\tabcolsep) * \real{0.3333}}
  >{\raggedright\arraybackslash}p{(\linewidth - 4\tabcolsep) * \real{0.3333}}@{}}
\toprule\noalign{}
\begin{minipage}[b]{\linewidth}\raggedright
HTTP
\end{minipage} & \begin{minipage}[b]{\linewidth}\raggedright
Significado declarado
\end{minipage} & \begin{minipage}[b]{\linewidth}\raggedright
Cuerpo y media type
\end{minipage} \\
\midrule\noalign{}
\endhead
\bottomrule\noalign{}
\endlastfoot
201 & Registrado & Sin cuerpo descrito \\
422 & Viola una regla de negocio &
\texorpdfstring{\softtext{application/problem+json:}}{application/problem+json:}
Error \\
401 & Sin sesión &
\texorpdfstring{\softtext{application/problem+json:}}{application/problem+json:}
Error \\
403 & Sin permiso para la sede o la acción &
\texorpdfstring{\softtext{application/problem+json:}}{application/problem+json:}
Error \\
503 & Dependencia no disponible; una respuesta perdida puede requerir
comprobar resultado con la misma clave. &
\texorpdfstring{\softtext{application/problem+json:}}{application/problem+json:}
Error \\
400 & Forma o parámetro inválido &
\texorpdfstring{\softtext{application/problem+json:}}{application/problem+json:}
Error \\
\end{longtable}
}

\section{\texorpdfstring{GET
\texorpdfstring{\softtext{/sedes/\{sedeId\}/liquidaciones/\{liquidacionId\}}}{/sedes/\{sedeId\}/liquidaciones/\{liquidacionId\}}}{GET }}\label{get-sedessedeidliquidacionesliquidacionid}

Detalle liquidacion

\textbf{Identificador:}
\texorpdfstring{\softcode{detalleLiquidacion}}{detalleLiquidacion}.

{\def\LTcaptype{none} % do not increment counter
\begin{longtable}[]{@{}lll@{}}
\toprule\noalign{}
Parámetro y ubicación & Tipo & Requerido y propósito \\
\midrule\noalign{}
\endhead
\bottomrule\noalign{}
\endlastfoot
\texorpdfstring{\softcode{sedeId}}{sedeId} (path) & string (uuid) &
Sí. \\
\texorpdfstring{\softcode{liquidacionId}}{liquidacionId} (path) & string
(uuid) & Sí. \\
\end{longtable}
}

{\def\LTcaptype{none} % do not increment counter
\begin{longtable}[]{@{}
  >{\raggedright\arraybackslash}p{(\linewidth - 4\tabcolsep) * \real{0.3333}}
  >{\raggedright\arraybackslash}p{(\linewidth - 4\tabcolsep) * \real{0.3333}}
  >{\raggedright\arraybackslash}p{(\linewidth - 4\tabcolsep) * \real{0.3333}}@{}}
\toprule\noalign{}
\begin{minipage}[b]{\linewidth}\raggedright
HTTP
\end{minipage} & \begin{minipage}[b]{\linewidth}\raggedright
Significado declarado
\end{minipage} & \begin{minipage}[b]{\linewidth}\raggedright
Cuerpo y media type
\end{minipage} \\
\midrule\noalign{}
\endhead
\bottomrule\noalign{}
\endlastfoot
200 & Snapshot, ajustes, aplicaciones y saldo &
\texorpdfstring{\softtext{application/json:}}{application/json:}
Liquidacion \\
404 & No existe o no es accesible (política uniforme) &
\texorpdfstring{\softtext{application/problem+json:}}{application/problem+json:}
Error \\
401 & Sin sesión &
\texorpdfstring{\softtext{application/problem+json:}}{application/problem+json:}
Error \\
403 & Sin permiso para la sede o la acción &
\texorpdfstring{\softtext{application/problem+json:}}{application/problem+json:}
Error \\
503 & Dependencia no disponible; una respuesta perdida puede requerir
comprobar resultado con la misma clave. &
\texorpdfstring{\softtext{application/problem+json:}}{application/problem+json:}
Error \\
400 & Forma o parámetro inválido &
\texorpdfstring{\softtext{application/problem+json:}}{application/problem+json:}
Error \\
\end{longtable}
}

\section{\texorpdfstring{POST
\texorpdfstring{\softtext{/sedes/\{sedeId\}/liquidaciones/\{liquidacionId\}/ajustes}}{/sedes/\{sedeId\}/liquidaciones/\{liquidacionId\}/ajustes}}{POST }}\label{post-sedessedeidliquidacionesliquidacionidajustes}

Ajustar liquidacion

\textbf{Identificador:}
\texorpdfstring{\softcode{ajustarLiquidacion}}{ajustarLiquidacion}.

{\def\LTcaptype{none} % do not increment counter
\begin{longtable}[]{@{}lll@{}}
\toprule\noalign{}
Parámetro y ubicación & Tipo & Requerido y propósito \\
\midrule\noalign{}
\endhead
\bottomrule\noalign{}
\endlastfoot
\texorpdfstring{\softcode{sedeId}}{sedeId} (path) & string (uuid) &
Sí. \\
\texorpdfstring{\softcode{liquidacionId}}{liquidacionId} (path) & string
(uuid) & Sí. \\
\texorpdfstring{\softcode{Idempotency-Key}}{Idempotency-Key} (header) &
string & Sí. \\
\texorpdfstring{\softcode{X-XSRF-TOKEN}}{X-XSRF-TOKEN} (header) & string
& Sí. \\
\end{longtable}
}

\textbf{Cuerpo
\texorpdfstring{\softtext{application/json:}}{application/json:}}
object. \texorpdfstring{\softtext{Obligatorio.}}{Obligatorio.}

{\def\LTcaptype{none} % do not increment counter
\begin{longtable}[]{@{}
  >{\raggedright\arraybackslash}p{(\linewidth - 4\tabcolsep) * \real{0.3333}}
  >{\raggedright\arraybackslash}p{(\linewidth - 4\tabcolsep) * \real{0.3333}}
  >{\raggedright\arraybackslash}p{(\linewidth - 4\tabcolsep) * \real{0.3333}}@{}}
\toprule\noalign{}
\begin{minipage}[b]{\linewidth}\raggedright
HTTP
\end{minipage} & \begin{minipage}[b]{\linewidth}\raggedright
Significado declarado
\end{minipage} & \begin{minipage}[b]{\linewidth}\raggedright
Cuerpo y media type
\end{minipage} \\
\midrule\noalign{}
\endhead
\bottomrule\noalign{}
\endlastfoot
201 & Ajuste registrado &
\texorpdfstring{\softtext{application/json:}}{application/json:}
Liquidacion \\
422 & Viola una regla de negocio &
\texorpdfstring{\softtext{application/problem+json:}}{application/problem+json:}
Error \\
401 & Sin sesión &
\texorpdfstring{\softtext{application/problem+json:}}{application/problem+json:}
Error \\
403 & Sin permiso para la sede o la acción &
\texorpdfstring{\softtext{application/problem+json:}}{application/problem+json:}
Error \\
503 & Dependencia no disponible; una respuesta perdida puede requerir
comprobar resultado con la misma clave. &
\texorpdfstring{\softtext{application/problem+json:}}{application/problem+json:}
Error \\
400 & Forma o parámetro inválido &
\texorpdfstring{\softtext{application/problem+json:}}{application/problem+json:}
Error \\
\end{longtable}
}

\section{\texorpdfstring{GET
\texorpdfstring{\softtext{/sedes/\{sedeId\}/liquidaciones/\{liquidacionId\}/comprobante}}{/sedes/\{sedeId\}/liquidaciones/\{liquidacionId\}/comprobante}}{GET }}\label{get-sedessedeidliquidacionesliquidacionidcomprobante}

Comprobante interno

\textbf{Identificador:}
\texorpdfstring{\softcode{comprobanteInterno}}{comprobanteInterno}. PDF
interno (no es CFDI). No genera ningún hecho
\texorpdfstring{\softtext{económico.}}{económico.}

{\def\LTcaptype{none} % do not increment counter
\begin{longtable}[]{@{}lll@{}}
\toprule\noalign{}
Parámetro y ubicación & Tipo & Requerido y propósito \\
\midrule\noalign{}
\endhead
\bottomrule\noalign{}
\endlastfoot
\texorpdfstring{\softcode{sedeId}}{sedeId} (path) & string (uuid) &
Sí. \\
\texorpdfstring{\softcode{liquidacionId}}{liquidacionId} (path) & string
(uuid) & Sí. \\
\end{longtable}
}

{\def\LTcaptype{none} % do not increment counter
\begin{longtable}[]{@{}
  >{\raggedright\arraybackslash}p{(\linewidth - 4\tabcolsep) * \real{0.3333}}
  >{\raggedright\arraybackslash}p{(\linewidth - 4\tabcolsep) * \real{0.3333}}
  >{\raggedright\arraybackslash}p{(\linewidth - 4\tabcolsep) * \real{0.3333}}@{}}
\toprule\noalign{}
\begin{minipage}[b]{\linewidth}\raggedright
HTTP
\end{minipage} & \begin{minipage}[b]{\linewidth}\raggedright
Significado declarado
\end{minipage} & \begin{minipage}[b]{\linewidth}\raggedright
Cuerpo y media type
\end{minipage} \\
\midrule\noalign{}
\endhead
\bottomrule\noalign{}
\endlastfoot
200 & PDF &
\texorpdfstring{\softtext{application/pdf:}}{application/pdf:} string
(binary) \\
401 & Sin sesión &
\texorpdfstring{\softtext{application/problem+json:}}{application/problem+json:}
Error \\
403 & Sin permiso para la sede o la acción &
\texorpdfstring{\softtext{application/problem+json:}}{application/problem+json:}
Error \\
503 & Dependencia no disponible; una respuesta perdida puede requerir
comprobar resultado con la misma clave. &
\texorpdfstring{\softtext{application/problem+json:}}{application/problem+json:}
Error \\
400 & Forma o parámetro inválido &
\texorpdfstring{\softtext{application/problem+json:}}{application/problem+json:}
Error \\
\end{longtable}
}

\section{\texorpdfstring{POST
\texorpdfstring{\softtext{/sedes/\{sedeId\}/gastos}}{/sedes/\{sedeId\}/gastos}}{POST }}\label{post-sedessedeidgastos}

Registrar gasto

\textbf{Identificador:}
\texorpdfstring{\softcode{registrarGasto}}{registrarGasto}.

{\def\LTcaptype{none} % do not increment counter
\begin{longtable}[]{@{}lll@{}}
\toprule\noalign{}
Parámetro y ubicación & Tipo & Requerido y propósito \\
\midrule\noalign{}
\endhead
\bottomrule\noalign{}
\endlastfoot
\texorpdfstring{\softcode{sedeId}}{sedeId} (path) & string (uuid) &
Sí. \\
\texorpdfstring{\softcode{Idempotency-Key}}{Idempotency-Key} (header) &
string & Sí. \\
\texorpdfstring{\softcode{X-XSRF-TOKEN}}{X-XSRF-TOKEN} (header) & string
& Sí. \\
\end{longtable}
}

\textbf{Cuerpo
\texorpdfstring{\softtext{application/json:}}{application/json:}}
\texorpdfstring{\softtext{GastoNuevo.}}{GastoNuevo.}
\texorpdfstring{\softtext{Obligatorio.}}{Obligatorio.}

{\def\LTcaptype{none} % do not increment counter
\begin{longtable}[]{@{}
  >{\raggedright\arraybackslash}p{(\linewidth - 4\tabcolsep) * \real{0.3333}}
  >{\raggedright\arraybackslash}p{(\linewidth - 4\tabcolsep) * \real{0.3333}}
  >{\raggedright\arraybackslash}p{(\linewidth - 4\tabcolsep) * \real{0.3333}}@{}}
\toprule\noalign{}
\begin{minipage}[b]{\linewidth}\raggedright
HTTP
\end{minipage} & \begin{minipage}[b]{\linewidth}\raggedright
Significado declarado
\end{minipage} & \begin{minipage}[b]{\linewidth}\raggedright
Cuerpo y media type
\end{minipage} \\
\midrule\noalign{}
\endhead
\bottomrule\noalign{}
\endlastfoot
201 & Gasto registrado &
\texorpdfstring{\softtext{application/json:}}{application/json:}
Gasto \\
401 & Sin sesión &
\texorpdfstring{\softtext{application/problem+json:}}{application/problem+json:}
Error \\
403 & Sin permiso para la sede o la acción &
\texorpdfstring{\softtext{application/problem+json:}}{application/problem+json:}
Error \\
503 & Dependencia no disponible; una respuesta perdida puede requerir
comprobar resultado con la misma clave. &
\texorpdfstring{\softtext{application/problem+json:}}{application/problem+json:}
Error \\
422 & Viola una regla de negocio &
\texorpdfstring{\softtext{application/problem+json:}}{application/problem+json:}
Error \\
400 & Forma o parámetro inválido &
\texorpdfstring{\softtext{application/problem+json:}}{application/problem+json:}
Error \\
\end{longtable}
}

\section{\texorpdfstring{GET
\texorpdfstring{\softtext{/sedes/\{sedeId\}/gastos}}{/sedes/\{sedeId\}/gastos}}{GET }}\label{get-sedessedeidgastos}

\texorpdfstring{\softtext{listarGastos}}{listarGastos}

\textbf{Identificador:}
\texorpdfstring{\softcode{listarGastos}}{listarGastos}.

{\def\LTcaptype{none} % do not increment counter
\begin{longtable}[]{@{}
  >{\raggedright\arraybackslash}p{(\linewidth - 4\tabcolsep) * \real{0.3333}}
  >{\raggedright\arraybackslash}p{(\linewidth - 4\tabcolsep) * \real{0.3333}}
  >{\raggedright\arraybackslash}p{(\linewidth - 4\tabcolsep) * \real{0.3333}}@{}}
\toprule\noalign{}
\begin{minipage}[b]{\linewidth}\raggedright
Parámetro y ubicación
\end{minipage} & \begin{minipage}[b]{\linewidth}\raggedright
Tipo
\end{minipage} & \begin{minipage}[b]{\linewidth}\raggedright
Requerido y propósito
\end{minipage} \\
\midrule\noalign{}
\endhead
\bottomrule\noalign{}
\endlastfoot
\texorpdfstring{\softcode{sedeId}}{sedeId} (path) & string (uuid) &
Sí. \\
\texorpdfstring{\softcode{pagina}}{pagina} (query) & integer & No. \\
\texorpdfstring{\softcode{tamano}}{tamano} (query) & integer & No. \\
\texorpdfstring{\softcode{q}}{q} (query) & string &
\texorpdfstring{\softtext{No. Búsqueda}}{No. Búsqueda} literal acotada;
\% y \_ no actúan como comodines. \\
\end{longtable}
}

{\def\LTcaptype{none} % do not increment counter
\begin{longtable}[]{@{}
  >{\raggedright\arraybackslash}p{(\linewidth - 4\tabcolsep) * \real{0.3333}}
  >{\raggedright\arraybackslash}p{(\linewidth - 4\tabcolsep) * \real{0.3333}}
  >{\raggedright\arraybackslash}p{(\linewidth - 4\tabcolsep) * \real{0.3333}}@{}}
\toprule\noalign{}
\begin{minipage}[b]{\linewidth}\raggedright
HTTP
\end{minipage} & \begin{minipage}[b]{\linewidth}\raggedright
Significado declarado
\end{minipage} & \begin{minipage}[b]{\linewidth}\raggedright
Cuerpo y media type
\end{minipage} \\
\midrule\noalign{}
\endhead
\bottomrule\noalign{}
\endlastfoot
200 & OK &
\texorpdfstring{\softtext{application/json:}}{application/json:}
objeto \\
401 & Sin sesión &
\texorpdfstring{\softtext{application/problem+json:}}{application/problem+json:}
Error \\
403 & Sin permiso para la sede o la acción &
\texorpdfstring{\softtext{application/problem+json:}}{application/problem+json:}
Error \\
503 & Dependencia no disponible; una respuesta perdida puede requerir
comprobar resultado con la misma clave. &
\texorpdfstring{\softtext{application/problem+json:}}{application/problem+json:}
Error \\
400 & Forma o parámetro inválido &
\texorpdfstring{\softtext{application/problem+json:}}{application/problem+json:}
Error \\
\end{longtable}
}

\section{\texorpdfstring{GET
\texorpdfstring{\softtext{/sedes/\{sedeId\}/gastos/\{gastoId\}/reparto}}{/sedes/\{sedeId\}/gastos/\{gastoId\}/reparto}}{GET }}\label{get-sedessedeidgastosgastoidreparto}

Previsualizar reparto

\textbf{Identificador:}
\texorpdfstring{\softcode{previsualizarReparto}}{previsualizarReparto}.

{\def\LTcaptype{none} % do not increment counter
\begin{longtable}[]{@{}lll@{}}
\toprule\noalign{}
Parámetro y ubicación & Tipo & Requerido y propósito \\
\midrule\noalign{}
\endhead
\bottomrule\noalign{}
\endlastfoot
\texorpdfstring{\softcode{sedeId}}{sedeId} (path) & string (uuid) &
Sí. \\
\texorpdfstring{\softcode{gastoId}}{gastoId} (path) & string (uuid) &
Sí. \\
\texorpdfstring{\softcode{metodo}}{metodo} (query) & string & Sí. \\
\texorpdfstring{\softcode{pagina}}{pagina} (query) & integer & No. \\
\texorpdfstring{\softcode{tamano}}{tamano} (query) & integer & No. \\
\end{longtable}
}

{\def\LTcaptype{none} % do not increment counter
\begin{longtable}[]{@{}
  >{\raggedright\arraybackslash}p{(\linewidth - 4\tabcolsep) * \real{0.3333}}
  >{\raggedright\arraybackslash}p{(\linewidth - 4\tabcolsep) * \real{0.3333}}
  >{\raggedright\arraybackslash}p{(\linewidth - 4\tabcolsep) * \real{0.3333}}@{}}
\toprule\noalign{}
\begin{minipage}[b]{\linewidth}\raggedright
HTTP
\end{minipage} & \begin{minipage}[b]{\linewidth}\raggedright
Significado declarado
\end{minipage} & \begin{minipage}[b]{\linewidth}\raggedright
Cuerpo y media type
\end{minipage} \\
\midrule\noalign{}
\endhead
\bottomrule\noalign{}
\endlastfoot
200 & Asignaciones propuestas (sin persistir) &
\texorpdfstring{\softtext{application/json:}}{application/json:}
objeto \\
401 & Sin sesión &
\texorpdfstring{\softtext{application/problem+json:}}{application/problem+json:}
Error \\
403 & Sin permiso para la sede o la acción &
\texorpdfstring{\softtext{application/problem+json:}}{application/problem+json:}
Error \\
503 & Dependencia no disponible; una respuesta perdida puede requerir
comprobar resultado con la misma clave. &
\texorpdfstring{\softtext{application/problem+json:}}{application/problem+json:}
Error \\
400 & Forma o parámetro inválido &
\texorpdfstring{\softtext{application/problem+json:}}{application/problem+json:}
Error \\
\end{longtable}
}

\section{\texorpdfstring{POST
\texorpdfstring{\softtext{/sedes/\{sedeId\}/gastos/\{gastoId\}/reparto}}{/sedes/\{sedeId\}/gastos/\{gastoId\}/reparto}}{POST }}\label{post-sedessedeidgastosgastoidreparto}

Aprobar reparto

\textbf{Identificador:}
\texorpdfstring{\softcode{aprobarReparto}}{aprobarReparto}. Aprueba el
gasto y crea asignaciones y asientos según destino económico (incluye
costos tardíos).

{\def\LTcaptype{none} % do not increment counter
\begin{longtable}[]{@{}lll@{}}
\toprule\noalign{}
Parámetro y ubicación & Tipo & Requerido y propósito \\
\midrule\noalign{}
\endhead
\bottomrule\noalign{}
\endlastfoot
\texorpdfstring{\softcode{sedeId}}{sedeId} (path) & string (uuid) &
Sí. \\
\texorpdfstring{\softcode{gastoId}}{gastoId} (path) & string (uuid) &
Sí. \\
\texorpdfstring{\softcode{Idempotency-Key}}{Idempotency-Key} (header) &
string & Sí. \\
\texorpdfstring{\softcode{X-XSRF-TOKEN}}{X-XSRF-TOKEN} (header) & string
& Sí. \\
\end{longtable}
}

\textbf{Cuerpo
\texorpdfstring{\softtext{application/json:}}{application/json:}}
\texorpdfstring{\softtext{DistribucionInput.}}{DistribucionInput.}
\texorpdfstring{\softtext{Obligatorio.}}{Obligatorio.}

{\def\LTcaptype{none} % do not increment counter
\begin{longtable}[]{@{}
  >{\raggedright\arraybackslash}p{(\linewidth - 4\tabcolsep) * \real{0.3333}}
  >{\raggedright\arraybackslash}p{(\linewidth - 4\tabcolsep) * \real{0.3333}}
  >{\raggedright\arraybackslash}p{(\linewidth - 4\tabcolsep) * \real{0.3333}}@{}}
\toprule\noalign{}
\begin{minipage}[b]{\linewidth}\raggedright
HTTP
\end{minipage} & \begin{minipage}[b]{\linewidth}\raggedright
Significado declarado
\end{minipage} & \begin{minipage}[b]{\linewidth}\raggedright
Cuerpo y media type
\end{minipage} \\
\midrule\noalign{}
\endhead
\bottomrule\noalign{}
\endlastfoot
201 & Aprobado & Sin cuerpo descrito \\
422 & Viola una regla de negocio &
\texorpdfstring{\softtext{application/problem+json:}}{application/problem+json:}
Error \\
401 & Sin sesión &
\texorpdfstring{\softtext{application/problem+json:}}{application/problem+json:}
Error \\
403 & Sin permiso para la sede o la acción &
\texorpdfstring{\softtext{application/problem+json:}}{application/problem+json:}
Error \\
503 & Dependencia no disponible; una respuesta perdida puede requerir
comprobar resultado con la misma clave. &
\texorpdfstring{\softtext{application/problem+json:}}{application/problem+json:}
Error \\
400 & Forma o parámetro inválido &
\texorpdfstring{\softtext{application/problem+json:}}{application/problem+json:}
Error \\
\end{longtable}
}

\section{\texorpdfstring{POST
\texorpdfstring{\softtext{/sedes/\{sedeId\}/pagos}}{/sedes/\{sedeId\}/pagos}}{POST }}\label{post-sedessedeidpagos}

Registrar pago

\textbf{Identificador:}
\texorpdfstring{\softcode{registrarPago}}{registrarPago}.

{\def\LTcaptype{none} % do not increment counter
\begin{longtable}[]{@{}lll@{}}
\toprule\noalign{}
Parámetro y ubicación & Tipo & Requerido y propósito \\
\midrule\noalign{}
\endhead
\bottomrule\noalign{}
\endlastfoot
\texorpdfstring{\softcode{sedeId}}{sedeId} (path) & string (uuid) &
Sí. \\
\texorpdfstring{\softcode{Idempotency-Key}}{Idempotency-Key} (header) &
string & Sí. \\
\texorpdfstring{\softcode{X-XSRF-TOKEN}}{X-XSRF-TOKEN} (header) & string
& Sí. \\
\end{longtable}
}

\textbf{Cuerpo
\texorpdfstring{\softtext{application/json:}}{application/json:}}
PagoNuevo. \texorpdfstring{\softtext{Obligatorio.}}{Obligatorio.}

{\def\LTcaptype{none} % do not increment counter
\begin{longtable}[]{@{}
  >{\raggedright\arraybackslash}p{(\linewidth - 4\tabcolsep) * \real{0.3333}}
  >{\raggedright\arraybackslash}p{(\linewidth - 4\tabcolsep) * \real{0.3333}}
  >{\raggedright\arraybackslash}p{(\linewidth - 4\tabcolsep) * \real{0.3333}}@{}}
\toprule\noalign{}
\begin{minipage}[b]{\linewidth}\raggedright
HTTP
\end{minipage} & \begin{minipage}[b]{\linewidth}\raggedright
Significado declarado
\end{minipage} & \begin{minipage}[b]{\linewidth}\raggedright
Cuerpo y media type
\end{minipage} \\
\midrule\noalign{}
\endhead
\bottomrule\noalign{}
\endlastfoot
201 & Registrado &
\texorpdfstring{\softtext{application/json:}}{application/json:} Pago \\
401 & Sin sesión &
\texorpdfstring{\softtext{application/problem+json:}}{application/problem+json:}
Error \\
403 & Sin permiso para la sede o la acción &
\texorpdfstring{\softtext{application/problem+json:}}{application/problem+json:}
Error \\
503 & Dependencia no disponible; una respuesta perdida puede requerir
comprobar resultado con la misma clave. &
\texorpdfstring{\softtext{application/problem+json:}}{application/problem+json:}
Error \\
422 & Viola una regla de negocio &
\texorpdfstring{\softtext{application/problem+json:}}{application/problem+json:}
Error \\
400 & Forma o parámetro inválido &
\texorpdfstring{\softtext{application/problem+json:}}{application/problem+json:}
Error \\
\end{longtable}
}

\section{\texorpdfstring{GET
\texorpdfstring{\softtext{/sedes/\{sedeId\}/pagos}}{/sedes/\{sedeId\}/pagos}}{GET }}\label{get-sedessedeidpagos}

\texorpdfstring{\softtext{listarPagos}}{listarPagos}

\textbf{Identificador:}
\texorpdfstring{\softcode{listarPagos}}{listarPagos}.

{\def\LTcaptype{none} % do not increment counter
\begin{longtable}[]{@{}
  >{\raggedright\arraybackslash}p{(\linewidth - 4\tabcolsep) * \real{0.3333}}
  >{\raggedright\arraybackslash}p{(\linewidth - 4\tabcolsep) * \real{0.3333}}
  >{\raggedright\arraybackslash}p{(\linewidth - 4\tabcolsep) * \real{0.3333}}@{}}
\toprule\noalign{}
\begin{minipage}[b]{\linewidth}\raggedright
Parámetro y ubicación
\end{minipage} & \begin{minipage}[b]{\linewidth}\raggedright
Tipo
\end{minipage} & \begin{minipage}[b]{\linewidth}\raggedright
Requerido y propósito
\end{minipage} \\
\midrule\noalign{}
\endhead
\bottomrule\noalign{}
\endlastfoot
\texorpdfstring{\softcode{sedeId}}{sedeId} (path) & string (uuid) &
Sí. \\
\texorpdfstring{\softcode{pagina}}{pagina} (query) & integer & No. \\
\texorpdfstring{\softcode{tamano}}{tamano} (query) & integer & No. \\
\texorpdfstring{\softcode{q}}{q} (query) & string &
\texorpdfstring{\softtext{No. Búsqueda}}{No. Búsqueda} literal acotada;
\% y \_ no actúan como comodines. \\
\texorpdfstring{\softcode{estado}}{estado} (query) & string & No. \\
\texorpdfstring{\softcode{tipo}}{tipo} (query) & string & No. \\
\end{longtable}
}

{\def\LTcaptype{none} % do not increment counter
\begin{longtable}[]{@{}
  >{\raggedright\arraybackslash}p{(\linewidth - 4\tabcolsep) * \real{0.3333}}
  >{\raggedright\arraybackslash}p{(\linewidth - 4\tabcolsep) * \real{0.3333}}
  >{\raggedright\arraybackslash}p{(\linewidth - 4\tabcolsep) * \real{0.3333}}@{}}
\toprule\noalign{}
\begin{minipage}[b]{\linewidth}\raggedright
HTTP
\end{minipage} & \begin{minipage}[b]{\linewidth}\raggedright
Significado declarado
\end{minipage} & \begin{minipage}[b]{\linewidth}\raggedright
Cuerpo y media type
\end{minipage} \\
\midrule\noalign{}
\endhead
\bottomrule\noalign{}
\endlastfoot
200 & OK &
\texorpdfstring{\softtext{application/json:}}{application/json:}
objeto \\
401 & Sin sesión &
\texorpdfstring{\softtext{application/problem+json:}}{application/problem+json:}
Error \\
403 & Sin permiso para la sede o la acción &
\texorpdfstring{\softtext{application/problem+json:}}{application/problem+json:}
Error \\
503 & Dependencia no disponible; una respuesta perdida puede requerir
comprobar resultado con la misma clave. &
\texorpdfstring{\softtext{application/problem+json:}}{application/problem+json:}
Error \\
400 & Forma o parámetro inválido &
\texorpdfstring{\softtext{application/problem+json:}}{application/problem+json:}
Error \\
\end{longtable}
}

\section{\texorpdfstring{POST
\texorpdfstring{\softtext{/sedes/\{sedeId\}/pagos/\{pagoId\}/aplicaciones}}{/sedes/\{sedeId\}/pagos/\{pagoId\}/aplicaciones}}{POST }}\label{post-sedessedeidpagospagoidaplicaciones}

Aplicar pago

\textbf{Identificador:}
\texorpdfstring{\softcode{aplicarPago}}{aplicarPago}.

{\def\LTcaptype{none} % do not increment counter
\begin{longtable}[]{@{}lll@{}}
\toprule\noalign{}
Parámetro y ubicación & Tipo & Requerido y propósito \\
\midrule\noalign{}
\endhead
\bottomrule\noalign{}
\endlastfoot
\texorpdfstring{\softcode{sedeId}}{sedeId} (path) & string (uuid) &
Sí. \\
\texorpdfstring{\softcode{pagoId}}{pagoId} (path) & string (uuid) &
Sí. \\
\texorpdfstring{\softcode{Idempotency-Key}}{Idempotency-Key} (header) &
string & Sí. \\
\texorpdfstring{\softcode{X-XSRF-TOKEN}}{X-XSRF-TOKEN} (header) & string
& Sí. \\
\end{longtable}
}

\textbf{Cuerpo
\texorpdfstring{\softtext{application/json:}}{application/json:}}
object. \texorpdfstring{\softtext{Obligatorio.}}{Obligatorio.}

{\def\LTcaptype{none} % do not increment counter
\begin{longtable}[]{@{}
  >{\raggedright\arraybackslash}p{(\linewidth - 4\tabcolsep) * \real{0.3333}}
  >{\raggedright\arraybackslash}p{(\linewidth - 4\tabcolsep) * \real{0.3333}}
  >{\raggedright\arraybackslash}p{(\linewidth - 4\tabcolsep) * \real{0.3333}}@{}}
\toprule\noalign{}
\begin{minipage}[b]{\linewidth}\raggedright
HTTP
\end{minipage} & \begin{minipage}[b]{\linewidth}\raggedright
Significado declarado
\end{minipage} & \begin{minipage}[b]{\linewidth}\raggedright
Cuerpo y media type
\end{minipage} \\
\midrule\noalign{}
\endhead
\bottomrule\noalign{}
\endlastfoot
201 & Aplicado &
\texorpdfstring{\softtext{application/json:}}{application/json:} Pago \\
422 & Excede disponible o saldo (RN-029) &
\texorpdfstring{\softtext{application/problem+json:}}{application/problem+json:}
Error \\
409 & Conflicto de concurrencia,
\texorpdfstring{\softtext{version\_lock}}{version\_lock} o idempotencia
(misma clave, payload distinto) &
\texorpdfstring{\softtext{application/problem+json:}}{application/problem+json:}
Error \\
401 & Sin sesión &
\texorpdfstring{\softtext{application/problem+json:}}{application/problem+json:}
Error \\
403 & Sin permiso para la sede o la acción &
\texorpdfstring{\softtext{application/problem+json:}}{application/problem+json:}
Error \\
503 & Dependencia no disponible; una respuesta perdida puede requerir
comprobar resultado con la misma clave. &
\texorpdfstring{\softtext{application/problem+json:}}{application/problem+json:}
Error \\
400 & Forma o parámetro inválido &
\texorpdfstring{\softtext{application/problem+json:}}{application/problem+json:}
Error \\
\end{longtable}
}

\section{\texorpdfstring{POST
\texorpdfstring{\softtext{/sedes/\{sedeId\}/aplicaciones/\{aplicacionId\}/reverso}}{/sedes/\{sedeId\}/aplicaciones/\{aplicacionId\}/reverso}}{POST }}\label{post-sedessedeidaplicacionesaplicacionidreverso}

Revertir aplicacion

\textbf{Identificador:}
\texorpdfstring{\softcode{revertirAplicacion}}{revertirAplicacion}.

{\def\LTcaptype{none} % do not increment counter
\begin{longtable}[]{@{}lll@{}}
\toprule\noalign{}
Parámetro y ubicación & Tipo & Requerido y propósito \\
\midrule\noalign{}
\endhead
\bottomrule\noalign{}
\endlastfoot
\texorpdfstring{\softcode{sedeId}}{sedeId} (path) & string (uuid) &
Sí. \\
\texorpdfstring{\softcode{aplicacionId}}{aplicacionId} (path) & string
(uuid) & Sí. \\
\texorpdfstring{\softcode{Idempotency-Key}}{Idempotency-Key} (header) &
string & Sí. \\
\texorpdfstring{\softcode{X-XSRF-TOKEN}}{X-XSRF-TOKEN} (header) & string
& Sí. \\
\end{longtable}
}

\textbf{Cuerpo
\texorpdfstring{\softtext{application/json:}}{application/json:}}
object. \texorpdfstring{\softtext{Obligatorio.}}{Obligatorio.}

{\def\LTcaptype{none} % do not increment counter
\begin{longtable}[]{@{}
  >{\raggedright\arraybackslash}p{(\linewidth - 4\tabcolsep) * \real{0.3333}}
  >{\raggedright\arraybackslash}p{(\linewidth - 4\tabcolsep) * \real{0.3333}}
  >{\raggedright\arraybackslash}p{(\linewidth - 4\tabcolsep) * \real{0.3333}}@{}}
\toprule\noalign{}
\begin{minipage}[b]{\linewidth}\raggedright
HTTP
\end{minipage} & \begin{minipage}[b]{\linewidth}\raggedright
Significado declarado
\end{minipage} & \begin{minipage}[b]{\linewidth}\raggedright
Cuerpo y media type
\end{minipage} \\
\midrule\noalign{}
\endhead
\bottomrule\noalign{}
\endlastfoot
201 & Reverso completo registrado & Sin cuerpo descrito \\
409 & Conflicto de concurrencia,
\texorpdfstring{\softtext{version\_lock}}{version\_lock} o idempotencia
(misma clave, payload distinto) &
\texorpdfstring{\softtext{application/problem+json:}}{application/problem+json:}
Error \\
401 & Sin sesión &
\texorpdfstring{\softtext{application/problem+json:}}{application/problem+json:}
Error \\
403 & Sin permiso para la sede o la acción &
\texorpdfstring{\softtext{application/problem+json:}}{application/problem+json:}
Error \\
503 & Dependencia no disponible; una respuesta perdida puede requerir
comprobar resultado con la misma clave. &
\texorpdfstring{\softtext{application/problem+json:}}{application/problem+json:}
Error \\
422 & Viola una regla de negocio &
\texorpdfstring{\softtext{application/problem+json:}}{application/problem+json:}
Error \\
400 & Forma o parámetro inválido &
\texorpdfstring{\softtext{application/problem+json:}}{application/problem+json:}
Error \\
\end{longtable}
}

\section{\texorpdfstring{POST
\texorpdfstring{\softtext{/sedes/\{sedeId\}/pagos/\{pagoId\}/reembolsos}}{/sedes/\{sedeId\}/pagos/\{pagoId\}/reembolsos}}{POST }}\label{post-sedessedeidpagospagoidreembolsos}

Registrar reembolso

\textbf{Identificador:}
\texorpdfstring{\softcode{registrarReembolso}}{registrarReembolso}.

{\def\LTcaptype{none} % do not increment counter
\begin{longtable}[]{@{}lll@{}}
\toprule\noalign{}
Parámetro y ubicación & Tipo & Requerido y propósito \\
\midrule\noalign{}
\endhead
\bottomrule\noalign{}
\endlastfoot
\texorpdfstring{\softcode{sedeId}}{sedeId} (path) & string (uuid) &
Sí. \\
\texorpdfstring{\softcode{pagoId}}{pagoId} (path) & string (uuid) &
Sí. \\
\texorpdfstring{\softcode{Idempotency-Key}}{Idempotency-Key} (header) &
string & Sí. \\
\texorpdfstring{\softcode{X-XSRF-TOKEN}}{X-XSRF-TOKEN} (header) & string
& Sí. \\
\end{longtable}
}

\textbf{Cuerpo
\texorpdfstring{\softtext{application/json:}}{application/json:}}
object. \texorpdfstring{\softtext{Obligatorio.}}{Obligatorio.}

{\def\LTcaptype{none} % do not increment counter
\begin{longtable}[]{@{}
  >{\raggedright\arraybackslash}p{(\linewidth - 4\tabcolsep) * \real{0.3333}}
  >{\raggedright\arraybackslash}p{(\linewidth - 4\tabcolsep) * \real{0.3333}}
  >{\raggedright\arraybackslash}p{(\linewidth - 4\tabcolsep) * \real{0.3333}}@{}}
\toprule\noalign{}
\begin{minipage}[b]{\linewidth}\raggedright
HTTP
\end{minipage} & \begin{minipage}[b]{\linewidth}\raggedright
Significado declarado
\end{minipage} & \begin{minipage}[b]{\linewidth}\raggedright
Cuerpo y media type
\end{minipage} \\
\midrule\noalign{}
\endhead
\bottomrule\noalign{}
\endlastfoot
201 & Reembolso registrado &
\texorpdfstring{\softtext{application/json:}}{application/json:} Pago \\
422 & Viola una regla de negocio &
\texorpdfstring{\softtext{application/problem+json:}}{application/problem+json:}
Error \\
401 & Sin sesión &
\texorpdfstring{\softtext{application/problem+json:}}{application/problem+json:}
Error \\
403 & Sin permiso para la sede o la acción &
\texorpdfstring{\softtext{application/problem+json:}}{application/problem+json:}
Error \\
503 & Dependencia no disponible; una respuesta perdida puede requerir
comprobar resultado con la misma clave. &
\texorpdfstring{\softtext{application/problem+json:}}{application/problem+json:}
Error \\
400 & Forma o parámetro inválido &
\texorpdfstring{\softtext{application/problem+json:}}{application/problem+json:}
Error \\
\end{longtable}
}

\section{\texorpdfstring{GET
\texorpdfstring{\softtext{/sedes/\{sedeId\}/reportes/cartera}}{/sedes/\{sedeId\}/reportes/cartera}}{GET }}\label{get-sedessedeidreportescartera}

Reporte cartera

\textbf{Identificador:}
\texorpdfstring{\softcode{reporteCartera}}{reporteCartera}.

{\def\LTcaptype{none} % do not increment counter
\begin{longtable}[]{@{}lll@{}}
\toprule\noalign{}
Parámetro y ubicación & Tipo & Requerido y propósito \\
\midrule\noalign{}
\endhead
\bottomrule\noalign{}
\endlastfoot
\texorpdfstring{\softcode{sedeId}}{sedeId} (path) & string (uuid) &
Sí. \\
\texorpdfstring{\softcode{pagina}}{pagina} (query) & integer & No. \\
\texorpdfstring{\softcode{tamano}}{tamano} (query) & integer & No. \\
\texorpdfstring{\softcode{direccion}}{direccion} (query) & string &
No. \\
\texorpdfstring{\softcode{vencida}}{vencida} (query) & boolean & No. \\
\texorpdfstring{\softcode{q}}{q} (query) & string &
\texorpdfstring{\softtext{No. Filtro}}{No. Filtro} de búsqueda
acotado. \\
\end{longtable}
}

{\def\LTcaptype{none} % do not increment counter
\begin{longtable}[]{@{}
  >{\raggedright\arraybackslash}p{(\linewidth - 4\tabcolsep) * \real{0.3333}}
  >{\raggedright\arraybackslash}p{(\linewidth - 4\tabcolsep) * \real{0.3333}}
  >{\raggedright\arraybackslash}p{(\linewidth - 4\tabcolsep) * \real{0.3333}}@{}}
\toprule\noalign{}
\begin{minipage}[b]{\linewidth}\raggedright
HTTP
\end{minipage} & \begin{minipage}[b]{\linewidth}\raggedright
Significado declarado
\end{minipage} & \begin{minipage}[b]{\linewidth}\raggedright
Cuerpo y media type
\end{minipage} \\
\midrule\noalign{}
\endhead
\bottomrule\noalign{}
\endlastfoot
200 & Cartera con corte &
\texorpdfstring{\softtext{application/json:}}{application/json:}
\texorpdfstring{\softtext{ReporteCartera}}{ReporteCartera} \\
401 & Sin sesión &
\texorpdfstring{\softtext{application/problem+json:}}{application/problem+json:}
Error \\
403 & Sin permiso para la sede o la acción &
\texorpdfstring{\softtext{application/problem+json:}}{application/problem+json:}
Error \\
503 & Dependencia no disponible; una respuesta perdida puede requerir
comprobar resultado con la misma clave. &
\texorpdfstring{\softtext{application/problem+json:}}{application/problem+json:}
Error \\
400 & Forma o parámetro inválido &
\texorpdfstring{\softtext{application/problem+json:}}{application/problem+json:}
Error \\
\end{longtable}
}

\section{\texorpdfstring{GET
\texorpdfstring{\softtext{/sedes/\{sedeId\}/reportes/margen}}{/sedes/\{sedeId\}/reportes/margen}}{GET }}\label{get-sedessedeidreportesmargen}

Reporte margen

\textbf{Identificador:}
\texorpdfstring{\softcode{reporteMargen}}{reporteMargen}.

{\def\LTcaptype{none} % do not increment counter
\begin{longtable}[]{@{}lll@{}}
\toprule\noalign{}
Parámetro y ubicación & Tipo & Requerido y propósito \\
\midrule\noalign{}
\endhead
\bottomrule\noalign{}
\endlastfoot
\texorpdfstring{\softcode{sedeId}}{sedeId} (path) & string (uuid) &
Sí. \\
\texorpdfstring{\softcode{operacion\_id}}{operacion\_id} (query) &
string (uuid) & No. \\
\texorpdfstring{\softcode{periodo\_id}}{periodo\_id} (query) & string
(uuid) & No. \\
\texorpdfstring{\softcode{q}}{q} (query) & string &
\texorpdfstring{\softtext{No. Filtro}}{No. Filtro} de búsqueda
acotado. \\
\end{longtable}
}

{\def\LTcaptype{none} % do not increment counter
\begin{longtable}[]{@{}
  >{\raggedright\arraybackslash}p{(\linewidth - 4\tabcolsep) * \real{0.3333}}
  >{\raggedright\arraybackslash}p{(\linewidth - 4\tabcolsep) * \real{0.3333}}
  >{\raggedright\arraybackslash}p{(\linewidth - 4\tabcolsep) * \real{0.3333}}@{}}
\toprule\noalign{}
\begin{minipage}[b]{\linewidth}\raggedright
HTTP
\end{minipage} & \begin{minipage}[b]{\linewidth}\raggedright
Significado declarado
\end{minipage} & \begin{minipage}[b]{\linewidth}\raggedright
Cuerpo y media type
\end{minipage} \\
\midrule\noalign{}
\endhead
\bottomrule\noalign{}
\endlastfoot
200 & Margen etiquetado
\texorpdfstring{\softtext{PROVISIONAL/ACTUALIZADO/CERRADO}}{PROVISIONAL/ACTUALIZADO/CERRADO}
& \texorpdfstring{\softtext{application/json:}}{application/json:}
\texorpdfstring{\softtext{ReporteMargen}}{ReporteMargen} \\
401 & Sin sesión &
\texorpdfstring{\softtext{application/problem+json:}}{application/problem+json:}
Error \\
403 & Sin permiso para la sede o la acción &
\texorpdfstring{\softtext{application/problem+json:}}{application/problem+json:}
Error \\
503 & Dependencia no disponible; una respuesta perdida puede requerir
comprobar resultado con la misma clave. &
\texorpdfstring{\softtext{application/problem+json:}}{application/problem+json:}
Error \\
400 & Forma o parámetro inválido &
\texorpdfstring{\softtext{application/problem+json:}}{application/problem+json:}
Error \\
\end{longtable}
}

\section{\texorpdfstring{POST
\texorpdfstring{\softtext{/sedes/\{sedeId\}/periodos/\{periodoId\}/cierre}}{/sedes/\{sedeId\}/periodos/\{periodoId\}/cierre}}{POST }}\label{post-sedessedeidperiodosperiodoidcierre}

Cerrar periodo

\textbf{Identificador:}
\texorpdfstring{\softcode{cerrarPeriodo}}{cerrarPeriodo}.

{\def\LTcaptype{none} % do not increment counter
\begin{longtable}[]{@{}lll@{}}
\toprule\noalign{}
Parámetro y ubicación & Tipo & Requerido y propósito \\
\midrule\noalign{}
\endhead
\bottomrule\noalign{}
\endlastfoot
\texorpdfstring{\softcode{sedeId}}{sedeId} (path) & string (uuid) &
Sí. \\
\texorpdfstring{\softcode{periodoId}}{periodoId} (path) & string (uuid)
& Sí. \\
\texorpdfstring{\softcode{Idempotency-Key}}{Idempotency-Key} (header) &
string & Sí. \\
\texorpdfstring{\softcode{X-XSRF-TOKEN}}{X-XSRF-TOKEN} (header) & string
& Sí. \\
\end{longtable}
}

\textbf{Cuerpo
\texorpdfstring{\softtext{application/json:}}{application/json:}}
object. \texorpdfstring{\softtext{Obligatorio.}}{Obligatorio.}

{\def\LTcaptype{none} % do not increment counter
\begin{longtable}[]{@{}
  >{\raggedright\arraybackslash}p{(\linewidth - 4\tabcolsep) * \real{0.3333}}
  >{\raggedright\arraybackslash}p{(\linewidth - 4\tabcolsep) * \real{0.3333}}
  >{\raggedright\arraybackslash}p{(\linewidth - 4\tabcolsep) * \real{0.3333}}@{}}
\toprule\noalign{}
\begin{minipage}[b]{\linewidth}\raggedright
HTTP
\end{minipage} & \begin{minipage}[b]{\linewidth}\raggedright
Significado declarado
\end{minipage} & \begin{minipage}[b]{\linewidth}\raggedright
Cuerpo y media type
\end{minipage} \\
\midrule\noalign{}
\endhead
\bottomrule\noalign{}
\endlastfoot
201 & Versión de cierre creada & Sin cuerpo descrito \\
409 & Conflicto de concurrencia,
\texorpdfstring{\softtext{version\_lock}}{version\_lock} o idempotencia
(misma clave, payload distinto) &
\texorpdfstring{\softtext{application/problem+json:}}{application/problem+json:}
Error \\
401 & Sin sesión &
\texorpdfstring{\softtext{application/problem+json:}}{application/problem+json:}
Error \\
403 & Sin permiso para la sede o la acción &
\texorpdfstring{\softtext{application/problem+json:}}{application/problem+json:}
Error \\
503 & Dependencia no disponible; una respuesta perdida puede requerir
comprobar resultado con la misma clave. &
\texorpdfstring{\softtext{application/problem+json:}}{application/problem+json:}
Error \\
422 & Viola una regla de negocio &
\texorpdfstring{\softtext{application/problem+json:}}{application/problem+json:}
Error \\
400 & Forma o parámetro inválido &
\texorpdfstring{\softtext{application/problem+json:}}{application/problem+json:}
Error \\
\end{longtable}
}

\section{\texorpdfstring{POST
\texorpdfstring{\softtext{/sedes/\{sedeId\}/documentos}}{/sedes/\{sedeId\}/documentos}}{POST }}\label{post-sedessedeiddocumentos}

Subir documento

\textbf{Identificador:}
\texorpdfstring{\softcode{subirDocumento}}{subirDocumento}.

{\def\LTcaptype{none} % do not increment counter
\begin{longtable}[]{@{}lll@{}}
\toprule\noalign{}
Parámetro y ubicación & Tipo & Requerido y propósito \\
\midrule\noalign{}
\endhead
\bottomrule\noalign{}
\endlastfoot
\texorpdfstring{\softcode{sedeId}}{sedeId} (path) & string (uuid) &
Sí. \\
\texorpdfstring{\softcode{Idempotency-Key}}{Idempotency-Key} (header) &
string & Sí. \\
\texorpdfstring{\softcode{X-XSRF-TOKEN}}{X-XSRF-TOKEN} (header) & string
& Sí. \\
\end{longtable}
}

\textbf{Cuerpo
\texorpdfstring{\softtext{multipart/form-data:}}{multipart/form-data:}}
object. \texorpdfstring{\softtext{Obligatorio.}}{Obligatorio.}

{\def\LTcaptype{none} % do not increment counter
\begin{longtable}[]{@{}
  >{\raggedright\arraybackslash}p{(\linewidth - 4\tabcolsep) * \real{0.3333}}
  >{\raggedright\arraybackslash}p{(\linewidth - 4\tabcolsep) * \real{0.3333}}
  >{\raggedright\arraybackslash}p{(\linewidth - 4\tabcolsep) * \real{0.3333}}@{}}
\toprule\noalign{}
\begin{minipage}[b]{\linewidth}\raggedright
HTTP
\end{minipage} & \begin{minipage}[b]{\linewidth}\raggedright
Significado declarado
\end{minipage} & \begin{minipage}[b]{\linewidth}\raggedright
Cuerpo y media type
\end{minipage} \\
\midrule\noalign{}
\endhead
\bottomrule\noalign{}
\endlastfoot
201 & Documento (metadatos y sha256) &
\texorpdfstring{\softtext{application/json:}}{application/json:}
Documento \\
422 & Viola una regla de negocio &
\texorpdfstring{\softtext{application/problem+json:}}{application/problem+json:}
Error \\
401 & Sin sesión &
\texorpdfstring{\softtext{application/problem+json:}}{application/problem+json:}
Error \\
403 & Sin permiso para la sede o la acción &
\texorpdfstring{\softtext{application/problem+json:}}{application/problem+json:}
Error \\
503 & Dependencia no disponible; una respuesta perdida puede requerir
comprobar resultado con la misma clave. &
\texorpdfstring{\softtext{application/problem+json:}}{application/problem+json:}
Error \\
413 & Archivo excede 20 MiB &
\texorpdfstring{\softtext{application/problem+json:}}{application/problem+json:}
Error \\
415 & Contenido real de archivo no admitido &
\texorpdfstring{\softtext{application/problem+json:}}{application/problem+json:}
Error \\
400 & Forma o parámetro inválido &
\texorpdfstring{\softtext{application/problem+json:}}{application/problem+json:}
Error \\
\end{longtable}
}

\section{\texorpdfstring{GET
\texorpdfstring{\softtext{/sedes/\{sedeId\}/documentos}}{/sedes/\{sedeId\}/documentos}}{GET }}\label{get-sedessedeiddocumentos}

Consultar versiones de documentos privadas

\textbf{Identificador:}
\texorpdfstring{\softcode{listarDocumentos}}{listarDocumentos}.

{\def\LTcaptype{none} % do not increment counter
\begin{longtable}[]{@{}
  >{\raggedright\arraybackslash}p{(\linewidth - 4\tabcolsep) * \real{0.3333}}
  >{\raggedright\arraybackslash}p{(\linewidth - 4\tabcolsep) * \real{0.3333}}
  >{\raggedright\arraybackslash}p{(\linewidth - 4\tabcolsep) * \real{0.3333}}@{}}
\toprule\noalign{}
\begin{minipage}[b]{\linewidth}\raggedright
Parámetro y ubicación
\end{minipage} & \begin{minipage}[b]{\linewidth}\raggedright
Tipo
\end{minipage} & \begin{minipage}[b]{\linewidth}\raggedright
Requerido y propósito
\end{minipage} \\
\midrule\noalign{}
\endhead
\bottomrule\noalign{}
\endlastfoot
\texorpdfstring{\softcode{sedeId}}{sedeId} (path) & string (uuid) &
Sí. \\
\texorpdfstring{\softcode{q}}{q} (query) & string &
\texorpdfstring{\softtext{No. Búsqueda}}{No. Búsqueda} literal acotada;
\% y \_ no actúan como comodines. \\
\end{longtable}
}

{\def\LTcaptype{none} % do not increment counter
\begin{longtable}[]{@{}
  >{\raggedright\arraybackslash}p{(\linewidth - 4\tabcolsep) * \real{0.3333}}
  >{\raggedright\arraybackslash}p{(\linewidth - 4\tabcolsep) * \real{0.3333}}
  >{\raggedright\arraybackslash}p{(\linewidth - 4\tabcolsep) * \real{0.3333}}@{}}
\toprule\noalign{}
\begin{minipage}[b]{\linewidth}\raggedright
HTTP
\end{minipage} & \begin{minipage}[b]{\linewidth}\raggedright
Significado declarado
\end{minipage} & \begin{minipage}[b]{\linewidth}\raggedright
Cuerpo y media type
\end{minipage} \\
\midrule\noalign{}
\endhead
\bottomrule\noalign{}
\endlastfoot
200 & Resultado persistido del servidor &
\texorpdfstring{\softtext{application/json:}}{application/json:}
objeto \\
401 & Sin sesión &
\texorpdfstring{\softtext{application/problem+json:}}{application/problem+json:}
Error \\
403 & Sin permiso para la sede o la acción &
\texorpdfstring{\softtext{application/problem+json:}}{application/problem+json:}
Error \\
503 & Dependencia no disponible; una respuesta perdida puede requerir
comprobar resultado con la misma clave. &
\texorpdfstring{\softtext{application/problem+json:}}{application/problem+json:}
Error \\
400 & Forma o parámetro inválido &
\texorpdfstring{\softtext{application/problem+json:}}{application/problem+json:}
Error \\
\end{longtable}
}

\section{\texorpdfstring{GET
\texorpdfstring{\softtext{/sedes/\{sedeId\}/documentos/\{documentoId\}/contenido}}{/sedes/\{sedeId\}/documentos/\{documentoId\}/contenido}}{GET }}\label{get-sedessedeiddocumentosdocumentoidcontenido}

Descargar documento

\textbf{Identificador:}
\texorpdfstring{\softcode{descargarDocumento}}{descargarDocumento}.

{\def\LTcaptype{none} % do not increment counter
\begin{longtable}[]{@{}lll@{}}
\toprule\noalign{}
Parámetro y ubicación & Tipo & Requerido y propósito \\
\midrule\noalign{}
\endhead
\bottomrule\noalign{}
\endlastfoot
\texorpdfstring{\softcode{sedeId}}{sedeId} (path) & string (uuid) &
Sí. \\
\texorpdfstring{\softcode{documentoId}}{documentoId} (path) & string
(uuid) & Sí. \\
\end{longtable}
}

{\def\LTcaptype{none} % do not increment counter
\begin{longtable}[]{@{}
  >{\raggedright\arraybackslash}p{(\linewidth - 4\tabcolsep) * \real{0.3333}}
  >{\raggedright\arraybackslash}p{(\linewidth - 4\tabcolsep) * \real{0.3333}}
  >{\raggedright\arraybackslash}p{(\linewidth - 4\tabcolsep) * \real{0.3333}}@{}}
\toprule\noalign{}
\begin{minipage}[b]{\linewidth}\raggedright
HTTP
\end{minipage} & \begin{minipage}[b]{\linewidth}\raggedright
Significado declarado
\end{minipage} & \begin{minipage}[b]{\linewidth}\raggedright
Cuerpo y media type
\end{minipage} \\
\midrule\noalign{}
\endhead
\bottomrule\noalign{}
\endlastfoot
200 & Contenido autorizado &
\texorpdfstring{\softtext{application/octet-stream:}}{application/octet-stream:}
string (binary) \\
404 & No existe o no es accesible (política uniforme) &
\texorpdfstring{\softtext{application/problem+json:}}{application/problem+json:}
Error \\
401 & Sin sesión &
\texorpdfstring{\softtext{application/problem+json:}}{application/problem+json:}
Error \\
403 & Sin permiso para la sede o la acción &
\texorpdfstring{\softtext{application/problem+json:}}{application/problem+json:}
Error \\
503 & Dependencia no disponible; una respuesta perdida puede requerir
comprobar resultado con la misma clave. &
\texorpdfstring{\softtext{application/problem+json:}}{application/problem+json:}
Error \\
400 & Forma o parámetro inválido &
\texorpdfstring{\softtext{application/problem+json:}}{application/problem+json:}
Error \\
\end{longtable}
}

\section{\texorpdfstring{GET
\texorpdfstring{\softtext{/sedes/\{sedeId\}/integraciones/proveedor/tarifas}}{/sedes/\{sedeId\}/integraciones/proveedor/tarifas}}{GET }}\label{get-sedessedeidintegracionesproveedortarifas}

Consultar tarifas proveedor

\textbf{Identificador:}
\texorpdfstring{\softcode{consultarTarifasProveedor}}{consultarTarifasProveedor}.
Consulta y captura técnica local de proveedor simulado por el dueño de
la sede; sin locks
\texorpdfstring{\softtext{financieros.}}{financieros.} No-store. Una
captura no constituye todavía un gasto.

{\def\LTcaptype{none} % do not increment counter
\begin{longtable}[]{@{}lll@{}}
\toprule\noalign{}
Parámetro y ubicación & Tipo & Requerido y propósito \\
\midrule\noalign{}
\endhead
\bottomrule\noalign{}
\endlastfoot
\texorpdfstring{\softcode{sedeId}}{sedeId} (path) & string (uuid) &
Sí. \\
\texorpdfstring{\softcode{ruta}}{ruta} (query) & string & No. \\
\texorpdfstring{\softcode{pagina}}{pagina} (query) & integer & No. \\
\texorpdfstring{\softcode{tamano}}{tamano} (query) & integer & No. \\
\end{longtable}
}

{\def\LTcaptype{none} % do not increment counter
\begin{longtable}[]{@{}
  >{\raggedright\arraybackslash}p{(\linewidth - 4\tabcolsep) * \real{0.3333}}
  >{\raggedright\arraybackslash}p{(\linewidth - 4\tabcolsep) * \real{0.3333}}
  >{\raggedright\arraybackslash}p{(\linewidth - 4\tabcolsep) * \real{0.3333}}@{}}
\toprule\noalign{}
\begin{minipage}[b]{\linewidth}\raggedright
HTTP
\end{minipage} & \begin{minipage}[b]{\linewidth}\raggedright
Significado declarado
\end{minipage} & \begin{minipage}[b]{\linewidth}\raggedright
Cuerpo y media type
\end{minipage} \\
\midrule\noalign{}
\endhead
\bottomrule\noalign{}
\endlastfoot
200 & Tarifas &
\texorpdfstring{\softtext{application/json:}}{application/json:}
objeto \\
503 & Dependencia no disponible; una respuesta perdida puede requerir
comprobar resultado con la misma clave. &
\texorpdfstring{\softtext{application/problem+json:}}{application/problem+json:}
Error \\
401 & Sin sesión &
\texorpdfstring{\softtext{application/problem+json:}}{application/problem+json:}
Error \\
403 & Sin permiso para la sede o la acción &
\texorpdfstring{\softtext{application/problem+json:}}{application/problem+json:}
Error \\
400 & Forma o parámetro inválido &
\texorpdfstring{\softtext{application/problem+json:}}{application/problem+json:}
Error \\
\end{longtable}
}

\section{\texorpdfstring{GET
\texorpdfstring{\softtext{/sesion/csrf}}{/sesion/csrf}}{GET }}\label{get-sesioncsrf}

\texorpdfstring{\softtext{obtenerCsrf}}{obtenerCsrf}

\textbf{Identificador:}
\texorpdfstring{\softcode{obtenerCsrf}}{obtenerCsrf}.

{\def\LTcaptype{none} % do not increment counter
\begin{longtable}[]{@{}
  >{\raggedright\arraybackslash}p{(\linewidth - 4\tabcolsep) * \real{0.3333}}
  >{\raggedright\arraybackslash}p{(\linewidth - 4\tabcolsep) * \real{0.3333}}
  >{\raggedright\arraybackslash}p{(\linewidth - 4\tabcolsep) * \real{0.3333}}@{}}
\toprule\noalign{}
\begin{minipage}[b]{\linewidth}\raggedright
HTTP
\end{minipage} & \begin{minipage}[b]{\linewidth}\raggedright
Significado declarado
\end{minipage} & \begin{minipage}[b]{\linewidth}\raggedright
Cuerpo y media type
\end{minipage} \\
\midrule\noalign{}
\endhead
\bottomrule\noalign{}
\endlastfoot
200 & OK &
\texorpdfstring{\softtext{application/json:}}{application/json:} Csrf \\
503 & Dependencia no disponible; una respuesta perdida puede requerir
comprobar resultado con la misma clave. &
\texorpdfstring{\softtext{application/problem+json:}}{application/problem+json:}
Error \\
\end{longtable}
}

\section{\texorpdfstring{GET
\texorpdfstring{\softtext{/sedes/\{sedeId\}/liquidaciones}}{/sedes/\{sedeId\}/liquidaciones}}{GET }}\label{get-sedessedeidliquidaciones}

\texorpdfstring{\softtext{listarLiquidaciones}}{listarLiquidaciones}

\textbf{Identificador:}
\texorpdfstring{\softcode{listarLiquidaciones}}{listarLiquidaciones}.

{\def\LTcaptype{none} % do not increment counter
\begin{longtable}[]{@{}
  >{\raggedright\arraybackslash}p{(\linewidth - 4\tabcolsep) * \real{0.3333}}
  >{\raggedright\arraybackslash}p{(\linewidth - 4\tabcolsep) * \real{0.3333}}
  >{\raggedright\arraybackslash}p{(\linewidth - 4\tabcolsep) * \real{0.3333}}@{}}
\toprule\noalign{}
\begin{minipage}[b]{\linewidth}\raggedright
Parámetro y ubicación
\end{minipage} & \begin{minipage}[b]{\linewidth}\raggedright
Tipo
\end{minipage} & \begin{minipage}[b]{\linewidth}\raggedright
Requerido y propósito
\end{minipage} \\
\midrule\noalign{}
\endhead
\bottomrule\noalign{}
\endlastfoot
\texorpdfstring{\softcode{sedeId}}{sedeId} (path) & string (uuid) &
Sí. \\
\texorpdfstring{\softcode{pagina}}{pagina} (query) & integer & No. \\
\texorpdfstring{\softcode{tamano}}{tamano} (query) & integer & No. \\
\texorpdfstring{\softcode{q}}{q} (query) & string &
\texorpdfstring{\softtext{No. Búsqueda}}{No. Búsqueda} literal acotada;
\% y \_ no actúan como comodines. \\
\end{longtable}
}

{\def\LTcaptype{none} % do not increment counter
\begin{longtable}[]{@{}
  >{\raggedright\arraybackslash}p{(\linewidth - 4\tabcolsep) * \real{0.3333}}
  >{\raggedright\arraybackslash}p{(\linewidth - 4\tabcolsep) * \real{0.3333}}
  >{\raggedright\arraybackslash}p{(\linewidth - 4\tabcolsep) * \real{0.3333}}@{}}
\toprule\noalign{}
\begin{minipage}[b]{\linewidth}\raggedright
HTTP
\end{minipage} & \begin{minipage}[b]{\linewidth}\raggedright
Significado declarado
\end{minipage} & \begin{minipage}[b]{\linewidth}\raggedright
Cuerpo y media type
\end{minipage} \\
\midrule\noalign{}
\endhead
\bottomrule\noalign{}
\endlastfoot
200 & OK &
\texorpdfstring{\softtext{application/json:}}{application/json:}
objeto \\
401 & Sin sesión &
\texorpdfstring{\softtext{application/problem+json:}}{application/problem+json:}
Error \\
403 & Sin permiso para la sede o la acción &
\texorpdfstring{\softtext{application/problem+json:}}{application/problem+json:}
Error \\
503 & Dependencia no disponible; una respuesta perdida puede requerir
comprobar resultado con la misma clave. &
\texorpdfstring{\softtext{application/problem+json:}}{application/problem+json:}
Error \\
400 & Forma o parámetro inválido &
\texorpdfstring{\softtext{application/problem+json:}}{application/problem+json:}
Error \\
\end{longtable}
}

\section{\texorpdfstring{POST
\texorpdfstring{\softtext{/sedes/\{sedeId\}/pagos/\{pagoId\}/validar}}{/sedes/\{sedeId\}/pagos/\{pagoId\}/validar}}{POST }}\label{post-sedessedeidpagospagoidvalidar}

Valida dinero registrado y su periodo

\textbf{Identificador:}
\texorpdfstring{\softcode{validarPago}}{validarPago}.

{\def\LTcaptype{none} % do not increment counter
\begin{longtable}[]{@{}lll@{}}
\toprule\noalign{}
Parámetro y ubicación & Tipo & Requerido y propósito \\
\midrule\noalign{}
\endhead
\bottomrule\noalign{}
\endlastfoot
\texorpdfstring{\softcode{sedeId}}{sedeId} (path) & string (uuid) &
Sí. \\
\texorpdfstring{\softcode{pagoId}}{pagoId} (path) & string (uuid) &
Sí. \\
\texorpdfstring{\softcode{Idempotency-Key}}{Idempotency-Key} (header) &
string & Sí. \\
\texorpdfstring{\softcode{X-XSRF-TOKEN}}{X-XSRF-TOKEN} (header) & string
& Sí. \\
\end{longtable}
}

\textbf{Cuerpo
\texorpdfstring{\softtext{application/json:}}{application/json:}}
\texorpdfstring{\softtext{Reconocimiento.}}{Reconocimiento.}
\texorpdfstring{\softtext{Obligatorio.}}{Obligatorio.}

{\def\LTcaptype{none} % do not increment counter
\begin{longtable}[]{@{}
  >{\raggedright\arraybackslash}p{(\linewidth - 4\tabcolsep) * \real{0.3333}}
  >{\raggedright\arraybackslash}p{(\linewidth - 4\tabcolsep) * \real{0.3333}}
  >{\raggedright\arraybackslash}p{(\linewidth - 4\tabcolsep) * \real{0.3333}}@{}}
\toprule\noalign{}
\begin{minipage}[b]{\linewidth}\raggedright
HTTP
\end{minipage} & \begin{minipage}[b]{\linewidth}\raggedright
Significado declarado
\end{minipage} & \begin{minipage}[b]{\linewidth}\raggedright
Cuerpo y media type
\end{minipage} \\
\midrule\noalign{}
\endhead
\bottomrule\noalign{}
\endlastfoot
200 & OK &
\texorpdfstring{\softtext{application/json:}}{application/json:} Pago \\
401 & Sin sesión &
\texorpdfstring{\softtext{application/problem+json:}}{application/problem+json:}
Error \\
403 & Sin permiso para la sede o la acción &
\texorpdfstring{\softtext{application/problem+json:}}{application/problem+json:}
Error \\
503 & Dependencia no disponible; una respuesta perdida puede requerir
comprobar resultado con la misma clave. &
\texorpdfstring{\softtext{application/problem+json:}}{application/problem+json:}
Error \\
422 & Viola una regla de negocio &
\texorpdfstring{\softtext{application/problem+json:}}{application/problem+json:}
Error \\
400 & Forma o parámetro inválido &
\texorpdfstring{\softtext{application/problem+json:}}{application/problem+json:}
Error \\
\end{longtable}
}

\section{\texorpdfstring{GET
\texorpdfstring{\softtext{/sedes/\{sedeId\}/comandos/\{comando\}/\{clave\}}}{/sedes/\{sedeId\}/comandos/\{comando\}/\{clave\}}}{GET }}\label{get-sedessedeidcomandoscomandoclave}

Recuperar un resultado propio comprometido

\textbf{Identificador:}
\texorpdfstring{\softcode{consultarResultadoComando}}{consultarResultadoComando}.
Un 404 no prueba rollback de una petición anterior. Reintentar conserva
clave, contenido y ETag.

{\def\LTcaptype{none} % do not increment counter
\begin{longtable}[]{@{}lll@{}}
\toprule\noalign{}
Parámetro y ubicación & Tipo & Requerido y propósito \\
\midrule\noalign{}
\endhead
\bottomrule\noalign{}
\endlastfoot
\texorpdfstring{\softcode{sedeId}}{sedeId} (path) & string (uuid) &
Sí. \\
\texorpdfstring{\softcode{comando}}{comando} (path) & string & Sí. \\
\texorpdfstring{\softcode{clave}}{clave} (path) & string & Sí. \\
\end{longtable}
}

{\def\LTcaptype{none} % do not increment counter
\begin{longtable}[]{@{}
  >{\raggedright\arraybackslash}p{(\linewidth - 4\tabcolsep) * \real{0.3333}}
  >{\raggedright\arraybackslash}p{(\linewidth - 4\tabcolsep) * \real{0.3333}}
  >{\raggedright\arraybackslash}p{(\linewidth - 4\tabcolsep) * \real{0.3333}}@{}}
\toprule\noalign{}
\begin{minipage}[b]{\linewidth}\raggedright
HTTP
\end{minipage} & \begin{minipage}[b]{\linewidth}\raggedright
Significado declarado
\end{minipage} & \begin{minipage}[b]{\linewidth}\raggedright
Cuerpo y media type
\end{minipage} \\
\midrule\noalign{}
\endhead
\bottomrule\noalign{}
\endlastfoot
200 & Resultado comprometido &
\texorpdfstring{\softtext{application/json:}}{application/json:}
object \\
404 & Sin resultado visible en este instante & Sin cuerpo descrito \\
401 & Sin sesión &
\texorpdfstring{\softtext{application/problem+json:}}{application/problem+json:}
Error \\
403 & Sin permiso para la sede o la acción &
\texorpdfstring{\softtext{application/problem+json:}}{application/problem+json:}
Error \\
503 & Dependencia no disponible; una respuesta perdida puede requerir
comprobar resultado con la misma clave. &
\texorpdfstring{\softtext{application/problem+json:}}{application/problem+json:}
Error \\
400 & Forma o parámetro inválido &
\texorpdfstring{\softtext{application/problem+json:}}{application/problem+json:}
Error \\
\end{longtable}
}

\section{\texorpdfstring{GET
\texorpdfstring{\softtext{/sedes/\{sedeId\}/usuarios}}{/sedes/\{sedeId\}/usuarios}}{GET }}\label{get-sedessedeidusuarios}

Consultar usuarios

\textbf{Identificador:}
\texorpdfstring{\softcode{listarUsuarios}}{listarUsuarios}.

{\def\LTcaptype{none} % do not increment counter
\begin{longtable}[]{@{}
  >{\raggedright\arraybackslash}p{(\linewidth - 4\tabcolsep) * \real{0.3333}}
  >{\raggedright\arraybackslash}p{(\linewidth - 4\tabcolsep) * \real{0.3333}}
  >{\raggedright\arraybackslash}p{(\linewidth - 4\tabcolsep) * \real{0.3333}}@{}}
\toprule\noalign{}
\begin{minipage}[b]{\linewidth}\raggedright
Parámetro y ubicación
\end{minipage} & \begin{minipage}[b]{\linewidth}\raggedright
Tipo
\end{minipage} & \begin{minipage}[b]{\linewidth}\raggedright
Requerido y propósito
\end{minipage} \\
\midrule\noalign{}
\endhead
\bottomrule\noalign{}
\endlastfoot
\texorpdfstring{\softcode{sedeId}}{sedeId} (path) & string (uuid) &
Sí. \\
\texorpdfstring{\softcode{pagina}}{pagina} (query) & integer & No. \\
\texorpdfstring{\softcode{tamano}}{tamano} (query) & integer & No. \\
\texorpdfstring{\softcode{q}}{q} (query) & string &
\texorpdfstring{\softtext{No. Búsqueda}}{No. Búsqueda} literal acotada;
\% y \_ no actúan como comodines. \\
\end{longtable}
}

{\def\LTcaptype{none} % do not increment counter
\begin{longtable}[]{@{}
  >{\raggedright\arraybackslash}p{(\linewidth - 4\tabcolsep) * \real{0.3333}}
  >{\raggedright\arraybackslash}p{(\linewidth - 4\tabcolsep) * \real{0.3333}}
  >{\raggedright\arraybackslash}p{(\linewidth - 4\tabcolsep) * \real{0.3333}}@{}}
\toprule\noalign{}
\begin{minipage}[b]{\linewidth}\raggedright
HTTP
\end{minipage} & \begin{minipage}[b]{\linewidth}\raggedright
Significado declarado
\end{minipage} & \begin{minipage}[b]{\linewidth}\raggedright
Cuerpo y media type
\end{minipage} \\
\midrule\noalign{}
\endhead
\bottomrule\noalign{}
\endlastfoot
200 & Resultado persistido del servidor &
\texorpdfstring{\softtext{application/json:}}{application/json:}
objeto \\
401 & Sin sesión &
\texorpdfstring{\softtext{application/problem+json:}}{application/problem+json:}
Error \\
403 & Sin permiso para la sede o la acción &
\texorpdfstring{\softtext{application/problem+json:}}{application/problem+json:}
Error \\
503 & Dependencia no disponible; una respuesta perdida puede requerir
comprobar resultado con la misma clave. &
\texorpdfstring{\softtext{application/problem+json:}}{application/problem+json:}
Error \\
400 & Forma o parámetro inválido &
\texorpdfstring{\softtext{application/problem+json:}}{application/problem+json:}
Error \\
\end{longtable}
}

\section{\texorpdfstring{POST
\texorpdfstring{\softtext{/sedes/\{sedeId\}/usuarios}}{/sedes/\{sedeId\}/usuarios}}{POST }}\label{post-sedessedeidusuarios}

Crear usuario por ADMIN

\textbf{Identificador:}
\texorpdfstring{\softcode{crearUsuario}}{crearUsuario}.

{\def\LTcaptype{none} % do not increment counter
\begin{longtable}[]{@{}lll@{}}
\toprule\noalign{}
Parámetro y ubicación & Tipo & Requerido y propósito \\
\midrule\noalign{}
\endhead
\bottomrule\noalign{}
\endlastfoot
\texorpdfstring{\softcode{sedeId}}{sedeId} (path) & string (uuid) &
Sí. \\
\texorpdfstring{\softcode{Idempotency-Key}}{Idempotency-Key} (header) &
string & Sí. \\
\texorpdfstring{\softcode{X-XSRF-TOKEN}}{X-XSRF-TOKEN} (header) & string
& Sí. \\
\end{longtable}
}

\textbf{Cuerpo
\texorpdfstring{\softtext{application/json:}}{application/json:}}
\texorpdfstring{\softtext{UsuarioNuevo.}}{UsuarioNuevo.}
\texorpdfstring{\softtext{Obligatorio.}}{Obligatorio.}

{\def\LTcaptype{none} % do not increment counter
\begin{longtable}[]{@{}
  >{\raggedright\arraybackslash}p{(\linewidth - 4\tabcolsep) * \real{0.3333}}
  >{\raggedright\arraybackslash}p{(\linewidth - 4\tabcolsep) * \real{0.3333}}
  >{\raggedright\arraybackslash}p{(\linewidth - 4\tabcolsep) * \real{0.3333}}@{}}
\toprule\noalign{}
\begin{minipage}[b]{\linewidth}\raggedright
HTTP
\end{minipage} & \begin{minipage}[b]{\linewidth}\raggedright
Significado declarado
\end{minipage} & \begin{minipage}[b]{\linewidth}\raggedright
Cuerpo y media type
\end{minipage} \\
\midrule\noalign{}
\endhead
\bottomrule\noalign{}
\endlastfoot
201 & Resultado persistido del servidor &
\texorpdfstring{\softtext{application/json:}}{application/json:}
\texorpdfstring{\softtext{ResultadoComando}}{ResultadoComando} \\
401 & Sin sesión &
\texorpdfstring{\softtext{application/problem+json:}}{application/problem+json:}
Error \\
403 & Sin permiso para la sede o la acción &
\texorpdfstring{\softtext{application/problem+json:}}{application/problem+json:}
Error \\
503 & Dependencia no disponible; una respuesta perdida puede requerir
comprobar resultado con la misma clave. &
\texorpdfstring{\softtext{application/problem+json:}}{application/problem+json:}
Error \\
409 & Conflicto de concurrencia,
\texorpdfstring{\softtext{version\_lock}}{version\_lock} o idempotencia
(misma clave, payload distinto) &
\texorpdfstring{\softtext{application/problem+json:}}{application/problem+json:}
Error \\
422 & Viola una regla de negocio &
\texorpdfstring{\softtext{application/problem+json:}}{application/problem+json:}
Error \\
400 & Forma o parámetro inválido &
\texorpdfstring{\softtext{application/problem+json:}}{application/problem+json:}
Error \\
\end{longtable}
}

\section{\texorpdfstring{GET
\texorpdfstring{\softtext{/sedes/\{sedeId\}/periodos}}{/sedes/\{sedeId\}/periodos}}{GET }}\label{get-sedessedeidperiodos}

Consultar periodos

\textbf{Identificador:}
\texorpdfstring{\softcode{listarPeriodos}}{listarPeriodos}.

{\def\LTcaptype{none} % do not increment counter
\begin{longtable}[]{@{}
  >{\raggedright\arraybackslash}p{(\linewidth - 4\tabcolsep) * \real{0.3333}}
  >{\raggedright\arraybackslash}p{(\linewidth - 4\tabcolsep) * \real{0.3333}}
  >{\raggedright\arraybackslash}p{(\linewidth - 4\tabcolsep) * \real{0.3333}}@{}}
\toprule\noalign{}
\begin{minipage}[b]{\linewidth}\raggedright
Parámetro y ubicación
\end{minipage} & \begin{minipage}[b]{\linewidth}\raggedright
Tipo
\end{minipage} & \begin{minipage}[b]{\linewidth}\raggedright
Requerido y propósito
\end{minipage} \\
\midrule\noalign{}
\endhead
\bottomrule\noalign{}
\endlastfoot
\texorpdfstring{\softcode{sedeId}}{sedeId} (path) & string (uuid) &
Sí. \\
\texorpdfstring{\softcode{pagina}}{pagina} (query) & integer & No. \\
\texorpdfstring{\softcode{tamano}}{tamano} (query) & integer & No. \\
\texorpdfstring{\softcode{q}}{q} (query) & string &
\texorpdfstring{\softtext{No. Búsqueda}}{No. Búsqueda} literal acotada;
\% y \_ no actúan como comodines. \\
\end{longtable}
}

{\def\LTcaptype{none} % do not increment counter
\begin{longtable}[]{@{}
  >{\raggedright\arraybackslash}p{(\linewidth - 4\tabcolsep) * \real{0.3333}}
  >{\raggedright\arraybackslash}p{(\linewidth - 4\tabcolsep) * \real{0.3333}}
  >{\raggedright\arraybackslash}p{(\linewidth - 4\tabcolsep) * \real{0.3333}}@{}}
\toprule\noalign{}
\begin{minipage}[b]{\linewidth}\raggedright
HTTP
\end{minipage} & \begin{minipage}[b]{\linewidth}\raggedright
Significado declarado
\end{minipage} & \begin{minipage}[b]{\linewidth}\raggedright
Cuerpo y media type
\end{minipage} \\
\midrule\noalign{}
\endhead
\bottomrule\noalign{}
\endlastfoot
200 & Resultado persistido del servidor &
\texorpdfstring{\softtext{application/json:}}{application/json:}
objeto \\
401 & Sin sesión &
\texorpdfstring{\softtext{application/problem+json:}}{application/problem+json:}
Error \\
403 & Sin permiso para la sede o la acción &
\texorpdfstring{\softtext{application/problem+json:}}{application/problem+json:}
Error \\
503 & Dependencia no disponible; una respuesta perdida puede requerir
comprobar resultado con la misma clave. &
\texorpdfstring{\softtext{application/problem+json:}}{application/problem+json:}
Error \\
400 & Forma o parámetro inválido &
\texorpdfstring{\softtext{application/problem+json:}}{application/problem+json:}
Error \\
\end{longtable}
}

\section{\texorpdfstring{POST
\texorpdfstring{\softtext{/sedes/\{sedeId\}/periodos}}{/sedes/\{sedeId\}/periodos}}{POST }}\label{post-sedessedeidperiodos}

Abrir mes calendario consecutivo

\textbf{Identificador:}
\texorpdfstring{\softcode{crearPeriodo}}{crearPeriodo}.

{\def\LTcaptype{none} % do not increment counter
\begin{longtable}[]{@{}lll@{}}
\toprule\noalign{}
Parámetro y ubicación & Tipo & Requerido y propósito \\
\midrule\noalign{}
\endhead
\bottomrule\noalign{}
\endlastfoot
\texorpdfstring{\softcode{sedeId}}{sedeId} (path) & string (uuid) &
Sí. \\
\texorpdfstring{\softcode{Idempotency-Key}}{Idempotency-Key} (header) &
string & Sí. \\
\texorpdfstring{\softcode{X-XSRF-TOKEN}}{X-XSRF-TOKEN} (header) & string
& Sí. \\
\end{longtable}
}

\textbf{Cuerpo
\texorpdfstring{\softtext{application/json:}}{application/json:}}
object. \texorpdfstring{\softtext{Obligatorio.}}{Obligatorio.}

{\def\LTcaptype{none} % do not increment counter
\begin{longtable}[]{@{}
  >{\raggedright\arraybackslash}p{(\linewidth - 4\tabcolsep) * \real{0.3333}}
  >{\raggedright\arraybackslash}p{(\linewidth - 4\tabcolsep) * \real{0.3333}}
  >{\raggedright\arraybackslash}p{(\linewidth - 4\tabcolsep) * \real{0.3333}}@{}}
\toprule\noalign{}
\begin{minipage}[b]{\linewidth}\raggedright
HTTP
\end{minipage} & \begin{minipage}[b]{\linewidth}\raggedright
Significado declarado
\end{minipage} & \begin{minipage}[b]{\linewidth}\raggedright
Cuerpo y media type
\end{minipage} \\
\midrule\noalign{}
\endhead
\bottomrule\noalign{}
\endlastfoot
201 & Resultado persistido del servidor &
\texorpdfstring{\softtext{application/json:}}{application/json:}
\texorpdfstring{\softtext{ResultadoComando}}{ResultadoComando} \\
401 & Sin sesión &
\texorpdfstring{\softtext{application/problem+json:}}{application/problem+json:}
Error \\
403 & Sin permiso para la sede o la acción &
\texorpdfstring{\softtext{application/problem+json:}}{application/problem+json:}
Error \\
503 & Dependencia no disponible; una respuesta perdida puede requerir
comprobar resultado con la misma clave. &
\texorpdfstring{\softtext{application/problem+json:}}{application/problem+json:}
Error \\
409 & Conflicto de concurrencia,
\texorpdfstring{\softtext{version\_lock}}{version\_lock} o idempotencia
(misma clave, payload distinto) &
\texorpdfstring{\softtext{application/problem+json:}}{application/problem+json:}
Error \\
422 & Viola una regla de negocio &
\texorpdfstring{\softtext{application/problem+json:}}{application/problem+json:}
Error \\
400 & Forma o parámetro inválido &
\texorpdfstring{\softtext{application/problem+json:}}{application/problem+json:}
Error \\
\end{longtable}
}

\section{\texorpdfstring{GET
\texorpdfstring{\softtext{/sedes/\{sedeId\}/predios}}{/sedes/\{sedeId\}/predios}}{GET }}\label{get-sedessedeidpredios}

Consultar predios

\textbf{Identificador:}
\texorpdfstring{\softcode{listarPredios}}{listarPredios}.

{\def\LTcaptype{none} % do not increment counter
\begin{longtable}[]{@{}
  >{\raggedright\arraybackslash}p{(\linewidth - 4\tabcolsep) * \real{0.3333}}
  >{\raggedright\arraybackslash}p{(\linewidth - 4\tabcolsep) * \real{0.3333}}
  >{\raggedright\arraybackslash}p{(\linewidth - 4\tabcolsep) * \real{0.3333}}@{}}
\toprule\noalign{}
\begin{minipage}[b]{\linewidth}\raggedright
Parámetro y ubicación
\end{minipage} & \begin{minipage}[b]{\linewidth}\raggedright
Tipo
\end{minipage} & \begin{minipage}[b]{\linewidth}\raggedright
Requerido y propósito
\end{minipage} \\
\midrule\noalign{}
\endhead
\bottomrule\noalign{}
\endlastfoot
\texorpdfstring{\softcode{sedeId}}{sedeId} (path) & string (uuid) &
Sí. \\
\texorpdfstring{\softcode{pagina}}{pagina} (query) & integer & No. \\
\texorpdfstring{\softcode{tamano}}{tamano} (query) & integer & No. \\
\texorpdfstring{\softcode{q}}{q} (query) & string &
\texorpdfstring{\softtext{No. Búsqueda}}{No. Búsqueda} literal acotada;
\% y \_ no actúan como comodines. \\
\end{longtable}
}

{\def\LTcaptype{none} % do not increment counter
\begin{longtable}[]{@{}
  >{\raggedright\arraybackslash}p{(\linewidth - 4\tabcolsep) * \real{0.3333}}
  >{\raggedright\arraybackslash}p{(\linewidth - 4\tabcolsep) * \real{0.3333}}
  >{\raggedright\arraybackslash}p{(\linewidth - 4\tabcolsep) * \real{0.3333}}@{}}
\toprule\noalign{}
\begin{minipage}[b]{\linewidth}\raggedright
HTTP
\end{minipage} & \begin{minipage}[b]{\linewidth}\raggedright
Significado declarado
\end{minipage} & \begin{minipage}[b]{\linewidth}\raggedright
Cuerpo y media type
\end{minipage} \\
\midrule\noalign{}
\endhead
\bottomrule\noalign{}
\endlastfoot
200 & Resultado persistido del servidor &
\texorpdfstring{\softtext{application/json:}}{application/json:}
objeto \\
401 & Sin sesión &
\texorpdfstring{\softtext{application/problem+json:}}{application/problem+json:}
Error \\
403 & Sin permiso para la sede o la acción &
\texorpdfstring{\softtext{application/problem+json:}}{application/problem+json:}
Error \\
503 & Dependencia no disponible; una respuesta perdida puede requerir
comprobar resultado con la misma clave. &
\texorpdfstring{\softtext{application/problem+json:}}{application/problem+json:}
Error \\
400 & Forma o parámetro inválido &
\texorpdfstring{\softtext{application/problem+json:}}{application/problem+json:}
Error \\
\end{longtable}
}

\section{\texorpdfstring{POST
\texorpdfstring{\softtext{/sedes/\{sedeId\}/predios}}{/sedes/\{sedeId\}/predios}}{POST }}\label{post-sedessedeidpredios}

Crear predio declarado

\textbf{Identificador:}
\texorpdfstring{\softcode{crearPredio}}{crearPredio}.

{\def\LTcaptype{none} % do not increment counter
\begin{longtable}[]{@{}lll@{}}
\toprule\noalign{}
Parámetro y ubicación & Tipo & Requerido y propósito \\
\midrule\noalign{}
\endhead
\bottomrule\noalign{}
\endlastfoot
\texorpdfstring{\softcode{sedeId}}{sedeId} (path) & string (uuid) &
Sí. \\
\texorpdfstring{\softcode{Idempotency-Key}}{Idempotency-Key} (header) &
string & Sí. \\
\texorpdfstring{\softcode{X-XSRF-TOKEN}}{X-XSRF-TOKEN} (header) & string
& Sí. \\
\end{longtable}
}

\textbf{Cuerpo
\texorpdfstring{\softtext{application/json:}}{application/json:}}
\texorpdfstring{\softtext{PredioNuevo.}}{PredioNuevo.}
\texorpdfstring{\softtext{Obligatorio.}}{Obligatorio.}

{\def\LTcaptype{none} % do not increment counter
\begin{longtable}[]{@{}
  >{\raggedright\arraybackslash}p{(\linewidth - 4\tabcolsep) * \real{0.3333}}
  >{\raggedright\arraybackslash}p{(\linewidth - 4\tabcolsep) * \real{0.3333}}
  >{\raggedright\arraybackslash}p{(\linewidth - 4\tabcolsep) * \real{0.3333}}@{}}
\toprule\noalign{}
\begin{minipage}[b]{\linewidth}\raggedright
HTTP
\end{minipage} & \begin{minipage}[b]{\linewidth}\raggedright
Significado declarado
\end{minipage} & \begin{minipage}[b]{\linewidth}\raggedright
Cuerpo y media type
\end{minipage} \\
\midrule\noalign{}
\endhead
\bottomrule\noalign{}
\endlastfoot
201 & Resultado persistido del servidor &
\texorpdfstring{\softtext{application/json:}}{application/json:}
\texorpdfstring{\softtext{ResultadoComando}}{ResultadoComando} \\
401 & Sin sesión &
\texorpdfstring{\softtext{application/problem+json:}}{application/problem+json:}
Error \\
403 & Sin permiso para la sede o la acción &
\texorpdfstring{\softtext{application/problem+json:}}{application/problem+json:}
Error \\
503 & Dependencia no disponible; una respuesta perdida puede requerir
comprobar resultado con la misma clave. &
\texorpdfstring{\softtext{application/problem+json:}}{application/problem+json:}
Error \\
409 & Conflicto de concurrencia,
\texorpdfstring{\softtext{version\_lock}}{version\_lock} o idempotencia
(misma clave, payload distinto) &
\texorpdfstring{\softtext{application/problem+json:}}{application/problem+json:}
Error \\
422 & Viola una regla de negocio &
\texorpdfstring{\softtext{application/problem+json:}}{application/problem+json:}
Error \\
400 & Forma o parámetro inválido &
\texorpdfstring{\softtext{application/problem+json:}}{application/problem+json:}
Error \\
\end{longtable}
}

\section{\texorpdfstring{PATCH
\texorpdfstring{\softtext{/sedes/\{sedeId\}/usuarios/\{usuarioId\}}}{/sedes/\{sedeId\}/usuarios/\{usuarioId\}}}{PATCH }}\label{patch-sedessedeidusuariosusuarioid}

Cambiar nombre o vigencia del acceso

\textbf{Identificador:}
\texorpdfstring{\softcode{editarUsuario}}{editarUsuario}.

{\def\LTcaptype{none} % do not increment counter
\begin{longtable}[]{@{}lll@{}}
\toprule\noalign{}
Parámetro y ubicación & Tipo & Requerido y propósito \\
\midrule\noalign{}
\endhead
\bottomrule\noalign{}
\endlastfoot
\texorpdfstring{\softcode{sedeId}}{sedeId} (path) & string (uuid) &
Sí. \\
\texorpdfstring{\softcode{usuarioId}}{usuarioId} (path) & string (uuid)
& Sí. \\
\texorpdfstring{\softcode{Idempotency-Key}}{Idempotency-Key} (header) &
string & Sí. \\
\texorpdfstring{\softcode{X-XSRF-TOKEN}}{X-XSRF-TOKEN} (header) & string
& Sí. \\
\end{longtable}
}

\textbf{Cuerpo
\texorpdfstring{\softtext{application/json:}}{application/json:}}
object. \texorpdfstring{\softtext{Obligatorio.}}{Obligatorio.}

{\def\LTcaptype{none} % do not increment counter
\begin{longtable}[]{@{}
  >{\raggedright\arraybackslash}p{(\linewidth - 4\tabcolsep) * \real{0.3333}}
  >{\raggedright\arraybackslash}p{(\linewidth - 4\tabcolsep) * \real{0.3333}}
  >{\raggedright\arraybackslash}p{(\linewidth - 4\tabcolsep) * \real{0.3333}}@{}}
\toprule\noalign{}
\begin{minipage}[b]{\linewidth}\raggedright
HTTP
\end{minipage} & \begin{minipage}[b]{\linewidth}\raggedright
Significado declarado
\end{minipage} & \begin{minipage}[b]{\linewidth}\raggedright
Cuerpo y media type
\end{minipage} \\
\midrule\noalign{}
\endhead
\bottomrule\noalign{}
\endlastfoot
200 & Resultado persistido del servidor &
\texorpdfstring{\softtext{application/json:}}{application/json:}
\texorpdfstring{\softtext{ResultadoComando}}{ResultadoComando} \\
401 & Sin sesión &
\texorpdfstring{\softtext{application/problem+json:}}{application/problem+json:}
Error \\
403 & Sin permiso para la sede o la acción &
\texorpdfstring{\softtext{application/problem+json:}}{application/problem+json:}
Error \\
503 & Dependencia no disponible; una respuesta perdida puede requerir
comprobar resultado con la misma clave. &
\texorpdfstring{\softtext{application/problem+json:}}{application/problem+json:}
Error \\
409 & Conflicto de concurrencia,
\texorpdfstring{\softtext{version\_lock}}{version\_lock} o idempotencia
(misma clave, payload distinto) &
\texorpdfstring{\softtext{application/problem+json:}}{application/problem+json:}
Error \\
422 & Viola una regla de negocio &
\texorpdfstring{\softtext{application/problem+json:}}{application/problem+json:}
Error \\
400 & Forma o parámetro inválido &
\texorpdfstring{\softtext{application/problem+json:}}{application/problem+json:}
Error \\
\end{longtable}
}

\section{\texorpdfstring{POST
\texorpdfstring{\softtext{/sedes/\{sedeId\}/usuarios/\{usuarioId\}/membresias}}{/sedes/\{sedeId\}/usuarios/\{usuarioId\}/membresias}}{POST }}\label{post-sedessedeidusuariosusuarioidmembresias}

Asignar o retirar rol en la sede autorizada

\textbf{Identificador:}
\texorpdfstring{\softcode{asignarMembresia}}{asignarMembresia}.

{\def\LTcaptype{none} % do not increment counter
\begin{longtable}[]{@{}lll@{}}
\toprule\noalign{}
Parámetro y ubicación & Tipo & Requerido y propósito \\
\midrule\noalign{}
\endhead
\bottomrule\noalign{}
\endlastfoot
\texorpdfstring{\softcode{sedeId}}{sedeId} (path) & string (uuid) &
Sí. \\
\texorpdfstring{\softcode{usuarioId}}{usuarioId} (path) & string (uuid)
& Sí. \\
\texorpdfstring{\softcode{Idempotency-Key}}{Idempotency-Key} (header) &
string & Sí. \\
\texorpdfstring{\softcode{X-XSRF-TOKEN}}{X-XSRF-TOKEN} (header) & string
& Sí. \\
\end{longtable}
}

\textbf{Cuerpo
\texorpdfstring{\softtext{application/json:}}{application/json:}}
object. \texorpdfstring{\softtext{Obligatorio.}}{Obligatorio.}

{\def\LTcaptype{none} % do not increment counter
\begin{longtable}[]{@{}
  >{\raggedright\arraybackslash}p{(\linewidth - 4\tabcolsep) * \real{0.3333}}
  >{\raggedright\arraybackslash}p{(\linewidth - 4\tabcolsep) * \real{0.3333}}
  >{\raggedright\arraybackslash}p{(\linewidth - 4\tabcolsep) * \real{0.3333}}@{}}
\toprule\noalign{}
\begin{minipage}[b]{\linewidth}\raggedright
HTTP
\end{minipage} & \begin{minipage}[b]{\linewidth}\raggedright
Significado declarado
\end{minipage} & \begin{minipage}[b]{\linewidth}\raggedright
Cuerpo y media type
\end{minipage} \\
\midrule\noalign{}
\endhead
\bottomrule\noalign{}
\endlastfoot
200 & Resultado persistido del servidor &
\texorpdfstring{\softtext{application/json:}}{application/json:}
\texorpdfstring{\softtext{ResultadoComando}}{ResultadoComando} \\
401 & Sin sesión &
\texorpdfstring{\softtext{application/problem+json:}}{application/problem+json:}
Error \\
403 & Sin permiso para la sede o la acción &
\texorpdfstring{\softtext{application/problem+json:}}{application/problem+json:}
Error \\
503 & Dependencia no disponible; una respuesta perdida puede requerir
comprobar resultado con la misma clave. &
\texorpdfstring{\softtext{application/problem+json:}}{application/problem+json:}
Error \\
409 & Conflicto de concurrencia,
\texorpdfstring{\softtext{version\_lock}}{version\_lock} o idempotencia
(misma clave, payload distinto) &
\texorpdfstring{\softtext{application/problem+json:}}{application/problem+json:}
Error \\
422 & Viola una regla de negocio &
\texorpdfstring{\softtext{application/problem+json:}}{application/problem+json:}
Error \\
400 & Forma o parámetro inválido &
\texorpdfstring{\softtext{application/problem+json:}}{application/problem+json:}
Error \\
\end{longtable}
}

\section{\texorpdfstring{POST
\texorpdfstring{\softtext{/sedes/\{sedeId\}/periodos/\{periodoId\}/reabrir}}{/sedes/\{sedeId\}/periodos/\{periodoId\}/reabrir}}{POST }}\label{post-sedessedeidperiodosperiodoidreabrir}

Reabrir último periodo sin periodos posteriores

\textbf{Identificador:}
\texorpdfstring{\softcode{reabrirPeriodo}}{reabrirPeriodo}.

{\def\LTcaptype{none} % do not increment counter
\begin{longtable}[]{@{}lll@{}}
\toprule\noalign{}
Parámetro y ubicación & Tipo & Requerido y propósito \\
\midrule\noalign{}
\endhead
\bottomrule\noalign{}
\endlastfoot
\texorpdfstring{\softcode{sedeId}}{sedeId} (path) & string (uuid) &
Sí. \\
\texorpdfstring{\softcode{periodoId}}{periodoId} (path) & string (uuid)
& Sí. \\
\texorpdfstring{\softcode{Idempotency-Key}}{Idempotency-Key} (header) &
string & Sí. \\
\texorpdfstring{\softcode{X-XSRF-TOKEN}}{X-XSRF-TOKEN} (header) & string
& Sí. \\
\end{longtable}
}

\textbf{Cuerpo
\texorpdfstring{\softtext{application/json:}}{application/json:}}
object. \texorpdfstring{\softtext{Obligatorio.}}{Obligatorio.}

{\def\LTcaptype{none} % do not increment counter
\begin{longtable}[]{@{}
  >{\raggedright\arraybackslash}p{(\linewidth - 4\tabcolsep) * \real{0.3333}}
  >{\raggedright\arraybackslash}p{(\linewidth - 4\tabcolsep) * \real{0.3333}}
  >{\raggedright\arraybackslash}p{(\linewidth - 4\tabcolsep) * \real{0.3333}}@{}}
\toprule\noalign{}
\begin{minipage}[b]{\linewidth}\raggedright
HTTP
\end{minipage} & \begin{minipage}[b]{\linewidth}\raggedright
Significado declarado
\end{minipage} & \begin{minipage}[b]{\linewidth}\raggedright
Cuerpo y media type
\end{minipage} \\
\midrule\noalign{}
\endhead
\bottomrule\noalign{}
\endlastfoot
200 & Resultado persistido del servidor &
\texorpdfstring{\softtext{application/json:}}{application/json:}
\texorpdfstring{\softtext{ResultadoComando}}{ResultadoComando} \\
401 & Sin sesión &
\texorpdfstring{\softtext{application/problem+json:}}{application/problem+json:}
Error \\
403 & Sin permiso para la sede o la acción &
\texorpdfstring{\softtext{application/problem+json:}}{application/problem+json:}
Error \\
503 & Dependencia no disponible; una respuesta perdida puede requerir
comprobar resultado con la misma clave. &
\texorpdfstring{\softtext{application/problem+json:}}{application/problem+json:}
Error \\
409 & Conflicto de concurrencia,
\texorpdfstring{\softtext{version\_lock}}{version\_lock} o idempotencia
(misma clave, payload distinto) &
\texorpdfstring{\softtext{application/problem+json:}}{application/problem+json:}
Error \\
422 & Viola una regla de negocio &
\texorpdfstring{\softtext{application/problem+json:}}{application/problem+json:}
Error \\
400 & Forma o parámetro inválido &
\texorpdfstring{\softtext{application/problem+json:}}{application/problem+json:}
Error \\
\end{longtable}
}

\section{\texorpdfstring{GET
\texorpdfstring{\softtext{/sedes/\{sedeId\}/periodos/\{periodoId\}/versiones}}{/sedes/\{sedeId\}/periodos/\{periodoId\}/versiones}}{GET }}\label{get-sedessedeidperiodosperiodoidversiones}

Consultar versiones históricas conservadas

\textbf{Identificador:}
\texorpdfstring{\softcode{consultarVersionesCierre}}{consultarVersionesCierre}.

{\def\LTcaptype{none} % do not increment counter
\begin{longtable}[]{@{}lll@{}}
\toprule\noalign{}
Parámetro y ubicación & Tipo & Requerido y propósito \\
\midrule\noalign{}
\endhead
\bottomrule\noalign{}
\endlastfoot
\texorpdfstring{\softcode{sedeId}}{sedeId} (path) & string (uuid) &
Sí. \\
\texorpdfstring{\softcode{periodoId}}{periodoId} (path) & string (uuid)
& Sí. \\
\end{longtable}
}

{\def\LTcaptype{none} % do not increment counter
\begin{longtable}[]{@{}
  >{\raggedright\arraybackslash}p{(\linewidth - 4\tabcolsep) * \real{0.3333}}
  >{\raggedright\arraybackslash}p{(\linewidth - 4\tabcolsep) * \real{0.3333}}
  >{\raggedright\arraybackslash}p{(\linewidth - 4\tabcolsep) * \real{0.3333}}@{}}
\toprule\noalign{}
\begin{minipage}[b]{\linewidth}\raggedright
HTTP
\end{minipage} & \begin{minipage}[b]{\linewidth}\raggedright
Significado declarado
\end{minipage} & \begin{minipage}[b]{\linewidth}\raggedright
Cuerpo y media type
\end{minipage} \\
\midrule\noalign{}
\endhead
\bottomrule\noalign{}
\endlastfoot
200 & Resultado persistido del servidor &
\texorpdfstring{\softtext{application/json:}}{application/json:}
object \\
401 & Sin sesión &
\texorpdfstring{\softtext{application/problem+json:}}{application/problem+json:}
Error \\
403 & Sin permiso para la sede o la acción &
\texorpdfstring{\softtext{application/problem+json:}}{application/problem+json:}
Error \\
503 & Dependencia no disponible; una respuesta perdida puede requerir
comprobar resultado con la misma clave. &
\texorpdfstring{\softtext{application/problem+json:}}{application/problem+json:}
Error \\
400 & Forma o parámetro inválido &
\texorpdfstring{\softtext{application/problem+json:}}{application/problem+json:}
Error \\
\end{longtable}
}

\section{\texorpdfstring{POST
\texorpdfstring{\softtext{/sedes/\{sedeId\}/operaciones/\{operacionId\}/acordar}}{/sedes/\{sedeId\}/operaciones/\{operacionId\}/acordar}}{POST }}\label{post-sedessedeidoperacionesoperacionidacordar}

Acordar operación conservando versiones

\textbf{Identificador:}
\texorpdfstring{\softcode{acordarOperacion}}{acordarOperacion}.

{\def\LTcaptype{none} % do not increment counter
\begin{longtable}[]{@{}lll@{}}
\toprule\noalign{}
Parámetro y ubicación & Tipo & Requerido y propósito \\
\midrule\noalign{}
\endhead
\bottomrule\noalign{}
\endlastfoot
\texorpdfstring{\softcode{sedeId}}{sedeId} (path) & string (uuid) &
Sí. \\
\texorpdfstring{\softcode{operacionId}}{operacionId} (path) & string
(uuid) & Sí. \\
\texorpdfstring{\softcode{Idempotency-Key}}{Idempotency-Key} (header) &
string & Sí. \\
\texorpdfstring{\softcode{X-XSRF-TOKEN}}{X-XSRF-TOKEN} (header) & string
& Sí. \\
\texorpdfstring{\softcode{If-Match}}{If-Match} (header) & string & Sí.
\texorpdfstring{\softtext{version\_lock}}{version\_lock} recibido en
ETag \\
\end{longtable}
}

\textbf{Cuerpo
\texorpdfstring{\softtext{application/json:}}{application/json:}}
object. \texorpdfstring{\softtext{Obligatorio.}}{Obligatorio.}

{\def\LTcaptype{none} % do not increment counter
\begin{longtable}[]{@{}
  >{\raggedright\arraybackslash}p{(\linewidth - 4\tabcolsep) * \real{0.3333}}
  >{\raggedright\arraybackslash}p{(\linewidth - 4\tabcolsep) * \real{0.3333}}
  >{\raggedright\arraybackslash}p{(\linewidth - 4\tabcolsep) * \real{0.3333}}@{}}
\toprule\noalign{}
\begin{minipage}[b]{\linewidth}\raggedright
HTTP
\end{minipage} & \begin{minipage}[b]{\linewidth}\raggedright
Significado declarado
\end{minipage} & \begin{minipage}[b]{\linewidth}\raggedright
Cuerpo y media type
\end{minipage} \\
\midrule\noalign{}
\endhead
\bottomrule\noalign{}
\endlastfoot
200 & Resultado persistido del servidor &
\texorpdfstring{\softtext{application/json:}}{application/json:}
\texorpdfstring{\softtext{ResultadoComando}}{ResultadoComando} \\
401 & Sin sesión &
\texorpdfstring{\softtext{application/problem+json:}}{application/problem+json:}
Error \\
403 & Sin permiso para la sede o la acción &
\texorpdfstring{\softtext{application/problem+json:}}{application/problem+json:}
Error \\
503 & Dependencia no disponible; una respuesta perdida puede requerir
comprobar resultado con la misma clave. &
\texorpdfstring{\softtext{application/problem+json:}}{application/problem+json:}
Error \\
409 & Conflicto de concurrencia,
\texorpdfstring{\softtext{version\_lock}}{version\_lock} o idempotencia
(misma clave, payload distinto) &
\texorpdfstring{\softtext{application/problem+json:}}{application/problem+json:}
Error \\
422 & Viola una regla de negocio &
\texorpdfstring{\softtext{application/problem+json:}}{application/problem+json:}
Error \\
412 & If-Match no corresponde a la versión revisada &
\texorpdfstring{\softtext{application/problem+json:}}{application/problem+json:}
Error \\
428 & Falta If-Match obligatorio &
\texorpdfstring{\softtext{application/problem+json:}}{application/problem+json:}
Error \\
400 & Forma o parámetro inválido &
\texorpdfstring{\softtext{application/problem+json:}}{application/problem+json:}
Error \\
\end{longtable}
}

\section{\texorpdfstring{POST
\texorpdfstring{\softtext{/sedes/\{sedeId\}/operaciones/\{operacionId\}/reabrir}}{/sedes/\{sedeId\}/operaciones/\{operacionId\}/reabrir}}{POST }}\label{post-sedessedeidoperacionesoperacionidreabrir}

Reabrir operación conservando versiones

\textbf{Identificador:}
\texorpdfstring{\softcode{reabrirOperacion}}{reabrirOperacion}.

{\def\LTcaptype{none} % do not increment counter
\begin{longtable}[]{@{}lll@{}}
\toprule\noalign{}
Parámetro y ubicación & Tipo & Requerido y propósito \\
\midrule\noalign{}
\endhead
\bottomrule\noalign{}
\endlastfoot
\texorpdfstring{\softcode{sedeId}}{sedeId} (path) & string (uuid) &
Sí. \\
\texorpdfstring{\softcode{operacionId}}{operacionId} (path) & string
(uuid) & Sí. \\
\texorpdfstring{\softcode{Idempotency-Key}}{Idempotency-Key} (header) &
string & Sí. \\
\texorpdfstring{\softcode{X-XSRF-TOKEN}}{X-XSRF-TOKEN} (header) & string
& Sí. \\
\texorpdfstring{\softcode{If-Match}}{If-Match} (header) & string & Sí.
\texorpdfstring{\softtext{version\_lock}}{version\_lock} recibido en
ETag \\
\end{longtable}
}

\textbf{Cuerpo
\texorpdfstring{\softtext{application/json:}}{application/json:}}
object. \texorpdfstring{\softtext{Obligatorio.}}{Obligatorio.}

{\def\LTcaptype{none} % do not increment counter
\begin{longtable}[]{@{}
  >{\raggedright\arraybackslash}p{(\linewidth - 4\tabcolsep) * \real{0.3333}}
  >{\raggedright\arraybackslash}p{(\linewidth - 4\tabcolsep) * \real{0.3333}}
  >{\raggedright\arraybackslash}p{(\linewidth - 4\tabcolsep) * \real{0.3333}}@{}}
\toprule\noalign{}
\begin{minipage}[b]{\linewidth}\raggedright
HTTP
\end{minipage} & \begin{minipage}[b]{\linewidth}\raggedright
Significado declarado
\end{minipage} & \begin{minipage}[b]{\linewidth}\raggedright
Cuerpo y media type
\end{minipage} \\
\midrule\noalign{}
\endhead
\bottomrule\noalign{}
\endlastfoot
200 & Resultado persistido del servidor &
\texorpdfstring{\softtext{application/json:}}{application/json:}
\texorpdfstring{\softtext{ResultadoComando}}{ResultadoComando} \\
401 & Sin sesión &
\texorpdfstring{\softtext{application/problem+json:}}{application/problem+json:}
Error \\
403 & Sin permiso para la sede o la acción &
\texorpdfstring{\softtext{application/problem+json:}}{application/problem+json:}
Error \\
503 & Dependencia no disponible; una respuesta perdida puede requerir
comprobar resultado con la misma clave. &
\texorpdfstring{\softtext{application/problem+json:}}{application/problem+json:}
Error \\
409 & Conflicto de concurrencia,
\texorpdfstring{\softtext{version\_lock}}{version\_lock} o idempotencia
(misma clave, payload distinto) &
\texorpdfstring{\softtext{application/problem+json:}}{application/problem+json:}
Error \\
422 & Viola una regla de negocio &
\texorpdfstring{\softtext{application/problem+json:}}{application/problem+json:}
Error \\
412 & If-Match no corresponde a la versión revisada &
\texorpdfstring{\softtext{application/problem+json:}}{application/problem+json:}
Error \\
428 & Falta If-Match obligatorio &
\texorpdfstring{\softtext{application/problem+json:}}{application/problem+json:}
Error \\
400 & Forma o parámetro inválido &
\texorpdfstring{\softtext{application/problem+json:}}{application/problem+json:}
Error \\
\end{longtable}
}

\section{\texorpdfstring{POST
\texorpdfstring{\softtext{/sedes/\{sedeId\}/operaciones/\{operacionId\}/cancelar}}{/sedes/\{sedeId\}/operaciones/\{operacionId\}/cancelar}}{POST }}\label{post-sedessedeidoperacionesoperacionidcancelar}

Cancelar operación conservando versiones

\textbf{Identificador:}
\texorpdfstring{\softcode{cancelarOperacion}}{cancelarOperacion}.

{\def\LTcaptype{none} % do not increment counter
\begin{longtable}[]{@{}lll@{}}
\toprule\noalign{}
Parámetro y ubicación & Tipo & Requerido y propósito \\
\midrule\noalign{}
\endhead
\bottomrule\noalign{}
\endlastfoot
\texorpdfstring{\softcode{sedeId}}{sedeId} (path) & string (uuid) &
Sí. \\
\texorpdfstring{\softcode{operacionId}}{operacionId} (path) & string
(uuid) & Sí. \\
\texorpdfstring{\softcode{Idempotency-Key}}{Idempotency-Key} (header) &
string & Sí. \\
\texorpdfstring{\softcode{X-XSRF-TOKEN}}{X-XSRF-TOKEN} (header) & string
& Sí. \\
\texorpdfstring{\softcode{If-Match}}{If-Match} (header) & string & Sí.
\texorpdfstring{\softtext{version\_lock}}{version\_lock} recibido en
ETag \\
\end{longtable}
}

\textbf{Cuerpo
\texorpdfstring{\softtext{application/json:}}{application/json:}}
object. \texorpdfstring{\softtext{Obligatorio.}}{Obligatorio.}

{\def\LTcaptype{none} % do not increment counter
\begin{longtable}[]{@{}
  >{\raggedright\arraybackslash}p{(\linewidth - 4\tabcolsep) * \real{0.3333}}
  >{\raggedright\arraybackslash}p{(\linewidth - 4\tabcolsep) * \real{0.3333}}
  >{\raggedright\arraybackslash}p{(\linewidth - 4\tabcolsep) * \real{0.3333}}@{}}
\toprule\noalign{}
\begin{minipage}[b]{\linewidth}\raggedright
HTTP
\end{minipage} & \begin{minipage}[b]{\linewidth}\raggedright
Significado declarado
\end{minipage} & \begin{minipage}[b]{\linewidth}\raggedright
Cuerpo y media type
\end{minipage} \\
\midrule\noalign{}
\endhead
\bottomrule\noalign{}
\endlastfoot
200 & Resultado persistido del servidor &
\texorpdfstring{\softtext{application/json:}}{application/json:}
\texorpdfstring{\softtext{ResultadoComando}}{ResultadoComando} \\
401 & Sin sesión &
\texorpdfstring{\softtext{application/problem+json:}}{application/problem+json:}
Error \\
403 & Sin permiso para la sede o la acción &
\texorpdfstring{\softtext{application/problem+json:}}{application/problem+json:}
Error \\
503 & Dependencia no disponible; una respuesta perdida puede requerir
comprobar resultado con la misma clave. &
\texorpdfstring{\softtext{application/problem+json:}}{application/problem+json:}
Error \\
409 & Conflicto de concurrencia,
\texorpdfstring{\softtext{version\_lock}}{version\_lock} o idempotencia
(misma clave, payload distinto) &
\texorpdfstring{\softtext{application/problem+json:}}{application/problem+json:}
Error \\
422 & Viola una regla de negocio &
\texorpdfstring{\softtext{application/problem+json:}}{application/problem+json:}
Error \\
412 & If-Match no corresponde a la versión revisada &
\texorpdfstring{\softtext{application/problem+json:}}{application/problem+json:}
Error \\
428 & Falta If-Match obligatorio &
\texorpdfstring{\softtext{application/problem+json:}}{application/problem+json:}
Error \\
400 & Forma o parámetro inválido &
\texorpdfstring{\softtext{application/problem+json:}}{application/problem+json:}
Error \\
\end{longtable}
}

\section{\texorpdfstring{POST
\texorpdfstring{\softtext{/sedes/\{sedeId\}/operaciones/\{operacionId\}/anular}}{/sedes/\{sedeId\}/operaciones/\{operacionId\}/anular}}{POST }}\label{post-sedessedeidoperacionesoperacionidanular}

Anular operación conservando versiones

\textbf{Identificador:}
\texorpdfstring{\softcode{anularOperacion}}{anularOperacion}.

{\def\LTcaptype{none} % do not increment counter
\begin{longtable}[]{@{}lll@{}}
\toprule\noalign{}
Parámetro y ubicación & Tipo & Requerido y propósito \\
\midrule\noalign{}
\endhead
\bottomrule\noalign{}
\endlastfoot
\texorpdfstring{\softcode{sedeId}}{sedeId} (path) & string (uuid) &
Sí. \\
\texorpdfstring{\softcode{operacionId}}{operacionId} (path) & string
(uuid) & Sí. \\
\texorpdfstring{\softcode{Idempotency-Key}}{Idempotency-Key} (header) &
string & Sí. \\
\texorpdfstring{\softcode{X-XSRF-TOKEN}}{X-XSRF-TOKEN} (header) & string
& Sí. \\
\texorpdfstring{\softcode{If-Match}}{If-Match} (header) & string & Sí.
\texorpdfstring{\softtext{version\_lock}}{version\_lock} recibido en
ETag \\
\end{longtable}
}

\textbf{Cuerpo
\texorpdfstring{\softtext{application/json:}}{application/json:}}
\texorpdfstring{\softtext{Reconocimiento.}}{Reconocimiento.}
\texorpdfstring{\softtext{Obligatorio.}}{Obligatorio.}

{\def\LTcaptype{none} % do not increment counter
\begin{longtable}[]{@{}
  >{\raggedright\arraybackslash}p{(\linewidth - 4\tabcolsep) * \real{0.3333}}
  >{\raggedright\arraybackslash}p{(\linewidth - 4\tabcolsep) * \real{0.3333}}
  >{\raggedright\arraybackslash}p{(\linewidth - 4\tabcolsep) * \real{0.3333}}@{}}
\toprule\noalign{}
\begin{minipage}[b]{\linewidth}\raggedright
HTTP
\end{minipage} & \begin{minipage}[b]{\linewidth}\raggedright
Significado declarado
\end{minipage} & \begin{minipage}[b]{\linewidth}\raggedright
Cuerpo y media type
\end{minipage} \\
\midrule\noalign{}
\endhead
\bottomrule\noalign{}
\endlastfoot
200 & Resultado persistido del servidor &
\texorpdfstring{\softtext{application/json:}}{application/json:}
\texorpdfstring{\softtext{ResultadoComando}}{ResultadoComando} \\
401 & Sin sesión &
\texorpdfstring{\softtext{application/problem+json:}}{application/problem+json:}
Error \\
403 & Sin permiso para la sede o la acción &
\texorpdfstring{\softtext{application/problem+json:}}{application/problem+json:}
Error \\
503 & Dependencia no disponible; una respuesta perdida puede requerir
comprobar resultado con la misma clave. &
\texorpdfstring{\softtext{application/problem+json:}}{application/problem+json:}
Error \\
409 & Conflicto de concurrencia,
\texorpdfstring{\softtext{version\_lock}}{version\_lock} o idempotencia
(misma clave, payload distinto) &
\texorpdfstring{\softtext{application/problem+json:}}{application/problem+json:}
Error \\
422 & Viola una regla de negocio &
\texorpdfstring{\softtext{application/problem+json:}}{application/problem+json:}
Error \\
412 & If-Match no corresponde a la versión revisada &
\texorpdfstring{\softtext{application/problem+json:}}{application/problem+json:}
Error \\
428 & Falta If-Match obligatorio &
\texorpdfstring{\softtext{application/problem+json:}}{application/problem+json:}
Error \\
400 & Forma o parámetro inválido &
\texorpdfstring{\softtext{application/problem+json:}}{application/problem+json:}
Error \\
\end{longtable}
}

\section{\texorpdfstring{GET
\texorpdfstring{\softtext{/sedes/\{sedeId\}/operaciones/\{operacionId\}/simulador}}{/sedes/\{sedeId\}/operaciones/\{operacionId\}/simulador}}{GET }}\label{get-sedessedeidoperacionesoperacionidsimulador}

Previsualizar sin reconocer hechos

\textbf{Identificador:}
\texorpdfstring{\softcode{simularOperacion}}{simularOperacion}.

{\def\LTcaptype{none} % do not increment counter
\begin{longtable}[]{@{}lll@{}}
\toprule\noalign{}
Parámetro y ubicación & Tipo & Requerido y propósito \\
\midrule\noalign{}
\endhead
\bottomrule\noalign{}
\endlastfoot
\texorpdfstring{\softcode{sedeId}}{sedeId} (path) & string (uuid) &
Sí. \\
\texorpdfstring{\softcode{operacionId}}{operacionId} (path) & string
(uuid) & Sí. \\
\end{longtable}
}

{\def\LTcaptype{none} % do not increment counter
\begin{longtable}[]{@{}
  >{\raggedright\arraybackslash}p{(\linewidth - 4\tabcolsep) * \real{0.3333}}
  >{\raggedright\arraybackslash}p{(\linewidth - 4\tabcolsep) * \real{0.3333}}
  >{\raggedright\arraybackslash}p{(\linewidth - 4\tabcolsep) * \real{0.3333}}@{}}
\toprule\noalign{}
\begin{minipage}[b]{\linewidth}\raggedright
HTTP
\end{minipage} & \begin{minipage}[b]{\linewidth}\raggedright
Significado declarado
\end{minipage} & \begin{minipage}[b]{\linewidth}\raggedright
Cuerpo y media type
\end{minipage} \\
\midrule\noalign{}
\endhead
\bottomrule\noalign{}
\endlastfoot
200 & Resultado persistido del servidor &
\texorpdfstring{\softtext{application/json:}}{application/json:}
object \\
401 & Sin sesión &
\texorpdfstring{\softtext{application/problem+json:}}{application/problem+json:}
Error \\
403 & Sin permiso para la sede o la acción &
\texorpdfstring{\softtext{application/problem+json:}}{application/problem+json:}
Error \\
503 & Dependencia no disponible; una respuesta perdida puede requerir
comprobar resultado con la misma clave. &
\texorpdfstring{\softtext{application/problem+json:}}{application/problem+json:}
Error \\
400 & Forma o parámetro inválido &
\texorpdfstring{\softtext{application/problem+json:}}{application/problem+json:}
Error \\
\end{longtable}
}

\section{\texorpdfstring{GET
\texorpdfstring{\softtext{/sedes/\{sedeId\}/operaciones/\{operacionId\}/expediente}}{/sedes/\{sedeId\}/operaciones/\{operacionId\}/expediente}}{GET }}\label{get-sedessedeidoperacionesoperacionidexpediente}

Consultar hechos y documentos de la operación

\textbf{Identificador:}
\texorpdfstring{\softcode{consultarExpediente}}{consultarExpediente}.

{\def\LTcaptype{none} % do not increment counter
\begin{longtable}[]{@{}lll@{}}
\toprule\noalign{}
Parámetro y ubicación & Tipo & Requerido y propósito \\
\midrule\noalign{}
\endhead
\bottomrule\noalign{}
\endlastfoot
\texorpdfstring{\softcode{sedeId}}{sedeId} (path) & string (uuid) &
Sí. \\
\texorpdfstring{\softcode{operacionId}}{operacionId} (path) & string
(uuid) & Sí. \\
\end{longtable}
}

{\def\LTcaptype{none} % do not increment counter
\begin{longtable}[]{@{}
  >{\raggedright\arraybackslash}p{(\linewidth - 4\tabcolsep) * \real{0.3333}}
  >{\raggedright\arraybackslash}p{(\linewidth - 4\tabcolsep) * \real{0.3333}}
  >{\raggedright\arraybackslash}p{(\linewidth - 4\tabcolsep) * \real{0.3333}}@{}}
\toprule\noalign{}
\begin{minipage}[b]{\linewidth}\raggedright
HTTP
\end{minipage} & \begin{minipage}[b]{\linewidth}\raggedright
Significado declarado
\end{minipage} & \begin{minipage}[b]{\linewidth}\raggedright
Cuerpo y media type
\end{minipage} \\
\midrule\noalign{}
\endhead
\bottomrule\noalign{}
\endlastfoot
200 & Resultado persistido del servidor &
\texorpdfstring{\softtext{application/json:}}{application/json:}
object \\
401 & Sin sesión &
\texorpdfstring{\softtext{application/problem+json:}}{application/problem+json:}
Error \\
403 & Sin permiso para la sede o la acción &
\texorpdfstring{\softtext{application/problem+json:}}{application/problem+json:}
Error \\
503 & Dependencia no disponible; una respuesta perdida puede requerir
comprobar resultado con la misma clave. &
\texorpdfstring{\softtext{application/problem+json:}}{application/problem+json:}
Error \\
400 & Forma o parámetro inválido &
\texorpdfstring{\softtext{application/problem+json:}}{application/problem+json:}
Error \\
\end{longtable}
}

\section{\texorpdfstring{GET
\texorpdfstring{\softtext{/sedes/\{sedeId\}/pagos/\{pagoId\}}}{/sedes/\{sedeId\}/pagos/\{pagoId\}}}{GET }}\label{get-sedessedeidpagospagoid}

Consultar dinero disponible canónico

\textbf{Identificador:}
\texorpdfstring{\softcode{consultarPago}}{consultarPago}.

{\def\LTcaptype{none} % do not increment counter
\begin{longtable}[]{@{}lll@{}}
\toprule\noalign{}
Parámetro y ubicación & Tipo & Requerido y propósito \\
\midrule\noalign{}
\endhead
\bottomrule\noalign{}
\endlastfoot
\texorpdfstring{\softcode{sedeId}}{sedeId} (path) & string (uuid) &
Sí. \\
\texorpdfstring{\softcode{pagoId}}{pagoId} (path) & string (uuid) &
Sí. \\
\end{longtable}
}

{\def\LTcaptype{none} % do not increment counter
\begin{longtable}[]{@{}
  >{\raggedright\arraybackslash}p{(\linewidth - 4\tabcolsep) * \real{0.3333}}
  >{\raggedright\arraybackslash}p{(\linewidth - 4\tabcolsep) * \real{0.3333}}
  >{\raggedright\arraybackslash}p{(\linewidth - 4\tabcolsep) * \real{0.3333}}@{}}
\toprule\noalign{}
\begin{minipage}[b]{\linewidth}\raggedright
HTTP
\end{minipage} & \begin{minipage}[b]{\linewidth}\raggedright
Significado declarado
\end{minipage} & \begin{minipage}[b]{\linewidth}\raggedright
Cuerpo y media type
\end{minipage} \\
\midrule\noalign{}
\endhead
\bottomrule\noalign{}
\endlastfoot
200 & Resultado persistido del servidor &
\texorpdfstring{\softtext{application/json:}}{application/json:} Pago \\
401 & Sin sesión &
\texorpdfstring{\softtext{application/problem+json:}}{application/problem+json:}
Error \\
403 & Sin permiso para la sede o la acción &
\texorpdfstring{\softtext{application/problem+json:}}{application/problem+json:}
Error \\
503 & Dependencia no disponible; una respuesta perdida puede requerir
comprobar resultado con la misma clave. &
\texorpdfstring{\softtext{application/problem+json:}}{application/problem+json:}
Error \\
400 & Forma o parámetro inválido &
\texorpdfstring{\softtext{application/problem+json:}}{application/problem+json:}
Error \\
\end{longtable}
}

\section{\texorpdfstring{GET
\texorpdfstring{\softtext{/sedes/\{sedeId\}/gastos/\{gastoId\}}}{/sedes/\{sedeId\}/gastos/\{gastoId\}}}{GET }}\label{get-sedessedeidgastosgastoid}

Consultar gasto persistido y procedencia

\textbf{Identificador:}
\texorpdfstring{\softcode{consultarGasto}}{consultarGasto}.

{\def\LTcaptype{none} % do not increment counter
\begin{longtable}[]{@{}lll@{}}
\toprule\noalign{}
Parámetro y ubicación & Tipo & Requerido y propósito \\
\midrule\noalign{}
\endhead
\bottomrule\noalign{}
\endlastfoot
\texorpdfstring{\softcode{sedeId}}{sedeId} (path) & string (uuid) &
Sí. \\
\texorpdfstring{\softcode{gastoId}}{gastoId} (path) & string (uuid) &
Sí. \\
\end{longtable}
}

{\def\LTcaptype{none} % do not increment counter
\begin{longtable}[]{@{}
  >{\raggedright\arraybackslash}p{(\linewidth - 4\tabcolsep) * \real{0.3333}}
  >{\raggedright\arraybackslash}p{(\linewidth - 4\tabcolsep) * \real{0.3333}}
  >{\raggedright\arraybackslash}p{(\linewidth - 4\tabcolsep) * \real{0.3333}}@{}}
\toprule\noalign{}
\begin{minipage}[b]{\linewidth}\raggedright
HTTP
\end{minipage} & \begin{minipage}[b]{\linewidth}\raggedright
Significado declarado
\end{minipage} & \begin{minipage}[b]{\linewidth}\raggedright
Cuerpo y media type
\end{minipage} \\
\midrule\noalign{}
\endhead
\bottomrule\noalign{}
\endlastfoot
200 & Resultado persistido del servidor &
\texorpdfstring{\softtext{application/json:}}{application/json:}
Gasto \\
401 & Sin sesión &
\texorpdfstring{\softtext{application/problem+json:}}{application/problem+json:}
Error \\
403 & Sin permiso para la sede o la acción &
\texorpdfstring{\softtext{application/problem+json:}}{application/problem+json:}
Error \\
503 & Dependencia no disponible; una respuesta perdida puede requerir
comprobar resultado con la misma clave. &
\texorpdfstring{\softtext{application/problem+json:}}{application/problem+json:}
Error \\
400 & Forma o parámetro inválido &
\texorpdfstring{\softtext{application/problem+json:}}{application/problem+json:}
Error \\
\end{longtable}
}

\section{\texorpdfstring{GET
\texorpdfstring{\softtext{/sedes/\{sedeId\}/documentos/\{documentoId\}}}{/sedes/\{sedeId\}/documentos/\{documentoId\}}}{GET }}\label{get-sedessedeiddocumentosdocumentoid}

Consultar metadatos y revisión exacta

\textbf{Identificador:}
\texorpdfstring{\softcode{consultarDocumento}}{consultarDocumento}.

{\def\LTcaptype{none} % do not increment counter
\begin{longtable}[]{@{}lll@{}}
\toprule\noalign{}
Parámetro y ubicación & Tipo & Requerido y propósito \\
\midrule\noalign{}
\endhead
\bottomrule\noalign{}
\endlastfoot
\texorpdfstring{\softcode{sedeId}}{sedeId} (path) & string (uuid) &
Sí. \\
\texorpdfstring{\softcode{documentoId}}{documentoId} (path) & string
(uuid) & Sí. \\
\end{longtable}
}

{\def\LTcaptype{none} % do not increment counter
\begin{longtable}[]{@{}
  >{\raggedright\arraybackslash}p{(\linewidth - 4\tabcolsep) * \real{0.3333}}
  >{\raggedright\arraybackslash}p{(\linewidth - 4\tabcolsep) * \real{0.3333}}
  >{\raggedright\arraybackslash}p{(\linewidth - 4\tabcolsep) * \real{0.3333}}@{}}
\toprule\noalign{}
\begin{minipage}[b]{\linewidth}\raggedright
HTTP
\end{minipage} & \begin{minipage}[b]{\linewidth}\raggedright
Significado declarado
\end{minipage} & \begin{minipage}[b]{\linewidth}\raggedright
Cuerpo y media type
\end{minipage} \\
\midrule\noalign{}
\endhead
\bottomrule\noalign{}
\endlastfoot
200 & Resultado persistido del servidor &
\texorpdfstring{\softtext{application/json:}}{application/json:}
Documento \\
401 & Sin sesión &
\texorpdfstring{\softtext{application/problem+json:}}{application/problem+json:}
Error \\
403 & Sin permiso para la sede o la acción &
\texorpdfstring{\softtext{application/problem+json:}}{application/problem+json:}
Error \\
503 & Dependencia no disponible; una respuesta perdida puede requerir
comprobar resultado con la misma clave. &
\texorpdfstring{\softtext{application/problem+json:}}{application/problem+json:}
Error \\
400 & Forma o parámetro inválido &
\texorpdfstring{\softtext{application/problem+json:}}{application/problem+json:}
Error \\
\end{longtable}
}

\section{\texorpdfstring{POST
\texorpdfstring{\softtext{/sedes/\{sedeId\}/documentos/\{documentoId\}/revisiones}}{/sedes/\{sedeId\}/documentos/\{documentoId\}/revisiones}}{POST }}\label{post-sedessedeiddocumentosdocumentoidrevisiones}

Registrar revisión administrativa de esta versión

\textbf{Identificador:}
\texorpdfstring{\softcode{revisarDocumento}}{revisarDocumento}.

{\def\LTcaptype{none} % do not increment counter
\begin{longtable}[]{@{}lll@{}}
\toprule\noalign{}
Parámetro y ubicación & Tipo & Requerido y propósito \\
\midrule\noalign{}
\endhead
\bottomrule\noalign{}
\endlastfoot
\texorpdfstring{\softcode{sedeId}}{sedeId} (path) & string (uuid) &
Sí. \\
\texorpdfstring{\softcode{documentoId}}{documentoId} (path) & string
(uuid) & Sí. \\
\texorpdfstring{\softcode{Idempotency-Key}}{Idempotency-Key} (header) &
string & Sí. \\
\texorpdfstring{\softcode{X-XSRF-TOKEN}}{X-XSRF-TOKEN} (header) & string
& Sí. \\
\end{longtable}
}

\textbf{Cuerpo
\texorpdfstring{\softtext{application/json:}}{application/json:}}
object. \texorpdfstring{\softtext{Obligatorio.}}{Obligatorio.}

{\def\LTcaptype{none} % do not increment counter
\begin{longtable}[]{@{}
  >{\raggedright\arraybackslash}p{(\linewidth - 4\tabcolsep) * \real{0.3333}}
  >{\raggedright\arraybackslash}p{(\linewidth - 4\tabcolsep) * \real{0.3333}}
  >{\raggedright\arraybackslash}p{(\linewidth - 4\tabcolsep) * \real{0.3333}}@{}}
\toprule\noalign{}
\begin{minipage}[b]{\linewidth}\raggedright
HTTP
\end{minipage} & \begin{minipage}[b]{\linewidth}\raggedright
Significado declarado
\end{minipage} & \begin{minipage}[b]{\linewidth}\raggedright
Cuerpo y media type
\end{minipage} \\
\midrule\noalign{}
\endhead
\bottomrule\noalign{}
\endlastfoot
200 & Resultado persistido del servidor &
\texorpdfstring{\softtext{application/json:}}{application/json:}
\texorpdfstring{\softtext{ResultadoComando}}{ResultadoComando} \\
401 & Sin sesión &
\texorpdfstring{\softtext{application/problem+json:}}{application/problem+json:}
Error \\
403 & Sin permiso para la sede o la acción &
\texorpdfstring{\softtext{application/problem+json:}}{application/problem+json:}
Error \\
503 & Dependencia no disponible; una respuesta perdida puede requerir
comprobar resultado con la misma clave. &
\texorpdfstring{\softtext{application/problem+json:}}{application/problem+json:}
Error \\
409 & Conflicto de concurrencia,
\texorpdfstring{\softtext{version\_lock}}{version\_lock} o idempotencia
(misma clave, payload distinto) &
\texorpdfstring{\softtext{application/problem+json:}}{application/problem+json:}
Error \\
422 & Viola una regla de negocio &
\texorpdfstring{\softtext{application/problem+json:}}{application/problem+json:}
Error \\
400 & Forma o parámetro inválido &
\texorpdfstring{\softtext{application/problem+json:}}{application/problem+json:}
Error \\
\end{longtable}
}

\section{\texorpdfstring{GET
\texorpdfstring{\softtext{/sedes/\{sedeId\}/pesajes/comparacion}}{/sedes/\{sedeId\}/pesajes/comparacion}}{GET }}\label{get-sedessedeidpesajescomparacion}

Comparar exclusivamente conjuntos compatibles

\textbf{Identificador:}
\texorpdfstring{\softcode{compararPesajes}}{compararPesajes}.

{\def\LTcaptype{none} % do not increment counter
\begin{longtable}[]{@{}lll@{}}
\toprule\noalign{}
Parámetro y ubicación & Tipo & Requerido y propósito \\
\midrule\noalign{}
\endhead
\bottomrule\noalign{}
\endlastfoot
\texorpdfstring{\softcode{sedeId}}{sedeId} (path) & string (uuid) &
Sí. \\
\texorpdfstring{\softcode{inicial\_id}}{inicial\_id} (query) & string
(uuid) & Sí. \\
\texorpdfstring{\softcode{final\_id}}{final\_id} (query) & string (uuid)
& Sí. \\
\end{longtable}
}

{\def\LTcaptype{none} % do not increment counter
\begin{longtable}[]{@{}
  >{\raggedright\arraybackslash}p{(\linewidth - 4\tabcolsep) * \real{0.3333}}
  >{\raggedright\arraybackslash}p{(\linewidth - 4\tabcolsep) * \real{0.3333}}
  >{\raggedright\arraybackslash}p{(\linewidth - 4\tabcolsep) * \real{0.3333}}@{}}
\toprule\noalign{}
\begin{minipage}[b]{\linewidth}\raggedright
HTTP
\end{minipage} & \begin{minipage}[b]{\linewidth}\raggedright
Significado declarado
\end{minipage} & \begin{minipage}[b]{\linewidth}\raggedright
Cuerpo y media type
\end{minipage} \\
\midrule\noalign{}
\endhead
\bottomrule\noalign{}
\endlastfoot
200 & Resultado persistido del servidor &
\texorpdfstring{\softtext{application/json:}}{application/json:}
object \\
401 & Sin sesión &
\texorpdfstring{\softtext{application/problem+json:}}{application/problem+json:}
Error \\
403 & Sin permiso para la sede o la acción &
\texorpdfstring{\softtext{application/problem+json:}}{application/problem+json:}
Error \\
503 & Dependencia no disponible; una respuesta perdida puede requerir
comprobar resultado con la misma clave. &
\texorpdfstring{\softtext{application/problem+json:}}{application/problem+json:}
Error \\
400 & Forma o parámetro inválido &
\texorpdfstring{\softtext{application/problem+json:}}{application/problem+json:}
Error \\
\end{longtable}
}

\section{\texorpdfstring{GET
\texorpdfstring{\softtext{/sedes/\{sedeId\}/animales/\{animalId\}/qr}}{/sedes/\{sedeId\}/animales/\{animalId\}/qr}}{GET }}\label{get-sedessedeidanimalesanimalidqr}

Exportar recurso autenticado

\textbf{Identificador:} \texorpdfstring{\softcode{qrAnimal}}{qrAnimal}.

{\def\LTcaptype{none} % do not increment counter
\begin{longtable}[]{@{}lll@{}}
\toprule\noalign{}
Parámetro y ubicación & Tipo & Requerido y propósito \\
\midrule\noalign{}
\endhead
\bottomrule\noalign{}
\endlastfoot
\texorpdfstring{\softcode{sedeId}}{sedeId} (path) & string (uuid) &
Sí. \\
\texorpdfstring{\softcode{animalId}}{animalId} (path) & string (uuid) &
Sí. \\
\end{longtable}
}

{\def\LTcaptype{none} % do not increment counter
\begin{longtable}[]{@{}
  >{\raggedright\arraybackslash}p{(\linewidth - 4\tabcolsep) * \real{0.3333}}
  >{\raggedright\arraybackslash}p{(\linewidth - 4\tabcolsep) * \real{0.3333}}
  >{\raggedright\arraybackslash}p{(\linewidth - 4\tabcolsep) * \real{0.3333}}@{}}
\toprule\noalign{}
\begin{minipage}[b]{\linewidth}\raggedright
HTTP
\end{minipage} & \begin{minipage}[b]{\linewidth}\raggedright
Significado declarado
\end{minipage} & \begin{minipage}[b]{\linewidth}\raggedright
Cuerpo y media type
\end{minipage} \\
\midrule\noalign{}
\endhead
\bottomrule\noalign{}
\endlastfoot
200 & Resultado persistido del servidor &
\texorpdfstring{\softtext{image/svg+xml:}}{image/svg+xml:} string \\
401 & Sin sesión &
\texorpdfstring{\softtext{application/problem+json:}}{application/problem+json:}
Error \\
403 & Sin permiso para la sede o la acción &
\texorpdfstring{\softtext{application/problem+json:}}{application/problem+json:}
Error \\
503 & Dependencia no disponible; una respuesta perdida puede requerir
comprobar resultado con la misma clave. &
\texorpdfstring{\softtext{application/problem+json:}}{application/problem+json:}
Error \\
400 & Forma o parámetro inválido &
\texorpdfstring{\softtext{application/problem+json:}}{application/problem+json:}
Error \\
\end{longtable}
}

\section{\texorpdfstring{GET
\texorpdfstring{\softtext{/sedes/\{sedeId\}/reportes/\{tipo\}/csv}}{/sedes/\{sedeId\}/reportes/\{tipo\}/csv}}{GET }}\label{get-sedessedeidreportestipocsv}

Exportar recurso autenticado

\textbf{Identificador:}
\texorpdfstring{\softcode{exportarCsv}}{exportarCsv}.

{\def\LTcaptype{none} % do not increment counter
\begin{longtable}[]{@{}lll@{}}
\toprule\noalign{}
Parámetro y ubicación & Tipo & Requerido y propósito \\
\midrule\noalign{}
\endhead
\bottomrule\noalign{}
\endlastfoot
\texorpdfstring{\softcode{sedeId}}{sedeId} (path) & string (uuid) &
Sí. \\
\texorpdfstring{\softcode{tipo}}{tipo} (path) & string & Sí. \\
\end{longtable}
}

{\def\LTcaptype{none} % do not increment counter
\begin{longtable}[]{@{}
  >{\raggedright\arraybackslash}p{(\linewidth - 4\tabcolsep) * \real{0.3333}}
  >{\raggedright\arraybackslash}p{(\linewidth - 4\tabcolsep) * \real{0.3333}}
  >{\raggedright\arraybackslash}p{(\linewidth - 4\tabcolsep) * \real{0.3333}}@{}}
\toprule\noalign{}
\begin{minipage}[b]{\linewidth}\raggedright
HTTP
\end{minipage} & \begin{minipage}[b]{\linewidth}\raggedright
Significado declarado
\end{minipage} & \begin{minipage}[b]{\linewidth}\raggedright
Cuerpo y media type
\end{minipage} \\
\midrule\noalign{}
\endhead
\bottomrule\noalign{}
\endlastfoot
200 & Resultado persistido del servidor & text/csv: string \\
401 & Sin sesión &
\texorpdfstring{\softtext{application/problem+json:}}{application/problem+json:}
Error \\
403 & Sin permiso para la sede o la acción &
\texorpdfstring{\softtext{application/problem+json:}}{application/problem+json:}
Error \\
503 & Dependencia no disponible; una respuesta perdida puede requerir
comprobar resultado con la misma clave. &
\texorpdfstring{\softtext{application/problem+json:}}{application/problem+json:}
Error \\
400 & Forma o parámetro inválido &
\texorpdfstring{\softtext{application/problem+json:}}{application/problem+json:}
Error \\
\end{longtable}
}

\section{\texorpdfstring{GET
\texorpdfstring{\softtext{/sedes/\{sedeId\}/sincronizacion}}{/sedes/\{sedeId\}/sincronizacion}}{GET }}\label{get-sedessedeidsincronizacion}

Consultar incidencias del propietario

\textbf{Identificador:}
\texorpdfstring{\softcode{consultarSincronizacion}}{consultarSincronizacion}.

{\def\LTcaptype{none} % do not increment counter
\begin{longtable}[]{@{}lll@{}}
\toprule\noalign{}
Parámetro y ubicación & Tipo & Requerido y propósito \\
\midrule\noalign{}
\endhead
\bottomrule\noalign{}
\endlastfoot
\texorpdfstring{\softcode{sedeId}}{sedeId} (path) & string (uuid) &
Sí. \\
\end{longtable}
}

{\def\LTcaptype{none} % do not increment counter
\begin{longtable}[]{@{}
  >{\raggedright\arraybackslash}p{(\linewidth - 4\tabcolsep) * \real{0.3333}}
  >{\raggedright\arraybackslash}p{(\linewidth - 4\tabcolsep) * \real{0.3333}}
  >{\raggedright\arraybackslash}p{(\linewidth - 4\tabcolsep) * \real{0.3333}}@{}}
\toprule\noalign{}
\begin{minipage}[b]{\linewidth}\raggedright
HTTP
\end{minipage} & \begin{minipage}[b]{\linewidth}\raggedright
Significado declarado
\end{minipage} & \begin{minipage}[b]{\linewidth}\raggedright
Cuerpo y media type
\end{minipage} \\
\midrule\noalign{}
\endhead
\bottomrule\noalign{}
\endlastfoot
200 & Resultado persistido del servidor &
\texorpdfstring{\softtext{application/json:}}{application/json:}
object \\
401 & Sin sesión &
\texorpdfstring{\softtext{application/problem+json:}}{application/problem+json:}
Error \\
403 & Sin permiso para la sede o la acción &
\texorpdfstring{\softtext{application/problem+json:}}{application/problem+json:}
Error \\
503 & Dependencia no disponible; una respuesta perdida puede requerir
comprobar resultado con la misma clave. &
\texorpdfstring{\softtext{application/problem+json:}}{application/problem+json:}
Error \\
400 & Forma o parámetro inválido &
\texorpdfstring{\softtext{application/problem+json:}}{application/problem+json:}
Error \\
\end{longtable}
}

\section{\texorpdfstring{POST
\texorpdfstring{\softtext{/sedes/\{sedeId\}/integraciones/proveedor/tarifas/aceptaciones}}{/sedes/\{sedeId\}/integraciones/proveedor/tarifas/aceptaciones}}{POST }}\label{post-sedessedeidintegracionesproveedortarifasaceptaciones}

Aceptar captura local sin consultar de nuevo el proveedor

\textbf{Identificador:}
\texorpdfstring{\softcode{aceptarTarifa}}{aceptarTarifa}.

{\def\LTcaptype{none} % do not increment counter
\begin{longtable}[]{@{}lll@{}}
\toprule\noalign{}
Parámetro y ubicación & Tipo & Requerido y propósito \\
\midrule\noalign{}
\endhead
\bottomrule\noalign{}
\endlastfoot
\texorpdfstring{\softcode{sedeId}}{sedeId} (path) & string (uuid) &
Sí. \\
\texorpdfstring{\softcode{Idempotency-Key}}{Idempotency-Key} (header) &
string & Sí. \\
\texorpdfstring{\softcode{X-XSRF-TOKEN}}{X-XSRF-TOKEN} (header) & string
& Sí. \\
\end{longtable}
}

\textbf{Cuerpo
\texorpdfstring{\softtext{application/json:}}{application/json:}}
object. \texorpdfstring{\softtext{Obligatorio.}}{Obligatorio.}

{\def\LTcaptype{none} % do not increment counter
\begin{longtable}[]{@{}
  >{\raggedright\arraybackslash}p{(\linewidth - 4\tabcolsep) * \real{0.3333}}
  >{\raggedright\arraybackslash}p{(\linewidth - 4\tabcolsep) * \real{0.3333}}
  >{\raggedright\arraybackslash}p{(\linewidth - 4\tabcolsep) * \real{0.3333}}@{}}
\toprule\noalign{}
\begin{minipage}[b]{\linewidth}\raggedright
HTTP
\end{minipage} & \begin{minipage}[b]{\linewidth}\raggedright
Significado declarado
\end{minipage} & \begin{minipage}[b]{\linewidth}\raggedright
Cuerpo y media type
\end{minipage} \\
\midrule\noalign{}
\endhead
\bottomrule\noalign{}
\endlastfoot
201 & Captura aceptada para registrar gasto en el mismo actor/sede &
\texorpdfstring{\softtext{application/json:}}{application/json:}
\texorpdfstring{\softtext{ResultadoComando}}{ResultadoComando} \\
401 & Sin sesión &
\texorpdfstring{\softtext{application/problem+json:}}{application/problem+json:}
Error \\
403 & Sin permiso para la sede o la acción &
\texorpdfstring{\softtext{application/problem+json:}}{application/problem+json:}
Error \\
409 & Conflicto de concurrencia,
\texorpdfstring{\softtext{version\_lock}}{version\_lock} o idempotencia
(misma clave, payload distinto) &
\texorpdfstring{\softtext{application/problem+json:}}{application/problem+json:}
Error \\
422 & Viola una regla de negocio &
\texorpdfstring{\softtext{application/problem+json:}}{application/problem+json:}
Error \\
503 & Dependencia no disponible; una respuesta perdida puede requerir
comprobar resultado con la misma clave. &
\texorpdfstring{\softtext{application/problem+json:}}{application/problem+json:}
Error \\
\end{longtable}
}

\section{\texorpdfstring{GET
\texorpdfstring{\softtext{/sedes/\{sedeId\}/integraciones/diagnostico}}{/sedes/\{sedeId\}/integraciones/diagnostico}}{GET }}\label{get-sedessedeidintegracionesdiagnostico}

Diagnóstico técnico acotado de ADMIN

\textbf{Identificador:}
\texorpdfstring{\softcode{consultarDiagnostico}}{consultarDiagnostico}.
Cache-Control:no-store; sin URLs, credenciales ni SQL. Duraciones
acumuladas de proceso, sin histogramas ni
\texorpdfstring{\softtext{certificaciónSLO.}}{certificaciónSLO.}

{\def\LTcaptype{none} % do not increment counter
\begin{longtable}[]{@{}lll@{}}
\toprule\noalign{}
Parámetro y ubicación & Tipo & Requerido y propósito \\
\midrule\noalign{}
\endhead
\bottomrule\noalign{}
\endlastfoot
\texorpdfstring{\softcode{sedeId}}{sedeId} (path) & string (uuid) &
Sí. \\
\end{longtable}
}

{\def\LTcaptype{none} % do not increment counter
\begin{longtable}[]{@{}
  >{\raggedright\arraybackslash}p{(\linewidth - 4\tabcolsep) * \real{0.3333}}
  >{\raggedright\arraybackslash}p{(\linewidth - 4\tabcolsep) * \real{0.3333}}
  >{\raggedright\arraybackslash}p{(\linewidth - 4\tabcolsep) * \real{0.3333}}@{}}
\toprule\noalign{}
\begin{minipage}[b]{\linewidth}\raggedright
HTTP
\end{minipage} & \begin{minipage}[b]{\linewidth}\raggedright
Significado declarado
\end{minipage} & \begin{minipage}[b]{\linewidth}\raggedright
Cuerpo y media type
\end{minipage} \\
\midrule\noalign{}
\endhead
\bottomrule\noalign{}
\endlastfoot
200 & Contadores del proceso y pools reales, sólo ADMIN; reinician al
reiniciar API. &
\texorpdfstring{\softtext{application/json:}}{application/json:}
\texorpdfstring{\softtext{DiagnosticoProceso}}{DiagnosticoProceso} \\
401 & Sin sesión &
\texorpdfstring{\softtext{application/problem+json:}}{application/problem+json:}
Error \\
403 & Sin permiso para la sede o la acción &
\texorpdfstring{\softtext{application/problem+json:}}{application/problem+json:}
Error \\
503 & Dependencia no disponible; una respuesta perdida puede requerir
comprobar resultado con la misma clave. &
\texorpdfstring{\softtext{application/problem+json:}}{application/problem+json:}
Error \\
\end{longtable}
}

\section{GET /salud}\label{get-salud}

Salud del proceso y conectividad de componentes

\textbf{Identificador:}
\texorpdfstring{\softcode{consultarSaludProceso}}{consultarSaludProceso}.

{\def\LTcaptype{none} % do not increment counter
\begin{longtable}[]{@{}
  >{\raggedright\arraybackslash}p{(\linewidth - 4\tabcolsep) * \real{0.3333}}
  >{\raggedright\arraybackslash}p{(\linewidth - 4\tabcolsep) * \real{0.3333}}
  >{\raggedright\arraybackslash}p{(\linewidth - 4\tabcolsep) * \real{0.3333}}@{}}
\toprule\noalign{}
\begin{minipage}[b]{\linewidth}\raggedright
HTTP
\end{minipage} & \begin{minipage}[b]{\linewidth}\raggedright
Significado declarado
\end{minipage} & \begin{minipage}[b]{\linewidth}\raggedright
Cuerpo y media type
\end{minipage} \\
\midrule\noalign{}
\endhead
\bottomrule\noalign{}
\endlastfoot
200 & El proceso sigue activo aunque un nodo falle; no sustituye
disponibilidad \texorpdfstring{\softtext{autenticada.}}{autenticada.} &
\texorpdfstring{\softtext{application/json:}}{application/json:}
object \\
\end{longtable}
}

\chapter{Diccionario completo de DTO y tipos
API}\label{diccionario-completo-de-dto-y-tipos-api}

Un DTO es el objeto que cruza HTTP. Un
\texorpdfstring{\softcode{\$ref}}{\$ref} reutiliza la definición de un
tipo; no es una FK. \texorpdfstring{\softcode{required}}{required}
indica presencia obligatoria,
\texorpdfstring{\softcode{nullable}}{nullable} o un tipo que incluye
null indica ausencia admitida. Un número transportado como string
preserva su precisión y escala; el formato visual no cambia el valor
canónico.

\section{Decimal2}\label{decimal2}

Definición contractual reutilizable de tipo Decimal2.

\textbf{Tipo:} string.

\textbf{Patrón:}
\texorpdfstring{\softcode{\textasciicircum{}-?\textbackslash{}d\{1,16\}\textbackslash{}.\textbackslash{}d\{2\}\$}}{\textasciicircum{}-?\textbackslash{}d\{1,16\}\textbackslash{}.\textbackslash{}d\{2\}\$}.

\section{Decimal3}\label{decimal3}

Definición contractual reutilizable de tipo Decimal3.

\textbf{Tipo:} string.

\textbf{Patrón:}
\texorpdfstring{\softcode{\textasciicircum{}\textbackslash{}d\{1,11\}\textbackslash{}.\textbackslash{}d\{3\}\$}}{\textasciicircum{}\textbackslash{}d\{1,11\}\textbackslash{}.\textbackslash{}d\{3\}\$}.

\section{Decimal4}\label{decimal4}

Definición contractual reutilizable de tipo Decimal4.

\textbf{Tipo:} string.

\textbf{Patrón:}
\texorpdfstring{\softcode{\textasciicircum{}\textbackslash{}d\{1,10\}\textbackslash{}.\textbackslash{}d\{4\}\$}}{\textasciicircum{}\textbackslash{}d\{1,10\}\textbackslash{}.\textbackslash{}d\{4\}\$}.

\section{Dinero}\label{dinero}

Definición contractual reutilizable de tipo Dinero.

\textbf{Tipo:} object.

{\def\LTcaptype{none} % do not increment counter
\begin{longtable}[]{@{}lll@{}}
\toprule\noalign{}
Campo & Tipo y presencia & Restricción o significado \\
\midrule\noalign{}
\endhead
\bottomrule\noalign{}
\endlastfoot
\texorpdfstring{\softcode{importe}}{importe} & Decimal2; requerido & Ver
definición Decimal2 \\
\texorpdfstring{\softcode{moneda}}{moneda} & string; requerido &
enum={[}`MXN'{]} \\
\end{longtable}
}

\section{Peso}\label{peso}

Definición contractual reutilizable de tipo Peso.

\textbf{Tipo:} object.

{\def\LTcaptype{none} % do not increment counter
\begin{longtable}[]{@{}lll@{}}
\toprule\noalign{}
Campo & Tipo y presencia & Restricción o significado \\
\midrule\noalign{}
\endhead
\bottomrule\noalign{}
\endlastfoot
\texorpdfstring{\softcode{kg}}{kg} & Decimal3; requerido & Ver
definición Decimal3 \\
\end{longtable}
}

\section{Error}\label{error}

Definición contractual reutilizable de tipo Error.

\textbf{Tipo:} object.

{\def\LTcaptype{none} % do not increment counter
\begin{longtable}[]{@{}lll@{}}
\toprule\noalign{}
Campo & Tipo y presencia & Restricción o significado \\
\midrule\noalign{}
\endhead
\bottomrule\noalign{}
\endlastfoot
\texorpdfstring{\softcode{codigo}}{codigo} & string; requerido & \\
\texorpdfstring{\softcode{mensaje}}{mensaje} & string; requerido & \\
\texorpdfstring{\softcode{regla\_id}}{regla\_id} & string; opcional & \\
\texorpdfstring{\softcode{correlation\_id}}{correlation\_id} & string
(uuid); requerido & \\
\texorpdfstring{\softcode{campos}}{campos} & colección de object;
opcional & \\
\end{longtable}
}

\section{Pagina}\label{pagina}

Definición contractual reutilizable de tipo Pagina.

\textbf{Tipo:} object.

{\def\LTcaptype{none} % do not increment counter
\begin{longtable}[]{@{}lll@{}}
\toprule\noalign{}
Campo & Tipo y presencia & Restricción o significado \\
\midrule\noalign{}
\endhead
\bottomrule\noalign{}
\endlastfoot
\texorpdfstring{\softcode{pagina}}{pagina} & integer; requerido & \\
\texorpdfstring{\softcode{tamano}}{tamano} & integer; requerido & \\
\texorpdfstring{\softcode{total}}{total} & integer; requerido & \\
\end{longtable}
}

\section{\texorpdfstring{\texorpdfstring{\softtext{SesionActual}}{SesionActual}}{}}\label{sesionactual}

Definición contractual reutilizable de tipo
\texorpdfstring{\softtext{SesionActual.}}{SesionActual.}

\textbf{Tipo:} object.

{\def\LTcaptype{none} % do not increment counter
\begin{longtable}[]{@{}lll@{}}
\toprule\noalign{}
Campo & Tipo y presencia & Restricción o significado \\
\midrule\noalign{}
\endhead
\bottomrule\noalign{}
\endlastfoot
\texorpdfstring{\softcode{usuario\_id}}{usuario\_id} & string (uuid);
requerido & \\
\texorpdfstring{\softcode{nombre}}{nombre} & string; requerido & \\
\texorpdfstring{\softcode{sedes}}{sedes} & colección de object;
requerido & \\
\end{longtable}
}

\section{\texorpdfstring{\texorpdfstring{\softtext{ContraparteNueva}}{ContraparteNueva}}{}}\label{contrapartenueva}

Definición contractual reutilizable de tipo
\texorpdfstring{\softtext{ContraparteNueva.}}{ContraparteNueva.}

\textbf{Tipo:} object.

{\def\LTcaptype{none} % do not increment counter
\begin{longtable}[]{@{}lll@{}}
\toprule\noalign{}
Campo & Tipo y presencia & Restricción o significado \\
\midrule\noalign{}
\endhead
\bottomrule\noalign{}
\endlastfoot
\texorpdfstring{\softcode{tipo\_persona}}{tipo\_persona} & string;
requerido & enum={[}`FISICA', `MORAL'{]} \\
\texorpdfstring{\softcode{nombre}}{nombre} & string; requerido &
\texorpdfstring{\softtext{maxLength=200}}{maxLength=200} \\
\texorpdfstring{\softcode{rfc}}{rfc} & string; opcional &
\texorpdfstring{\softtext{maxLength=13}}{maxLength=13} \\
\texorpdfstring{\softcode{telefono}}{telefono} & string; opcional &
\texorpdfstring{\softtext{maxLength=20}}{maxLength=20} \\
\end{longtable}
}

\section{Contraparte}\label{contraparte}

Definición contractual reutilizable de tipo
\texorpdfstring{\softtext{Contraparte.}}{Contraparte.}

\textbf{Tipo:} objeto.

\textbf{allOf:}
\texorpdfstring{\softtext{ContraparteNueva,}}{ContraparteNueva,} object.

\section{Animal}\label{animal}

Definición contractual reutilizable de tipo Animal.

\textbf{Tipo:} object.

{\def\LTcaptype{none} % do not increment counter
\begin{longtable}[]{@{}
  >{\raggedright\arraybackslash}p{(\linewidth - 4\tabcolsep) * \real{0.3333}}
  >{\raggedright\arraybackslash}p{(\linewidth - 4\tabcolsep) * \real{0.3333}}
  >{\raggedright\arraybackslash}p{(\linewidth - 4\tabcolsep) * \real{0.3333}}@{}}
\toprule\noalign{}
\begin{minipage}[b]{\linewidth}\raggedright
Campo
\end{minipage} & \begin{minipage}[b]{\linewidth}\raggedright
Tipo y presencia
\end{minipage} & \begin{minipage}[b]{\linewidth}\raggedright
Restricción o significado
\end{minipage} \\
\midrule\noalign{}
\endhead
\bottomrule\noalign{}
\endlastfoot
\texorpdfstring{\softcode{id}}{id} & string (uuid); opcional & \\
\texorpdfstring{\softcode{sede\_operadora\_id}}{sede\_operadora\_id} &
string (uuid); opcional & \\
\texorpdfstring{\softcode{situacion}}{situacion} & string; opcional &
enum={[}`VIGENTE',
`\texorpdfstring{\softtext{IDENTIDAD\_EN\_REVISION}}{IDENTIDAD\_EN\_REVISION}',
`MUERTO'{]} \\
\texorpdfstring{\softcode{estado\_convergencia}}{estado\_convergencia} &
string; opcional & enum={[}`VISIBLE',
`\texorpdfstring{\softtext{PENDIENTE\_CONVERGENCIA}}{PENDIENTE\_CONVERGENCIA}'{]} \\
\texorpdfstring{\softcode{identificaciones}}{identificaciones} &
colección de object; opcional & \\
\end{longtable}
}

\section{\texorpdfstring{\texorpdfstring{\softtext{FichaAnimal}}{FichaAnimal}}{}}\label{fichaanimal}

Definición contractual reutilizable de tipo
\texorpdfstring{\softtext{FichaAnimal.}}{FichaAnimal.}

\textbf{Tipo:} objeto.

\textbf{allOf:} Animal, object.

\section{\texorpdfstring{\texorpdfstring{\softtext{EstadoOperacion}}{EstadoOperacion}}{}}\label{estadooperacion}

Definición contractual reutilizable de tipo
\texorpdfstring{\softtext{EstadoOperacion.}}{EstadoOperacion.}

\textbf{Tipo:} string.

\textbf{Valores permitidos:}
\texorpdfstring{\softcode{BORRADOR}}{BORRADOR},
\texorpdfstring{\softcode{ACORDADA}}{ACORDADA},
\texorpdfstring{\softcode{CONFIRMADA}}{CONFIRMADA},
\texorpdfstring{\softcode{CANCELADA}}{CANCELADA},
\texorpdfstring{\softcode{ANULADA}}{ANULADA}.

\section{\texorpdfstring{\texorpdfstring{\softtext{EstadoPosicion}}{EstadoPosicion}}{}}\label{estadoposicion}

Definición contractual reutilizable de tipo
\texorpdfstring{\softtext{EstadoPosicion.}}{EstadoPosicion.}

\textbf{Tipo:} string.

\textbf{Valores permitidos:}
\texorpdfstring{\softcode{PENDIENTE\_RECEPCION}}{PENDIENTE\_RECEPCION},
\texorpdfstring{\softcode{DISPONIBLE}}{DISPONIBLE},
\texorpdfstring{\softcode{VENDIDA\_PENDIENTE\_ENTREGA}}{VENDIDA\_PENDIENTE\_ENTREGA},
\texorpdfstring{\softcode{CERRADA\_VENTA}}{CERRADA\_VENTA},
\texorpdfstring{\softcode{BAJA}}{BAJA}.

\section{\texorpdfstring{\texorpdfstring{\softtext{OperacionResumen}}{OperacionResumen}}{}}\label{operacionresumen}

Definición contractual reutilizable de tipo
\texorpdfstring{\softtext{OperacionResumen.}}{OperacionResumen.}

\textbf{Tipo:} object.

{\def\LTcaptype{none} % do not increment counter
\begin{longtable}[]{@{}
  >{\raggedright\arraybackslash}p{(\linewidth - 4\tabcolsep) * \real{0.3333}}
  >{\raggedright\arraybackslash}p{(\linewidth - 4\tabcolsep) * \real{0.3333}}
  >{\raggedright\arraybackslash}p{(\linewidth - 4\tabcolsep) * \real{0.3333}}@{}}
\toprule\noalign{}
\begin{minipage}[b]{\linewidth}\raggedright
Campo
\end{minipage} & \begin{minipage}[b]{\linewidth}\raggedright
Tipo y presencia
\end{minipage} & \begin{minipage}[b]{\linewidth}\raggedright
Restricción o significado
\end{minipage} \\
\midrule\noalign{}
\endhead
\bottomrule\noalign{}
\endlastfoot
\texorpdfstring{\softcode{id}}{id} & string (uuid); opcional & \\
\texorpdfstring{\softcode{tipo}}{tipo} & string; opcional &
enum={[}`COMPRA', `VENTA'{]} \\
\texorpdfstring{\softcode{folio}}{folio} & string; opcional & \\
\texorpdfstring{\softcode{contraparte}}{contraparte} & string; opcional
& \\
\texorpdfstring{\softcode{estado}}{estado} &
\texorpdfstring{\softtext{EstadoOperacion;}}{EstadoOperacion;} opcional
& Ver definición
\texorpdfstring{\softtext{EstadoOperacion}}{EstadoOperacion} \\
\texorpdfstring{\softcode{estado\_entrega}}{estado\_entrega} & string;
opcional & enum={[}`PENDIENTE', `PARCIAL', `COMPLETA'{]} \\
\texorpdfstring{\softcode{pendientes}}{pendientes} & colección de
string; opcional & \\
\end{longtable}
}

\section{Operacion}\label{operacion}

Definición contractual reutilizable de tipo Operacion.

\textbf{Tipo:} objeto.

\textbf{allOf:}
\texorpdfstring{\softtext{OperacionResumen,}}{OperacionResumen,} object.

\section{Version}\label{version}

Definición contractual reutilizable de tipo Version.

\textbf{Tipo:} object.

{\def\LTcaptype{none} % do not increment counter
\begin{longtable}[]{@{}
  >{\raggedright\arraybackslash}p{(\linewidth - 4\tabcolsep) * \real{0.3333}}
  >{\raggedright\arraybackslash}p{(\linewidth - 4\tabcolsep) * \real{0.3333}}
  >{\raggedright\arraybackslash}p{(\linewidth - 4\tabcolsep) * \real{0.3333}}@{}}
\toprule\noalign{}
\begin{minipage}[b]{\linewidth}\raggedright
Campo
\end{minipage} & \begin{minipage}[b]{\linewidth}\raggedright
Tipo y presencia
\end{minipage} & \begin{minipage}[b]{\linewidth}\raggedright
Restricción o significado
\end{minipage} \\
\midrule\noalign{}
\endhead
\bottomrule\noalign{}
\endlastfoot
\texorpdfstring{\softcode{numero\_version}}{numero\_version} & integer;
opcional & \\
\texorpdfstring{\softcode{estado}}{estado} & string; opcional &
enum={[}`PROPUESTA', `ACEPTADA', `SUPERADA'{]} \\
\texorpdfstring{\softcode{condiciones}}{condiciones} & colección de
object; opcional & \\
\end{longtable}
}

\section{\texorpdfstring{\texorpdfstring{\softtext{EdicionBorrador}}{EdicionBorrador}}{}}\label{edicionborrador}

Definición contractual reutilizable de tipo
\texorpdfstring{\softtext{EdicionBorrador.}}{EdicionBorrador.}

\textbf{Tipo:} object.

{\def\LTcaptype{none} % do not increment counter
\begin{longtable}[]{@{}lll@{}}
\toprule\noalign{}
Campo & Tipo y presencia & Restricción o significado \\
\midrule\noalign{}
\endhead
\bottomrule\noalign{}
\endlastfoot
\texorpdfstring{\softcode{lotes}}{lotes} & colección de objeto;
requerido & minItems=1 \\
\texorpdfstring{\softcode{acordar}}{acordar} & boolean; opcional & \\
\texorpdfstring{\softcode{requiere\_recepcion}}{requiere\_recepcion} &
boolean; opcional & \\
\end{longtable}
}

\section{\texorpdfstring{\texorpdfstring{\softtext{PesajeNuevo}}{PesajeNuevo}}{}}\label{pesajenuevo}

Definición contractual reutilizable de tipo
\texorpdfstring{\softtext{PesajeNuevo.}}{PesajeNuevo.}

\textbf{Tipo:} object.

{\def\LTcaptype{none} % do not increment counter
\begin{longtable}[]{@{}lll@{}}
\toprule\noalign{}
Campo & Tipo y presencia & Restricción o significado \\
\midrule\noalign{}
\endhead
\bottomrule\noalign{}
\endlastfoot
\texorpdfstring{\softcode{numero\_lote}}{numero\_lote} & integer;
requerido & \\
\texorpdfstring{\softcode{modo}}{modo} & string; requerido &
enum={[}`BRUTO\_TARA',
`\texorpdfstring{\softtext{NETO\_REPORTADO}}{NETO\_REPORTADO}'{]} \\
\texorpdfstring{\softcode{bruto}}{bruto} & Peso; opcional & Ver
definición Peso \\
\texorpdfstring{\softcode{tara}}{tara} & Peso; opcional & Ver definición
Peso \\
\texorpdfstring{\softcode{neto}}{neto} & Peso; requerido & Ver
definición Peso \\
\texorpdfstring{\softcode{ticket\_folio}}{ticket\_folio} & string;
opcional & \\
\texorpdfstring{\softcode{pesado\_en}}{pesado\_en} & string (date-time);
requerido & \\
\texorpdfstring{\softcode{documento\_id}}{documento\_id} & string
(uuid); opcional & \\
\texorpdfstring{\softcode{animales}}{animales} & colección de object;
requerido & \\
\end{longtable}
}

\section{\texorpdfstring{\texorpdfstring{\softtext{PesajeCorreccion}}{PesajeCorreccion}}{}}\label{pesajecorreccion}

Definición contractual reutilizable de tipo
\texorpdfstring{\softtext{PesajeCorreccion.}}{PesajeCorreccion.}

\textbf{Tipo:} objeto.

\textbf{allOf:} \texorpdfstring{\softtext{PesajeNuevo,}}{PesajeNuevo,}
object.

\section{Pesaje}\label{pesaje}

Definición contractual reutilizable de tipo Pesaje.

\textbf{Tipo:} objeto.

\textbf{allOf:} \texorpdfstring{\softtext{PesajeNuevo,}}{PesajeNuevo,}
object.

\section{Reserva}\label{reserva}

Definición contractual reutilizable de tipo Reserva.

\textbf{Tipo:} object.

{\def\LTcaptype{none} % do not increment counter
\begin{longtable}[]{@{}
  >{\raggedright\arraybackslash}p{(\linewidth - 4\tabcolsep) * \real{0.3333}}
  >{\raggedright\arraybackslash}p{(\linewidth - 4\tabcolsep) * \real{0.3333}}
  >{\raggedright\arraybackslash}p{(\linewidth - 4\tabcolsep) * \real{0.3333}}@{}}
\toprule\noalign{}
\begin{minipage}[b]{\linewidth}\raggedright
Campo
\end{minipage} & \begin{minipage}[b]{\linewidth}\raggedright
Tipo y presencia
\end{minipage} & \begin{minipage}[b]{\linewidth}\raggedright
Restricción o significado
\end{minipage} \\
\midrule\noalign{}
\endhead
\bottomrule\noalign{}
\endlastfoot
\texorpdfstring{\softcode{id}}{id} & string (uuid); opcional & \\
\texorpdfstring{\softcode{posicion\_id}}{posicion\_id} & string (uuid);
opcional & \\
\texorpdfstring{\softcode{estado}}{estado} & string; opcional &
enum={[}`ACTIVA', `CONSUMIDA', `LIBERADA', `EXPIRADA'{]} \\
\texorpdfstring{\softcode{vence\_en}}{vence\_en} & string (date-time);
opcional & \\
\end{longtable}
}

\section{Posicion}\label{posicion}

Definición contractual reutilizable de tipo Posicion.

\textbf{Tipo:} object.

{\def\LTcaptype{none} % do not increment counter
\begin{longtable}[]{@{}lll@{}}
\toprule\noalign{}
Campo & Tipo y presencia & Restricción o significado \\
\midrule\noalign{}
\endhead
\bottomrule\noalign{}
\endlastfoot
\texorpdfstring{\softcode{id}}{id} & string (uuid); opcional & \\
\texorpdfstring{\softcode{animal\_id}}{animal\_id} & string (uuid);
opcional & \\
\texorpdfstring{\softcode{numero\_ciclo}}{numero\_ciclo} & integer;
opcional & \\
\texorpdfstring{\softcode{estado}}{estado} &
\texorpdfstring{\softtext{EstadoPosicion;}}{EstadoPosicion;} opcional &
Ver definición
\texorpdfstring{\softtext{EstadoPosicion}}{EstadoPosicion} \\
\texorpdfstring{\softcode{reservada}}{reservada} & boolean; opcional &
Derivado de una reserva ACTIVA \\
\texorpdfstring{\softcode{costo\_vigente}}{costo\_vigente} & Dinero;
opcional & Ver definición Dinero \\
\end{longtable}
}

\section{Liquidacion}\label{liquidacion}

Definición contractual reutilizable de tipo
\texorpdfstring{\softtext{Liquidacion.}}{Liquidacion.}

\textbf{Tipo:} object.

{\def\LTcaptype{none} % do not increment counter
\begin{longtable}[]{@{}lll@{}}
\toprule\noalign{}
Campo & Tipo y presencia & Restricción o significado \\
\midrule\noalign{}
\endhead
\bottomrule\noalign{}
\endlastfoot
\texorpdfstring{\softcode{id}}{id} & string (uuid); opcional & \\
\texorpdfstring{\softcode{operacion\_id}}{operacion\_id} & string
(uuid); opcional & \\
\texorpdfstring{\softcode{direccion}}{direccion} & string; opcional &
enum={[}`POR\_PAGAR', `POR\_COBRAR'{]} \\
\texorpdfstring{\softcode{total}}{total} & Dinero; opcional & Ver
definición Dinero \\
\texorpdfstring{\softcode{ajustes\_cargo}}{ajustes\_cargo} & Dinero;
opcional & Ver definición Dinero \\
\texorpdfstring{\softcode{ajustes\_abono}}{ajustes\_abono} & Dinero;
opcional & Ver definición Dinero \\
\texorpdfstring{\softcode{aplicaciones\_netas}}{aplicaciones\_netas} &
Dinero; opcional & Ver definición Dinero \\
\texorpdfstring{\softcode{saldo}}{saldo} & Dinero; opcional & Ver
definición Dinero \\
\texorpdfstring{\softcode{vence\_el}}{vence\_el} & string (date);
opcional & \\
\texorpdfstring{\softcode{conceptos}}{conceptos} & colección de object;
opcional & \\
\end{longtable}
}

\section{GastoNuevo}\label{gastonuevo}

Definición contractual reutilizable de tipo
\texorpdfstring{\softtext{GastoNuevo.}}{GastoNuevo.}

\textbf{Tipo:} object.

{\def\LTcaptype{none} % do not increment counter
\begin{longtable}[]{@{}
  >{\raggedright\arraybackslash}p{(\linewidth - 4\tabcolsep) * \real{0.3333}}
  >{\raggedright\arraybackslash}p{(\linewidth - 4\tabcolsep) * \real{0.3333}}
  >{\raggedright\arraybackslash}p{(\linewidth - 4\tabcolsep) * \real{0.3333}}@{}}
\toprule\noalign{}
\begin{minipage}[b]{\linewidth}\raggedright
Campo
\end{minipage} & \begin{minipage}[b]{\linewidth}\raggedright
Tipo y presencia
\end{minipage} & \begin{minipage}[b]{\linewidth}\raggedright
Restricción o significado
\end{minipage} \\
\midrule\noalign{}
\endhead
\bottomrule\noalign{}
\endlastfoot
\texorpdfstring{\softcode{operacion\_id}}{operacion\_id} & string
(uuid); requerido & \\
\texorpdfstring{\softcode{naturaleza}}{naturaleza} & string; requerido &
enum={[}`COSTO', `BONIFICACION'{]} \\
\texorpdfstring{\softcode{clasificacion}}{clasificacion} & string;
requerido & enum={[}`CAPITALIZABLE',
`\texorpdfstring{\softtext{DIRECTO\_VENTA}}{DIRECTO\_VENTA}',
`OTRO'{]} \\
\texorpdfstring{\softcode{importe}}{importe} & Dinero; requerido & Ver
definición Dinero \\
\texorpdfstring{\softcode{concepto}}{concepto} & string; requerido &
\texorpdfstring{\softtext{maxLength=160}}{maxLength=160} \\
\texorpdfstring{\softcode{proveedor\_parte\_id}}{proveedor\_parte\_id} &
string / null (uuid); opcional & \\
\texorpdfstring{\softcode{referencia\_origen}}{referencia\_origen} &
string / null; opcional &
\texorpdfstring{\softtext{maxLength=60}}{maxLength=60} \\
\texorpdfstring{\softcode{bonifica\_gasto\_id}}{bonifica\_gasto\_id} &
string / null (uuid); opcional & \\
\texorpdfstring{\softcode{fecha\_economica}}{fecha\_economica} & string
(date); requerido & \\
\texorpdfstring{\softcode{periodo\_id}}{periodo\_id} & string (uuid);
requerido & \\
\texorpdfstring{\softcode{fecha\_reconocimiento}}{fecha\_reconocimiento}
& string (date); requerido & \\
\texorpdfstring{\softcode{snapshot\_id}}{snapshot\_id} & string (uuid);
opcional & \\
\end{longtable}
}

\section{Asignacion}\label{asignacion}

Definición contractual reutilizable de tipo
\texorpdfstring{\softtext{Asignacion.}}{Asignacion.}

\textbf{Tipo:} object.

{\def\LTcaptype{none} % do not increment counter
\begin{longtable}[]{@{}
  >{\raggedright\arraybackslash}p{(\linewidth - 4\tabcolsep) * \real{0.3333}}
  >{\raggedright\arraybackslash}p{(\linewidth - 4\tabcolsep) * \real{0.3333}}
  >{\raggedright\arraybackslash}p{(\linewidth - 4\tabcolsep) * \real{0.3333}}@{}}
\toprule\noalign{}
\begin{minipage}[b]{\linewidth}\raggedright
Campo
\end{minipage} & \begin{minipage}[b]{\linewidth}\raggedright
Tipo y presencia
\end{minipage} & \begin{minipage}[b]{\linewidth}\raggedright
Restricción o significado
\end{minipage} \\
\midrule\noalign{}
\endhead
\bottomrule\noalign{}
\endlastfoot
\texorpdfstring{\softcode{posicion\_id}}{posicion\_id} & string (uuid);
opcional & \\
\texorpdfstring{\softcode{base\_valor}}{base\_valor} & string; opcional
& \\
\texorpdfstring{\softcode{importe}}{importe} & Dinero; opcional & Ver
definición Dinero \\
\texorpdfstring{\softcode{destino}}{destino} & string; opcional &
enum={[}`INVENTARIO',
`\texorpdfstring{\softtext{COSTO\_VENTA}}{COSTO\_VENTA}',
`PERDIDA'{]} \\
\end{longtable}
}

\section{PagoNuevo}\label{pagonuevo}

Definición contractual reutilizable de tipo PagoNuevo.

\textbf{Tipo:} object.

{\def\LTcaptype{none} % do not increment counter
\begin{longtable}[]{@{}
  >{\raggedright\arraybackslash}p{(\linewidth - 4\tabcolsep) * \real{0.3333}}
  >{\raggedright\arraybackslash}p{(\linewidth - 4\tabcolsep) * \real{0.3333}}
  >{\raggedright\arraybackslash}p{(\linewidth - 4\tabcolsep) * \real{0.3333}}@{}}
\toprule\noalign{}
\begin{minipage}[b]{\linewidth}\raggedright
Campo
\end{minipage} & \begin{minipage}[b]{\linewidth}\raggedright
Tipo y presencia
\end{minipage} & \begin{minipage}[b]{\linewidth}\raggedright
Restricción o significado
\end{minipage} \\
\midrule\noalign{}
\endhead
\bottomrule\noalign{}
\endlastfoot
\texorpdfstring{\softcode{parte\_id}}{parte\_id} & string (uuid);
requerido & \\
\texorpdfstring{\softcode{tipo}}{tipo} & string; requerido &
enum={[}`COBRO', `PAGO'{]} \\
\texorpdfstring{\softcode{importe}}{importe} & Dinero; requerido & Ver
definición Dinero \\
\texorpdfstring{\softcode{metodo}}{metodo} & string; requerido &
enum={[}`TRANSFERENCIA', `EFECTIVO', `CHEQUE', `OTRO'{]} \\
\texorpdfstring{\softcode{referencia}}{referencia} & string / null;
opcional & \texorpdfstring{\softtext{maxLength=60}}{maxLength=60} \\
\texorpdfstring{\softcode{recibido\_el}}{recibido\_el} & string (date);
requerido & \\
\end{longtable}
}

\section{Pago}\label{pago}

Definición contractual reutilizable de tipo Pago.

\textbf{Tipo:} objeto.

\textbf{allOf:} PagoNuevo, object.

\section{Documento}\label{documento}

Definición contractual reutilizable de tipo Documento.

\textbf{Tipo:} object.

{\def\LTcaptype{none} % do not increment counter
\begin{longtable}[]{@{}lll@{}}
\toprule\noalign{}
Campo & Tipo y presencia & Restricción o significado \\
\midrule\noalign{}
\endhead
\bottomrule\noalign{}
\endlastfoot
\texorpdfstring{\softcode{id}}{id} & string (uuid); opcional & \\
\texorpdfstring{\softcode{tipo\_documental}}{tipo\_documental} & string;
opcional & \\
\texorpdfstring{\softcode{sha256}}{sha256} & string; opcional &
pattern=\footnote{0-9a-f}\{64\}\$ \\
\texorpdfstring{\softcode{tamano\_bytes}}{tamano\_bytes} & integer;
opcional & \\
\texorpdfstring{\softcode{mime}}{mime} & string; opcional & \\
\texorpdfstring{\softcode{version}}{version} & integer; opcional & \\
\end{longtable}
}

\section{\texorpdfstring{\texorpdfstring{\softtext{ReporteCartera}}{ReporteCartera}}{}}\label{reportecartera}

Definición contractual reutilizable de tipo
\texorpdfstring{\softtext{ReporteCartera.}}{ReporteCartera.}

\textbf{Tipo:} object.

{\def\LTcaptype{none} % do not increment counter
\begin{longtable}[]{@{}
  >{\raggedright\arraybackslash}p{(\linewidth - 4\tabcolsep) * \real{0.3333}}
  >{\raggedright\arraybackslash}p{(\linewidth - 4\tabcolsep) * \real{0.3333}}
  >{\raggedright\arraybackslash}p{(\linewidth - 4\tabcolsep) * \real{0.3333}}@{}}
\toprule\noalign{}
\begin{minipage}[b]{\linewidth}\raggedright
Campo
\end{minipage} & \begin{minipage}[b]{\linewidth}\raggedright
Tipo y presencia
\end{minipage} & \begin{minipage}[b]{\linewidth}\raggedright
Restricción o significado
\end{minipage} \\
\midrule\noalign{}
\endhead
\bottomrule\noalign{}
\endlastfoot
\texorpdfstring{\softcode{corte\_en}}{corte\_en} & string (date-time);
requerido & \\
\texorpdfstring{\softcode{etiqueta}}{etiqueta} & string; requerido &
enum={[}`ACTUALIZADO'{]} \\
\texorpdfstring{\softcode{fuente\_consulta}}{fuente\_consulta} & string;
requerido & enum={[}`PROPIETARIO', `REPLICA'{]} \\
\texorpdfstring{\softcode{antiguedad\_replica\_segundos}}{antiguedad\_replica\_segundos}
& integer / null; requerido & 0 para lectura propietaria; null si la
medida de réplica es
\texorpdfstring{\softtext{desconocida.}}{desconocida.} \\
\texorpdfstring{\softcode{pagina}}{pagina} & integer; requerido &
minimum=1 \\
\texorpdfstring{\softcode{tamano}}{tamano} & integer; requerido &
minimum=1; maximum=100 \\
\texorpdfstring{\softcode{total}}{total} & integer; requerido &
minimum=0 \\
\texorpdfstring{\softcode{filas}}{filas} & colección de object;
requerido & \\
\end{longtable}
}

\section{\texorpdfstring{\texorpdfstring{\softtext{ReporteMargen}}{ReporteMargen}}{}}\label{reportemargen}

Definición contractual reutilizable de tipo
\texorpdfstring{\softtext{ReporteMargen.}}{ReporteMargen.}

\textbf{Tipo:} object.

{\def\LTcaptype{none} % do not increment counter
\begin{longtable}[]{@{}
  >{\raggedright\arraybackslash}p{(\linewidth - 4\tabcolsep) * \real{0.3333}}
  >{\raggedright\arraybackslash}p{(\linewidth - 4\tabcolsep) * \real{0.3333}}
  >{\raggedright\arraybackslash}p{(\linewidth - 4\tabcolsep) * \real{0.3333}}@{}}
\toprule\noalign{}
\begin{minipage}[b]{\linewidth}\raggedright
Campo
\end{minipage} & \begin{minipage}[b]{\linewidth}\raggedright
Tipo y presencia
\end{minipage} & \begin{minipage}[b]{\linewidth}\raggedright
Restricción o significado
\end{minipage} \\
\midrule\noalign{}
\endhead
\bottomrule\noalign{}
\endlastfoot
\texorpdfstring{\softcode{corte\_en}}{corte\_en} & string (date-time);
requerido & \\
\texorpdfstring{\softcode{etiqueta}}{etiqueta} & string; requerido &
enum={[}`PROVISIONAL', `ACTUALIZADO', `CERRADO'{]} \\
\texorpdfstring{\softcode{fuente\_consulta}}{fuente\_consulta} & string;
requerido & enum={[}`PROPIETARIO', `REPLICA'{]} \\
\texorpdfstring{\softcode{antiguedad\_replica\_segundos}}{antiguedad\_replica\_segundos}
& integer / null; requerido & 0 para lectura propietaria; null si la
medida de réplica es
\texorpdfstring{\softtext{desconocida.}}{desconocida.} \\
\texorpdfstring{\softcode{venta\_neta}}{venta\_neta} & Dinero; requerido
& Ver definición Dinero \\
\texorpdfstring{\softcode{costo\_vendido}}{costo\_vendido} & Dinero;
requerido & Ver definición Dinero \\
\texorpdfstring{\softcode{gastos\_directos}}{gastos\_directos} & Dinero;
requerido & Ver definición Dinero \\
\texorpdfstring{\softcode{margen}}{margen} & Dinero; requerido & Ver
definición Dinero \\
\texorpdfstring{\softcode{perdidas}}{perdidas} & Dinero; requerido & Ver
definición Dinero \\
\texorpdfstring{\softcode{faltantes}}{faltantes} & colección de string;
requerido & \\
\end{longtable}
}

\section{\texorpdfstring{\texorpdfstring{\softtext{TarifaExterna}}{TarifaExterna}}{}}\label{tarifaexterna}

Definición contractual reutilizable de tipo
\texorpdfstring{\softtext{TarifaExterna.}}{TarifaExterna.}

\textbf{Tipo:} object.

{\def\LTcaptype{none} % do not increment counter
\begin{longtable}[]{@{}lll@{}}
\toprule\noalign{}
Campo & Tipo y presencia & Restricción o significado \\
\midrule\noalign{}
\endhead
\bottomrule\noalign{}
\endlastfoot
\texorpdfstring{\softcode{fuente}}{fuente} & string; opcional & \\
\texorpdfstring{\softcode{ruta}}{ruta} & string; opcional & \\
\texorpdfstring{\softcode{precio}}{precio} & Decimal4; opcional & Ver
definición Decimal4 \\
\texorpdfstring{\softcode{unidad}}{unidad} & string; opcional & \\
\texorpdfstring{\softcode{vigente\_desde}}{vigente\_desde} & string
(date); opcional & \\
\texorpdfstring{\softcode{version}}{version} & integer; opcional & \\
\end{longtable}
}

\section{Csrf}\label{csrf}

Definición contractual reutilizable de tipo Csrf.

\textbf{Tipo:} object.

{\def\LTcaptype{none} % do not increment counter
\begin{longtable}[]{@{}lll@{}}
\toprule\noalign{}
Campo & Tipo y presencia & Restricción o significado \\
\midrule\noalign{}
\endhead
\bottomrule\noalign{}
\endlastfoot
\texorpdfstring{\softcode{token}}{token} & string; requerido & \\
\texorpdfstring{\softcode{headerName}}{headerName} & string; requerido
& \\
\end{longtable}
}

\section{Reconocimiento}\label{reconocimiento}

Definición contractual reutilizable de tipo
\texorpdfstring{\softtext{Reconocimiento.}}{Reconocimiento.}

\textbf{Tipo:} object.

{\def\LTcaptype{none} % do not increment counter
\begin{longtable}[]{@{}lll@{}}
\toprule\noalign{}
Campo & Tipo y presencia & Restricción o significado \\
\midrule\noalign{}
\endhead
\bottomrule\noalign{}
\endlastfoot
\texorpdfstring{\softcode{periodo\_id}}{periodo\_id} & string (uuid);
requerido & \\
\texorpdfstring{\softcode{fecha\_reconocimiento}}{fecha\_reconocimiento}
& string (date); requerido & \\
\end{longtable}
}

\section{\texorpdfstring{\texorpdfstring{\softtext{LoteCompraInput}}{LoteCompraInput}}{}}\label{lotecomprainput}

Definición contractual reutilizable de tipo
\texorpdfstring{\softtext{LoteCompraInput.}}{LoteCompraInput.}

\textbf{Tipo:} object.

{\def\LTcaptype{none} % do not increment counter
\begin{longtable}[]{@{}
  >{\raggedright\arraybackslash}p{(\linewidth - 4\tabcolsep) * \real{0.3333}}
  >{\raggedright\arraybackslash}p{(\linewidth - 4\tabcolsep) * \real{0.3333}}
  >{\raggedright\arraybackslash}p{(\linewidth - 4\tabcolsep) * \real{0.3333}}@{}}
\toprule\noalign{}
\begin{minipage}[b]{\linewidth}\raggedright
Campo
\end{minipage} & \begin{minipage}[b]{\linewidth}\raggedright
Tipo y presencia
\end{minipage} & \begin{minipage}[b]{\linewidth}\raggedright
Restricción o significado
\end{minipage} \\
\midrule\noalign{}
\endhead
\bottomrule\noalign{}
\endlastfoot
\texorpdfstring{\softcode{numero\_lote}}{numero\_lote} & integer;
requerido & minimum=1 \\
\texorpdfstring{\softcode{descripcion}}{descripcion} & string; opcional
& \texorpdfstring{\softtext{maxLength=120}}{maxLength=120} \\
\texorpdfstring{\softcode{animales}}{animales} & colección de string
(uuid); requerido & minItems=1 \\
\texorpdfstring{\softcode{unidad\_precio}}{unidad\_precio} & string;
requerido & enum={[}`KG', `CABEZA'{]} \\
\texorpdfstring{\softcode{precio\_unitario}}{precio\_unitario} &
Decimal4; requerido & Ver definición Decimal4 \\
\texorpdfstring{\softcode{pesaje\_base\_id}}{pesaje\_base\_id} & string
(uuid); opcional & \\
\texorpdfstring{\softcode{deduccion\_pct}}{deduccion\_pct} & string;
opcional & pattern=\^{}\d{1,2}.\d{1,6}\$ \\
\texorpdfstring{\softcode{cargos\_pactados}}{cargos\_pactados} &
Decimal2; opcional & Ver definición Decimal2 \\
\texorpdfstring{\softcode{bonificaciones\_pactadas}}{bonificaciones\_pactadas}
& Decimal2; opcional & Ver definición Decimal2 \\
\end{longtable}
}

\section{\texorpdfstring{\texorpdfstring{\softtext{LoteVentaInput}}{LoteVentaInput}}{}}\label{loteventainput}

Definición contractual reutilizable de tipo
\texorpdfstring{\softtext{LoteVentaInput.}}{LoteVentaInput.}

\textbf{Tipo:} object.

{\def\LTcaptype{none} % do not increment counter
\begin{longtable}[]{@{}
  >{\raggedright\arraybackslash}p{(\linewidth - 4\tabcolsep) * \real{0.3333}}
  >{\raggedright\arraybackslash}p{(\linewidth - 4\tabcolsep) * \real{0.3333}}
  >{\raggedright\arraybackslash}p{(\linewidth - 4\tabcolsep) * \real{0.3333}}@{}}
\toprule\noalign{}
\begin{minipage}[b]{\linewidth}\raggedright
Campo
\end{minipage} & \begin{minipage}[b]{\linewidth}\raggedright
Tipo y presencia
\end{minipage} & \begin{minipage}[b]{\linewidth}\raggedright
Restricción o significado
\end{minipage} \\
\midrule\noalign{}
\endhead
\bottomrule\noalign{}
\endlastfoot
\texorpdfstring{\softcode{numero\_lote}}{numero\_lote} & integer;
requerido & minimum=1 \\
\texorpdfstring{\softcode{descripcion}}{descripcion} & string; opcional
& \texorpdfstring{\softtext{maxLength=120}}{maxLength=120} \\
\texorpdfstring{\softcode{unidad\_precio}}{unidad\_precio} & string;
requerido & enum={[}`KG', `CABEZA'{]} \\
\texorpdfstring{\softcode{precio\_unitario}}{precio\_unitario} &
Decimal4; requerido & Ver definición Decimal4 \\
\texorpdfstring{\softcode{pesaje\_base\_id}}{pesaje\_base\_id} & string
(uuid); opcional & \\
\texorpdfstring{\softcode{deduccion\_pct}}{deduccion\_pct} & string;
opcional & pattern=\^{}\d{1,2}.\d{1,6}\$ \\
\texorpdfstring{\softcode{cargos\_pactados}}{cargos\_pactados} &
Decimal2; opcional & Ver definición Decimal2 \\
\texorpdfstring{\softcode{bonificaciones\_pactadas}}{bonificaciones\_pactadas}
& Decimal2; opcional & Ver definición Decimal2 \\
\texorpdfstring{\softcode{posiciones}}{posiciones} & colección de string
(uuid); requerido & minItems=1 \\
\end{longtable}
}

\section{\texorpdfstring{\texorpdfstring{\softtext{DistribucionInput}}{DistribucionInput}}{}}\label{distribucioninput}

Definición contractual reutilizable de tipo
\texorpdfstring{\softtext{DistribucionInput.}}{DistribucionInput.}

\textbf{Tipo:} object.

{\def\LTcaptype{none} % do not increment counter
\begin{longtable}[]{@{}
  >{\raggedright\arraybackslash}p{(\linewidth - 4\tabcolsep) * \real{0.3333}}
  >{\raggedright\arraybackslash}p{(\linewidth - 4\tabcolsep) * \real{0.3333}}
  >{\raggedright\arraybackslash}p{(\linewidth - 4\tabcolsep) * \real{0.3333}}@{}}
\toprule\noalign{}
\begin{minipage}[b]{\linewidth}\raggedright
Campo
\end{minipage} & \begin{minipage}[b]{\linewidth}\raggedright
Tipo y presencia
\end{minipage} & \begin{minipage}[b]{\linewidth}\raggedright
Restricción o significado
\end{minipage} \\
\midrule\noalign{}
\endhead
\bottomrule\noalign{}
\endlastfoot
\texorpdfstring{\softcode{metodo}}{metodo} & string; requerido &
enum={[}`POR\_CABEZA',
`\texorpdfstring{\softtext{POR\_PESO\_MEDIDO}}{POR\_PESO\_MEDIDO}',
`\texorpdfstring{\softtext{IMPORTE\_INDIVIDUAL}}{IMPORTE\_INDIVIDUAL}'{]} \\
\texorpdfstring{\softcode{periodo\_id}}{periodo\_id} & string (uuid);
requerido & \\
\texorpdfstring{\softcode{fecha\_reconocimiento}}{fecha\_reconocimiento}
& string (date); requerido & \\
\texorpdfstring{\softcode{asignaciones}}{asignaciones} & colección de
object; opcional & \\
\end{longtable}
}

\section{Gasto}\label{gasto}

Definición contractual reutilizable de tipo Gasto.

\textbf{Tipo:} objeto.

\textbf{allOf:} \texorpdfstring{\softtext{GastoNuevo,}}{GastoNuevo,}
object.

\section{\texorpdfstring{\texorpdfstring{\softtext{ResultadoComando}}{ResultadoComando}}{}}\label{resultadocomando}

Definición contractual reutilizable de tipo
\texorpdfstring{\softtext{ResultadoComando.}}{ResultadoComando.}

\textbf{Tipo:} object.

{\def\LTcaptype{none} % do not increment counter
\begin{longtable}[]{@{}lll@{}}
\toprule\noalign{}
Campo & Tipo y presencia & Restricción o significado \\
\midrule\noalign{}
\endhead
\bottomrule\noalign{}
\endlastfoot
\texorpdfstring{\softcode{id}}{id} & string (uuid); opcional & \\
\texorpdfstring{\softcode{version\_lock}}{version\_lock} & integer;
opcional & minimum=0 \\
\end{longtable}
}

\section{\texorpdfstring{\texorpdfstring{\softtext{OperacionNueva}}{OperacionNueva}}{}}\label{operacionnueva}

Definición contractual reutilizable de tipo
\texorpdfstring{\softtext{OperacionNueva.}}{OperacionNueva.}

\textbf{Tipo:} object.

{\def\LTcaptype{none} % do not increment counter
\begin{longtable}[]{@{}
  >{\raggedright\arraybackslash}p{(\linewidth - 4\tabcolsep) * \real{0.3333}}
  >{\raggedright\arraybackslash}p{(\linewidth - 4\tabcolsep) * \real{0.3333}}
  >{\raggedright\arraybackslash}p{(\linewidth - 4\tabcolsep) * \real{0.3333}}@{}}
\toprule\noalign{}
\begin{minipage}[b]{\linewidth}\raggedright
Campo
\end{minipage} & \begin{minipage}[b]{\linewidth}\raggedright
Tipo y presencia
\end{minipage} & \begin{minipage}[b]{\linewidth}\raggedright
Restricción o significado
\end{minipage} \\
\midrule\noalign{}
\endhead
\bottomrule\noalign{}
\endlastfoot
\texorpdfstring{\softcode{tipo}}{tipo} & string; requerido &
enum={[}`COMPRA', `VENTA'{]} \\
\texorpdfstring{\softcode{parte\_id}}{parte\_id} & string (uuid);
requerido & \\
\texorpdfstring{\softcode{fecha\_operacion}}{fecha\_operacion} & string
(date); requerido & \\
\texorpdfstring{\softcode{requiere\_recepcion}}{requiere\_recepcion} &
boolean; opcional & En compra: conserva la posición pendiente de
recepción física cuando es true. default=False \\
\texorpdfstring{\softcode{lotes}}{lotes} & colección de objeto; opcional
& minItems=1 \\
\end{longtable}
}

\section{Usuario}\label{usuario}

Definición contractual reutilizable de tipo Usuario.

\textbf{Tipo:} object.

{\def\LTcaptype{none} % do not increment counter
\begin{longtable}[]{@{}lll@{}}
\toprule\noalign{}
Campo & Tipo y presencia & Restricción o significado \\
\midrule\noalign{}
\endhead
\bottomrule\noalign{}
\endlastfoot
\texorpdfstring{\softcode{id}}{id} & string (uuid); requerido & \\
\texorpdfstring{\softcode{nombre}}{nombre} & string; requerido & \\
\texorpdfstring{\softcode{login}}{login} & string; requerido & \\
\texorpdfstring{\softcode{estado}}{estado} & string; requerido &
enum={[}`ACTIVO', `BLOQUEADO'{]} \\
\texorpdfstring{\softcode{roles}}{roles} & colección de string; opcional
& \\
\end{longtable}
}

\section{\texorpdfstring{\texorpdfstring{\softtext{UsuarioNuevo}}{UsuarioNuevo}}{}}\label{usuarionuevo}

Definición contractual reutilizable de tipo
\texorpdfstring{\softtext{UsuarioNuevo.}}{UsuarioNuevo.}

\textbf{Tipo:} object.

{\def\LTcaptype{none} % do not increment counter
\begin{longtable}[]{@{}
  >{\raggedright\arraybackslash}p{(\linewidth - 4\tabcolsep) * \real{0.3333}}
  >{\raggedright\arraybackslash}p{(\linewidth - 4\tabcolsep) * \real{0.3333}}
  >{\raggedright\arraybackslash}p{(\linewidth - 4\tabcolsep) * \real{0.3333}}@{}}
\toprule\noalign{}
\begin{minipage}[b]{\linewidth}\raggedright
Campo
\end{minipage} & \begin{minipage}[b]{\linewidth}\raggedright
Tipo y presencia
\end{minipage} & \begin{minipage}[b]{\linewidth}\raggedright
Restricción o significado
\end{minipage} \\
\midrule\noalign{}
\endhead
\bottomrule\noalign{}
\endlastfoot
\texorpdfstring{\softcode{login}}{login} & string; requerido &
\texorpdfstring{\softtext{minLength=1;}}{minLength=1;}
\texorpdfstring{\softtext{maxLength=100}}{maxLength=100} \\
\texorpdfstring{\softcode{nombre}}{nombre} & string; requerido &
\texorpdfstring{\softtext{minLength=1;}}{minLength=1;}
\texorpdfstring{\softtext{maxLength=200}}{maxLength=200} \\
\texorpdfstring{\softcode{contrasena}}{contrasena} & string; requerido &
Máximo 72 bytes UTF-8 para bcrypt; no se devuelve.
\texorpdfstring{\softtext{minLength=1;}}{minLength=1;}
\texorpdfstring{\softtext{maxLength=72}}{maxLength=72} \\
\end{longtable}
}

\section{Periodo}\label{periodo}

Definición contractual reutilizable de tipo Periodo.

\textbf{Tipo:} object.

{\def\LTcaptype{none} % do not increment counter
\begin{longtable}[]{@{}lll@{}}
\toprule\noalign{}
Campo & Tipo y presencia & Restricción o significado \\
\midrule\noalign{}
\endhead
\bottomrule\noalign{}
\endlastfoot
\texorpdfstring{\softcode{id}}{id} & string (uuid); requerido & \\
\texorpdfstring{\softcode{fecha\_inicio}}{fecha\_inicio} & string
(date); opcional & \\
\texorpdfstring{\softcode{fecha\_fin}}{fecha\_fin} & string (date);
opcional & \\
\texorpdfstring{\softcode{estado}}{estado} & string; requerido &
enum={[}`ABIERTO', `CERRADO'{]} \\
\texorpdfstring{\softcode{desde}}{desde} & string (date); requerido & \\
\texorpdfstring{\softcode{hasta}}{hasta} & string (date); requerido & \\
\end{longtable}
}

\section{Predio}\label{predio}

Definición contractual reutilizable de tipo Predio.

\textbf{Tipo:} object.

{\def\LTcaptype{none} % do not increment counter
\begin{longtable}[]{@{}lll@{}}
\toprule\noalign{}
Campo & Tipo y presencia & Restricción o significado \\
\midrule\noalign{}
\endhead
\bottomrule\noalign{}
\endlastfoot
\texorpdfstring{\softcode{id}}{id} & string (uuid); requerido & \\
\texorpdfstring{\softcode{nombre}}{nombre} & string; requerido & \\
\texorpdfstring{\softcode{municipio}}{municipio} & string; opcional & \\
\texorpdfstring{\softcode{entidad\_federativa}}{entidad\_federativa} &
string; opcional & \\
\texorpdfstring{\softcode{parte\_responsable\_id}}{parte\_responsable\_id}
& string (uuid); opcional & \\
\end{longtable}
}

\section{\texorpdfstring{\texorpdfstring{\softtext{PredioNuevo}}{PredioNuevo}}{}}\label{predionuevo}

Definición contractual reutilizable de tipo
\texorpdfstring{\softtext{PredioNuevo.}}{PredioNuevo.}

\textbf{Tipo:} object.

{\def\LTcaptype{none} % do not increment counter
\begin{longtable}[]{@{}lll@{}}
\toprule\noalign{}
Campo & Tipo y presencia & Restricción o significado \\
\midrule\noalign{}
\endhead
\bottomrule\noalign{}
\endlastfoot
\texorpdfstring{\softcode{nombre}}{nombre} & string; requerido & \\
\texorpdfstring{\softcode{municipio}}{municipio} & string; opcional & \\
\texorpdfstring{\softcode{entidad\_federativa}}{entidad\_federativa} &
string; opcional & \\
\texorpdfstring{\softcode{parte\_responsable\_id}}{parte\_responsable\_id}
& string (uuid); opcional & \\
\texorpdfstring{\softcode{clave\_upp\_declarada}}{clave\_upp\_declarada}
& string; opcional & \\
\end{longtable}
}

\section{\texorpdfstring{\texorpdfstring{\softtext{CierreVersion}}{CierreVersion}}{}}\label{cierreversion}

Definición contractual reutilizable de tipo
\texorpdfstring{\softtext{CierreVersion.}}{CierreVersion.}

\textbf{Tipo:} object.

{\def\LTcaptype{none} % do not increment counter
\begin{longtable}[]{@{}lll@{}}
\toprule\noalign{}
Campo & Tipo y presencia & Restricción o significado \\
\midrule\noalign{}
\endhead
\bottomrule\noalign{}
\endlastfoot
\texorpdfstring{\softcode{numero\_version}}{numero\_version} & integer;
requerido & \\
\texorpdfstring{\softcode{periodo\_id}}{periodo\_id} & string (uuid);
requerido & \\
\texorpdfstring{\softcode{totales}}{totales} & object; opcional & \\
\texorpdfstring{\softcode{corte\_en}}{corte\_en} & string (date-time);
opcional & \\
\texorpdfstring{\softcode{tipo}}{tipo} & string; opcional &
enum={[}`CIERRE', `REAPERTURA'{]} \\
\end{longtable}
}

\section{\texorpdfstring{\texorpdfstring{\softtext{CalculoLote}}{CalculoLote}}{}}\label{calculolote}

Definición contractual reutilizable de tipo
\texorpdfstring{\softtext{CalculoLote.}}{CalculoLote.}

\textbf{Tipo:} object.

{\def\LTcaptype{none} % do not increment counter
\begin{longtable}[]{@{}lll@{}}
\toprule\noalign{}
Campo & Tipo y presencia & Restricción o significado \\
\midrule\noalign{}
\endhead
\bottomrule\noalign{}
\endlastfoot
\texorpdfstring{\softcode{numero\_lote}}{numero\_lote} & integer;
opcional & \\
\texorpdfstring{\softcode{cantidad}}{cantidad} & string; opcional & \\
\texorpdfstring{\softcode{unidad}}{unidad} & string; opcional &
enum={[}`KG', `CABEZA'{]} \\
\texorpdfstring{\softcode{precio}}{precio} & Decimal4; opcional & Ver
definición Decimal4 \\
\texorpdfstring{\softcode{base}}{base} & Dinero; opcional & Ver
definición Dinero \\
\texorpdfstring{\softcode{total}}{total} & Dinero; opcional & Ver
definición Dinero \\
\end{longtable}
}

\section{\texorpdfstring{\texorpdfstring{\softtext{TarifaCapturada}}{TarifaCapturada}}{}}\label{tarifacapturada}

Definición contractual reutilizable de tipo
\texorpdfstring{\softtext{TarifaCapturada.}}{TarifaCapturada.}

\textbf{Tipo:} objeto.

\textbf{allOf:}
\texorpdfstring{\softtext{TarifaExterna,}}{TarifaExterna,} object.

\section{\texorpdfstring{\texorpdfstring{\softtext{DiagnosticoProceso}}{DiagnosticoProceso}}{}}\label{diagnosticoproceso}

Definición contractual reutilizable de tipo
\texorpdfstring{\softtext{DiagnosticoProceso.}}{DiagnosticoProceso.}

\textbf{Tipo:} object.

{\def\LTcaptype{none} % do not increment counter
\begin{longtable}[]{@{}
  >{\raggedright\arraybackslash}p{(\linewidth - 4\tabcolsep) * \real{0.3333}}
  >{\raggedright\arraybackslash}p{(\linewidth - 4\tabcolsep) * \real{0.3333}}
  >{\raggedright\arraybackslash}p{(\linewidth - 4\tabcolsep) * \real{0.3333}}@{}}
\toprule\noalign{}
\begin{minipage}[b]{\linewidth}\raggedright
Campo
\end{minipage} & \begin{minipage}[b]{\linewidth}\raggedright
Tipo y presencia
\end{minipage} & \begin{minipage}[b]{\linewidth}\raggedright
Restricción o significado
\end{minipage} \\
\midrule\noalign{}
\endhead
\bottomrule\noalign{}
\endlastfoot
\texorpdfstring{\softcode{proceso}}{proceso} & object; requerido & \\
\texorpdfstring{\softcode{comandos}}{comandos} & object; requerido & \\
\texorpdfstring{\softcode{consultas}}{consultas} & object; requerido
& \\
\texorpdfstring{\softcode{pools}}{pools} & object; requerido & \\
\texorpdfstring{\softcode{sincronizacion}}{sincronizacion} & object;
requerido & Estado SQL observado; antigüedad desconocida se mantiene
null. Si falla dependencia:
\texorpdfstring{\softtext{estadoNO\_DISPONIBLE}}{estadoNO\_DISPONIBLE} y
código público. \\
\texorpdfstring{\softcode{capturado\_en}}{capturado\_en} & string
(date-time); requerido & \\
\end{longtable}
}

\chapter{Definiciones canónicas
completas}\label{definiciones-canuxf3nicas-completas}

Los siguientes anexos preservan reglas, campos, casos de uso y
decisiones de \texorpdfstring{\softtext{arquitectura.}}{arquitectura.}
Complementan la explicación, permiten buscar
\texorpdfstring{\softtext{RN/ENT/UC/ADR}}{RN/ENT/UC/ADR} y mantienen los
límites del alcance. Los ADR anteriores deben interpretarse con sus
notas de sustitución por decisiones
\texorpdfstring{\softtext{posteriores.}}{posteriores.} La autoridad del
esquema físico es la última definición aplicada por las migraciones; no
se ejecuta SQL desde estos párrafos.

\section{\texorpdfstring{Reglas de negocio
\texorpdfstring{\softtext{RN-001..RN-045}}{RN-001..RN-045}}{Reglas de negocio }}\label{reglas-de-negocio-rn-001..rn-045}

\textbf{Fuente:}
\texorpdfstring{\softcode{docs/domain/business-rules.md}}{docs/domain/business-rules.md}.

Fuente canónica de las reglas. Origen: \textbf{MP} = megaprompt §5
(propuesta del producto); \textbf{DEC-xx} = decisión registrada en
\texorpdfstring{\softcode{docs/research/evidence-register.md}}{docs/research/evidence-register.md}.
Línea base de diseño:
\texorpdfstring{\softtext{2026-10-10.}}{2026-10-10.} Al 2026-10-11
existen implementación SQL/Java y pruebas locales por incremento; el
resultado y alcance de cada ejecución se registran en evidencia
(referencia del repositorio:
\texorpdfstring{\softcode{../testing/evidence.md}}{../testing/evidence.md}).
La presencia de una RN o un TEST en esta tabla no acredita por sí sola
su cumplimiento completo.

\subsubsection{Convención de errores}\label{convenciuxf3n-de-errores}

Errores de dominio levantados por PL/pgSQL con
\texorpdfstring{\softcode{RAISE EXCEPTION ... USING ERRCODE = 'P0001', DETAIL = json}}{RAISE EXCEPTION ... USING ERRCODE = 'P0001', DETAIL = json}
(objeto validado con \texorpdfstring{\softcode{codigo}}{codigo},
\texorpdfstring{\softcode{regla\_id}}{regla\_id} y campos públicos) y
mapeados por la API: \texorpdfstring{\softcode{BV0xx}}{BV0xx} → 422
(regla de dominio), \texorpdfstring{\softcode{BV4xx}}{BV4xx} → 409
\texorpdfstring{\softtext{(conflicto/concurrencia/idempotencia),}}{(conflicto/concurrencia/idempotencia),}
\texorpdfstring{\softcode{BV5xx}}{BV5xx} → 503 (dueño no
\texorpdfstring{\softtext{disponible).}}{disponible).} Errores nativos
relevantes: \texorpdfstring{\softcode{23505}}{23505}
\texorpdfstring{\softtext{unique\_violation}}{unique\_violation} → 409,
\texorpdfstring{\softcode{23503}}{23503}
\texorpdfstring{\softtext{foreign\_key\_violation}}{foreign\_key\_violation}
→ 422, \texorpdfstring{\softcode{23514}}{23514}
\texorpdfstring{\softtext{check\_violation}}{check\_violation} → 422,
\texorpdfstring{\softcode{42501}}{42501}
\texorpdfstring{\softtext{insufficient\_privilege}}{insufficient\_privilege}
→ 403 (sólo pruebas de rol),
\texorpdfstring{\softcode{40001}}{40001}/\texorpdfstring{\softcode{40P01}}{40P01}
→ reintento acotado.

Actores: \textbf{ADM} administrador · \textbf{OPE} operador ·
\textbf{CAJ} caja · \textbf{AUD} auditor · \textbf{SIS} sistema
\texorpdfstring{\softtext{(comando/trigger).}}{(comando/trigger).}

\begin{itemize}
\item
  \textbf{ID:} RN-001. \textbf{Regla:} Aislamiento por empresa: un
  usuario sólo accede a empresas y sedes autorizadas; referencias
  cruzadas entre empresas se rechazan también en BD. \textbf{Razón:}
  Evitar fuga y mezcla de datos. \textbf{Actor:} Todos.
  \textbf{Precondición:} Sesión válida y
  \texorpdfstring{\softtext{membresía.}}{membresía.} \textbf{Transición
  / efecto:} Ninguna: guarda. \textbf{Datos:}
  \texorpdfstring{\softtext{membresia\_usuario,}}{membresia\_usuario,}
  todas las FK \texorpdfstring{\softtext{compuestas.}}{compuestas.}
  \textbf{Error:} 403 / \texorpdfstring{\softcode{23503}}{23503}.
  \textbf{Implementación prevista:} Spring Security + FK compuestas
  (E,\ldots). \textbf{Prueba:} TEST-14, TEST-15. \textbf{Origen:} MP.
\item
  \textbf{ID:} RN-002. \textbf{Regla:} Cada animal conserva identidad
  estable; identificaciones oficiales y UUID interno son campos
  distintos. \textbf{Razón:} Recompra e historial. \textbf{Actor:} OPE.
  \textbf{Precondición:} -. \textbf{Transición / efecto:} Alta de
  \texorpdfstring{\softcode{animal}}{animal} independiente de sus
  códigos. \textbf{Datos:} animal,
  \texorpdfstring{\softtext{animal\_identificacion.}}{animal\_identificacion.}
  \textbf{Error:} -. \textbf{Implementación prevista:} Tablas separadas.
  \textbf{Prueba:} TEST-05. \textbf{Origen:} MP.
\item
  \textbf{ID:} RN-003. \textbf{Regla:} Identidad maestra e
  identificadores se registran en A; unicidad por empresa verificada
  allí. \textbf{Razón:} Evitar dos UUID para el mismo arete.
  \textbf{Actor:} OPE en A. \textbf{Precondición:} Nodo = A.
  \textbf{Transición / efecto:}
  \texorpdfstring{\softtext{RegistrarIdentidadAnimal}}{RegistrarIdentidadAnimal}
  inserta en A. \textbf{Datos:}
  \texorpdfstring{\softtext{animal\_identificacion}}{animal\_identificacion}
  \texorpdfstring{\softtext{UK(E,sistema,codigo\_normalizado).}}{UK(E,sistema,codigo\_normalizado).}
  \textbf{Error:} \texorpdfstring{\softcode{23505}}{23505}→409
  \texorpdfstring{\softcode{BV401}}{BV401}. \textbf{Implementación
  prevista:} Comando en A + UNIQUE. \textbf{Prueba:} TEST-13.
  \textbf{Origen:} MP.
\item
  \textbf{ID:} RN-004. \textbf{Regla:} Sede operadora del animal
  inmutable en el MVP; propietario, custodio, ubicación y nodo son
  conceptos distintos. \textbf{Razón:} Escritor único por agregado.
  \textbf{Actor:} SIS. \textbf{Precondición:} -. \textbf{Transición /
  efecto:} UPDATE de
  \texorpdfstring{\softcode{sede\_operadora\_id}}{sede\_operadora\_id}
  rechazado. \textbf{Datos:} animal. \textbf{Error:}
  \texorpdfstring{\softcode{BV010}}{BV010}. \textbf{Implementación
  prevista:} Trigger de columna inmutable. \textbf{Prueba:} TEST-17.
  \textbf{Origen:} MP.
\item
  \textbf{ID:} RN-005. \textbf{Regla:} Animal y posición son distintos;
  recompra permitida tras cerrar el ciclo previo; no doble venta del
  ciclo vigente. \textbf{Razón:} Ciclos de costo
  \texorpdfstring{\softtext{independientes.}}{independientes.}
  \textbf{Actor:} OPE. \textbf{Precondición:} Ciclo previo
  \texorpdfstring{\softtext{CERRADA\_VENTA.}}{CERRADA\_VENTA.}
  \textbf{Transición / efecto:} Nueva posición
  \texorpdfstring{\softtext{numero\_ciclo+1.}}{numero\_ciclo+1.}
  \textbf{Datos:}
  \texorpdfstring{\softtext{posicion\_inventario.}}{posicion\_inventario.}
  \textbf{Error:} \texorpdfstring{\softcode{BV011}}{BV011}.
  \textbf{Implementación prevista:}
  \texorpdfstring{\softtext{ConfirmarCompra.}}{ConfirmarCompra.}
  \textbf{Prueba:} TEST-05. \textbf{Origen:} MP.
\item
  \textbf{ID:} RN-006. \textbf{Regla:} Sólo una posición activa por
  animal en la empresa; toda posición deriva de una adquisición
  \texorpdfstring{\softtext{registrada.}}{registrada.} \textbf{Razón:}
  Inventario sin saldos manuales. \textbf{Actor:} SIS.
  \textbf{Precondición:} -. \textbf{Transición / efecto:} -.
  \textbf{Datos:}
  \texorpdfstring{\softtext{posicion\_inventario}}{posicion\_inventario}
  UNIQUE parcial
  \texorpdfstring{\softtext{(E,animal\_id)}}{(E,animal\_id)} WHERE
  estado activo. \textbf{Error:}
  \texorpdfstring{\softcode{23505}}{23505}→409. \textbf{Implementación
  prevista:} Índice único parcial; sin INSERT directo para app.
  \textbf{Prueba:} TEST-05. \textbf{Origen:} MP.
\item
  \textbf{ID:} RN-007. \textbf{Regla:} Un lote agrupa
  \texorpdfstring{\softtext{animales/posiciones}}{animales/posiciones}
  con historia; una compra alimenta varias ventas y una venta usa
  posiciones de varias compras. \textbf{Razón:} Realidad comercial.
  \textbf{Actor:} OPE. \textbf{Precondición:} -. \textbf{Transición /
  efecto:} -. \textbf{Datos:}
  \texorpdfstring{\softtext{lote\_operacion,}}{lote\_operacion,}
  \texorpdfstring{\softtext{seleccion\_*.}}{seleccion\_*.}
  \textbf{Error:} -. \textbf{Implementación prevista:} Modelo N:M.
  \textbf{Prueba:} TEST-02. \textbf{Origen:} MP.
\item
  \textbf{ID:} RN-008. \textbf{Regla:} La selección de un
  \texorpdfstring{\softtext{animal/posición}}{animal/posición} es única
  entre lotes de una misma operación; composición confirmada
  \texorpdfstring{\softtext{histórica.}}{histórica.} \textbf{Razón:}
  Evitar duplicar un animal en un trato. \textbf{Actor:} OPE.
  \textbf{Precondición:} Operación BORRADOR para editar.
  \textbf{Transición / efecto:} -. \textbf{Datos:}
  \texorpdfstring{\softtext{seleccion\_compra}}{seleccion\_compra}
  UK(E,S,op,animal),
  \texorpdfstring{\softtext{seleccion\_venta}}{seleccion\_venta}
  \texorpdfstring{\softtext{UK(E,S,op,posicion).}}{UK(E,S,op,posicion).}
  \textbf{Error:} \texorpdfstring{\softcode{23505}}{23505}→409.
  \textbf{Implementación prevista:} UNIQUE + trigger inmutable tras
  confirmar. \textbf{Prueba:} TEST-17. \textbf{Origen:} MP.
\item
  \textbf{ID:} RN-009. \textbf{Regla:} Toda operación tiene empresa,
  sede, tipo, contraparte, moneda, versión y estado; no admite animales
  de otra sede escritora. \textbf{Razón:} Propiedad de escritura.
  \textbf{Actor:} OPE. \textbf{Precondición:} -. \textbf{Transición /
  efecto:} -. \textbf{Datos:} operacion NOT NULL; FK (E,S,animal) contra
  \texorpdfstring{\softtext{UK(E,sede\_operadora\_id,id).}}{UK(E,sede\_operadora\_id,id).}
  \textbf{Error:} \texorpdfstring{\softcode{23503}}{23503}.
  \textbf{Implementación prevista:} FK compuesta. \textbf{Prueba:}
  TEST-14. \textbf{Origen:} MP.
\item
  \textbf{ID:} RN-010. \textbf{Regla:} Compra y venta por cuenta propia
  son los únicos modos; un tercero puede ser comprador y vendedor.
  \textbf{Razón:} Alcance MVP. \textbf{Actor:} OPE.
  \textbf{Precondición:} -. \textbf{Transición / efecto:} -.
  \textbf{Datos:}
  \texorpdfstring{\softtext{operacion.tipo}}{operacion.tipo} CHECK IN
  (`COMPRA',`VENTA'); parte sin rol fijo. \textbf{Error:}
  \texorpdfstring{\softcode{23514}}{23514}. \textbf{Implementación
  prevista:} CHECK + sin ISA disjunta. \textbf{Prueba:} TEST-30.
  \textbf{Origen:} MP.
\item
  \textbf{ID:} RN-011. \textbf{Regla:} Condiciones de precio, pesaje,
  ajustes y redondeo se versionan; una versión nueva no recalcula
  \texorpdfstring{\softtext{confirmadas.}}{confirmadas.} \textbf{Razón:}
  Snapshot \texorpdfstring{\softtext{histórico.}}{histórico.}
  \textbf{Actor:} OPE/ADM. \textbf{Precondición:} -. \textbf{Transición
  / efecto:} Nueva
  \texorpdfstring{\softtext{contrato\_version}}{contrato\_version}
  numero+1. \textbf{Datos:}
  \texorpdfstring{\softtext{contrato\_version,}}{contrato\_version,}
  \texorpdfstring{\softtext{condicion\_lote.}}{condicion\_lote.}
  \textbf{Error:} \texorpdfstring{\softcode{BV012}}{BV012} si se edita
  aceptada. \textbf{Implementación prevista:} Trigger inmutable en
  versión aceptada. \textbf{Prueba:} TEST-17. \textbf{Origen:} MP.
\item
  \textbf{ID:} RN-012. \textbf{Regla:} Pesos en kg con precisión
  declarada; tara ≤ bruto y neto = bruto − tara. \textbf{Razón:}
  Integridad de medición. \textbf{Actor:} OPE. \textbf{Precondición:} -.
  \textbf{Transición / efecto:} -. \textbf{Datos:} pesaje CHECK.
  \textbf{Error:} \texorpdfstring{\softcode{23514}}{23514}.
  \textbf{Implementación prevista:} CHECK de fila + no NaN/∞.
  \textbf{Prueba:} TEST-18. \textbf{Origen:} MP.
\item
  \textbf{ID:} RN-013. \textbf{Regla:} Un pesaje de lote conserva la
  lista exacta de animales; no se inventan pesos
  \texorpdfstring{\softtext{individuales.}}{individuales.}
  \textbf{Razón:} Evidencia real. \textbf{Actor:} OPE.
  \textbf{Precondición:} -. \textbf{Transición / efecto:}
  \texorpdfstring{\softtext{AprobarPesaje}}{AprobarPesaje} congela
  \texorpdfstring{\softtext{composición.}}{composición.} \textbf{Datos:}
  \texorpdfstring{\softtext{pesaje\_animal.peso\_individual\_kg}}{pesaje\_animal.peso\_individual\_kg}
  NULL si no medido. \textbf{Error:}
  \texorpdfstring{\softcode{BV013}}{BV013}. \textbf{Implementación
  prevista:} Comando + trigger inmutable. \textbf{Prueba:} TEST-18.
  \textbf{Origen:} MP.
\item
  \textbf{ID:} RN-014. \textbf{Regla:} Sólo se comparan pesajes del
  mismo conjunto; si difiere la composición → ``no comparable''.
  \textbf{Razón:} Evitar merma falsa. \textbf{Actor:} SIS.
  \textbf{Precondición:} -. \textbf{Transición / efecto:} Comparación
  devuelve estado
  \texorpdfstring{\softtext{NO\_COMPARABLE.}}{NO\_COMPARABLE.}
  \textbf{Datos:}
  \texorpdfstring{\softtext{pesaje\_animal.}}{pesaje\_animal.}
  \textbf{Error:} -. \textbf{Implementación prevista:} Función de
  comparación por igualdad de conjuntos. \textbf{Prueba:} TEST-19.
  \textbf{Origen:} MP.
\item
  \textbf{ID:} RN-015. \textbf{Regla:} Variación de peso y deducción
  contractual son independientes; no descontar dos veces.
  \textbf{Razón:} Precio justo. \textbf{Actor:} SIS.
  \textbf{Precondición:} -. \textbf{Transición / efecto:} Liquidación
  usa sólo base y deducciones pactadas. \textbf{Datos:}
  \texorpdfstring{\softtext{condicion\_lote.}}{condicion\_lote.}
  \textbf{Error:} -. \textbf{Implementación prevista:} Función canónica
  de \texorpdfstring{\softtext{liquidación.}}{liquidación.}
  \textbf{Prueba:} TEST-01. \textbf{Origen:} MP.
\item
  \textbf{ID:} RN-016. \textbf{Regla:} OCR/CSV generan propuestas; un
  responsable confirma antes de usarlas. \textbf{Razón:} Control humano.
  \textbf{Actor:} OPE. \textbf{Precondición:} -. \textbf{Transición /
  efecto:} Estado PROPUESTA →
  \texorpdfstring{\softtext{CONFIRMADA.}}{CONFIRMADA.} \textbf{Datos:}
  \texorpdfstring{\softtext{pesaje.estado,}}{pesaje.estado,}
  \texorpdfstring{\softtext{infra.importacion\_fuente.}}{infra.importacion\_fuente.}
  \textbf{Error:} \texorpdfstring{\softcode{BV014}}{BV014}.
  \textbf{Implementación prevista:} Estado separado. \textbf{Prueba:}
  TEST-18. \textbf{Origen:} MP.
\item
  \textbf{ID:} RN-017. \textbf{Regla:} Un borrador puede estar
  incompleto; confirmar exige partidas, contraparte, condiciones
  válidas, pesos cuando procedan y
  \texorpdfstring{\softtext{autorización.}}{autorización.}
  \textbf{Razón:} Calidad del cierre. \textbf{Actor:} OPE.
  \textbf{Precondición:}
  \texorpdfstring{\softtext{BORRADOR/ACORDADA.}}{BORRADOR/ACORDADA.}
  \textbf{Transición / efecto:} →
  \texorpdfstring{\softtext{CONFIRMADA.}}{CONFIRMADA.} \textbf{Datos:}
  operacion,
  \texorpdfstring{\softtext{condicion\_lote,}}{condicion\_lote,} pesaje.
  \textbf{Error:} \texorpdfstring{\softcode{BV015}}{BV015}.
  \textbf{Implementación prevista:} Validaciones en
  \texorpdfstring{\softtext{ConfirmarCompra/Venta.}}{ConfirmarCompra/Venta.}
  \textbf{Prueba:} TEST-01. \textbf{Origen:} MP.
\item
  \textbf{ID:} RN-018. \textbf{Regla:} Reserva atómica en el nodo
  propietario, bloqueando posiciones en orden estable; una reserva
  activa por posición. \textbf{Razón:} Evitar doble venta.
  \textbf{Actor:} OPE. \textbf{Precondición:} Posición DISPONIBLE,
  operación VENTA no
  \texorpdfstring{\softtext{confirmada.}}{confirmada.}
  \textbf{Transición / efecto:} Crea reserva ACTIVA. \textbf{Datos:}
  \texorpdfstring{\softtext{reserva\_venta}}{reserva\_venta} UNIQUE
  parcial
  \texorpdfstring{\softtext{(E,S,posicion\_id)}}{(E,S,posicion\_id)}
  WHERE estado=`ACTIVA'. \textbf{Error:}
  \texorpdfstring{\softcode{23505}}{23505}/\texorpdfstring{\softcode{BV402}}{BV402}→409.
  \textbf{Implementación prevista:}
  \texorpdfstring{\softcode{SELECT … FOR UPDATE ORDER BY id}}{SELECT … FOR UPDATE ORDER BY id}
  + índice parcial. \textbf{Prueba:} TEST-10. \textbf{Origen:} MP,
  DEC-11.
\item
  \textbf{ID:} RN-019. \textbf{Regla:} Reservas con vencimiento
  explícito; confirmar tras vencer revalida
  \texorpdfstring{\softtext{disponibilidad.}}{disponibilidad.}
  \textbf{Razón:} No retener inventario
  \texorpdfstring{\softtext{indefinidamente.}}{indefinidamente.}
  \textbf{Actor:} SIS. \textbf{Precondición:} -. \textbf{Transición /
  efecto:} \texorpdfstring{\softtext{ExpirarReserva:}}{ExpirarReserva:}
  \texorpdfstring{\softtext{ACTIVA→EXPIRADA.}}{ACTIVA→EXPIRADA.}
  \textbf{Datos:}
  \texorpdfstring{\softtext{reserva\_venta.vence\_en.}}{reserva\_venta.vence\_en.}
  \textbf{Error:} \texorpdfstring{\softcode{BV403}}{BV403}.
  \textbf{Implementación prevista:} Comando (no índice con now()).
  \textbf{Prueba:} TEST-20. \textbf{Origen:} MP.
\item
  \textbf{ID:} RN-020. \textbf{Regla:} Venta, entrega física y cobro son
  ejes separados. \textbf{Razón:} Un solo evento
  \texorpdfstring{\softtext{económico.}}{económico.} \textbf{Actor:}
  SIS. \textbf{Precondición:} -. \textbf{Transición / efecto:} -.
  \textbf{Datos:}
  \texorpdfstring{\softtext{operacion.estado,}}{operacion.estado,}
  \texorpdfstring{\softtext{movimiento\_inventario,}}{movimiento\_inventario,}
  \texorpdfstring{\softtext{aplicacion\_pago.}}{aplicacion\_pago.}
  \textbf{Error:} -. \textbf{Implementación prevista:} Estados
  \texorpdfstring{\softtext{independientes.}}{independientes.}
  \textbf{Prueba:} TEST-32. \textbf{Origen:} MP.
\item
  \textbf{ID:} RN-021. \textbf{Regla:}
  \texorpdfstring{\softtext{ConfirmarCompra}}{ConfirmarCompra} crea
  posiciones, valoración inicial, liquidación principal confirmada y
  obligación en una transacción local. \textbf{Razón:}
  \texorpdfstring{\softtext{Atomicidad.}}{Atomicidad.} \textbf{Actor:}
  OPE. \textbf{Precondición:} RN-017. \textbf{Transición / efecto:}
  Operación → CONFIRMADA; posiciones
  \texorpdfstring{\softtext{PENDIENTE\_RECEPCION/DISPONIBLE.}}{PENDIENTE\_RECEPCION/DISPONIBLE.}
  \textbf{Datos:}
  \texorpdfstring{\softtext{posicion\_inventario,}}{posicion\_inventario,}
  liquidacion,
  \texorpdfstring{\softtext{concepto\_liquidacion,}}{concepto\_liquidacion,}
  \texorpdfstring{\softtext{asignacion\_costo,}}{asignacion\_costo,}
  \texorpdfstring{\softtext{asiento\_valoracion,}}{asiento\_valoracion,}
  \texorpdfstring{\softtext{auditoria\_evento,}}{auditoria\_evento,}
  \texorpdfstring{\softtext{evento\_outbox.}}{evento\_outbox.}
  \textbf{Error:} \texorpdfstring{\softcode{BV015}}{BV015},
  \texorpdfstring{\softcode{BV404}}{BV404}. \textbf{Implementación
  prevista:} Función
  \texorpdfstring{\softcode{cmd\_confirmar\_compra}}{cmd\_confirmar\_compra}.
  \textbf{Prueba:} TEST-01, TEST-02. \textbf{Origen:} MP, DEC-02.
\item
  \textbf{ID:} RN-022. \textbf{Regla:}
  \texorpdfstring{\softtext{ConfirmarVenta}}{ConfirmarVenta} compromete
  posiciones, consume reservas, registra costo de venta, liquidación y
  obligación; la entrega posterior no duplica. \textbf{Razón:}
  Atomicidad y reconocimiento único. \textbf{Actor:} OPE.
  \textbf{Precondición:} Reservas ACTIVAS vigentes. \textbf{Transición /
  efecto:} Posiciones →
  \texorpdfstring{\softtext{VENDIDA\_PENDIENTE\_ENTREGA.}}{VENDIDA\_PENDIENTE\_ENTREGA.}
  \textbf{Datos:} idem +
  \texorpdfstring{\softtext{reserva\_venta.}}{reserva\_venta.}
  \textbf{Error:} \texorpdfstring{\softcode{BV403}}{BV403},
  \texorpdfstring{\softcode{BV402}}{BV402}. \textbf{Implementación
  prevista:} Función
  \texorpdfstring{\softcode{cmd\_confirmar\_venta}}{cmd\_confirmar\_venta}.
  \textbf{Prueba:} TEST-10, TEST-32. \textbf{Origen:} MP.
\item
  \textbf{ID:} RN-023. \textbf{Regla:} Liquidación confirmada inmutable;
  \texorpdfstring{\softtext{ajustes/devoluciones/reversos}}{ajustes/devoluciones/reversos}
  conservan referencia, motivo, actor y fecha. \textbf{Razón:}
  \texorpdfstring{\softtext{Auditoría.}}{Auditoría.} \textbf{Actor:}
  SIS. \textbf{Precondición:} -. \textbf{Transición / efecto:}
  \texorpdfstring{\softtext{UPDATE/DELETE}}{UPDATE/DELETE} rechazado.
  \textbf{Datos:} liquidacion,
  \texorpdfstring{\softtext{concepto\_liquidacion.}}{concepto\_liquidacion.}
  \textbf{Error:} \texorpdfstring{\softcode{BV016}}{BV016}.
  \textbf{Implementación prevista:} Trigger +
  \texorpdfstring{\softtext{privilegios.}}{privilegios.}
  \textbf{Prueba:} TEST-17. \textbf{Origen:} MP.
\item
  \textbf{ID:} RN-024. \textbf{Regla:} Tarifa nueva, tercero renombrado
  o documento posterior no alteran el snapshot de una
  \texorpdfstring{\softtext{liquidación.}}{liquidación.} \textbf{Razón:}
  Historia. \textbf{Actor:} SIS. \textbf{Precondición:} -.
  \textbf{Transición / efecto:} -. \textbf{Datos:}
  \texorpdfstring{\softtext{liquidacion.snapshot\_*.}}{liquidacion.snapshot\_*.}
  \textbf{Error:} -. \textbf{Implementación prevista:} Columnas snapshot
  sin FK en cascada de UPDATE. \textbf{Prueba:} TEST-17.
  \textbf{Origen:} MP.
\item
  \textbf{ID:} RN-025. \textbf{Regla:} Cada gasto: costo de adquisición
  asignable, gasto directo de venta u otro; no se suma dos veces.
  \textbf{Razón:} Margen correcto. \textbf{Actor:} OPE.
  \textbf{Precondición:} -. \textbf{Transición / efecto:} -.
  \textbf{Datos:}
  \texorpdfstring{\softtext{gasto\_operacion.clasificacion.}}{gasto\_operacion.clasificacion.}
  \textbf{Error:} \texorpdfstring{\softcode{BV017}}{BV017}.
  \textbf{Implementación prevista:} CHECK + validación de duplicado por
  referencia de origen. \textbf{Prueba:} TEST-01. \textbf{Origen:} MP,
  DEC-08.
\item
  \textbf{ID:} RN-026. \textbf{Regla:} Asignaciones de un gasto suman
  exactamente el total; método y residuo
  \texorpdfstring{\softtext{deterministas.}}{deterministas.}
  \textbf{Razón:} Cuadre al centavo. \textbf{Actor:} SIS.
  \textbf{Precondición:} Gasto aprobado. \textbf{Transición / efecto:}
  \texorpdfstring{\softtext{DistribuirCosto.}}{DistribuirCosto.}
  \textbf{Datos:}
  \texorpdfstring{\softtext{asignacion\_costo.}}{asignacion\_costo.}
  \textbf{Error:} \texorpdfstring{\softcode{BV018}}{BV018}.
  \textbf{Implementación prevista:} Función de reparto con residuo por
  orden de id. \textbf{Prueba:} TEST-03. \textbf{Origen:} MP.
\item
  \textbf{ID:} RN-027. \textbf{Regla:} Costo tardío modifica el margen
  mediante ajuste trazable, sin editar liquidación ni periodo cerrado.
  \textbf{Razón:} Historia. \textbf{Actor:} OPE. \textbf{Precondición:}
  Periodo destino ABIERTO. \textbf{Transición / efecto:}
  \texorpdfstring{\softtext{RegistrarCostoTardio}}{RegistrarCostoTardio}
  crea asignaciones y asientos. \textbf{Datos:}
  \texorpdfstring{\softtext{asignacion\_costo,}}{asignacion\_costo,}
  \texorpdfstring{\softtext{asiento\_valoracion.}}{asiento\_valoracion.}
  \textbf{Error:} \texorpdfstring{\softcode{BV019}}{BV019}.
  \textbf{Implementación prevista:} Comando. \textbf{Prueba:} TEST-06.
  \textbf{Origen:} MP.
\item
  \textbf{ID:} RN-028. \textbf{Regla:} El pago existe separado de su
  aplicación; N:M
  \texorpdfstring{\softtext{pago-obligación.}}{pago-obligación.}
  \textbf{Razón:} Anticipos y pagos parciales. \textbf{Actor:} CAJ.
  \textbf{Precondición:} -. \textbf{Transición / efecto:} -.
  \textbf{Datos:} pago,
  \texorpdfstring{\softtext{aplicacion\_pago.}}{aplicacion\_pago.}
  \textbf{Error:} -. \textbf{Implementación prevista:} Tabla
  \texorpdfstring{\softtext{asociativa.}}{asociativa.} \textbf{Prueba:}
  TEST-04. \textbf{Origen:} MP.
\item
  \textbf{ID:} RN-029. \textbf{Regla:} Aplicaciones no exceden
  disponible del pago ni saldo de la
  \texorpdfstring{\softtext{obligación.}}{obligación.} \textbf{Razón:}
  Sin \texorpdfstring{\softtext{sobreaplicación.}}{sobreaplicación.}
  \textbf{Actor:} CAJ. \textbf{Precondición:} Locks de pago y
  \texorpdfstring{\softtext{liquidación.}}{liquidación.}
  \textbf{Transición / efecto:}
  \texorpdfstring{\softtext{AplicarPago.}}{AplicarPago.} \textbf{Datos:}
  \texorpdfstring{\softtext{aplicacion\_pago.}}{aplicacion\_pago.}
  \textbf{Error:} \texorpdfstring{\softcode{BV020}}{BV020}.
  \textbf{Implementación prevista:} Comando con FOR UPDATE en orden
  estable. \textbf{Prueba:} TEST-04, TEST-11. \textbf{Origen:} MP.
\item
  \textbf{ID:} RN-030. \textbf{Regla:} Pago y aplicación: misma empresa,
  sede, contraparte, moneda y dirección
  \texorpdfstring{\softtext{compatible.}}{compatible.} \textbf{Razón:}
  \texorpdfstring{\softtext{Coherencia.}}{Coherencia.} \textbf{Actor:}
  CAJ. \textbf{Precondición:} -. \textbf{Transición / efecto:} -.
  \textbf{Datos:} pago.tipo vs
  \texorpdfstring{\softtext{operacion.tipo.}}{operacion.tipo.}
  \textbf{Error:} \texorpdfstring{\softcode{BV021}}{BV021}.
  \textbf{Implementación prevista:} FK compuestas +
  \texorpdfstring{\softtext{validación.}}{validación.} \textbf{Prueba:}
  TEST-11. \textbf{Origen:} MP, DEC-06.
\item
  \textbf{ID:} RN-031. \textbf{Regla:} Importe no aplicado queda como
  \texorpdfstring{\softtext{anticipo/saldo}}{anticipo/saldo} a favor; no
  es venta adicional. \textbf{Razón:} Separar caja y margen.
  \textbf{Actor:} SIS. \textbf{Precondición:} -. \textbf{Transición /
  efecto:} -. \textbf{Datos:} saldo derivado del pago. \textbf{Error:}
  -. \textbf{Implementación prevista:} Vista de saldos. \textbf{Prueba:}
  TEST-04. \textbf{Origen:} MP.
\item
  \textbf{ID:} RN-032. \textbf{Regla:} Misma clave + mismo payload
  devuelve el resultado previo; misma clave + payload distinto =
  conflicto. \textbf{Razón:}
  \texorpdfstring{\softtext{Idempotencia.}}{Idempotencia.}
  \textbf{Actor:} SIS. \textbf{Precondición:} -. \textbf{Transición /
  efecto:} Reclamo de clave. \textbf{Datos:}
  \texorpdfstring{\softtext{clave\_idempotencia.}}{clave\_idempotencia.}
  \textbf{Error:} \texorpdfstring{\softcode{BV405}}{BV405}→409.
  \textbf{Implementación prevista:} Tabla + hash SHA-256 del payload
  canónico. \textbf{Prueba:} TEST-12. \textbf{Origen:} MP.
\item
  \textbf{ID:} RN-033. \textbf{Regla:} Reversar conserva original y crea
  contramovimiento; se verifica efecto en aplicaciones
  \texorpdfstring{\softtext{relacionadas.}}{relacionadas.}
  \textbf{Razón:} Historia. \textbf{Actor:} CAJ/ADM.
  \textbf{Precondición:} Original no revertido. \textbf{Transición /
  efecto:}
  \texorpdfstring{\softtext{RevertirAplicacion.}}{RevertirAplicacion.}
  \textbf{Datos:}
  \texorpdfstring{\softtext{aplicacion\_pago.reverso\_de\_id}}{aplicacion\_pago.reverso\_de\_id}
  UNIQUE. \textbf{Error:} \texorpdfstring{\softcode{BV022}}{BV022}.
  \textbf{Implementación prevista:} Comando. \textbf{Prueba:} TEST-07.
  \textbf{Origen:} MP, DEC-09.
\item
  \textbf{ID:} RN-034. \textbf{Regla:} Importes con NUMERIC y
  \texorpdfstring{\softtext{BigDecimal;}}{BigDecimal;} nunca
  \texorpdfstring{\softtext{float/double.}}{float/double.}
  \textbf{Razón:} Exactitud. \textbf{Actor:} -. \textbf{Precondición:}
  -. \textbf{Transición / efecto:} -. \textbf{Datos:} todos los
  importes. \textbf{Error:} -. \textbf{Implementación prevista:} Tipos;
  prueba de lint. \textbf{Prueba:} TEST-03. \textbf{Origen:} MP.
\item
  \textbf{ID:} RN-035. \textbf{Regla:} Margen = ventas reconocidas −
  costo de posiciones vendidas − gastos directos; cobros no afectan
  margen. \textbf{Razón:} Separar devengo y caja. \textbf{Actor:} SIS.
  \textbf{Precondición:} -. \textbf{Transición / efecto:} -.
  \textbf{Datos:} vistas de reporte. \textbf{Error:} -.
  \textbf{Implementación prevista:} Vista
  \texorpdfstring{\softcode{v\_margen\_operacion}}{v\_margen\_operacion}.
  \textbf{Prueba:} TEST-01. \textbf{Origen:} MP.
\item
  \textbf{ID:} RN-036. \textbf{Regla:} Si faltan costos o revisiones,
  margen PROVISIONAL; nunca ``utilidad neta''. \textbf{Razón:}
  Honestidad del reporte. \textbf{Actor:} SIS. \textbf{Precondición:} -.
  \textbf{Transición / efecto:} Etiqueta. \textbf{Datos:} vista de
  margen. \textbf{Error:} -. \textbf{Implementación prevista:} Columna
  derivada. \textbf{Prueba:} TEST-01. \textbf{Origen:} MP.
\item
  \textbf{ID:} RN-037. \textbf{Regla:} Evidencias con versión, hash,
  tamaño, tipo y control de acceso; hash detecta cambios, no prueba
  autoría. \textbf{Razón:} Integridad de archivos. \textbf{Actor:} OPE.
  \textbf{Precondición:} -. \textbf{Transición / efecto:} -.
  \textbf{Datos:} documento. \textbf{Error:}
  \texorpdfstring{\softcode{BV023}}{BV023}. \textbf{Implementación
  prevista:} SHA-256 + almacenamiento privado. \textbf{Prueba:} TEST-34.
  \textbf{Origen:} MP.
\item
  \textbf{ID:} RN-038. \textbf{Regla:} Cancelación, pérdida, muerte,
  devolución y venta son motivos distintos con evidencia y autorización;
  no borrar animales. \textbf{Razón:}
  \texorpdfstring{\softtext{Trazabilidad.}}{Trazabilidad.}
  \textbf{Actor:} ADM. \textbf{Precondición:} -. \textbf{Transición /
  efecto:}
  \texorpdfstring{\softtext{RegistrarBajaInventario}}{RegistrarBajaInventario}
  → BAJA. \textbf{Datos:}
  \texorpdfstring{\softtext{posicion\_inventario.motivo\_baja,}}{posicion\_inventario.motivo\_baja,}
  \texorpdfstring{\softtext{animal.situacion.}}{animal.situacion.}
  \textbf{Error:} \texorpdfstring{\softcode{BV024}}{BV024}.
  \textbf{Implementación prevista:} Comando; sin DELETE.
  \textbf{Prueba:} TEST-08, TEST-33. \textbf{Origen:} MP, DEC-03/04.
\item
  \textbf{ID:} RN-039. \textbf{Regla:} Periodo cerrado impide ediciones
  retroactivas; reapertura con permiso y bitácora. \textbf{Razón:}
  Cierres \texorpdfstring{\softtext{confiables.}}{confiables.}
  \textbf{Actor:} ADM. \textbf{Precondición:} -. \textbf{Transición /
  efecto:}
  \texorpdfstring{\softtext{Cerrar/ReabrirPeriodo.}}{Cerrar/ReabrirPeriodo.}
  \textbf{Datos:}
  \texorpdfstring{\softtext{periodo\_cierre,}}{periodo\_cierre,}
  \texorpdfstring{\softtext{cierre\_periodo\_version.}}{cierre\_periodo\_version.}
  \textbf{Error:} \texorpdfstring{\softcode{BV025}}{BV025}.
  \textbf{Implementación prevista:} Lock de fila de periodo.
  \textbf{Prueba:} TEST-09. \textbf{Origen:} MP.
\item
  \textbf{ID:} RN-040. \textbf{Regla:} Una réplica retrasada no autoriza
  operaciones críticas; decide el escritor
  \texorpdfstring{\softtext{propietario.}}{propietario.} \textbf{Razón:}
  \texorpdfstring{\softtext{Consistencia.}}{Consistencia.}
  \textbf{Actor:} SIS. \textbf{Precondición:} -. \textbf{Transición /
  efecto:} -. \textbf{Datos:} -. \textbf{Error:}
  \texorpdfstring{\softcode{BV501}}{BV501}→503. \textbf{Implementación
  prevista:}
  \texorpdfstring{\softtext{OwnerDatabaseRouter.}}{OwnerDatabaseRouter.}
  \textbf{Prueba:} TEST-22. \textbf{Origen:} MP.
\item
  \textbf{ID:} RN-041. \textbf{Regla:} Si el dueño no está disponible la
  respuesta puede quedar pendiente de comprobar; no hay cambio
  automático de dueño. \textbf{Razón:} Evitar
  \texorpdfstring{\softtext{split-brain.}}{split-brain.} \textbf{Actor:}
  SIS. \textbf{Precondición:} -. \textbf{Transición / efecto:} -.
  \textbf{Datos:} -. \textbf{Error:}
  \texorpdfstring{\softcode{BV501}}{BV501}→503. \textbf{Implementación
  prevista:} Router sin failover. \textbf{Prueba:} TEST-22.
  \textbf{Origen:} MP.
\item
  \textbf{ID:} RN-042. \textbf{Regla:} Dato del proveedor FDW no es
  histórico por sí mismo: se guarda snapshot aceptado
  \texorpdfstring{\softtext{localmente.}}{localmente.} \textbf{Razón:}
  \texorpdfstring{\softtext{Reproducibilidad.}}{Reproducibilidad.}
  \textbf{Actor:} OPE. \textbf{Precondición:} -. \textbf{Transición /
  efecto:}
  \texorpdfstring{\softtext{AceptarCotizacion}}{AceptarCotizacion} crea
  snapshot. \textbf{Datos:}
  \texorpdfstring{\softtext{gasto\_operacion/documento}}{gasto\_operacion/documento}
  snapshot. \textbf{Error:} -. \textbf{Implementación prevista:} Copia
  local con hash. \textbf{Prueba:} TEST-25. \textbf{Origen:} MP.
\item
  \textbf{ID:} RN-043. \textbf{Regla:} Confirmar una operación y su
  auditoría es atómico; evento de integración en outbox en la misma
  \texorpdfstring{\softtext{transacción.}}{transacción.} \textbf{Razón:}
  \texorpdfstring{\softtext{Consistencia.}}{Consistencia.}
  \textbf{Actor:} SIS. \textbf{Precondición:} -. \textbf{Transición /
  efecto:} -. \textbf{Datos:}
  \texorpdfstring{\softtext{auditoria\_evento,}}{auditoria\_evento,}
  \texorpdfstring{\softtext{evento\_outbox.}}{evento\_outbox.}
  \textbf{Error:} -. \textbf{Implementación prevista:} Insert en el
  comando. \textbf{Prueba:} TEST-24. \textbf{Origen:} MP.
\item
  \textbf{ID:} RN-044. \textbf{Regla:} Entrega externa admite
  reintentos; outbox no garantiza exactamente una vez. \textbf{Razón:}
  Honestidad técnica. \textbf{Actor:} SIS. \textbf{Precondición:} -.
  \textbf{Transición / efecto:} -. \textbf{Datos:}
  \texorpdfstring{\softtext{infra.entrega\_outbox,}}{infra.entrega\_outbox,}
  \texorpdfstring{\softtext{mensaje\_inbox.}}{mensaje\_inbox.}
  \textbf{Error:} -. \textbf{Implementación prevista:} Inbox deduplica.
  \textbf{Prueba:} TEST-24. \textbf{Origen:} MP.
\item
  \textbf{ID:} RN-045. \textbf{Regla:} Datos de demostración
  identificados como sintéticos; sin documentos oficiales aparentando
  \texorpdfstring{\softtext{autenticidad.}}{autenticidad.}
  \textbf{Razón:} Ética y legal. \textbf{Actor:} Todos.
  \textbf{Precondición:} -. \textbf{Transición / efecto:} -.
  \textbf{Datos:} seed con prefijo DEMO. \textbf{Error:} -.
  \textbf{Implementación prevista:} Seed marcado. \textbf{Prueba:}
  revisión manual. \textbf{Origen:} MP.
\end{itemize}

\section{Glosario}\label{glosario}

\textbf{Fuente:}
\texorpdfstring{\softcode{docs/domain/glossary.md}}{docs/domain/glossary.md}.

Nombre canónico en código entre
\texorpdfstring{\softcode{backticks}}{backticks}. Datos de ejemplo
\textbf{sintéticos}.

\begin{itemize}
\item
  \textbf{Término:} Empresa cliente
  (\texorpdfstring{\softcode{empresa}}{empresa}). \textbf{Definición:}
  Organización que contrata
  \texorpdfstring{\softtext{BoviCierre;}}{BoviCierre;} todo dato
  pertenece a una. \textbf{Ejemplo:} ``Acopio Demo SA''
  \texorpdfstring{\softtext{(sintética).}}{(sintética).} \textbf{No
  confundir con:} La contraparte con la que
  \texorpdfstring{\softtext{compra/vende.}}{compra/vende.}
\item
  \textbf{Término:} Sede
  (\texorpdfstring{\softcode{sucursal}}{sucursal}). \textbf{Definición:}
  Centro operativo de la empresa; cada operación pertenece a una.
  \textbf{Ejemplo:} Sede A ``Centro'', Sede B ``Norte''. \textbf{No
  confundir con:}
  \texorpdfstring{\softtext{Predio/establecimiento;}}{Predio/establecimiento;}
  nodo de BD.
\item
  \textbf{Término:} Nodo escritor. \textbf{Definición:} Instancia
  PostgreSQL que acepta escrituras de un agregado: A (maestros + sede A)
  o B (sede B). \textbf{Ejemplo:} Una venta de la sede B se escribe sólo
  en B. \textbf{No confundir con:} Sede (una sede \emph{tiene} un nodo
  dueño); Pgpool.
\item
  \textbf{Término:} Contraparte
  (\texorpdfstring{\softcode{parte}}{parte}). \textbf{Definición:}
  Persona u organización con la que se compra o vende; puede ser ambas
  cosas en momentos distintos. \textbf{Ejemplo:} ``Productor Sintético
  01''. \textbf{No confundir con:} Empresa cliente; usuario.
\item
  \textbf{Término:} Establecimiento
  (\texorpdfstring{\softcode{predio}}{predio}). \textbf{Definición:}
  Lugar físico declarado de origen, destino o custodia.
  \textbf{Ejemplo:} ``Rancho Demo El Llano''. \textbf{No confundir con:}
  Sede; UPP.
\item
  \textbf{Término:} UPP. \textbf{Definición:} Unidad de Producción
  Pecuaria: clave de registro oficial de un establecimiento productivo;
  se guarda como texto declarado con evidencia. \textbf{Ejemplo:}
  \texorpdfstring{\softcode{UPP-DEMO-0001}}{UPP-DEMO-0001}
  \texorpdfstring{\softtext{(ficticia).}}{(ficticia).} \textbf{No
  confundir con:} Validación de autenticidad
  \texorpdfstring{\softtext{(BoviCierre}}{(BoviCierre} no la
  \texorpdfstring{\softtext{certifica).}}{certifica).}
\item
  \textbf{Término:} PSG. \textbf{Definición:} Prestador de Servicios
  Ganaderos: registro oficial para quien
  \texorpdfstring{\softtext{compra/acopia/comercializa;}}{compra/acopia/comercializa;}
  se guarda como texto declarado. \textbf{Ejemplo:}
  \texorpdfstring{\softcode{PSG-DEMO-01}}{PSG-DEMO-01}. \textbf{No
  confundir con:} Permiso de
  \texorpdfstring{\softtext{movilización.}}{movilización.}
\item
  \textbf{Término:} Animal (\texorpdfstring{\softcode{animal}}{animal}).
  \textbf{Definición:} Identidad permanente de un bovino dentro de la
  empresa; UUID técnico. \textbf{Ejemplo:}
  \texorpdfstring{\softcode{animal 7b1…}}{animal 7b1…}. \textbf{No
  confundir con:} Posición de
  \texorpdfstring{\softtext{inventario.}}{inventario.}
\item
  \textbf{Término:} Identificación oficial
  (\texorpdfstring{\softcode{animal\_identificacion}}{animal\_identificacion}).
  \textbf{Definición:} Código de
  \texorpdfstring{\softtext{arete/registro}}{arete/registro}
  \texorpdfstring{\softtext{(SINIIGA/SINIDA}}{(SINIIGA/SINIDA}
  {[}F04{]}) asociado a un animal, con sistema, vigencia y evidencia.
  \textbf{Ejemplo:}
  \texorpdfstring{\softcode{SINIIGA DEMO-000123}}{SINIIGA DEMO-000123}.
  \textbf{No confundir con:} UUID; QR interno (no sustituye arete).
\item
  \textbf{Término:} Lote
  (\texorpdfstring{\softcode{lote\_operacion}}{lote\_operacion}).
  \textbf{Definición:} Grupo comercial de
  \texorpdfstring{\softtext{animales/posiciones}}{animales/posiciones}
  dentro de una operación, numerado. \textbf{Ejemplo:} Lote 1 de la
  compra C-0001. \textbf{No confundir con:} Pesaje; corral físico.
\item
  \textbf{Término:} Membresía de lote. \textbf{Definición:} Pertenencia
  de un animal (compra) o posición (venta) a un lote:
  \texorpdfstring{\softcode{seleccion\_compra}}{seleccion\_compra} /
  \texorpdfstring{\softcode{seleccion\_venta}}{seleccion\_venta}.
  \textbf{Ejemplo:} Animal X en lote 2. \textbf{No confundir con:}
  Composición de un pesaje.
\item
  \textbf{Término:} Posición / ciclo de inventario
  (\texorpdfstring{\softcode{posicion\_inventario}}{posicion\_inventario}).
  \textbf{Definición:} Tenencia de un animal entre una adquisición y su
  salida; tiene costo propio. \textbf{Ejemplo:} Ciclo 2 del animal X
  tras recompra. \textbf{No confundir con:} Animal.
\item
  \textbf{Término:} Custodia. \textbf{Definición:} Quién tiene
  físicamente al animal; se registra en
  \texorpdfstring{\softtext{movimientos.}}{movimientos.}
  \textbf{Ejemplo:} Transportista durante el viaje. \textbf{No confundir
  con:} Propiedad comercial; nodo dueño.
\item
  \textbf{Término:} Operación
  (\texorpdfstring{\softcode{operacion}}{operacion}).
  \textbf{Definición:} Compra o venta por cuenta propia con una
  contraparte; tipo inmutable. \textbf{Ejemplo:} Compra C-0001.
  \textbf{No confundir con:} Liquidación; pago.
\item
  \textbf{Término:} Acuerdo / versión
  (\texorpdfstring{\softcode{contrato\_version}}{contrato\_version}).
  \textbf{Definición:} Conjunto numerado de condiciones aceptadas para
  una \texorpdfstring{\softtext{operación.}}{operación.}
  \textbf{Ejemplo:} v2 aceptada. \textbf{No confundir con:} Condición
  por lote.
\item
  \textbf{Término:} Condición por lote
  (\texorpdfstring{\softcode{condicion\_lote}}{condicion\_lote}).
  \textbf{Definición:} Precio, unidad
  \texorpdfstring{\softtext{(kg/cabeza),}}{(kg/cabeza),} pesaje base y
  deducciones pactadas para un lote en una versión. \textbf{Ejemplo:}
  48.0000 MXN/kg, base pesaje de origen. \textbf{No confundir con:}
  Tarifa de catálogo.
\item
  \textbf{Término:} Pesaje (\texorpdfstring{\softcode{pesaje}}{pesaje}).
  \textbf{Definición:} Medición de báscula de un conjunto exacto de
  animales. \textbf{Ejemplo:} Ticket 4471, 5,000.000 kg. \textbf{No
  confundir con:} Peso estimado.
\item
  \textbf{Término:} Bruto / tara / neto. \textbf{Definición:} Neto =
  bruto − tara; si el ticket sólo trae neto: modo
  \texorpdfstring{\softtext{NETO\_REPORTADO.}}{NETO\_REPORTADO.}
  \textbf{Ejemplo:} 12,000 − 7,000 = 5,000 kg. \textbf{No confundir
  con:} -.
\item
  \textbf{Término:} Variación física. \textbf{Definición:} Diferencia
  observada entre dos pesajes compatibles del mismo conjunto.
  \textbf{Ejemplo:} 100 kg (2\%). \textbf{No confundir con:} Deducción
  \texorpdfstring{\softtext{contractual.}}{contractual.}
\item
  \textbf{Término:} Ajuste
  \texorpdfstring{\softtext{contractual.}}{contractual.}
  \textbf{Definición:} Deducción o cargo pactado expresamente en la
  \texorpdfstring{\softtext{condición.}}{condición.} \textbf{Ejemplo:}
  1\% de deducción pactada. \textbf{No confundir con:} Merma
  ``habitual'' supuesta.
\item
  \textbf{Término:} Liquidación comercial
  (\texorpdfstring{\softcode{liquidacion}}{liquidacion}).
  \textbf{Definición:} Snapshot inmutable del cálculo de una operación
  confirmada; es la obligación con la contraparte (DEC-05).
  \textbf{Ejemplo:} 240,000.00 por pagar. \textbf{No confundir con:}
  CFDI; PDF.
\item
  \textbf{Término:} \texorpdfstring{\softtext{Obligación.}}{Obligación.}
  \textbf{Definición:} Saldo por cobrar (venta) o pagar (compra) de una
  liquidación + ajustes −
  \texorpdfstring{\softtext{aplicaciones.}}{aplicaciones.}
  \textbf{Ejemplo:} 244,000 por cobrar. \textbf{No confundir con:} Pago.
\item
  \textbf{Término:} Pago (\texorpdfstring{\softcode{pago}}{pago}).
  \textbf{Definición:} Movimiento de dinero registrado (COBRO o PAGO);
  existe aunque no esté aplicado. \textbf{Ejemplo:} Transferencia de
  50,000. \textbf{No confundir con:}
  \texorpdfstring{\softtext{Aplicación.}}{Aplicación.}
\item
  \textbf{Término:} Aplicación
  (\texorpdfstring{\softcode{aplicacion\_pago}}{aplicacion\_pago}).
  \textbf{Definición:} Asignación de un importe de un pago a una
  \texorpdfstring{\softtext{liquidación.}}{liquidación.}
  \textbf{Ejemplo:} 50,000 a la venta V-0001. \textbf{No confundir con:}
  Pago.
\item
  \textbf{Término:} Reverso. \textbf{Definición:} Contramovimiento que
  anula un registro conservando el original (aplicación, ajuste,
  asiento). \textbf{Ejemplo:} Reverso de aplicación 400. \textbf{No
  confundir con:} Corrección (nueva versión) / reembolso (dinero real).
\item
  \textbf{Término:} \texorpdfstring{\softtext{Corrección.}}{Corrección.}
  \textbf{Definición:} Nuevo registro que sustituye a otro por error de
  captura, enlazado al original. \textbf{Ejemplo:} Pesaje corregido.
  \textbf{No confundir con:} Reverso.
\item
  \textbf{Término:} Reembolso / reintegro. \textbf{Definición:}
  Devolución real de dinero que consume disponibilidad del pago origen.
  \textbf{Ejemplo:} Devolver 100 al comprador. \textbf{No confundir
  con:} Reverso de \texorpdfstring{\softtext{aplicación.}}{aplicación.}
\item
  \textbf{Término:} Costo de
  \texorpdfstring{\softtext{adquisición.}}{adquisición.}
  \textbf{Definición:} Precio de compra + gastos capitalizables
  asignados a cada posición. \textbf{Ejemplo:} 1,100 por posición
  (GOLD-02). \textbf{No confundir con:} Costo de venta.
\item
  \textbf{Término:} Costo de venta. \textbf{Definición:} Costo de las
  posiciones cuya venta se
  \texorpdfstring{\softtext{reconoció.}}{reconoció.} \textbf{Ejemplo:}
  4,400. \textbf{No confundir con:} Gasto directo de venta.
\item
  \textbf{Término:} Gasto directo de venta. \textbf{Definición:} Gasto
  atribuible a una venta que no se capitaliza en
  \texorpdfstring{\softtext{inventario.}}{inventario.} \textbf{Ejemplo:}
  Flete de entrega 200. \textbf{No confundir con:} Costo de
  \texorpdfstring{\softtext{adquisición.}}{adquisición.}
\item
  \textbf{Término:} Margen
  \texorpdfstring{\softtext{provisional.}}{provisional.}
  \textbf{Definición:} Margen con costos o revisiones faltantes; siempre
  \texorpdfstring{\softtext{etiquetado.}}{etiquetado.} \textbf{Ejemplo:}
  ``Provisional: falta flete''. \textbf{No confundir con:} Margen
  cerrado.
\item
  \textbf{Término:} Margen cerrado. \textbf{Definición:} Margen incluido
  en un cierre de periodo publicado (versión
  \texorpdfstring{\softtext{inmutable).}}{inmutable).} \textbf{Ejemplo:}
  Cierre oct v1. \textbf{No confundir con:} Margen acumulado
  \texorpdfstring{\softtext{actualizado.}}{actualizado.}
\end{itemize}

\subsubsection{Sinónimos aceptados}\label{sinuxf3nimos-aceptados}

sede = sucursal · contraparte = tercero = parte · establecimiento =
predio · acuerdo =
\texorpdfstring{\softtext{contrato\_version}}{contrato\_version} · ciclo
= posición.

\subsubsection{Denominaciones oficiales}\label{denominaciones-oficiales}

La
\texorpdfstring{\softtext{NOM-001-SAG/GAN-2015}}{NOM-001-SAG/GAN-2015}
usa \textbf{SINIDA}; en la operación del sistema aparece
\textbf{SINIIGA} {[}F04-F05{]}. Se guarda el sistema declarado en
\texorpdfstring{\softcode{animal\_identificacion.sistema}}{animal\_identificacion.sistema}.
Un establecimiento registrado, el propietario comercial y la sede que
escribe en PostgreSQL son conceptos distintos.

\subsubsection{Unidades y precisión}\label{unidades-y-precisiuxf3n}

kg NUMERIC(14,3) · tarifas NUMERIC(14,4) · importes NUMERIC(18,2) ·
tasas NUMERIC(9,6) · moneda MXN · cantidades enteras positivas ·
redondeo HALF\_UP a 2 decimales en importes.

\section{Diccionario de datos}\label{diccionario-de-datos}

\textbf{Fuente:}
\texorpdfstring{\softcode{docs/database/data-dictionary.md}}{docs/database/data-dictionary.md}.

Convenciones comunes (no se repiten por tabla):

\begin{itemize}
\tightlist
\item
  \texorpdfstring{\softcode{empresa\_id uuid NOT NULL}}{empresa\_id uuid NOT NULL},
  \texorpdfstring{\softcode{sucursal\_id uuid NOT NULL}}{sucursal\_id uuid NOT NULL}
  en tablas locales;
  \texorpdfstring{\softcode{id uuid NOT NULL DEFAULT uuidv7()}}{id uuid NOT NULL DEFAULT uuidv7()}.
\item
  \texorpdfstring{\softcode{nodo\_origen char(1) NOT NULL CHECK (nodo\_origen IN ('A','B'))}}{nodo\_origen char(1) NOT NULL CHECK (nodo\_origen IN ('A','B'))},
  lo deriva el comando; trigger
  \texorpdfstring{\softcode{trg\_guarda\_propietario}}{trg\_guarda\_propietario}
  rechaza \texorpdfstring{\softtext{INSERT/UPDATE}}{INSERT/UPDATE}
  cuando \texorpdfstring{\softcode{nodo\_origen}}{nodo\_origen} ≠
  \texorpdfstring{\softcode{infra.nodo\_config.nodo}}{infra.nodo\_config.nodo}
  (no corre en apply porque
  \texorpdfstring{\softcode{session\_replication\_role = replica}}{session\_replication\_role = replica}).
\item
  Importes \texorpdfstring{\softcode{numeric(18,2)}}{numeric(18,2)}, kg
  \texorpdfstring{\softcode{numeric(14,3)}}{numeric(14,3)}, tarifas
  \texorpdfstring{\softcode{numeric(14,4)}}{numeric(14,4)}, tasas
  \texorpdfstring{\softcode{numeric(9,6)}}{numeric(9,6)}; todos con
  \texorpdfstring{\softcode{CHECK (x NOT IN ('NaN','Infinity','-Infinity'))}}{CHECK (x NOT IN ('NaN','Infinity','-Infinity'))}
  como defensa explícita (comportamiento de NaN/∞ con typmod a verificar
  en TASK-A-009 con prueba negativa).
\item
  Instantes \texorpdfstring{\softcode{timestamptz}}{timestamptz}; fechas
  civiles \texorpdfstring{\softcode{date}}{date}; presentación en
  \texorpdfstring{\softtext{America/Mexico\_City.}}{America/Mexico\_City.}
\item
  Hechos comprometidos:
  \texorpdfstring{\softcode{ON UPDATE RESTRICT ON DELETE RESTRICT}}{ON UPDATE RESTRICT ON DELETE RESTRICT};
  la app no tiene privilegio DELETE en tablas de hechos.
\item
  Sensibilidad: \textbf{P} personal (LFPDPPP {[}F10{]}), \textbf{C}
  comercial confidencial, \textbf{T} técnica, \textbf{S} secreta.
\item
  Replicación: \textbf{A→B} maestros originados en A; \textbf{L↔} local
  del dueño (A→B y B→A según
  \texorpdfstring{\softtext{nodo\_origen);}}{nodo\_origen);}
  \textbf{local} no se publica.
\end{itemize}

\begin{center}\rule{0.5\linewidth}{0.5pt}\end{center}

\subsubsection{ENT-01 empresa}\label{ent-01-empresa}

Propósito: cliente del SaaS; raíz del aislamiento (RN-001). Dueño: A ·
Replicación A→B · Sensibilidad C. - \textbf{Columna:} id. \textbf{Tipo:}
uuid. \textbf{Nulo:} no. \textbf{Default:} uuidv7(). \textbf{Dominio /
nota:} PK.

\begin{itemize}
\item
  \textbf{Columna:}
  \texorpdfstring{\softtext{razon\_social.}}{razon\_social.}
  \textbf{Tipo:}
  \texorpdfstring{\softtext{varchar(200).}}{varchar(200).}
  \textbf{Nulo:} no. \textbf{Default:} . \textbf{Dominio / nota:} .
\item
  \textbf{Columna:} rfc. \textbf{Tipo:}
  \texorpdfstring{\softtext{varchar(13).}}{varchar(13).} \textbf{Nulo:}
  sí. \textbf{Default:} . \textbf{Dominio / nota:} P/C.
\item
  \textbf{Columna:} estado. \textbf{Tipo:}
  \texorpdfstring{\softtext{varchar(12).}}{varchar(12).} \textbf{Nulo:}
  no. \textbf{Default:} `ACTIVA'. \textbf{Dominio / nota:} ACTIVA,
  \texorpdfstring{\softtext{SUSPENDIDA.}}{SUSPENDIDA.}
\item
  \textbf{Columna:} creado\_en. \textbf{Tipo:}
  \texorpdfstring{\softtext{timestamptz.}}{timestamptz.} \textbf{Nulo:}
  no. \textbf{Default:} now(). \textbf{Dominio / nota:} .
\end{itemize}

Válido:
\texorpdfstring{\softcode{('Acopio Demo SA', NULL)}}{('Acopio Demo SA', NULL)}.
Rechazado: estado `BORRADA'.

\subsubsection{ENT-02 sucursal}\label{ent-02-sucursal}

Propósito: sede y su nodo dueño (RN-004, RN-040). PK(E,id);
\texorpdfstring{\softtext{UK(E,codigo).}}{UK(E,codigo).} Dueño A · A→B ·
C. - \textbf{Columna:} codigo. \textbf{Tipo:}
\texorpdfstring{\softtext{varchar(10).}}{varchar(10).} \textbf{Nulo:}
no. \textbf{Default:} . \textbf{Nota:}
\texorpdfstring{\softcode{\textasciicircum{}[A-Z0-9-]+\$}}{\textasciicircum{}[A-Z0-9-]+\$}.

\begin{itemize}
\item
  \textbf{Columna:} nombre. \textbf{Tipo:}
  \texorpdfstring{\softtext{varchar(120).}}{varchar(120).}
  \textbf{Nulo:} no. \textbf{Default:} . \textbf{Nota:} .
\item
  \textbf{Columna:}
  \texorpdfstring{\softtext{nodo\_dueno.}}{nodo\_dueno.} \textbf{Tipo:}
  char(1). \textbf{Nulo:} no. \textbf{Default:} . \textbf{Nota:} A, B;
  inmutable (trigger).
\item
  \textbf{Columna:}
  \texorpdfstring{\softtext{zona\_horaria.}}{zona\_horaria.}
  \textbf{Tipo:} \texorpdfstring{\softtext{varchar(40).}}{varchar(40).}
  \textbf{Nulo:} no. \textbf{Default:}
  `\texorpdfstring{\softtext{America/Mexico\_City}}{America/Mexico\_City}'.
  \textbf{Nota:} .
\end{itemize}

Índices: PK, UK. Válido: (E,`CEN',`A'). Rechazado: dos sedes con código
`CEN' en la misma empresa.

\subsubsection{ENT-03 parte}\label{ent-03-parte}

Propósito: contraparte que compra y/o vende (RN-010). PK(E,id); UK
\texorpdfstring{\softtext{parcial(E,rfc).}}{parcial(E,rfc).} A · A→B ·
\textbf{P}. - \textbf{Columna:}
\texorpdfstring{\softtext{tipo\_persona.}}{tipo\_persona.}
\textbf{Tipo:} \texorpdfstring{\softtext{varchar(6).}}{varchar(6).}
\textbf{Nulo:} no. \textbf{Default:} . \textbf{Nota:} FISICA, MORAL.

\begin{itemize}
\item
  \textbf{Columna:} nombre. \textbf{Tipo:}
  \texorpdfstring{\softtext{varchar(200).}}{varchar(200).}
  \textbf{Nulo:} no. \textbf{Default:} . \textbf{Nota:} .
\item
  \textbf{Columna:} rfc. \textbf{Tipo:}
  \texorpdfstring{\softtext{varchar(13).}}{varchar(13).} \textbf{Nulo:}
  sí. \textbf{Default:} . \textbf{Nota:} normalizado
  \texorpdfstring{\softtext{mayúsculas.}}{mayúsculas.}
\item
  \textbf{Columna:} telefono. \textbf{Tipo:}
  \texorpdfstring{\softtext{varchar(20).}}{varchar(20).} \textbf{Nulo:}
  sí. \textbf{Default:} . \textbf{Nota:} P.
\item
  \textbf{Columna:}
  \texorpdfstring{\softtext{es\_sintetico.}}{es\_sintetico.}
  \textbf{Tipo:} boolean. \textbf{Nulo:} no. \textbf{Default:} false.
  \textbf{Nota:} RN-045.
\end{itemize}

Índices:
\texorpdfstring{\softcode{(E, lower(nombre))}}{(E, lower(nombre))}
búsqueda. Rechazado: mismo RFC dos veces en la empresa.

\subsubsection{ENT-04 predio}\label{ent-04-predio}

Propósito:
\texorpdfstring{\softtext{origen/destino/custodia}}{origen/destino/custodia}
\texorpdfstring{\softtext{declarados.}}{declarados.} PK(E,id);
\texorpdfstring{\softtext{FK(E,parte\_responsable\_id).}}{FK(E,parte\_responsable\_id).}
A · A→B · C/P. Columnas: nombre varchar(160) NOT NULL;
\texorpdfstring{\softtext{clave\_upp\_declarada}}{clave\_upp\_declarada}
varchar(40) NULL (texto declarado, no certificado); municipio
varchar(80);
\texorpdfstring{\softtext{entidad\_federativa}}{entidad\_federativa}
varchar(40);
\texorpdfstring{\softtext{parte\_responsable\_id}}{parte\_responsable\_id}
uuid NULL. Rechazado:
\texorpdfstring{\softtext{parte\_responsable\_id}}{parte\_responsable\_id}
de otra empresa (FK compuesta falla, TEST-14).

\subsubsection{ENT-05 usuario}\label{ent-05-usuario}

Propósito: usuario del producto (no rol
\texorpdfstring{\softtext{PostgreSQL).}}{PostgreSQL).} PK(E,id);
\texorpdfstring{\softtext{UK(E,login\_normalizado).}}{UK(E,login\_normalizado).}
A · A→B · \textbf{S} (hash). Columnas:
\texorpdfstring{\softtext{login\_normalizado}}{login\_normalizado}
varchar(80) NOT NULL
(\texorpdfstring{\softcode{lower(trim())}}{lower(trim())}); nombre
varchar(120);
\texorpdfstring{\softtext{hash\_contrasena}}{hash\_contrasena}
varchar(255) NOT NULL (formato
\texorpdfstring{\softcode{\{bcrypt\}}}{\{bcrypt\}}/\texorpdfstring{\softcode{\{argon2\}}}{\{argon2\}}
de Spring Security
\texorpdfstring{\softtext{DelegatingPasswordEncoder);}}{DelegatingPasswordEncoder);}
estado varchar(10) DEFAULT `ACTIVO' (ACTIVO, BLOQUEADO); creado\_en.
Nunca se loguea ni se indexa en el grafo el hash.

\subsubsection{\texorpdfstring{ENT-06
\texorpdfstring{\softtext{rol\_aplicacion}}{rol\_aplicacion}}{ENT-06 }}\label{ent-06-rol_aplicacion}

PK(codigo varchar(12));
\texorpdfstring{\softtext{descripcion.}}{descripcion.} Valores: ADMIN,
OPERADOR, CAJA, AUDITOR. A→B · T.

\subsubsection{\texorpdfstring{ENT-07
\texorpdfstring{\softtext{membresia\_usuario}}{membresia\_usuario}}{ENT-07 }}\label{ent-07-membresia_usuario}

\texorpdfstring{\softtext{PK(E,S,usuario\_id,rol\_codigo);}}{PK(E,S,usuario\_id,rol\_codigo);}
vigente boolean DEFAULT true. Autoriza por sede (RN-001). A · A→B · T.
Índice:
\texorpdfstring{\softcode{(E, usuario\_id) WHERE vigente}}{(E, usuario\_id) WHERE vigente}.
Rechazado: rol \texorpdfstring{\softtext{inexistente.}}{inexistente.}

\subsubsection{ENT-08 animal}\label{ent-08-animal}

Propósito: identidad permanente (RN-002, RN-004). PK(E,id);
\texorpdfstring{\softtext{UK(E,sede\_operadora\_id,id).}}{UK(E,sede\_operadora\_id,id).}
A · A→B · C. - \textbf{Columna:}
\texorpdfstring{\softtext{sede\_operadora\_id.}}{sede\_operadora\_id.}
\textbf{Tipo:} uuid. \textbf{Nulo:} no. \textbf{Default:} .
\textbf{Nota:} FK (E,·) → sucursal; inmutable (trigger).

\begin{itemize}
\item
  \textbf{Columna:} sexo. \textbf{Tipo:} char(1). \textbf{Nulo:} no.
  \textbf{Default:} . \textbf{Nota:} M, H.
\item
  \textbf{Columna:}
  \texorpdfstring{\softtext{fecha\_nacimiento\_declarada.}}{fecha\_nacimiento\_declarada.}
  \textbf{Tipo:} date. \textbf{Nulo:} sí. \textbf{Default:} .
  \textbf{Nota:} declarada.
\item
  \textbf{Columna:}
  \texorpdfstring{\softtext{raza\_declarada.}}{raza\_declarada.}
  \textbf{Tipo:} \texorpdfstring{\softtext{varchar(60).}}{varchar(60).}
  \textbf{Nulo:} sí. \textbf{Default:} . \textbf{Nota:} .
\item
  \textbf{Columna:} situacion. \textbf{Tipo:}
  \texorpdfstring{\softtext{varchar(24).}}{varchar(24).} \textbf{Nulo:}
  no. \textbf{Default:} `VIGENTE'. \textbf{Nota:} VIGENTE,
  \texorpdfstring{\softtext{IDENTIDAD\_EN\_REVISION,}}{IDENTIDAD\_EN\_REVISION,}
  MUERTO (DEC-04).
\item
  \textbf{Columna:} creado\_en. \textbf{Tipo:}
  \texorpdfstring{\softtext{timestamptz.}}{timestamptz.} \textbf{Nulo:}
  no. \textbf{Default:} now(). \textbf{Nota:} .
\end{itemize}

Rechazado:
\texorpdfstring{\softcode{UPDATE animal SET sede\_operadora\_id = …}}{UPDATE animal SET sede\_operadora\_id = …}
(BV010).

\subsubsection{\texorpdfstring{ENT-36
\texorpdfstring{\softtext{animal\_identificacion}}{animal\_identificacion}}{ENT-36 }}\label{ent-36-animal_identificacion}

Propósito: códigos
\texorpdfstring{\softtext{oficiales/internos}}{oficiales/internos} con
historia (RN-003). PK(E,id);
\texorpdfstring{\softtext{UK(E,sistema,codigo\_normalizado).}}{UK(E,sistema,codigo\_normalizado).}
A · A→B · C. - \textbf{Columna:} animal\_id. \textbf{Tipo:} uuid.
\textbf{Nulo:} no. \textbf{Nota:} FK (E,·) → animal.

\begin{itemize}
\item
  \textbf{Columna:} sistema. \textbf{Tipo:}
  \texorpdfstring{\softtext{varchar(10).}}{varchar(10).} \textbf{Nulo:}
  no. \textbf{Nota:} SINIIGA, SINIDA, INTERNO, OTRO.
\item
  \textbf{Columna:}
  \texorpdfstring{\softtext{codigo\_declarado.}}{codigo\_declarado.}
  \textbf{Tipo:} \texorpdfstring{\softtext{varchar(40).}}{varchar(40).}
  \textbf{Nulo:} no. \textbf{Nota:} tal como se capturó.
\item
  \textbf{Columna:}
  \texorpdfstring{\softtext{codigo\_normalizado.}}{codigo\_normalizado.}
  \textbf{Tipo:} \texorpdfstring{\softtext{varchar(40).}}{varchar(40).}
  \textbf{Nulo:} no. \textbf{Nota:}
  \texorpdfstring{\softcode{regexp\_replace(upper(codigo),'[ -]','','g'); INTERNO admite ASCII alfanumérico, espacio/guion; rechaza vacío o caracteres fuera de política}}{regexp\_replace(upper(codigo),'[ -]','','g'); INTERNO admite ASCII alfanumérico, espacio/guion; rechaza vacío o caracteres fuera de política}.
\item
  \textbf{Columna:} estado. \textbf{Tipo:}
  \texorpdfstring{\softtext{varchar(12).}}{varchar(12).} \textbf{Nulo:}
  no. \textbf{Nota:} VIGENTE, INACTIVA,
  \texorpdfstring{\softtext{EN\_REVISION.}}{EN\_REVISION.}
\item
  \textbf{Columna:}
  \texorpdfstring{\softtext{vigente\_desde}}{vigente\_desde} /
  \texorpdfstring{\softtext{vigente\_hasta.}}{vigente\_hasta.}
  \textbf{Tipo:} date. \textbf{Nulo:} no / sí. \textbf{Nota:} CHECK
  hasta ≥ desde.
\item
  \textbf{Columna:}
  \texorpdfstring{\softtext{evidencia\_sucursal\_id,}}{evidencia\_sucursal\_id,}
  \texorpdfstring{\softtext{evidencia\_documento\_id.}}{evidencia\_documento\_id.}
  \textbf{Tipo:} uuid. \textbf{Nulo:} sí. \textbf{Nota:} FK documento
  (ambas o ninguna).
\end{itemize}

Índice parcial: UNIQUE(E, animal\_id, sistema) WHERE estado=`VIGENTE'.
Rechazado: mismo código normalizado en otro animal (TEST-13).

\subsubsection{ENT-09 operacion}\label{ent-09-operacion}

Propósito: compra o venta por cuenta propia (RN-009, RN-010).
PK(E,S,id); UK(E,S,tipo,folio);
\texorpdfstring{\softtext{UK(E,S,id,tipo).}}{UK(E,S,id,tipo).} Dueño:
nodo de la sede · L↔ · C. - \textbf{Columna:} tipo. \textbf{Tipo:}
\texorpdfstring{\softtext{varchar(6).}}{varchar(6).} \textbf{Nulo:} no.
\textbf{Default:} . \textbf{Nota:} COMPRA, VENTA; inmutable.

\begin{itemize}
\item
  \textbf{Columna:} folio. \textbf{Tipo:}
  \texorpdfstring{\softtext{varchar(20).}}{varchar(20).} \textbf{Nulo:}
  no. \textbf{Default:} . \textbf{Nota:}
  \texorpdfstring{\softcode{C-2026-000123}}{C-2026-000123}; de
  \texorpdfstring{\softtext{folio\_contador;}}{folio\_contador;} no
  fiscal.
\item
  \textbf{Columna:} parte\_id. \textbf{Tipo:} uuid. \textbf{Nulo:} no.
  \textbf{Default:} . \textbf{Nota:} → parte.
\item
  \textbf{Columna:}
  \texorpdfstring{\softtext{predio\_origen\_id}}{predio\_origen\_id} /
  \texorpdfstring{\softtext{predio\_destino\_id.}}{predio\_destino\_id.}
  \textbf{Tipo:} uuid. \textbf{Nulo:} sí. \textbf{Default:} .
  \textbf{Nota:} → predio.
\item
  \textbf{Columna:} moneda. \textbf{Tipo:} char(3). \textbf{Nulo:} no.
  \textbf{Default:} `MXN'. \textbf{Nota:} CHECK = `MXN' en MVP.
\item
  \textbf{Columna:} estado. \textbf{Tipo:}
  \texorpdfstring{\softtext{varchar(12).}}{varchar(12).} \textbf{Nulo:}
  no. \textbf{Default:} `BORRADOR'. \textbf{Nota:} BORRADOR, ACORDADA,
  CONFIRMADA, CANCELADA, ANULADA.
\item
  \textbf{Columna:}
  \texorpdfstring{\softtext{estado\_entrega.}}{estado\_entrega.}
  \textbf{Tipo:} \texorpdfstring{\softtext{varchar(10).}}{varchar(10).}
  \textbf{Nulo:} no. \textbf{Default:} `PENDIENTE'. \textbf{Nota:}
  PENDIENTE, PARCIAL, COMPLETA (mantenido por comando de
  \texorpdfstring{\softtext{movimientos).}}{movimientos).}
\item
  \textbf{Columna:}
  \texorpdfstring{\softtext{numero\_version\_aceptada.}}{numero\_version\_aceptada.}
  \textbf{Tipo:} smallint. \textbf{Nulo:} sí. \textbf{Default:} .
  \textbf{Nota:} NOT NULL si CONFIRMADA (CHECK).
\item
  \textbf{Columna:}
  \texorpdfstring{\softtext{fecha\_operacion.}}{fecha\_operacion.}
  \textbf{Tipo:} date. \textbf{Nulo:} no. \textbf{Default:} .
  \textbf{Nota:} .
\item
  \textbf{Columna:}
  \texorpdfstring{\softtext{version\_lock.}}{version\_lock.}
  \textbf{Tipo:} integer. \textbf{Nulo:} no. \textbf{Default:} 0.
  \textbf{Nota:} concurrencia optimista de borrador.
\item
  \textbf{Columna:} creado\_por / creado\_en /
  \texorpdfstring{\softtext{confirmada\_en.}}{confirmada\_en.}
  \textbf{Tipo:} uuid / timestamptz /
  \texorpdfstring{\softtext{timestamptz.}}{timestamptz.} \textbf{Nulo:}
  no/no/sí. \textbf{Default:} . \textbf{Nota:} .
\end{itemize}

Índices:
\texorpdfstring{\softcode{(E,S,estado,fecha\_operacion)}}{(E,S,estado,fecha\_operacion)}
bandeja; \texorpdfstring{\softcode{(E,parte\_id)}}{(E,parte\_id)}
historial de \texorpdfstring{\softtext{contraparte.}}{contraparte.}
Rechazado: tipo `CONSIGNA'.

\subsubsection{\texorpdfstring{ENT-10
\texorpdfstring{\softtext{lote\_operacion}}{lote\_operacion}}{ENT-10 }}\label{ent-10-lote_operacion}

\texorpdfstring{\softtext{PK(E,S,operacion\_id,numero\_lote}}{PK(E,S,operacion\_id,numero\_lote}
smallint CHECK \textgreater{} 0); descripcion
\texorpdfstring{\softtext{varchar(120).}}{varchar(120).} Débil. L↔.
Editable sólo si la operación está en BORRADOR (trigger).

\subsubsection{\texorpdfstring{ENT-11
\texorpdfstring{\softtext{seleccion\_compra}}{seleccion\_compra}}{ENT-11 }}\label{ent-11-seleccion_compra}

\texorpdfstring{\softtext{PK(E,S,operacion\_id,numero\_lote,animal\_id);}}{PK(E,S,operacion\_id,numero\_lote,animal\_id);}
\texorpdfstring{\softtext{UK(E,S,operacion\_id,animal\_id);}}{UK(E,S,operacion\_id,animal\_id);}
\texorpdfstring{\softcode{tipo\_operacion varchar(6) NOT NULL DEFAULT 'COMPRA' CHECK (= 'COMPRA')}}{tipo\_operacion varchar(6) NOT NULL DEFAULT 'COMPRA' CHECK (= 'COMPRA')}.
FKs a lote, a operacion(E,S,id,tipo) y a
\texorpdfstring{\softtext{animal(E,sede\_operadora\_id,id).}}{animal(E,sede\_operadora\_id,id).}
L↔. Rechazado: animal de la sede A en compra de sede B (FK falla,
RN-009); mismo animal en dos lotes de la misma compra (UK).

\subsubsection{\texorpdfstring{ENT-12
\texorpdfstring{\softtext{seleccion\_venta}}{seleccion\_venta}}{ENT-12 }}\label{ent-12-seleccion_venta}

Igual que ENT-11 con
\texorpdfstring{\softcode{posicion\_id}}{posicion\_id} y
\texorpdfstring{\softcode{tipo\_operacion='VENTA'}}{tipo\_operacion='VENTA'};
\texorpdfstring{\softtext{UK(E,S,operacion\_id,posicion\_id).}}{UK(E,S,operacion\_id,posicion\_id).}
L↔.

\subsubsection{\texorpdfstring{ENT-13
\texorpdfstring{\softtext{posicion\_inventario}}{posicion\_inventario}}{ENT-13 }}\label{ent-13-posicion_inventario}

Propósito: ciclo de inventario (RN-005, RN-006).
\texorpdfstring{\softtext{PK(E,S,id).}}{PK(E,S,id).} L↔ · C. -
\textbf{Columna:} animal\_id. \textbf{Tipo:} uuid. \textbf{Nulo:} no.
\textbf{Nota:} →
\texorpdfstring{\softtext{animal(E,sede\_operadora\_id,id).}}{animal(E,sede\_operadora\_id,id).}

\begin{itemize}
\item
  \textbf{Columna:}
  \texorpdfstring{\softtext{numero\_ciclo.}}{numero\_ciclo.}
  \textbf{Tipo:} smallint. \textbf{Nulo:} no. \textbf{Nota:}
  \textgreater{} 0; max+1 por animal.
\item
  \textbf{Columna:}
  \texorpdfstring{\softtext{compra\_operacion\_id,}}{compra\_operacion\_id,}
  \texorpdfstring{\softtext{compra\_numero\_lote.}}{compra\_numero\_lote.}
  \textbf{Tipo:} uuid, smallint. \textbf{Nulo:} no. \textbf{Nota:} →
  \texorpdfstring{\softtext{seleccion\_compra}}{seleccion\_compra} (con
  \texorpdfstring{\softtext{animal\_id).}}{animal\_id).}
\item
  \textbf{Columna:} estado. \textbf{Tipo:}
  \texorpdfstring{\softtext{varchar(26).}}{varchar(26).} \textbf{Nulo:}
  no. \textbf{Nota:}
  \texorpdfstring{\softtext{PENDIENTE\_RECEPCION,}}{PENDIENTE\_RECEPCION,}
  DISPONIBLE,
  \texorpdfstring{\softtext{VENDIDA\_PENDIENTE\_ENTREGA,}}{VENDIDA\_PENDIENTE\_ENTREGA,}
  \texorpdfstring{\softtext{CERRADA\_VENTA,}}{CERRADA\_VENTA,} BAJA.
\item
  \textbf{Columna:}
  \texorpdfstring{\softtext{motivo\_baja.}}{motivo\_baja.}
  \textbf{Tipo:} \texorpdfstring{\softtext{varchar(10).}}{varchar(10).}
  \textbf{Nulo:} sí. \textbf{Nota:} PERDIDA, MUERTE; NOT NULL ⇔ estado =
  BAJA.
\item
  \textbf{Columna:} abierta\_en /
  \texorpdfstring{\softtext{cerrada\_en.}}{cerrada\_en.} \textbf{Tipo:}
  \texorpdfstring{\softtext{timestamptz.}}{timestamptz.} \textbf{Nulo:}
  no / sí. \textbf{Nota:} cerrada\_en NOT NULL ⇔ estado terminal.
\end{itemize}

Índices: UNIQUE parcial
\texorpdfstring{\softcode{(E, animal\_id) WHERE estado IN (activos)}}{(E, animal\_id) WHERE estado IN (activos)};
\texorpdfstring{\softcode{(E,S,estado)}}{(E,S,estado)} inventario
\texorpdfstring{\softtext{disponible.}}{disponible.} Sin
\texorpdfstring{\softtext{INSERT/UPDATE}}{INSERT/UPDATE} directos de la
app (sólo comandos). Rechazado: segunda posición activa del mismo animal
(TEST-05).

\subsubsection{\texorpdfstring{ENT-14
\texorpdfstring{\softtext{movimiento\_inventario}}{movimiento\_inventario}}{ENT-14 }}\label{ent-14-movimiento_inventario}

Hecho físico inmutable (RN-020). Columnas:
\texorpdfstring{\softtext{posicion\_id}}{posicion\_id} NOT NULL;
\texorpdfstring{\softtext{operacion\_id}}{operacion\_id} NULL; tipo
varchar(12) (RECEPCION, ENTREGA, BAJA, REVERSO);
\texorpdfstring{\softtext{ocurrido\_en}}{ocurrido\_en} timestamptz;
predio\_id,
\texorpdfstring{\softtext{custodio\_parte\_id,}}{custodio\_parte\_id,}
\texorpdfstring{\softtext{documento\_id}}{documento\_id} NULL;
\texorpdfstring{\softtext{reverso\_de\_id}}{reverso\_de\_id} UNIQUE
NULL; \texorpdfstring{\softtext{registrado\_por.}}{registrado\_por.}
Índice
\texorpdfstring{\softcode{(E,S,posicion\_id, ocurrido\_en)}}{(E,S,posicion\_id, ocurrido\_en)}.
L↔.

\subsubsection{\texorpdfstring{ENT-15
\texorpdfstring{\softtext{reserva\_venta}}{reserva\_venta}}{ENT-15 }}\label{ent-15-reserva_venta}

Columnas: \texorpdfstring{\softtext{operacion\_id,}}{operacion\_id,}
\texorpdfstring{\softtext{numero\_lote,}}{numero\_lote,}
\texorpdfstring{\softtext{posicion\_id}}{posicion\_id} (→
\texorpdfstring{\softtext{seleccion\_venta);}}{seleccion\_venta);}
estado varchar(10) (ACTIVA, CONSUMIDA, LIBERADA, EXPIRADA); creada\_en;
vence\_en timestamptz NOT NULL CHECK \textgreater{} creada\_en;
cerrada\_en NULL. Índice UNIQUE parcial
\texorpdfstring{\softcode{(E,S,posicion\_id) WHERE estado='ACTIVA'}}{(E,S,posicion\_id) WHERE estado='ACTIVA'}
(RN-018);
\texorpdfstring{\softcode{(E,S,vence\_en) WHERE estado='ACTIVA'}}{(E,S,vence\_en) WHERE estado='ACTIVA'}
para el expirador. L↔.

\subsubsection{\texorpdfstring{ENT-16
\texorpdfstring{\softtext{contrato\_version}}{contrato\_version}}{ENT-16 }}\label{ent-16-contrato_version}

\texorpdfstring{\softtext{PK(E,S,operacion\_id,numero\_version}}{PK(E,S,operacion\_id,numero\_version}
smallint \textgreater{} 0). estado (PROPUESTA, ACEPTADA, SUPERADA);
\texorpdfstring{\softtext{politica\_redondeo}}{politica\_redondeo}
varchar(12) DEFAULT `HALF\_UP\_2';
\texorpdfstring{\softtext{requiere\_recepcion}}{requiere\_recepcion}
boolean DEFAULT false;
\texorpdfstring{\softtext{documento\_id}}{documento\_id} NULL;
\texorpdfstring{\softtext{aceptada\_por/aceptada\_en}}{aceptada\_por/aceptada\_en}
NULL. Inmutable una vez ACEPTADA (RN-011). Pueden conservarse varias
versiones ACEPTADAS históricas:
\texorpdfstring{\softcode{operacion.numero\_version\_aceptada}}{operacion.numero\_version\_aceptada}
identifica la que referencia la
\texorpdfstring{\softtext{operación.}}{operación.} La restricción
parcial propuesta inicialmente para permitir sólo una ACEPTADA no existe
en el DDL vigente de
\texorpdfstring{\softcode{V1\_\_esquema.sql}}{V1\_\_esquema.sql}; la PK
distingue cada versión sin modificar las aceptadas
\texorpdfstring{\softtext{anteriores.}}{anteriores.}

\subsubsection{\texorpdfstring{ENT-17
\texorpdfstring{\softtext{condicion\_lote}}{condicion\_lote}}{ENT-17 }}\label{ent-17-condicion_lote}

\begin{itemize}
\item
  \textbf{Columna:}
  \texorpdfstring{\softtext{unidad\_precio.}}{unidad\_precio.}
  \textbf{Tipo:} \texorpdfstring{\softtext{varchar(6).}}{varchar(6).}
  \textbf{Nulo:} no. \textbf{Default:} . \textbf{Nota:} KG, CABEZA.
\item
  \textbf{Columna:}
  \texorpdfstring{\softtext{precio\_unitario.}}{precio\_unitario.}
  \textbf{Tipo:}
  \texorpdfstring{\softtext{numeric(14,4).}}{numeric(14,4).}
  \textbf{Nulo:} no. \textbf{Default:} . \textbf{Nota:} \textgreater{}
  0.
\item
  \textbf{Columna:}
  \texorpdfstring{\softtext{base\_pesaje.}}{base\_pesaje.}
  \textbf{Tipo:} \texorpdfstring{\softtext{varchar(10).}}{varchar(10).}
  \textbf{Nulo:} no. \textbf{Default:} `NINGUNA'. \textbf{Nota:} NINGUNA
  (por cabeza), PESAJE.
\item
  \textbf{Columna:}
  \texorpdfstring{\softtext{pesaje\_base\_id.}}{pesaje\_base\_id.}
  \textbf{Tipo:} uuid. \textbf{Nulo:} sí. \textbf{Default:} .
  \textbf{Nota:} NOT NULL ⇔
  \texorpdfstring{\softtext{base\_pesaje}}{base\_pesaje} = `PESAJE'.
\item
  \textbf{Columna:}
  \texorpdfstring{\softtext{deduccion\_pct.}}{deduccion\_pct.}
  \textbf{Tipo:}
  \texorpdfstring{\softtext{numeric(9,6).}}{numeric(9,6).}
  \textbf{Nulo:} no. \textbf{Default:} 0. \textbf{Nota:} 0 ≤ x
  \textless{} 100; sólo pactada (RN-015).
\item
  \textbf{Columna:}
  \texorpdfstring{\softtext{cargos\_pactados}}{cargos\_pactados} /
  \texorpdfstring{\softtext{bonificaciones\_pactadas.}}{bonificaciones\_pactadas.}
  \textbf{Tipo:}
  \texorpdfstring{\softtext{numeric(18,2).}}{numeric(18,2).}
  \textbf{Nulo:} no. \textbf{Default:} 0. \textbf{Nota:} ≥ 0.
\end{itemize}

Rechazado: unidad KG sin pesaje base.

\subsubsection{ENT-18 pesaje}\label{ent-18-pesaje}

\begin{itemize}
\item
  \textbf{Columna:}
  \texorpdfstring{\softtext{operacion\_id,}}{operacion\_id,}
  \texorpdfstring{\softtext{numero\_lote.}}{numero\_lote.}
  \textbf{Tipo:} uuid, smallint. \textbf{Nulo:} no. \textbf{Nota:} →
  \texorpdfstring{\softtext{lote\_operacion.}}{lote\_operacion.}
\item
  \textbf{Columna:} modo. \textbf{Tipo:}
  \texorpdfstring{\softtext{varchar(14).}}{varchar(14).} \textbf{Nulo:}
  no. \textbf{Nota:}
  \texorpdfstring{\softtext{BRUTO\_TARA,}}{BRUTO\_TARA,}
  \texorpdfstring{\softtext{NETO\_REPORTADO.}}{NETO\_REPORTADO.}
\item
  \textbf{Columna:} bruto\_kg, tara\_kg. \textbf{Tipo:}
  \texorpdfstring{\softtext{numeric(14,3).}}{numeric(14,3).}
  \textbf{Nulo:} sí. \textbf{Nota:} NOT NULL ⇔
  \texorpdfstring{\softtext{BRUTO\_TARA;}}{BRUTO\_TARA;} tara ≤ bruto.
\item
  \textbf{Columna:} neto\_kg. \textbf{Tipo:}
  \texorpdfstring{\softtext{numeric(14,3).}}{numeric(14,3).}
  \textbf{Nulo:} no. \textbf{Nota:} \textgreater{} 0; = bruto − tara si
  \texorpdfstring{\softtext{BRUTO\_TARA.}}{BRUTO\_TARA.}
\item
  \textbf{Columna:}
  \texorpdfstring{\softtext{ticket\_folio.}}{ticket\_folio.}
  \textbf{Tipo:} \texorpdfstring{\softtext{varchar(40).}}{varchar(40).}
  \textbf{Nulo:} sí. \textbf{Nota:} .
\item
  \textbf{Columna:} bascula. \textbf{Tipo:}
  \texorpdfstring{\softtext{varchar(80).}}{varchar(80).} \textbf{Nulo:}
  sí. \textbf{Nota:} descripción declarada.
\item
  \textbf{Columna:} pesado\_en. \textbf{Tipo:}
  \texorpdfstring{\softtext{timestamptz.}}{timestamptz.} \textbf{Nulo:}
  no. \textbf{Nota:} .
\item
  \textbf{Columna:} estado. \textbf{Tipo:}
  \texorpdfstring{\softtext{varchar(10).}}{varchar(10).} \textbf{Nulo:}
  no. \textbf{Nota:} PROPUESTA, APROBADO, CORREGIDO.
\item
  \textbf{Columna:} fuente. \textbf{Tipo:}
  \texorpdfstring{\softtext{varchar(8).}}{varchar(8).} \textbf{Nulo:}
  no. \textbf{Nota:} MANUAL, CSV, OCR, FDW (RN-016).
\item
  \textbf{Columna:}
  \texorpdfstring{\softtext{corrige\_a\_id.}}{corrige\_a\_id.}
  \textbf{Tipo:} uuid. \textbf{Nulo:} sí. \textbf{Nota:} UNIQUE →
  pesaje.
\end{itemize}

Rechazado: bruto 5,000, tara 6,000 (TEST);
\texorpdfstring{\softtext{NETO\_REPORTADO}}{NETO\_REPORTADO} con bruto
informado.

\subsubsection{\texorpdfstring{ENT-19
\texorpdfstring{\softtext{pesaje\_animal}}{pesaje\_animal}}{ENT-19 }}\label{ent-19-pesaje_animal}

\texorpdfstring{\softtext{PK(E,S,pesaje\_id,animal\_id);}}{PK(E,S,pesaje\_id,animal\_id);}
\texorpdfstring{\softtext{peso\_individual\_kg}}{peso\_individual\_kg}
numeric(14,3) NULL (\textgreater{} 0 si existe; sólo medido, RN-013).
Inmutable tras APROBADO.

\subsubsection{\texorpdfstring{ENT-20
\texorpdfstring{\softtext{gasto\_operacion}}{gasto\_operacion}}{ENT-20 }}\label{ent-20-gasto_operacion}

\begin{itemize}
\item
  \textbf{Columna:}
  \texorpdfstring{\softtext{operacion\_id.}}{operacion\_id.}
  \textbf{Tipo:} uuid. \textbf{Nulo:} no. \textbf{Nota:} .
\item
  \textbf{Columna:} \texorpdfstring{\softtext{naturaleza.}}{naturaleza.}
  \textbf{Tipo:} \texorpdfstring{\softtext{varchar(12).}}{varchar(12).}
  \textbf{Nulo:} no. \textbf{Nota:} COSTO,
  \texorpdfstring{\softtext{BONIFICACION.}}{BONIFICACION.}
\item
  \textbf{Columna:}
  \texorpdfstring{\softtext{clasificacion.}}{clasificacion.}
  \textbf{Tipo:} \texorpdfstring{\softtext{varchar(14).}}{varchar(14).}
  \textbf{Nulo:} no. \textbf{Nota:} CAPITALIZABLE,
  \texorpdfstring{\softtext{DIRECTO\_VENTA,}}{DIRECTO\_VENTA,} OTRO
  (RN-025).
\item
  \textbf{Columna:} importe. \textbf{Tipo:}
  \texorpdfstring{\softtext{numeric(18,2).}}{numeric(18,2).}
  \textbf{Nulo:} no. \textbf{Nota:} \textgreater{} 0 (signo derivado de
  \texorpdfstring{\softtext{naturaleza).}}{naturaleza).}
\item
  \textbf{Columna:} concepto. \textbf{Tipo:}
  \texorpdfstring{\softtext{varchar(160).}}{varchar(160).}
  \textbf{Nulo:} no. \textbf{Nota:} .
\item
  \textbf{Columna:}
  \texorpdfstring{\softtext{proveedor\_parte\_id.}}{proveedor\_parte\_id.}
  \textbf{Tipo:} uuid. \textbf{Nulo:} sí. \textbf{Nota:} .
\item
  \textbf{Columna:}
  \texorpdfstring{\softtext{referencia\_origen.}}{referencia\_origen.}
  \textbf{Tipo:} \texorpdfstring{\softtext{varchar(60).}}{varchar(60).}
  \textbf{Nulo:} sí. \textbf{Nota:} UK parcial con proveedor: detecta
  \texorpdfstring{\softtext{duplicados.}}{duplicados.}
\item
  \textbf{Columna:}
  \texorpdfstring{\softtext{bonifica\_gasto\_id.}}{bonifica\_gasto\_id.}
  \textbf{Tipo:} uuid. \textbf{Nulo:} sí. \textbf{Nota:} NOT NULL ⇔
  \texorpdfstring{\softtext{BONIFICACION.}}{BONIFICACION.}
\item
  \textbf{Columna:} estado. \textbf{Tipo:}
  \texorpdfstring{\softtext{varchar(10).}}{varchar(10).} \textbf{Nulo:}
  no. \textbf{Nota:} BORRADOR, APROBADO.
\item
  \textbf{Columna:}
  \texorpdfstring{\softtext{fecha\_economica.}}{fecha\_economica.}
  \textbf{Tipo:} date. \textbf{Nulo:} no. \textbf{Nota:} .
\item
  \textbf{Columna:}
  \texorpdfstring{\softtext{snapshot\_fuente\_externa.}}{snapshot\_fuente\_externa.}
  \textbf{Tipo:} jsonb. \textbf{Nulo:} sí. \textbf{Nota:} tarifa FDW
  aceptada: fuente, ref, fecha, unidad, hash (RN-042).
\end{itemize}

\subsubsection{\texorpdfstring{ENT-21
\texorpdfstring{\softtext{asignacion\_costo}}{asignacion\_costo}}{ENT-21 }}\label{ent-21-asignacion_costo}

Columnas: \texorpdfstring{\softtext{posicion\_id;}}{posicion\_id;}
\textbf{exactamente una} de
\texorpdfstring{\softtext{(fuente\_liquidacion\_id}}{(fuente\_liquidacion\_id}
+
\texorpdfstring{\softtext{fuente\_numero\_concepto)}}{fuente\_numero\_concepto)}
/ \texorpdfstring{\softtext{fuente\_gasto\_id}}{fuente\_gasto\_id} /
\texorpdfstring{\softtext{fuente\_ajuste\_id}}{fuente\_ajuste\_id}
(\texorpdfstring{\softcode{CHECK (num\_nonnulls(fuente\_liquidacion\_id, fuente\_gasto\_id, fuente\_ajuste\_id) = 1)}}{CHECK (num\_nonnulls(fuente\_liquidacion\_id, fuente\_gasto\_id, fuente\_ajuste\_id) = 1)});
metodo
\texorpdfstring{\softtext{(IMPORTE\_INDIVIDUAL,}}{(IMPORTE\_INDIVIDUAL,}
\texorpdfstring{\softtext{POR\_CABEZA,}}{POR\_CABEZA,}
\texorpdfstring{\softtext{POR\_PESO\_MEDIDO);}}{POR\_PESO\_MEDIDO);}
base\_valor numeric(14,3); importe numeric(18,2) firmado (≠ 0);
\texorpdfstring{\softtext{lote\_reparto\_id}}{lote\_reparto\_id} uuid
(agrupa un reparto);
\texorpdfstring{\softtext{periodo\_id.}}{periodo\_id.} Índice
\texorpdfstring{\softcode{(E,S,lote\_reparto\_id)}}{(E,S,lote\_reparto\_id)}
para verificar Σ = fuente.

\subsubsection{ENT-22 liquidacion}\label{ent-22-liquidacion}

\begin{itemize}
\item
  \textbf{Columna:}
  \texorpdfstring{\softtext{operacion\_id.}}{operacion\_id.}
  \textbf{Tipo:} uuid. \textbf{Nulo:} no. \textbf{Nota:}
  \texorpdfstring{\softtext{UK(E,S,operacion\_id)}}{UK(E,S,operacion\_id)}
  principal.
\item
  \textbf{Columna:}
  \texorpdfstring{\softtext{numero\_version.}}{numero\_version.}
  \textbf{Tipo:} smallint. \textbf{Nulo:} no. \textbf{Nota:} versión
  aceptada.
\item
  \textbf{Columna:} direccion. \textbf{Tipo:}
  \texorpdfstring{\softtext{varchar(10).}}{varchar(10).} \textbf{Nulo:}
  no. \textbf{Nota:} POR\_PAGAR (compra), POR\_COBRAR (venta).
\item
  \textbf{Columna:} total. \textbf{Tipo:}
  \texorpdfstring{\softtext{numeric(18,2).}}{numeric(18,2).}
  \textbf{Nulo:} no. \textbf{Nota:} ≥ 0; = Σ conceptos (validado en
  comando).
\item
  \textbf{Columna:}
  \texorpdfstring{\softtext{snapshot\_contraparte\_nombre}}{snapshot\_contraparte\_nombre}
  / \_rfc. \textbf{Tipo:} varchar. \textbf{Nulo:} no / sí.
  \textbf{Nota:} RN-024.
\item
  \textbf{Columna:}
  \texorpdfstring{\softtext{snapshot\_condiciones.}}{snapshot\_condiciones.}
  \textbf{Tipo:} jsonb. \textbf{Nulo:} no. \textbf{Nota:} condiciones
  aceptadas tal cual.
\item
  \textbf{Columna:} vence\_el. \textbf{Tipo:} date. \textbf{Nulo:} sí.
  \textbf{Nota:} cartera vencida.
\item
  \textbf{Columna:}
  \texorpdfstring{\softtext{periodo\_id.}}{periodo\_id.} \textbf{Tipo:}
  uuid. \textbf{Nulo:} no. \textbf{Nota:} periodo de
  \texorpdfstring{\softtext{reconocimiento.}}{reconocimiento.}
\item
  \textbf{Columna:}
  \texorpdfstring{\softtext{confirmada\_por}}{confirmada\_por} /
  \texorpdfstring{\softtext{confirmada\_en.}}{confirmada\_en.}
  \textbf{Tipo:} . \textbf{Nulo:} no. \textbf{Nota:} .
\end{itemize}

Inmutable (trigger BV016). Índice
\texorpdfstring{\softcode{(E,S,direccion,vence\_el)}}{(E,S,direccion,vence\_el)}
cartera.

\subsubsection{\texorpdfstring{ENT-23
\texorpdfstring{\softtext{concepto\_liquidacion}}{concepto\_liquidacion}}{ENT-23 }}\label{ent-23-concepto_liquidacion}

\texorpdfstring{\softtext{PK(E,S,liquidacion\_id,numero\_concepto}}{PK(E,S,liquidacion\_id,numero\_concepto}
smallint). tipo (BASE, CARGO, BONIFICACION, DEDUCCION); descripcion;
\texorpdfstring{\softtext{numero\_lote}}{numero\_lote} NULL; cantidad
numeric(14,3); unidad (KG, CABEZA, MXN);
\texorpdfstring{\softtext{precio\_unitario}}{precio\_unitario}
numeric(14,4) NULL; importe numeric(18,2) firmado (negativo para
\texorpdfstring{\softtext{BONIFICACION/DEDUCCION,}}{BONIFICACION/DEDUCCION,}
CHECK por tipo); pesaje\_id NULL. Inmutable.

\subsubsection{\texorpdfstring{ENT-24
\texorpdfstring{\softtext{ajuste\_liquidacion}}{ajuste\_liquidacion}}{ENT-24 }}\label{ent-24-ajuste_liquidacion}

naturaleza (CARGO, ABONO); monto numeric(18,2) \textgreater{} 0; motivo
varchar(200) NOT NULL; fuente
\texorpdfstring{\softtext{(PRECIO\_GANADO,}}{(PRECIO\_GANADO,}
\texorpdfstring{\softtext{BONIFICACION\_COMERCIAL,}}{BONIFICACION\_COMERCIAL,}
ANULACION, OTRO);
\texorpdfstring{\softtext{documento\_id}}{documento\_id} NULL;
\texorpdfstring{\softtext{revierte\_ajuste\_id}}{revierte\_ajuste\_id}
UNIQUE NULL; periodo\_id NOT NULL;
\texorpdfstring{\softtext{fecha\_reconocimiento}}{fecha\_reconocimiento}
date NOT NULL dentro del periodo abierto;
\texorpdfstring{\softtext{registrado\_por/en.}}{registrado\_por/en.}

\subsubsection{ENT-25 pago}\label{ent-25-pago}

\begin{itemize}
\item
  \textbf{Columna:} parte\_id. \textbf{Tipo:} uuid. \textbf{Nulo:} no.
  \textbf{Nota:} .
\item
  \textbf{Columna:} tipo. \textbf{Tipo:}
  \texorpdfstring{\softtext{varchar(10).}}{varchar(10).} \textbf{Nulo:}
  no. \textbf{Nota:} COBRO (entra), PAGO (sale); naturaleza separada
  \texorpdfstring{\softtext{NORMAL/REEMBOLSO.}}{NORMAL/REEMBOLSO.}
\item
  \textbf{Columna:} moneda. \textbf{Tipo:} char(3). \textbf{Nulo:} no.
  \textbf{Nota:} MXN.
\item
  \textbf{Columna:} importe. \textbf{Tipo:}
  \texorpdfstring{\softtext{numeric(18,2).}}{numeric(18,2).}
  \textbf{Nulo:} no. \textbf{Nota:} \textgreater{} 0.
\item
  \textbf{Columna:} metodo. \textbf{Tipo:}
  \texorpdfstring{\softtext{varchar(14).}}{varchar(14).} \textbf{Nulo:}
  no. \textbf{Nota:} TRANSFERENCIA, EFECTIVO, CHEQUE, OTRO.
\item
  \textbf{Columna:} \texorpdfstring{\softtext{referencia.}}{referencia.}
  \textbf{Tipo:} \texorpdfstring{\softtext{varchar(60).}}{varchar(60).}
  \textbf{Nulo:} sí. \textbf{Nota:} .
\item
  \textbf{Columna:}
  \texorpdfstring{\softtext{recibido\_el.}}{recibido\_el.}
  \textbf{Tipo:} date. \textbf{Nulo:} no. \textbf{Nota:} .
\item
  \textbf{Columna:} estado. \textbf{Tipo:}
  \texorpdfstring{\softtext{varchar(10).}}{varchar(10).} \textbf{Nulo:}
  no. \textbf{Nota:} REGISTRADO, VALIDADO, REVERTIDO.
\item
  \textbf{Columna:} \texorpdfstring{\softtext{naturaleza.}}{naturaleza.}
  \textbf{Tipo:} \texorpdfstring{\softtext{varchar(10).}}{varchar(10).}
  \textbf{Nulo:} no. \textbf{Nota:} NORMAL / REEMBOLSO; no sustituye
  dirección \texorpdfstring{\softtext{COBRO/PAGO.}}{COBRO/PAGO.}
\item
  \textbf{Columna:}
  \texorpdfstring{\softtext{periodo\_id.}}{periodo\_id.} \textbf{Tipo:}
  uuid. \textbf{Nulo:} sí. \textbf{Nota:} NULL en registro; obligatorio
  al validar y para reembolso.
\item
  \textbf{Columna:}
  \texorpdfstring{\softtext{fecha\_reconocimiento.}}{fecha\_reconocimiento.}
  \textbf{Tipo:} date. \textbf{Nulo:} sí. \textbf{Nota:} NULL en
  registro; dentro del periodo al validar.
\item
  \textbf{Columna:}
  \texorpdfstring{\softtext{pago\_origen\_id.}}{pago\_origen\_id.}
  \textbf{Tipo:} uuid. \textbf{Nulo:} sí. \textbf{Nota:} NOT NULL ⇔
  naturaleza REEMBOLSO.
\end{itemize}

Rechazado: importe 0; REEMBOLSO sin origen.

\subsubsection{\texorpdfstring{ENT-26
\texorpdfstring{\softtext{aplicacion\_pago}}{aplicacion\_pago}}{ENT-26 }}\label{ent-26-aplicacion_pago}

pago\_id, \texorpdfstring{\softtext{liquidacion\_id}}{liquidacion\_id}
NOT NULL; importe numeric(18,2) \textgreater{} 0; naturaleza
(APLICACION, REVERSO);
\texorpdfstring{\softtext{reverso\_de\_id}}{reverso\_de\_id} UNIQUE
NULL; \texorpdfstring{\softtext{aplicada\_en;}}{aplicada\_en;}
\texorpdfstring{\softtext{registrado\_por;}}{registrado\_por;}
periodo\_id uuid y
\texorpdfstring{\softtext{fecha\_reconocimiento}}{fecha\_reconocimiento}
date NOT NULL. El registro de pago no cambia cartera;
\texorpdfstring{\softtext{validación/aplicación/reverso/reembolso}}{validación/aplicación/reverso/reembolso}
identifican reconocimiento
\texorpdfstring{\softtext{explícito.}}{explícito.} Índices
\texorpdfstring{\softcode{(E,S,liquidacion\_id)}}{(E,S,liquidacion\_id)},
\texorpdfstring{\softcode{(E,S,pago\_id)}}{(E,S,pago\_id)}.

\subsubsection{ENT-27 documento}\label{ent-27-documento}

\begin{itemize}
\item
  \textbf{Columna:}
  \texorpdfstring{\softtext{tipo\_documental.}}{tipo\_documental.}
  \textbf{Tipo:} \texorpdfstring{\softtext{varchar(24).}}{varchar(24).}
  \textbf{Nulo:} no. \textbf{Nota:}
  \texorpdfstring{\softtext{TICKET\_BASCULA,}}{TICKET\_BASCULA,}
  CONTRATO,
  \texorpdfstring{\softtext{GUIA\_MOVILIZACION,}}{GUIA\_MOVILIZACION,}
  \texorpdfstring{\softtext{COMPROBANTE\_PAGO,}}{COMPROBANTE\_PAGO,}
  IDENTIFICACION, OTRO.
\item
  \textbf{Columna:}
  \texorpdfstring{\softtext{clave\_almacenamiento.}}{clave\_almacenamiento.}
  \textbf{Tipo:} \texorpdfstring{\softtext{varchar(80).}}{varchar(80).}
  \textbf{Nulo:} no. \textbf{Nota:} aleatoria; UK; nunca ruta del
  cliente.
\item
  \textbf{Columna:} sha256. \textbf{Tipo:} char(64). \textbf{Nulo:} no.
  \textbf{Nota:}
  \texorpdfstring{\softcode{\textasciicircum{}[0-9a-f]\{64\}\$}}{\textasciicircum{}[0-9a-f]\{64\}\$}.
\item
  \textbf{Columna:}
  \texorpdfstring{\softtext{tamano\_bytes.}}{tamano\_bytes.}
  \textbf{Tipo:} bigint. \textbf{Nulo:} no. \textbf{Nota:} 1..20 MiB.
\item
  \textbf{Columna:} mime. \textbf{Tipo:}
  \texorpdfstring{\softtext{varchar(60).}}{varchar(60).} \textbf{Nulo:}
  no. \textbf{Nota:}
  \texorpdfstring{\softtext{application/pdf,}}{application/pdf,}
  \texorpdfstring{\softtext{image/jpeg,}}{image/jpeg,} image/png.
\item
  \textbf{Columna:} version. \textbf{Tipo:} smallint. \textbf{Nulo:} no.
  \textbf{Nota:} 1 + versión previa.
\item
  \textbf{Columna:}
  \texorpdfstring{\softtext{documento\_previo\_id.}}{documento\_previo\_id.}
  \textbf{Tipo:} uuid. \textbf{Nulo:} sí. \textbf{Nota:} UNIQUE.
\item
  \textbf{Columna:} emisor, folio,
  \texorpdfstring{\softtext{fecha\_emision,}}{fecha\_emision,}
  \texorpdfstring{\softtext{vigente\_hasta,}}{vigente\_hasta,}
  \texorpdfstring{\softtext{origen\_texto,}}{origen\_texto,}
  \texorpdfstring{\softtext{destino\_texto,}}{destino\_texto,}
  finalidad. \textbf{Tipo:} varios. \textbf{Nulo:} sí. \textbf{Nota:}
  según tipo (§4.2).
\item
  \textbf{Columna:}
  \texorpdfstring{\softtext{es\_sintetico.}}{es\_sintetico.}
  \textbf{Tipo:} boolean. \textbf{Nulo:} no. \textbf{Nota:} RN-045.
\end{itemize}

Sensibilidad C/P. Bytes fuera de la BD
\texorpdfstring{\softtext{(DocumentStorage).}}{(DocumentStorage).}

\subsubsection{\texorpdfstring{ENT-28
\texorpdfstring{\softtext{operacion\_documento}}{operacion\_documento}}{ENT-28 }}\label{ent-28-operacion_documento}

\texorpdfstring{\softtext{PK(E,S,operacion\_id,documento\_id);}}{PK(E,S,operacion\_id,documento\_id);}
rol \texorpdfstring{\softtext{varchar(16).}}{varchar(16).}

\subsubsection{\texorpdfstring{ENT-29
\texorpdfstring{\softtext{auditoria\_evento}}{auditoria\_evento}}{ENT-29 }}\label{ent-29-auditoria_evento}

\texorpdfstring{\softtext{ocurrido\_en;}}{ocurrido\_en;} actor\_id NULL;
comando varchar(60);
\texorpdfstring{\softtext{correlation\_id}}{correlation\_id} uuid;
entidad varchar(40); entidad\_id uuid; cambio jsonb (sólo campos
pertinentes, sin secretos). Insertado por el comando en el nodo origen;
replicado como dato; nunca generado por apply. Índice
\texorpdfstring{\softcode{(E,S,entidad,entidad\_id)}}{(E,S,entidad,entidad\_id)}.

\subsubsection{\texorpdfstring{ENT-30
\texorpdfstring{\softtext{evento\_outbox}}{evento\_outbox}}{ENT-30 }}\label{ent-30-evento_outbox}

\texorpdfstring{\softtext{tipo\_evento}}{tipo\_evento} varchar(60);
\texorpdfstring{\softtext{version\_payload}}{version\_payload} smallint;
payload jsonb;
\texorpdfstring{\softtext{propietario\_despacho}}{propietario\_despacho}
char(1); creado\_en. Sólo
\texorpdfstring{\softcode{propietario\_despacho = nodo local}}{propietario\_despacho = nodo local}
se despacha.

\subsubsection{\texorpdfstring{ENT-31
\texorpdfstring{\softtext{mensaje\_inbox}}{mensaje\_inbox}
(local)}{ENT-31  (local)}}\label{ent-31-mensaje_inbox-local}

PK(consumidor varchar(40), fuente\_id varchar(40), event\_id uuid);
\texorpdfstring{\softtext{empresa\_id,}}{empresa\_id,}
\texorpdfstring{\softtext{sucursal\_id}}{sucursal\_id} NULL;
\texorpdfstring{\softtext{recibido\_en;}}{recibido\_en;} resultado
\texorpdfstring{\softtext{varchar(10).}}{varchar(10).}

\subsubsection{\texorpdfstring{ENT-32
\texorpdfstring{\softtext{clave\_idempotencia}}{clave\_idempotencia}
(local)}{ENT-32  (local)}}\label{ent-32-clave_idempotencia-local}

\texorpdfstring{\softtext{PK(E,S,principal\_id,comando,clave}}{PK(E,S,principal\_id,comando,clave}
varchar(80)); \texorpdfstring{\softtext{hash\_peticion}}{hash\_peticion}
char(64); estado (EN\_CURSO, COMPLETADA); resultado jsonb; creada\_en;
deduplicación conservada mientras existan los hechos financieros; no
expira a los 30 días.

\subsubsection{\texorpdfstring{ENT-33
\texorpdfstring{\softtext{infra.nodo\_config}}{infra.nodo\_config}
(local)}{ENT-33  (local)}}\label{ent-33-infra.nodo_config-local}

nodo char(1) PK CHECK IN (`A',`B');
\texorpdfstring{\softtext{empresa\_id;}}{empresa\_id;}
\texorpdfstring{\softtext{descripcion.}}{descripcion.} Una fila. App
sólo SELECT.

\subsubsection{\texorpdfstring{ENT-34
\texorpdfstring{\softtext{infra.entrega\_outbox}}{infra.entrega\_outbox}
(local)}{ENT-34  (local)}}\label{ent-34-infra.entrega_outbox-local}

\texorpdfstring{\softtext{PK(evento\_id,}}{PK(evento\_id,} destino);
intentos int DEFAULT 0;
\texorpdfstring{\softtext{proximo\_intento}}{proximo\_intento}
timestamptz; \texorpdfstring{\softtext{confirmado\_en}}{confirmado\_en}
NULL; \texorpdfstring{\softtext{ultimo\_error}}{ultimo\_error} text
NULL.

\subsubsection{\texorpdfstring{ENT-35
\texorpdfstring{\softtext{infra.importacion\_fuente}}{infra.importacion\_fuente}
(local)}{ENT-35  (local)}}\label{ent-35-infra.importacion_fuente-local}

fuente varchar(20) (CSV, OCR, FDW); corte\_en; sha256;
\texorpdfstring{\softtext{filas\_leidas,}}{filas\_leidas,}
\texorpdfstring{\softtext{filas\_error}}{filas\_error} int; actor\_id.
\texorpdfstring{\softtext{UK(empresa\_id,}}{UK(empresa\_id,} fuente,
sha256) evita reimportar el mismo archivo.

\subsubsection{\texorpdfstring{ENT-37
\texorpdfstring{\softtext{periodo\_cierre}}{periodo\_cierre}}{ENT-37 }}\label{ent-37-periodo_cierre}

desde, hasta date (hasta ≥ desde); estado (ABIERTO, CERRADO);
\texorpdfstring{\softtext{version\_lock.}}{version\_lock.}
\texorpdfstring{\softcode{EXCLUDE USING gist (empresa\_id WITH =, sucursal\_id WITH =, daterange(desde, hasta, '[]') WITH \&\&)}}{EXCLUDE USING gist (empresa\_id WITH =, sucursal\_id WITH =, daterange(desde, hasta, '[]') WITH \&\&)}
(extensión \texorpdfstring{\softtext{btree\_gist,}}{btree\_gist,}
justificada: no solapamiento
\texorpdfstring{\softtext{declarativo).}}{declarativo).}

\subsubsection{\texorpdfstring{ENT-38
\texorpdfstring{\softtext{cierre\_periodo\_version}}{cierre\_periodo\_version}}{ENT-38 }}\label{ent-38-cierre_periodo_version}

\texorpdfstring{\softtext{PK(E,S,periodo\_id,numero\_version);}}{PK(E,S,periodo\_id,numero\_version);}
tipo (CIERRE, REAPERTURA); actor\_id; motivo; corte\_en;
\texorpdfstring{\softtext{hash\_conciliacion}}{hash\_conciliacion}
char(64); totales jsonb (inventario, costo de venta, pérdida, cartera).
Inmutable.

\subsubsection{\texorpdfstring{ENT-39
\texorpdfstring{\softtext{asiento\_valoracion}}{asiento\_valoracion}}{ENT-39 }}\label{ent-39-asiento_valoracion}

clase (ADQUISICION, VENTA,
\texorpdfstring{\softtext{COSTO\_TARDIO\_INVENTARIO,}}{COSTO\_TARDIO\_INVENTARIO,}
\texorpdfstring{\softtext{COSTO\_TARDIO\_VENTA,}}{COSTO\_TARDIO\_VENTA,}
\texorpdfstring{\softtext{COSTO\_TARDIO\_PERDIDA,}}{COSTO\_TARDIO\_PERDIDA,}
BAJA, \texorpdfstring{\softtext{ANULACION\_VENTA);}}{ANULACION\_VENTA);}
sentido (AUMENTO, DISMINUCION) - DISMINUCION sólo para bonificaciones de
costo tardío; importe numeric(18,2) \textgreater{} 0; columnas
\texorpdfstring{\softcode{GENERATED ALWAYS AS (...) STORED}}{GENERATED ALWAYS AS (...) STORED}:
\texorpdfstring{\softtext{delta\_inventario,}}{delta\_inventario,}
\texorpdfstring{\softtext{delta\_costo\_venta,}}{delta\_costo\_venta,}
\texorpdfstring{\softtext{delta\_perdida}}{delta\_perdida} según la
tabla de
\texorpdfstring{\softcode{states-and-calculations.md}}{states-and-calculations.md}
§3; \texorpdfstring{\softtext{asignacion\_id,}}{asignacion\_id,}
\texorpdfstring{\softtext{operacion\_id}}{operacion\_id} NULL;
\texorpdfstring{\softtext{fecha\_economica;}}{fecha\_economica;}
\texorpdfstring{\softtext{periodo\_id.}}{periodo\_id.} Índice
\texorpdfstring{\softcode{(E,S,posicion\_id)}}{(E,S,posicion\_id)}.

\subsubsection{\texorpdfstring{ENT-40
\texorpdfstring{\softtext{revision\_administrativa}}{revision\_administrativa}}{ENT-40 }}\label{ent-40-revision_administrativa}

\texorpdfstring{\softtext{operacion\_id}}{operacion\_id} XOR
\texorpdfstring{\softtext{identificacion\_id;}}{identificacion\_id;}
criterio varchar(200); resultado
\texorpdfstring{\softtext{(SIN\_OBSERVACIONES,}}{(SIN\_OBSERVACIONES,}
\texorpdfstring{\softtext{CON\_OBSERVACIONES,}}{CON\_OBSERVACIONES,}
RECHAZADA);
\texorpdfstring{\softtext{revision\_anterior\_id}}{revision\_anterior\_id}
NULL; \texorpdfstring{\softtext{responsable\_id;}}{responsable\_id;}
\texorpdfstring{\softtext{revisado\_en.}}{revisado\_en.} La UI rotula
``Revisión administrativa registrada'', nunca ``autorizado''.

\subsubsection{\texorpdfstring{ENT-41
\texorpdfstring{\softtext{revision\_documento}}{revision\_documento}}{ENT-41 }}\label{ent-41-revision_documento}

\texorpdfstring{\softtext{PK(E,S,revision\_id,documento\_id);}}{PK(E,S,revision\_id,documento\_id);}
resultado; observacion
\texorpdfstring{\softtext{varchar(500).}}{varchar(500).}

\subsubsection{\texorpdfstring{Soporte:
\texorpdfstring{\softtext{folio\_contador}}{folio\_contador} (local del
dueño,
L↔)}{Soporte:  (local del dueño, L↔)}}\label{soporte-folio_contador-local-del-dueuxf1o-l}

PK(E,S,tipo,anio); ultimo integer. Bloqueado FOR UPDATE al asignar
folio.

\subsubsection{Concordancia física durante
implementación}\label{concordancia-fuxedsica-durante-implementaciuxf3n}

Autoridad ejecutable:
\texorpdfstring{\softcode{db/migrations/}}{db/migrations/}. Se
incorporaron
\texorpdfstring{\softcode{periodo\_id}}{periodo\_id}/\texorpdfstring{\softcode{fecha\_reconocimiento}}{fecha\_reconocimiento}
en gasto, pago y aplicación, y fecha de reconocimiento en ajuste
(ADR-015). Pago admite ambos nulos mientras está REGISTRADO; validación
exige ambos y comprueba periodo abierto. Gasto y aplicación los
conservan no nulos. Son tres FK de periodo adicionales: metadata 42
tablas/83 FK. Sesiones JDBC son dos tablas técnicas locales, fuera del
MER de negocio y de réplica.

\texorpdfstring{\softcode{python3 tools/esquema/comparar\_catalogo.py}}{python3 tools/esquema/comparar\_catalogo.py}
contrasta tipos, columnas, nullabilidad, PK, UK declaradas y FK contra
\texorpdfstring{\softtext{PostgreSQL.}}{PostgreSQL.} No sustituye
pruebas de CHECK, índices parciales, triggers, permisos ni revisión
\texorpdfstring{\softtext{conceptual.}}{conceptual.}

\section{Dependencias funcionales y
normalización}\label{dependencias-funcionales-y-normalizaciuxf3n}

\textbf{Fuente:}
\texorpdfstring{\softcode{docs/database/normalization.md}}{docs/database/normalization.md}.

Referencia: F15. Estado: análisis de las 42 tablas lógicas, contrastadas
en catálogo A/B; restricciones no expresadas como DF se prueban por
separado. Las tablas técnicas adicionales
\texorpdfstring{\softtext{JDBC/Flyway/captura}}{JDBC/Flyway/captura} no
son entidades \texorpdfstring{\softtext{comerciales.}}{comerciales.}

\subsubsection{1. Dependencias funcionales y claves candidatas por
tabla}\label{dependencias-funcionales-y-claves-candidatas-por-tabla}

Notación: X → Y. ``CK'' = claves
\texorpdfstring{\softtext{candidatas.}}{candidatas.} Las claves técnicas
conservan el scope E/S de su UNIQUE o PK: una UUID aislada no se presume
única. Se listan además las claves
\texorpdfstring{\softtext{comerciales.}}{comerciales.}

\begin{itemize}
\item
  \textbf{Tabla:} empresa. \textbf{CK:} \{id\}. \textbf{DF no triviales
  relevantes:} id →
  \texorpdfstring{\softtext{razon\_social,}}{razon\_social,} rfc,
  estado. \textbf{Forma:} BCNF.
\item
  \textbf{Tabla:} sucursal. \textbf{CK:} \{E,id\},
  \texorpdfstring{\softtext{\{E,codigo\}.}}{\{E,codigo\}.} \textbf{DF no
  triviales relevantes:} E,id → codigo, nombre,
  \texorpdfstring{\softtext{nodo\_dueno;}}{nodo\_dueno;} E,codigo → id.
  \textbf{Forma:} BCNF (todo determinante es CK).
\item
  \textbf{Tabla:} parte. \textbf{CK:} \{E,id\}, \{E,rfc\} (si rfc no
  nulo). \textbf{DF no triviales relevantes:} E,id → nombre, rfc,
  \ldots. \textbf{Forma:} BCNF.
\item
  \textbf{Tabla:} predio. \textbf{CK:} \{E,id\}. \textbf{DF no triviales
  relevantes:} E,id → nombre,
  \texorpdfstring{\softtext{parte\_responsable\_id.}}{parte\_responsable\_id.}
  \textbf{Forma:} BCNF.
\item
  \textbf{Tabla:} usuario. \textbf{CK:} \{E,id\},
  \texorpdfstring{\softtext{\{E,login\_normalizado\}.}}{\{E,login\_normalizado\}.}
  \textbf{DF no triviales relevantes:} -. \textbf{Forma:} BCNF.
\item
  \textbf{Tabla:}
  \texorpdfstring{\softtext{membresia\_usuario.}}{membresia\_usuario.}
  \textbf{CK:}
  \texorpdfstring{\softtext{\{E,S,usuario\_id,rol\_codigo\}.}}{\{E,S,usuario\_id,rol\_codigo\}.}
  \textbf{DF no triviales relevantes:} → vigente. \textbf{Forma:} BCNF.
\item
  \textbf{Tabla:} animal. \textbf{CK:} \{E,id\} (UK(E,sede,id) es
  superclave usada para FK). \textbf{DF no triviales relevantes:} E,id →
  \texorpdfstring{\softtext{sede\_operadora\_id,}}{sede\_operadora\_id,}
  sexo, situacion. \textbf{Forma:} BCNF.
\item
  \textbf{Tabla:}
  \texorpdfstring{\softtext{animal\_identificacion.}}{animal\_identificacion.}
  \textbf{CK:} \{E,id\},
  \texorpdfstring{\softtext{\{E,sistema,codigo\_normalizado\}.}}{\{E,sistema,codigo\_normalizado\}.}
  \textbf{DF no triviales relevantes:}
  \texorpdfstring{\softtext{codigo\_declarado}}{codigo\_declarado} →
  \texorpdfstring{\softtext{codigo\_normalizado}}{codigo\_normalizado}
  (función determinista, ver §4). \textbf{Forma:} Desnormalización
  derivada declarada; no afirmar 3FN/BCNF sin tratar
  \texorpdfstring{\softtext{codigo\_declarado}}{codigo\_declarado} →
  \texorpdfstring{\softtext{codigo\_normalizado.}}{codigo\_normalizado.}
\item
  \textbf{Tabla:} operacion. \textbf{CK:} \{E,S,id\},
  \texorpdfstring{\softtext{\{E,S,tipo,folio\}.}}{\{E,S,tipo,folio\}.}
  \textbf{DF no triviales relevantes:} E,S,id → tipo, parte\_id, estado,
  \ldots; E,S,tipo,folio → id. \textbf{Forma:} BCNF.
\item
  \textbf{Tabla:}
  \texorpdfstring{\softtext{lote\_operacion.}}{lote\_operacion.}
  \textbf{CK:}
  \texorpdfstring{\softtext{\{E,S,operacion\_id,numero\_lote\}.}}{\{E,S,operacion\_id,numero\_lote\}.}
  \textbf{DF no triviales relevantes:} →
  \texorpdfstring{\softtext{descripcion.}}{descripcion.} \textbf{Forma:}
  BCNF.
\item
  \textbf{Tabla:}
  \texorpdfstring{\softtext{contrato\_version.}}{contrato\_version.}
  \textbf{CK:}
  \texorpdfstring{\softtext{\{E,S,operacion\_id,numero\_version\}.}}{\{E,S,operacion\_id,numero\_version\}.}
  \textbf{DF no triviales relevantes:} → estado,
  \texorpdfstring{\softtext{politica\_redondeo.}}{politica\_redondeo.}
  \textbf{Forma:} BCNF.
\item
  \textbf{Tabla:}
  \texorpdfstring{\softtext{condicion\_lote.}}{condicion\_lote.}
  \textbf{CK:}
  \texorpdfstring{\softtext{\{E,S,operacion\_id,numero\_version,numero\_lote\}.}}{\{E,S,operacion\_id,numero\_version,numero\_lote\}.}
  \textbf{DF no triviales relevantes:} →
  \texorpdfstring{\softtext{precio\_unitario,}}{precio\_unitario,}
  \texorpdfstring{\softtext{unidad\_precio,}}{unidad\_precio,} \ldots.
  \textbf{Forma:} BCNF.
\item
  \textbf{Tabla:}
  \texorpdfstring{\softtext{seleccion\_compra.}}{seleccion\_compra.}
  \textbf{CK:} \{E,S,op,animal\}; PK física \{E,S,op,lote,animal\} no
  mínima. \textbf{DF no triviales relevantes:} E,S,op,animal → lote.
  \textbf{Forma:} BCNF respecto a la DF declarada: el determinante
  mínimo es
  \texorpdfstring{\softtext{\{E,S,op,animal\}.}}{\{E,S,op,animal\}.}
\item
  \textbf{Tabla:}
  \texorpdfstring{\softtext{posicion\_inventario.}}{posicion\_inventario.}
  \textbf{CK:} \{E,S,id\},
  \texorpdfstring{\softtext{\{E,S,animal\_id,numero\_ciclo\},}}{\{E,S,animal\_id,numero\_ciclo\},}
  \texorpdfstring{\softtext{\{E,S,compra\_op,compra\_lote,animal\_id\}.}}{\{E,S,compra\_op,compra\_lote,animal\_id\}.}
  \textbf{DF no triviales relevantes:} id → estado, \ldots;
  \texorpdfstring{\softtext{animal\_id,numero\_ciclo}}{animal\_id,numero\_ciclo}
  → id. \textbf{Forma:} BCNF.
\item
  \textbf{Tabla:}
  \texorpdfstring{\softtext{reserva\_venta.}}{reserva\_venta.}
  \textbf{CK:} \{E,S,id\}. \textbf{DF no triviales relevantes:} id →
  \texorpdfstring{\softtext{posicion\_id,}}{posicion\_id,} estado,
  vence\_en. \textbf{Forma:} BCNF.
\item
  \textbf{Tabla:} pesaje. \textbf{CK:} \{E,S,id\}. \textbf{DF no
  triviales relevantes:} id → modo, bruto, tara, neto; (bruto,tara) →
  neto \textbf{cuando modo = BRUTO\_TARA}. \textbf{Forma:} Dependencia
  condicional por modo; neto conservado con CHECK, no prueba automática
  de 3FN.
\item
  \textbf{Tabla:}
  \texorpdfstring{\softtext{pesaje\_animal.}}{pesaje\_animal.}
  \textbf{CK:}
  \texorpdfstring{\softtext{\{E,S,pesaje\_id,animal\_id\}.}}{\{E,S,pesaje\_id,animal\_id\}.}
  \textbf{DF no triviales relevantes:} →
  \texorpdfstring{\softtext{peso\_individual\_kg.}}{peso\_individual\_kg.}
  \textbf{Forma:} BCNF.
\item
  \textbf{Tabla:} \texorpdfstring{\softtext{liquidacion.}}{liquidacion.}
  \textbf{CK:} \{E,S,id\},
  \texorpdfstring{\softtext{\{E,S,operacion\_id\}.}}{\{E,S,operacion\_id\}.}
  \textbf{DF no triviales relevantes:} id → total, snapshots.
  \textbf{Forma:} BCNF (snapshots dependen de la liquidación, ver §5).
\item
  \textbf{Tabla:}
  \texorpdfstring{\softtext{concepto\_liquidacion.}}{concepto\_liquidacion.}
  \textbf{CK:}
  \texorpdfstring{\softtext{\{E,S,liquidacion\_id,numero\_concepto\}.}}{\{E,S,liquidacion\_id,numero\_concepto\}.}
  \textbf{DF no triviales relevantes:} → tipo, cantidad, precio,
  importe; (cantidad, precio) → importe salvo redondeo. \textbf{Forma:}
  Importe derivado materializado: excepción a normalización estricta,
  hecho redondeado inmutable.
\item
  \textbf{Tabla:} pago. \textbf{CK:} \{E,S,id\}. \textbf{DF no triviales
  relevantes:} id → parte\_id, importe, tipo. \textbf{Forma:} BCNF.
\item
  \textbf{Tabla:}
  \texorpdfstring{\softtext{aplicacion\_pago.}}{aplicacion\_pago.}
  \textbf{CK:} \{E,S,id\},
  \texorpdfstring{\softtext{\{E,S,reverso\_de\_id\}}}{\{E,S,reverso\_de\_id\}}
  (si no nulo). \textbf{DF no triviales relevantes:} id → pago\_id,
  \texorpdfstring{\softtext{liquidacion\_id,}}{liquidacion\_id,}
  importe. \textbf{Forma:} BCNF.
\item
  \textbf{Tabla:}
  \texorpdfstring{\softtext{asignacion\_costo.}}{asignacion\_costo.}
  \textbf{CK:} \{E,S,id\}. \textbf{DF no triviales relevantes:} id →
  posicion, fuente, importe. \textbf{Forma:} BCNF.
\item
  \textbf{Tabla:}
  \texorpdfstring{\softtext{asiento\_valoracion.}}{asiento\_valoracion.}
  \textbf{CK:} \{E,S,id\}. \textbf{DF no triviales relevantes:} (clase,
  sentido, importe) → delta\_*. \textbf{Forma:} deltas generados (no
  redundancia editable).
\item
  \textbf{Tabla:} documento. \textbf{CK:} \{E,S,id\},
  \texorpdfstring{\softtext{\{clave\_almacenamiento\}.}}{\{clave\_almacenamiento\}.}
  \textbf{DF no triviales relevantes:} id → sha256, tamano, \ldots.
  \textbf{Forma:} BCNF.
\item
  \textbf{Tabla:}
  \texorpdfstring{\softtext{periodo\_cierre.}}{periodo\_cierre.}
  \textbf{CK:} \{E,S,id\}; además (E,S,fecha) determina el periodo por
  la \texorpdfstring{\softtext{exclusión.}}{exclusión.} \textbf{DF no
  triviales relevantes:} -. \textbf{Forma:} BCNF.
\end{itemize}

\subsubsection{2. Ejemplo A - operación /
contraparte}\label{ejemplo-a---operaciuxf3n-contraparte}

\textbf{No normalizado} (planilla típica de un comercializador):

\begin{itemize}
\item
  \textbf{folio:} C-01. \textbf{fecha:}
  \texorpdfstring{\softtext{2026-10-02.}}{2026-10-02.}
  \textbf{contraparte:} Productor Demo 01. \textbf{rfc:}
  \texorpdfstring{\softtext{DEMO010101AAA.}}{DEMO010101AAA.}
  \textbf{teléfono:} 555-0101. \textbf{lote:} 1. \textbf{arete:} DEMO-1,
  DEMO-2. \textbf{precio\_kg:} 48.
\item
  \textbf{1FN}: ``DEMO-1, DEMO-2'' es no atómico → una fila por arete:
  (folio, lote, arete, \ldots). La clave pasa a \{folio, arete\}.
\item
  \textbf{2FN}: con clave \{folio, arete\},
  \texorpdfstring{\softcode{fecha}}{fecha},
  \texorpdfstring{\softcode{contraparte}}{contraparte},
  \texorpdfstring{\softcode{rfc}}{rfc},
  \texorpdfstring{\softcode{teléfono}}{teléfono} dependen sólo de
  \texorpdfstring{\softcode{folio}}{folio} (dependencia parcial) → se
  separa
  \texorpdfstring{\softcode{operacion(folio, fecha, parte)}}{operacion(folio, fecha, parte)}
  de
  \texorpdfstring{\softcode{seleccion(folio, lote, arete)}}{seleccion(folio, lote, arete)}.
\item
  \textbf{3FN}: en \texorpdfstring{\softcode{operacion}}{operacion},
  \texorpdfstring{\softcode{rfc}}{rfc} y
  \texorpdfstring{\softcode{teléfono}}{teléfono} dependen de
  \texorpdfstring{\softcode{contraparte}}{contraparte}, no del folio
  (transitiva folio → contraparte → rfc) →
  \texorpdfstring{\softcode{parte(id, nombre, rfc, telefono)}}{parte(id, nombre, rfc, telefono)}
  y
  \texorpdfstring{\softcode{operacion.parte\_id}}{operacion.parte\_id}.
\item
  \texorpdfstring{\softcode{precio\_kg}}{precio\_kg} depende de (folio,
  lote, \textbf{versión} aceptada), no del arete →
  \texorpdfstring{\softcode{condicion\_lote}}{condicion\_lote} (ejemplo
  B).
\item
  Resultado: \texorpdfstring{\softcode{parte}}{parte},
  \texorpdfstring{\softcode{operacion}}{operacion},
  \texorpdfstring{\softcode{lote\_operacion}}{lote\_operacion},
  \texorpdfstring{\softcode{seleccion\_compra}}{seleccion\_compra},
  \texorpdfstring{\softcode{animal\_identificacion}}{animal\_identificacion}
  (el arete identifica al animal, no a la
  \texorpdfstring{\softtext{selección).}}{selección).}
\end{itemize}

\subsubsection{3. Ejemplo B - versión / condición por
lote}\label{ejemplo-b---versiuxf3n-condiciuxf3n-por-lote}

Planteamiento ingenuo:
\texorpdfstring{\softcode{operacion(id, precio\_kg, deduccion, version)}}{operacion(id, precio\_kg, deduccion, version)}
con UPDATE al \texorpdfstring{\softtext{renegociar.}}{renegociar.}

\begin{itemize}
\tightlist
\item
  Viola el hecho histórico: la versión 1 deja de existir (no es un
  problema de forma normal, es de modelado temporal).
\item
  Al versionar:
  \texorpdfstring{\softcode{condicion(operacion, version, lote, precio, unidad, deduccion, politica\_redondeo)}}{condicion(operacion, version, lote, precio, unidad, deduccion, politica\_redondeo)}.
  DF: (op, version) →
  \texorpdfstring{\softtext{politica\_redondeo;}}{politica\_redondeo;}
  (op, version, lote) → precio, unidad, deduccion.
  \texorpdfstring{\softcode{politica\_redondeo}}{politica\_redondeo}
  depende de parte de la clave → \textbf{viola 2FN} → se mueve a
  \texorpdfstring{\softcode{contrato\_version(op, version, politica\_redondeo, estado)}}{contrato\_version(op, version, politica\_redondeo, estado)}.
\item
  \texorpdfstring{\softcode{condicion\_lote(op, version, lote, precio, unidad, base\_pesaje, deduccion)}}{condicion\_lote(op, version, lote, precio, unidad, base\_pesaje, deduccion)}:
  toda DF tiene como determinante la clave completa → 3FN y BCNF.
\item
  Preservación de dependencias: las dos DF quedan cada una en una tabla;
  la unión
  \texttt{contrato\_version }\ensuremath{\bowtie}\texttt{ condicion\_lote}
  sobre (op, version) es sin pérdida porque (op, version) es clave de
  \texorpdfstring{\softcode{contrato\_version}}{contrato\_version}.
\end{itemize}

\subsubsection{4. Ejemplo C - liquidación /
aplicaciones}\label{ejemplo-c---liquidaciuxf3n-aplicaciones}

Planteamiento ingenuo:
\texorpdfstring{\softcode{liquidacion(id, total, pagado, saldo, pago1\_ref, pago1\_importe, pago2\_ref, …)}}{liquidacion(id, total, pagado, saldo, pago1\_ref, pago1\_importe, pago2\_ref, …)}.

\begin{itemize}
\tightlist
\item
  \textbf{1FN}: grupo repetido
  \texorpdfstring{\softtext{pago1/pago2}}{pago1/pago2} → tabla aparte.
\item
  \texorpdfstring{\softcode{saldo}}{saldo} = total + ajustes −
  aplicaciones: atributo derivado editable en paralelo → riesgo de
  inconsistencia; se elimina y se calcula en vista
  (\texorpdfstring{\softcode{v\_saldo\_liquidacion}}{v\_saldo\_liquidacion}).
\item
  Un pago cubre varias liquidaciones y una liquidación recibe varios
  pagos (N:M) →
  \texorpdfstring{\softcode{pago(id, parte, importe, …)}}{pago(id, parte, importe, …)}
  +
  \texorpdfstring{\softcode{aplicacion\_pago(id, pago\_id, liquidacion\_id, importe, naturaleza, reverso\_de\_id)}}{aplicacion\_pago(id, pago\_id, liquidacion\_id, importe, naturaleza, reverso\_de\_id)}.
\item
  DF en \texorpdfstring{\softcode{aplicacion\_pago}}{aplicacion\_pago}:
  id → (pago\_id,
  \texorpdfstring{\softtext{liquidacion\_id,}}{liquidacion\_id,}
  importe). No hay DF (pago\_id,
  \texorpdfstring{\softtext{liquidacion\_id)}}{liquidacion\_id)} →
  importe porque puede haber varias aplicaciones y reversos entre el
  mismo par: por eso la aplicación tiene identidad propia.
\item
  Restricción entre filas Σ aplicaciones ≤ importe del pago: no es una
  DF ni puede ser CHECK (F20); se garantiza en el comando
  \texorpdfstring{\softtext{AplicarPago}}{AplicarPago} con locks.
\item
  Resultado: BCNF; ninguna sumatoria
  \texorpdfstring{\softtext{almacenada.}}{almacenada.}
\end{itemize}

\subsubsection{5. Excepciones
deliberadas}\label{excepciones-deliberadas}

{\def\LTcaptype{none} % do not increment counter
\begin{longtable}[]{@{}
  >{\raggedright\arraybackslash}p{(\linewidth - 2\tabcolsep) * \real{0.5000}}
  >{\raggedright\arraybackslash}p{(\linewidth - 2\tabcolsep) * \real{0.5000}}@{}}
\toprule\noalign{}
\begin{minipage}[b]{\linewidth}\raggedright
Caso
\end{minipage} & \begin{minipage}[b]{\linewidth}\raggedright
Por qué no es redundancia a eliminar
\end{minipage} \\
\midrule\noalign{}
\endhead
\bottomrule\noalign{}
\endlastfoot
\texorpdfstring{\softcode{liquidacion.snapshot\_*}}{liquidacion.snapshot\_*}
& Hecho temporal ``así se aceptó''; depende de la liquidación, no de la
parte actual (RN-024) \\
\texorpdfstring{\softcode{concepto\_liquidacion.importe}}{concepto\_liquidacion.importe}
& Valor redondeado HALF\_UP que se acordó; recalcular podría diferir por
política futura \\
\texorpdfstring{\softcode{seleccion\_compra.tipo\_operacion}}{seleccion\_compra.tipo\_operacion}
& Constante para FK compuesta que garantiza ``sólo COMPRA''
declarativamente \\
\texorpdfstring{\softcode{pesaje.neto\_kg}}{pesaje.neto\_kg} con
bruto/tara & En
\texorpdfstring{\softtext{NETO\_REPORTADO}}{NETO\_REPORTADO} es dato
primario; en BRUTO\_TARA se fuerza igualdad con CHECK \\
\texorpdfstring{\softcode{animal\_identificacion.codigo\_normalizado}}{animal\_identificacion.codigo\_normalizado}
& Necesario para la UNIQUE; determinado por función inmutable sobre
\texorpdfstring{\softcode{codigo\_declarado}}{codigo\_declarado} (se
podría generar; se mantiene CHECK de igualdad) \\
\texorpdfstring{\softcode{asiento\_valoracion.delta\_*}}{asiento\_valoracion.delta\_*}
& Columnas generadas: no editables \\
\end{longtable}
}

\subsubsection{6. BCNF y preservación de
dependencias}\label{bcnf-y-preservaciuxf3n-de-dependencias}

La selección de compra tiene la clave candidata mínima
\{E,S,op,animal\}; añadir lote produce una superclave no mínima. Su DF
determina lote y no viola BCNF. La tabla anterior analiza las
dependencias declaradas de las tablas enumeradas; no demuestra BCNF del
modelo completo. El inventario completo siguiente explicita
determinantes mínimos conocidos y excepciones; las restricciones físicas
se contrastaron con el catálogo ejecutado, sin confundir una comparación
con una demostración exhaustiva de reglas (TEST-30). Los totales
materializados y snapshots conservan autoridad e inmutabilidad
explícitas; no son campos editables paralelos.

\subsubsection{7. Inventario completo de determinantes
conocidos}\label{inventario-completo-de-determinantes-conocidos}

Para cada fila, K → todos los atributos propios del registro. Las claves
de abajo provienen de PK/UK completas declaradas, eliminando superclaves
que contienen otra clave más pequeña. UK con NULL y unicidad parcial por
estado no son claves candidatas del conjunto completo. Todas las DF
abreviadas id → \ldots{} de la tabla inicial significan la clave
E{[},S{]},id correspondiente, nunca UUID global.

Los datos
\texorpdfstring{\softtext{escalares/composición}}{escalares/composición}
separados sostienen 1FN; los atributos comerciales de
\texorpdfstring{\softtext{contraparte/animal}}{contraparte/animal} se
separan de selecciones y hechos (2FN/3FN). Las constantes para FK de
tipo, valores calculados, JSON snapshots y efectos históricos son
excepciones \texorpdfstring{\softtext{explícitas.}}{explícitas.} Una
comprobación de claves no descubre automáticamente todas las DF
semánticas; este inventario analiza las reconocidas por el dominio, sin
una proclamación global falsa de BCNF.

\begin{itemize}
\item
  \textbf{Tabla:} empresa. \textbf{Determinantes mínimos NOT NULL:}
  \{id\}. \textbf{UK condicional conocida:} -. \textbf{Análisis /
  excepción:} \texorpdfstring{\softtext{1FN/2FN/3FN}}{1FN/2FN/3FN} y
  BCNF respecto a DF declaradas: determinantes son claves; no se presume
  otra unicidad comercial.
\item
  \textbf{Tabla:} sucursal. \textbf{Determinantes mínimos NOT NULL:}
  \texorpdfstring{\softtext{\{empresa\_id,id\},}}{\{empresa\_id,id\},}
  \texorpdfstring{\softtext{\{empresa\_id,codigo\}.}}{\{empresa\_id,codigo\}.}
  \textbf{UK condicional conocida:} -. \textbf{Análisis / excepción:}
  \texorpdfstring{\softtext{1FN/2FN/3FN}}{1FN/2FN/3FN} y BCNF respecto a
  DF declaradas: determinantes son claves; no se presume otra unicidad
  comercial.
\item
  \textbf{Tabla:}
  \texorpdfstring{\softtext{rol\_aplicacion.}}{rol\_aplicacion.}
  \textbf{Determinantes mínimos NOT NULL:} \{codigo\}. \textbf{UK
  condicional conocida:} -. \textbf{Análisis / excepción:}
  \texorpdfstring{\softtext{1FN/2FN/3FN}}{1FN/2FN/3FN} y BCNF respecto a
  DF declaradas: determinantes son claves; no se presume otra unicidad
  comercial.
\item
  \textbf{Tabla:} usuario. \textbf{Determinantes mínimos NOT NULL:}
  \texorpdfstring{\softtext{\{empresa\_id,id\},}}{\{empresa\_id,id\},}
  \texorpdfstring{\softtext{\{empresa\_id,login\_normalizado\}.}}{\{empresa\_id,login\_normalizado\}.}
  \textbf{UK condicional conocida:} -. \textbf{Análisis / excepción:}
  \texorpdfstring{\softtext{1FN/2FN/3FN}}{1FN/2FN/3FN} y BCNF respecto a
  DF declaradas: determinantes son claves; no se presume otra unicidad
  comercial.
\item
  \textbf{Tabla:}
  \texorpdfstring{\softtext{membresia\_usuario.}}{membresia\_usuario.}
  \textbf{Determinantes mínimos NOT NULL:}
  \texorpdfstring{\softtext{\{empresa\_id,sucursal\_id,usuario\_id,rol\_codigo\}.}}{\{empresa\_id,sucursal\_id,usuario\_id,rol\_codigo\}.}
  \textbf{UK condicional conocida:} -. \textbf{Análisis / excepción:}
  \texorpdfstring{\softtext{1FN/2FN/3FN}}{1FN/2FN/3FN} y BCNF respecto a
  DF declaradas: determinantes son claves; no se presume otra unicidad
  comercial.
\item
  \textbf{Tabla:} parte. \textbf{Determinantes mínimos NOT NULL:}
  \texorpdfstring{\softtext{\{empresa\_id,id\}.}}{\{empresa\_id,id\}.}
  \textbf{UK condicional conocida:} -. \textbf{Análisis / excepción:}
  \texorpdfstring{\softtext{1FN/2FN/3FN}}{1FN/2FN/3FN} y BCNF respecto a
  DF declaradas: determinantes son claves; no se presume otra unicidad
  comercial.
\item
  \textbf{Tabla:} predio. \textbf{Determinantes mínimos NOT NULL:}
  \texorpdfstring{\softtext{\{empresa\_id,id\}.}}{\{empresa\_id,id\}.}
  \textbf{UK condicional conocida:} -. \textbf{Análisis / excepción:}
  \texorpdfstring{\softtext{1FN/2FN/3FN}}{1FN/2FN/3FN} y BCNF respecto a
  DF declaradas: determinantes son claves; no se presume otra unicidad
  comercial.
\item
  \textbf{Tabla:} animal. \textbf{Determinantes mínimos NOT NULL:}
  \texorpdfstring{\softtext{\{empresa\_id,id\}.}}{\{empresa\_id,id\}.}
  \textbf{UK condicional conocida:} -. \textbf{Análisis / excepción:}
  \texorpdfstring{\softtext{1FN/2FN/3FN}}{1FN/2FN/3FN} y BCNF respecto a
  DF declaradas: determinantes son claves; no se presume otra unicidad
  comercial.
\item
  \textbf{Tabla:}
  \texorpdfstring{\softtext{animal\_identificacion.}}{animal\_identificacion.}
  \textbf{Determinantes mínimos NOT NULL:}
  \texorpdfstring{\softtext{\{empresa\_id,id\},}}{\{empresa\_id,id\},}
  \texorpdfstring{\softtext{\{empresa\_id,sistema,codigo\_normalizado\}.}}{\{empresa\_id,sistema,codigo\_normalizado\}.}
  \textbf{UK condicional conocida:} -. \textbf{Análisis / excepción:}
  normalización derivada del código; véase §5, no BCNF estricta.
\item
  \textbf{Tabla:} operacion. \textbf{Determinantes mínimos NOT NULL:}
  \texorpdfstring{\softtext{\{empresa\_id,sucursal\_id,id\},}}{\{empresa\_id,sucursal\_id,id\},}
  \texorpdfstring{\softtext{\{empresa\_id,sucursal\_id,tipo,folio\}.}}{\{empresa\_id,sucursal\_id,tipo,folio\}.}
  \textbf{UK condicional conocida:} -. \textbf{Análisis / excepción:}
  \texorpdfstring{\softtext{1FN/2FN/3FN}}{1FN/2FN/3FN} y BCNF respecto a
  DF declaradas: determinantes son claves; no se presume otra unicidad
  comercial.
\item
  \textbf{Tabla:}
  \texorpdfstring{\softtext{folio\_contador.}}{folio\_contador.}
  \textbf{Determinantes mínimos NOT NULL:}
  \texorpdfstring{\softtext{\{empresa\_id,sucursal\_id,tipo,anio\}.}}{\{empresa\_id,sucursal\_id,tipo,anio\}.}
  \textbf{UK condicional conocida:} -. \textbf{Análisis / excepción:}
  \texorpdfstring{\softtext{1FN/2FN/3FN}}{1FN/2FN/3FN} y BCNF respecto a
  DF declaradas: determinantes son claves; no se presume otra unicidad
  comercial.
\item
  \textbf{Tabla:}
  \texorpdfstring{\softtext{lote\_operacion.}}{lote\_operacion.}
  \textbf{Determinantes mínimos NOT NULL:}
  \texorpdfstring{\softtext{\{empresa\_id,sucursal\_id,operacion\_id,numero\_lote\}.}}{\{empresa\_id,sucursal\_id,operacion\_id,numero\_lote\}.}
  \textbf{UK condicional conocida:} -. \textbf{Análisis / excepción:}
  \texorpdfstring{\softtext{1FN/2FN/3FN}}{1FN/2FN/3FN} y BCNF respecto a
  DF declaradas: determinantes son claves; no se presume otra unicidad
  comercial.
\item
  \textbf{Tabla:}
  \texorpdfstring{\softtext{contrato\_version.}}{contrato\_version.}
  \textbf{Determinantes mínimos NOT NULL:}
  \texorpdfstring{\softtext{\{empresa\_id,sucursal\_id,operacion\_id,numero\_version\}.}}{\{empresa\_id,sucursal\_id,operacion\_id,numero\_version\}.}
  \textbf{UK condicional conocida:} -. \textbf{Análisis / excepción:}
  \texorpdfstring{\softtext{1FN/2FN/3FN}}{1FN/2FN/3FN} y BCNF respecto a
  DF declaradas: determinantes son claves; no se presume otra unicidad
  comercial.
\item
  \textbf{Tabla:}
  \texorpdfstring{\softtext{condicion\_lote.}}{condicion\_lote.}
  \textbf{Determinantes mínimos NOT NULL:}
  \texorpdfstring{\softtext{\{empresa\_id,sucursal\_id,operacion\_id,numero\_version,numero\_lote\}.}}{\{empresa\_id,sucursal\_id,operacion\_id,numero\_version,numero\_lote\}.}
  \textbf{UK condicional conocida:} -. \textbf{Análisis / excepción:}
  \texorpdfstring{\softtext{1FN/2FN/3FN}}{1FN/2FN/3FN} y BCNF respecto a
  DF declaradas: determinantes son claves; no se presume otra unicidad
  comercial.
\item
  \textbf{Tabla:} pesaje. \textbf{Determinantes mínimos NOT NULL:}
  \texorpdfstring{\softtext{\{empresa\_id,sucursal\_id,id\}.}}{\{empresa\_id,sucursal\_id,id\}.}
  \textbf{UK condicional conocida:} -. \textbf{Análisis / excepción:}
  neto derivado sólo en
  \texorpdfstring{\softtext{BRUTO\_TARA;}}{BRUTO\_TARA;} dependencia
  condicional con CHECK.
\item
  \textbf{Tabla:}
  \texorpdfstring{\softtext{pesaje\_animal.}}{pesaje\_animal.}
  \textbf{Determinantes mínimos NOT NULL:}
  \texorpdfstring{\softtext{\{empresa\_id,sucursal\_id,pesaje\_id,animal\_id\}.}}{\{empresa\_id,sucursal\_id,pesaje\_id,animal\_id\}.}
  \textbf{UK condicional conocida:} -. \textbf{Análisis / excepción:}
  \texorpdfstring{\softtext{1FN/2FN/3FN}}{1FN/2FN/3FN} y BCNF respecto a
  DF declaradas: determinantes son claves; no se presume otra unicidad
  comercial.
\item
  \textbf{Tabla:}
  \texorpdfstring{\softtext{seleccion\_compra.}}{seleccion\_compra.}
  \textbf{Determinantes mínimos NOT NULL:}
  \texorpdfstring{\softtext{\{empresa\_id,sucursal\_id,operacion\_id,animal\_id\}.}}{\{empresa\_id,sucursal\_id,operacion\_id,animal\_id\}.}
  \textbf{UK condicional conocida:} -. \textbf{Análisis / excepción:}
  \texorpdfstring{\softtext{1FN/2FN/3FN}}{1FN/2FN/3FN} y BCNF respecto a
  DF declaradas: determinantes son claves; no se presume otra unicidad
  comercial.
\item
  \textbf{Tabla:}
  \texorpdfstring{\softtext{posicion\_inventario.}}{posicion\_inventario.}
  \textbf{Determinantes mínimos NOT NULL:}
  \texorpdfstring{\softtext{\{empresa\_id,sucursal\_id,id\},}}{\{empresa\_id,sucursal\_id,id\},}
  \texorpdfstring{\softtext{\{empresa\_id,sucursal\_id,animal\_id,numero\_ciclo\},}}{\{empresa\_id,sucursal\_id,animal\_id,numero\_ciclo\},}
  \texorpdfstring{\softtext{\{empresa\_id,sucursal\_id,compra\_operacion\_id,compra\_numero\_lote,animal\_id\}.}}{\{empresa\_id,sucursal\_id,compra\_operacion\_id,compra\_numero\_lote,animal\_id\}.}
  \textbf{UK condicional conocida:} -. \textbf{Análisis / excepción:}
  \texorpdfstring{\softtext{1FN/2FN/3FN}}{1FN/2FN/3FN} y BCNF respecto a
  DF declaradas: determinantes son claves; no se presume otra unicidad
  comercial.
\item
  \textbf{Tabla:}
  \texorpdfstring{\softtext{seleccion\_venta.}}{seleccion\_venta.}
  \textbf{Determinantes mínimos NOT NULL:}
  \texorpdfstring{\softtext{\{empresa\_id,sucursal\_id,operacion\_id,posicion\_id\}.}}{\{empresa\_id,sucursal\_id,operacion\_id,posicion\_id\}.}
  \textbf{UK condicional conocida:} -. \textbf{Análisis / excepción:}
  selección de posición por operación mínima; estado derivado de
  \texorpdfstring{\softtext{reserva/posición,}}{reserva/posición,} no
  atributo repetido.
\item
  \textbf{Tabla:}
  \texorpdfstring{\softtext{reserva\_venta.}}{reserva\_venta.}
  \textbf{Determinantes mínimos NOT NULL:}
  \texorpdfstring{\softtext{\{empresa\_id,sucursal\_id,id\}.}}{\{empresa\_id,sucursal\_id,id\}.}
  \textbf{UK condicional conocida:} -. \textbf{Análisis / excepción:}
  \texorpdfstring{\softtext{1FN/2FN/3FN}}{1FN/2FN/3FN} y BCNF respecto a
  DF declaradas: determinantes son claves; no se presume otra unicidad
  comercial.
\item
  \textbf{Tabla:}
  \texorpdfstring{\softtext{movimiento\_inventario.}}{movimiento\_inventario.}
  \textbf{Determinantes mínimos NOT NULL:}
  \texorpdfstring{\softtext{\{empresa\_id,sucursal\_id,id\}.}}{\{empresa\_id,sucursal\_id,id\}.}
  \textbf{UK condicional conocida:} -. \textbf{Análisis / excepción:}
  \texorpdfstring{\softtext{1FN/2FN/3FN}}{1FN/2FN/3FN} y BCNF respecto a
  DF declaradas: determinantes son claves; no se presume otra unicidad
  comercial.
\item
  \textbf{Tabla:}
  \texorpdfstring{\softtext{periodo\_cierre.}}{periodo\_cierre.}
  \textbf{Determinantes mínimos NOT NULL:}
  \texorpdfstring{\softtext{\{empresa\_id,sucursal\_id,id\}.}}{\{empresa\_id,sucursal\_id,id\}.}
  \textbf{UK condicional conocida:} -. \textbf{Análisis / excepción:}
  \texorpdfstring{\softtext{1FN/2FN/3FN}}{1FN/2FN/3FN} y BCNF respecto a
  DF declaradas: determinantes son claves; no se presume otra unicidad
  comercial.
\item
  \textbf{Tabla:}
  \texorpdfstring{\softtext{cierre\_periodo\_version.}}{cierre\_periodo\_version.}
  \textbf{Determinantes mínimos NOT NULL:}
  \texorpdfstring{\softtext{\{empresa\_id,sucursal\_id,periodo\_id,numero\_version\}.}}{\{empresa\_id,sucursal\_id,periodo\_id,numero\_version\}.}
  \textbf{UK condicional conocida:} -. \textbf{Análisis / excepción:}
  \texorpdfstring{\softtext{totales/versiones}}{totales/versiones} y
  hash de corte materializados e
  \texorpdfstring{\softtext{inmutables.}}{inmutables.}
\item
  \textbf{Tabla:} \texorpdfstring{\softtext{liquidacion.}}{liquidacion.}
  \textbf{Determinantes mínimos NOT NULL:}
  \texorpdfstring{\softtext{\{empresa\_id,sucursal\_id,id\},}}{\{empresa\_id,sucursal\_id,id\},}
  \texorpdfstring{\softtext{\{empresa\_id,sucursal\_id,operacion\_id\}.}}{\{empresa\_id,sucursal\_id,operacion\_id\}.}
  \textbf{UK condicional conocida:} -. \textbf{Análisis / excepción:}
  snapshots temporales inmutables, no copia del maestro actual.
\item
  \textbf{Tabla:}
  \texorpdfstring{\softtext{concepto\_liquidacion.}}{concepto\_liquidacion.}
  \textbf{Determinantes mínimos NOT NULL:}
  \texorpdfstring{\softtext{\{empresa\_id,sucursal\_id,liquidacion\_id,numero\_concepto\}.}}{\{empresa\_id,sucursal\_id,liquidacion\_id,numero\_concepto\}.}
  \textbf{UK condicional conocida:} -. \textbf{Análisis / excepción:}
  importe redondeado materializado; no 3FN estricta respecto a fórmula.
\item
  \textbf{Tabla:}
  \texorpdfstring{\softtext{ajuste\_liquidacion.}}{ajuste\_liquidacion.}
  \textbf{Determinantes mínimos NOT NULL:}
  \texorpdfstring{\softtext{\{empresa\_id,sucursal\_id,id\}.}}{\{empresa\_id,sucursal\_id,id\}.}
  \textbf{UK condicional conocida:} -. \textbf{Análisis / excepción:}
  \texorpdfstring{\softtext{1FN/2FN/3FN}}{1FN/2FN/3FN} y BCNF respecto a
  DF declaradas: determinantes son claves; no se presume otra unicidad
  comercial.
\item
  \textbf{Tabla:}
  \texorpdfstring{\softtext{gasto\_operacion.}}{gasto\_operacion.}
  \textbf{Determinantes mínimos NOT NULL:}
  \texorpdfstring{\softtext{\{empresa\_id,sucursal\_id,id\}.}}{\{empresa\_id,sucursal\_id,id\}.}
  \textbf{UK condicional conocida:} -. \textbf{Análisis / excepción:}
  \texorpdfstring{\softtext{1FN/2FN/3FN}}{1FN/2FN/3FN} y BCNF respecto a
  DF declaradas: determinantes son claves; no se presume otra unicidad
  comercial.
\item
  \textbf{Tabla:}
  \texorpdfstring{\softtext{asignacion\_costo.}}{asignacion\_costo.}
  \textbf{Determinantes mínimos NOT NULL:}
  \texorpdfstring{\softtext{\{empresa\_id,sucursal\_id,id\}.}}{\{empresa\_id,sucursal\_id,id\}.}
  \textbf{UK condicional conocida:} -. \textbf{Análisis / excepción:}
  \texorpdfstring{\softtext{1FN/2FN/3FN}}{1FN/2FN/3FN} y BCNF respecto a
  DF declaradas: determinantes son claves; no se presume otra unicidad
  comercial.
\item
  \textbf{Tabla:}
  \texorpdfstring{\softtext{asiento\_valoracion.}}{asiento\_valoracion.}
  \textbf{Determinantes mínimos NOT NULL:}
  \texorpdfstring{\softtext{\{empresa\_id,sucursal\_id,id\}.}}{\{empresa\_id,sucursal\_id,id\}.}
  \textbf{UK condicional conocida:} -. \textbf{Análisis / excepción:}
  deltas generados, fórmulas no editables.
\item
  \textbf{Tabla:} pago. \textbf{Determinantes mínimos NOT NULL:}
  \texorpdfstring{\softtext{\{empresa\_id,sucursal\_id,id\}.}}{\{empresa\_id,sucursal\_id,id\}.}
  \textbf{UK condicional conocida:} -. \textbf{Análisis / excepción:}
  \texorpdfstring{\softtext{1FN/2FN/3FN}}{1FN/2FN/3FN} y BCNF respecto a
  DF declaradas: determinantes son claves; no se presume otra unicidad
  comercial.
\item
  \textbf{Tabla:}
  \texorpdfstring{\softtext{aplicacion\_pago.}}{aplicacion\_pago.}
  \textbf{Determinantes mínimos NOT NULL:}
  \texorpdfstring{\softtext{\{empresa\_id,sucursal\_id,id\}.}}{\{empresa\_id,sucursal\_id,id\}.}
  \textbf{UK condicional conocida:} -. \textbf{Análisis / excepción:}
  \texorpdfstring{\softtext{1FN/2FN/3FN}}{1FN/2FN/3FN} y BCNF respecto a
  DF declaradas: determinantes son claves; no se presume otra unicidad
  comercial.
\item
  \textbf{Tabla:} documento. \textbf{Determinantes mínimos NOT NULL:}
  \texorpdfstring{\softtext{\{empresa\_id,sucursal\_id,id\}.}}{\{empresa\_id,sucursal\_id,id\}.}
  \textbf{UK condicional conocida:} -. \textbf{Análisis / excepción:}
  sha256 no determina nombre
  \texorpdfstring{\softtext{original/versión/scope;}}{original/versión/scope;}
  múltiples metadatos de mismos bytes son válidos.
\item
  \textbf{Tabla:}
  \texorpdfstring{\softtext{operacion\_documento.}}{operacion\_documento.}
  \textbf{Determinantes mínimos NOT NULL:}
  \texorpdfstring{\softtext{\{empresa\_id,sucursal\_id,operacion\_id,documento\_id\}.}}{\{empresa\_id,sucursal\_id,operacion\_id,documento\_id\}.}
  \textbf{UK condicional conocida:} -. \textbf{Análisis / excepción:}
  \texorpdfstring{\softtext{1FN/2FN/3FN}}{1FN/2FN/3FN} y BCNF respecto a
  DF declaradas: determinantes son claves; no se presume otra unicidad
  comercial.
\item
  \textbf{Tabla:}
  \texorpdfstring{\softtext{revision\_administrativa.}}{revision\_administrativa.}
  \textbf{Determinantes mínimos NOT NULL:}
  \texorpdfstring{\softtext{\{empresa\_id,sucursal\_id,id\}.}}{\{empresa\_id,sucursal\_id,id\}.}
  \textbf{UK condicional conocida:} -. \textbf{Análisis / excepción:}
  \texorpdfstring{\softtext{1FN/2FN/3FN}}{1FN/2FN/3FN} y BCNF respecto a
  DF declaradas: determinantes son claves; no se presume otra unicidad
  comercial.
\item
  \textbf{Tabla:}
  \texorpdfstring{\softtext{revision\_documento.}}{revision\_documento.}
  \textbf{Determinantes mínimos NOT NULL:}
  \texorpdfstring{\softtext{\{empresa\_id,sucursal\_id,revision\_id,documento\_id\}.}}{\{empresa\_id,sucursal\_id,revision\_id,documento\_id\}.}
  \textbf{UK condicional conocida:} -. \textbf{Análisis / excepción:}
  \texorpdfstring{\softtext{1FN/2FN/3FN}}{1FN/2FN/3FN} y BCNF respecto a
  DF declaradas: determinantes son claves; no se presume otra unicidad
  comercial.
\item
  \textbf{Tabla:}
  \texorpdfstring{\softtext{auditoria\_evento.}}{auditoria\_evento.}
  \textbf{Determinantes mínimos NOT NULL:}
  \texorpdfstring{\softtext{\{empresa\_id,sucursal\_id,id\}.}}{\{empresa\_id,sucursal\_id,id\}.}
  \textbf{UK condicional conocida:} -. \textbf{Análisis / excepción:}
  \texorpdfstring{\softtext{1FN/2FN/3FN}}{1FN/2FN/3FN} y BCNF respecto a
  DF declaradas: determinantes son claves; no se presume otra unicidad
  comercial.
\item
  \textbf{Tabla:}
  \texorpdfstring{\softtext{evento\_outbox.}}{evento\_outbox.}
  \textbf{Determinantes mínimos NOT NULL:}
  \texorpdfstring{\softtext{\{empresa\_id,sucursal\_id,id\}.}}{\{empresa\_id,sucursal\_id,id\}.}
  \textbf{UK condicional conocida:} -. \textbf{Análisis / excepción:}
  \texorpdfstring{\softtext{1FN/2FN/3FN}}{1FN/2FN/3FN} y BCNF respecto a
  DF declaradas: determinantes son claves; no se presume otra unicidad
  comercial.
\item
  \textbf{Tabla:}
  \texorpdfstring{\softtext{clave\_idempotencia.}}{clave\_idempotencia.}
  \textbf{Determinantes mínimos NOT NULL:}
  \texorpdfstring{\softtext{\{empresa\_id,sucursal\_id,principal\_id,comando,clave\}.}}{\{empresa\_id,sucursal\_id,principal\_id,comando,clave\}.}
  \textbf{UK condicional conocida:} -. \textbf{Análisis / excepción:}
  resultado inmutable de intención;
  \texorpdfstring{\softtext{contenido/hash/resultado}}{contenido/hash/resultado}
  forman hecho temporal.
\item
  \textbf{Tabla:}
  \texorpdfstring{\softtext{mensaje\_inbox.}}{mensaje\_inbox.}
  \textbf{Determinantes mínimos NOT NULL:}
  \texorpdfstring{\softtext{\{consumidor,fuente\_id,event\_id\}.}}{\{consumidor,fuente\_id,event\_id\}.}
  \textbf{UK condicional conocida:} -. \textbf{Análisis / excepción:}
  \texorpdfstring{\softtext{1FN/2FN/3FN}}{1FN/2FN/3FN} y BCNF respecto a
  DF declaradas: determinantes son claves; no se presume otra unicidad
  comercial.
\item
  \textbf{Tabla:}
  \texorpdfstring{\softtext{infra.nodo\_config.}}{infra.nodo\_config.}
  \textbf{Determinantes mínimos NOT NULL:} \{nodo\}. \textbf{UK
  condicional conocida:} -. \textbf{Análisis / excepción:}
  \texorpdfstring{\softtext{1FN/2FN/3FN}}{1FN/2FN/3FN} y BCNF respecto a
  DF declaradas: determinantes son claves; no se presume otra unicidad
  comercial.
\item
  \textbf{Tabla:}
  \texorpdfstring{\softtext{infra.entrega\_outbox.}}{infra.entrega\_outbox.}
  \textbf{Determinantes mínimos NOT NULL:}
  \texorpdfstring{\softtext{\{evento\_id,destino\}.}}{\{evento\_id,destino\}.}
  \textbf{UK condicional conocida:} -. \textbf{Análisis / excepción:}
  \texorpdfstring{\softtext{1FN/2FN/3FN}}{1FN/2FN/3FN} y BCNF respecto a
  DF declaradas: determinantes son claves; no se presume otra unicidad
  comercial.
\item
  \textbf{Tabla:}
  \texorpdfstring{\softtext{infra.importacion\_fuente.}}{infra.importacion\_fuente.}
  \textbf{Determinantes mínimos NOT NULL:} \{id\}. \textbf{UK
  condicional conocida:} -. \textbf{Análisis / excepción:}
  \texorpdfstring{\softtext{1FN/2FN/3FN}}{1FN/2FN/3FN} y BCNF respecto a
  DF declaradas: determinantes son claves; no se presume otra unicidad
  comercial.
\end{itemize}

\paragraph{Tablas técnicas del
catálogo}\label{tablas-tuxe9cnicas-del-catuxe1logo}

\texorpdfstring{\softtext{SpringSession}}{SpringSession} (clave
\texorpdfstring{\softtext{PRIMARY\_ID;}}{PRIMARY\_ID;} SESSION\_ID
única) y attributes
\texorpdfstring{\softtext{(PRIMARY\_ID,ATTRIBUTE\_NAME)}}{(PRIMARY\_ID,ATTRIBUTE\_NAME)}
separan el grupo repetido; bytes de atributo son serialización técnica,
no una nueva entidad comercial. Flyway conserva historia de
\texorpdfstring{\softtext{versión/checksum.}}{versión/checksum.}
\texorpdfstring{\softcode{infra.tarifa\_snapshot.id}}{infra.tarifa\_snapshot.id}
determina
\texorpdfstring{\softtext{captura/actor/scope/contenido;}}{captura/actor/scope/contenido;}
sólo \texorpdfstring{\softtext{aceptada\_en}}{aceptada\_en} puede
transitar de NULL a timestamp una vez. Precio y hash son derivados
materializados deliberados; el snapshot aceptado del gasto conserva
hechos aunque cambie la fuente FDW. Ninguna se publica como maestro
comercial ni implica normalización de contenido
\texorpdfstring{\softtext{binario/JSON.}}{binario/JSON.}

\section{Casos de uso}\label{casos-de-uso}

\textbf{Fuente:}
\texorpdfstring{\softcode{docs/domain/use-cases.md}}{docs/domain/use-cases.md}.

Actores: Administrador (ADM), Operador (OPE), Caja (CAJ), Auditor (AUD),
Sistema (SIS). Todos los comandos con efecto exigen
\texorpdfstring{\softcode{Idempotency-Key}}{Idempotency-Key} y se
ejecutan en el nodo dueño
\texorpdfstring{\softtext{(OwnerDatabaseRouter).}}{(OwnerDatabaseRouter).}

\begin{itemize}
\item
  \textbf{UC:} UC-01. \textbf{Nombre:} Iniciar sesión y elegir sede.
  \textbf{Actor:} Todos. \textbf{Flujo normal:} Login → lista de sedes
  autorizadas → \texorpdfstring{\softtext{selección.}}{selección.}
  \textbf{Alternativos / errores:} Credenciales inválidas 401; sede no
  autorizada 403; A caído: todas las nuevas solicitudes autenticadas A/B
  fallan503. \textbf{Invariantes:} Sesión de servidor, cookie HttpOnly.
  \textbf{RN:} RN-001. \textbf{Comando:} -.
\item
  \textbf{UC:} UC-02. \textbf{Nombre:} Registrar
  \texorpdfstring{\softtext{contraparte/predio.}}{contraparte/predio.}
  \textbf{Actor:} OPE. \textbf{Flujo normal:} Captura en A, replicación
  a B. \textbf{Alternativos / errores:} Duplicado de
  \texorpdfstring{\softtext{RFC/código}}{RFC/código} → 409.
  \textbf{Invariantes:} Maestros sólo en A. \textbf{RN:} RN-001, RN-010.
  \textbf{Comando:} CRUD maestro.
\item
  \textbf{UC:} UC-03. \textbf{Nombre:} Registrar identidad animal.
  \textbf{Actor:} OPE. \textbf{Flujo normal:} Alta en A con código
  normalizado; si es para B, esperar visibilidad en B (estado
  \texorpdfstring{\softtext{PENDIENTE\_CONVERGENCIA).}}{PENDIENTE\_CONVERGENCIA).}
  \textbf{Alternativos / errores:} Colisión → revisión central, no
  fusión \texorpdfstring{\softtext{automática.}}{automática.}
  \textbf{Invariantes:}
  \texorpdfstring{\softtext{UK(E,sistema,codigo).}}{UK(E,sistema,codigo).}
  \textbf{RN:} \texorpdfstring{\softtext{RN-002..004.}}{RN-002..004.}
  \textbf{Comando:}
  \texorpdfstring{\softtext{RegistrarIdentidadAnimal.}}{RegistrarIdentidadAnimal.}
\item
  \textbf{UC:} UC-04. \textbf{Nombre:} Capturar compra. \textbf{Actor:}
  OPE. \textbf{Flujo normal:} Borrador → lotes → selección de animales →
  versión y condiciones → pesajes → revisión. \textbf{Alternativos /
  errores:} Edición concurrente →412 por If-Match obsoleto; falta →428.
  \textbf{Invariantes:} Un animal una vez por
  \texorpdfstring{\softtext{operación.}}{operación.} \textbf{RN:}
  \texorpdfstring{\softtext{RN-007..011,}}{RN-007..011,} RN-017.
  \textbf{Comando:} edición de borrador.
\item
  \textbf{UC:} UC-05. \textbf{Nombre:} Aprobar pesaje. \textbf{Actor:}
  OPE. \textbf{Flujo normal:} Adjuntar ticket → modo
  \texorpdfstring{\softtext{BRUTO\_TARA/NETO\_REPORTADO}}{BRUTO\_TARA/NETO\_REPORTADO}
  → composición → aprobar. \textbf{Alternativos / errores:} Composición
  distinta del lote → advertencia; corrección = nuevo pesaje.
  \textbf{Invariantes:} Composición congelada. \textbf{RN:}
  \texorpdfstring{\softtext{RN-012..016.}}{RN-012..016.}
  \textbf{Comando:}
  \texorpdfstring{\softtext{AprobarPesaje,}}{AprobarPesaje,}
  \texorpdfstring{\softtext{RegistrarCorreccionPesaje.}}{RegistrarCorreccionPesaje.}
\item
  \textbf{UC:} UC-06. \textbf{Nombre:} Confirmar compra. \textbf{Actor:}
  OPE. \textbf{Flujo normal:} Valida precondiciones → crea posiciones,
  liquidación, asignaciones, asientos, auditoría, outbox.
  \textbf{Alternativos / errores:} Precondición faltante 422; periodo
  cerrado 422; reintento devuelve mismo resultado. \textbf{Invariantes:}
  Todo o nada. \textbf{RN:} RN-021, RN-026, RN-043. \textbf{Comando:}
  \texorpdfstring{\softtext{ConfirmarCompra.}}{ConfirmarCompra.}
\item
  \textbf{UC:} UC-07. \textbf{Nombre:} Registrar
  \texorpdfstring{\softtext{recepción/entrega.}}{recepción/entrega.}
  \textbf{Actor:} OPE. \textbf{Flujo normal:} Movimiento por posición;
  recepción → DISPONIBLE; entrega →
  \texorpdfstring{\softtext{CERRADA\_VENTA.}}{CERRADA\_VENTA.}
  \textbf{Alternativos / errores:} Posición en otro estado 422.
  \textbf{Invariantes:} No genera venta ni costo. \textbf{RN:} RN-020.
  \textbf{Comando:}
  \texorpdfstring{\softtext{RegistrarRecepcion/Entrega.}}{RegistrarRecepcion/Entrega.}
\item
  \textbf{UC:} UC-08. \textbf{Nombre:} Reservar y confirmar venta.
  \textbf{Actor:} OPE. \textbf{Flujo normal:} Selección →
  \texorpdfstring{\softtext{ReservarPosiciones}}{ReservarPosiciones} →
  condiciones →
  \texorpdfstring{\softtext{ConfirmarVenta.}}{ConfirmarVenta.}
  \textbf{Alternativos / errores:} Posición reservada por otra venta
  409; reserva vencida → revalida. \textbf{Invariantes:} Una reserva
  activa por posición. \textbf{RN:} RN-018, RN-019, RN-022.
  \textbf{Comando:}
  \texorpdfstring{\softtext{ReservarPosiciones,}}{ReservarPosiciones,}
  \texorpdfstring{\softtext{ConfirmarVenta.}}{ConfirmarVenta.}
\item
  \textbf{UC:} UC-09. \textbf{Nombre:} Registrar y aplicar pago.
  \textbf{Actor:} CAJ. \textbf{Flujo normal:} Registrar pago → validar →
  aplicar a liquidaciones
  \texorpdfstring{\softtext{compatibles.}}{compatibles.}
  \textbf{Alternativos / errores:} Excede
  \texorpdfstring{\softtext{saldo/disponible}}{saldo/disponible} 422;
  carrera → 409. \textbf{Invariantes:} Σ aplicaciones ≤ pago y ≤ saldo.
  \textbf{RN:} \texorpdfstring{\softtext{RN-028..032.}}{RN-028..032.}
  \textbf{Comando:}
  \texorpdfstring{\softtext{Registrar/ConfirmarPago,}}{Registrar/ConfirmarPago,}
  \texorpdfstring{\softtext{AplicarPago.}}{AplicarPago.}
\item
  \textbf{UC:} UC-10. \textbf{Nombre:} Revertir aplicación /
  \texorpdfstring{\softtext{reembolsar.}}{reembolsar.} \textbf{Actor:}
  CAJ+ADM. \textbf{Flujo normal:} Reverso completo; reembolso consume
  \texorpdfstring{\softtext{disponible.}}{disponible.}
  \textbf{Alternativos / errores:} Doble reverso 409.
  \textbf{Invariantes:} Original
  \texorpdfstring{\softtext{conservado.}}{conservado.} \textbf{RN:}
  RN-033. \textbf{Comando:}
  \texorpdfstring{\softtext{RevertirAplicacion,}}{RevertirAplicacion,}
  \texorpdfstring{\softtext{RegistrarReembolso.}}{RegistrarReembolso.}
\item
  \textbf{UC:} UC-11. \textbf{Nombre:} Ajustar
  \texorpdfstring{\softtext{liquidación.}}{liquidación.} \textbf{Actor:}
  ADM. \textbf{Flujo normal:} Ajuste
  \texorpdfstring{\softtext{CARGO/ABONO}}{CARGO/ABONO} con motivo y
  fuente. \textbf{Alternativos / errores:} Bonificación sobre pagada →
  guía GOLD-07. \textbf{Invariantes:} Liquidación inmutable.
  \textbf{RN:} RN-023, RN-027. \textbf{Comando:}
  \texorpdfstring{\softtext{AjustarLiquidacion.}}{AjustarLiquidacion.}
\item
  \textbf{UC:} UC-12. \textbf{Nombre:} Registrar costo tardío.
  \textbf{Actor:} OPE. \textbf{Flujo normal:} Gasto → previsualizar
  reparto → aprobar → asientos por destino
  \texorpdfstring{\softtext{económico.}}{económico.}
  \textbf{Alternativos / errores:} Periodo cerrado → BV025 sin efecto;
  seleccionar explícitamente periodo abierto con nueva
  \texorpdfstring{\softtext{intención.}}{intención.}
  \textbf{Invariantes:} Σ asignaciones = gasto. \textbf{RN:}
  \texorpdfstring{\softtext{RN-025..027.}}{RN-025..027.}
  \textbf{Comando:}
  \texorpdfstring{\softtext{DistribuirCosto,}}{DistribuirCosto,}
  \texorpdfstring{\softtext{RegistrarCostoTardio.}}{RegistrarCostoTardio.}
\item
  \textbf{UC:} UC-13. \textbf{Nombre:} Dar de baja posición.
  \textbf{Actor:} ADM. \textbf{Flujo normal:} Motivo
  \texorpdfstring{\softtext{(PERDIDA/MUERTE)}}{(PERDIDA/MUERTE)} +
  evidencia → BAJA. \textbf{Alternativos / errores:} Posición vendida
  422. \textbf{Invariantes:} No borrar. \textbf{RN:} RN-038.
  \textbf{Comando:}
  \texorpdfstring{\softtext{RegistrarBajaInventario.}}{RegistrarBajaInventario.}
\item
  \textbf{UC:} UC-14. \textbf{Nombre:}
  \texorpdfstring{\softtext{Cerrar/reabrir}}{Cerrar/reabrir} periodo.
  \textbf{Actor:} ADM. \textbf{Flujo normal:} Lock periodo →
  conciliación → versión de cierre con hash. \textbf{Alternativos /
  errores:} Hechos concurrentes esperan el lock. \textbf{Invariantes:}
  Cierre publicado inmutable. \textbf{RN:} RN-039. \textbf{Comando:}
  \texorpdfstring{\softtext{CerrarPeriodo,}}{CerrarPeriodo,}
  \texorpdfstring{\softtext{ReabrirPeriodo.}}{ReabrirPeriodo.}
\item
  \textbf{UC:} UC-15. \textbf{Nombre:} Consultar reportes.
  \textbf{Actor:} Todos/AUD. \textbf{Flujo normal:} Cartera, margen,
  inventario con corte y
  \texorpdfstring{\softtext{antigüedad.}}{antigüedad.}
  \textbf{Alternativos / errores:} Réplica atrasada → marca de
  \texorpdfstring{\softtext{antigüedad.}}{antigüedad.}
  \textbf{Invariantes:} Margen provisional
  \texorpdfstring{\softtext{etiquetado.}}{etiquetado.} \textbf{RN:}
  RN-035, RN-036. \textbf{Comando:} consultas.
\item
  \textbf{UC:} UC-16. \textbf{Nombre:} Consultar proveedor FDW y aceptar
  tarifa. \textbf{Actor:} OPE. \textbf{Flujo normal:} Consulta fuera de
  transacción con locks → aceptar → snapshot local. \textbf{Alternativos
  / errores:} Fuente caída: mensaje y reintento. \textbf{Invariantes:}
  No media \texorpdfstring{\softtext{liquidación.}}{liquidación.}
  \textbf{RN:} RN-042. \textbf{Comando:}
  \texorpdfstring{\softtext{AceptarTarifaExterna.}}{AceptarTarifaExterna.}
\item
  \textbf{UC:} UC-17. \textbf{Nombre:} Revisión administrativa de
  \texorpdfstring{\softtext{documentos.}}{documentos.} \textbf{Actor:}
  ADM. \textbf{Flujo normal:} Revisión de versión exacta con resultado y
  \texorpdfstring{\softtext{observación.}}{observación.}
  \textbf{Alternativos / errores:} Nueva versión no altera revisión
  previa. \textbf{Invariantes:} ``Revisión administrativa registrada'',
  nunca ``autorizado''. \textbf{RN:} RN-037. \textbf{Comando:}
  \texorpdfstring{\softtext{RegistrarRevision.}}{RegistrarRevision.}
\item
  \textbf{UC:} UC-18. \textbf{Nombre:} Validar conciliación del
  \texorpdfstring{\softtext{propietario.}}{propietario.} \textbf{Actor:}
  ADM/SIS. \textbf{Flujo normal:} Conteos por dueño, huérfanos,
  \texorpdfstring{\softtext{sumatorias.}}{sumatorias.}
  \textbf{Alternativos / errores:} Diferencias → incidente.
  \textbf{Invariantes:} -. \textbf{RN:} RN-040, RN-043.
  \textbf{Comando:}
  \texorpdfstring{\softtext{ValidarConciliacionDelPropietario.}}{ValidarConciliacionDelPropietario.}
\item
  \textbf{UC:} UC-19. \textbf{Nombre:} Simular margen. \textbf{Actor:}
  OPE. \textbf{Flujo normal:} Usa la función canónica en modo preview
  sin mutar. \textbf{Alternativos / errores:} -. \textbf{Invariantes:}
  No persiste. \textbf{RN:} RN-035. \textbf{Comando:} preview.
\end{itemize}

\section{ADR-001: Monolito modular frente a
microservicios}\label{adr-001-monolito-modular-frente-a-microservicios}

\textbf{Fuente:}
\texorpdfstring{\softcode{docs/architecture/adr/ADR-001-monolito-modular.md}}{docs/architecture/adr/ADR-001-monolito-modular.md}.

\begin{itemize}
\tightlist
\item
  Estado: propuesto · 2026-10-10 · Requisito: AC-08 \#\#\#\# Contexto
  Equipo de 2 personas, 21 días, transacciones locales que deben cubrir
  operación + liquidación + inventario +
  \texorpdfstring{\softtext{auditoría.}}{auditoría.} \#\#\#\#
  Alternativas
\end{itemize}

\begin{enumerate}
\def\labelenumi{\arabic{enumi}.}
\tightlist
\item
  Microservicios por dominio - transacciones distribuidas, más
  despliegue y
  \texorpdfstring{\softtext{observabilidad.}}{observabilidad.}
\item
  Monolito sin módulos - rápido al inicio, acoplamiento alto.
\item
  \textbf{Monolito modular} con paquetes por dominio y puertos. \#\#\#\#
  Decisión Opción 3: un backend Spring Boot desplegable con 10 módulos
  (\texorpdfstring{\softcode{overview.md}}{overview.md} §3). \#\#\#\#
  Consecuencias
\end{enumerate}

\begin{itemize}
\tightlist
\item
  Una transacción local por caso de uso; despliegue simple. − Disciplina
  de dependencias entre módulos sin herramienta que la imponga (se
  revisará con ArchUnit si hay tiempo). Sin
  \texorpdfstring{\softtext{Kafka/Kubernetes.}}{Kafka/Kubernetes.}
\end{itemize}

\section{ADR-002: SQL explícito, cálculo canónico y migraciones
versionadas}\label{adr-002-sql-expluxedcito-cuxe1lculo-canuxf3nico-y-migraciones-versionadas}

\textbf{Fuente:}
\texorpdfstring{\softcode{docs/architecture/adr/ADR-002-sql-explicito-calculo-canonico.md}}{docs/architecture/adr/ADR-002-sql-explicito-calculo-canonico.md}.

\begin{itemize}
\tightlist
\item
  Estado: propuesto · 2026-10-10 · Requisitos: AC-03, AC-04, AC-05
  \#\#\#\# Contexto La materia evalúa SQL, PL/pgSQL y transacciones; los
  importes deben ser reproducibles al centavo. \#\#\#\# Alternativas
\end{itemize}

\begin{enumerate}
\def\labelenumi{\arabic{enumi}.}
\tightlist
\item
  \texorpdfstring{\softtext{JPA/Hibernate}}{JPA/Hibernate} con
  \texorpdfstring{\softcode{ddl-auto}}{ddl-auto} - oculta SQL y no sirve
  como fuente del diseño.
\item
  \textbf{\texorpdfstring{\softtext{JdbcTemplate}}{JdbcTemplate} +
  Flyway + funciones PL/pgSQL
  \texorpdfstring{\softtext{canónicas.}}{canónicas.}}
\item
  jOOQ - buena opción, añade generación de código. \#\#\#\# Decisión
  Opción 2. Importes confirmados se calculan en
  \texorpdfstring{\softcode{fn\_calcular\_liquidacion}}{fn\_calcular\_liquidacion}
  (SQL) y se congelan; la previsualización llama la misma función sin
  persistir. Java usa BigDecimal y valida; no recalcula con otra
  fórmula. \#\#\#\# Consecuencias
\end{enumerate}

\begin{itemize}
\tightlist
\item
  SQL visible y evaluable; una sola fórmula. − Más SQL a mano; mapeos
  \texorpdfstring{\softtext{explícitos.}}{explícitos.} Pruebas con
  PostgreSQL real (Testcontainers), nunca H2.
\end{itemize}

\section{ADR-003: Animal permanente y posición de inventario por
ciclo}\label{adr-003-animal-permanente-y-posiciuxf3n-de-inventario-por-ciclo}

\textbf{Fuente:}
\texorpdfstring{\softcode{docs/architecture/adr/ADR-003-animal-y-posicion.md}}{docs/architecture/adr/ADR-003-animal-y-posicion.md}.

\begin{itemize}
\tightlist
\item
  Estado: propuesto · 2026-10-10 · Reglas: RN-002, RN-005, RN-006
  \#\#\#\# Contexto Un animal puede comprarse, venderse y recomprarse;
  cada tenencia tiene su propio costo. \#\#\#\# Alternativas
\end{itemize}

\begin{enumerate}
\def\labelenumi{\arabic{enumi}.}
\tightlist
\item
  Estado y costo en \texorpdfstring{\softcode{animal}}{animal} - pierde
  historia en recompra.
\item
  \textbf{\texorpdfstring{\softcode{animal}}{animal} (identidad) +
  \texorpdfstring{\softcode{posicion\_inventario}}{posicion\_inventario}
  (ciclo)} con índice único parcial de ciclo activo. \#\#\#\# Decisión
  Opción 2. Estados de posición:
  \texorpdfstring{\softtext{PENDIENTE\_RECEPCION,}}{PENDIENTE\_RECEPCION,}
  DISPONIBLE,
  \texorpdfstring{\softtext{VENDIDA\_PENDIENTE\_ENTREGA,}}{VENDIDA\_PENDIENTE\_ENTREGA,}
  \texorpdfstring{\softtext{CERRADA\_VENTA,}}{CERRADA\_VENTA,} BAJA;
  ``reservada'' derivada de reserva ACTIVA
  \texorpdfstring{\softtext{(DEC-01..03).}}{(DEC-01..03).}
  \texorpdfstring{\softcode{animal.situacion}}{animal.situacion} bloquea
  recompra de MUERTO (DEC-04). \#\#\#\# Consecuencias
\end{enumerate}

\begin{itemize}
\tightlist
\item
  GOLD-05/09 naturales. − Consultas de ficha requieren join
  \texorpdfstring{\softtext{animal-posiciones.}}{animal-posiciones.}
\end{itemize}

\section{ADR-004: Reconocimiento gerencial y costo por
posición}\label{adr-004-reconocimiento-gerencial-y-costo-por-posiciuxf3n}

\textbf{Fuente:}
\texorpdfstring{\softcode{docs/architecture/adr/ADR-004-reconocimiento-y-costo.md}}{docs/architecture/adr/ADR-004-reconocimiento-y-costo.md}.

\begin{itemize}
\tightlist
\item
  Estado: propuesto · 2026-10-10 · Reglas: RN-021, RN-022, RN-035
  \#\#\#\# Contexto Entrega, cobro y PDF no deben generar una segunda
  venta, deuda o costo. \#\#\#\# Alternativas
\end{itemize}

\begin{enumerate}
\def\labelenumi{\arabic{enumi}.}
\tightlist
\item
  Reconocer al entregar.
\item
  Reconocer al cobrar.
\item
  \textbf{Reconocer al CONFIRMAR} la operación (animales, base, versión
  y costos \texorpdfstring{\softtext{determinados).}}{determinados).}
  \#\#\#\# Decisión Opción 3 como política operativa versionada (no
  dictamen contable). Costo identificado por posición mediante
  \texorpdfstring{\softcode{asignacion\_costo}}{asignacion\_costo}
  (soporte) y
  \texorpdfstring{\softcode{asiento\_valoracion}}{asiento\_valoracion}
  (efecto), con la tabla de destino económico de
  \texorpdfstring{\softcode{states-and-calculations.md}}{states-and-calculations.md}
  §3. \#\#\#\# Consecuencias
\end{enumerate}

\begin{itemize}
\tightlist
\item
  Un único evento económico; margen y cartera
  \texorpdfstring{\softtext{conciliables.}}{conciliables.} − Cambiar la
  política exige migración de modelo y reportes.
\end{itemize}

\section{ADR-005: Maestro de identidad en A y propiedad de escritura
inmutable}\label{adr-005-maestro-de-identidad-en-a-y-propiedad-de-escritura-inmutable}

\textbf{Fuente:}
\texorpdfstring{\softcode{docs/architecture/adr/ADR-005-maestro-y-propiedad.md}}{docs/architecture/adr/ADR-005-maestro-y-propiedad.md}.

\begin{itemize}
\tightlist
\item
  Estado: propuesto · 2026-10-10 · Requisito: AC-06 · Reglas: RN-003,
  RN-004, RN-040, RN-041 \#\#\#\# Contexto Dos escritores independientes
  podrían dar dos UUID al mismo arete; la replicación lógica no
  garantiza unicidad global. \#\#\#\# Alternativas
\end{itemize}

\begin{enumerate}
\def\labelenumi{\arabic{enumi}.}
\tightlist
\item
  Unicidad por nodo - permite duplicados entre nodos.
\item
  Escritura libre en ambos + resolución de conflictos - no existe
  resolución de negocio segura para
  \texorpdfstring{\softtext{inventario/dinero.}}{inventario/dinero.}
\item
  \textbf{Maestros sólo en A; agregados locales sólo en el nodo de su
  sede; sede del animal inmutable.} \#\#\#\# Decisión Opción 3, con
  \texorpdfstring{\softcode{nodo\_origen}}{nodo\_origen} en cada fila y
  \texorpdfstring{\softcode{trg\_guarda\_propietario}}{trg\_guarda\_propietario}.
  Sin failover automático de dueño. Transferir un animal de sede queda
  fuera (requiere fencing y otro ADR). \#\#\#\# Consecuencias
\end{enumerate}

\begin{itemize}
\tightlist
\item
  Sin conflictos de negocio por diseño. − Si A cae no hay maestros
  nuevos ni login nuevo.
\end{itemize}

\section{ADR-006: Replicación bidireccional, bootstrap y
recuperación}\label{adr-006-replicaciuxf3n-bidireccional-bootstrap-y-recuperaciuxf3n}

\textbf{Fuente:}
\texorpdfstring{\softcode{docs/architecture/adr/ADR-006-replicacion-bidireccional.md}}{docs/architecture/adr/ADR-006-replicacion-bidireccional.md}.

\begin{itemize}
\tightlist
\item
  Estado: propuesto · 2026-10-10 · Requisito: AC-06 \#\#\#\# Contexto Se
  requiere M-M con dos nodos escritores; PostgreSQL 18 ofrece
  replicación lógica con
  \texorpdfstring{\softcode{origin = none}}{origin = none}. \#\#\#\#
  Alternativas
\end{itemize}

\begin{enumerate}
\def\labelenumi{\arabic{enumi}.}
\tightlist
\item
  Streaming físico (un solo escritor) - no cumple M-M.
\item
  Extensiones de terceros (pglogical, BDR/PGD) -
  \texorpdfstring{\softtext{licencia/soporte,}}{licencia/soporte,} más
  \texorpdfstring{\softtext{complejidad.}}{complejidad.}
\item
  \textbf{Replicación lógica nativa bidireccional con publicaciones
  explícitas y
  \texorpdfstring{\softcode{origin = none}}{origin = none}.} \#\#\#\#
  Decisión Opción 3 con el bootstrap de
  \texorpdfstring{\softcode{docs/database/replication-and-fdw.md}}{docs/database/replication-and-fdw.md}
  §3 (\texorpdfstring{\softcode{copy\_data = false}}{copy\_data = false}
  sólo en bases vacías). Recuperación: slots monitoreados,
  \texorpdfstring{\softcode{max\_slot\_wal\_keep\_size}}{max\_slot\_wal\_keep\_size},
  runbook de
  \texorpdfstring{\softtext{resincronización.}}{resincronización.}
  \#\#\#\# Consecuencias
\end{enumerate}

\begin{itemize}
\tightlist
\item
  Nativo y \texorpdfstring{\softtext{demostrable.}}{demostrable.} − DDL
  coordinado; \texorpdfstring{\softtext{triggers/FK}}{triggers/FK} no se
  verifican en apply → conciliación
  \texorpdfstring{\softtext{obligatoria.}}{obligatoria.}
\end{itemize}

\section{ADR-007: Pgpool-II en modo raw y presupuesto de
conexiones}\label{adr-007-pgpool-ii-en-modo-raw-y-presupuesto-de-conexiones}

\textbf{Fuente:}
\texorpdfstring{\softcode{docs/architecture/adr/ADR-007-pgpool-raw.md}}{docs/architecture/adr/ADR-007-pgpool-raw.md}.

\begin{itemize}
\tightlist
\item
  Estado: propuesto · 2026-10-10 · Requisito: AC-09 \#\#\#\# Contexto La
  rúbrica exige Pgpool-II en el dibujo; los modos de balanceo asumen un
  primario único. \#\#\#\# Alternativas
\end{itemize}

\begin{enumerate}
\def\labelenumi{\arabic{enumi}.}
\tightlist
\item
  Pgpool
  \texorpdfstring{\softtext{streaming/replication}}{streaming/replication}
  mode delante de ambos - enrutaría escrituras a un nodo
  \texorpdfstring{\softtext{equivocado.}}{equivocado.}
\item
  \textbf{Una instancia Pgpool raw por nodo}, la app elige el Pgpool del
  dueño.
\item
  Sin Pgpool - no cumple AC-09. \#\#\#\# Decisión Opción 2 (Pgpool
  4.7.x). Presupuesto por nodo: Hikari 8,
  \texorpdfstring{\softcode{num\_init\_children}}{num\_init\_children}
  16, \texorpdfstring{\softcode{max\_pool}}{max\_pool} 1, + replicación
  y admin ≤
  \texorpdfstring{\softcode{max\_connections}}{max\_connections} 50.
  \#\#\#\# Consecuencias
\end{enumerate}

\begin{itemize}
\tightlist
\item
  Cumple AC-09 sin atribuirle
  \texorpdfstring{\softtext{sincronización.}}{sincronización.} − Pgpool
  añade un salto sin balanceo real.
\end{itemize}

\section{ADR-008: Autenticación, permisos y aislamiento del
piloto}\label{adr-008-autenticaciuxf3n-permisos-y-aislamiento-del-piloto}

\textbf{Fuente:}
\texorpdfstring{\softcode{docs/architecture/adr/ADR-008-autenticacion-y-aislamiento.md}}{docs/architecture/adr/ADR-008-autenticacion-y-aislamiento.md}.

\begin{itemize}
\tightlist
\item
  Estado: propuesto · 2026-10-10 · Reglas: RN-001 \#\#\#\# Alternativas
\end{itemize}

\begin{enumerate}
\def\labelenumi{\arabic{enumi}.}
\tightlist
\item
  JWT en \texorpdfstring{\softtext{localStorage}}{localStorage} -
  expuesto a XSS, revocación difícil.
\item
  \textbf{Sesión de servidor (Spring Security) con cookie
  \texorpdfstring{\softtext{HttpOnly/Secure/SameSite=Lax}}{HttpOnly/Secure/SameSite=Lax}
  + CSRF, mismo origen.}
\item
  OIDC externo - no necesario para el piloto. \#\#\#\# Decisión Opción
  2. Sesiones locales no replicadas: si cae el nodo de sesiones se
  \texorpdfstring{\softtext{reautentica.}}{reautentica.} Contraseñas con
  \texorpdfstring{\softcode{DelegatingPasswordEncoder}}{DelegatingPasswordEncoder}.
  Piloto: una empresa por despliegue; empresa\_id se conserva para
  crecer, sin vender multiempresa compartido hasta probar casos
  negativos. \#\#\#\# Consecuencias
\end{enumerate}

\begin{itemize}
\tightlist
\item
  Revocación inmediata. − Escritura en B requiere sesión válida (no se
  promete disponibilidad total si A cae).
\end{itemize}

\section{ADR-009: Documentos privados, evidencia y revisión
administrativa}\label{adr-009-documentos-privados-evidencia-y-revisiuxf3n-administrativa}

\textbf{Fuente:}
\texorpdfstring{\softcode{docs/architecture/adr/ADR-009-documentos-y-revision.md}}{docs/architecture/adr/ADR-009-documentos-y-revision.md}.

\begin{itemize}
\tightlist
\item
  Estado: propuesto · 2026-10-10 · Reglas: RN-037 \#\#\#\# Alternativas
\end{itemize}

\begin{enumerate}
\def\labelenumi{\arabic{enumi}.}
\tightlist
\item
  Bytes en la BD (bytea) - infla respaldos y réplica.
\item
  \textbf{Puerto
  \texorpdfstring{\softcode{DocumentStorage}}{DocumentStorage} con
  adaptador de volumen local privado}
  \texorpdfstring{\softtext{(dev/academic)}}{(dev/academic)} y
  almacenamiento de objetos privado en piloto; metadatos + SHA-256 en
  BD. \#\#\#\# Decisión Opción 2. Nombres internos aleatorios,
  validación de tamaño y MIME real, descarga autorizada; nueva versión =
  nuevo registro con
  \texorpdfstring{\softcode{documento\_previo\_id}}{documento\_previo\_id}.
  Revisión apunta a versión exacta; rótulo ``revisión administrativa
  registrada'', nunca ``autorizado por SENASICA''. \#\#\#\#
  Consecuencias
\end{enumerate}

\begin{itemize}
\tightlist
\item
  Integridad \texorpdfstring{\softtext{verificable.}}{verificable.} − El
  respaldo debe incluir el volumen de documentos además de la BD.
\end{itemize}

\section{ADR-010: Despliegue, respaldos y límites de
disponibilidad}\label{adr-010-despliegue-respaldos-y-luxedmites-de-disponibilidad}

\textbf{Fuente:}
\texorpdfstring{\softcode{docs/architecture/adr/ADR-010-despliegue-respaldos.md}}{docs/architecture/adr/ADR-010-despliegue-respaldos.md}.

\begin{itemize}
\tightlist
\item
  Estado: propuesto · 2026-10-10 \#\#\#\# Alternativas
\end{itemize}

\begin{enumerate}
\def\labelenumi{\arabic{enumi}.}
\tightlist
\item
  Kubernetes -
  \texorpdfstring{\softtext{sobredimensionado.}}{sobredimensionado.}
\item
  \textbf{Docker Compose + Caddy en un VPS}, respaldos fuera del host.
  \#\#\#\# Decisión Opción 2. Demo: dumps lógicos con corte y
  conciliación + restauración aislada medida. Piloto: backup base +
  archivo WAL por nodo antes de prometer RPO 15 min / RTO 2 h. \#\#\#\#
  Consecuencias
\end{enumerate}

\begin{itemize}
\tightlist
\item
  Reproducible y barato. − Un solo host es un único dominio de fallo; no
  hay alta \texorpdfstring{\softtext{disponibilidad.}}{disponibilidad.}
\end{itemize}

\section{ADR-011: Graphify y mantenimiento del
contexto}\label{adr-011-graphify-y-mantenimiento-del-contexto}

\textbf{Fuente:}
\texorpdfstring{\softcode{docs/architecture/adr/ADR-011-graphify.md}}{docs/architecture/adr/ADR-011-graphify.md}.

\begin{itemize}
\tightlist
\item
  Estado: propuesto · 2026-10-10 \#\#\#\# Contexto El equipo quiere
  contexto navegable para Claude y personas. Instalado: paquete
  \texorpdfstring{\softcode{graphifyy}}{graphifyy} 0.9.49. \#\#\#\#
  Decisión Grafo de desarrollo (no dependencia del runtime) con
  \texorpdfstring{\softcode{/graphify .}}{/graphify .} y
  \texorpdfstring{\softcode{/graphify . --update}}{/graphify . --update}
  tras cambios de documentos;
  \texorpdfstring{\softcode{graphify update .}}{graphify update .} sólo
  para cambios de código.
  \texorpdfstring{\softcode{.graphifyignore}}{.graphifyignore} excluye
  secretos, datos, dumps, dependencias y
  \texorpdfstring{\softcode{graphify-out/}}{graphify-out/}. Salidas por
  checkout; \texorpdfstring{\softcode{graphify-out/}}{graphify-out/} no
  se versiona; sin Graphify Cloud ni servidor público. Estado en
  \texorpdfstring{\softcode{docs/knowledge/graph-status.md}}{docs/knowledge/graph-status.md}.
  \#\#\#\# Consecuencias
\item
  Navegación de
  \texorpdfstring{\softtext{reglas↔tablas↔pruebas.}}{reglas↔tablas↔pruebas.}
  − Aristas inferidas deben verificarse contra fuentes; el procesamiento
  semántico puede usar el modelo configurado (no es offline
  \texorpdfstring{\softtext{garantizado).}}{garantizado).}
\end{itemize}

\section{ADR-012: Alcance aplazado y condiciones para
comercializar}\label{adr-012-alcance-aplazado-y-condiciones-para-comercializar}

\textbf{Fuente:}
\texorpdfstring{\softcode{docs/architecture/adr/ADR-012-alcance-aplazado.md}}{docs/architecture/adr/ADR-012-alcance-aplazado.md}.

\begin{itemize}
\tightlist
\item
  Estado: propuesto · 2026-10-10 \#\#\#\# Fuera del MVP Exportación,
  subastas,
  \texorpdfstring{\softtext{sacrificio/cortes,}}{sacrificio/cortes,}
  leche, reproducción, nutrición, diagnósticos, prescripciones, crédito,
  custodia de dinero, timbrado CFDI, Carta Porte, básculas en tiempo
  real, cambio de sede escritora, transacciones entre sedes, edición
  concurrente de la misma entidad en ambos nodos, devoluciones físicas
  posteriores a la entrega, migración de inventario histórico, cuentas
  por pagar a
  \texorpdfstring{\softtext{transportistas.}}{transportistas.} \#\#\#\#
  Condiciones para comercializar
\end{itemize}

\begin{enumerate}
\def\labelenumi{\arabic{enumi}.}
\tightlist
\item
  Problema y precio validados
  \texorpdfstring{\softtext{(validation-plan).}}{(validation-plan).} 2.
  Contrato con cancelación, exportación, retención y aviso de
  \texorpdfstring{\softtext{privacidad.}}{privacidad.} 3. Respaldo con
  WAL y restauración ensayada. 4. Facturación del servicio vía PAC. 5.
  Soporte \texorpdfstring{\softtext{dimensionado.}}{dimensionado.}
  \#\#\#\# Consecuencias Las funciones aplazadas se bloquean
  explícitamente en API/UI (no botones genéricos como ``restaurar
  inventario'').
\end{enumerate}

\section{ADR-013: Sesiones centrales y
disponibilidad}\label{adr-013-sesiones-centrales-y-disponibilidad}

\textbf{Fuente:}
\texorpdfstring{\softcode{docs/architecture/adr/ADR-013-sesiones-disponibilidad.md}}{docs/architecture/adr/ADR-013-sesiones-disponibilidad.md}.

\begin{itemize}
\tightlist
\item
  Estado: aceptado para implementación local por encargo del usuario ·
  2026-10-11
\end{itemize}

\subsubsection{Decisión}\label{decisiuxf3n}

Spring Session JDBC y autorización vigente residen en A. Toda nueva
petición autenticada de A/B depende de A; su caída produce 503
\texorpdfstring{\softtext{controlado.}}{controlado.} Una transacción B
ya autorizada puede finalizar. Sustituye cualquier promesa de
continuidad de sesiones B del ADR-008. Cookie
\texorpdfstring{\softtext{HttpOnly/SameSite=Lax,}}{HttpOnly/SameSite=Lax,}
CSRF también en \texorpdfstring{\softtext{login/logout;}}{login/logout;}
Secure obligatorio fuera del perfil HTTP local.

\subsubsection{Validación}\label{validaciuxf3n}

Implementación y pruebas pendientes al registrar esta decisión; ver
evidencia (referencia del repositorio:
\texorpdfstring{\softcode{../../testing/evidence.md}}{../../testing/evidence.md}).
La revisión humana permanece pendiente.

\section{ADR-014: Muerte en B e identidad en
A}\label{adr-014-muerte-en-b-e-identidad-en-a}

\textbf{Fuente:}
\texorpdfstring{\softcode{docs/architecture/adr/ADR-014-muerte-entre-duenos.md}}{docs/architecture/adr/ADR-014-muerte-entre-duenos.md}.

\begin{itemize}
\tightlist
\item
  Estado: aceptado para implementación local por encargo del usuario ·
  2026-10-11
\end{itemize}

\subsubsection{Decisión}\label{decisiuxf3n-1}

B reconoce \texorpdfstring{\softtext{baja/pérdida}}{baja/pérdida} y
\texorpdfstring{\softtext{MuerteAnimalSolicitada.v1}}{MuerteAnimalSolicitada.v1}
\texorpdfstring{\softtext{atómicamente.}}{atómicamente.} Worker reclama
lease local corto, termina la transacción y entrega a inbox idempotente
A. A actualiza identidad sin volver a reconocer pérdida. Acuse perdido
recupera resultado inbox; estado aplicado en A y visibilidad en B son
distintos. Sustituye la actualización directa del maestro desde el
comando B.

\subsubsection{Validación}\label{validaciuxf3n-1}

Implementación y pruebas pendientes al registrar esta decisión; ver
evidencia (referencia del repositorio:
\texorpdfstring{\softcode{../../testing/evidence.md}}{../../testing/evidence.md}).
La revisión humana permanece pendiente.

\section{ADR-015: Reconocimiento y cierres
históricos}\label{adr-015-reconocimiento-y-cierres-histuxf3ricos}

\textbf{Fuente:}
\texorpdfstring{\softcode{docs/architecture/adr/ADR-015-periodos-cartera.md}}{docs/architecture/adr/ADR-015-periodos-cartera.md}.

\begin{itemize}
\tightlist
\item
  Estado: aceptado para implementación local por encargo del usuario ·
  2026-10-11
\end{itemize}

\subsubsection{Decisión}\label{decisiuxf3n-2}

Periodos mensuales calendario por
\texorpdfstring{\softtext{empresa/sede}}{empresa/sede}
\texorpdfstring{\softtext{America/Mexico\_City,}}{America/Mexico\_City,}
uno abierto y rangos consecutivos sin
\texorpdfstring{\softtext{solapamiento.}}{solapamiento.} Todo hecho que
cambia
\texorpdfstring{\softtext{valoración/cartera}}{valoración/cartera}
conserva periodo y fecha de reconocimiento, además de fecha documental y
registro UTC. Lock común antes de recursos
\texorpdfstring{\softtext{financieros.}}{financieros.}
\texorpdfstring{\softtext{Cierre/versiones}}{Cierre/versiones}
\texorpdfstring{\softtext{inmutables.}}{inmutables.} Reapertura sólo
último cerrado sin periodo posterior ni abierto. Perdedor frente al
cierre recibe BV025 sin efecto, cambia periodo sólo mediante intención
explícita nueva.

\subsubsection{Validación}\label{validaciuxf3n-2}

Implementación y pruebas pendientes al registrar esta decisión; ver
evidencia (referencia del repositorio:
\texorpdfstring{\softcode{../../testing/evidence.md}}{../../testing/evidence.md}).
La revisión humana permanece pendiente.

\section{ADR-016: Autoridad del esquema
físico}\label{adr-016-autoridad-del-esquema-fuxedsico}

\textbf{Fuente:}
\texorpdfstring{\softcode{docs/architecture/adr/ADR-016-autoridad-esquema.md}}{docs/architecture/adr/ADR-016-autoridad-esquema.md}.

\begin{itemize}
\tightlist
\item
  Estado: aceptado para implementación local por encargo del usuario ·
  2026-10-11
\end{itemize}

\subsubsection{Decisión}\label{decisiuxf3n-3}

\texorpdfstring{\softtext{db/migrations}}{db/migrations} es autoridad
física \texorpdfstring{\softtext{ejecutable.}}{ejecutable.}
\texorpdfstring{\softtext{tools/esquema/esquema.py}}{tools/esquema/esquema.py}
es representación lógica documental que debe contrastarse con catálogo
\texorpdfstring{\softtext{PostgreSQL;}}{PostgreSQL;} no se genera DDL
\texorpdfstring{\softtext{ciegamente.}}{ciegamente.} Migrador separado,
runtime noDDL. Tablas técnicas locales y excepciones de publicación
\texorpdfstring{\softtext{explícitas.}}{explícitas.} P0001 DETAIL JSON
validado transporta códigos BV, no SQLSTATE inventado.

\subsubsection{Validación}\label{validaciuxf3n-3}

Implementación y pruebas pendientes al registrar esta decisión; ver
evidencia (referencia del repositorio:
\texorpdfstring{\softcode{../../testing/evidence.md}}{../../testing/evidence.md}).
La revisión humana permanece pendiente.

\section{ADR-017: Despliegue inicial
preparado}\label{adr-017-despliegue-inicial-preparado}

\textbf{Fuente:}
\texorpdfstring{\softcode{docs/architecture/adr/ADR-017-despliegue-local.md}}{docs/architecture/adr/ADR-017-despliegue-local.md}.

\begin{itemize}
\tightlist
\item
  Estado: aceptado para implementación local por encargo del usuario ·
  2026-10-11
\end{itemize}

\subsubsection{Decisión}\label{decisiuxf3n-4}

Una empresa por despliegue, VM Linux
\texorpdfstring{\softtext{DockerCompose+Caddy}}{DockerCompose+Caddy}
preparada; ejecución actual localhost. Dos nodos en una máquina son
topología académica, no redundancia física. Sin
\texorpdfstring{\softtext{publicación/push/contratación.}}{publicación/push/contratación.}
Secretos externos privados, puertosloopback, archivos
\texorpdfstring{\softtext{fuerawebroot.}}{fuerawebroot.} Restauración
aislada ejecutada antes de afirmar
\texorpdfstring{\softtext{recuperación.}}{recuperación.}

\subsubsection{Validación}\label{validaciuxf3n-4}

Implementación y pruebas pendientes al registrar esta decisión; ver
evidencia (referencia del repositorio:
\texorpdfstring{\softcode{../../testing/evidence.md}}{../../testing/evidence.md}).
La revisión humana permanece pendiente.

\chapter{Estado de tareas y alcance de la
entrega}\label{estado-de-tareas-y-alcance-de-la-entrega}

Este estado proviene del backlog canónico del corte del software.
Conserva responsables humanos aunque los agentes hayan ejecutado
trabajo. Las52 tareas verificadas localmente no están aprobadas
automáticamente por \texorpdfstring{\softtext{Aldo/Jared:}}{Aldo/Jared:}
todas las55 revisiones humanas permanecen
\texorpdfstring{\softtext{pendientes.}}{pendientes.} A028 contiene
preparación local pero depende de
\texorpdfstring{\softtext{staging/TLS/CI}}{staging/TLS/CI} externos;
A031/J020 son OCR aplazado. La publicación GitHub main se realizó
después del corte de algunos documentos y se confirmó en92 d3 d41.

{\def\LTcaptype{none} % do not increment counter
\begin{longtable}[]{@{}
  >{\raggedright\arraybackslash}p{(\linewidth - 2\tabcolsep) * \real{0.5000}}
  >{\raggedright\arraybackslash}p{(\linewidth - 2\tabcolsep) * \real{0.5000}}@{}}
\toprule\noalign{}
\begin{minipage}[b]{\linewidth}\raggedright
Tarea y responsable
\end{minipage} & \begin{minipage}[b]{\linewidth}\raggedright
Entregable y estado técnico
\end{minipage} \\
\midrule\noalign{}
\endhead
\bottomrule\noalign{}
\endlastfoot
\texorpdfstring{\softcode{TASK-A-001}}{TASK-A-001} - Aldo & Inventario
del \texorpdfstring{\softtext{repo/herramientas}}{repo/herramientas} y
matriz AC.
\textbf{\texorpdfstring{\softtext{verificado\_local}}{verificado\_local}}.
Revisión humana pendiente. \\
\texorpdfstring{\softcode{TASK-A-002}}{TASK-A-002} - Aldo & Negocio,
glosario, reglas y GOLD.
\textbf{\texorpdfstring{\softtext{verificado\_local}}{verificado\_local}}.
Revisión humana pendiente. \\
\texorpdfstring{\softcode{TASK-A-003}}{TASK-A-003} - Aldo & MER Chen y
matriz de \texorpdfstring{\softtext{relaciones.}}{relaciones.}
\textbf{\texorpdfstring{\softtext{verificado\_local}}{verificado\_local}}.
Revisión humana pendiente. \\
\texorpdfstring{\softcode{TASK-A-004}}{TASK-A-004} - Aldo & Modelo
relacional, diccionario y
\texorpdfstring{\softtext{normalización.}}{normalización.}
\textbf{\texorpdfstring{\softtext{verificado\_local}}{verificado\_local}}.
Revisión humana pendiente. \\
\texorpdfstring{\softcode{TASK-A-005}}{TASK-A-005} - Aldo &
Arquitectura, ADR y plan de
\texorpdfstring{\softtext{despliegue.}}{despliegue.}
\textbf{\texorpdfstring{\softtext{verificado\_local}}{verificado\_local}}.
Revisión humana pendiente. \\
\texorpdfstring{\softcode{TASK-A-006}}{TASK-A-006} - Aldo & Contexto
Markdown, README y grafo inicial.
\textbf{\texorpdfstring{\softtext{verificado\_local}}{verificado\_local}}.
Revisión humana pendiente. \\
\texorpdfstring{\softcode{TASK-A-007}}{TASK-A-007} - Aldo & OpenAPI,
errores y ejemplos para Jared.
\textbf{\texorpdfstring{\softtext{verificado\_local}}{verificado\_local}}.
Revisión humana pendiente. \\
\texorpdfstring{\softcode{TASK-A-008}}{TASK-A-008} - Aldo & Backend
modular y build \texorpdfstring{\softtext{reproducible.}}{reproducible.}
\textbf{\texorpdfstring{\softtext{verificado\_local}}{verificado\_local}}.
Revisión humana pendiente. \\
\texorpdfstring{\softcode{TASK-A-009}}{TASK-A-009} - Aldo &
\texorpdfstring{\softtext{DDL/Flyway,}}{DDL/Flyway,} claves, índices y
seed \texorpdfstring{\softtext{sintético.}}{sintético.}
\textbf{\texorpdfstring{\softtext{verificado\_local}}{verificado\_local}}.
Revisión humana pendiente. \\
\texorpdfstring{\softcode{TASK-A-010}}{TASK-A-010} - Aldo & Roles
\texorpdfstring{\softtext{PostgreSQL,}}{PostgreSQL,} grants y guardas de
dueño.
\textbf{\texorpdfstring{\softtext{verificado\_local}}{verificado\_local}}.
Revisión humana pendiente. \\
\texorpdfstring{\softcode{TASK-A-011}}{TASK-A-011} - Aldo & Compose A/B
+ Pgpool y experimento M-M mínimo.
\textbf{\texorpdfstring{\softtext{verificado\_local}}{verificado\_local}}.
Revisión humana pendiente. \\
\texorpdfstring{\softcode{TASK-A-012}}{TASK-A-012} - Aldo & Tercero
simulado y FDW sólo lectura.
\textbf{\texorpdfstring{\softtext{verificado\_local}}{verificado\_local}}.
Revisión humana pendiente. \\
\texorpdfstring{\softcode{TASK-A-013}}{TASK-A-013} - Aldo &
Autenticación y permisos por sede.
\textbf{\texorpdfstring{\softtext{verificado\_local}}{verificado\_local}}.
Revisión humana pendiente. \\
\texorpdfstring{\softcode{TASK-A-014}}{TASK-A-014} - Aldo & Maestros,
animal e identificaciones centrales.
\textbf{\texorpdfstring{\softtext{verificado\_local}}{verificado\_local}}.
Revisión humana pendiente. \\
\texorpdfstring{\softcode{TASK-A-015}}{TASK-A-015} - Aldo & Operaciones,
lotes y versiones
\texorpdfstring{\softtext{contractuales.}}{contractuales.}
\textbf{\texorpdfstring{\softtext{verificado\_local}}{verificado\_local}}.
Revisión humana pendiente. \\
\texorpdfstring{\softcode{TASK-A-016}}{TASK-A-016} - Aldo & Documentos,
revisión y pesajes.
\textbf{\texorpdfstring{\softtext{verificado\_local}}{verificado\_local}}.
Revisión humana pendiente. \\
\texorpdfstring{\softcode{TASK-A-017}}{TASK-A-017} - Aldo &
\texorpdfstring{\softtext{ConfirmarCompra}}{ConfirmarCompra} y asignar
costo inicial.
\textbf{\texorpdfstring{\softtext{verificado\_local}}{verificado\_local}}.
Revisión humana pendiente. \\
\texorpdfstring{\softcode{TASK-A-018}}{TASK-A-018} - Aldo & Movimientos,
recepción, entrega y baja.
\textbf{\texorpdfstring{\softtext{verificado\_local}}{verificado\_local}}.
Revisión humana pendiente. \\
\texorpdfstring{\softcode{TASK-A-019}}{TASK-A-019} - Aldo & Reservas,
expiración y
\texorpdfstring{\softtext{ConfirmarVenta.}}{ConfirmarVenta.}
\textbf{\texorpdfstring{\softtext{verificado\_local}}{verificado\_local}}.
Revisión humana pendiente. \\
\texorpdfstring{\softcode{TASK-A-020}}{TASK-A-020} - Aldo & Pagos,
anticipos y \texorpdfstring{\softtext{aplicaciones.}}{aplicaciones.}
\textbf{\texorpdfstring{\softtext{verificado\_local}}{verificado\_local}}.
Revisión humana pendiente. \\
\texorpdfstring{\softcode{TASK-A-021}}{TASK-A-021} - Aldo & Ajustes,
reversos y \texorpdfstring{\softtext{reembolsos.}}{reembolsos.}
\textbf{\texorpdfstring{\softtext{verificado\_local}}{verificado\_local}}.
Revisión humana pendiente. \\
\texorpdfstring{\softcode{TASK-A-022}}{TASK-A-022} - Aldo & Costo
tardío, valoración y periodos.
\textbf{\texorpdfstring{\softtext{verificado\_local}}{verificado\_local}}.
Revisión humana pendiente. \\
\texorpdfstring{\softcode{TASK-A-023}}{TASK-A-023} - Aldo & Auditoría,
\texorpdfstring{\softtext{outbox/inbox}}{outbox/inbox} y
\texorpdfstring{\softtext{reintentos.}}{reintentos.}
\textbf{\texorpdfstring{\softtext{verificado\_local}}{verificado\_local}}.
Revisión humana pendiente. \\
\texorpdfstring{\softcode{TASK-A-024}}{TASK-A-024} - Aldo & Consultas de
cartera, margen e \texorpdfstring{\softtext{inventario.}}{inventario.}
\textbf{\texorpdfstring{\softtext{verificado\_local}}{verificado\_local}}.
Revisión humana pendiente. \\
\texorpdfstring{\softcode{TASK-A-025}}{TASK-A-025} - Aldo & Integrar
aceptación de fuente FDW y QR.
\textbf{\texorpdfstring{\softtext{verificado\_local}}{verificado\_local}}.
Revisión humana pendiente. \\
\texorpdfstring{\softcode{TASK-A-026}}{TASK-A-026} - Aldo & Pruebas M-M
completas, fallas y
\texorpdfstring{\softtext{conciliación.}}{conciliación.}
\textbf{\texorpdfstring{\softtext{verificado\_local}}{verificado\_local}}.
Revisión humana pendiente. \\
\texorpdfstring{\softcode{TASK-A-027}}{TASK-A-027} - Aldo & Pruebas
acumulativas de
\texorpdfstring{\softtext{dominio/SQL/API}}{dominio/SQL/API} y
\texorpdfstring{\softtext{concurrencia.}}{concurrencia.}
\textbf{\texorpdfstring{\softtext{verificado\_local}}{verificado\_local}}.
Revisión humana pendiente. \\
\texorpdfstring{\softcode{TASK-A-028}}{TASK-A-028} - Aldo & Pipeline,
staging, TLS y
\texorpdfstring{\softtext{configuración.}}{configuración.}
\textbf{\texorpdfstring{\softtext{pendiente\_externo}}{pendiente\_externo}}.
Revisión humana pendiente. \\
\texorpdfstring{\softcode{TASK-A-029}}{TASK-A-029} - Aldo & Respaldo,
restauración y runbooks.
\textbf{\texorpdfstring{\softtext{verificado\_local}}{verificado\_local}}.
Revisión humana pendiente. \\
\texorpdfstring{\softcode{TASK-A-030}}{TASK-A-030} - Aldo & Rendimiento,
conexiones y
\texorpdfstring{\softtext{observabilidad.}}{observabilidad.}
\textbf{\texorpdfstring{\softtext{verificado\_local}}{verificado\_local}}.
Revisión humana pendiente. \\
\texorpdfstring{\softcode{TASK-A-031}}{TASK-A-031} - Aldo & OCR
supervisado si se aprueba el núcleo. \textbf{diferido}. Revisión humana
pendiente. \\
\texorpdfstring{\softcode{TASK-A-032}}{TASK-A-032} - Aldo & Dossier
académico, manual SQL y
\texorpdfstring{\softtext{exposición.}}{exposición.}
\textbf{\texorpdfstring{\softtext{verificado\_local}}{verificado\_local}}.
Revisión humana pendiente. \\
\texorpdfstring{\softcode{TASK-A-033}}{TASK-A-033} - Aldo & Fuentes,
README y Graphify final
\texorpdfstring{\softtext{consolidado.}}{consolidado.}
\textbf{\texorpdfstring{\softtext{verificado\_local}}{verificado\_local}}.
Revisión humana pendiente. \\
\texorpdfstring{\softcode{TASK-J-001}}{TASK-J-001} - Jared & Entender
negocio, actores y términos.
\textbf{\texorpdfstring{\softtext{verificado\_local}}{verificado\_local}}.
Revisión humana pendiente. \\
\texorpdfstring{\softcode{TASK-J-002}}{TASK-J-002} - Jared & Wireframes
y arquitectura de \texorpdfstring{\softtext{información.}}{información.}
\textbf{\texorpdfstring{\softtext{verificado\_local}}{verificado\_local}}.
Revisión humana pendiente. \\
\texorpdfstring{\softcode{TASK-J-003}}{TASK-J-003} - Jared & Tokens y
componentes visuales
\texorpdfstring{\softtext{accesibles.}}{accesibles.}
\textbf{\texorpdfstring{\softtext{verificado\_local}}{verificado\_local}}.
Revisión humana pendiente. \\
\texorpdfstring{\softcode{TASK-J-004}}{TASK-J-004} - Jared &
\texorpdfstring{\softtext{React/TS/Vite,}}{React/TS/Vite,} router y
build.
\textbf{\texorpdfstring{\softtext{verificado\_local}}{verificado\_local}}.
Revisión humana pendiente. \\
\texorpdfstring{\softcode{TASK-J-005}}{TASK-J-005} - Jared & Cliente
tipado, mocks y errores del contrato.
\textbf{\texorpdfstring{\softtext{verificado\_local}}{verificado\_local}}.
Revisión humana pendiente. \\
\texorpdfstring{\softcode{TASK-J-006}}{TASK-J-006} - Jared & Sesión,
navegación y sede \texorpdfstring{\softtext{autorizada.}}{autorizada.}
\textbf{\texorpdfstring{\softtext{verificado\_local}}{verificado\_local}}.
Revisión humana pendiente. \\
\texorpdfstring{\softcode{TASK-J-007}}{TASK-J-007} - Jared &
Contrapartes y predios.
\textbf{\texorpdfstring{\softtext{verificado\_local}}{verificado\_local}}.
Revisión humana pendiente. \\
\texorpdfstring{\softcode{TASK-J-008}}{TASK-J-008} - Jared & Ficha
animal, ciclos e \texorpdfstring{\softtext{inventario.}}{inventario.}
\textbf{\texorpdfstring{\softtext{verificado\_local}}{verificado\_local}}.
Revisión humana pendiente. \\
\texorpdfstring{\softcode{TASK-J-009}}{TASK-J-009} - Jared & Captura de
compra y condiciones por lote.
\textbf{\texorpdfstring{\softtext{verificado\_local}}{verificado\_local}}.
Revisión humana pendiente. \\
\texorpdfstring{\softcode{TASK-J-010}}{TASK-J-010} - Jared & Pesajes,
archivos y revisión.
\textbf{\texorpdfstring{\softtext{verificado\_local}}{verificado\_local}}.
Revisión humana pendiente. \\
\texorpdfstring{\softcode{TASK-J-011}}{TASK-J-011} - Jared & Selección
de venta y reservas.
\textbf{\texorpdfstring{\softtext{verificado\_local}}{verificado\_local}}.
Revisión humana pendiente. \\
\texorpdfstring{\softcode{TASK-J-012}}{TASK-J-012} - Jared & Simulador y
detalle explicable de
\texorpdfstring{\softtext{liquidación.}}{liquidación.}
\textbf{\texorpdfstring{\softtext{verificado\_local}}{verificado\_local}}.
Revisión humana pendiente. \\
\texorpdfstring{\softcode{TASK-J-013}}{TASK-J-013} - Jared &
\texorpdfstring{\softtext{Cobros/pagos}}{Cobros/pagos} y aplicación
guiada.
\textbf{\texorpdfstring{\softtext{verificado\_local}}{verificado\_local}}.
Revisión humana pendiente. \\
\texorpdfstring{\softcode{TASK-J-014}}{TASK-J-014} - Jared & Ajustes,
devoluciones de dinero y periodos.
\textbf{\texorpdfstring{\softtext{verificado\_local}}{verificado\_local}}.
Revisión humana pendiente. \\
\texorpdfstring{\softcode{TASK-J-015}}{TASK-J-015} - Jared &
Recepciones, entregas y expediente temporal.
\textbf{\texorpdfstring{\softtext{verificado\_local}}{verificado\_local}}.
Revisión humana pendiente. \\
\texorpdfstring{\softcode{TASK-J-016}}{TASK-J-016} - Jared & Cartera,
margen y QR interno.
\textbf{\texorpdfstring{\softtext{verificado\_local}}{verificado\_local}}.
Revisión humana pendiente. \\
\texorpdfstring{\softcode{TASK-J-017}}{TASK-J-017} - Jared & Estados de
fuente
\texorpdfstring{\softtext{caída/sincronización/conflicto.}}{caída/sincronización/conflicto.}
\textbf{\texorpdfstring{\softtext{verificado\_local}}{verificado\_local}}.
Revisión humana pendiente. \\
\texorpdfstring{\softcode{TASK-J-018}}{TASK-J-018} - Jared & Responsive,
teclado y \texorpdfstring{\softtext{accesibilidad.}}{accesibilidad.}
\textbf{\texorpdfstring{\softtext{verificado\_local}}{verificado\_local}}.
Revisión humana pendiente. \\
\texorpdfstring{\softcode{TASK-J-019}}{TASK-J-019} - Jared & Pruebas E2E
integradas con Aldo.
\textbf{\texorpdfstring{\softtext{verificado\_local}}{verificado\_local}}.
Revisión humana pendiente. \\
\texorpdfstring{\softcode{TASK-J-020}}{TASK-J-020} - Jared & UI de OCR
sólo si el extra está
\texorpdfstring{\softtext{implementado.}}{implementado.}
\textbf{diferido}. Revisión humana pendiente. \\
\texorpdfstring{\softcode{TASK-J-021}}{TASK-J-021} - Jared & Manual,
capturas auténticas y demo.
\textbf{\texorpdfstring{\softtext{verificado\_local}}{verificado\_local}}.
Revisión humana pendiente. \\
\texorpdfstring{\softcode{TASK-J-022}}{TASK-J-022} - Jared & README,
contexto frontend y actualización de grafo.
\textbf{\texorpdfstring{\softtext{verificado\_local}}{verificado\_local}}.
Revisión humana pendiente. \\
\end{longtable}
}

\chapter{Matriz de pruebas por
escenario}\label{matriz-de-pruebas-por-escenario}

El siguiente registro enlaza cada TEST con evidencia SQL/API/UI y
conserva sus límites. Es el alcance de las ejecuciones registradas, no
pruebas nuevas lanzadas por generar el manual.

\subsection{Matriz de pruebas SQL e infraestructura
ejecutadas}\label{matriz-de-pruebas-sql-e-infraestructura-ejecutadas}

Esta matriz distingue aserciones SQL, infraestructura y solicitudes
HTTP. \texorpdfstring{\softcode{PASS\_SQL}}{PASS\_SQL} nunca certifica
por sí solo API/UI, capacidad, despliegue o35/35 completos. Las
mutaciones SQL usan namespaces QA; GOLD revierte su
\texorpdfstring{\softtext{transacción.}}{transacción.} Main conserva sus
hechos y sólo se
\texorpdfstring{\softtext{consultó/autenticó}}{consultó/autenticó}
durante fallos.

\begin{itemize}
\item
  \textbf{ID:} TEST-01. \textbf{Estado:} PASS\_SQL. \textbf{Evidencia y
  resultado:} GOLD01 costo252000, margen34000, saldo244000;
  \texorpdfstring{\softtext{variación100kg/2\%}}{variación100kg/2\%} sin
  descuento automático
  \texorpdfstring{\softtext{gold-final-all.log}}{gold-final-all.log}
  (referencia del repositorio:
  \texorpdfstring{\softcode{../../../db/tests/gold-final-all.log}}{../../../db/tests/gold-final-all.log}).
  \textbf{Límite:} HTTP/UI se certifican con evidencia de sus
  \texorpdfstring{\softtext{ejecutores..}}{ejecutores..}
\item
  \textbf{ID:} TEST-02. \textbf{Estado:} PASS\_SQL. \textbf{Evidencia y
  resultado:} GOLD02 costo4400, margen1800,
  \texorpdfstring{\softtext{inventario6600.}}{inventario6600.}
  \texorpdfstring{\softtext{gold-final-all.log}}{gold-final-all.log}
  (referencia del repositorio:
  \texorpdfstring{\softcode{../../../db/tests/gold-final-all.log}}{../../../db/tests/gold-final-all.log}).
  \textbf{Límite:} -.
\item
  \textbf{ID:} TEST-03. \textbf{Estado:} PASS\_SQL. \textbf{Evidencia y
  resultado:} GOLD03 reparto
  \texorpdfstring{\softtext{exacto0.34/0.33/0.33;}}{exacto0.34/0.33/0.33;}
  NUMERIC real.
  \texorpdfstring{\softtext{gold-final-all.log}}{gold-final-all.log}
  (referencia del repositorio:
  \texorpdfstring{\softcode{../../../db/tests/gold-final-all.log}}{../../../db/tests/gold-final-all.log}).
  \textbf{Límite:} No sustituye prueba de unidad Java del transporte
  decimal..
\item
  \textbf{ID:} TEST-04. \textbf{Estado:} PASS\_SQL. \textbf{Evidencia y
  resultado:} GOLD04A
  \texorpdfstring{\softtext{deuda1000/pago600/aplicaciones400+200}}{deuda1000/pago600/aplicaciones400+200}
  y replayexacto; GOLD04B
  \texorpdfstring{\softtext{deuda1000/pago1200/aplicado1000/anticipo200.}}{deuda1000/pago1200/aplicado1000/anticipo200.}
  \texorpdfstring{\softtext{gold-final-all.log}}{gold-final-all.log}
  (referencia del repositorio:
  \texorpdfstring{\softcode{../../../db/tests/gold-final-all.log}}{../../../db/tests/gold-final-all.log}).
  \textbf{Límite:} -.
\item
  \textbf{ID:} TEST-05. \textbf{Estado:} PASS\_SQL. \textbf{Evidencia y
  resultado:} Entrega cierra primer ciclo; recompra crea ciclo2; costo
  histórico1200 \texorpdfstring{\softtext{conservado.}}{conservado.}
  \texorpdfstring{\softtext{gold-final-all.log}}{gold-final-all.log}
  (referencia del repositorio:
  \texorpdfstring{\softcode{../../../db/tests/gold-final-all.log}}{../../../db/tests/gold-final-all.log}).
  \textbf{Límite:} -.
\item
  \textbf{ID:} TEST-06. \textbf{Estado:} PASS\_SQL. \textbf{Evidencia y
  resultado:} GOLD06
  \texorpdfstring{\softtext{tardío4600/6900/1600;}}{tardío4600/6900/1600;}
  \texorpdfstring{\softtext{bonificación4520/6780/1680;}}{bonificación4520/6780/1680;}
  deuda inicial inmutable.
  \texorpdfstring{\softtext{gold-final-all.log}}{gold-final-all.log}
  (referencia del repositorio:
  \texorpdfstring{\softcode{../../../db/tests/gold-final-all.log}}{../../../db/tests/gold-final-all.log}).
  \textbf{Límite:} -.
\item
  \textbf{ID:} TEST-07. \textbf{Estado:} PASS\_SQL. \textbf{Evidencia y
  resultado:} GOLD07 ajuste atómico:
  \texorpdfstring{\softtext{reverso1000/reaplicación900/abono100;}}{reverso1000/reaplicación900/abono100;}
  disponible100,
  \texorpdfstring{\softtext{reembolso100→0.}}{reembolso100→0.}
  \texorpdfstring{\softtext{gold-final-all.log}}{gold-final-all.log}
  (referencia del repositorio:
  \texorpdfstring{\softcode{../../../db/tests/gold-final-all.log}}{../../../db/tests/gold-final-all.log}).
  \textbf{Límite:} -.
\item
  \textbf{ID:} TEST-08. \textbf{Estado:} PASS\_SQL. \textbf{Evidencia y
  resultado:} GOLD08 baja1100 y costo tardío50 producen
  \texorpdfstring{\softtext{pérdida1150/inventario0.}}{pérdida1150/inventario0.}
  \texorpdfstring{\softtext{gold-final-all.log}}{gold-final-all.log}
  (referencia del repositorio:
  \texorpdfstring{\softcode{../../../db/tests/gold-final-all.log}}{../../../db/tests/gold-final-all.log}).
  \textbf{Límite:} -.
\item
  \textbf{ID:} TEST-09. \textbf{Estado:} PASS\_SQL. \textbf{Evidencia y
  resultado:} Ambos órdenes: carrera cierre primero y coste primero
  forzado con \texorpdfstring{\softtext{barrera/espera}}{barrera/espera}
  Lock observada.
  \texorpdfstring{\softtext{Inventario200→207/cierre207}}{Inventario200→207/cierre207}
  con
  \texorpdfstring{\softtext{documental2027-02-01/reconocimiento2026-10-11;}}{documental2027-02-01/reconocimiento2026-10-11;}
  versión anterior intacta
  \texorpdfstring{\softtext{trasreabrir.}}{trasreabrir.}
  \texorpdfstring{\softtext{academic-result.json}}{academic-result.json}
  (referencia del repositorio:
  \texorpdfstring{\softcode{../../../db/tests/academic-result.json}}{../../../db/tests/academic-result.json}),
  \texorpdfstring{\softtext{academic.log}}{academic.log} (referencia del
  repositorio:
  \texorpdfstring{\softcode{../../../db/tests/academic.log}}{../../../db/tests/academic.log}),
  \texorpdfstring{\softtext{closure-order-result.json}}{closure-order-result.json}
  (referencia del repositorio:
  \texorpdfstring{\softcode{../../../db/tests/closure-order-result.json}}{../../../db/tests/closure-order-result.json}),
  \texorpdfstring{\softtext{closure-order.log}}{closure-order.log}
  (referencia del repositorio:
  \texorpdfstring{\softcode{../../../db/tests/closure-order.log}}{../../../db/tests/closure-order.log}).
  \textbf{Límite:} -.
\item
  \textbf{ID:} TEST-10. \textbf{Estado:} PASS\_SQL. \textbf{Evidencia y
  resultado:} Dos
  \texorpdfstring{\softtext{conexiones/barrera:}}{conexiones/barrera:}
  una \texorpdfstring{\softtext{reserva/venta}}{reserva/venta} gana;
  otra23505; sin segunda
  \texorpdfstring{\softtext{liquidación.}}{liquidación.}
  \texorpdfstring{\softtext{academic-result.json}}{academic-result.json}
  (referencia del repositorio:
  \texorpdfstring{\softcode{../../../db/tests/academic-result.json}}{../../../db/tests/academic-result.json}),
  \texorpdfstring{\softtext{academic.log}}{academic.log} (referencia del
  repositorio:
  \texorpdfstring{\softcode{../../../db/tests/academic.log}}{../../../db/tests/academic.log}).
  \textbf{Límite:} Mapeo HTTP409 pertenece a API..
\item
  \textbf{ID:} TEST-11. \textbf{Estado:} PASS\_SQL. \textbf{Evidencia y
  resultado:} Dos
  \texorpdfstring{\softtext{conexiones/barrera}}{conexiones/barrera}
  compiten por pago600: sólo aplica400, otraBV020,
  \texorpdfstring{\softtext{disponible200.}}{disponible200.}
  \texorpdfstring{\softtext{academic-result.json}}{academic-result.json}
  (referencia del repositorio:
  \texorpdfstring{\softcode{../../../db/tests/academic-result.json}}{../../../db/tests/academic-result.json}),
  \texorpdfstring{\softtext{academic.log}}{academic.log} (referencia del
  repositorio:
  \texorpdfstring{\softcode{../../../db/tests/academic.log}}{../../../db/tests/academic.log}).
  \textbf{Límite:} -.
\item
  \textbf{ID:} TEST-12. \textbf{Estado:} PASS\_SQL. \textbf{Evidencia y
  resultado:} Misma clave/hash replayexacto; payload
  \texorpdfstring{\softtext{distintoBV405}}{distintoBV405} y
  \texorpdfstring{\softtext{ETag/version}}{ETag/version}
  \texorpdfstring{\softtext{obsoletaBV412.}}{obsoletaBV412.}
  \texorpdfstring{\softtext{Cancelar/anular/acordar/reabrir:}}{Cancelar/anular/acordar/reabrir:}
  \texorpdfstring{\softtext{stale412sinidem/audit;}}{stale412sinidem/audit;}
  replayexacto trascommit
  \texorpdfstring{\softtext{antesguardversionactual.}}{antesguardversionactual.}
  \texorpdfstring{\softtext{gold-final-all.log}}{gold-final-all.log}
  (referencia del repositorio:
  \texorpdfstring{\softcode{../../../db/tests/gold-final-all.log}}{../../../db/tests/gold-final-all.log}),
  \texorpdfstring{\softtext{version-guards-result.json}}{version-guards-result.json}
  (referencia del repositorio:
  \texorpdfstring{\softcode{../../../db/tests/version-guards-result.json}}{../../../db/tests/version-guards-result.json}).
  \textbf{Límite:} SQL no certifica por sí solo
  \texorpdfstring{\softtext{HTTP409/412..}}{HTTP409/412..}
\item
  \textbf{ID:} TEST-13. \textbf{Estado:} PASS\_SQL. \textbf{Evidencia y
  resultado:} Dos
  \texorpdfstring{\softtext{conexiones/barrera,}}{conexiones/barrera,}
  scopes de sedesA/B hacia escritor
  \texorpdfstring{\softtext{canónicoA:}}{canónicoA:} un identificador
  normalizado único; otro23505.
  \texorpdfstring{\softtext{negatives-result.json}}{negatives-result.json}
  (referencia del repositorio:
  \texorpdfstring{\softcode{../../../db/tests/negatives-result.json}}{../../../db/tests/negatives-result.json}),
  \texorpdfstring{\softtext{negatives.log}}{negatives.log} (referencia
  del repositorio:
  \texorpdfstring{\softcode{../../../db/tests/negatives.log}}{../../../db/tests/negatives.log}).
  \textbf{Límite:} -.
\item
  \textbf{ID:} TEST-14. \textbf{Estado:} PASS\_SQL. \textbf{Evidencia y
  resultado:} FK otra sede23503; empresa ajena42501; ningún hecho
  cruzado.
  \texorpdfstring{\softtext{negatives-result.json}}{negatives-result.json}
  (referencia del repositorio:
  \texorpdfstring{\softcode{../../../db/tests/negatives-result.json}}{../../../db/tests/negatives-result.json}),
  \texorpdfstring{\softtext{negatives.log}}{negatives.log} (referencia
  del repositorio:
  \texorpdfstring{\softcode{../../../db/tests/negatives.log}}{../../../db/tests/negatives.log}).
  \textbf{Límite:} -.
\item
  \textbf{ID:} TEST-15. \textbf{Estado:}
  \texorpdfstring{\softtext{PASS\_SQL\_API\_UI.}}{PASS\_SQL\_API\_UI.}
  \textbf{Evidencia y resultado:} Scope empresa ajena42501 y token de
  otro \texorpdfstring{\softtext{actor/sede404;}}{actor/sede404;}
  control canónico.
  \texorpdfstring{\softtext{negatives-result.json}}{negatives-result.json}
  (referencia del repositorio:
  \texorpdfstring{\softcode{../../../db/tests/negatives-result.json}}{../../../db/tests/negatives-result.json}),
  \texorpdfstring{\softtext{negatives.log}}{negatives.log} (referencia
  del repositorio:
  \texorpdfstring{\softcode{../../../db/tests/negatives.log}}{../../../db/tests/negatives.log}),
  \texorpdfstring{\softtext{owner-capture-result.json}}{owner-capture-result.json}
  (referencia del repositorio:
  \texorpdfstring{\softcode{../../../db/tests/owner-capture-result.json}}{../../../db/tests/owner-capture-result.json}).
  \textbf{Límite:} API valida scope y AUDITOR403; UI crea
  \texorpdfstring{\softtext{auditor/membresía}}{auditor/membresía}
  reales y comprueba rechazo. Backend (referencia del repositorio:
  \texorpdfstring{\softcode{../../../backend/verification-result.json}}{../../../backend/verification-result.json}),
  UI (referencia del repositorio:
  \texorpdfstring{\softcode{ui/bovicierre-ui-final-run.json}}{ui/bovicierre-ui-final-run.json})..
\item
  \textbf{ID:} TEST-16. \textbf{Estado:} PASS\_SQL. \textbf{Evidencia y
  resultado:} Rol app
  \texorpdfstring{\softtext{sinDML/DDL/SETsession\_replication\_role;}}{sinDML/DDL/SETsession\_replication\_role;}
  intento desactivar trigger42501; PUBLIC EXECUTE revocado. DML real
  como \texorpdfstring{\softtext{rolbovi\_auditor}}{rolbovi\_auditor}
  \texorpdfstring{\softtext{también42501.}}{también42501.}
  \texorpdfstring{\softtext{negatives-result.json}}{negatives-result.json}
  (referencia del repositorio:
  \texorpdfstring{\softcode{../../../db/tests/negatives-result.json}}{../../../db/tests/negatives-result.json}),
  \texorpdfstring{\softtext{negatives.log}}{negatives.log} (referencia
  del repositorio:
  \texorpdfstring{\softcode{../../../db/tests/negatives.log}}{../../../db/tests/negatives.log}),
  \texorpdfstring{\softtext{security.sql}}{security.sql} (referencia del
  repositorio:
  \texorpdfstring{\softcode{../../../db/tests/security.sql}}{../../../db/tests/security.sql}),
  \texorpdfstring{\softtext{security-final.log}}{security-final.log}
  (referencia del repositorio:
  \texorpdfstring{\softcode{../../../db/tests/security-final.log}}{../../../db/tests/security-final.log}).
  \textbf{Límite:} La prueba usa roles técnicos; políticas de auditor
  humano se verifican en API/UI..
\item
  \textbf{ID:} TEST-17. \textbf{Estado:} PASS\_SQL. \textbf{Evidencia y
  resultado:} Como rol dueño de funciones, UPDATE a
  \texorpdfstring{\softtext{liquidación/contrato}}{liquidación/contrato}
  \texorpdfstring{\softtext{aceptado/condición/pesaje}}{aceptado/condición/pesaje}
  aprobado \texorpdfstring{\softtext{produceBV016:}}{produceBV016:}
  protección independiente del ACL app.
  \texorpdfstring{\softtext{negatives-result.json}}{negatives-result.json}
  (referencia del repositorio:
  \texorpdfstring{\softcode{../../../db/tests/negatives-result.json}}{../../../db/tests/negatives-result.json}),
  \texorpdfstring{\softtext{negatives.log}}{negatives.log} (referencia
  del repositorio:
  \texorpdfstring{\softcode{../../../db/tests/negatives.log}}{../../../db/tests/negatives.log}).
  \textbf{Límite:} -.
\item
  \textbf{ID:} TEST-18. \textbf{Estado:} PASS\_SQL. \textbf{Evidencia y
  resultado:} Pesaje grupal no fabrica pesos individuales; reparto por
  peso \texorpdfstring{\softtext{ausenteBV018.}}{ausenteBV018.}
  \texorpdfstring{\softtext{gold-final-all.log}}{gold-final-all.log}
  (referencia del repositorio:
  \texorpdfstring{\softcode{../../../db/tests/gold-final-all.log}}{../../../db/tests/gold-final-all.log}).
  \textbf{Límite:} -.
\item
  \textbf{ID:} TEST-19. \textbf{Estado:} PASS\_SQL. \textbf{Evidencia y
  resultado:} \texorpdfstring{\softtext{DobleEXCEPT}}{DobleEXCEPT} de
  composiciones distintas
  \texorpdfstring{\softtext{produceNO\_COMPARABLE;}}{produceNO\_COMPARABLE;}
  variación comparable exacta enGOLD01.
  \texorpdfstring{\softtext{gold-final-all.log}}{gold-final-all.log}
  (referencia del repositorio:
  \texorpdfstring{\softcode{../../../db/tests/gold-final-all.log}}{../../../db/tests/gold-final-all.log}).
  \textbf{Límite:} -.
\item
  \textbf{ID:} TEST-20. \textbf{Estado:} PASS\_SQL. \textbf{Evidencia y
  resultado:} Dos
  \texorpdfstring{\softtext{conexiones/barrera,}}{conexiones/barrera,}
  esperaLock observada, tx empezó antes de vencimiento y desbloqueó
  después:BV403, sin
  \texorpdfstring{\softtext{liquidación/idempotencia}}{liquidación/idempotencia}
  parcial.
  \texorpdfstring{\softtext{expiry-result.json}}{expiry-result.json}
  (referencia del repositorio:
  \texorpdfstring{\softcode{../../../db/tests/expiry-result.json}}{../../../db/tests/expiry-result.json}),
  expiry.log (referencia del repositorio:
  \texorpdfstring{\softcode{../../../db/tests/expiry.log}}{../../../db/tests/expiry.log}).
  \textbf{Límite:} -.
\item
  \textbf{ID:} TEST-21. \textbf{Estado:}
  \texorpdfstring{\softtext{PASS\_INFRA\_QA.}}{PASS\_INFRA\_QA.}
  \textbf{Evidencia y resultado:} Publicaciones bidireccionales reales;
  \texorpdfstring{\softtext{DISABLE/ENABLE}}{DISABLE/ENABLE} de
  suscripción QA y convergencia; conteos
  \texorpdfstring{\softtext{audit/outbox}}{audit/outbox} sin
  \texorpdfstring{\softtext{duplicación.}}{duplicación.}
  \texorpdfstring{\softtext{academic-result.json}}{academic-result.json}
  (referencia del repositorio:
  \texorpdfstring{\softcode{../../../db/tests/academic-result.json}}{../../../db/tests/academic-result.json}),
  \texorpdfstring{\softtext{academic.log}}{academic.log} (referencia del
  repositorio:
  \texorpdfstring{\softcode{../../../db/tests/academic.log}}{../../../db/tests/academic.log}).
  \textbf{Límite:} Partición se simuló en stream apply QA; no se
  interrumpió red física..
\item
  \textbf{ID:} TEST-22. \textbf{Estado:}
  \texorpdfstring{\softtext{PASS\_SQL\_INFRA\_HTTP.}}{PASS\_SQL\_INFRA\_HTTP.}
  \textbf{Evidencia y resultado:} A source parado real: HTTP autenticado
  nuevas
  \texorpdfstring{\softtext{solicitudesA503/B503;}}{solicitudesA503/B503;}
  B transacción autorizada antes de caída confirma hechos propios; A
  reanudado enfinally.
  \texorpdfstring{\softtext{outages-result.json}}{outages-result.json}
  (referencia del repositorio:
  \texorpdfstring{\softcode{../../../db/tests/outages-result.json}}{../../../db/tests/outages-result.json}),
  \texorpdfstring{\softtext{outages.log}}{outages.log} (referencia del
  repositorio:
  \texorpdfstring{\softcode{../../../db/tests/outages.log}}{../../../db/tests/outages.log}).
  \textbf{Límite:} SQL financiero de prueba usa
  \texorpdfstring{\softtext{namespaceQA;}}{namespaceQA;}
  \texorpdfstring{\softtext{solicitudesHTTPmain}}{solicitudesHTTPmain}
  fueron \texorpdfstring{\softtext{consultas/auth.}}{consultas/auth.} No
  \texorpdfstring{\softtext{promoción..}}{promoción..}
\item
  \textbf{ID:} TEST-23. \textbf{Estado:} PASS\_SQL. \textbf{Evidencia y
  resultado:} Conciliación explícita en ambos dueños:0 huérfanos y0
  liquidaciones desbalanceadas sobreFK
  \texorpdfstring{\softtext{comerciales.}}{comerciales.}
  \texorpdfstring{\softtext{academic-result.json}}{academic-result.json}
  (referencia del repositorio:
  \texorpdfstring{\softcode{../../../db/tests/academic-result.json}}{../../../db/tests/academic-result.json}),
  \texorpdfstring{\softtext{academic.log}}{academic.log} (referencia del
  repositorio:
  \texorpdfstring{\softcode{../../../db/tests/academic.log}}{../../../db/tests/academic.log}).
  \textbf{Límite:} -.
\item
  \textbf{ID:} TEST-24. \textbf{Estado:}
  \texorpdfstring{\softtext{PASS\_SQL\_INFRA.}}{PASS\_SQL\_INFRA.}
  \textbf{Evidencia y resultado:} A caída real: fallo entrega guarda
  \texorpdfstring{\softtext{intento/backoff}}{intento/backoff} y limpia
  lease; lease abandonada30s recuperada por worker nuevo;
  \texorpdfstring{\softtext{ACKantiguaBV402.}}{ACKantiguaBV402.}
  ConsumeA dos \texorpdfstring{\softtext{veces/inbox1,}}{veces/inbox1,}
  acuseB; perdida1100
  \texorpdfstring{\softtext{+costetardío1.}}{+costetardío1.}
  \texorpdfstring{\softtext{academic-result.json}}{academic-result.json}
  (referencia del repositorio:
  \texorpdfstring{\softcode{../../../db/tests/academic-result.json}}{../../../db/tests/academic-result.json}),
  \texorpdfstring{\softtext{outages-result.json}}{outages-result.json}
  (referencia del repositorio:
  \texorpdfstring{\softcode{../../../db/tests/outages-result.json}}{../../../db/tests/outages-result.json}),
  \texorpdfstring{\softtext{outages.log}}{outages.log} (referencia del
  repositorio:
  \texorpdfstring{\softcode{../../../db/tests/outages.log}}{../../../db/tests/outages.log}).
  \textbf{Límite:} Worker abandonado simulado dejando lease vencer; no
  se mata un \texorpdfstring{\softtext{procesoJava.}}{procesoJava.}
  ACKperdido simulado omitiendo primer acuse..
\item
  \textbf{ID:} TEST-25. \textbf{Estado:}
  \texorpdfstring{\softtext{PASS\_SQL\_INFRA\_HTTP.}}{PASS\_SQL\_INFRA\_HTTP.}
  \textbf{Evidencia y resultado:} FDW realA viaPgpool
  \texorpdfstring{\softtext{SELECT12000/UPDATE42501;}}{SELECT12000/UPDATE42501;}
  B sin servidor proveedor. Proveedor source parado
  \texorpdfstring{\softtext{devuelveHTTP503}}{devuelveHTTP503} y gastoB
  usa captura aceptada local sin red.
  \texorpdfstring{\softtext{Frontera12000.0050→Money12000.01}}{Frontera12000.0050→Money12000.01}
  aceptada, raw4/hash
  \texorpdfstring{\softtext{preservados.}}{preservados.}
  \texorpdfstring{\softtext{academic.log}}{academic.log} (referencia del
  repositorio:
  \texorpdfstring{\softcode{../../../db/tests/academic.log}}{../../../db/tests/academic.log}),
  \texorpdfstring{\softtext{provider.log}}{provider.log} (referencia del
  repositorio:
  \texorpdfstring{\softcode{../../../db/tests/provider.log}}{../../../db/tests/provider.log}),
  \texorpdfstring{\softtext{owner-capture-result.json}}{owner-capture-result.json}
  (referencia del repositorio:
  \texorpdfstring{\softcode{../../../db/tests/owner-capture-result.json}}{../../../db/tests/owner-capture-result.json}),
  \texorpdfstring{\softtext{provider-precision-result.json}}{provider-precision-result.json}
  (referencia del repositorio:
  \texorpdfstring{\softcode{../../../db/tests/provider-precision-result.json}}{../../../db/tests/provider-precision-result.json}),
  \texorpdfstring{\softtext{outages-result.json}}{outages-result.json}
  (referencia del repositorio:
  \texorpdfstring{\softcode{../../../db/tests/outages-result.json}}{../../../db/tests/outages-result.json}).
  \textbf{Límite:} FuenteQA modificada para fronteras y
  restauradafinally; stop real afecta sólo container propio
  \texorpdfstring{\softtext{proveedor..}}{proveedor..}
\item
  \textbf{ID:} TEST-26. \textbf{Estado:}
  \texorpdfstring{\softtext{PASS\_RECUPERACION\_HTTP.}}{PASS\_RECUPERACION\_HTTP.}
  \textbf{Evidencia y resultado:} Coordinador: restauración final V13
  aislada28,616s,
  \texorpdfstring{\softtext{datos/hash/archivos/dedup/inbox,0sesiones;}}{datos/hash/archivos/dedup/inbox,0sesiones;}
  \texorpdfstring{\softtext{compra1000/replay/lookup/descarga}}{compra1000/replay/lookup/descarga}
  HTTP. Registro (referencia del repositorio:
  \texorpdfstring{\softcode{recovery/restore-v13.json}}{recovery/restore-v13.json}).
  \textbf{Límite:} Mismo host; no PITR ni backup fuera del host..
\item
  \textbf{ID:} TEST-27. \textbf{Estado:}
  \texorpdfstring{\softtext{PASS\_INFRA.}}{PASS\_INFRA.}
  \textbf{Evidencia y resultado:} Recreación real sólo servicio pg-a,
  IDnuevo, mismo volumen; firmas completas datos main+QA
  \texorpdfstring{\softtext{antes/después}}{antes/después} idénticas,
  \texorpdfstring{\softtext{incluyehistory/localcapturas/dedup/inbox/leases.}}{incluyehistory/localcapturas/dedup/inbox/leases.}
  \texorpdfstring{\softtext{persistence-result.json}}{persistence-result.json}
  (referencia del repositorio:
  \texorpdfstring{\softcode{../../../db/tests/persistence-result.json}}{../../../db/tests/persistence-result.json}),
  \texorpdfstring{\softtext{persistence.log}}{persistence.log}
  (referencia del repositorio:
  \texorpdfstring{\softcode{../../../db/tests/persistence.log}}{../../../db/tests/persistence.log}).
  \textbf{Límite:} Se recrea nodoA; se excluyen sesiones auth efímeras
  de firma. \texorpdfstring{\softtext{Tiempo1.744s,}}{Tiempo1.744s,} sin
  pérdida de datos..
\item
  \textbf{ID:} TEST-28. \textbf{Estado:}
  \texorpdfstring{\softtext{PASS\_MEDICION\_LOCAL.}}{PASS\_MEDICION\_LOCAL.}
  \textbf{Evidencia y resultado:} 4EXPLAINANALYZEBUFFERS
  \texorpdfstring{\softtext{y100muestras/consulta}}{y100muestras/consulta}
  tras5warmups:
  \texorpdfstring{\softtext{inventario/cartera/liquidación/expediente,}}{inventario/cartera/liquidación/expediente,}
  dataset declarado.
  \texorpdfstring{\softtext{performance-result.json}}{performance-result.json}
  (referencia del repositorio:
  \texorpdfstring{\softcode{../../../db/tests/performance-result.json}}{../../../db/tests/performance-result.json}),
  \texorpdfstring{\softtext{performance.log}}{performance.log}
  (referencia del repositorio:
  \texorpdfstring{\softcode{../../../db/tests/performance.log}}{../../../db/tests/performance.log}).
  \textbf{Límite:} Dataset
  \texorpdfstring{\softtext{mayor:1901comandos/replay,800consultas}}{mayor:1901comandos/replay,800consultas}
  canónicas \texorpdfstring{\softtext{mediantePgpool}}{mediantePgpool}
  \texorpdfstring{\softtext{con1/4clientes;}}{con1/4clientes;} benchmark
  (referencia del repositorio:
  \texorpdfstring{\softcode{../../../db/tests/benchmark-result.json}}{../../../db/tests/benchmark-result.json}).
  Caché caliente,
  \texorpdfstring{\softtext{sinHTTP/carga}}{sinHTTP/carga} escritora; no
  \texorpdfstring{\softtext{certificaciónSLO/capacidad..}}{certificaciónSLO/capacidad..}
\item
  \textbf{ID:} TEST-29. \textbf{Estado:}
  \texorpdfstring{\softtext{PASS\_UI\_ACOTADO.}}{PASS\_UI\_ACOTADO.}
  \textbf{Evidencia y resultado:} Dieciséis E2E:15integrados más caída
  física de proveedor
  \texorpdfstring{\softtext{reales,13unit/build;390/768/1440}}{reales,13unit/build;390/768/1440}
  y teclado, Intel real. Registro (referencia del repositorio:
  \texorpdfstring{\softcode{ui/bovicierre-ui-final-run.json}}{ui/bovicierre-ui-final-run.json}),
  \texorpdfstring{\softtext{capturas/alcance}}{capturas/alcance}
  (referencia del repositorio:
  \texorpdfstring{\softcode{ui/frontend-ui-tests.md}}{ui/frontend-ui-tests.md}).
  \textbf{Límite:} No auditoría WCAG completa ni teléfono físico..
\item
  \textbf{ID:} TEST-30. \textbf{Estado:}
  \texorpdfstring{\softtext{PASS\_CATALOGO\_VISUAL.}}{PASS\_CATALOGO\_VISUAL.}
  \textbf{Evidencia y resultado:} Baseline3 comparado antes adopción;
  V13, metadatos42tablas83FK vsA/B sin diferencias; cinco tests de
  diagramas y PDF/SVG
  \texorpdfstring{\softtext{MER/relacional}}{MER/relacional} revisados.
  Adopción (referencia del repositorio:
  \texorpdfstring{\softcode{../../../db/tests/flyway-adoption.json}}{../../../db/tests/flyway-adoption.json}),
  diagramas (referencia del repositorio:
  \texorpdfstring{\softcode{../../diagrams/README.md}}{../../diagrams/README.md}).
  \textbf{Límite:} No demuestra todas las RN; revisión semántica humana
  \texorpdfstring{\softtext{pendiente..}}{pendiente..}
\item
  \textbf{ID:} TEST-31. \textbf{Estado:}
  \texorpdfstring{\softtext{PASS\_GRAFO\_FUENTES\_CAMBIADAS.}}{PASS\_GRAFO\_FUENTES\_CAMBIADAS.}
  \textbf{Evidencia y resultado:} AST real y53documentos canónicos
  reinterpretados; rutas
  \texorpdfstring{\softtext{Pantalla8/API,}}{Pantalla8/API,}
  \texorpdfstring{\softtext{RN029/comandoV13/TEST11}}{RN029/comandoV13/TEST11}
  y \texorpdfstring{\softtext{tabla/FK/REL20}}{tabla/FK/REL20}
  \texorpdfstring{\softtext{contrastadas.}}{contrastadas.} Registro
  (referencia del repositorio:
  \texorpdfstring{\softcode{graph-final.json}}{graph-final.json}).
  \textbf{Límite:} Históricos sin cambios conservan extracción anterior;
  aristas no sustituyen SQL ni revisión humana..
\item
  \textbf{ID:} TEST-32. \textbf{Estado:} PASS\_SQL. \textbf{Evidencia y
  resultado:} Snapshot de
  \texorpdfstring{\softtext{conteo/liquidaciones/obligaciones/conceptos/asientos/inventario/costo}}{conteo/liquidaciones/obligaciones/conceptos/asientos/inventario/costo}
  sin cambio tras
  \texorpdfstring{\softtext{registrar/validar/aplicar}}{registrar/validar/aplicar}
  pago, documento y entrega.
  \texorpdfstring{\softtext{gold-final-all.log}}{gold-final-all.log}
  (referencia del repositorio:
  \texorpdfstring{\softcode{../../../db/tests/gold-final-all.log}}{../../../db/tests/gold-final-all.log}).
  \textbf{Límite:} Documento SQL usa metadatos sintéticos;
  \texorpdfstring{\softtext{descarga/generaciónPDF}}{descarga/generaciónPDF}
  real se certifica aparte en API/UI..
\item
  \textbf{ID:} TEST-33. \textbf{Estado:} PASS\_SQL. \textbf{Evidencia y
  resultado:} Muerte centralA converge aB; recompra del animal
  \texorpdfstring{\softtext{muertoBV011}}{muertoBV011} sin
  \texorpdfstring{\softtext{liquidación.}}{liquidación.}
  \texorpdfstring{\softtext{academic-result.json}}{academic-result.json}
  (referencia del repositorio:
  \texorpdfstring{\softcode{../../../db/tests/academic-result.json}}{../../../db/tests/academic-result.json}),
  \texorpdfstring{\softtext{academic.log}}{academic.log} (referencia del
  repositorio:
  \texorpdfstring{\softcode{../../../db/tests/academic.log}}{../../../db/tests/academic.log}).
  \textbf{Límite:} -.
\item
  \textbf{ID:} TEST-34. \textbf{Estado:} PASS\_SQL. \textbf{Evidencia y
  resultado:} Documento v2 no altera v1/hash ni revisión anterior;
  revisiones vinculadas a versión exacta.
  \texorpdfstring{\softtext{gold-final-all.log}}{gold-final-all.log}
  (referencia del repositorio:
  \texorpdfstring{\softcode{../../../db/tests/gold-final-all.log}}{../../../db/tests/gold-final-all.log}).
  \textbf{Límite:} SQL metadata only; bytes
  \texorpdfstring{\softtext{privados/versionado}}{privados/versionado}
  archivo probados por API/UI..
\item
  \textbf{ID:} TEST-35. \textbf{Estado:} PASS\_SQL. \textbf{Evidencia y
  resultado:} Anulación antes entrega con tardío restaura1200 a
  inventario; costo venta neto0; ciclo histórico
  \texorpdfstring{\softtext{conservado.}}{conservado.}
  \texorpdfstring{\softtext{gold-final-all.log}}{gold-final-all.log}
  (referencia del repositorio:
  \texorpdfstring{\softcode{../../../db/tests/gold-final-all.log}}{../../../db/tests/gold-final-all.log}).
  \textbf{Límite:} -.
\end{itemize}

\subsubsection{Ejecuciones y alcance}\label{ejecuciones-y-alcance}

GOLD último:157 llamadas a comandos en
\texorpdfstring{\softtext{PostgreSQL18}}{PostgreSQL18}
\texorpdfstring{\softtext{real/NUMERIC;}}{real/NUMERIC;} rollback
completo. Carrera coste-primero:7
\texorpdfstring{\softtext{comandos/0.107}}{comandos/0.107} s. Guardas
versiones:22 \texorpdfstring{\softtext{comandos/0.140}}{comandos/0.140}
s. Outages:9
\texorpdfstring{\softtext{comandos/52.177}}{comandos/52.177} s;
\texorpdfstring{\softtext{fuentesA/B/proveedor}}{fuentesA/B/proveedor}
reanudadas enfinally.
\texorpdfstring{\softtext{RecreaciónPG}}{RecreaciónPG} A:1.744 s y
\texorpdfstring{\softtext{firmasmain/QA}}{firmasmain/QA} iguales. Estas
cifras no expresan cobertura niSLO.

\texorpdfstring{\softtext{TEST15/26/29/30}}{TEST15/26/29/30} incorporan
evidencia definitiva de
\texorpdfstring{\softtext{coordinador/backend/frontend;}}{coordinador/backend/frontend;}
TEST31 integra los53 documentos canónicos cambiados, conservando
extracciones históricas sin cambios. TEST28 incorpora cuatro planes y800
consultas canónicas con uno y cuatro clientes en QA; caché caliente, sin
HTTP ni escrituras concurrentes y sin capacidad
\texorpdfstring{\softtext{certificada.}}{certificada.} La restauración
final V13 fue ejecutada por el coordinador y concilió datos y archivos;
origen reanudado antes del ensayo.

\texorpdfstring{\softtext{Fallos/recreación}}{Fallos/recreación} se
ejecutaron en ventana coordinada sin otros escritores, verificando
labels
\texorpdfstring{\softtext{exactasbovicierre-academic.}}{exactasbovicierre-academic.}
El harness no administra contenedores de otros proyectos. Capturas y
metadatos documentales SQL son sintéticos;
\texorpdfstring{\softtext{bytes/documentos}}{bytes/documentos} reales se
acreditan con API/UI.

\chapter{Recorrido visual de la
aplicación}\label{recorrido-visual-de-la-aplicaciuxf3n}

Las capturas siguientes proceden de las pruebas de la aplicación
integrada con datos sintéticos
\texorpdfstring{\softtext{persistidos.}}{persistidos.} Son evidencia del
corte anterior, no maquetas ni una visita nueva a una empresa. La
explicación detallada del frontend identifica sus campos, eventos y
llamadas.

\section{Bandeja de operaciones}\label{bandeja-de-operaciones}

Permite localizar operaciones y
\texorpdfstring{\softtext{pendientes.}}{pendientes.} El estado de una
venta se distingue del saldo y del estado físico.

\begin{figure}[htbp]
\centering
\includegraphics[width=.95\linewidth,height=.70\textheight,keepaspectratio]{\detokenize{@@BOVI_ROOT@@/docs/testing/evidence/ui/bovicierre-ui-bandeja-final-1440.png}}
\caption{Bandeja de operaciones con datos DEMO.}
\end{figure}
\clearpage

\section{Recorrido de operación}\label{recorrido-de-operaciuxf3n}

Los datos provienen de backend y
\texorpdfstring{\softtext{PostgreSQL.}}{PostgreSQL.} Las condiciones por
lote y sus selecciones deben revisarse antes de confirmar.

\begin{figure}[htbp]
\centering
\includegraphics[width=.95\linewidth,height=.70\textheight,keepaspectratio]{\detokenize{@@BOVI_ROOT@@/docs/testing/evidence/ui/bovicierre-ui-flow-1440.png}}
\caption{Recorrido de operación con datos DEMO.}
\end{figure}
\clearpage

\section{Corrección de pesajes}\label{correcciuxf3n-de-pesajes}

La corrección conserva el conjunto exacto y exige motivo. Los pesos
individuales desconocidos permanecen ausentes.

\begin{figure}[htbp]
\centering
\includegraphics[width=.95\linewidth,height=.70\textheight,keepaspectratio]{\detokenize{@@BOVI_ROOT@@/docs/testing/evidence/ui/bovicierre-ui-pesajes-correccion-1440.png}}
\caption{Corrección de pesajes con datos DEMO.}
\end{figure}
\clearpage

\section{Cartera}\label{cartera}

La consulta informa corte y fuente. Saldo, vencimiento y disponible de
pago tienen significados distintos.

\begin{figure}[htbp]
\centering
\includegraphics[width=.95\linewidth,height=.70\textheight,keepaspectratio]{\detokenize{@@BOVI_ROOT@@/docs/testing/evidence/ui/bovicierre-ui-reporte-cartera.png}}
\caption{Cartera con datos DEMO.}
\end{figure}
\clearpage

\section{Margen}\label{margen}

El resultado contiene venta, costo, gastos y pérdida con
\texorpdfstring{\softtext{procedencia.}}{procedencia.} Un cierre
histórico se distingue del resultado actual.

\begin{figure}[htbp]
\centering
\includegraphics[width=.95\linewidth,height=.70\textheight,keepaspectratio]{\detokenize{@@BOVI_ROOT@@/docs/testing/evidence/ui/bovicierre-ui-reporte-margen.png}}
\caption{Margen con datos DEMO.}
\end{figure}
\clearpage

\section{Periodos y cierres}\label{periodos-y-cierres}

El reconocimiento utiliza un periodo abierto
\texorpdfstring{\softtext{explícito.}}{explícito.} Las versiones de
cierre conservan el resultado
\texorpdfstring{\softtext{histórico.}}{histórico.}

\begin{figure}[htbp]
\centering
\includegraphics[width=.95\linewidth,height=.70\textheight,keepaspectratio]{\detokenize{@@BOVI_ROOT@@/docs/testing/evidence/ui/bovicierre-ui-periodos.png}}
\caption{Periodos y cierres con datos DEMO.}
\end{figure}
\clearpage

\section{Operación a390 px con
teclado}\label{operaciuxf3n-a390-px-con-teclado}

El flujo de compra, venta y aplicación se comprobó con teclado. La
captura no certifica lector de pantalla ni dispositivo físico.

\begin{figure}[htbp]
\centering
\includegraphics[width=.95\linewidth,height=.70\textheight,keepaspectratio]{\detokenize{@@BOVI_ROOT@@/docs/testing/evidence/ui/bovicierre-ui-keyboard390.png}}
\caption{Operación a390px con teclado con datos DEMO.}
\end{figure}
\clearpage

\chapter{Láminas de arquitectura de
datos}\label{luxe1minas-de-arquitectura-de-datos}

Las láminas vectoriales siguientes reúnen el MER Chen y el modelo
relacional completo, seguidos de sus vistas por dominio. El dibujo
completo permite comprender el conjunto; las vistas permiten leer las
relaciones y atributos. Puede ampliarse el PDF sin convertir el texto
del diagrama en una captura borrosa. La matriz y el diccionario anexos
proporcionan la referencia textual. Revisión humana de cardinalidades y
aceptación del profesor siguen
\texorpdfstring{\softtext{pendientes.}}{pendientes.}

\includepdf[pages=-,landscape=true,scale=0.89,pagecommand={\thispagestyle{fancy}}]{\detokenize{@@BOVI_BUILD@@/assets/diagram-01.pdf}}

\includepdf[pages=-,landscape=true,scale=0.89,pagecommand={\thispagestyle{fancy}}]{\detokenize{@@BOVI_BUILD@@/assets/diagram-02.pdf}}

\includepdf[pages=-,landscape=true,scale=0.89,pagecommand={\thispagestyle{fancy}}]{\detokenize{@@BOVI_BUILD@@/assets/diagram-03.pdf}}

\includepdf[pages=-,landscape=true,scale=0.89,pagecommand={\thispagestyle{fancy}}]{\detokenize{@@BOVI_BUILD@@/assets/diagram-04.pdf}}

\includepdf[pages=-,landscape=true,scale=0.89,pagecommand={\thispagestyle{fancy}}]{\detokenize{@@BOVI_BUILD@@/assets/diagram-05.pdf}}

\includepdf[pages=-,landscape=true,scale=0.89,pagecommand={\thispagestyle{fancy}}]{\detokenize{@@BOVI_BUILD@@/assets/diagram-06.pdf}}

\includepdf[pages=-,landscape=true,scale=0.89,pagecommand={\thispagestyle{fancy}}]{\detokenize{@@BOVI_BUILD@@/assets/diagram-07.pdf}}

\includepdf[pages=-,landscape=true,scale=0.89,pagecommand={\thispagestyle{fancy}}]{\detokenize{@@BOVI_BUILD@@/assets/diagram-08.pdf}}

\includepdf[pages=-,landscape=true,scale=0.89,pagecommand={\thispagestyle{fancy}}]{\detokenize{@@BOVI_BUILD@@/assets/diagram-09.pdf}}

\includepdf[pages=-,landscape=true,scale=0.89,pagecommand={\thispagestyle{fancy}}]{\detokenize{@@BOVI_BUILD@@/assets/diagram-10.pdf}}

\includepdf[pages=-,landscape=true,scale=0.89,pagecommand={\thispagestyle{fancy}}]{\detokenize{@@BOVI_BUILD@@/assets/diagram-11.pdf}}

\includepdf[pages=-,landscape=true,scale=0.89,pagecommand={\thispagestyle{fancy}}]{\detokenize{@@BOVI_BUILD@@/assets/diagram-12.pdf}}

\chapter{Actualización de la campaña UI/UX y QA del 12 de octubre de
2026}\label{actualizacion-campana-uiux-qa}

Este capítulo documenta, sin reescribir los anteriores, todo lo que
cambió en el software durante la campaña del 12 de octubre de 2026,
ejecutada en la rama local \texttt{feature/uiux-qa-claude} a partir de
la revisión \texttt{38ccdc3}. Los capítulos previos describen el corte
\texttt{92d3d41} con migraciones V1--V13; aquí se explica qué se añadió
encima (migraciones V14--V18, endurecimiento de seguridad, rediseño
visual, contrato y pruebas), por qué se decidió así y cómo se
comprueba. Cuando este capítulo y uno anterior difieran, prevalece este
capítulo para el estado actual.

Alcance honesto: todo lo descrito es verificación técnica local con
evidencia reproducible. No es aceptación humana, despliegue ni
certificación de seguridad. La revisión de Aldo y Jared sigue
pendiente. Nada se publicó en un remoto.

\section{Resumen de cambios por capa}\label{campana-resumen}

\begin{longtable}[]{@{}>{\raggedright\arraybackslash}p{0.17\linewidth}>{\raggedright\arraybackslash}p{0.48\linewidth}>{\raggedright\arraybackslash}p{0.29\linewidth}@{}}
\toprule
Capa & Cambio & Identificadores \\
\midrule
\endhead
Interfaz & Identidad visual bovina con tokens, fuente Inter local,
marca e iconos SVG propios, dos fotografías de dominio público,
insignias de estado con color y forma, saldo destacado, UUID en
monoespaciado y tablas legibles en móvil. & S1, UXD-01..11 \\
Backend & Autorización de administración por sede, límite de intentos
de inicio de sesión, tope de sesiones anónimas, UNIQUE
\texttt{*\_key} traducido a 409. & SF-01, SF-02, SF-03, DIV-01/SF-20 \\
Borde HTTP & Cabeceras CSP, \texttt{X-Frame-Options},
\texttt{Permissions-Policy} y límites de tamaño de cuerpo en Caddy. &
SF-04, SF-05 \\
Base de datos & V14 estados terminales inmutables y
\texttt{search\_path} con \texttt{pg\_temp}; V15 segundo reverso con
BV022; V16/V17 liquidación cuadrada y un solo periodo abierto; V18 vencimiento máximo de reserva. &
DIV-04, SF-08, DIV-08, DIV-05, DIV-14 \\
Contrato & OpenAPI de movimientos y anulación con los campos que la
base exige; cuatro rutas reales documentadas y alias obsoleto marcado; campos de \texttt{Pago} que nunca se devolvían retirados; tipos TypeScript regenerados; comprobación automática de cobertura de rutas. & DIV-11, DIV-12, DIV-13 \\
Documentación & Catálogo de códigos BV regenerado desde el código;
fórmula canónica de margen; RN-034 y ADR-002 corregidos; nombres reales de triggers, enrutamiento y comandos; ENT-42; retención de idempotencia; alcance de BV402. & DIV-06,
DIV-07, DIV-12b, DIV-15, DIV-16, DIV-17 \\
Pruebas & Puerta local en serie, inmutabilidad ampliada, seguridad
dinámica, carga k6, lint de tipos monetarios, pesaje BRUTO\_TARA,
BV021, BV010 y posición activa única. & DIV-02, DIV-03, DIV-09,
DIV-10 \\
Recuperación & \texttt{restore-check} reaplica el ACL de base que
\texttt{pg\_dump} no transporta. & SF-08 \\
\bottomrule
\end{longtable}

\section{Rediseño visual (S1)}\label{campana-rediseno}

El rediseño se aplicó sobre la aplicación real, sin maquetas paralelas
ni cambios de flujo. Sus piezas son:

\begin{itemize}
\item \textbf{Tokens} en \texttt{frontend/src/styles.css}: escalas
\texttt{--green-*} y \texttt{--earth-*}, espaciado \texttt{--s*},
radios \texttt{--r-*} y tipografía \texttt{--fs-*}. Se conservaron
todas las clases existentes para no romper componentes.
\item \textbf{Fuente Inter Variable} servida localmente
(\texttt{@fontsource-variable/inter}, licencia OFL). No hay
peticiones a CDN, lo que también permite la CSP estricta.
\item \textbf{Marca e iconos} en \texttt{frontend/src/brand.tsx}
(\texttt{BrandMark}, \texttt{Brand}, \texttt{NavIcon},
\texttt{initials}), dibujados como SVG originales.
\item \textbf{Fotografías} de Wikimedia Commons en dominio público, sólo
en acceso y selección de sede. Licencia, autor y SHA-256 constan en
\texttt{docs/ui/asset-register.md}.
\item \textbf{Insignias de estado} en \texttt{frontend/src/tone.ts}:
\texttt{toneOf} clasifica cada valor como positivo, en curso, atención,
crítico o neutro, y \texttt{statusColumns} decide qué columnas se
pintan. Cada tono combina color \emph{y} forma, de modo que no depende
sólo del color (WCAG 1.4.1). Estado comercial, estado de entrega y
estado de pago aparecen como insignias separadas.
\item \textbf{Contraste}: \texttt{tools/ui/contrast.py} comprueba 20
pares de tokens. El borde de controles se subió de 2.07:1 a 3.69:1
(\texttt{\#8A8574}).
\item \textbf{Móvil}: tablas con ancho mínimo de 680 px dentro de un
contenedor desplazable y \texttt{overflow-wrap:break-word}; objetivos
de navegación de 44 px. Un primer intento partía las palabras de las
tablas y se corrigió.
\end{itemize}

\begin{figure}[htbp]
\centering
\includegraphics[width=.95\linewidth,height=.42\textheight,keepaspectratio]{\detokenize{@@BOVI_ROOT@@/docs/testing/evidence/ui-redesign/20261012T0214/after/1440-login.jpg}}
\caption{Acceso rediseñado a 1440 px, con la fotografía de dominio
público y su crédito.}
\end{figure}

\begin{figure}[htbp]
\centering
\includegraphics[width=.95\linewidth,height=.42\textheight,keepaspectratio]{\detokenize{@@BOVI_ROOT@@/docs/testing/evidence/ui-redesign/20261012T0214/after/1440-operaciones.jpg}}
\caption{Bandeja de operaciones con insignias de estado comercial, de
entrega y de pago separadas.}
\end{figure}

\begin{figure}[htbp]
\centering
\includegraphics[width=.95\linewidth,height=.45\textheight,keepaspectratio]{\detokenize{@@BOVI_ROOT@@/docs/testing/evidence/ui-redesign/20261012T0214/compare/1440-liquidacion-detalle.jpg}}
\caption{Detalle de liquidación antes y después del rediseño: el saldo
pasa a ser el dato dominante.}
\end{figure}

\begin{figure}[htbp]
\centering
\includegraphics[width=.40\linewidth,height=.55\textheight,keepaspectratio]{\detokenize{@@BOVI_ROOT@@/docs/testing/evidence/ui-redesign/20261012T0214/after/390-operaciones.jpg}}
\caption{La misma bandeja a 390 px: la tabla se desplaza dentro de su
contenedor y la página no tiene desbordamiento horizontal.}
\end{figure}
\clearpage

Las 45 capturas posteriores (15 vistas a 390, 768 y 1440 px), las previas y las comparaciones están en
\texttt{docs/testing/evidence/ui-redesign/20261012T0214/}. La prueba
\texttt{frontend/tests\_e2e/visual.spec.ts} recorre esas vistas sin
mutar datos y falla si \texttt{scrollWidth} supera \texttt{clientWidth}.
No se ejecutaron lector de pantalla, axe ni teléfono físico.

\section{Seguridad endurecida}\label{campana-seguridad}

La revisión combinó un modelo de amenazas
(\texttt{docs/security/threat-model.md}), revisión estática y 63 casos
dinámicos (\softcode{tests/security/sec\_dynamic.py}) contra la instancia
local con datos DEMO. La primera ejecución dio 48/56: cuatro fallos eran
defectos del propio test y cuatro eran hallazgos reales.

\begin{longtable}[]{@{}>{\raggedright\arraybackslash}p{0.10\linewidth}>{\raggedright\arraybackslash}p{0.40\linewidth}>{\raggedright\arraybackslash}p{0.44\linewidth}@{}}
\toprule
ID & Problema & Corrección \\
\midrule
\endhead
SF-01 & Un ADMIN de una sede podía bloquear usuarios de otra. &
\texttt{AccessService.member} y \texttt{UserController}: 404 si el
usuario no pertenece a la sede activa; 400 ante autobloqueo. \\
SF-02 & El inicio de sesión no tenía límite de intentos. &
\texttt{LoginThrottle}: 10 fallos por login y 50 por origen en 15
minutos con ventana deslizante; responde 429 \texttt{AUTH429} con
\texttt{Retry-After}. El éxito reinicia el contador del login. \\
SF-03 & Pedir el token CSRF sin sesión creaba una sesión de 30 minutos,
lo que permitía llenar \texttt{infra.spring\_session}. & Las sesiones
anónimas duran 300 s y hay un tope de 600 nuevas por origen; la sesión
autenticada vuelve a 1800 s tras rotar su identificador. \\
SF-04 & Faltaban CSP y protección contra incrustación. & Caddy envía
\texttt{default-src 'self'}, \texttt{frame-ancestors 'none'},
\texttt{X-Frame-Options DENY} y \texttt{Permissions-Policy}. \\
SF-05 & Cuerpos JSON sin límite. & 1 MB para JSON y 22 MB para subida
de documentos. \\
SF-08 & \texttt{search\_path=pg\_catalog} sin \texttt{pg\_temp} y TEMP
para PUBLIC. & V14 (sección siguiente). \\
SF-17 & Las expresiones del mod de campaña podían marcar como
verificado un comando que sólo contenía la palabra. & Expresiones
ancladas al inicio del comando. \\
SF-20 & UNIQUE autogenerada \texttt{*\_key} respondía 500. &
\texttt{ApiErrors}: \texttt{uk\_*} o \texttt{*\_key} → 409 BV401. \\
SF-21 & \texttt{vite} 8.0.3 con avisos de seguridad en el servidor de
desarrollo. & \texttt{vite} 8.0.16 exacto. \\
\bottomrule
\end{longtable}

Riesgos que siguen abiertos: SF-09 (media) requiere decisión humana,
porque la credencial \texttt{bovi\_app\_*} lee todo el esquema
\texttt{bovi} y \texttt{p\_actor} lo envía el llamador; SF-19 se aceptó
porque purgar intentos pendientes al cerrar sesión podría duplicar
comandos económicos; SF-06, SF-07, SF-10 y SF-11 son bajas estáticas.
No se ejecutaron ZAP, Semgrep, Gitleaks, Trivy ni un escáner de Maven.

\section{Cambios en la base de datos: V14 a V17}\label{campana-migraciones}

Ninguna migración V1--V13 se reescribió. Cada cambio se probó primero en
rojo (la prueba nueva fallaba contra el esquema anterior), después se
aplicó con Flyway en los nodos A y B mediante
\texttt{db/bootstrap/migrate.sh}, y finalmente se repitió la puerta
local y el ensayo de restauración.

\subsection{V14: estados terminales inmutables y search\_path seguro}

\textbf{Qué hace.} Crea \texttt{bovi.trg\_estado\_terminal} y el trigger
\texttt{estado\_terminal\_inmutable} BEFORE UPDATE OR DELETE en
\texttt{operacion}, \texttt{gasto\_operacion}, \texttt{pago} y
\texttt{posicion\_inventario}. Además fija
\texttt{search\_path = pg\_catalog, pg\_temp} en las 32 funciones de
\texttt{bovi} e \texttt{infra} y revoca TEMP a PUBLIC sobre la base.

\textbf{Por qué.} RN-023 exige que los hechos confirmados no se
reescriban. Antes de V14, esos cuatro agregados sólo estaban protegidos
por el código de \texttt{cmd\_ejecutar} y por la falta de DML del rol de
aplicación. La fase roja demostró la brecha: los 13 intentos sobre
estados terminales pasaban o sólo los frenaba una FK, y borrar un pago
VALIDADO se permitía porque ninguna FK lo referenciaba. Para SF-08, si
\texttt{pg\_temp} no aparece explícitamente en el
\texttt{search\_path}, PostgreSQL lo busca \emph{primero}, de modo que
una tabla temporal del llamador puede sombrear un catálogo no
calificado.

\begin{longtable}[]{@{}>{\raggedright\arraybackslash}p{0.22\linewidth}>{\raggedright\arraybackslash}p{0.30\linewidth}>{\raggedright\arraybackslash}p{0.42\linewidth}@{}}
\toprule
Tabla & Estados terminales & Únicos cambios permitidos \\
\midrule
\endhead
\texttt{operacion} & CONFIRMADA, ANULADA, CANCELADA &
CONFIRMADA→ANULADA (\texttt{anular\_venta}, sólo sin entregas);
\texttt{estado\_entrega} mientras está CONFIRMADA; avance de
\texttt{version\_lock}. Todo lo demás exige reverso o ajuste. \\
\texttt{gasto\_operacion} & APROBADO & Ninguno. \\
\texttt{pago} & VALIDADO, REVERTIDO & Ninguno; el reembolso es otra
fila. \\
\softcode{posicion\_inventario} & CERRADA\_VENTA, BAJA & Ninguno; una
recompra abre un ciclo nuevo. \\
\bottomrule
\end{longtable}

\textbf{Hallazgo de recuperación.} El primer \texttt{restore-check}
posterior a V14 concilió datos y hashes, pero la base restaurada
recuperaba TEMP para PUBLIC: el ACL de la base no viaja en
\texttt{pg\_dump} sin \texttt{--create}. \texttt{scripts/recovery.py}
ahora reaplica \texttt{REVOKE TEMPORARY} tras \texttt{pg\_restore} y lo
afirma con \texttt{has\_database\_privilege}.

\subsection{V15: un segundo reverso responde BV022}

\textbf{Qué hace.} \texttt{bovi.trg\_reverso\_unico} (BEFORE INSERT en
\texttt{aplicacion\_pago} cuando \texttt{reverso\_de\_id} no es nulo)
lanza BV022 (RN-033) si la aplicación ya fue revertida.

\textbf{Por qué.} RN-033 documenta BV022 como «original no revertido»,
pero un segundo reverso chocaba con
\texttt{uk\_aplicacion\_pago\_reverso\_de\_id} y salía como 409 BV401
genérico. \texttt{cmd\_ejecutar} bloquea el pago con
\texttt{FOR UPDATE} antes de insertar, así que dos reversos simultáneos
se serializan y el segundo ve el primero ya confirmado. La UNIQUE queda
como respaldo.

\subsection{V16 y V17: liquidación cuadrada y un solo periodo abierto}

\textbf{Qué hace.} V16 crea el constraint trigger
\texttt{liquidacion\_cuadrada}, DEFERRABLE INITIALLY DEFERRED, en
\texttt{liquidacion} y \texttt{concepto\_liquidacion}: al confirmar la
transacción comprueba
\[
\texttt{liquidacion.total} \;=\; \sum \texttt{concepto\_liquidacion.importe}
\]
y responde BV015 (RN-023) si no se cumple. También crea el índice único
parcial \texttt{uk\_periodo\_un\_abierto} sobre
\texttt{periodo\_cierre (empresa\_id, sucursal\_id)} donde
\texttt{estado = 'ABIERTO'}.

\textbf{Por qué.} Antes la igualdad sólo se detectaba después, en
\texttt{fn\_validar\_conciliacion\_propietario} (V3), y la unicidad del
periodo abierto dependía de un advisory lock dentro del comando. Ahora
la base impide ambos estados inválidos aunque alguien escriba fuera del
comando.

\textbf{V17.} V16 tenía un defecto: en PL/pgSQL, \texttt{NEW.id} se
resuelve aunque la rama del \texttt{CASE} no se ejecute, y
\texttt{concepto\_liquidacion} no tiene columna \texttt{id}, así que
cualquier concepto nuevo habría fallado al confirmar. La fase verde lo
detectó antes de cualquier commit. Como V16 ya estaba aplicada en A y B,
se corrigió hacia delante con V17 (\texttt{CREATE OR REPLACE}, leyendo
la columna mediante \texttt{to\_jsonb(NEW)}) en lugar de editar una
migración aplicada y reescribir el historial de Flyway.

\subsection{V18: vencimiento máximo de una reserva}

\textbf{Qué hace.} Añade el CHECK
\texttt{reserva\_venta\_vence\_max\_check}:
\texttt{vence\_en <= creada\_en + 30 días}. Una violación sale como 23514
y la API la traduce a 422 BV015.

\textbf{Por qué.} RN-019 justifica el vencimiento con «no retener
inventario indefinidamente», pero \texttt{vence\_en} lo escribe el
usuario en la pantalla de reserva sin ningún tope. Las reservas reales
del sistema duran como máximo unas 10 horas, así que 30 días es un tope
holgado que no bloquea el uso actual. Es un parámetro de negocio que Aldo
y Jared pueden ajustar (DEC-24). RN-019 se reescribió para describir la
expiración real: no existe un comando \texttt{ExpirarReserva}; reservar,
liberar y confirmar marcan EXPIRADA las reservas vencidas antes de
operar.

\section{Contrato HTTP y backend}\label{campana-contrato}

\begin{itemize}
\item \textbf{Movimientos de inventario}
(\texttt{registrarMovimiento}): el esquema declara ahora
\texttt{operacion\_id}, \texttt{periodo\_id} y
\texttt{fecha\_reconocimiento}, y con \texttt{if/then} de OpenAPI 3.1
exige \texttt{motivo\_baja}, \texttt{documento\_id}, \texttt{periodo\_id}
y \texttt{fecha\_reconocimiento} cuando \texttt{tipo = BAJA}, igual que
la función SQL (BV025 si faltan).
\item \textbf{Anulación} (\texttt{anularOperacion}): el cuerpo combina
\texttt{Reconocimiento} con \texttt{motivo} obligatorio, porque el
ajuste de anulación escribe \texttt{ajuste\_liquidacion.motivo NOT NULL}.
\item \texttt{frontend/src/generated/api-schema.ts} se regeneró con
\texttt{npm --prefix frontend run generate:api}; Redocly valida el
documento.
\item \textbf{Cobertura de rutas} (DIV-11): se documentaron
\texttt{GET /operaciones/\{id\}/pesajes}, \texttt{GET /reportes/inventario},
\texttt{GET /auditoria} y \texttt{POST /operaciones/\{id\}/revisiones}, y
el alias \texttt{POST /periodos/\{id\}/cerrar} quedó marcado como
\texttt{deprecated}. \texttt{tools/api/route\_coverage.py}, dentro de
\texttt{make verify-docs}, falla si una ruta de un controlador Java no
está en el contrato; hoy cubre las 71 rutas.
\item \textbf{Pago} (DIV-13): se retiraron \texttt{aplicado} y
\texttt{reembolsado}, que ninguna consulta devolvía.
\item \textbf{RN-034} se cumple con una prueba real:
\texttt{MoneyTypesLintTest} falla si aparece \texttt{double},
\texttt{float}, \texttt{Double} o \texttt{Float} en
\texttt{backend/src/main/java}, o \texttt{parseFloat}/\texttt{toFixed}
en \texttt{frontend/src}, y comprueba que recorrió archivos reales para
no pasar en vacío.
\end{itemize}

\section{Decisiones semánticas razonadas}\label{campana-decisiones}

Las decisiones siguientes resolvieron divergencias entre reglas,
modelo, SQL, API y documentación. Se registran como DEC-12 a DEC-25 en
\texttt{docs/research/evidence-register.md}. Cada una indica la
alternativa descartada y la razón.

\begin{description}
\item[DEC-12.] \textbf{Margen canónico (DIV-06).} La fórmula oficial es
\texttt{fn\_reporte} con el recurso \texttt{reportes\_margen}, la que consume el
endpoint vía \texttt{consultar}. Se descartó declarar canónica la vista
\texttt{v\_margen\_operacion} porque fija \texttt{estado\_margen} a
PROVISIONAL, ignora el periodo de reconocimiento y las pérdidas, y sólo
la usaban las pruebas. Se descartó eliminar la vista porque romper las
cifras GOLD no aporta valor; en su lugar \texttt{gold.py} afirma que
vista y reporte dan las mismas cifras en los escenarios GOLD.
\item[DEC-13.] \textbf{Transiciones permitidas en estados terminales (DIV-04).}
Se derivaron literalmente de las sentencias \texttt{UPDATE} vigentes en
\texttt{cmd\_ejecutar} V13, no de la documentación, para no romper un
flujo real. CANCELADA es terminal porque cancelar sólo procede desde
BORRADOR o ACORDADA. REVERTIDO es terminal porque la reversión crea
filas nuevas. La anulación conserva \texttt{estado\_entrega} porque sólo
procede sin entregas.
\item[DEC-14.] \textbf{Inmutabilidad en la base y no sólo en el comando.} Se
eligieron triggers aunque el rol de aplicación ya no tenga DML, porque
la defensa en profundidad debe resistir un error futuro de grants, una
función SECURITY DEFINER nueva o una intervención manual. El costo es
una comparación de \texttt{jsonb} por fila modificada.
\item[DEC-15.] \textbf{BV022 por trigger y no reescribiendo
\texttt{cmd\_ejecutar} (DIV-08).} Reescribir la función completa para
añadir una línea habría copiado unas 350 líneas y elevado el riesgo de
regresión. El trigger es localizado y el bloqueo del pago ya serializa
la carrera.
\item[DEC-16.] \textbf{Liquidación cuadrada diferida con BV015 y RN-023
(DIV-05).} Diferida porque el comando inserta la liquidación antes que
sus conceptos. BV015 porque es una violación de forma de los datos y no
un conflicto de concurrencia. RN-023 porque protege la integridad de la
liquidación confirmada. Se reutilizó la expresión exacta de la
conciliación V3 para que detección y prevención no diverjan.
\item[DEC-17.] \textbf{Corregir hacia delante (V17).} Una migración aplicada no
se edita: cambiar su checksum rompe \texttt{flyway validate} y oculta
historia. Corregir con V17 deja el defecto y su arreglo trazables.
\item[DEC-18.] \textbf{Catálogo de códigos (DIV-07).} El catálogo se regeneró a
partir de lo que emiten \texttt{pg\_proc} y \texttt{ApiException}. Los
códigos documentados y nunca emitidos (BV012, BV014, BV017, BV019)
quedan reservados y no se reasignan, para que un cliente antiguo no
interprete un significado distinto. La convención se corrigió:
BV404→404, BV412→412 y BV428→428 son excepciones a BV4xx→409.
\item[DEC-19.] \textbf{\texttt{motivo} obligatorio en la anulación (DIV-12).}
La base sólo lo necesita si \texttt{total\_ajustado > 0}, pero un
contrato condicional por saldo sería opaco para el cliente. Como no hay
pantalla que anule y las pruebas siempre lo envían, exigirlo siempre es
más simple y seguro.
\item[DEC-20.] \textbf{ACL de base reaplicado al restaurar (SF-08).} Se prefirió
reaplicar y afirmar el \texttt{REVOKE} en \texttt{restore-check} antes
que volcar con \texttt{--create}, que intentaría recrear la base que el
contenedor ya creó.
\item[DEC-21.] \textbf{Grafo semántico sin capturas de evidencia.} Las 151
capturas de \texttt{docs/testing/evidence} no se extrajeron
semánticamente: costarían un agente por imagen y repiten lo que ya
describen los documentos. Quedan sin sellar para una pasada posterior.
\item[DEC-22.] \textbf{Rutas sin contrato (DIV-11).} Documentar las
cuatro rutas que la aplicación usa y marcar el alias \texttt{/cerrar}
como obsoleto en lugar de retirarlo: ningún cliente lo usa, pero
retirarlo rompería la compatibilidad sin beneficio.
\item[DEC-23.] \textbf{Campos fantasma de Pago (DIV-13).} Quitarlos del
contrato. Calcularlos sería una función nueva que nadie pidió.
\item[DEC-24.] \textbf{Tope de vencimiento de reserva (DIV-14).} 30 días,
explicado en la sección de V18.
\item[DEC-25.] \textbf{BV402 compartido (DIV-17).} Sólo documentar. El
lease del outbox nunca sale por HTTP y su única prueba es una caída real
de nodos fuera de la puerta; renombrar el código dejaría un cambio sin
verificar.
\end{description}

\section{Pruebas nuevas y puerta local}\label{campana-pruebas}

\texttt{make verify-release-local} (\texttt{scripts/verify\_release\_local.py})
ejecuta en serie docs, frontend, backend con HTTP real y Testcontainers,
harness SQL, inmutabilidad, E2E y seguridad dinámica, y opcionalmente
carga k6. Falla si una etapa obligatoria falla, se omite, reporta casos
omitidos o declara cero casos ejecutados. Su primera ejecución fue
rechazada por la propia puerta (la etapa de inmutabilidad no declaraba
casos), y ese rechazo se conserva como evidencia.

\begin{longtable}[]{@{}>{\raggedright\arraybackslash}p{0.27\linewidth}>{\raggedright\arraybackslash}p{0.67\linewidth}@{}}
\toprule
Prueba & Qué demuestra \\
\midrule
\endhead
\softcode{db/tests/immutability.py} & UPDATE y DELETE bloqueados en las 13
tablas con \texttt{hecho\_inmutable}; 13 bloqueos y 5 transiciones
legítimas de DIV-04; liquidación descuadrada y segundo periodo abierto
(DIV-05); sede operadora y tipo de operación inmutables (BV010) y una
sola posición activa por animal (DIV-02); reserva con vencimiento mayor a 30 días (DIV-14). Todo en BEGIN…ROLLBACK y con
conteos intactos. \\
\softcode{db/tests/gold.py} & Además de las cifras GOLD: segundo reverso
→ BV022 (DIV-08); cobro aplicado a una obligación por pagar y cobro de
otra contraparte → BV021 sin efectos (DIV-09); vista de margen igual al
reporte canónico (DIV-06). \\
\softcode{db/tests/weighing\_reports.py} & Pesaje BRUTO\_TARA: el neto se
calcula en el servidor (1250.500 − 950.250 = 300.250) ignorando un neto
falso enviado; tara mayor que bruto → 23514 (API BV015) (DIV-10). \\
\softcode{db/tests/security.sql} & Todas las funciones con
\texttt{pg\_catalog, pg\_temp}; PUBLIC y roles de ejecución sin TEMP;
\texttt{CREATE TEMP TABLE pg\_tables} como \texttt{bovi\_auditor} →
42501. \\
\softcode{tests/security/sec\_dynamic.py} & 63 casos: inyección SQL,
acceso horizontal, escalada, CSRF, sesión, XSS almacenado, CSV, QR,
abuso de login, cabeceras y 409 de UNIQUE. \\
\softcode{tests/performance/load.js} & k6 de 1 a 20 sesiones propias
(70\% lectura, 20\% escritura, 10\% descarga). Con V16/V17: 8,214
peticiones sin fallos, p95 22.9 ms y 125 operaciones confirmadas, todas
con liquidación cuadrada. No satura el sistema: no prueba capacidad
sostenida ni volumen de 10,000 posiciones. \\
\softcode{MoneyTypesLintTest} & Ausencia de coma flotante en importes
(RN-034). \\
\softcode{frontend/tests/tone.test.ts}, \softcode{visual.spec.ts} &
Clasificación de tonos y ausencia de desbordamiento a tres anchos. \\
\bottomrule
\end{longtable}

\textbf{Pruebas frágiles corregidas.} \texttt{keyboard.spec.ts} dependía
del idioma del sistema (es\_ES invierte día y mes en
\texttt{datetime-local}) y fijaba octubre de 2026; ahora fija
\texttt{en-US} y calcula la fecha esperada. \texttt{ApiPostgresTest}
buscaba la posición nueva sólo en la primera página de 100; ahora
filtra con \texttt{q}.

\textbf{Recuperación.} Se tomó un respaldo antes de V14 y se repitió
\texttt{make backup} y \texttt{make restore-check} después de V14, V15,
V17 y V18. El último ensayo concilió firma, hashes y sesiones en 31.1 s,
con 7 triggers de inmutabilidad, \texttt{uk\_periodo\_un\_abierto},
\texttt{reserva\_venta\_vence\_max\_check} y TEMP revocado en la base
restaurada. No acredita PITR ni pérdida del host.

\section{Herramientas de campaña}\label{campana-herramientas}

\begin{itemize}
\item \texttt{tools/claude/checkpoint.py} y
\texttt{tools/claude/resume-check.py} guardan el estado de la campaña
fuera de Git y verifican, antes de continuar, el remoto, el marcador, la
ausencia de \texttt{STOP}, el estado no terminal y el lock.
\texttt{docs/continuity/NEXT\_SESSION.md} indica la siguiente tarea; un
archivo Markdown no ejecuta nada por sí mismo.
\item El plugin \texttt{tools/claude/bovicierre-workbench} añade el
comando \texttt{/bovi-status}, avisos de cuota y de checkpoint, y marca
áreas pendientes de verificación tras editar. Su carga se verificó en
sesión.
\item El grafo de conocimiento (\texttt{graphify-out/}) se actualizó con
extracción AST del código y semántica de los documentos de la campaña:
2,871 nodos, 4,758 aristas y 298 comunidades, sin aristas colgantes.
\end{itemize}

\section{Pendientes después de esta actualización}\label{campana-pendientes}

\begin{itemize}
\item Decisión humana sobre SF-09: protección por fila o contexto de
actor en la base, o un ADR que acepte el riesgo por escrito.
\item Las 17 divergencias de la matriz quedaron corregidas, cubiertas o resueltas; las que se resolvieron sólo en documentación (DIV-16, DIV-17) lo indican con su razón.
\item Escáneres SCA y DAST, carga sostenida de 15 a 20 minutos y dataset
de 10,000 posiciones.
\item Revisión de los diagramas MER y relacional contra la matriz de
coherencia; las migraciones V14--V18 no añaden tablas ni columnas, sólo
triggers, un índice, un CHECK y privilegios.
\item Lector de pantalla, axe y teléfono físico para el rediseño.
\item Revisión humana de Aldo y Jared de todo lo anterior.
\end{itemize}

\chapter{Fuentes y trazabilidad del
manual}\label{fuentes-y-trazabilidad-del-manual}

Revisión de código
\texorpdfstring{\softcode{92d3d4103a174eb11f4ccf21eff2deed8cc77682}}{92d3d4103a174eb11f4ccf21eff2deed8cc77682}.
Los hashes identifican los archivos públicos utilizados, no incluyen
directorios privados, contraseñas, cookies ni dumps. La documentación de
capítulos aporta referencias adicionales de evidencia e
\texorpdfstring{\softtext{infraestructura.}}{infraestructura.}

{\def\LTcaptype{none} % do not increment counter
\begin{longtable}[]{@{}
  >{\raggedright\arraybackslash}p{(\linewidth - 2\tabcolsep) * \real{0.5000}}
  >{\raggedright\arraybackslash}p{(\linewidth - 2\tabcolsep) * \real{0.5000}}@{}}
\toprule\noalign{}
\begin{minipage}[b]{\linewidth}\raggedright
Archivo
\end{minipage} & \begin{minipage}[b]{\linewidth}\raggedright
SHA-256
\end{minipage} \\
\midrule\noalign{}
\endhead
\bottomrule\noalign{}
\endlastfoot
\texorpdfstring{\softcode{backend/src/main/java/mx/bovicierre/BoviCierreApplication.java}}{backend/src/main/java/mx/bovicierre/BoviCierreApplication.java}
&
\texorpdfstring{\softcode{d980c168ba1cdfcaf12a34ff85e5c3ac9fbac7d4dbcfb5a9ffb4bd3b334d2fd8}}{d980c168ba1cdfcaf12a34ff85e5c3ac9fbac7d4dbcfb5a9ffb4bd3b334d2fd8} \\
\texorpdfstring{\softcode{backend/src/main/java/mx/bovicierre/acceso/AccessService.java}}{backend/src/main/java/mx/bovicierre/acceso/AccessService.java}
&
\texorpdfstring{\softcode{0304ac667578d051ff55484e9972c438d99d1ce82ec77d0fb6a419adf0f021a8}}{0304ac667578d051ff55484e9972c438d99d1ce82ec77d0fb6a419adf0f021a8} \\
\texorpdfstring{\softcode{backend/src/main/java/mx/bovicierre/acceso/AvailabilityFilter.java}}{backend/src/main/java/mx/bovicierre/acceso/AvailabilityFilter.java}
&
\texorpdfstring{\softcode{5806a5edd039a20d529c07331985ccdd8b1cf14a4d911684f5db25262e08718c}}{5806a5edd039a20d529c07331985ccdd8b1cf14a4d911684f5db25262e08718c} \\
\texorpdfstring{\softcode{backend/src/main/java/mx/bovicierre/acceso/SecurityConfiguration.java}}{backend/src/main/java/mx/bovicierre/acceso/SecurityConfiguration.java}
&
\texorpdfstring{\softcode{10c96e1cae65befebfeacb9f3bc539464cfaec4fc93e4af6503c3dbee8b5ac22}}{10c96e1cae65befebfeacb9f3bc539464cfaec4fc93e4af6503c3dbee8b5ac22} \\
\texorpdfstring{\softcode{backend/src/main/java/mx/bovicierre/acceso/SessionController.java}}{backend/src/main/java/mx/bovicierre/acceso/SessionController.java}
&
\texorpdfstring{\softcode{5754cfe69035116b05a3dc46abbd0ddfa8a8c5b69867c19e0342bb09fb8c64f7}}{5754cfe69035116b05a3dc46abbd0ddfa8a8c5b69867c19e0342bb09fb8c64f7} \\
\texorpdfstring{\softcode{backend/src/main/java/mx/bovicierre/acceso/UserController.java}}{backend/src/main/java/mx/bovicierre/acceso/UserController.java}
&
\texorpdfstring{\softcode{c12a7a0cf0baaf2c1aa9ee0ded3bf45979109c8d2e6e8d08b163a84aacfd6896}}{c12a7a0cf0baaf2c1aa9ee0ded3bf45979109c8d2e6e8d08b163a84aacfd6896} \\
\texorpdfstring{\softcode{backend/src/main/java/mx/bovicierre/evidencia/AnimalQrController.java}}{backend/src/main/java/mx/bovicierre/evidencia/AnimalQrController.java}
&
\texorpdfstring{\softcode{0b6a6e3e58b074f1f5a824fc9230322c00d11750e97f35889fb54c0cf2f9ff47}}{0b6a6e3e58b074f1f5a824fc9230322c00d11750e97f35889fb54c0cf2f9ff47} \\
\texorpdfstring{\softcode{backend/src/main/java/mx/bovicierre/evidencia/CleanupMain.java}}{backend/src/main/java/mx/bovicierre/evidencia/CleanupMain.java}
&
\texorpdfstring{\softcode{0397be2c5e180c65bebe8a040abd58b03619909d4d3f0068f7928bca62268772}}{0397be2c5e180c65bebe8a040abd58b03619909d4d3f0068f7928bca62268772} \\
\texorpdfstring{\softcode{backend/src/main/java/mx/bovicierre/evidencia/DocumentController.java}}{backend/src/main/java/mx/bovicierre/evidencia/DocumentController.java}
&
\texorpdfstring{\softcode{bb3ea4e99dd3cbd3ee1fb3cb0dc568d26f6929a27390f66e3f9a154cc29ce11d}}{bb3ea4e99dd3cbd3ee1fb3cb0dc568d26f6929a27390f66e3f9a154cc29ce11d} \\
\texorpdfstring{\softcode{backend/src/main/java/mx/bovicierre/evidencia/PrivateDocuments.java}}{backend/src/main/java/mx/bovicierre/evidencia/PrivateDocuments.java}
&
\texorpdfstring{\softcode{1ba2a68c44b82f52befbd895646e6abd9dd30b62475f29455c0229854cbaf75f}}{1ba2a68c44b82f52befbd895646e6abd9dd30b62475f29455c0229854cbaf75f} \\
\texorpdfstring{\softcode{backend/src/main/java/mx/bovicierre/evidencia/SettlementController.java}}{backend/src/main/java/mx/bovicierre/evidencia/SettlementController.java}
&
\texorpdfstring{\softcode{7c052b65f2a025c8939a7a1b4be095ebb21732b58993107425a0af94e41a31f5}}{7c052b65f2a025c8939a7a1b4be095ebb21732b58993107425a0af94e41a31f5} \\
\texorpdfstring{\softcode{backend/src/main/java/mx/bovicierre/evidencia/SettlementPdf.java}}{backend/src/main/java/mx/bovicierre/evidencia/SettlementPdf.java}
&
\texorpdfstring{\softcode{8dbbcb11d87bc00b9e74e50d3d99d8664d3af111be0066cf44463dd7f21ef698}}{8dbbcb11d87bc00b9e74e50d3d99d8664d3af111be0066cf44463dd7f21ef698} \\
\texorpdfstring{\softcode{backend/src/main/java/mx/bovicierre/evidencia/application/DocumentStorage.java}}{backend/src/main/java/mx/bovicierre/evidencia/application/DocumentStorage.java}
&
\texorpdfstring{\softcode{5ac197356bde15fc6ce659d6440e2a039a651c5cd9b3b4c09ba5448b1630a810}}{5ac197356bde15fc6ce659d6440e2a039a651c5cd9b3b4c09ba5448b1630a810} \\
\texorpdfstring{\softcode{backend/src/main/java/mx/bovicierre/integraciones/DiagnosticsController.java}}{backend/src/main/java/mx/bovicierre/integraciones/DiagnosticsController.java}
&
\texorpdfstring{\softcode{50f1358ceb808e0c8e4ef871196a5b4c9441052840efb734b09278ca08679b87}}{50f1358ceb808e0c8e4ef871196a5b4c9441052840efb734b09278ca08679b87} \\
\texorpdfstring{\softcode{backend/src/main/java/mx/bovicierre/integraciones/OutboxWorker.java}}{backend/src/main/java/mx/bovicierre/integraciones/OutboxWorker.java}
&
\texorpdfstring{\softcode{342557eef2b4ee80ae5a8a4abe93e649aa905a7bf5fb5071a8df4bf0a144b5a6}}{342557eef2b4ee80ae5a8a4abe93e649aa905a7bf5fb5071a8df4bf0a144b5a6} \\
\texorpdfstring{\softcode{backend/src/main/java/mx/bovicierre/integraciones/ProviderController.java}}{backend/src/main/java/mx/bovicierre/integraciones/ProviderController.java}
&
\texorpdfstring{\softcode{178136f71e4f0195f8d8a8fd84007f0bb9d306e7b33ec729b04dd91791aa39a1}}{178136f71e4f0195f8d8a8fd84007f0bb9d306e7b33ec729b04dd91791aa39a1} \\
\texorpdfstring{\softcode{backend/src/main/java/mx/bovicierre/operaciones/CoreController.java}}{backend/src/main/java/mx/bovicierre/operaciones/CoreController.java}
&
\texorpdfstring{\softcode{b1dfb33002e08028c46057fcfe8cc2ef8103c6a0fe72c13f0a528f347ab5c261}}{b1dfb33002e08028c46057fcfe8cc2ef8103c6a0fe72c13f0a528f347ab5c261} \\
\texorpdfstring{\softcode{backend/src/main/java/mx/bovicierre/reportes/CsvController.java}}{backend/src/main/java/mx/bovicierre/reportes/CsvController.java}
&
\texorpdfstring{\softcode{3e47b279a3e96034c186d92ac8cb9943acbca898cce440a2662a0c5e49646195}}{3e47b279a3e96034c186d92ac8cb9943acbca898cce440a2662a0c5e49646195} \\
\texorpdfstring{\softcode{backend/src/main/java/mx/bovicierre/shared/ApiErrors.java}}{backend/src/main/java/mx/bovicierre/shared/ApiErrors.java}
&
\texorpdfstring{\softcode{412a052933aa19717d87c29d5339e4d40a8633e1998f211528aa5d2352f1a77b}}{412a052933aa19717d87c29d5339e4d40a8633e1998f211528aa5d2352f1a77b} \\
\texorpdfstring{\softcode{backend/src/main/java/mx/bovicierre/shared/ApiException.java}}{backend/src/main/java/mx/bovicierre/shared/ApiException.java}
&
\texorpdfstring{\softcode{9ec9372a5a0cbd0b76540fa371038a71afaba3705f3dae75e691867f2b53dfc4}}{9ec9372a5a0cbd0b76540fa371038a71afaba3705f3dae75e691867f2b53dfc4} \\
\texorpdfstring{\softcode{backend/src/main/java/mx/bovicierre/shared/CommandService.java}}{backend/src/main/java/mx/bovicierre/shared/CommandService.java}
&
\texorpdfstring{\softcode{1155245ac57caf95b6c85e014711ab91b3168c7b550ed9526f27b65fe3a2b171}}{1155245ac57caf95b6c85e014711ab91b3168c7b550ed9526f27b65fe3a2b171} \\
\texorpdfstring{\softcode{backend/src/main/java/mx/bovicierre/shared/DatabaseConfiguration.java}}{backend/src/main/java/mx/bovicierre/shared/DatabaseConfiguration.java}
&
\texorpdfstring{\softcode{a3836302471740162a1274873080ad1e93b2c36db20ac14b00065898fdbc174f}}{a3836302471740162a1274873080ad1e93b2c36db20ac14b00065898fdbc174f} \\
\texorpdfstring{\softcode{backend/src/main/java/mx/bovicierre/shared/JsonSupport.java}}{backend/src/main/java/mx/bovicierre/shared/JsonSupport.java}
&
\texorpdfstring{\softcode{2080fcd9843efadb72229c4fccc29ae5e8e0766f0b7b14ae87d130812f7ade20}}{2080fcd9843efadb72229c4fccc29ae5e8e0766f0b7b14ae87d130812f7ade20} \\
\texorpdfstring{\softcode{backend/src/main/java/mx/bovicierre/shared/NodeHealth.java}}{backend/src/main/java/mx/bovicierre/shared/NodeHealth.java}
&
\texorpdfstring{\softcode{84ffffa377961c6113844d45666b8ca66ad2ed870c987fdfe523dc2843f7f48d}}{84ffffa377961c6113844d45666b8ca66ad2ed870c987fdfe523dc2843f7f48d} \\
\texorpdfstring{\softcode{backend/src/main/resources/application-production.yml}}{backend/src/main/resources/application-production.yml}
&
\texorpdfstring{\softcode{bfcc733378cf611c3ac96f728c53e7931fbddcfb0b8a8da91eb02271f9d92124}}{bfcc733378cf611c3ac96f728c53e7931fbddcfb0b8a8da91eb02271f9d92124} \\
\texorpdfstring{\softcode{backend/src/main/resources/application.yml}}{backend/src/main/resources/application.yml}
&
\texorpdfstring{\softcode{4c4d625e70f97278baff22b1d560dae926aa14ce34488380ac23feb5c1da2987}}{4c4d625e70f97278baff22b1d560dae926aa14ce34488380ac23feb5c1da2987} \\
\texorpdfstring{\softcode{db/bootstrap/bootstrap.py}}{db/bootstrap/bootstrap.py}
&
\texorpdfstring{\softcode{c848592d82e067a1eb6462219ec194b4487870bbffb15576aea7914199a63fed}}{c848592d82e067a1eb6462219ec194b4487870bbffb15576aea7914199a63fed} \\
\texorpdfstring{\softcode{db/migrations/V10\_\_filtro\_posiciones\_vendibles.sql}}{db/migrations/V10\_\_filtro\_posiciones\_vendibles.sql}
&
\texorpdfstring{\softcode{fa88f9c089a4cc52ae031aa99b55999315be2151031cef23f202531bf3b2b1ab}}{fa88f9c089a4cc52ae031aa99b55999315be2151031cef23f202531bf3b2b1ab} \\
\texorpdfstring{\softcode{db/migrations/V11\_\_cancelacion\_anulacion\_version.sql}}{db/migrations/V11\_\_cancelacion\_anulacion\_version.sql}
&
\texorpdfstring{\softcode{85cb7f45ac83d6df429fcaa5f8362294dcceb632b94378cb662a1dd536aeedff}}{85cb7f45ac83d6df429fcaa5f8362294dcceb632b94378cb662a1dd536aeedff} \\
\texorpdfstring{\softcode{db/migrations/V12\_\_pesajes\_reportes\_canonicos.sql}}{db/migrations/V12\_\_pesajes\_reportes\_canonicos.sql}
&
\texorpdfstring{\softcode{ba0af087d801fafaf088f7d38012412cd4dd1b1a2cb513728b624341890ccdf9}}{ba0af087d801fafaf088f7d38012412cd4dd1b1a2cb513728b624341890ccdf9} \\
\texorpdfstring{\softcode{db/migrations/V13\_\_correccion\_composicion\_exacta.sql}}{db/migrations/V13\_\_correccion\_composicion\_exacta.sql}
&
\texorpdfstring{\softcode{1c58fc42da30d1cdea67b339f8ece408764fa7f0103eda20d9792eb2b30d743c}}{1c58fc42da30d1cdea67b339f8ece408764fa7f0103eda20d9792eb2b30d743c} \\
\texorpdfstring{\softcode{db/migrations/V1\_\_esquema.sql}}{db/migrations/V1\_\_esquema.sql}
&
\texorpdfstring{\softcode{8412272f53a2aaeb0453c4acfd8af9d36a6fe359ee9fb663da05a674cf637e5e}}{8412272f53a2aaeb0453c4acfd8af9d36a6fe359ee9fb663da05a674cf637e5e} \\
\texorpdfstring{\softcode{db/migrations/V2\_\_comandos.sql}}{db/migrations/V2\_\_comandos.sql}
&
\texorpdfstring{\softcode{7022f6430a72603d6a1781407705631ad884ef7430f8b6f2db961461d597412c}}{7022f6430a72603d6a1781407705631ad884ef7430f8b6f2db961461d597412c} \\
\texorpdfstring{\softcode{db/migrations/V3\_\_consultas\_sesiones\_grants.sql}}{db/migrations/V3\_\_consultas\_sesiones\_grants.sql}
&
\texorpdfstring{\softcode{e26bc752bccd1dea5a22ba7776ef3d5869ab85856f46dac0412192c52a97ae3b}}{e26bc752bccd1dea5a22ba7776ef3d5869ab85856f46dac0412192c52a97ae3b} \\
\texorpdfstring{\softcode{db/migrations/V4\_\_capturas\_proveedor\_dto.sql}}{db/migrations/V4\_\_capturas\_proveedor\_dto.sql}
&
\texorpdfstring{\softcode{32f075db76d99245c8bb48e373c4e300ced46ad473bd3f1fad2d76c1a3d3f378}}{32f075db76d99245c8bb48e373c4e300ced46ad473bd3f1fad2d76c1a3d3f378} \\
\texorpdfstring{\softcode{db/migrations/V5\_\_corte\_por\_periodo\_reconocimiento.sql}}{db/migrations/V5\_\_corte\_por\_periodo\_reconocimiento.sql}
&
\texorpdfstring{\softcode{7e2c049246ba3ecb0f2e1c1b5cfec621d79cddbd86d6d4763e884f71bfdfcbfc}}{7e2c049246ba3ecb0f2e1c1b5cfec621d79cddbd86d6d4763e884f71bfdfcbfc} \\
\texorpdfstring{\softcode{db/migrations/V6\_\_tarifa\_money\_paginacion.sql}}{db/migrations/V6\_\_tarifa\_money\_paginacion.sql}
&
\texorpdfstring{\softcode{9fe7e7e27c008a0ffeb585ea32a317344f60656105bc496cc23f0ab619329759}}{9fe7e7e27c008a0ffeb585ea32a317344f60656105bc496cc23f0ab619329759} \\
\texorpdfstring{\softcode{db/migrations/V7\_\_filtros\_consultas.sql}}{db/migrations/V7\_\_filtros\_consultas.sql}
&
\texorpdfstring{\softcode{7f6f3226a72d23ea7964783f717267524a3bf998037fe7b625d457730b502479}}{7f6f3226a72d23ea7964783f717267524a3bf998037fe7b625d457730b502479} \\
\texorpdfstring{\softcode{db/migrations/V8\_\_borrador\_kg\_y\_expiracion.sql}}{db/migrations/V8\_\_borrador\_kg\_y\_expiracion.sql}
&
\texorpdfstring{\softcode{cb9d7b36eead9672bcf60c3b9a341442a959e0fdea7be1ed856443cee90f5367}}{cb9d7b36eead9672bcf60c3b9a341442a959e0fdea7be1ed856443cee90f5367} \\
\texorpdfstring{\softcode{db/migrations/V9\_\_lector\_a\_captura\_owner\_fecha.sql}}{db/migrations/V9\_\_lector\_a\_captura\_owner\_fecha.sql}
&
\texorpdfstring{\softcode{a7dc89b0656e14dc62c28f6653dad4954005b7750cd6b37954483adea208bb1d}}{a7dc89b0656e14dc62c28f6653dad4954005b7750cd6b37954483adea208bb1d} \\
\texorpdfstring{\softcode{db/seed/demo.py}}{db/seed/demo.py} &
\texorpdfstring{\softcode{d1c5aac0d7f173e67e2ae856565d45f13669ff13f6d54477ab8b2cc7cdb7d3b9}}{d1c5aac0d7f173e67e2ae856565d45f13669ff13f6d54477ab8b2cc7cdb7d3b9} \\
\texorpdfstring{\softcode{db/tests/benchmark.py}}{db/tests/benchmark.py}
&
\texorpdfstring{\softcode{567196ae28024ca38a8516260cfcd306b73b1b42a8cb92aeba2bcec6e741d02b}}{567196ae28024ca38a8516260cfcd306b73b1b42a8cb92aeba2bcec6e741d02b} \\
\texorpdfstring{\softcode{docs/api/openapi.yaml}}{docs/api/openapi.yaml}
&
\texorpdfstring{\softcode{d635fb00d126833d3bd9a6f243ab2ec21a59488f304ce3bac1a9cdc412fddad7}}{d635fb00d126833d3bd9a6f243ab2ec21a59488f304ce3bac1a9cdc412fddad7} \\
\texorpdfstring{\softcode{docs/architecture/adr/ADR-001-monolito-modular.md}}{docs/architecture/adr/ADR-001-monolito-modular.md}
&
\texorpdfstring{\softcode{4d72eeec9291a6073ae1a631fb88b42bed004218039e6e40033ecc93ecfdc6c3}}{4d72eeec9291a6073ae1a631fb88b42bed004218039e6e40033ecc93ecfdc6c3} \\
\texorpdfstring{\softcode{docs/architecture/adr/ADR-002-sql-explicito-calculo-canonico.md}}{docs/architecture/adr/ADR-002-sql-explicito-calculo-canonico.md}
&
\texorpdfstring{\softcode{550f9d2167b1e62adc92a46a12b0a36d43498dc65cfadc26cfcf1d97b054ecf1}}{550f9d2167b1e62adc92a46a12b0a36d43498dc65cfadc26cfcf1d97b054ecf1} \\
\texorpdfstring{\softcode{docs/architecture/adr/ADR-003-animal-y-posicion.md}}{docs/architecture/adr/ADR-003-animal-y-posicion.md}
&
\texorpdfstring{\softcode{2152cf4798f1142471f3fcb7bf059563ae2af341e6b1593844ada8f0fb96f202}}{2152cf4798f1142471f3fcb7bf059563ae2af341e6b1593844ada8f0fb96f202} \\
\texorpdfstring{\softcode{docs/architecture/adr/ADR-004-reconocimiento-y-costo.md}}{docs/architecture/adr/ADR-004-reconocimiento-y-costo.md}
&
\texorpdfstring{\softcode{fc7d95d6edb5423e9ebf71f54e0fdc74dca379c23fcd405f5e088a7de0a0326d}}{fc7d95d6edb5423e9ebf71f54e0fdc74dca379c23fcd405f5e088a7de0a0326d} \\
\texorpdfstring{\softcode{docs/architecture/adr/ADR-005-maestro-y-propiedad.md}}{docs/architecture/adr/ADR-005-maestro-y-propiedad.md}
&
\texorpdfstring{\softcode{ea5ba9327d9ec434d552644dea391509837bfd9cf303dcc97119e569e7dd5d71}}{ea5ba9327d9ec434d552644dea391509837bfd9cf303dcc97119e569e7dd5d71} \\
\texorpdfstring{\softcode{docs/architecture/adr/ADR-006-replicacion-bidireccional.md}}{docs/architecture/adr/ADR-006-replicacion-bidireccional.md}
&
\texorpdfstring{\softcode{6a80728dd848a2a01838c1f013f3101e317725234733da840d5f80297d1fbab1}}{6a80728dd848a2a01838c1f013f3101e317725234733da840d5f80297d1fbab1} \\
\texorpdfstring{\softcode{docs/architecture/adr/ADR-007-pgpool-raw.md}}{docs/architecture/adr/ADR-007-pgpool-raw.md}
&
\texorpdfstring{\softcode{6728367b5a398bd9b200532d4d978d3d92c6c1e44451d7f37d9f5fe72b25714c}}{6728367b5a398bd9b200532d4d978d3d92c6c1e44451d7f37d9f5fe72b25714c} \\
\texorpdfstring{\softcode{docs/architecture/adr/ADR-008-autenticacion-y-aislamiento.md}}{docs/architecture/adr/ADR-008-autenticacion-y-aislamiento.md}
&
\texorpdfstring{\softcode{20804ffaba0c97de9d7492a53c911f1e186f40f1a2d9c66f7ab4dbae519cb370}}{20804ffaba0c97de9d7492a53c911f1e186f40f1a2d9c66f7ab4dbae519cb370} \\
\texorpdfstring{\softcode{docs/architecture/adr/ADR-009-documentos-y-revision.md}}{docs/architecture/adr/ADR-009-documentos-y-revision.md}
&
\texorpdfstring{\softcode{b06908e8e11d08c1da88a957171e152311120784b89e5fadc0ad3c8b0294ceb7}}{b06908e8e11d08c1da88a957171e152311120784b89e5fadc0ad3c8b0294ceb7} \\
\texorpdfstring{\softcode{docs/architecture/adr/ADR-010-despliegue-respaldos.md}}{docs/architecture/adr/ADR-010-despliegue-respaldos.md}
&
\texorpdfstring{\softcode{afc5ca3b528d63864d3631e45515dcd71b08d638616acc92269f1710f7fe51b3}}{afc5ca3b528d63864d3631e45515dcd71b08d638616acc92269f1710f7fe51b3} \\
\texorpdfstring{\softcode{docs/architecture/adr/ADR-011-graphify.md}}{docs/architecture/adr/ADR-011-graphify.md}
&
\texorpdfstring{\softcode{540ff40c28a3301be41a4319ea3d6273166f63aa6c6f6023a305523984073545}}{540ff40c28a3301be41a4319ea3d6273166f63aa6c6f6023a305523984073545} \\
\texorpdfstring{\softcode{docs/architecture/adr/ADR-012-alcance-aplazado.md}}{docs/architecture/adr/ADR-012-alcance-aplazado.md}
&
\texorpdfstring{\softcode{3ffbb4312018355ef5e26ecc336067461fa70417da59be4bd67a2f8a8663ef55}}{3ffbb4312018355ef5e26ecc336067461fa70417da59be4bd67a2f8a8663ef55} \\
\texorpdfstring{\softcode{docs/architecture/adr/ADR-013-sesiones-disponibilidad.md}}{docs/architecture/adr/ADR-013-sesiones-disponibilidad.md}
&
\texorpdfstring{\softcode{ca5ea5e8ad10f271cf797fe6c2eefb9ac51249314a19e8c4094d817c38148a2f}}{ca5ea5e8ad10f271cf797fe6c2eefb9ac51249314a19e8c4094d817c38148a2f} \\
\texorpdfstring{\softcode{docs/architecture/adr/ADR-014-muerte-entre-duenos.md}}{docs/architecture/adr/ADR-014-muerte-entre-duenos.md}
&
\texorpdfstring{\softcode{018c9c1055690b7974ed0bd410f4e1258311316ce8659d75717e40868e861f41}}{018c9c1055690b7974ed0bd410f4e1258311316ce8659d75717e40868e861f41} \\
\texorpdfstring{\softcode{docs/architecture/adr/ADR-015-periodos-cartera.md}}{docs/architecture/adr/ADR-015-periodos-cartera.md}
&
\texorpdfstring{\softcode{5f804df10834dad175b406d4561fdd47bb34c48398a19a0244ffe17ba8f6bbb1}}{5f804df10834dad175b406d4561fdd47bb34c48398a19a0244ffe17ba8f6bbb1} \\
\texorpdfstring{\softcode{docs/architecture/adr/ADR-016-autoridad-esquema.md}}{docs/architecture/adr/ADR-016-autoridad-esquema.md}
&
\texorpdfstring{\softcode{d7e7585da3ad5dd1748d9196aea81876e35231a2ebbe29752c2afb596424d9c9}}{d7e7585da3ad5dd1748d9196aea81876e35231a2ebbe29752c2afb596424d9c9} \\
\texorpdfstring{\softcode{docs/architecture/adr/ADR-017-despliegue-local.md}}{docs/architecture/adr/ADR-017-despliegue-local.md}
&
\texorpdfstring{\softcode{a14fb7d1a048369cec2bed892417fdde5d817838ca3b6f0edf97f59a1b5fcede}}{a14fb7d1a048369cec2bed892417fdde5d817838ca3b6f0edf97f59a1b5fcede} \\
\texorpdfstring{\softcode{docs/database/data-dictionary.md}}{docs/database/data-dictionary.md}
&
\texorpdfstring{\softcode{4062189a8d17e4c04837507e045d317fc1ea6f4e18fcbb0ffd57bec2a6406202}}{4062189a8d17e4c04837507e045d317fc1ea6f4e18fcbb0ffd57bec2a6406202} \\
\texorpdfstring{\softcode{docs/database/normalization.md}}{docs/database/normalization.md}
&
\texorpdfstring{\softcode{41e9b66e398822eb5196e69e75f2e768c6b3d91f9500a3ccd52e42370b763ee6}}{41e9b66e398822eb5196e69e75f2e768c6b3d91f9500a3ccd52e42370b763ee6} \\
\texorpdfstring{\softcode{docs/domain/business-rules.md}}{docs/domain/business-rules.md}
&
\texorpdfstring{\softcode{bcda2f84af489ec68fd5c7712dffbf3c79433b7dc3a61dd87c68255dd3f2796a}}{bcda2f84af489ec68fd5c7712dffbf3c79433b7dc3a61dd87c68255dd3f2796a} \\
\texorpdfstring{\softcode{docs/domain/glossary.md}}{docs/domain/glossary.md}
&
\texorpdfstring{\softcode{32af005f9ba92301617325c1ffc68b2970c9f92e9bdd3e3d3921c47fc93bf7a0}}{32af005f9ba92301617325c1ffc68b2970c9f92e9bdd3e3d3921c47fc93bf7a0} \\
\texorpdfstring{\softcode{docs/domain/use-cases.md}}{docs/domain/use-cases.md}
&
\texorpdfstring{\softcode{43f4a1a0f3f92dcc80cf480794a5c6565d04b34fb98926960f914f85f0b76f7f}}{43f4a1a0f3f92dcc80cf480794a5c6565d04b34fb98926960f914f85f0b76f7f} \\
\texorpdfstring{\softcode{frontend/src/App.tsx}}{frontend/src/App.tsx} &
\texorpdfstring{\softcode{f0325f17a2e5bc541be695a59c1afd7211f19be1de3738c7b550ba1b57225367}}{f0325f17a2e5bc541be695a59c1afd7211f19be1de3738c7b550ba1b57225367} \\
\texorpdfstring{\softcode{frontend/src/api.ts}}{frontend/src/api.ts} &
\texorpdfstring{\softcode{33a8ed401e986157f05586bcb769360dd5da57046bc4eddbfe9e3ccb8234d787}}{33a8ed401e986157f05586bcb769360dd5da57046bc4eddbfe9e3ccb8234d787} \\
\texorpdfstring{\softcode{frontend/src/domain.ts}}{frontend/src/domain.ts}
&
\texorpdfstring{\softcode{7410af2f7a5db646f020ab5bd93f45d520b62eb1fde48f0a020e7ce26781b5d6}}{7410af2f7a5db646f020ab5bd93f45d520b62eb1fde48f0a020e7ce26781b5d6} \\
\texorpdfstring{\softcode{frontend/src/generated/api-schema.ts}}{frontend/src/generated/api-schema.ts}
&
\texorpdfstring{\softcode{24f01aa65cecf591cb8adcf83352d6ae859da28c1e465227b8305c83a933f97b}}{24f01aa65cecf591cb8adcf83352d6ae859da28c1e465227b8305c83a933f97b} \\
\texorpdfstring{\softcode{frontend/src/main.tsx}}{frontend/src/main.tsx}
&
\texorpdfstring{\softcode{e34b50223ac942759bc00c81a3d886fa77fc39645c31a32bd14841cbaceb2300}}{e34b50223ac942759bc00c81a3d886fa77fc39645c31a32bd14841cbaceb2300} \\
\texorpdfstring{\softcode{frontend/src/screens.tsx}}{frontend/src/screens.tsx}
&
\texorpdfstring{\softcode{db38246a91b46b34c692cbbb256537a8c9e33306a8ef611bb5c9a1f927d92eae}}{db38246a91b46b34c692cbbb256537a8c9e33306a8ef611bb5c9a1f927d92eae} \\
\texorpdfstring{\softcode{frontend/src/styles.css}}{frontend/src/styles.css}
&
\texorpdfstring{\softcode{d59390ef19345dedab9c2f9c517f728e47d26c3f0bd045863d2ca9e67231dd1a}}{d59390ef19345dedab9c2f9c517f728e47d26c3f0bd045863d2ca9e67231dd1a} \\
\texorpdfstring{\softcode{frontend/src/ui.tsx}}{frontend/src/ui.tsx} &
\texorpdfstring{\softcode{18619ae8fcc27bffd235ca8734be12ba92d22acdd44b1538a7823d311e1a1a0a}}{18619ae8fcc27bffd235ca8734be12ba92d22acdd44b1538a7823d311e1a1a0a} \\
\texorpdfstring{\softcode{frontend/src/workflow.ts}}{frontend/src/workflow.ts}
&
\texorpdfstring{\softcode{20e030ed1ddf17df4604576c053adef3517fbfe14bc82f9210f6c614205403b7}}{20e030ed1ddf17df4604576c053adef3517fbfe14bc82f9210f6c614205403b7} \\
\texorpdfstring{\softcode{scripts/academic.py}}{scripts/academic.py} &
\texorpdfstring{\softcode{c88e9852db350f5476e070d08f88602168606ca69c7e39fc490e577c58839b4e}}{c88e9852db350f5476e070d08f88602168606ca69c7e39fc490e577c58839b4e} \\
\texorpdfstring{\softcode{scripts/dev.py}}{scripts/dev.py} &
\texorpdfstring{\softcode{3747b730bb044601a8be5f14e6253e7fe3a91d98288ea5366de02b9ed699020b}}{3747b730bb044601a8be5f14e6253e7fe3a91d98288ea5366de02b9ed699020b} \\
\texorpdfstring{\softcode{scripts/doctor.py}}{scripts/doctor.py} &
\texorpdfstring{\softcode{68bc5b701db50d83b4ac04dcf3f38cddfb6151f57bade4f92f3ac5ed52414f97}}{68bc5b701db50d83b4ac04dcf3f38cddfb6151f57bade4f92f3ac5ed52414f97} \\
\texorpdfstring{\softcode{scripts/recovery.py}}{scripts/recovery.py} &
\texorpdfstring{\softcode{f965704aea82f2fb262054b769f8a5d42cfd1b74bdcc575c12ec856356d00900}}{f965704aea82f2fb262054b769f8a5d42cfd1b74bdcc575c12ec856356d00900} \\
\texorpdfstring{\softcode{scripts/sync-task-tables.py}}{scripts/sync-task-tables.py}
&
\texorpdfstring{\softcode{efac9ff97c57233fb41b39fa70b544ff9af5cd59e403b574cfd55504a8ea29b0}}{efac9ff97c57233fb41b39fa70b544ff9af5cd59e403b574cfd55504a8ea29b0} \\
\texorpdfstring{\softcode{tools/dev/check\_api\_contract.py}}{tools/dev/check\_api\_contract.py}
&
\texorpdfstring{\softcode{58a7df485501aaf23cf8ceaeba610addbea2c5602799d5f443f31e34d68053ef}}{58a7df485501aaf23cf8ceaeba610addbea2c5602799d5f443f31e34d68053ef} \\
\texorpdfstring{\softcode{tools/esquema/comparar\_catalogo.py}}{tools/esquema/comparar\_catalogo.py}
&
\texorpdfstring{\softcode{2ce2511a2ef240a8e3bb8797799eb159f93ef85ba963b399fbd7e410d31e2698}}{2ce2511a2ef240a8e3bb8797799eb159f93ef85ba963b399fbd7e410d31e2698} \\
\end{longtable}
}

\end{document}
